\documentclass[11pt]{article}
\usepackage[a4paper,margin=24mm]{geometry}
\usepackage{fontspec}
\usepackage[AutoFallBack=true]{xeCJK}
\setmainfont{Noto Serif}
\setCJKmainfont{Noto Serif CJK SC}
\setCJKfallbackfamilyfont{rm}{SimSun-ExtB}
\xeCJKDeclareCharClass{CJK}{"25A1, "25CB}
\usepackage{xcolor}
\usepackage[hidelinks]{hyperref}
\setlength{\parindent}{0pt}
\setlength{\parskip}{0.4em}
\newcommand{\zhsmall}{\footnotesize\color{black!70}}
\newcommand{\unitlabel}[1]{\marginpar{\raggedleft\tiny\color{gray}#1}}
\newenvironment{unitzh}{\par\CJKfontspec{Noto Serif CJK SC}[AutoFallBack=true]\large\linespread{1.25}\selectfont}{\par}
\newenvironment{uniten}{\par\small\color{black!85}}{\par\medskip}
\sloppy
\begin{document}
\begin{titlepage}\centering
{\fontsize{34}{40}\selectfont 數書九章}\par\vspace{1em}
{\LARGE Mathematical Treatise in Nine Sections}\par\vspace{0.6em}
{\large by Qin Jiushao (秦九韶), 1247}\par\vspace{2em}
{\normalsize Chinese text with modern punctuation, and a complete English translation}\par\vfill
\begin{minipage}{0.86\textwidth}\small
Chinese text: Siku quanshu (四庫全書) edition, titled 數學九章 there, as transcribed on Chinese Wikisource (the pages and revisions used are listed at the end).
Punctuation, division into units and English translation: AI-produced by Claude Opus 5.5 (Anthropic), 24--25 October 2026.
The punctuated Chinese was checked by program against the Wikisource text: every character is reproduced, in order,
and small print (commentary) stays small print. In the translation, commentary is given in square brackets with its
author. The labels in the margin identify the units of the machine-readable file units.jsonl.
\end{minipage}
\end{titlepage}
\tableofcontents\clearpage
\section*{數書九章}\addcontentsline{toc}{section}{數書九章}

\begin{unitzh}\unitlabel{shushu-001.1}欽定四庫全書 子部六\end{unitzh}
\begin{uniten}Imperially Commissioned Complete Library of the Four Treasuries (Qinding Siku quanshu) — Masters Branch (zibu), 6\end{uniten}

\begin{unitzh}\unitlabel{shushu-001.2}數學九章 天文算法類二{\zhsmall 算書之屬}提要\end{unitzh}
\begin{uniten}Shuxue jiuzhang (Mathematics in Nine Sections) — Astronomy and Mathematics Category (tianwen suanfa lei), 2 [Subcategory of mathematical books] — Abstract (tiyao)\end{uniten}

\begin{unitzh}\unitlabel{shushu-001.3}{\zhsmall 臣}等謹案：《數學九章》十八卷，宋秦九韶撰。九韶始末未詳，惟據原序自稱其籍曰魯郡；然序題淳祐七年，魯郡已久入於元，九韶葢述其祖貫，未詳實為何許人也。是書分為九類：一曰大衍，以奇零求總數，為九類之綱；二曰天時，以步氣朔、晷影及五星伏見；三曰田域，以推方圓羃積；四曰測望，以推髙深廣逺；五曰賦役，以均租稅力役；六曰錢榖，以權輕重出入；七曰營建，以度土功；八曰軍旅，以定行陣；九曰市易，以治交易。雖以「九章」為名，而與古《九章》門目迥别，葢古法設其術，九韶則别其用耳。\end{unitzh}
\begin{uniten}[Your servant] and others respectfully note: the Shuxue jiuzhang, in 18 juan (卷), was composed by Qin Jiushao of the Song. The particulars of Jiushao's life are not known; only according to his original preface does he give his native place as Lu commandery (Lujun). But the preface is dated the 7th year of Chunyou, when Lu commandery had long since passed into the Yuan; Jiushao was presumably stating his ancestral home, and it is not known where he actually came from. This book is divided into 9 categories. The first is called Great Extension (dayan); it seeks the total number from the odd remainders (qiling) and is the guiding principle (gang) of the 9 categories. The second is called Heavenly Seasons (tianshi), for computing the solar terms (qi) and new moons (shuo), the gnomon shadows, and the disappearances and reappearances (fujian) of the five planets. The third is called Fields and Domains (tianyu), for deducing the areas (miji 羃積) of squares and circles. The fourth is called Surveying (cewang), for deducing heights, depths, breadths and distances. The fifth is called Taxes and Corvée (fuyi), for equalizing land rents, taxes and labor service. The sixth is called Money and Grain (qiangu), for weighing light and heavy, outgoing and incoming. The seventh is called Construction (yingjian), for measuring earthworks. The eighth is called Military Affairs (junlü), for determining marching formations. The ninth is called Markets and Trade (shiyi), for managing exchange. Although it takes jiuzhang ('nine chapters') as its name, its divisions and headings are entirely different from those of the ancient Nine Chapters; for the ancient methods set out their techniques, whereas Jiushao merely distinguishes them by their uses.\end{uniten}

\begin{unitzh}\unitlabel{shushu-001.4}宋代諸儒尚虚談而薄實用，數雖聖門六藝之一，亦鄙之不言；即有談數學者，亦不過推演河洛之奇偶，於人事無闗。故樂屢争而不決，厯亦每變而愈舛，豈非算術不明、惟憑臆斷之故歟？數百年中，惟沈括究心是事，而自《夢溪筆談》以外，未有成書。九韶當宋末造，獨崛起而明絶學。其中如大衍類「蓍卦發微」，欲以新術改《周易》揲蓍之法，殊乖古義；「古厯㑹積」，題數既誤，且謂設問以明大衍之理，初不計前後多少之厯過，尤非實據；天時類「綴術推星」，本非方程法，而術曰方程，復於草中多設一數以合方程行列，更為牽合；所載皆平氣平朔，凡晷影長短、五星遲疾，皆設數加減，不過得其大㮣，較今之定氣定朔用三角形推算者，亦為未密。\end{unitzh}
\begin{uniten}The Confucian scholars of the Song esteemed empty talk and slighted practical use; although number is one of the six arts of the Sage's school, they too despised it and did not speak of it; and even those who did discuss the learning of number did no more than elaborate the odd and even [numbers] of the River and the Luo, which have no bearing on human affairs. Hence music was repeatedly disputed without being settled, and the calendar, with every change, grew more erroneous: was this not because the art of calculation (suanshu) was not understood and they relied only on conjecture? In several hundred years only Shen Kuo devoted his mind to this matter, and apart from the Mengxi bitan (Brush Talks from Dream Brook) there was no finished book. Jiushao, in the last days of the Song, alone rose up and brought the extinct learning to light. Among its contents, for example: in the Great Extension category, 'Revealing the Subtleties of Milfoil and Hexagrams' (Shigua fawei) seeks to change the Zhou Changes' (Zhouyi) method of counting off the milfoil stalks by a new technique, which greatly departs from the ancient meaning; in 'Accumulated Total of the Ancient Calendars' (Guli huiji) the numbers of the problem are wrong, and moreover he says that the question was set to illuminate the principle of the Great Extension, without at all reckoning how much had elapsed (厯過) before and after, which is especially without real basis; in the Heavenly Seasons category, 'Computing the Planets by the Zhui Technique' (Zhuishu tuixing) is fundamentally not the rectangular arrays (fangcheng) method, yet the Method says 'rectangular arrays', and moreover in the Working an extra number is set up to fit the rows and columns of the rectangular arrays, which is even more forced. What it records are all mean solar terms (pingqi) and mean new moons (pingshuo); for the lengths of the gnomon shadows and the slowness and speed of the five planets, assumed numbers are in every case added and subtracted, giving no more than a rough approximation; compared with today's true solar terms (dingqi) and true new moons (dingshuo), computed with triangles, it is also not precise.\end{uniten}

\begin{unitzh}\unitlabel{shushu-001.5}然自秦漢以來，成法相傳，未有言其立法之意者；惟此書大衍術中所載立天元一法，能舉立法之意而言之。其用雖僅一端，而以零數推總數，足以盡奇偶和較之變，至為精妙。苟得其意而用之，凡諸法所不能得者，皆隨所用而無不通。後元郭守敬用之於弧矢，李冶用之於勾股方圓；歐邏巴新法易其名曰借根方，用之於九章、八線。其源實開自九韶，亦可云有功於算術者矣。至於田域、測望、賦役、錢榖、營建、軍旅、市易七類，皆擴充古法，取事命題，雖條目紛紜，曲折往復，不免瑕瑜互見，而其精確者居多。今即《永樂大典》所載，於其誤者正之，疎者辨之，顛倒者次序之，各加案語於下，庶得失不掩，俾算家有所稽考焉。乾隆四十九年十月恭校上。\end{unitzh}
\begin{uniten}Yet since the Qin and Han, established methods have been handed down, but no one has stated the intent behind setting up the methods; only the method of setting up the celestial unknown (tianyuan) recorded in the Great Extension technique of this book is able to bring out the intent behind setting up the method and state it. Although its use is only one aspect, yet, deducing the total number from the remainders (lingshu), it suffices to exhaust the variations of odd and even and of sums and differences (hejiao), and is supremely refined and subtle. If one grasps its intent and uses it, then whatever the various methods cannot obtain is in every case, according to the use made of it, penetrated without exception. Later, in the Yuan, Guo Shoujing applied it to arcs and sagittae (hushi), and Li Ye applied it to base and altitude (gougu) and to squares and circles; the new methods of Europe changed its name to 'borrowed root and powers' (jiegenfang) and applied it to the Nine Chapters and the eight lines (baxian). Its source was in fact opened up by Jiushao, and he too may be said to have rendered meritorious service to the art of calculation. As for the 7 categories Fields and Domains, Surveying, Taxes and Corvée, Money and Grain, Construction, Military Affairs and Markets and Trade, they all expand the ancient methods, taking up affairs to frame the problems; although the items are numerous and confused, winding and roundabout, so that flaws and fine jade inevitably appear side by side, yet the precise and exact ones are in the majority. Now, on the basis of what is contained in the Yongle dadian, we have corrected what is erroneous in it, discussed what is careless and put in order what is misarranged, adding notes (anyu) below in each case, so that its merits and faults may not be concealed and mathematicians may have something to examine. Respectfully collated and presented in the 10th month of the 49th year of Qianlong.\end{uniten}

\begin{unitzh}\unitlabel{shushu-001.6}總纂官：{\zhsmall 臣}紀昀、{\zhsmall 臣}陸錫熊、{\zhsmall 臣}孫士毅。總校官：{\zhsmall 臣}陸費墀。\end{unitzh}
\begin{uniten}Chief Compilers (zongzuanguan): [your servant] Ji Yun, [your servant] Lu Xixiong, [your servant] Sun Shiyi. Chief Collator (zongjiaoguan): [your servant] Lufei Chi.\end{uniten}

\begin{unitzh}\unitlabel{shushu-001.7}《數學九章》序\end{unitzh}
\begin{uniten}Preface to the Shuxue jiuzhang\end{uniten}

\begin{unitzh}\unitlabel{shushu-001.8}周教六藝，數實成之，學士大夫所從事尚矣。其用本太虛生一，而周流無窮；大可以通神明、順性命，小足以經世務、類萬物，距容以淺近窺哉？若昔推䇿以迎日，定律而和氣，髀距濬川，土圭度咎，天地之大，囿焉而不能外，况其間總總者乎？爰自河圖、洛書闓發幽祕，八卦、九疇錯綜精㣲，極而至於大衍、皇極之用，而人事之變無不該，鬼神之情莫能隱矣。聖人神之，言而遺其粗；常人昧之，由而莫之覺。要其歸，則數與道非二本也。漢去古未逺，有張蒼、許商、馬延年、耿夀昌、鄭𤣥、張衡、劉洪之倫，或明天道而法傳於後，或計功䇿而效驗於時。後世學者自高，鄙不之講，此學殆絕。惟治厯疇人能為乘除，而弗通於開方衍變。若官府會事，則府史一二絫之；算家位置，素所不識；上之人亦委而聽焉。持算者惟若人，則鄙之也宜矣。嗚呼！樂有制氏，僅記鏗鏘，而謂與天地同和者止於是，可乎？\end{unitzh}
\begin{uniten}Of the six arts (liuyi) taught by the Zhou, number (shu) truly completes them, and it has long been what scholars and grandees occupied themselves with. Its use is rooted in the Great Void giving birth to the One, and it flows round everywhere without end; in great things it can penetrate the spirits and accord with nature and destiny, in small things it suffices to order the affairs of the world and to classify the myriad things: how could it permit being peered at in a shallow and near-sighted way? As of old, when they reckoned with the stalks (ce) to welcome the sun, fixed the pitch-pipes (律) and so harmonized the qi, [used] the gnomon and distances (髀距) to dredge the rivers, and the earth gauge (tugui) to measure the 'fault' (咎) [Translator's note: 咎 'fault' gives no sense here; the text appears corrupt.]: the vastness of heaven and earth is enclosed by it and cannot lie outside it—how much more the multitudinous things between them! Ever since the River Chart (Hetu) and the Luo Writing (Luoshu) opened up the hidden secrets, and the Eight Trigrams (bagua) and the Nine Divisions (jiuchou) interwove the subtle and minute, reaching their utmost in the uses of the Great Extension and the August Ultimate (huangji), no change in human affairs has been left unencompassed, and the dispositions of ghosts and spirits cannot be hidden. The sages held it divine, and spoke of it while leaving aside its coarser parts; ordinary people are in the dark about it, and follow it without being aware of it. To sum up where it leads: number and the Way are not two roots. The Han was not far removed from antiquity, and there were men such as Zhang Cang, Xu Shang, Ma Yannian, Geng Shouchang, Zheng Xuan, Zhang Heng and Liu Hong: some understood the Way of Heaven, and their methods were handed down to later ages; others calculated works and stratagems (gongce), and their efficacy was proven in their own time. Scholars of later ages held themselves high, despised it and would not study it, so that this learning was nearly extinct. Only the hereditary experts (chouren) who manage the calendar could perform multiplication and division, but they did not master root extraction (kaifang) and its extended variations. As for the reckoning of affairs in government offices, the clerks (fushi) pile [the numbers] up one by one; the placements (weizhi) of the mathematicians are something they have never known; and those above simply leave it to them and accept [their results]. If those who hold the calculations are only such men, then it is fitting that [the art] is despised. Alas! In music there was the Zhi clan, who merely recorded the clanging and jingling; but to say that what shares in the harmony of heaven and earth stops at this—is that acceptable?\end{uniten}

\begin{unitzh}\unitlabel{shushu-001.9}今數術之書尚三十餘家：天象厯度謂之綴術，太乙、壬、甲謂之三式，皆曰内算，言其祕也；《九章》所載及《周官》九數擊於方圓者為專術，皆曰外算，對内而言也。其用相通，不可岐二。獨大衍法不載《九章》，未有能推之者；厯家演法頗用之，以為方程者，誤也。且天下之事多矣。古之人先事而計，計定而行；仰觀俯察，人謀鬼謀，無所不用其謹，是以不愆於成，載籍章章可覆也。後世興事造始，鮮能考度，𣹰𣹰乎天紀人事之殽缺矣，可不求其故哉？\end{unitzh}
\begin{uniten}Today there are still more than 30 schools (jia) of books on the numerical arts (shushu). Celestial phenomena and calendrical degrees are called the zhui technique (zhuishu); Taiyi, Ren and Jia are called the Three Patterns (sanshi); all are called 'inner calculation' (neisuan), meaning that they are secret. What is recorded in the Nine Chapters, and the Nine Calculations (jiushu) of the Zhou Offices (Zhouguan), which strike (擊) at squares and circles, are specialized techniques (zhuanshu), all called 'outer calculation' (waisuan), so named in contrast to the inner. Their uses are interconnected and cannot be split in two. Only the Great Extension method is not recorded in the Nine Chapters, and no one has been able to work it out; calendar-makers, in developing their methods, make a good deal of use of it, and those who take it to be rectangular arrays (fangcheng) are mistaken. Moreover, the affairs of the world are many. The men of antiquity calculated before an undertaking, and acted once the calculation was settled; looking up to observe and looking down to examine, taking counsel with men and taking counsel with spirits, there was nothing in which they did not exercise their care; therefore they did not fail in their accomplishments, and the records clearly bear this out and can be checked. In later ages, in initiating undertakings and beginning works, few have been able to examine and measure: overwhelming (𣹰𣹰) is the confusion and deficiency of the celestial order and human affairs! Can one fail to seek the reason for it?\end{uniten}

\begin{unitzh}\unitlabel{shushu-001.10}九韶愚陋，不閑於藝。然早嵗侍親中都，因得訪習於太史；又嘗從隱君子受數學。時際兵難，歴嵗遥塞，不自意全於矢石之間；更險離憂，荏苒十禩，心槁氣落，信知失物莫不有數也。乃肆意其間，旁諏方能，探索杳𦕈，粗若有得焉。所謂通神明、順性命，固膚末於見；若其小者，竊嘗設為問荅，以擬於用。積多而惜其棄，因取八十一題，釐為九類，立術具草，間以圖發之。恐或可備博學多識君子之餘觀，曲藝可遂也，願進之於道。儻曰藝成而下，是惟疇人府史流也，烏足盡天下之用，亦無瞢焉。時淳祐七年九月，魯郡秦九韶叙。\end{unitzh}
\begin{uniten}I, Jiushao, am foolish and crude and unskilled in the arts. But in my early years I attended on my parents in the capital (zhongdu), and so was able to inquire and study with the Grand Astrologer (taishi); I also once received mathematics from a recluse gentleman. Then came the calamities of war; for years I was on the distant frontier, and did not expect to come through whole among the arrows and stones. Passing through dangers and suffering griefs, 10 years slipped away; my heart withered and my spirit sank, and truly I came to know that no lost thing is without its number (shu). So I gave free rein to my mind in it, consulted widely with men of skill, searched into the obscure and the faint, and roughly seemed to gain something. As for what is called 'penetrating the spirits and according with nature and destiny', my insight is surely only skin-deep; as for its small [matters], I have privately ventured to compose questions and answers, modelled on practical use. When they had accumulated to many, I was loath to discard them; so I took 81 problems, arranged them into 9 categories, set up the Methods and furnished the Workings, and here and there brought them out with diagrams. Perhaps they may serve for the leisure perusal of erudite and well-informed gentlemen; if a minor art may thus be carried through, I wish to advance it to the Way. Should someone say, 'When an art is perfected it ranks low; this is merely the business of hereditary experts and clerks; how could it suffice to fulfil the uses of the world?', then I too feel no shame about it. Prefaced in the 9th month of the 7th year of Chunyou by Qin Jiushao of Lu commandery.\end{uniten}

\begin{unitzh}\unitlabel{shushu-001.11}且系之曰：昆崙旁薄，道本虚一。聖有大衍，㣲寓於《易》。竒餘取䇿，郡數皆捐。衍而䆒之，探隱之原。數術之傳，以實為體。其書《九章》，惟兹弗紀。厯家雖用，用而不知。小試經世，姑推所為。述大衍第一。\end{unitzh}
\begin{uniten}And [he] appended to it, saying: Undifferentiated and all-pervading (昆崙旁薄), the Way is at root void and one. The sages had the Great Extension, subtly lodged in the Changes. The odd remainders are taken as the stalks (ce); the 'commandery' (郡) numbers are all discarded. [Translator's note: 郡 'commandery' gives no sense here; the text appears corrupt.] Extend it and investigate it, and one probes the source of the hidden. In the transmission of the numerical arts, the concrete is taken as the substance. Their book is the Nine Chapters; this alone it does not record. Though calendar-makers use it, they use it without understanding. As a small trial in ordering the world, I tentatively work out what it does. Setting forth Great Extension, [Section] 1.\end{uniten}

\begin{unitzh}\unitlabel{shushu-001.12}七精迴穹，人事之紀。追綴而求，宵星畫晷。厯乆則踈，惟智能革。不尋天道，模襲何益？三農務穯，厥施自天。以滋以生，雨膏雪零。司牧閔馬，尺寸驗之。積以噐移，憂喜皆非。述天時第二。\end{unitzh}
\begin{uniten}The seven luminaries (qijing) revolve round the vault; they are the regulators of human affairs. Pursuing and piecing together (追綴) to find them: the stars by night, the traced (畫) shadows. When a calendar has run long it grows inexact; only the wise can reform it. If one does not search out the Way of Heaven, what benefit is there in imitating and copying? The three [classes of] farmers devote themselves to sowing; what is bestowed comes from Heaven. Moistening and engendering, the rain enriches and the snow falls. The shepherds [of the people] are anxious about the horses (閔馬), and test it by chi (尺) and cun (寸). [Translator's note: 閔馬 'anxious about the horses' gives no clear sense; the text appears corrupt.] As the accumulation shifts with the vessel, worry and joy are both mistaken. Setting forth Heavenly Seasons, [Section] 2.\end{uniten}

\begin{unitzh}\unitlabel{shushu-001.13}魁隗粒民，甄度四海。蒼姫井之，仁政攸在。代逺庶蕃，懇菑日廣。步度庀賦，版圖是掌。方圓異狀，袤窳殊形。專術精㣲，熟䆒厥真。差之毫釐，謬乃千里。公私共𡚁，蓋謹其籍。述田域第三。\end{unitzh}
\begin{uniten}Kuiwei fed the people with grain and examined and measured the four seas. The Azure Ji (Cangji) laid them out in wells (jing): therein lay benevolent government. [Translator's note: Kuiwei 魁隗 is an epithet of Shennong, the Divine Farmer; the Azure Ji are the Zhou royal house, surnamed Ji.] As the ages grew remote and the multitudes multiplied, the diligently opened new fields (懇菑) grew wider by the day. Pacing and measuring to arrange the levies, they kept charge of the registers and maps. Square and round differ in shape; long and sunken (袤窳) differ in form. The specialized techniques are subtle; one must thoroughly investigate their truth. Err by a hao (毫) or a li (釐), and the mistake becomes 1000 li (里). Public and private alike suffer the harm; hence care must be taken with the registers. Setting forth Fields and Domains, [Section] 3.\end{uniten}

\begin{unitzh}\unitlabel{shushu-001.14}莫高匪山，莫濬匪川。神禹奠之，積矩攸傳。智創巧述，重差夕桀。求之既詳，揆之岡越。崇深廣逺，度則靡容。形格勢禁，㓂壘仇墉。欲知其數，先望以表。因差施術，坐悉㣲𣺌。述測望第四。\end{unitzh}
\begin{uniten}Nothing is high but the mountains; nothing is deep but the rivers. The divine Yu settled them, and the [method of] accumulated squares (jiju) was thereby handed down. The wise created and the skilful transmitted: Double Differences (chongcha) and xijie (夕桀). Having sought them out in detail, one gauges them across the ridges (岡越). Height and depth, breadth and distance: measuring them is not permitted. The terrain obstructs and circumstances forbid: bandits' ramparts and enemies' walls. Wishing to know their numbers, first sight them with gnomons (biao). Applying the techniques according to the differences, one knows the minute and fine while sitting still. Setting forth Surveying, [Section] 4.\end{uniten}

\begin{unitzh}\unitlabel{shushu-001.15}邦國之賦，以待百事。畦田經入，取之有度。未免力役，先商厥功。以衰以率，逸乃同。漢猶近古，租稅以算。調均錢榖，何菑之扞？惟仁隱民，猶已溺饑。賦役不均，寜得勿思？述賦役第五。\end{unitzh}
\begin{uniten}The levies of the state serve to provide for the hundred affairs. From the plots of fields comes the regular revenue; it must be taken with measure. Labor service being unavoidable, first deliberate on its work. By graded shares (cui) and by rates (lü), ease is thus made equal. [Translator's note: this verse line has only three characters; one appears to be missing.] The Han was still near antiquity; rents and taxes were levied by reckoning (suan). Adjust and equalize money and grain, and what calamity is there to ward off? Only the humane feel pity for the people, as though they were already drowning or starving. When levies and labor service are unequal, how can one not take thought? Setting forth Taxes and Corvée, [Section] 5.\end{uniten}

\begin{unitzh}\unitlabel{shushu-001.16}物等歛賦，式時府庾。粒粟寸絲，褐夫紅女。商征邊糴，後世多端。立縁為欺，上下俱殫。我聞理財，如智治水。澄原濬流，惟其深矣。彼昧弗察，𢡖急煩刑。去理益逺，吁嗟不仁！述錢榖第六。\end{unitzh}
\begin{uniten}Goods are graded and levies collected; by regulation and season [they go to] the treasuries and granaries. Every grain of millet and every cun of silk [comes from] the men in coarse cloth and the women at their needlework. Merchant duties and frontier grain purchases take many forms in later ages; pretexts are set up for deceit, and those above and below are alike exhausted. I have heard that managing wealth is like the wise man's controlling of water: clear the source and dredge the stream—only so does it run deep. Those who are benighted and do not examine [this] are harsh and hasty, with vexatious punishments; departing ever farther from principle—alas, how inhumane! Setting forth Money and Grain, [Section] 6.\end{uniten}

\begin{unitzh}\unitlabel{shushu-001.17}斯城斯池，乃棟乃宇。宅坐寄命，以保以聚。鳩功雉制，竹箇本章。匪䆒匪度，財蠧力傷。圍蔡而裁，如子西素。匠計露臺，俾漢文懼。惟武圖功，惟儉昭德。有國有家，兹焉取則。述營建第七。\end{unitzh}
\begin{uniten}These walls and these moats, these ridgepoles and these eaves: in them one dwells and sits (宅坐) and entrusts one's life, for protection and for gathering together. Assembling the labor and regulating by the zhi (雉); bamboo by the piece (箇) and trunks by the zhang (章): without investigation and without measurement, wealth is eaten away and strength injured. When Cai was besieged, [the works] were measured out as Zixi had planned beforehand; the craftsmen's estimate for the open terrace made [Emperor] Wen of Han apprehensive. Only for the martial does one plan works; only frugality makes virtue shine. Those who have a state or a house take their rule from this. Setting forth Construction, [Section] 7.\end{uniten}

\begin{unitzh}\unitlabel{shushu-001.18}天生五材，兵去未可。不教而戰，維上之過。堂堂之陣，鵝鸛為行。營應規矩，其將莫當。師中之吉，惟智仁勇。夜算軍書，先計攸重。我聞在昔，輕則寡謀。殄民以幸，亦孔之憂。述軍旅第八。\end{unitzh}
\begin{uniten}Heaven produced the five materials; to do away with weapons is not yet possible. To fight without training [the people] is the fault of those above. Imposing are the battle arrays, with goose and crane [formations] for ranks; when the camp conforms to compass and square, its general is irresistible. The good fortune of being in the midst of the army lies only in wisdom, humaneness and courage. Calculating the military documents by night, it is the prior reckoning that weighs most. I have heard that in former times the rash had little strategy; to destroy the people for the sake of luck is a great cause for sorrow. Setting forth Military Affairs, [Section] 8.\end{uniten}

\begin{unitzh}\unitlabel{shushu-001.19}日中而市，萬民所資。賈貿墆鬻，利析銖錙。蹛財役貧，封君低首。豕末兼并，非國之厚。述市易第九。\end{unitzh}
\begin{uniten}Holding the market at midday is what the myriad people depend upon. Merchants trade, hoard and sell, splitting their profits down to the zhu (銖) and zi (錙). Those who hoard wealth make the poor serve them, and enfeoffed lords bow their heads. Swine (豕), the branch occupations and engrossment are not what makes a state prosperous. [Translator's note: 豕 'swine' gives no sense here; the text appears corrupt.] Setting forth Markets and Trade, [Section] 9.\end{uniten}

\begin{unitzh}\unitlabel{shushu-001.20}全覽。卷一上、卷一下、卷二上、卷二下、卷三上、卷三下、卷四上、卷四下、卷五上、卷五下、卷六上、卷六下、卷七上、卷七下、巻八上、卷八下、卷九上、卷九下。\end{unitzh}
\begin{uniten}Full view (quanlan). Juan 1, upper; Juan 1, lower; Juan 2, upper; Juan 2, lower; Juan 3, upper; Juan 3, lower; Juan 4, upper; Juan 4, lower; Juan 5, upper; Juan 5, lower; Juan 6, upper; Juan 6, lower; Juan 7, upper; Juan 7, lower; Juan 8, upper; Juan 8, lower; Juan 9, upper; Juan 9, lower.\end{uniten}

\section*{卷1上}\addcontentsline{toc}{section}{卷1上}

\begin{unitzh}\unitlabel{shushu-002.1}欽定四庫全書\end{unitzh}
\begin{uniten}Imperially Commissioned Complete Library of the Four Treasuries (Qinding Siku quanshu)\end{uniten}

\begin{unitzh}\unitlabel{shushu-002.2}數學九章卷一上 宋 秦九韶 撰\end{unitzh}
\begin{uniten}Mathematics in Nine Sections (Shuxue jiuzhang), juan 1, upper part. Composed by Qin Jiushao of the Song.\end{uniten}

\begin{unitzh}\unitlabel{shushu-002.3}大衍\end{unitzh}
\begin{uniten}Great Extension (dayan)\end{uniten}

\begin{unitzh}\unitlabel{shushu-002.4}按大衍術以各分數之竒零求各分數之總數，大而天行，小而物數，皆可御之。其法有求元、求定、求術、求竒、求乘、求用之目。大約以數之竒偶為根，而以諸數相度之盡不盡為用：有求彼此不能度盡之諸數者，元數、定數是也；有求諸數皆能度盡之一數者，衍母數是也；有求諸數皆能度盡而一數不能度盡之數者，各衍數是也，其不盡之數即竒數也；有求二數相度餘一之數者，乘數是也；有求二數相度餘一而諸數又能度盡之數者，用數是也。求元數、定數，初與約分法相似，終變二數，務使其等數為一。蓋以一為等數，始能度盡二數，是他數俱不能度盡二數，而二數相度益不能盡也。以定數竒數求乘數之法，名曰大衍求一，中有「立天元一於左上」之語，下載立天元一算式。\end{unitzh}
\begin{uniten}Note: The Great Extension method uses the odd remainders (jiling) of the various dividing numbers (fenshu) to find the total number (zongshu) of the various dividing numbers; great—the course of the heavens—and small—the numbers of things—all can be governed by it. Its method has the headings: finding the original (yuan), finding the definite (ding), finding the method (shu) [Translator's note: 術 here is probably an error for 衍 "extension"], finding the remainder (ji), finding the multiplier (cheng), finding the use (yong). Broadly, it takes the oddness and evenness of numbers as its root, and whether or not numbers measure (du) one another exhaustively as its use. There is the finding of numbers that cannot measure one another exhaustively: the original numbers (yuanshu) and definite numbers (dingshu) are these. There is the finding of the one number that all the numbers can measure exhaustively: the extension mother (yanmu) number is this. There is the finding of numbers that all the numbers can measure exhaustively but one number cannot measure exhaustively: the respective extension numbers (yanshu) are these, and the number not exhausted is the remainder number (jishu). There is the finding of a number such that two numbers measuring each other leave 1: the multiplier number (chengshu) is this. There is the finding of a number such that two numbers measuring each other leave 1 and that the numbers can moreover measure exhaustively: the use number (yongshu) is this. Finding the original numbers and definite numbers is at first similar to the method of reducing fractions (yuefen), but in the end it alters the two numbers, striving to make their common measure (dengshu) 1. For only when 1 is the common measure can [it] measure the two numbers exhaustively; this means that no other number can measure the two numbers exhaustively, and still less can the two numbers measuring each other be exhausted. The method of finding the multiplier number from the definite number and the remainder number is called Great Extension seeking one (dayan qiuyi); in it there is the phrase "set up the celestial unknown at the upper left", and below is recorded a calculation layout (suanshi) of setting up the celestial unknown.\end{uniten}

\begin{unitzh}\unitlabel{shushu-002.5}按立天元一法，見於元郭守敬之《厯源》、李冶之《測圓海鏡》，及四海之借根方者，皆虛設所求之數為一，與所有實數反覆推求，歸於少廣諸乘方，得其積數與邊數或正負亷隅數而止，次用除法或開方法得所求數。此數命定數為一，與竒數反覆商較，至餘一實數而止，其竒數所積即為乘數。蓋其用不同，而法則無二也。然其極和較之用，窮竒偶之情，則有為元法、西法所未及者。但原本法解煩雜，圖式譌舛，今詳加改定，並釋其義，俾學者易見焉。\end{unitzh}
\begin{uniten}Note: The method of setting up the celestial unknown (tianyuan), as seen in the Liyuan (厯源) of Guo Shoujing of the Yuan and the Ceyuan haijing (測圓海鏡) of Li Ye, and the borrowed root and powers (jiegenfang) of the Four Seas [Translator's note: 四海 perhaps stands for 西洋 "the West"; cf. 西法 below]—all of these hypothetically set the number sought as 1, work it out back and forth with all the actual numbers (shishu) given, reduce it to the various powers (chengfang) of lesser breadth (shaoguang), and stop when they have obtained its area (ji) number and side number, or the positive and negative edge (lian) and corner (yu) numbers; next they use division or root extraction to obtain the number sought. This reckoning designates the definite number as 1, estimates the quotient (shang) and compares back and forth with the remainder number, and stops when 1 remains as the dividend (shi); what its remainder number accumulates is then the multiplier number. For their uses differ, but in method they are not two. Yet in carrying the use of sums and differences (hejiao) to the utmost and exhausting the properties of odd and even, it has [points] that the Yuan method and the Western method have not reached. But in the original the explanations of the methods are cumbersome and confused, and the diagrams and layouts are erroneous; now we have carefully revised and fixed them and explained their meaning as well, so that students may see it easily.\end{uniten}

\begin{unitzh}\unitlabel{shushu-002.6}大衍數術\end{unitzh}
\begin{uniten}The Great Extension Number Method\end{uniten}

\begin{unitzh}\unitlabel{shushu-002.7}大衍總數術曰：置諸問數{\zhsmall 類名有四}：一曰元數{\zhsmall 謂尾位見單零者，本門揲蓍、酒息、斛糶、砌甎、失米之類是也}；二曰牧數{\zhsmall 謂尾位見分釐者，假令冬至三百六十五日二十五刻，欲與甲子六十日為一會，而求積日之類}；三曰通數{\zhsmall 謂諸數各有分子母者，本門問一會積年是也}；四曰復數{\zhsmall 謂尾位見十或百及千以上者，本門築隄并急足之類是也}。\end{unitzh}
\begin{uniten}The Great Extension total-number method (dayan zongshu shu) says: Set out the numbers asked (wenshu) [Note: The category names are four]. First, original numbers [Note: That is, those whose last place shows single units (danling); in this section, counting off the yarrow stalks, wine interest, selling grain by the hu (斛), laying bricks, lost rice and the like are these]. Second, mu (牧) numbers [Translator's note: 牧 is evidently an error for 收 shou, "collected"; cf. 收數 below] [Note: That is, those whose last place shows fen (分) and li (釐); for instance, the winter solstice being 365 days 25 ke (刻), one wishes to make one meeting (hui) with the jiazi of 60 days and seeks the accumulated days, and the like]. Third, converted numbers (tongshu) [Note: That is, numbers that each have numerator and denominator; in this section, the question of the accumulated years of one meeting is this]. Fourth, repeated numbers (fushu) [Note: That is, those whose last place shows tens, or hundreds, or thousands and above; in this section, building dikes, express couriers and the like are these].\end{uniten}

\begin{unitzh}\unitlabel{shushu-002.8}按此言問題有是四類。\end{unitzh}
\begin{uniten}Note: This says that problems are of these four kinds.\end{uniten}

\begin{unitzh}\unitlabel{shushu-002.9}元數者，先以兩連環求等，約竒弗約偶{\zhsmall 或約得五而彼有十，乃約偶而弗約竒}。或元數俱偶，約畢可存一位見偶；或皆約而猶有類數存，姑置之，俟與其他約徧而後，乃與姑置者求等約之；或諸數皆不可盡類，則以諸元數命曰復{\zhsmall 按復應作定}數，以復{\zhsmall 按同上}數格入之。\end{unitzh}
\begin{uniten}As for original numbers: first, two by two in linked rings (lianhuan), seek the common measure (deng); reduce (yue) the odd and do not reduce the even [Note: Or if reducing gives 5 and the other has 10, then reduce the even and do not reduce the odd]. Or, if the original numbers are all even, when the reducing is finished one place may be kept showing even. Or, if all have been reduced and there still remain numbers of a kind, set them aside for the time being; wait until the reducing with the others is complete, and afterwards seek the common measure with those set aside and reduce them. Or, if none of the numbers can be exhausted as of a kind, then name the original numbers "repeated" [Siku editors: 復 "repeated" should be 定 "definite"] numbers, and enter them by the repeated [Siku editors: same as above] number frame (ge).\end{uniten}

\begin{unitzh}\unitlabel{shushu-002.10}按此以元數求定數法也。連環求等者，於諸數中逐次取二數相度，得一數可以度盡者為等數，為法只約一數。約竒弗約偶，專為等數為偶者言之；若等數為竒者，則約偶弗約竒；而等數為五與十者，又有或約竒或約偶者矣。皆約而猶有類，俟約徧求等約之者，逐條兩兩取約畢，猶有二數可約者，求得等數為續等{\zhsmall 見後}，續等約一數，必復乘一數。蓋等數為二數之較{\zhsmall 或二數展轉之較}，可約續等乃已約之較，不可𠕂約；不可約，而兩數又不可使有等，故約一數乘一數，猶之不約也。術内未詳。\end{unitzh}
\begin{uniten}Note: This is the method of finding the definite numbers from the original numbers. "Seeking the common measure in linked rings" means: among the numbers, take two numbers at a time in succession and let them measure each other; the one number obtained that can measure them exhaustively is the common measure; taken as divisor (fa), it reduces only one number. "Reduce the odd and do not reduce the even" is said specifically of the case where the common measure is even; if the common measure is odd, then reduce the even and do not reduce the odd; and where the common measure is 5 or 10, there are again cases of reducing the odd or reducing the even. As for "all having been reduced, there are still [numbers] of a kind; wait until the reducing is complete, seek the common measure and reduce them": when, after taking them item by item, two by two, the reducing is finished and there are still two numbers that can be reduced, the common measure found is the continued common measure (xudeng) [Note: see below]; when the continued common measure reduces one number, one must multiply back (fucheng) one number. For the common measure is the difference of the two numbers [Note: or the difference of the two numbers turned back and forth (zhanzhuan)]; the continued common measure that can reduce is a difference [of numbers] already reduced and cannot be reduced again; it cannot be reduced, and yet the two numbers must not be allowed to have a common measure; therefore reducing one number and multiplying one number is as though not reducing. In the method this is not made clear.\end{uniten}

\begin{unitzh}\unitlabel{shushu-002.11}收數者，乃命尾位分釐作單零，以進所問之數，定位訖，用元數格入之；或如意立數為母，收進分厘以從所問，用通數格入之。\end{unitzh}
\begin{uniten}As for collected numbers (shoushu): designate the fen and li of the last place as single units, so as to advance the numbers asked; when the places have been fixed, enter them by the original number frame. Or set up a number at will (ruyi) as denominator (mu), collect and advance the fen and li so as to follow what is asked, and enter them by the converted number frame.\end{uniten}

\begin{unitzh}\unitlabel{shushu-002.12}按收數者，單位下有竒零之數也。進位者，以竒零之末位為單位也。若立分母通之，反不如用原數為簡。\end{unitzh}
\begin{uniten}Note: Collected numbers are numbers that have odd remainders below the units place. "Advancing the places" means taking the last place of the odd remainder as the units place. If one sets up a denominator and converts them, it is, on the contrary, not as simple as using the original numbers.\end{uniten}

\begin{unitzh}\unitlabel{shushu-002.13}通數者，置問數，通分内子，互乘之，皆曰通數。求總等，不約一位，約衆位，得各原法數，用元數格入之。或諸母數繁，就分從省通之者，皆不用元，各母仍求總等，存一位，約衆位，亦各得原法數，亦用元法數格入之。\end{unitzh}
\begin{uniten}As for converted numbers: set out the numbers asked, convert to parts (tongfen) and put in the numerators (neizi), and cross-multiply (hucheng) them; all are called converted numbers. Seek the general common measure (zongdeng); do not reduce one place, reduce the many places, obtaining the respective original divisor numbers (yuanfa shu); enter them by the original number frame. Or, where the denominators are cumbersome and one converts them following the parts in a simplified way, in all cases one does not use the original; for the respective denominators still seek the general common measure, keep one place and reduce the many places; one likewise obtains the respective original divisor numbers, and likewise enters them by the original divisor number frame.\end{uniten}

\begin{unitzh}\unitlabel{shushu-002.14}按通數與收數相似，但單數有分母，竒零為分子耳。通分納子，即進尾數為單位之義，因加互乘一次，故加總等一約，然後為元數也。\end{unitzh}
\begin{uniten}Note: Converted numbers are similar to collected numbers, except that the units have a denominator and the odd remainder is the numerator. "Converting to parts and putting in the numerator" has the same sense as advancing the last digit to the units place; because one cross-multiplication is added, one reduction by the general common measure is added, and after that they are original numbers.\end{uniten}

\begin{unitzh}\unitlabel{shushu-002.15}又按求總等，不拘通數、復數，但題中有三數可以一等數度盡者，即可用總等法，存一數，約衆數，然後為元數。凡度之後等數仍可約者，此數必當存之。\end{unitzh}
\begin{uniten}Further note: Seeking the general common measure is not restricted to converted numbers and repeated numbers; whenever three numbers in a problem can be measured exhaustively by one common measure, the general common measure method can be used: keep one number and reduce the many numbers, and after that they are original numbers. Whenever, after measuring, the common measure can still reduce, this number must be kept.\end{uniten}

\begin{unitzh}\unitlabel{shushu-002.16}復數者，問數尾位見十以上者，以諸數求總等，存一位，約衆位，始得元數。兩兩連環求等，約竒弗約偶，復乘偶；或約偶或約竒，復乘竒{\zhsmall 按此四語有誤，應作「約竒弗約偶，復乘偶；或約偶弗約竒，弗乘竒」，然皆續等下用之，此處可省}。或彼此可約而猶有類數存者，又相减以求續等，以續等約彼，則必復乘此，乃得定數。所有元數、收數、通數三格，皆有復乘求定之理，悉可入之。\end{unitzh}
\begin{uniten}As for repeated numbers: where the last place of the numbers asked shows tens or above, seek the general common measure of the numbers, keep one place and reduce the many places; only then does one obtain the original numbers. Two by two in linked rings seek the common measure; reduce the odd, do not reduce the even, and multiply back the even; or reduce the even or reduce the odd, and multiply back the odd [Siku editors: These four phrases are in error; they should read "reduce the odd, do not reduce the even, multiply back the even; or reduce the even, do not reduce the odd, do not multiply the odd"; but all of these are used under the continued common measure and may be omitted here]. Or if they can be reduced by one another but numbers of a kind still remain, again subtract them from one another to find the continued common measure; if that one is reduced by the continued common measure, then this one must be multiplied back; thus one obtains the definite numbers. All three frames—original numbers, collected numbers and converted numbers—have the principle of multiplying back to find the definite [numbers]; they may all be entered into them.\end{uniten}

\begin{unitzh}\unitlabel{shushu-002.17}按復數者，諸問數皆至十或百或千而止也。右各叚皆云「以某格入之」，此又云「三格悉可入之」，大約古算必有其程式也。\end{unitzh}
\begin{uniten}Note: Repeated numbers are those in which the numbers asked all reach the tens, or hundreds, or thousands, and stop. Each of the preceding passages says "enter them by such-and-such frame"; this one further says "they may all be entered by the three frames". Presumably the ancient computations must have had their standard layouts (chengshi).\end{uniten}

\begin{unitzh}\unitlabel{shushu-002.18}求定數勿使兩位見偶，勿使見一太多，見一多則借用繁，不欲借則任得一。\end{unitzh}
\begin{uniten}In finding the definite numbers, do not let two places show even; do not let 1 appear too often; if 1 appears often, borrowing for use (jieyong) is cumbersome; if one does not wish to borrow, then let 1 be obtained as it may.\end{uniten}

\begin{unitzh}\unitlabel{shushu-002.19}按「勿使兩位見偶」者，蓋衆數連乘中有兩偶數，則所得總數以一偶數除之，必仍得偶數，不能求餘一之乘數也。「勿使見一太多，見一多則借數繁」者，蓋見一多因數本如此，且見一即不必推，乃云「勿使太多」，又云「借數」，皆塗人之耳目也。故曰「不欲借則任得一」。\end{unitzh}
\begin{uniten}Note: "Do not let two places show even": for if among the many numbers multiplied together in succession there are two even numbers, then the total number obtained, divided by one even number, must still give an even number, and one cannot find a multiplier number leaving 1. "Do not let 1 appear too often; if 1 appears often, the borrowed numbers (jieshu) are cumbersome": for if 1 appears often, it is because the numbers are originally thus; moreover, when 1 appears there is no need to work it out; yet it says "do not let it be too often" and further speaks of "borrowed numbers"—all of this muddles people's ears and eyes. Therefore it says, "if one does not wish to borrow, then let 1 be obtained as it may".\end{uniten}

\begin{unitzh}\unitlabel{shushu-002.20}以定相乘為衍母，以各定約衍母，得各衍數{\zhsmall 或列各定數於右行，各立天元一為子於左行，以母互乘子，亦得衍數}。諸衍數各滿定母去之，不滿曰竒，以竒與定用大衍求一入之，以求乘率{\zhsmall 或竒得一，即為乘率}。\end{unitzh}
\begin{uniten}Multiply the definite [numbers] together to make the extension mother; reduce the extension mother by each definite [number] to obtain the respective extension numbers [Note: Or arrange the definite numbers in the right column, and for each set up the celestial unknown as numerator (zi) in the left column; cross-multiply the numerators by the denominators (mu), and one also obtains the extension numbers]. For each of the extension numbers, as often as it fills the definite mother (dingmu), remove it; what does not fill is called the remainder (ji). With the remainder and the definite [number], use Great Extension seeking one and enter them in it, so as to find the multiplying rate (chenglü) [Note: Or if the remainder obtained is 1, it is the multiplying rate].\end{uniten}

\begin{unitzh}\unitlabel{shushu-002.21}按諸定數連乘為衍母，即為諸定數皆能度盡之數，亦為總數最大之限。凡總數在限内者，各定數之差皆不等；若過限外，則各定數之差有與限内相等者，其兩總數之差必為衍母之倍數。各衍母者，即諸數度盡、一數度不盡之數也。竒數者，定數度衍數不盡之數也。定數原為彼此不能度盡之數，衍數為他定數連乘之數，以此一定數度之，必不能盡也。\end{unitzh}
\begin{uniten}Note: The definite numbers multiplied together in succession make the extension mother; this is the number that all the definite numbers can measure exhaustively, and it is also the greatest limit of the total number. For every total number within the limit, the differences (cha) for the respective definite numbers are all unequal; if it passes beyond the limit, then there are differences for the respective definite numbers that are equal to [those] within the limit, and the difference of the two total numbers must be a multiple (beishu) of the extension mother. The respective extension mothers [Translator's note: 衍母 here is evidently for 衍數 "extension numbers"] are the numbers that the numbers measure exhaustively and one number does not measure exhaustively. The remainder numbers are what is not exhausted when the definite number measures the extension number. The definite numbers are originally numbers that cannot measure one another exhaustively; the extension number is the number of the other definite numbers multiplied together in succession; when it is measured by this one definite number, it certainly cannot be exhausted.\end{uniten}

\begin{unitzh}\unitlabel{shushu-002.22}大衍求一數云：置竒右上，定居右下，立天元一於左上。先以右上除右下，所得商數與左上一相生，入左下。然後乃以右行上下，以少除多，遞互除之，所得商數隨即遞互累乘，歸左行上下，須使右上末後竒一而止，乃驗左上所得，以為乘率。或竒數已見單一者，便為乘率{\zhsmall 按此二語重上}。\end{unitzh}
\begin{uniten}The Great Extension seeking one reckoning (shu) says: Place the remainder at the upper right; the definite [number] is placed at the lower right; set up the celestial unknown at the upper left. First divide the lower right by the upper right; the quotient number obtained mutually generates (xiangsheng) with the 1 at the upper left and is entered at the lower left. After that, with the upper and lower [numbers] of the right column, divide the greater by the lesser, dividing them alternately in turn; the quotient numbers obtained are immediately, alternately in turn, cumulatively multiplied and returned to the upper and lower [places] of the left column. One must make the upper right at the very end a remainder of 1, and stop; then examine what is obtained at the upper left and take it as the multiplying rate. Or if the remainder number already shows a single 1, it is the multiplying rate [Siku editors: These two phrases repeat the above].\end{uniten}

\begin{unitzh}\unitlabel{shushu-002.23}按此以定數竒數求乘數也。其法必使以定數度竒數僅餘一數，而竒數之倍數即乘數也。「置竒右上、定右下」者，初次以定為實、竒為法也。「立天元一於左上」者，以一為竒之倍數也。「得商數與左上相生入左下」者，以竒商定得商數，即竒之倍數，以乘天元一而書於下也。隨以竒數與商數相乘，以减定數，為餘實。次以竒為實，减餘為法，置前左下於左上，以法約實得商，乘左上，又併前之左上為左下，隨以法乘商，减實，又為餘實。次又以前餘為實，次餘為法，置前左下於左上，得商數乘左上，又倂前左上為左下，隨以法乘商减實。如此展轉相求，合兩次為一算，至餘實一，乃視左下天元數，即乘數也。若未至兩次餘實一者，仍以一為法，上餘數為實，實二則商一，實三則商二，如上求之，復得餘一，其天元數方為乘數。原文「遞互乘除」之語未詳。\end{unitzh}
\begin{uniten}Note: This is finding the multiplier number from the definite number and the remainder number. Its method must make it so that, when the remainder number is measured with the definite number, only 1 is left over; and the multiple of the remainder number is then the multiplier number. As for "place the remainder at the upper right, the definite at the lower right": the first time, the definite is taken as dividend (shi) and the remainder as divisor (fa). As for "set up the celestial unknown at the upper left": 1 is taken as the multiple of the remainder. As for "the quotient number obtained mutually generates with the upper left and is entered at the lower left": with the remainder estimate the quotient of the definite, obtaining the quotient number, which is the multiple of the remainder; with it multiply the celestial unknown and write [the result] below. Then multiply the remainder number and the quotient number together and subtract from the definite number, giving the remaining dividend (yushi). Next take the remainder as dividend and the remainder after subtraction (jianyu) as divisor; put the former lower left at the upper left; reduce the dividend by the divisor to obtain the quotient, multiply the upper left [by it], and further add the former upper left to make the lower left; then multiply the quotient by the divisor and subtract from the dividend, again giving the remaining dividend. Next, again take the former remainder as dividend and the next remainder as divisor; put the former lower left at the upper left; obtain the quotient number and multiply the upper left [by it], and further add the former upper left to make the lower left; then multiply the quotient by the divisor and subtract from the dividend. Seeking in this way, turning back and forth, combine two times into one computation; when the remaining dividend is 1, look at the celestial-unknown number at the lower left: it is the multiplier number. If the remaining dividend is 1 before two times have been reached, still take 1 as divisor and the upper remaining number as dividend; if the dividend is 2, the quotient is 1; if the dividend is 3, the quotient is 2; seeking it as above, one again obtains a remainder of 1, and only then is the celestial-unknown number the multiplier number. The wording "alternately multiplying and dividing in turn" in the original text is not made clear.\end{uniten}

\begin{unitzh}\unitlabel{shushu-002.24}置各乘率，對乘衍數，得泛用。併泛，課衍母，多一者為正用；或泛多衍母倍數者，驗元數竒偶同者，損其半倍{\zhsmall 或三處同類，以三約衍母，於三處損之}。同衍母者為無用，當驗元數同類者，而正用至多處借之：以元數兩位求等，以等約衍母為借數，以借數損有以益其無，為正用。或數處無者，如意立數為母，約衍母，所得以如意子乘之，均借補之。或欲從省，勿借，任之為空，可也。\end{unitzh}
\begin{uniten}Place the respective multiplying rates and multiply the corresponding extension numbers by them, obtaining the general use [numbers] (fanyong). Add the general [use numbers] and test (ke) them against the extension mother: if they exceed it by 1, they are the correct use [numbers] (zhengyong). Or if the general [use numbers] exceed by a multiple of the extension mother, examine the original numbers; for those alike in odd or even, diminish by half the multiple [Note: Or if three places are of the same kind, reduce the extension mother by 3 and diminish in the three places]. Those equal to the extension mother are "without use" (wuyong); one should examine the original numbers that are of the same kind and borrow from the place where the correct use is greatest: with two places of the original numbers seek the common measure, reduce the extension mother by the common measure to make the borrowed number, and with the borrowed number diminish the one that has, so as to increase the one that has none, giving the correct use [numbers]. Or if several places have none, set up a number at will as denominator and reduce the extension mother [by it]; multiply what is obtained by the at-will numerator, and borrow evenly to make them up. Or if one wishes to follow the simpler way, do not borrow; letting it be empty is acceptable.\end{uniten}

\begin{unitzh}\unitlabel{shushu-002.25}按此求各用數法也。其各乘率乘各衍數得用數者，即一數餘一、諸數度盡之數也。其云「併泛用過衍母倍數，驗元數同類損之」，此語似有誤，當云「驗問數同偶，而用數相併過衍母者損之」，蓋取用皆問數，非元數也。凡偶數减偶仍餘偶，减竒仍餘竒，其數有定；竒數减竒則餘偶，又或餘竒，减偶則餘竒，又或餘偶，其數無定。故惟偶數可驗也。定數一者即無用數，必虛為借數，未免徒滋煩擾。\end{unitzh}
\begin{uniten}Note: This is the method of finding the respective use numbers. The use numbers obtained by multiplying the respective extension numbers by the respective multiplying rates are the numbers that leave 1 for one number and that the numbers measure exhaustively. Where it says "the sum of the general uses exceeds a multiple of the extension mother; examine the original numbers of the same kind and diminish them", these words seem to be in error; it should say "examine the numbers asked that are alike even, and where the use numbers added together exceed the extension mother, diminish them"; for what is taken for use are all the numbers asked, not the original numbers. In general, an even number less an even still leaves even, and less an odd still leaves odd: the numbers are determinate. An odd number less an odd leaves even, or else leaves odd; less an even it leaves odd, or else leaves even: the numbers are indeterminate. Therefore only even numbers can be tested. Where a definite number is 1, it is a number without use; to insist on making a borrowed number for it in vain cannot avoid uselessly adding trouble.\end{uniten}

\begin{unitzh}\unitlabel{shushu-003.1}然後以其餘各乘正用為各總，併總，滿衍母去之，不滿為所求率數。\end{unitzh}
\begin{uniten}Then multiply each proper use number (zhengyong) by its remainder to make the respective totals; add the totals; remove the expansion mother (yanmu) as often as it is filled; what does not fill it is the rate (lü) number sought.\end{uniten}

\begin{unitzh}\unitlabel{shushu-003.2}按：此既得各用數，以題中所問之竒零求總數也。以各餘數乗各用數者，蓋用數為諸數度盡、一數餘一之數，以幾數乘之，必為諸數度盡、一數餘幾數之數也。併各條而以各數度之，必各數仍餘幾數也。餘數悉合，則總數必合矣。然衍母為諸數度盡之數，累加一衍母，衆餘數皆不變，故滿衍母去之，得在衍母内者，其數最小，為第一數。若大於此數者，遞加一衍母數，無不合者。\end{unitzh}
\begin{uniten}Siku editors' note: This is, the use numbers (yongshu) having been obtained, to seek the total number (zongshu) from the odd remnants (qiling) asked about in the problem. As for multiplying each use number by each remainder number: the use number is a number which the various numbers measure (du) exactly and which a single number measures with 1 left over; multiplied by some number, it must become a number which the various numbers measure exactly and which a single number measures with that many left over. When the items are added together and measured by each number, each number must still leave that many over. When the remainders all agree, the total number must agree. Now the expansion mother is a number which the various numbers measure exactly; when an expansion mother is added repeatedly, none of the remainders changes; therefore, removing the expansion mother as often as it is filled, one obtains the one lying within the expansion mother, whose number is the smallest: it is the first number. As for those larger than this number, by successively adding the number of an expansion mother, there are none that do not agree.\end{uniten}

\begin{unitzh}\unitlabel{shushu-003.3}按：右大衍本法也。原書入於「蓍䇿發㣲」題問荅之後，殊失其序，今修冠於卷首。\end{unitzh}
\begin{uniten}Siku editors' note: The above is the basic method of the Great Expansion (dayan). The original book put it after the question and answer of the problem "Revealing the Subtleties of the Yarrow Stalks" (Shice fawei), which greatly loses its order; now, in revising, we have put it at the head of the juan.\end{uniten}

\begin{unitzh}\unitlabel{shushu-003.4}蓍卦發㣲\end{unitzh}
\begin{uniten}Revealing the Subtleties of Yarrow and Hexagrams (Shigua fawei)\end{uniten}

\begin{unitzh}\unitlabel{shushu-003.5}問：《易》曰：「大衍之數五十，其用四十有九。」又曰：「分而為二以象兩，卦一以象三，揲之以四以象四時，三變而成爻，十有八變而成卦。」欲知所衍之術及其數各幾何。\end{unitzh}
\begin{uniten}Question: The Changes says: "The number of the Great Expansion is 50; its use is 49." It also says: "Divide them into 2 to represent the Two; hang 1 [Translator's note: 卦 here stands for 掛, 'hang'.] to represent the Three; count them off (she) by 4s to represent the 4 seasons; with 3 changes a line (yao) is completed; with 18 changes a hexagram (gua) is completed." One wishes to know the method of the expansion and how much each of its numbers is.\end{uniten}

\begin{unitzh}\unitlabel{shushu-003.6}按：揲蓍之法，載於《易傳》、《啓𫎇》，言之甚明。算術以竒偶相生，取名大衍，可也；竟欲以此易古法，則過矣。\end{unitzh}
\begin{uniten}Siku editors' note: The method of counting off the yarrow stalks (shi) is recorded in the Yi zhuan and the Qimeng, which explain it very clearly. That the art of calculation, with odd and even generating each other, takes the name "Great Expansion" is acceptable; but to go so far as to want to replace the ancient method with this is going too far.\end{uniten}

\begin{unitzh}\unitlabel{shushu-003.7}荅曰：衍母十二，衍法三。一元衍數二十四，二元衍數一十二，三元衍數八，四元衍數六，已上四位衍數計五十。一楪用數一十二，二楪用數二十四，三楪用數四，四揲用數九，已上四位用數四十九。\end{unitzh}
\begin{uniten}Answer: expansion mother 12; expansion divisor (yanfa) 3. Original (yuan) 1, expansion number (yanshu) 24; original 2, expansion number 12; original 3, expansion number 8; original 4, expansion number 6; the expansion numbers of the above 4 places total 50. Counting-off 1, use number 12; counting-off 2, use number 24; counting-off 3, use number 4; counting-off 4, use number 9; the use numbers of the above 4 places: 49.\end{uniten}

\begin{unitzh}\unitlabel{shushu-003.8}按：此附會五十、四十九之數，與本衍已牽强不合，觀後可知。\end{unitzh}
\begin{uniten}Siku editors' note: This makes a forced connection with the numbers 50 and 49; with respect to the basic expansion it is already strained and does not agree, as can be seen from what follows.\end{uniten}

\begin{unitzh}\unitlabel{shushu-003.9}水 火 木 金{\zhsmall 始此四數以揲} 隂陽象數圖 老陽 少隂 少陽 老隂{\zhsmall 終此四者為爻}\end{unitzh}
\begin{uniten}Water | Fire | Wood | Metal [Note: begin with these 4 numbers to count off] Diagram of the Image Numbers (xiangshu) of Yin and Yang. Old yang | Young yin | Young yang | Old yin [Note: end with these 4, which make the lines]\end{uniten}

\begin{unitzh}\unitlabel{shushu-003.10}按：此條與數無取義，可刪。\end{unitzh}
\begin{uniten}Siku editors' note: This item has no meaning to be taken with regard to the numbers; it may be deleted.\end{uniten}

\begin{unitzh}\unitlabel{shushu-003.11}本題術曰：置諸元數，兩兩連環求等，約竒弗約偶，徧約畢，乃變元數，皆曰定母，列右行；各立天元一為子，列左行。以諸定母互乘左行之子，各得，名曰衍數。次以各定母滿去衍數，各餘，名曰竒數。以竒數與定母用大衍術求一{\zhsmall 大衍求一術云：以竒於右上，定母於右下，立天元一於左上。先以右行上下兩位，以少除多，所得商數乃遞互乘，歸左行，使右上得一而止，左上為乘率。}，得各乘率，以乘衍數，各得用數。驗次所揲餘幾何，以其餘數乘諸用數，併之，名曰總數。滿衍母去之，不滿為所求數，以為實。《易》以三才為衍法，以法除實，所得為象數。如實有餘，或一或二，皆命作一，同為象數。其象數得一為老陽，得二為少隂，得三為少陽，得四為老隂。得老陽畫重爻，得少隂畫拆爻，得少陽畫單爻，得老隂畫交爻。凡六畫乃成卦。\end{unitzh}
\begin{uniten}Method for this problem: Set out the original numbers (yuanshu); seek the equal number (deng) in pairs, in linked rings (lianhuan); reduce the odd, do not reduce the even; when the reducing has been done throughout, the original numbers are changed, and all are called fixed mothers (dingmu); put them in the right column. For each, set up the celestial unknown (tianyuan) as the son (zi) [Translator's note: in this procedure the tianyuan is simply the unit 1.]; put them in the left column. Cross-multiply (hucheng) the sons of the left column by the fixed mothers; what each obtains is called the expansion number. Next, remove each fixed mother from the expansion number as often as it is filled; what each leaves is called the odd number (qishu). With the odd numbers and the fixed mothers, use the Great Expansion technique to seek one [Note: The Great Expansion technique of seeking one (dayan qiuyi) says: put the odd number at the upper right and the fixed mother at the lower right, and set up the celestial unknown at the upper left. First, with the 2 places, upper and lower, of the right column, divide the larger by the smaller; the quotients (shang) obtained are then successively cross-multiplied and returned to the left column, making the upper right obtain 1, and then stopping; the upper left is the multiplying rate (chenglü).], obtaining each multiplying rate; multiply the expansion numbers by them, and each obtains the use number. Examine how much remains from each counting-off in turn; multiply the use numbers by those remainders and add them; this is called the total number. Remove the expansion mother as often as it is filled; what does not fill it is the number sought; take it as the dividend (shi). The Changes takes the Three Powers (sancai) as the expansion divisor; divide the dividend by the divisor (fa); what is obtained is the image number. If the dividend has a remainder, whether 1 or 2, call it 1 in every case, together making the image number. If the image number obtained is 1, it is old yang; if 2, young yin; if 3, young yang; if 4, old yin. Obtaining old yang, draw a doubled (zhong) line; obtaining young yin, draw a split (chai) line; obtaining young yang, draw a single (dan) line; obtaining old yin, draw a crossed (jiao) line. In all, 6 drawings complete a hexagram.\end{uniten}

\begin{unitzh}\unitlabel{shushu-003.12}按：此即前大衍法，末以三歸取爻象，亦属附會。\end{unitzh}
\begin{uniten}Siku editors' note: This is just the preceding Great Expansion method; taking the images of the lines at the end by dividing by 3 (sangui) is also a forced connection.\end{uniten}

\begin{unitzh}\unitlabel{shushu-003.13}草曰：一、二、三、四列右行，立天元一列左行。\end{unitzh}
\begin{uniten}Working: Put 1, 2, 3, 4 in the right column; set up the celestial unknown and put it in the left column.\end{uniten}

\begin{unitzh}\unitlabel{shushu-003.14}元數右行 天元左行\end{unitzh}
\begin{uniten}Original numbers, right column. Celestial unknown, left column.\end{uniten}

\begin{unitzh}\unitlabel{shushu-003.15}以右行一、二、三、四互乘左行異子一，弗乘對位本子，各得衍數。\end{unitzh}
\begin{uniten}With the 1, 2, 3, 4 of the right column cross-multiply the other sons, 1, of the left column, not multiplying the own son in the corresponding position; each obtains an expansion number.\end{uniten}

\begin{unitzh}\unitlabel{shushu-003.16}元數右行 上 副 次 下 衍數左行。併之得五十。\end{unitzh}
\begin{uniten}Original numbers, right column. Upper | secondary (fu) | next (ci) | lower. Expansion numbers, left column. Adding them gives 50.\end{uniten}

\begin{unitzh}\unitlabel{shushu-003.17}乃併左行衍數四位，共計五十，故《易》曰：「大衍之數五十。」算理不可以此五十為用，葢分之為二，則左右手之數竒偶不同，見隂陽之伏數，必須復求用數。先名此曰衍數，以為限率。\end{unitzh}
\begin{uniten}Then add the expansion numbers of the 4 places of the left column: in all 50; therefore the Changes says: "The number of the Great Expansion is 50." By the principle of calculation (suanli) this 50 cannot be taken for use; for when they are divided into 2, the numbers of the left and right hands differ, odd and even, showing the hidden numbers of yin and yang; one must seek the use number anew. First call this the expansion number, taking it as the limiting rate (xianlü).\end{uniten}

\begin{unitzh}\unitlabel{shushu-003.18}遂乃復以一、二、三、四之元數求等數約定。按前術以兩兩相連環求等約之：先以一與二求等，一與三求等，一與四求等，皆得一，各約竒弗約偶，數不變；次以二與三求等，亦得一，約竒弗約偶，數亦不變；及以二與四求等，乃得二，此二只約副數二，變為一，而弗約四；次以三與四求等，亦得一，約竒亦不變。所得一、一、三、四，各為定數母，列右行；仍各立天元一為子，列左行。\end{unitzh}
\begin{uniten}Then, again, with the original numbers 1, 2, 3, 4, seek the equal numbers (dengshu) to reduce and fix them. Following the preceding method, seek the equal number in pairs, in mutually linked rings, and reduce by it. First seek the equal number of 1 and 2, of 1 and 3, and of 1 and 4: all give 1; each reduces the odd and does not reduce the even, and the numbers do not change. Next seek the equal number of 2 and 3: it also gives 1; reducing the odd and not the even, the numbers again do not change. When one seeks the equal number of 2 and 4, it gives 2; this 2 reduces only the secondary number 2, which changes into 1, and does not reduce the 4. Next seek the equal number of 3 and 4: it also gives 1; reducing the odd, again nothing changes. The 1, 1, 3, 4 obtained are each a fixed-number mother (dingshu mu); put them in the right column; again, for each, set up the celestial unknown as the son; put them in the left column.\end{uniten}

\begin{unitzh}\unitlabel{shushu-003.19}定母右行 天行左行\end{unitzh}
\begin{uniten}Fixed mothers, right column. Celestial column (tianxing), left column. [Translator's note: 天行 here evidently for 天元.]\end{uniten}

\begin{unitzh}\unitlabel{shushu-003.20}以右行定母一、一、三、四互乘左行各子一，惟不對乘本子，畢，左上得一十二，左副得一十二，左次得四，次下得三，皆曰衍數。\end{unitzh}
\begin{uniten}With the fixed mothers 1, 1, 3, 4 of the right column cross-multiply each son, 1, of the left column, only not multiplying the own son opposite; when finished, the upper left obtains 12, the secondary left obtains 12, the next left obtains 4, next the lower obtains 3; all are called expansion numbers.\end{uniten}

\begin{unitzh}\unitlabel{shushu-003.21}定母右行。以右定母滿去左 衍數左行。衍衍餘各為竒數。\end{unitzh}
\begin{uniten}Fixed mothers, right column. With the fixed mothers of the right, remove, as often as they are filled, from the left … Expansion numbers, left column. … expansion numbers; the remainders of the expansion numbers each become odd numbers.\end{uniten}

\begin{unitzh}\unitlabel{shushu-003.22}次以各母去衍數：其一母去衍一十二，竒一；其副母一亦去副子一十二，亦各竒一；其次母三去次衍四，亦竒一；其下母四欲去下子三，則不滿，便以三為左下竒數。\end{unitzh}
\begin{uniten}Next remove each mother from the expansion numbers. The first mother, removed from the expansion number 12, leaves an odd 1; the secondary mother 1, likewise removed from the secondary son 12, also each leaves an odd 1; the next mother 3, removed from the next expansion number 4, also leaves an odd 1; the lower mother 4, when one wants to remove it from the lower son 3, does not fill it, so take 3 as the odd number of the lower left.\end{uniten}

\begin{unitzh}\unitlabel{shushu-003.23}定母右行 竒數左行\end{unitzh}
\begin{uniten}Fixed mothers, right column. Odd numbers, left column.\end{uniten}

\begin{unitzh}\unitlabel{shushu-003.24}凡竒數得一者，便為乘率。今左下衍是三，乃與本母四用大衍求一術入之：列衍竒三於右上，定母於右下，立天元一於左上，空其左下。\end{unitzh}
\begin{uniten}Whenever the odd number obtained is 1, it is at once the multiplying rate. Now the lower-left expansion is 3; so, with its own mother 4, enter it into the Great Expansion technique of seeking one: put the expansion odd (yanqi) 3 at the upper right and the fixed mother at the lower right, set up the celestial unknown at the upper left, and leave its lower left empty.\end{uniten}

\begin{unitzh}\unitlabel{shushu-003.25}{\zhsmall 衍 定 商竒 母}{\zhsmall 天元}\end{unitzh}
\begin{uniten}[Diagram labels: expansion | fixed | quotient odd | mother] [Diagram label: celestial unknown]\end{uniten}

\begin{unitzh}\unitlabel{shushu-003.26}先以右上少數三除右下多數四，得一為商；以商一乘左上天元一，只得一，歸左下；其右下餘一。\end{unitzh}
\begin{uniten}First divide the larger number 4 at the lower right by the smaller number 3 at the upper right, obtaining 1 as the quotient; multiply the celestial unknown at the upper left by the quotient 1, obtaining only 1, and return it to the lower left; the lower right leaves 1.\end{uniten}

\begin{unitzh}\unitlabel{shushu-003.27}{\zhsmall 商 衍 定母竒 餘}○{\zhsmall 天 歸元 數}\end{unitzh}
\begin{uniten}[Diagram labels: quotient | expansion | fixed mother odd | remainder] ○ [Diagram labels: celestial | return unknown | number]\end{uniten}

\begin{unitzh}\unitlabel{shushu-003.28}次以右下少數一除右上多數三，須使右上必竒一算乃止，遂於右行最上商二，以除右衍，必竒一，乃以上商命右下定餘一除之，右衍餘一。\end{unitzh}
\begin{uniten}Next divide the larger number 3 at the upper right by the smaller number 1 at the lower right; one must make the upper right necessarily keep an odd 1 rod (suan), and only then stop. So at the very top of the right column put the quotient 2, to divide the right expansion, which must leave an odd 1; then with the quotient above command (ming) the fixed remainder 1 at the lower right and remove it; the right expansion leaves 1.\end{uniten}

\begin{unitzh}\unitlabel{shushu-003.29}{\zhsmall 商 衍竒 定母餘 餘}{\zhsmall 天 歸元 數}\end{unitzh}
\begin{uniten}[Diagram labels: quotient | expansion odd | fixed mother remainder | remainder] [Diagram labels: celestial | return unknown | number]\end{uniten}

\begin{unitzh}\unitlabel{shushu-003.30}次以商二與左下歸數相乘，得二，加入左上天元一内，共得三。\end{unitzh}
\begin{uniten}Next multiply the quotient 2 and the returned number (guishu) at the lower left together, obtaining 2; add it into the celestial unknown at the upper left, giving 3 in all.\end{uniten}

\begin{unitzh}\unitlabel{shushu-003.31}今驗右上衍餘得一，當止。乃以左上三為乘率，與前三者乘率各一，與衍定圖衍數對列之，通計三行。\end{unitzh}
\begin{uniten}Now examine: the expansion remainder at the upper right is 1, so one should stop. Then take the 3 at the upper left as the multiplying rate; together with the multiplying rates of the preceding 3, each 1, set them out opposite the expansion numbers of the expansion–fixed diagram; in all 3 columns.\end{uniten}

\begin{unitzh}\unitlabel{shushu-003.32}定母 衍數 乘率\end{unitzh}
\begin{uniten}Fixed mothers. Expansion numbers. Multiplying rates.\end{uniten}

\begin{unitzh}\unitlabel{shushu-003.33}以乘率對乘左行，畢，左上得一十二，左副得一十二，左次得四，左下得九，皆曰泛用數。\end{unitzh}
\begin{uniten}Multiply the left column by the multiplying rates, each by the one opposite (duicheng); when finished, the upper left obtains 12, the secondary left obtains 12, the next left obtains 4, the lower left obtains 9; all are called provisional use numbers (fanyong shu).\end{uniten}

\begin{unitzh}\unitlabel{shushu-003.34}定母 衍母 泛用\end{unitzh}
\begin{uniten}Fixed mothers. Expansion mother. Provisional uses.\end{uniten}

\begin{unitzh}\unitlabel{shushu-003.35}次以右行一、二、三、四相乘，得一十二，名曰衍母。復推元用等數二約副母二為一，今乃復歸之為二，遂用衍母一十二益於左副一十二内，共為二十四。\end{unitzh}
\begin{uniten}Next multiply together the 1, 2, 3, 4 of the right column, obtaining 12 [Translator's note: 1 × 2 × 3 × 4 is 24; the fixed mothers 1, 1, 3, 4 give 12.]; this is called the expansion mother. Further, reckoning back: originally the equal number 2 was used to reduce the secondary mother 2 to 1; now restore it again to 2, and accordingly add the expansion mother 12 into the secondary left 12; together this makes 24.\end{uniten}

\begin{unitzh}\unitlabel{shushu-003.36}元數 定用\end{unitzh}
\begin{uniten}Original numbers. Fixed uses (dingyong).\end{uniten}

\begin{unitzh}\unitlabel{shushu-003.37}今驗用數圖：右行之一、二、三、四，即是所揲之數；左行一十二併二十四及四與九，併之得四十九，名曰用數。用為蓍草數，故《易》曰「其用四十有九」是也。\end{unitzh}
\begin{uniten}Now examine the use-number diagram: the 1, 2, 3, 4 of the right column are the numbers by which one counts off; in the left column, 12 together with 24, and 4 and 9, added together give 49, which is called the use number. The use is the number of yarrow stalks; therefore the Changes says "its use is 49": this is it.\end{uniten}

\begin{unitzh}\unitlabel{shushu-003.38}假今用蓍四十九，信手分之為二，則左手竒，右手必偶；左手偶，右手必竒。欲使蓍數近大衍五十，非四十九或五十一不可。二數信手分之，必有一竒一偶。故所以用四十九，取七七之數，始有左副二十四内益十二。就其三十七泛為用數，但三十七無意義，兼蓍少太露，是以用四十有九。\end{unitzh}
\begin{uniten}Suppose now one uses 49 yarrow stalks and divides them at random into 2: then if the left hand is odd, the right hand must be even; if the left hand is even, the right hand must be odd. If one wants the number of stalks to be close to the Great Expansion's 50, it cannot be other than 49 or 51. These 2 numbers, divided at random, must give 1 odd and 1 even. Therefore the reason for using 49 is to take the number 7 times 7; only then is there the 24 of the secondary left, within which 12 was added. As for taking its 37 provisionally as the use number: 37, however, has no meaning, and moreover with few stalks it is too exposed; that is why 49 is used.\end{uniten}

\begin{unitzh}\unitlabel{shushu-003.39}凡揲蓍求一爻之數，欲得一、二、三、四岀於無餘，必令揲者不得知，故以四十九蓍分之為二，只用左手之數。假令左手分得三十三，自一一楪之，必竒一，故不繁楪，乃徑掛一，故《易》曰：「分而為二以象兩，掛一以象三。」\end{unitzh}
\begin{uniten}In general, when counting off the yarrow stalks to find the number of a single line, one wishes to obtain 1, 2, 3, 4 issuing from no remainder, and one must make sure the one counting off cannot know it; therefore the 49 stalks are divided into 2 and only the number in the left hand is used. Suppose the left hand receives 33 in the division: counting them off 1 by 1, there must be an odd 1; therefore one does not bother to count them off, but directly hangs 1; therefore the Changes says: "Divide them into 2 to represent the Two; hang 1 to represent the Three."\end{uniten}

\begin{unitzh}\unitlabel{shushu-003.40}次後又令筮人以二二揲之，其三十三亦竒一，故歸奇於㧅；又令之以三三揲之，其三十三必奇三，故歸奇於㧅；又令之以四四揲之，又奇一，亦歸奇於㧅。與前掛一併三度揲，通有四㧅，乃得一、一、三、一。其掛一者，乘用數圖左上用數一十二；其二揲㧅一者，乘左副用數二十四；其三揲㧅三者，乘左次用數四，得一十二；其四揲一者，乘左下用數九。\end{unitzh}
\begin{uniten}After that, again have the diviner count them off by 2s: the 33 also leaves an odd 1, therefore return the odd ones to the finger-clasp (le). Again have him count them off by 3s: the 33 must leave an odd 3, therefore return the odd ones to the finger-clasp. Again have him count them off by 4s: again an odd 1, and again return the odd ones to the finger-clasp. With the preceding hanging of 1 together with the 3 countings-off, there are 4 finger-clasps in all, giving 1, 1, 3, 1. The hung 1 multiplies the use number 12 at the upper left of the use-number diagram; the finger-clasp 1 of counting-off 2 multiplies the use number 24 of the secondary left; the finger-clasp 3 of counting-off 3 multiplies the use number 4 of the next left, giving 12; the 1 of counting-off 4 multiplies the use number 9 of the lower left.\end{uniten}

\begin{unitzh}\unitlabel{shushu-003.41}用數 左行三㧅，謂之三變。掛一得一十二，㧅一得二十四，㧅三得一十二，又㧅一得九，並為總數。\end{unitzh}
\begin{uniten}Use numbers. Left column: the 3 finger-clasps are called the 3 changes. The hung 1 gives 12; finger-clasp 1 gives 24; finger-clasp 3 gives 12; again finger-clasp 1 gives 9; all are total numbers.\end{uniten}

\begin{unitzh}\unitlabel{shushu-003.42}併此四總，得五十七，不問所握幾何，乃滿衍母一十二去之，得不滿者九{\zhsmall 或使知其所握五十七，亦滿衍母去之，亦只得九數}，以為實，用三才衍法約之，得三，乃畫少陽單爻{\zhsmall 或不滿得八、得七為實，皆命為三}。他皆倣此。\end{unitzh}
\begin{uniten}Add these 4 totals, obtaining 57; without asking how many were grasped, remove the expansion mother 12 as often as it is filled, obtaining the 9 that does not fill it [Note: Or, supposing one knows that what was grasped is 57, removing the expansion mother as often as it is filled likewise gives only the number 9]; take it as the dividend; reduce it by the expansion divisor of the Three Powers, obtaining 3; then draw the single line of young yang [Note: Or, if what does not fill it gives 8 or 7 as the dividend, call it 3 in every case]. All the others imitate this.\end{uniten}

\begin{unitzh}\unitlabel{shushu-003.43}術意謂揲二、楪三、楪四者，凡三度，復以三十三從頭數揲之，故曰「三變而成爻」；既卦有六爻，必有十八變，故曰「十有八變而成卦」。\end{unitzh}
\begin{uniten}The intent of the method is that counting off by 2, counting off by 3 and counting off by 4 are 3 times in all, each time again counting off the 33 from the beginning; therefore it says "with 3 changes a line is completed". Since a hexagram has 6 lines, there must be 18 changes; therefore it says "with 18 changes a hexagram is completed".\end{uniten}

\begin{unitzh}\unitlabel{shushu-003.44}按：此條强援蓍卦牽附衍數，致本法反晦。今以本法列於前，則其弊自見矣。\end{unitzh}
\begin{uniten}Siku editors' note: This item forcibly invokes the yarrow stalks and hexagrams and drags the expansion numbers into connection with them, so that the basic method is, on the contrary, obscured. Now that we have set out the basic method in front, its faults show for themselves.\end{uniten}

\begin{unitzh}\unitlabel{shushu-003.45}古厯會積\end{unitzh}
\begin{uniten}Accumulated Conjunction of the Ancient Calendars (Guli huiji)\end{uniten}

\begin{unitzh}\unitlabel{shushu-004.1}問：古厯冬至以三百六十五日四分日之一，朔䇿以二十九日九百四十分日之四百九十九，甲子以六十日，各為一周。假令至淳祐丙午十一月丙辰朔，初五日庚申冬至，初九日甲子，欲求古厯氣朔甲子一會積年、積月、積日及厯過、未至年數各幾何？\end{unitzh}
\begin{uniten}Question: In the ancient calendar (guli), the winter solstice (dongzhi) is taken as 365 1/4 days, the lunation (shuoce) as 29 499/940 days, and the jiazi as 60 days, each making one period (zhou). Suppose that, in the bingwu year of Chunyou, the new moon (shuo) of the 11th month is on a bingchen day, the 5th day, a gengshen day, is the winter solstice, and the 9th day is a jiazi day. One wishes to find, for the ancient calendar, the accumulated years, accumulated months and accumulated days of one meeting (hui) of the qi (氣), the new moon and the jiazi, and also the numbers of years elapsed (liguo) and not yet reached (weizhi). How much is each?\end{uniten}

\begin{unitzh}\unitlabel{shushu-004.2}按此題嵗實、朔䇿皆古法用數，淳祐丙午嵗合朔、冬至干支乃宋開禧法所步，題數已不相𫎇，即推算無誤，亦未合，况不能無誤耶？\end{unitzh}
\begin{uniten}[Siku editors: In this problem the year value (suishi) and the lunation are both numbers used in the ancient method, while the stems and branches (ganzhi) of the conjunction (heshuo) and the winter solstice in the bingwu year of Chunyou are what was computed by the Kaixi method of the Song. The numbers of the problem already do not fit each other, so even if the computation were without error, it would still not agree. How much less can it agree when the computation cannot be without error?]\end{uniten}

\begin{unitzh}\unitlabel{shushu-004.3}荅曰：一會積一萬八千二百四十年，二十二萬五千六百月，六百六十六萬二千一百六十日。厯過九千一百六十三年，未至九千七十七年。按荅數皆不合。\end{unitzh}
\begin{uniten}Answer: One meeting accumulates 18240 years, 225600 months, 6662160 days. Elapsed, 9163 years; not yet reached, 9077 years. [Siku editors: The numbers of the answer all do not agree.]\end{uniten}

\begin{unitzh}\unitlabel{shushu-004.4}術曰：同前。置問數，{\zhsmall 有分者通之，互乘之，得通數。}求總等，不約一位，約衆位，各得元法。連環求等，約竒弗約偶，各得定母。{\zhsmall 本題欲求一會，不復乘偶。}以定相乘為衍母，定除母得衍數，滿定去衍得竒，以大衍入之得乘率，以乘衍數得泛用數。併諸泛以課衍母，如泛内多倍數者損之。乃驗元數竒偶同類處，各損半倍，{\zhsmall 或三位同類者，三約衍母損泛。}各得正用。然後推氣朔不及，或所遇甲子日數，乘正用，加减之為總，滿衍去之，餘為所求厯過率實，如紀元法而一為厯過。以氣元法除衍母，得一會積年；以氣周日刻乘一會年，得一會積日；以朔元法除衍母，得一會積月數。\end{unitzh}
\begin{uniten}Method: As before. Set out the numbers of the question [Note: where there are fractions, convert them (tong) and cross-multiply them (hucheng), obtaining the converted numbers (tongshu)]. Find the total common divisor (zongdeng). Do not reduce one position, and reduce the many positions; each gives an original divisor (yuanfa). Find the common divisors in a linked chain (lianhuan qiudeng), reducing the odd and not reducing the even; each gives a fixed denominator (dingmu). [Note: In this problem one wishes to find one meeting, so one does not multiply the even again.]

Multiply the fixed (denominators) together to make the extended denominator (yanmu). Divide the denominator by each fixed (denominator) to obtain the extended number (yanshu). Remove the fixed (denominator) from the extended (number) as often as it fills it, obtaining the remainder (ji). Enter this into the great extension (dayan) to obtain the multiplying rate (chenglü). Multiply the extended number by it to obtain the provisional used number (fan yongshu).

Add the provisional (numbers) together in order to test them against the extended denominator. If within the provisional (sum) there are multiples (bei) in excess, reduce them. Then examine the original numbers (yuanshu): where odd and even are of the same kind, reduce each by half a multiple [Note: or, where three positions are of the same kind, divide the extended denominator by 3 to reduce the provisional (numbers)]. Each then gives the correct used number (zhengyong).

After that, work out the shortfalls (buji) of the qi and the new moon, or the numbers of days to the jiazi met with. Multiply the correct used numbers by them, and add or subtract them to make the total (zong). Remove the extended (denominator) as often as it fills it. The remainder is the dividend (shi) of the rate (lü) of the elapsed (years) sought; divide it by the cycle original divisor (jiyuan fa) to make the elapsed (years).

Divide the extended denominator by the qi original divisor (qiyuan fa) to obtain the accumulated years of one meeting. Multiply the years of one meeting by the days and ke (刻) of the qi period to obtain the accumulated days of one meeting. Divide the extended denominator by the new-moon original divisor (shuoyuan fa) to obtain the number of accumulated months of one meeting.\end{uniten}

\begin{unitzh}\unitlabel{shushu-004.5}按「如紀元法而一」、「以氣元法除衍母」二語皆誤，故得數不合，皆當以氣分為法，葢氣分即嵗實分也。\end{unitzh}
\begin{uniten}[Siku editors: The two phrases "divide it by the cycle original divisor" and "divide the extended denominator by the qi original divisor" are both wrong, and therefore the numbers obtained do not agree. In both cases one should take the qi parts (qifen) as divisor (fa), for the qi parts are just the parts of the year value.]\end{uniten}

\begin{unitzh}\unitlabel{shushu-004.6}右本題問氣朔甲子相距日數，係開禧厯推倒，或甲子日在氣朔之間，及非十一月前後者，其總數必滿母，贅去之，所得厯過年數，尾位雖倫，首位必異。今設問以明大衍之理，初不計其前多後少之厯過。\end{unitzh}
\begin{uniten}In the present problem above, the numbers of days separating the qi, the new moon and the jiazi that are asked about were worked out (tuidao) by the Kaixi calendar. Where the jiazi day lies between the qi and the new moon, or where it is not around the 11th month, the total will necessarily fill the denominator, and the superfluous is removed. In the number of elapsed years obtained, although the last places are in order, the first places will necessarily differ. Now the question is posed in order to make clear the principle of the great extension, and it does not at all reckon the elapsed (years) of "more before and less after".\end{uniten}

\begin{unitzh}\unitlabel{shushu-004.7}按此數語，葢因得數不合而自解之。然算家終以得數為凖，得數不合，則無以取信於人矣。\end{unitzh}
\begin{uniten}[Siku editors: These several phrases are presumably his own explanation, given because the numbers obtained do not agree. But mathematicians in the end take the numbers obtained as the standard; if the numbers obtained do not agree, there is nothing by which to win people's trust.]\end{uniten}

\begin{unitzh}\unitlabel{shushu-004.8}草曰：置問數，冬至三百六十五日四分日之一，朔䇿二十九日九百四十分日之四百九十九，甲子六十日，各通分内子，互乘之，列三等位，具圖如後。\end{unitzh}
\begin{uniten}Working: Set out the numbers of the question: winter solstice, 365 1/4 days; lunation, 29 499/940 days; jiazi, 60 days. For each, convert to parts and add in the numerator (tongfen nazi), then cross-multiply them and set them out in three equal positions; the diagram is as follows.\end{uniten}

\begin{unitzh}\unitlabel{shushu-004.9}冬至得一千四百六十一，朔實得二萬七千七百五十九，甲子無母，只是六十。列三行互乘之，具圖如後。{\zhsmall 氣通 母 總等不約紀分}\end{unitzh}
\begin{uniten}For the winter solstice one obtains 1461; the new-moon value (shuoshi) gives 27759; the jiazi has no denominator and is just 60. Set them out in three columns and cross-multiply them; the diagram is as follows. [Note: "qi, converted (tong)"; "denominator (mu)"; "total common divisor; do not reduce the cycle parts (jifen)".] [Translator's note: stray labels from the diagrams, which are not transcribed.]\end{uniten}

\begin{unitzh}\unitlabel{shushu-004.10}以三行互乘：右得一百三十七萬三千三百四十，為氣分；中得一十一萬一千三十六，為朔分；左得二十二萬五千六百，為紀分。先求總等，得一十二，{\zhsmall 按十二乃朔分、紀分所求等數，亦可為氣分等數，故為總等。}乃存紀分一位不約，只以等一十二約氣分，得一十一萬四千四百四十五，又約朔分，得九千二百五十三，皆為元法。\end{unitzh}
\begin{uniten}Cross-multiplying the three columns: on the right one obtains 1373340 as the qi parts; in the middle, 111036 as the new-moon parts (shuofen); on the left, 225600 as the cycle parts. First find the total common divisor, obtaining 12 [Siku editors: 12 is the common divisor found for the new-moon parts and the cycle parts, and it can also serve as a common divisor of the qi parts; therefore it is the total common divisor]. Then keep the single position of the cycle parts unreduced. With the common divisor 12 reduce only the qi parts, obtaining 114445, and the new-moon parts, obtaining 9253. Both are original divisors.\end{uniten}

\begin{unitzh}\unitlabel{shushu-004.11}乃以連環求等：次以紀元二十二萬五千六百與朔元九千二百五十三求等，得一，不約；又以紀元與氣元一十一萬四千四百四十五求等，得二百三十五，只約氣元，得四百八十七；次以氣元四百八十七與朔元九千二百五十三求等，得四百八十七，只約朔元九千二百五十三，得一十九。約偏畢，得四百八十七為氣定，得一十九為氣定，得二十二萬五千六百為紀定。以三定相乘，得二十億八千七百四十七萬六千八百，為衍母，具圖如後。\end{unitzh}
\begin{uniten}Then find the common divisors in a linked chain. Next, find the common divisor of the cycle original 225600 and the new-moon original 9253, obtaining 1: do not reduce. Again, find the common divisor of the cycle original and the qi original 114445, obtaining 235; reduce only the qi original, obtaining 487. Next, find the common divisor of the qi original 487 and the new-moon original 9253, obtaining 487; reduce only the new-moon original 9253, obtaining 19.

The reductions being completed throughout, one obtains 487 as the fixed (number) of the qi (qiding), 19 as the fixed (number) of the qi [Translator's note: this second 氣定 evidently stands for 朔定, the fixed number of the new moon], and 225600 as the fixed (number) of the cycle (jiding). Multiply the three fixed (numbers) together, obtaining 2087476800 as the extended denominator; the diagram is as follows.\end{uniten}

\begin{unitzh}\unitlabel{shushu-004.12}各以定數約衍母，各得衍數：氣得四百二十八萬六千四百，朔得一億九百八十六萬七千二百，紀得九千二百五十三，寄左行。各滿定數去之，各得竒數。\end{unitzh}
\begin{uniten}Reduce the extended denominator by each of the fixed numbers (dingshu), obtaining for each the extended number: for the qi, 4286400; for the new moon, 109867200; for the cycle, 9253. Lodge them in the left column. From each remove the fixed number as often as it fills it, obtaining for each the remainder number (jishu).\end{uniten}

\begin{unitzh}\unitlabel{shushu-004.13}氣竒得三百一十三，朔竒得四，紀竒得九千二百五十三，各與定數用大衍求一，各得乘數，列右行，對寄左行衍數，具圖如後。\end{unitzh}
\begin{uniten}The qi remainder is 313, the new-moon remainder 4, the cycle remainder 9253. For each, together with its fixed number, use the great-extension method of finding one (dayan qiuyi), obtaining for each the multiplying number (chengshu). Set them out in the right column, facing the extended numbers lodged in the left column; the diagram is as follows.\end{uniten}

\begin{unitzh}\unitlabel{shushu-004.14}各以大數入之：氣乘率得四百七十三，朔乘率得五，紀乘率得一十七萬二千七百一十七，對左行衍數。以右行乘率對乘左行衍數：氣泛得二十億二千七百四十六萬七千二百，朔泛得五億四千九百三十三萬六千，紀泛得一十五億九千八百一十五萬四百一十，具圖如後。\end{unitzh}
\begin{uniten}Enter each into the great number (dashu): the qi multiplying rate is 473, the new-moon multiplying rate 5, the cycle multiplying rate 172717, facing the extended numbers of the left column. Multiply the extended numbers of the left column by the multiplying rates of the right column facing them: the qi provisional is 2027467200, the new-moon provisional 549336000, the cycle provisional 1598150410; the diagram is as follows.\end{uniten}

\begin{unitzh}\unitlabel{shushu-004.15}右列用數併之，共得四十一億七千四百九十五萬三千六百一，為泛用數，與衍母二十億八千七百四十七萬六千八百驗之，在衍母以上，就以衍母除泛，得二，乃知泛内多一倍母數，當於各用内損去所多一倍。按術驗法元圖内諸元數，奇偶同類者，各損其半。今驗法元圖，氣元尾數是五，紀元尾數是六百，為俱五同類，乃以術母二十億八千七百四十七萬六千八百折半，得一十億四千三百七十三萬八千四百，以損泛用圖内氣泛、紀泛，畢，其朔泛不損，各得氣、朔、紀正用數：其氣正用得九億八千三百七十二萬八千八百，朔正用五億四千九百三十三萬六千，紀正用五億五千四百四十一萬二千一，列為正用圖在前。\end{unitzh}
\begin{uniten}Adding together the used numbers listed above, one obtains in all 4174953601 as the provisional used number. Tested against the extended denominator 2087476800, it lies above the extended denominator. Dividing the provisional by the extended denominator, one obtains 2, and so one knows that within the provisional there is one multiple of the denominator in excess. One should remove from the respective used numbers the one multiple in excess.

According to the method, examine the original numbers in the diagram of divisors and originals (fayuan tu): those whose odd and even are of the same kind are each reduced by half of it. Now, examining the diagram of divisors and originals, the last digit of the qi original is 5 and the last digits of the cycle original are 600, so both are of the same kind with 5.

Then halve the method denominator (shumu) 2087476800 [Translator's note: 術母 evidently stands for 衍母, the extended denominator], obtaining 1043738400, and with it reduce the qi provisional and the cycle provisional in the diagram of provisional used numbers. When this is done, the new-moon provisional is not reduced, and one obtains for each of the qi, the new moon and the cycle the correct used number: the qi correct used number is 983728800, the new-moon correct used number 549336000, and the cycle correct used number 554412001. These are set out as the diagram of correct used numbers, placed before.\end{uniten}

\begin{unitzh}\unitlabel{shushu-004.16}既得正用數，次驗問題：十一月朔日丙辰，冬至初五日庚申，初九日甲子。乃以初一减初九甲子，餘八日，為朔不及；次以初五亦减初九甲子，餘四日，為氣不及。以二不及各乘正用，得數，具圖如後。\end{unitzh}
\begin{uniten}Having obtained the correct used numbers, next examine the question: in the 11th month the new-moon day is bingchen; the winter solstice is the 5th day, gengshen; the 9th day is jiazi. Then subtract the 1st from the 9th, jiazi, leaving 8 days, as the new-moon shortfall. Next, likewise subtract the 5th from the 9th, jiazi, leaving 4 days, as the qi shortfall. Multiply the correct used numbers by the two shortfalls respectively to obtain the numbers; the diagram is as follows.\end{uniten}

\begin{unitzh}\unitlabel{shushu-004.17}先以氣不及甲子四日，以乘氣正用數九億八千三百七十二萬八千八百，得三十九億三千四百九十一萬五千二百，為氣總；次以朔不及甲子八日數，以乘其朔正用數五億五千九百三十三萬六千，得四十三億九千四百六十八萬八千，為朔總。併之，得八十三億二千九百六十萬三千二百，為總數，滿母二十億八千七百四十七萬六千八百去之，不滿二十億六千七百一十七萬二千八百，為所求率實，具圖如後。\end{unitzh}
\begin{uniten}First, with the qi shortfall to the jiazi, 4 days, multiply the qi correct used number 983728800, obtaining 3934915200 as the qi total. Next, with the new-moon shortfall to the jiazi, the number 8 days, multiply its new-moon correct used number 559336000, obtaining 4394688000 as the new-moon total. Adding them, one obtains 8329603200 as the total number. Remove the denominator 2087476800 as often as it fills it; what does not fill it, 2067172800, is the rate dividend sought. The diagram is as follows.\end{uniten}

\begin{unitzh}\unitlabel{shushu-004.18}按求積嵗，應以甲子距冬至前之日分乘紀用數為紀總，以合朔距冬至前之日分乘朔用數為朔總，併紀總、朔總，滿衍母去之，以嵗實分除之，即已過積年。草内以冬至距甲子前之日分乘氣用數，合朔距甲子前之日分乘朔用數，併之，乃求紀周法，非求嵗周法也，故不合。\end{unitzh}
\begin{uniten}[Siku editors: To find the accumulated years, one should multiply the cycle used number by the days and parts by which the jiazi lies before the winter solstice, making the cycle total. Multiply the new-moon used number by the days and parts by which the conjunction lies before the winter solstice, making the new-moon total. Add the cycle total and the new-moon total, remove the extended denominator as often as it fills it, and divide by the parts of the year value: that is the accumulated years already elapsed. In the working, the qi used number is multiplied by the days and parts by which the winter solstice lies before the jiazi, the new-moon used number by the days and parts by which the conjunction lies before the jiazi, and these are added. This is the method for finding the cycle periods, not the method for finding the year periods; therefore it does not agree.]\end{uniten}

\begin{unitzh}\unitlabel{shushu-004.19}置所得率實二十億六千七百一十七萬二千八百，如法元圖紀元法二十二萬五千六百而一，得九千一百六十三年，為厯過年數。次置衍母二十億八千七百四十七萬六千八百為實，如法元圖氣元一十一萬四千四百四十五為法而一，得一萬八千二百四十年，為氣朔甲子一會積年，内减厯過九千一百六十三年，餘九千七十七年，為未至年數。次以冬至周日三百六十五日二十五刻乘積一會年一萬八千二百四十，得六百六十六萬二千一百六十日，為一會積日。又以術母為實，如法元圖朔元法九千二百五十二而一，得二十二萬五千六百月，為一會積月。合問。\end{unitzh}
\begin{uniten}Set out the rate dividend obtained, 2067172800. Divide it by the cycle original divisor of the diagram of divisors and originals, 225600, obtaining 9163 years as the number of years elapsed. Next set out the extended denominator 2087476800 as dividend (shi). Divide it by the qi original of the diagram of divisors and originals, 114445, as divisor (fa), obtaining 18240 years as the accumulated years of one meeting of the qi, the new moon and the jiazi. From this subtract the elapsed 9163 years, leaving 9077 years as the number of years not yet reached. Next multiply the accumulated years of one meeting, 18240, by the period days of the winter solstice, 365 days 25 ke, obtaining 6662160 days as the accumulated days of one meeting. Again take the method denominator as dividend and divide it by the new-moon original divisor of the diagram of divisors and originals, 9252, obtaining 225600 months as the accumulated months of one meeting. This answers the question.\end{uniten}

\begin{unitzh}\unitlabel{shushu-004.20}按此紀元即紀分，以紀分除率實，乃紀周數，非已過年數也。求一會積年，當以氣分為法，以氣元為法亦誤。此二數既誤，餘數無是者矣。然題已不合，既法合，數亦不能合也。今改設一題於後，以明其法焉。\end{unitzh}
\begin{uniten}[Siku editors: This cycle original is just the cycle parts. Dividing the rate dividend by the cycle parts gives the number of cycle periods, not the number of years already elapsed. To find the accumulated years of one meeting one should take the qi parts as divisor; taking the qi original as divisor is also wrong. Since these two numbers are wrong, none of the remaining numbers is right. But the problem itself already does not agree, so even if the method agreed, the numbers could not agree either. Now a problem is set anew below, to make the method clear.]\end{uniten}

\begin{unitzh}\unitlabel{shushu-004.21}設古法嵗實三百六十五日四分日之一，朔䇿二十九日九百四十分日之四百九十九，甲子六十日。假令十一月平朔辛巳日四百七十分日之一百一十三，冬至癸夘日子正初刻。問距前後甲子日子正初刻合朔冬至之年數各幾何？荅曰：距前八百七十六年，距後六百四十四年。\end{unitzh}
\begin{uniten}Assume that in the ancient method the year value is 365 1/4 days, the lunation 29 499/940 days, and the jiazi 60 days. Suppose that in the 11th month the mean new moon (pingshuo) is on a xinsi day at 113/470 of a day, and the winter solstice is on a guimao day at the first ke of exact midnight (zizheng chuke). Question: how many years in each case is the distance, before and after, to the conjunction and the winter solstice at the first ke of exact midnight of a jiazi day? Answer: the distance before is 876 years; the distance after, 644 years.\end{uniten}

\begin{unitzh}\unitlabel{shushu-005.1}法：按前法求至正用，乃以冬至癸夘距甲子後三十九日為紀餘，以日法{\zhsmall 即氣分母、朔分母相乘之數}三千七百六十分通之，得十四萬六千六百四十為紀餘分，以乘紀正用，得紀總八十一兆二千九百八十九億七千五百八十二萬六千六百四十。次以平朔辛巳距甲子十七日又四百七十分之一百一十三，與冬至距甲子三十九日相减，得二十一日又四百七十分日之三百五十七，以日法三千七百六十通之，得八萬一千八百一十六為朔餘分，以乘朔正用，得朔總四十四兆九千四百四十四億七千四百一十七萬六千。併二總數，滿衍母去之，得率實十二億零三百零四萬五千八百四十為實，以嵗實一百三十七萬三千三百四十為法除之，得八百七十六年，為距前氣朔甲子會積之年數。又以衍母為實，以嵗實分為法除之，得一千五百二十年，為前會積距後會積之年數，减去距前會積之年數，餘六百四十四年，為距後會積之年數。\end{unitzh}
\begin{uniten}Rule: following the preceding rule, compute up to the correct use-numbers (zhengyong). Then take the 39 days by which the winter solstice, guimao, comes after jiazi as the cycle remainder (jiyu); convert (tong) it by the day divisor (rifa) [Note: that is, the number obtained by multiplying together the denominator of the qi parts and the denominator of the new-moon parts], 3760 parts, obtaining 146,640 as the cycle-remainder parts; with this multiply the correct use-number of the cycle, obtaining the cycle total (jizong) 81,298,975,826,640. Next take the mean new moon (pingshuo), xinsi, 17 113/470 days from jiazi, and the winter solstice, 39 days from jiazi, and subtract the one from the other, obtaining 21 357/470 days; convert this by the day divisor 3760, obtaining 81,816 as the new-moon-remainder parts; with this multiply the correct use-number of the new moon, obtaining the new-moon total 44,944,474,176,000. Add the two total numbers and cast out whole multiples of the extension modulus (yanmu), obtaining the rate-dividend (lüshi) 1,203,045,840 as the dividend (shi); take the year-length (suishi) 1,373,340 as the divisor (fa) and divide, obtaining 876 years as the number of years back to the preceding accumulated conjunction (huiji) of qi, new moon and jiazi. Further take the extension modulus as dividend and the year-length in parts as divisor and divide, obtaining 1520 years as the number of years from the preceding accumulated conjunction to the following accumulated conjunction; subtract the number of years back to the preceding accumulated conjunction, leaving 644 years as the number of years to the following accumulated conjunction.\end{uniten}

\begin{unitzh}\unitlabel{shushu-005.2}既得積年，若欲還原求題中干支時刻，則以前會之積年與嵗實相乘，得三十一萬九千九百五十九為積日，滿紀法六十去之，餘三十九日，自初日起甲子，得冬至為癸夘日子正初刻。又置積日，以朔䇿日分九百四十通之為實，以朔䇿通分納子為法除之，得一萬零八百三十四為積朔，餘二萬零四百五十四，又為實，以朔䇿日分九百四十為法除之，得二十一日又九百四十分之七百一十四，約之為四百七十分日之三百五十七，為距冬至前日數，與甲子距冬至前三十九日相减，得一十七日又四百七十分日之一百一十三，為距甲子後日數，自初日起甲子，得辛巳為平朔干支，悉與題合。\end{unitzh}
\begin{uniten}Having obtained the accumulated years, if one wishes to return to the original (huanyuan) and find the stems-and-branches (ganzhi) and times in the problem, then multiply the accumulated years of the preceding conjunction by the year-length, obtaining 319,959 as the accumulated days; cast out whole multiples of the cycle divisor (jifa), 60, leaving 39 days; counting from the first day as jiazi, the winter solstice is found to fall on the guimao day, at midnight (zizheng), initial ke (刻). Again set up the accumulated days and convert them by the day-parts of the lunation (shuoce), 940, to make the dividend; take the lunation converted into parts with the numerator added in (tongfen nazi) as divisor and divide, obtaining 10,834 as the accumulated new moons, with a remainder of 20,454, which again is made the dividend; take the day-parts of the lunation, 940, as divisor and divide, obtaining 21 714/940 days; reduce it to 357/470 of a day; this is the number of days before the winter solstice; subtract it from the 39 days by which jiazi is before the winter solstice, obtaining 17 113/470 days as the number of days after jiazi; counting from the first day as jiazi, xinsi is obtained as the stem-and-branch of the mean new moon. All agree with the problem.\end{uniten}

\begin{unitzh}\unitlabel{shushu-005.3}推庫額錢\end{unitzh}
\begin{uniten}Reckoning the quota cash of the treasuries\end{uniten}

\begin{unitzh}\unitlabel{shushu-005.4}問：有外邑七庫，日納息足錢適等，遞年成貫整納。近縁見錢稀少，聽各庫照當處市陌凖解舊會。其甲庫有零錢一十文，丁、庚二庫各零四文，戊庫零六文，餘庫無零錢。甲庫所在市陌一十二文，遞减一文，至庚庫而止。欲求諸庫日息原納足錢、展省及今納舊會并大小月分各幾何？\end{unitzh}
\begin{uniten}Question: There are seven treasuries in outlying towns whose daily deliveries of interest in full cash (zuqian) are exactly equal; year after year these were delivered whole, in complete guan (貫). Recently, because ready cash has been scarce, each treasury has been permitted to remit old notes (jiuhui) as the equivalent, according to the market string (shimo) of its locality. Treasury A (jia) has odd cash (lingqian) of 10 wen (文); the two Treasuries D (ding) and G (geng) each have an odd 4 wen; Treasury E (wu) has an odd 6 wen; the remaining treasuries have no odd cash. The market string where Treasury A is located is 12 wen, decreasing successively by 1 wen and stopping at Treasury G. It is desired to find the full cash of daily interest originally delivered by the treasuries, the expansion into short strings (zhansheng), and the old notes now delivered, together with the shares for the large and small months: how much is each?\end{uniten}

\begin{unitzh}\unitlabel{shushu-005.5}按：題意係七邑日納共錢同數，以各邑市陌數計之，或適足，或有餘，多寡不同：甲陌十二則餘十，乙陌十一、丙陌十則無餘，丁陌九則餘四，戊陌八則餘六，己陌七則無餘，庚陌六則餘四，以求共錢同數，此本術也。又問展省、舊會。按草中，展省乃官省陌以七十七為一百所展日息共錢之數；舊會乃以各陌數為一百所升日息共錢之數。二者在本術中已贅，且不明言展省、舊會用數求法，皆故為溟涬也。\end{unitzh}
\begin{uniten}Note: the meaning of the problem is that the seven towns deliver daily a common cash of the same number; reckoned by the market-string number of each town, it is either exactly sufficient or leaves a remainder, the amounts differing: A's string is 12, so the remainder is 10; B's (yi) string is 11 and C's (bing) string is 10, so there is no remainder; D's string is 9, so the remainder is 4; E's string is 8, so the remainder is 6; F's (ji) string is 7, so there is no remainder; G's string is 6, so the remainder is 4; from this one seeks the common cash of the same number. This is the basic method (benshu). It further asks for the expansion into short strings and the old notes. According to the working, the expansion into short strings is the amount of the common cash of daily interest expanded by the official short string (guan shengmo), which takes 77 as 100; the old notes are the amount of the common cash of daily interest raised by taking each string number as 100. These two are already superfluous in the basic method; moreover it does not state clearly the numbers used and the methods of finding the expansion into short strings and the old notes: all this is deliberately made obscure.\end{uniten}

\begin{unitzh}\unitlabel{shushu-005.6}荅曰：諸庫納日息足錢二十貫九百五十文。展省三十五貫文。甲庫日息舊會二百二十四貫五百一十文{\zhsmall 按應作五百文又六分文之五}，大月舊會六千七百三十七貫五百文{\zhsmall 按少二十五文}，小月舊會六千五百一十二貫九百二文{\zhsmall 按應作六千五百一十貫五百又六分文之一}。乙庫日息舊會二百四十五貫文，大月舊會七千三百五十貫文，小月舊會七千五百貫文{\zhsmall 按應作七千一百五貫}。丙庫日息舊會二百六十九貫五百文，大月舊會八千八十五貫文，小月舊會七千八百一十五貫五百文。\end{unitzh}
\begin{uniten}Answer: The full cash of daily interest delivered by the treasuries: 20 guan 950 wen [Translator's note: the figures below imply 26 guan 950 wen]. Expansion into short strings: 35 guan. Treasury A: daily interest in old notes 224 guan 510 wen [Siku editors: it should read 500 5/6 wen]; large month, old notes 6737 guan 500 wen [Siku editors: 25 wen short]; small month, old notes 6512 guan 902 wen [Siku editors: it should read 6510 guan 500 1/6 wen]. Treasury B: daily interest in old notes 245 guan; large month, old notes 7350 guan; small month, old notes 7500 guan [Siku editors: it should read 7105 guan]. Treasury C: daily interest in old notes 269 guan 500 wen; large month, old notes 8085 guan; small month, old notes 7815 guan 500 wen.\end{uniten}

\begin{unitzh}\unitlabel{shushu-005.7}丁庫日息舊會二百九十九貫四百四文{\zhsmall 按應作四百四十四文又九分文之四}，大月舊會八千九百八十三貫三百三文{\zhsmall 按少三十文又三分文之一}，小月舊會八千六百八十三貫八百八文{\zhsmall 按少八十文又九分文之八}。戊庫日息舊會三百三十六貫八百六文{\zhsmall 按應作七十五文}，大月舊會一萬一百六貫二百四文{\zhsmall 按應作二百五十文}，小月舊會九千七百六十九貫三百六文{\zhsmall 按應作七十五文}。己庫日息舊會三百八十五貫文，大月舊會一萬一千五百五十貫文，小月舊會一萬一千一百六十貫文{\zhsmall 按少五貫}。庚庫日息舊會四百四十九貫一百四文{\zhsmall 按應作一百六十六文又分三文之二}，大月舊會一萬三千四百七十五貫文，小月舊會一萬三千二十四貫八百二文{\zhsmall 按應作二萬三千二十五貫八百三十三文又三分文之一}。\end{unitzh}
\begin{uniten}Treasury D: daily interest in old notes 299 guan 404 wen [Siku editors: it should read 444 4/9 wen]; large month, old notes 8983 guan 303 wen [Siku editors: 30 1/3 wen short]; small month, old notes 8683 guan 808 wen [Siku editors: 80 8/9 wen short]. Treasury E: daily interest in old notes 336 guan 806 wen [Siku editors: it should read 75 wen]; large month, old notes 10,106 guan 204 wen [Siku editors: it should read 250 wen]; small month, old notes 9769 guan 306 wen [Siku editors: it should read 75 wen]. Treasury F: daily interest in old notes 385 guan; large month, old notes 11,550 guan; small month, old notes 11,160 guan [Siku editors: 5 guan short]. Treasury G: daily interest in old notes 449 guan 104 wen [Siku editors: it should read 166 2/3 wen]; large month, old notes 13,475 guan; small month, old notes 13,024 guan 802 wen [Siku editors: it should read 23,025 guan 833 1/3 wen].\end{uniten}

\begin{unitzh}\unitlabel{shushu-005.8}術曰：以大衍求之。置甲庫市陌，以庫减庫减之，各得諸庫原陌。連環求等，約竒弗約偶{\zhsmall 按此特為等數為偶者言之，若等數為竒者，則約偶弗約竒}，得定母。諸定相乘為衍母。以定約衍母，得衍數。衍數同衍母者，去之為無{\zhsmall 無者借之同類}。其各滿定母去，餘為竒數。以竒、定用大衍求乘率，乘衍數為用數。無者，則以原數同類者求等，約衍母，得數為借數。次置有零文庫零錢數，乘本用數，併為總數，滿衍母去之，不滿為諸庫日息足錢。各大小月日數乘之，各為實，各以原陌約為舊會。\end{unitzh}
\begin{uniten}Method: Find it by the great extension (dayan). Set up the market string of Treasury A and decrease it treasury by treasury, obtaining in each case the original strings (yuanmo) of the treasuries. Seek the common divisors in a linked chain (lianhuan qiudeng); reduce the odd, do not reduce the even [Siku editors: this is said only for the case where the common divisor (dengshu) is even; if the common divisor is odd, then reduce the even and do not reduce the odd], obtaining the definite moduli (dingmu). Multiply the definite moduli together to make the extension modulus. Reduce the extension modulus by each definite modulus, obtaining the extension numbers (yanshu). Where an extension number is the same as the extension modulus, remove it, making it "none" [Note: for those with none, borrow from the same kind]. For each of them, cast out whole multiples of the definite modulus; what remains is the excess number (qishu). With the excess number and the definite modulus, use the great extension to find the multiplying rate (chenglü); multiply the extension number by it to make the use-number (yongshu). For those with none, take the original numbers of the same kind, seek their common divisor, and reduce the extension modulus by it; the number obtained is the borrowed number (jieshu). Next set up the odd-cash amounts of the treasuries that have odd wen, multiply each by its own use-number, and add them to make the total number; cast out whole multiples of the extension modulus; what does not fill it is the full cash of daily interest of the treasuries. Multiply it by the numbers of days of the large and small months, each making a dividend; reduce each by its original string to make the old notes.\end{uniten}

\begin{unitzh}\unitlabel{shushu-005.9}草曰：置甲庫市陌一十二，遞减一，得一十一為乙庫陌，十一為丙庫陌，九為丁庫陌，八為戊庫陌，七為己庫陌，六為庚庫陌，得諸庫原陌。\end{unitzh}
\begin{uniten}Working: Set up the market string of Treasury A, 12; decreasing successively by 1, obtain 11 as the string of Treasury B, 11 as the string of Treasury C [Translator's note: by the successive decrease this should be 10], 9 as the string of Treasury D, 8 as the string of Treasury E, 7 as the string of Treasury F, and 6 as the string of Treasury G, obtaining the original strings of the treasuries.\end{uniten}

\begin{unitzh}\unitlabel{shushu-005.10}甲 乙 丙 丁 戊 己 庚\end{unitzh}
\begin{uniten}A B C D E F G\end{uniten}

\begin{unitzh}\unitlabel{shushu-005.11}以連環求等約訖，甲得一，乙得十一，丙得五，丁得九，戊得八，己得七，庚得一，各為定母。立各一為子。按：此法之要在於求定，而術中獨畧之，今詳其式於後。\end{unitzh}
\begin{uniten}Having reduced them by seeking common divisors in a linked chain, A obtains 1, B obtains 11, C obtains 5, D obtains 9, E obtains 8, F obtains 7, G obtains 1; each is a definite modulus. Set up 1 for each as the numerator (zi). Note: the key to this rule lies in finding the definite moduli, yet in the Method this alone is passed over briefly; now its pattern is set out in detail below.\end{uniten}

\begin{unitzh}\unitlabel{shushu-005.12}法：列七庫陌數於前。先以甲與乙相約，無等數；與丙數相約，得等數二{\zhsmall 偶}，約丙十得五{\zhsmall 竒}；與丁數相約，得等數三{\zhsmall 奇}，約甲十二得四{\zhsmall 偶}；與戊數相約，得等數四{\zhsmall 偶}，約甲四得一{\zhsmall 奇}。甲數既為一，不能𠕂約，即為與諸數徧約畢。次以乙與下五數相約，俱無等。次以丙與下四數相約，亦俱無等。次以丁與戊、己二數相約，俱無等；與庚數相約，得等數三{\zhsmall 奇}，約庚六得二{\zhsmall 偶}。次以戊與己相約，無等；與庚相約，得等數二{\zhsmall 偶}，約庚二得一{\zhsmall 奇}。庚既為一，己亦不能與之相約，乃為連環求等畢。得定數為甲一、乙十一、丙五、丁九、戊八、己七、庚一也。後凡求定數倣此。\end{unitzh}
\begin{uniten}Rule: arrange the string numbers of the seven treasuries in front. First reduce A and B mutually (xiangyue): there is no common divisor. With C's number, reducing mutually, the common divisor 2 is obtained [Note: even]; reducing C's 10 gives 5 [Note: odd]. With D's number, reducing mutually, the common divisor 3 is obtained [Note: odd]; reducing A's 12 gives 4 [Note: even]. With E's number, reducing mutually, the common divisor 4 is obtained [Note: even]; reducing A's 4 gives 1 [Note: odd]. A's number having become 1, it cannot be reduced again; this then completes its reduction with all the numbers in turn. Next reduce B mutually with the five numbers below: none has a common divisor. Next reduce C mutually with the four numbers below: likewise none has a common divisor. Next reduce D mutually with the two numbers E and F: neither has a common divisor; with G's number, reducing mutually, the common divisor 3 is obtained [Note: odd]; reducing G's 6 gives 2 [Note: even]. Next reduce E mutually with F: no common divisor; with G, reducing mutually, the common divisor 2 is obtained [Note: even]; reducing G's 2 gives 1 [Note: odd]. G having become 1, F likewise cannot be reduced with it; thus the seeking of common divisors in a linked chain is complete. The definite numbers obtained are A 1, B 11, C 5, D 9, E 8, F 7, G 1. Hereafter, whenever seeking definite numbers, imitate this.\end{uniten}

\begin{unitzh}\unitlabel{shushu-005.13}{\zhsmall 定母}\end{unitzh}
\begin{uniten}[Note: definite moduli]\end{uniten}

\begin{unitzh}\unitlabel{shushu-005.14}先以諸定相乘，得二萬七千七百二十為衍母。次以諸定互乘諸子，甲得二萬七千七百二十，乙得二千五百二十，丙得五千五百四十四，丁得三千八百，戊得三千四百六十五，己得三千九百六十，庚得二萬七千七百二十，各為衍數。\end{unitzh}
\begin{uniten}First multiply the definite moduli together, obtaining 27,720 as the extension modulus. Next cross-multiply (hucheng) the numerators by the definite moduli: A obtains 27,720, B obtains 2520, C obtains 5544, D obtains 3800 [Translator's note: 27,720 ÷ 9 = 3080], E obtains 3465, F obtains 3960, G obtains 27,720; each is an extension number.\end{uniten}

\begin{unitzh}\unitlabel{shushu-005.15}定母右行\end{unitzh}
\begin{uniten}Definite moduli, right column.\end{uniten}

\begin{unitzh}\unitlabel{shushu-005.16}次驗諸衍數，有同衍母者，皆去之，為無衍數。次各滿定母去各本衍，各得奇數：甲無，乙得一，丙得四，丁得二，戊得一，己得五，庚無，各為竒數。\end{unitzh}
\begin{uniten}Next examine the extension numbers: those that are the same as the extension modulus are all removed, becoming "no extension number". Next, in each case cast out whole multiples of the definite modulus from its own extension number, each obtaining an excess number: A none, B obtains 1, C obtains 4, D obtains 2, E obtains 1, F obtains 5, G none; each is an excess number.\end{uniten}

\begin{unitzh}\unitlabel{shushu-005.17}次驗有奇數者，得一，便以一為乘率；或得二數以上者，各以奇數於右上，定母於右下，立天元一於左上，用大衍求一之數入之，驗乘除至右上餘一而止，皆以左上所得為乘率。甲無，乙得一，丙得四，丁得五，戊得一，己得四，庚無，各為乘率。列右行，以對寄左衍數。\end{unitzh}
\begin{uniten}Next examine those that have excess numbers: where 1 is obtained, simply take 1 as the multiplying rate; where a number of 2 or more is obtained, in each case put the excess number at upper right and the definite modulus at lower right, set up the celestial unknown one (tianyuan yi) at upper left, enter them with the great-extension seeking-one (dayan qiuyi) reckoning, and check the multiplications and divisions until 1 remains at upper right, then stop; in every case take what is obtained at upper left as the multiplying rate. A none, B obtains 1, C obtains 4, D obtains 5, E obtains 1, F obtains 4 [Translator's note: the use-number 11,880 below is 3960 × 3], G none; each is a multiplying rate. Arrange them in the right column, facing the extension numbers lodged on the left.\end{uniten}

\begin{unitzh}\unitlabel{shushu-005.18}以兩行對乘之為用數：甲無，乙得二千五百二十，丙得二萬二千一百七十六，丁得一萬五千四百，戊得三千四百六十五，巳得一萬一千八百八十，庚無。\end{unitzh}
\begin{uniten}Multiply the two columns against each other (duicheng) to make the use-numbers: A none, B obtains 2520, C obtains 22,176, D obtains 15,400, E obtains 3465, F [Translator's note: the text has 巳 for 己] obtains 11,880, G none.\end{uniten}

\begin{unitzh}\unitlabel{shushu-005.19}次推無用數者，惟甲、庚，合於同類處借之。其同類，謂原陌列而視之。\end{unitzh}
\begin{uniten}Next examine those without use-numbers: only A and G; they should borrow from a place of the same kind. "The same kind" means arranging the original strings and looking at them.\end{uniten}

\begin{unitzh}\unitlabel{shushu-005.20}今視甲一十二、庚六，皆與丙一十、戊八俱偶，為同類。其戊用數三千四百六十五，其數少，不可借；惟丙一十之用數係二萬二千一百七十六，為最多，當以借之。乃以甲一十二、丙一十、庚六求等，得二。以等數二約衍母二萬七千七百二十，得一萬三千八百六十為借數。乃减丙用二萬二千一百七十六，餘八千三百一十六為丙用數。乃以所借岀之數一萬三千八百六十為實，以原等二為法除之，得六千九百三十為甲用數。以甲用數减借岀數，餘亦得六千九百三十為庚用數。今不欲甲、庚之借數同，乃驗得岀數一萬三千八百六十可用幾約如意，乃立三，取三分之一，得四千六百二十為甲用；取三分之二，得九千二百四十為庚用。列右行。\end{unitzh}
\begin{uniten}Now look at A 12 and G 6: together with C 10 and E 8 they are all even, and so are of the same kind. E's use-number is 3465; this number is small and cannot be borrowed from. Only the use-number of C 10, being 22,176, is the largest; one should borrow from it. Then seek the common divisor of A 12, C 10 and G 6, obtaining 2. Reduce the extension modulus 27,720 by the common divisor 2, obtaining 13,860 as the borrowed number. Then subtract it from C's use-number 22,176, leaving 8316 as C's use-number. Then take the number borrowed out, 13,860, as dividend and the original common divisor 2 as divisor, and divide, obtaining 6930 as A's use-number. Subtract A's use-number from the borrowed-out number; the remainder, likewise 6930, is G's use-number. Now, not wishing the borrowed numbers of A and G to be the same, one examines and finds that the out-number 13,860 can be reduced in several ways as desired; so set up 3: taking 1/3 gives 4620 as A's use-number; taking 2/3 gives 9240 as G's use-number. Arrange them in the right column.\end{uniten}

\begin{unitzh}\unitlabel{shushu-006.1}一 ○ ○ ○{\zhsmall 零數}左行\end{unitzh}
\begin{uniten}[Translator's note: remnant of a counting-rod diagram.] 1　0　0　0 [Note: leftover amounts (lingshu)] left column\end{uniten}

\begin{unitzh}\unitlabel{shushu-006.2}乃視諸庫有無零錢數，驗得乙丙己三庫無，先去其用數。乃以甲子戊庚四庫零錢列左行，對乘本用：甲得四萬六千二百，丁得六萬一千六百，戊得二萬七百九十，庚得三萬六千九百六十，合為總。\end{unitzh}
\begin{uniten}Then examine whether or not the treasuries have amounts of leftover cash (lingqian). On checking, the three treasuries B, C and F have none, so first discard their use numbers (yongshu). Then list the leftover cash of the four treasuries A, zi [Translator's note: zi 子 is evidently an error for ding 丁, D.], E and G in the left column, and multiply each by its own use number: A gives 46200, D gives 61600, E gives 20790, G gives 36960; together these are the totals (zong).\end{uniten}

\begin{unitzh}\unitlabel{shushu-006.3}併此四總，得一十六萬五千五百五十，滿衍母二萬七千七百二十去之，不滿二萬六千九百五十，為所求率。以貫約為二十六貫九千五十文，為諸庫日息等數。以官省七十七陌展，得三十五貫文{\zhsmall 按官省陌以七十七為一百，故二十六貫餘展為三十五貫}，各以其庫陌紐計{\zhsmall 按庫陌紐計，即以各陌數為一百}，各得舊會零錢。各以三十日乘，為大月息；以日息减大月息，餘為小月息。合問。\end{unitzh}
\begin{uniten}Add these four totals to get 165550. Remove from it the full multiples of the extended mother (yanmu), 27720; the remainder (buman), 26950, is the rate (lü) sought. Reduced (yue) by the guan (貫), this makes 26 guan 9050 wen (文) [Translator's note: sic; 26950 wen is 26 guan 950 wen.], which is the equal amount (dengshu) of the daily interest of all the treasuries. Expanding (zhan) by the official short (guansheng) mo (陌) of 77 gives 35 guan wen [Siku editors: The official short mo takes 77 as 100; therefore 26 guan and more is expanded to 35 guan.]. Each, reckoned in strings (niuji) by its own treasury's mo [Siku editors: Reckoning in strings by the treasury mo means taking each mo number as 100.], yields the leftover cash of the old hui (jiuhui) [Translator's note: hui, probably huizi paper money.]. Multiply each by 30 days to make the interest of a large month; subtract the daily interest from the large-month interest, and the remainder is the interest of a small month. This agrees with the question.\end{uniten}

\begin{unitzh}\unitlabel{shushu-006.4}分糶推原\end{unitzh}
\begin{uniten}Divided Grain Sales: Tracing Back the Original\end{uniten}

\begin{unitzh}\unitlabel{shushu-006.5}問：有上農三人，力田所收之米，係用足斗均分，各徃他處岀糶。甲糶本郡官塲，餘三斗二升；乙糶與吉安鄉民，餘七斗；丙糶與平江𭣄戸，餘三斗。欲知共米及三人所分、各糶石幾何？荅曰：共米七百三十八石；三人各分米二百四十六石；甲糶官斛二百九十六石，乙糶安吉斛二百二十三石，丙糶平江斛一百八十二石。\end{unitzh}
\begin{uniten}Question: There are three superior farmers. The rice harvested from their labour in the fields was divided equally using the full dou (斗), and each went to another place to sell it. A sold at the official market (guanchang) of his home commandery, with 3 dou 2 sheng (升) left over; B sold to the country people of Ji'an, with 7 dou left over; C sold to the contracting households (lanhu) of Pingjiang, with 3 dou left over. One wishes to know how many dan (石) there are in the total rice, in the shares of the three men, and in what each sold. Answer: The total rice is 738 dan; each of the three men's share of rice is 246 dan. A sold 296 dan by the official hu (斛); B sold 223 dan by the Anji hu; C sold 182 dan by the Pingjiang hu.\end{uniten}

\begin{unitzh}\unitlabel{shushu-006.6}術曰：以大衍求之。置官塲斛率、安吉鄉斛率、平江市斛率{\zhsmall 官私共知者：官斛八斗三升，安吉鄉斛一斗一升，平江市斛一石三斗五升}為原數，求總等，不約一位，約衆位；連環求等，約竒不約偶；或猶有類數存者，有求等約彼，為復乘此。各等定母相乘為衍母，互乘為衍數，滿定去之，得竒；大衍求一，得乘率，乘衍數為用數。以各餘米乘用，併之為總，滿衍母去之，不滿為所分；以原人數乘之，為共米。\end{unitzh}
\begin{uniten}Method: Seek it by the great extension (dayan). Set up the rate of the official-market hu, the rate of the Anji country hu and the rate of the Pingjiang market hu [Note: Known jointly to officials and private persons: the official hu is 8 dou 3 sheng, the Anji country hu 1 dou 1 sheng, the Pingjiang market hu 1 dan 3 dou 5 sheng.] as the original numbers (yuanshu). Seek the total common measure (zongdeng): do not reduce one position, reduce the many positions. Seek the common measures in linked rings (lianhuan qiudeng): reduce the odd, do not reduce the even. If numbers of a kind (leishu) still remain, again seek the common measure, reduce that one, and multiply this one back. Multiply together the fixed mothers (dingmu) of the respective common measures to make the extended mother; multiply them mutually (hucheng) to make the extended numbers (yanshu); remove from these the full multiples of the fixed mothers to get the odds (qi). By the great extension seeking one (dayan qiuyi) obtain the multiplying rates (chenglü), and multiply the extended numbers by them to make the use numbers. Multiply the use numbers by the respective leftover rice and add them to make the total; remove from it the full multiples of the extended mother; the remainder is the share. Multiply it by the original number of people to make the total rice.\end{uniten}

\begin{unitzh}\unitlabel{shushu-006.7}草曰：置文思院官斛八十三升，吉安州鄉一百一十升，平江府市斛一百三十升，各為其斛原率。\end{unitzh}
\begin{uniten}Working: Set up the official hu of the Wensi Yuan (文思院), 83 sheng; Ji'an prefecture country, 110 sheng; the market hu of Pingjiang prefecture, 130 sheng; each is the original rate (yuanlü) of its hu.\end{uniten}

\begin{unitzh}\unitlabel{shushu-006.8}先以三率求總等，得一，不約{\zhsmall 按此題只一數見十，不必用復數求總等}。次以連環求等：其安吉率一百一十與平江率一百三十五求等，得五，以約平江率，得二十七{\zhsmall 按五為中數，或約偶或約奇皆可，但不約可以𠕂約者}；餘皆求等，得一，不約，各得原數。\end{unitzh}
\begin{uniten}First seek the total common measure of the three rates: it is 1; do not reduce [Siku editors: In this problem only one number shows a ten; there is no need to use the repeated numbers (fushu) to seek the total common measure.]. Next seek the common measures in linked rings: seeking the common measure of the Anji rate 110 and the Pingjiang rate 135 gives 5; reducing the Pingjiang rate by it gives 27 [Siku editors: 5 is a middle number (zhongshu); reducing either the even or the odd is acceptable, but one does not reduce what can be reduced again.]; for all the rest, seeking the common measure gives 1; do not reduce; each obtains its original number.\end{uniten}

\begin{unitzh}\unitlabel{shushu-006.9}以定相乘，得二十四萬六千五百一十，為衍母。各以原率約之，得二千九百七十，為官斛衍數；得二千二百四十一，為安吉州衍數；得九千一百三十，為平江斛衍數。\end{unitzh}
\begin{uniten}Multiply the fixed ones together to get 246510 as the extended mother. Reduce it by each original rate: obtaining 2970 as the extended number of the official hu; obtaining 2241 as the extended number of Anji prefecture; obtaining 9130 as the extended number of the Pingjiang hu.\end{uniten}

\begin{unitzh}\unitlabel{shushu-006.10}官斛 安吉 平江 衍母\end{unitzh}
\begin{uniten}Official hu　Anji　Pingjiang　Extended mother\end{uniten}

\begin{unitzh}\unitlabel{shushu-006.11}次以定母滿法去衍數，得不滿六十五，為官斛奇；不滿四十一，為安吉奇；不滿四，為平江奇數。\end{unitzh}
\begin{uniten}Next, with the fixed mothers as divisors (fa), remove their full multiples from the extended numbers: obtaining the remainder 65 as the odd of the official hu; the remainder 41 as the odd of Anji; the remainder 4 as the odd number (qishu) of Pingjiang.\end{uniten}

\begin{unitzh}\unitlabel{shushu-006.12}定母奇數各以大衍入之，求得乘數：得二十三，為官斛乘率；得五十一，為安吉乘率；得七，為平江乘率。\end{unitzh}
\begin{uniten}Enter each fixed mother and odd number into the great extension and seek the multiplying numbers (chengshu): obtaining 23 as the multiplying rate of the official hu; obtaining 51 as the multiplying rate of Anji; obtaining 7 as the multiplying rate of Pingjiang.\end{uniten}

\begin{unitzh}\unitlabel{shushu-006.13}以乘率各乘寄左行衍數，得六萬八千三百一十，為官斛用數；得一十一萬四千二百九十一，為安吉用數；得六萬三千九百一十，為平江用數。\end{unitzh}
\begin{uniten}With the multiplying rates multiply the respective extended numbers lodged in the left column: obtaining 68310 as the use number of the official hu; obtaining 114291 as the use number of Anji; obtaining 63910 as the use number of Pingjiang.\end{uniten}

\begin{unitzh}\unitlabel{shushu-006.14}次以甲餘三十二升乘官斛用數六萬八千三百一十，得二百一十八萬五千九百二十升，於上；次以乙餘七十升乘安吉用數一十一萬四千二百九十一，得八百萬三百七十升，於中；次以丙餘三十升乘平江用數六萬三千九百一十，得一百一十九萬七千三百，於下。各為總，併之，得一千二百一十萬三千五百九十升，為總數。\end{unitzh}
\begin{uniten}Next, with A's leftover of 32 sheng multiply the use number of the official hu, 68310, to get 2185920 sheng, placed at the top; next, with B's leftover of 70 sheng multiply the Anji use number, 114291, to get 8000370 sheng, placed in the middle; next, with C's leftover of 30 sheng multiply the Pingjiang use number, 63910, to get 1197300 [Translator's note: sic; 63910 × 30 = 1917300.], placed at the bottom. Each is a total; add them to get 12103590 sheng as the total number (zongshu).\end{uniten}

\begin{unitzh}\unitlabel{shushu-006.15}滿衍母二十四萬六千五百一十升去之，不滿二萬四千六百升，為所求率，展為二百四十六石，為三人各分米。以兄弟三人因之，得七百三十八石，為共米。置分米二百四十六石，各以官斛八斗三升、安吉斛一石一斗、平江斛一石三斗五升約之：甲得二百九十六石，餘三斗二升；乙得二百二十三石，餘七斗；丙得一百八十二石，餘三斗。各為糶過及餘米。合問。\end{unitzh}
\begin{uniten}Remove from it the full multiples of the extended mother, 246510 sheng; the remainder, 24600 sheng, is the rate sought. Expanded, it is 246 dan, which is the rice each of the three men received as his share. Multiply (yin) it by the three brothers to get 738 dan as the total rice. Set up the share of rice, 246 dan, and reduce it respectively by the official hu of 8 dou 3 sheng, the Anji hu of 1 dan 1 dou and the Pingjiang hu of 1 dan 3 dou 5 sheng: A gets 296 dan, with 3 dou 2 sheng left over; B gets 223 dan, with 7 dou left over; C gets 182 dan, with 3 dou left over. These are respectively what was sold and the leftover rice. This agrees with the question.\end{uniten}

\begin{unitzh}\unitlabel{shushu-006.16}數學九章卷一上\end{unitzh}
\begin{uniten}Shuxue jiuzhang (Mathematical Treatise in Nine Sections), juan 1, part 1\end{uniten}

\begin{unitzh}\unitlabel{shushu-006.17}子部，天文算法類，算書之屬，數學九章\end{unitzh}
\begin{uniten}Masters (zi) division; Astronomy and Mathematics category; Mathematical Books subcategory; Shuxue jiuzhang\end{uniten}

\section*{卷1下}\addcontentsline{toc}{section}{卷1下}

\begin{unitzh}\unitlabel{shushu-007.1}欽定四庫全書\end{unitzh}
\begin{uniten}Imperially Commissioned Complete Library of the Four Treasuries (Qinding Siku quanshu)\end{uniten}

\begin{unitzh}\unitlabel{shushu-007.2}數學九章卷一下 宋 秦九韶 撰 大衍\end{unitzh}
\begin{uniten}Mathematics in Nine Sections (Shuxue jiuzhang), juan 1, second part. Written by Qin Jiushao of the Song. Great Extension (dayan).\end{uniten}

\begin{unitzh}\unitlabel{shushu-007.3}積足尋原\end{unitzh}
\begin{uniten}Accumulating to Sufficiency and Seeking the Origin\end{uniten}

\begin{unitzh}\unitlabel{shushu-007.4}問：欲砌基一叚，見管大小方甎、六門、城甎四色。令匠取便，或平或側，只用一色甎砌，湏要適足。匠以甎量地，計料稱：用大方料，廣多六寸，深少六寸{\zhsmall 按：即多七寸。}；用小方，廣多二寸，深少三寸{\zhsmall 按：即多八寸。}；用城甎長，廣多三寸，深少一寸{\zhsmall 按：即多一尺一寸。}；以闊，深少一寸{\zhsmall 按：即多五寸。}，廣多三寸；以厚，廣多五分，深多一寸；用六門甎長，廣三寸，深多一寸；以闊，廣多三寸，深多一寸；用厚，廣多一寸，深多一寸。皆不匼匝，未免修破甎料禆補。其四色甎：大方方一尺三寸，小方方一尺一寸，城甎長一尺二寸，闊六寸，厚二寸五分；六門長一尺，闊五寸，厚二寸。欲知基深、廣㡬何？\end{unitzh}
\begin{uniten}Question: One wishes to lay one section of foundation in brick. On hand are four kinds: large and small square bricks, liumen ("six-gate") [bricks], and city-wall bricks. The mason is told to do as is convenient, [laying them] either flat or on edge, using only one kind of brick for the masonry; it must fit exactly. The mason measured the ground with the bricks, reckoned the materials, and reported: using large-square material, the breadth is over by 6 cun (寸) and the depth short by 6 cun [Siku editors: That is, over by 7 cun.]; using small squares, the breadth is over by 2 cun and the depth short by 3 cun [Siku editors: That is, over by 8 cun.]; using the length of the city-wall brick, the breadth is over by 3 cun and the depth short by 1 cun [Siku editors: That is, over by 1 chi (尺) 1 cun.]; with the width, the depth is short by 1 cun [Siku editors: That is, over by 5 cun.] and the breadth over by 3 cun; with the thickness, the breadth is over by 5 fen (分) and the depth over by 1 cun; using the length of the liumen brick, the breadth 3 cun [Translator's note: 多 "over by" appears to be missing here.] and the depth over by 1 cun; with the width, the breadth is over by 3 cun and the depth over by 1 cun; using the thickness, the breadth is over by 1 cun and the depth over by 1 cun. None of these closes up exactly, and trimming and breaking brick material to patch [the gaps] cannot be avoided. Of the four kinds of brick: the large square is 1 chi 3 cun square; the small square is 1 chi 1 cun square; the city-wall brick is 1 chi 2 cun long, 6 cun wide and 2 cun 5 fen thick; the liumen [brick] is 1 chi long, 5 cun wide and 2 cun thick. One wishes to know how much the depth and breadth of the foundation are.\end{uniten}

\begin{unitzh}\unitlabel{shushu-007.5}按：題意謂以一尺三寸量基之廣，未餘六寸；以一尺一寸量之，餘二寸；以一尺二寸量之，餘三寸；以六寸量之，亦餘三寸；以二寸五分量之，餘五分；以一尺量之，餘三寸；以五寸量之，亦餘三寸；以二寸量之，餘一寸：以求廣也。其求深之意亦同。\end{unitzh}
\begin{uniten}Note: The meaning of the problem is this: measuring the breadth of the foundation with 1 chi 3 cun, 6 cun remain [Translator's note: the text reads 未餘 "not yet remain"; 未 is presumably an error for 末 "at the end".]; measuring it with 1 chi 1 cun, 2 cun remain; measuring it with 1 chi 2 cun, 3 cun remain; measuring it with 6 cun, 3 cun likewise remain; measuring it with 2 cun 5 fen, 5 fen remain; measuring it with 1 chi, 3 cun remain; measuring it with 5 cun, 3 cun likewise remain; measuring it with 2 cun, 1 cun remains: [from this] the breadth is to be found. The sense of finding the depth is the same.\end{uniten}

\begin{unitzh}\unitlabel{shushu-007.6}答曰：深三丈七尺一寸，廣一丈二尺三寸。\end{unitzh}
\begin{uniten}Answer: Depth 3 zhang (丈) 7 chi 1 cun; breadth 1 zhang 2 chi 3 cun.\end{uniten}

\begin{unitzh}\unitlabel{shushu-007.7}術曰：以大衍求之。置甎方、長、闊、厚為元數，以小者為單，起一。先求總等，存一位，約衆位{\zhsmall 列為多者，隨意立號。}，乃為元數。連環求等，約為定母。以定相乗為衍母。各定約衍母，得衍數。滿定去之，得竒。竒、定大衍，得乗率。以乗衍數，得用數。次置廣深多少數，多者乘用；少者減元數，餘以乘用。併為總，滿衍母去之，不滿得廣深。\end{unitzh}
\begin{uniten}Method: Solve it by the great extension. Set out the sides, lengths, widths and thicknesses of the bricks as the original numbers (yuanshu); take the small [unit] as the unit (dan), starting from one. First seek the overall common divisor (zongdeng); keep one place and reduce all the [other] places [Note: When the [items] set out are many, set up labels at will.]; these then are the original numbers. Seek the common divisors in a linked chain (lianhuan qiudeng) and reduce, to make the fixed denominators (dingmu). Multiply the fixed [denominators] together to make the extension denominator (yanmu). Reduce the extension denominator by each fixed [denominator] to get the extension numbers (yanshu). Remove full [multiples of] the fixed [denominator] from them to get the odd [remainders] (ji). [Apply] the great extension to the odd [remainder] and the fixed [denominator] to get the multiplying rate (chenglü). Multiply the extension number by it to get the use number (yongshu). Next set out the numbers by which breadth and depth are over or short: those that are over multiply the use [numbers]; those that are short are subtracted from the original numbers, and the remainders multiply the use [numbers]. Add them together to make the total; remove full [multiples of] the extension denominator from it; the unfilled [remainder] gives the breadth or depth.\end{uniten}

\begin{unitzh}\unitlabel{shushu-007.8}草曰：置四甎方、長、闊、厚，係八數。城甎厚有分，為小者，皆通之為單。大方得一百三十分，小方得一百一十分，城甎長得一百二十分，闊得六十分，厚得二十五分；六門甎長得一百分，闊得五十分，厚得二十分。\end{unitzh}
\begin{uniten}Working: Set out the sides, lengths, widths and thicknesses of the four bricks; these are 8 numbers. The thickness of the city-wall brick has fen; as this is the small [unit], convert everything to it as the unit. The large square gives 130 fen; the small square gives 110 fen; the city-wall brick length gives 120 fen, the width gives 60 fen, the thickness gives 25 fen; the liumen brick length gives 100 fen, the width gives 50 fen, the thickness gives 20 fen.\end{uniten}

\begin{unitzh}\unitlabel{shushu-007.9}錐行置之右。列位稍多，甎名相互，今假八音為號位。先以最少者自木二十與革二十五求等，得五，乃反約木二十為四。木四與土五十求等，得二，以約五十為二十五。木四與匏六十求等，得四，約六十為一十五。木四與竹一百求等，得四，約一百為二十五。木四與絲一百一十求等，得二，約一百一十為五十五。木四與石一百二十求等，得四，反約木四為一。以木一與金求等，得一，不約。為木與諸數求等約訖，為一變。得數具圖如後。\end{unitzh}
\begin{uniten}Set them out on the right in awl formation (zhuixing). The places set out are rather many and the brick names are intermingled, so now we borrow the eight sounds (bayin) as labels for the places [Translator's note: the eight classes of musical instruments: metal, stone, silk, bamboo, gourd, earth, hide, wood.]. First, taking the smallest, starting from wood 20, with hide 25 seek the common divisor (deng), getting 5; then, conversely, reduce wood 20 to 4. Wood 4 with earth 50: seek the common divisor, getting 2; with it reduce 50 to 25. Wood 4 with gourd 60: seek the common divisor, getting 4; reduce 60 to 15. Wood 4 with bamboo 100: seek the common divisor, getting 4; reduce 100 to 25. Wood 4 with silk 110: seek the common divisor, getting 2; reduce 110 to 55. Wood 4 with stone 120: seek the common divisor, getting 4; conversely reduce wood 4 to 1. Take wood 1 with metal and seek the common divisor, getting 1; do not reduce. When wood has been tested for common divisors with all the numbers and the reductions are finished, this makes the first transformation (bian). The resulting numbers are fully shown in the figure below.\end{uniten}

\begin{unitzh}\unitlabel{shushu-007.10}次以革二十五與土五十{\zhsmall 按：前巳約土為二十五，次變不應復用原數，然於得數却無碍。}求等，得二十五，約五十為二。以革二十五與匏一十五求等，得五，約匏一十五為三。以革二十五與竹二十五求等，得二十五，約竹二十五為一。又以革二十五與絲五十五求等，得五，約絲五十五得一十一。以革二十五與石一百二十求等，得五，約一百二十為二十四。以革二十五與金一百三十求等，得五，約金一百三十得二十六。革與諸數徧約訖，為二變。具圖如後。\end{unitzh}
\begin{uniten}Next take hide 25 with earth 50 [Siku editors: Earth was already reduced to 25 earlier; in the next transformation the original number should not be used again, but it does no harm to the resulting numbers.] and seek the common divisor, getting 25; reduce 50 to 2. Take hide 25 with gourd 15 and seek the common divisor, getting 5; reduce gourd 15 to 3. Take hide 25 with bamboo 25 and seek the common divisor, getting 25; reduce bamboo 25 to 1. Again take hide 25 with silk 55 and seek the common divisor, getting 5; reducing silk 55 gives 11. Take hide 25 with stone 120 and seek the common divisor, getting 5; reduce 120 to 24. Take hide 25 with metal 130 and seek the common divisor, getting 5; reducing metal 130 gives 26. Hide having been reduced against all the numbers throughout, this makes the second transformation. The figure is given below.\end{uniten}

\begin{unitzh}\unitlabel{shushu-007.11}乃以土二與匏三、竹一、絲一十一求等，皆得一，不約。以土二與石二十四求等，得二，及約土二得一。又不土一與金二十六求等，得一，不約。土與諸數約訖，為三變。具圖如後。\end{unitzh}
\begin{uniten}Then take earth 2 with gourd 3, bamboo 1 and silk 11 and seek the common divisors: all give 1; do not reduce. Take earth 2 with stone 24 and seek the common divisor, getting 2; and reduce earth 2, getting 1. Again "not" earth 1 with metal 26 [Translator's note: 不 "not" is evidently corrupt here, presumably for 以 "take".]: seek the common divisor, getting 1; do not reduce. Earth having been reduced against all the numbers, this makes the third transformation. The figure is given below.\end{uniten}

\begin{unitzh}\unitlabel{shushu-007.12}乃以匏三與竹一、絲十十一求等，皆得一。又以匏三與石二十四求等，得三，約石二十四為八。又匏三與金二十六求等，得一，不約。匏與諸數約訖，以為四變。\end{unitzh}
\begin{uniten}Then take gourd 3 with bamboo 1 and silk "十十一" [Translator's note: evidently for 一十一, 11] and seek the common divisors: all give 1. Again take gourd 3 with stone 24 and seek the common divisor, getting 3; reduce stone 24 to 8. Again, gourd 3 with metal 26: seek the common divisor, getting 1; do not reduce. Gourd having been reduced against all the numbers, this is taken as the fourth transformation.\end{uniten}

\begin{unitzh}\unitlabel{shushu-007.13}次以竹一與絲一十一、與石二十四{\zhsmall 按：巳約為八，云二十四，誤。}、與金二十六求等，皆得一。竹與諸數約訖，為五變。次以絲一十一與石二十四{\zhsmall 按：誤同上。}、與金二十六求等，皆得一，不約，為六變。復以石二{\zhsmall 十}四{\zhsmall 按：誤同上。}與金二十六求等，得二，約金二十六為一十三。至此七變，連環求等約俱畢，得數為定母。列圖如後。\end{unitzh}
\begin{uniten}Next take bamboo 1 with silk 11, with stone 24 [Siku editors: It has already been reduced to 8; saying 24 is an error.], and with metal 26, and seek the common divisors: all give 1. Bamboo having been reduced against all the numbers, this makes the fifth transformation. Next take silk 11 with stone 24 [Siku editors: The error is the same as above.] and with metal 26, and seek the common divisors: all give 1; do not reduce; this makes the sixth transformation. Once more take stone 24 [Translator's note: the 十 of 二十四 is in small print in the transcription.] [Siku editors: The error is the same as above.] with metal 26 and seek the common divisor, getting 2; reduce metal 26 to 13. At this point, after seven transformations, the linked-chain seeking of common divisors and the reductions are all completed, and the resulting numbers are the fixed denominators. The figure is set out below.\end{uniten}

\begin{unitzh}\unitlabel{shushu-007.14}右定母，列右行，以相乗，得八萬五千八百為衍母。以各定母約衍母，各得衍數。其竹、木、土定得一者為無。\end{unitzh}
\begin{uniten}The above fixed denominators are set in the right column and multiplied together, giving 85800 as the extension denominator. Reduce the extension denominator by each fixed denominator; each gives an extension number. Those among them, bamboo, wood and earth, whose fixed [denominators] come out as 1 are [counted as] none.\end{uniten}

\begin{unitzh}\unitlabel{shushu-007.15}金定一十三，得衍數六千六百；石定八，得衍數一萬七百二十五；絲定一十一，得衍數七千八百；竹定一，無衍；匏定三，得衍數二萬八千六百；土定一，無衍數；革定五五，得衍數三千四百三十二；木定一，無衍數。各滿定母去之，得竒數。\end{unitzh}
\begin{uniten}Metal: fixed [denominator] 13, giving extension number 6600. Stone: fixed 8, giving extension number 10725. Silk: fixed 11, giving extension number 7800. Bamboo: fixed 1, no extension. Gourd: fixed 3, giving extension number 28600. Earth: fixed 1, no extension number. Hide: fixed "五五" [Translator's note: evidently for 二十五, 25], giving extension number 3432. Wood: fixed 1, no extension number. From each remove full [multiples of] the fixed denominator, getting the odd numbers (jishu).\end{uniten}

\begin{unitzh}\unitlabel{shushu-007.16}金得竒九，石得竒四，絲得竒一，匏得竒一，草得竒七。其絲、匏得竒數一者，便以一為乗率。其金、石、革三處竒數，皆與本定母用大衍求一入之，各得乘率，列右行。\end{unitzh}
\begin{uniten}Metal gets odd 9; stone gets odd 4; silk gets odd 1; gourd gets odd 1; cao gets odd 7 [Translator's note: 草 is presumably an error for 革 "hide".]. Silk and gourd, which get the odd number 1, simply take 1 as the multiplying rate. The odd numbers in the three places metal, stone and hide are each entered together with their own fixed denominator into the great extension [method of] seeking one (dayan qiu yi); each gets its multiplying rate; set them in the right column.\end{uniten}

\begin{unitzh}\unitlabel{shushu-007.17}金得三，石得五，絲得一，革得一，革得一十八，各為乘率，寄左行衍數，各得為用數。\end{unitzh}
\begin{uniten}Metal gets 3; stone gets 5; silk gets 1; hide gets 1 [Translator's note: presumably an error for 匏 "gourd".]; hide gets 18. Each is a multiplying rate; lodge them against the extension numbers in the left column; what each gives is the use number.\end{uniten}

\begin{unitzh}\unitlabel{shushu-007.18}凡諸用數同類者，類必多，可互借以補無者。先驗革元數二十五與木元數二十為同類，求等得五，以等五約衍母八萬五千八百，得一萬七千一百六十，乃於革用數内減出，以補木位，為木用，餘四萬四千六百一十六為革用。次驗竹元數一百與土五十為同類，以求等得五十，以等五十約衍母八萬五千八百，得一千七百一十六，亦於革用内各借與竹、土為用數，革止餘四萬一千一百八十四為用，得諸定用數。\end{unitzh}
\begin{uniten}In general, where use numbers are of the same kind (tonglei), the kind is bound to be abundant, and they can lend to one another to supply those that have none. First check: the original number of hide, 25, and the original number of wood, 20, are of the same kind; seek the common divisor, getting 5. With the common divisor 5 reduce the extension denominator 85800, getting 17160; then subtract it out from within the use number of hide to supply the place of wood, as the use [number] of wood; the remaining 44616 is the use [number] of hide. Next check: the original number of bamboo, 100, and earth, 50, are of the same kind; seek the common divisor, getting 50. With the common divisor 50 reduce the extension denominator 85800, getting 1716; likewise, from within the use [number] of hide, lend [this] to each of bamboo and earth as their use numbers. Hide is left with only 41184 as its use [number]. [Thus] all the fixed use numbers (ding yongshu) are obtained.\end{uniten}

\begin{unitzh}\unitlabel{shushu-007.19}按：無用數則此條可省，借數轉生煩擾，非法也。其所以可用借補者，蓋以同類之元數，其較餘之竒偶必同，故一數可分用也。然惟元數同偶者為然，同竒則有不可用者。此題可用，因題中餘數未過小元數也。\end{unitzh}
\begin{uniten}Note: Where there is no use number, this item can be omitted; borrowing numbers instead gives rise to trouble and confusion, and is not the [proper] method. The reason why borrowing to supply can be used is presumably that, for original numbers of the same kind, the oddness or evenness of their remainders, compared, must be the same; therefore one number can be divided up for use. However, this holds only when the original numbers are both even; when they are both odd, there are cases where it cannot be used. In this problem it can be used because the remainders in the problem do not exceed the smaller original numbers.\end{uniten}

\begin{unitzh}\unitlabel{shushu-007.20}右行定用，始列錐行，假號求得，今照甎色遷次列之。\end{unitzh}
\begin{uniten}The fixed use [numbers] in the right column were at first set out in awl formation and obtained by means of borrowed labels; now, following the kinds of brick, shift their order and set them out.\end{uniten}

\begin{unitzh}\unitlabel{shushu-007.21}旣照甎次序列用數於右行，乃驗問題所謂大方甎砌廣多六寸，小方多二寸，城甎長多三寸，城甎闊多三寸，厚多五分，六門長多三寸，闊多三寸，厚多一寸，對本用列左行，各對乘之。具圖如後。\end{unitzh}
\begin{uniten}Having set out the use numbers in the right column following the order of the bricks, then check what the problem says: laying with large square bricks, the breadth is over by 6 cun; small squares, over by 2 cun; city-wall brick length, over by 3 cun; city-wall brick width, over by 3 cun; thickness, over by 5 fen; liumen length, over by 3 cun; width, over by 3 cun; thickness, over by 1 cun. Set these in the left column opposite their own use [numbers], and multiply each pair together. The figure is given below.\end{uniten}

\begin{unitzh}\unitlabel{shushu-007.22}兩行乗畢，金得一百一十八萬八千，絲得一十五萬六千，石得一百六十萬八千七百五十，匏得八十五萬八千，革得二十萬五千九百二十，竹得五萬一十四百八十，土亦得五萬一千四百八十，木得一十七萬一千六百。乃併前八位數，共得四百二十九萬一千二百三十分，為總。滿衍母八萬五千八百去之，不滿一千二百三十分，約之為一丈二尺三寸，為基元廣數。\end{unitzh}
\begin{uniten}When the multiplication of the two columns is finished: metal gets 1188000; silk gets 156000; stone gets 1608750; gourd gets 858000; hide gets 205920; bamboo gets "五萬一十四百八十" [Translator's note: evidently for 五萬一千四百八十, 51480]; earth likewise gets 51480; wood gets 171600. Then add the numbers of the preceding 8 places, getting in all 4291230 fen, as the total. Remove full [multiples of] the extension denominator 85800; the unfilled [remainder] is 1230 fen; reduced, it is 1 zhang 2 chi 3 cun, which is the original breadth of the foundation.\end{uniten}

\begin{unitzh}\unitlabel{shushu-007.23}乃求其深。驗問題：大方砌少六寸，小方砌少三寸，城甎長砌少一寸，闊砌少一寸，厚砌多一寸，六門長砌多一寸，六門闊砌多一寸，六門厚砌多一寸，列為中行。次置諸甎元數，列為左行，課減之。具圖如後。\end{unitzh}
\begin{uniten}Then seek its depth. Check the problem: laid with large squares, short by 6 cun; laid with small squares, short by 3 cun; laid with city-wall brick lengths, short by 1 cun; laid with widths, short by 1 cun; laid with thicknesses, over by 1 cun; laid with liumen lengths, over by 1 cun; laid with liumen widths, over by 1 cun; laid with liumen thicknesses, over by 1 cun. Set these out as the middle column. Next set out the original numbers of the bricks as the left column, and compare and subtract them (kejian). The figure is given below.\end{uniten}

\begin{unitzh}\unitlabel{shushu-007.24}今以中行多者存之，少者用減左行。存者，左行元數去之；所減者，左行餘數存之。金得七十，絲得八十，石得一百一十，匏得五十，革得一十，竹一十，土一十，本一十。具圖如後。\end{unitzh}
\begin{uniten}Now, in the middle column, keep those that are over; those that are short are used to subtract from the left column. For those kept, the original numbers in the left column are removed; for those subtracted, the remainders in the left column are kept. Metal gets 70; silk gets 80; stone gets 110; gourd gets 50; hide gets 10; bamboo 10; earth 10; ben 10 [Translator's note: 本 is presumably an error for 木 "wood".]. The figure is given below.\end{uniten}

\begin{unitzh}\unitlabel{shushu-007.25}列為左行，以對右行定用數。具圖如後。\end{unitzh}
\begin{uniten}Set these out as the left column, opposite the fixed use numbers of the right column. The figure is given below.\end{uniten}

\begin{unitzh}\unitlabel{shushu-007.26}以左行多餘數對乘右行用數：金得一百三十八萬六千，絲得六十二萬四千，石得五百八十九萬八千七百五十，匏得一百四十三萬，革得四十一萬一千八百四十，竹得一萬七千一百六十，土得一萬七千一百六十，木得一十七萬一千六百。具圖如後。\end{unitzh}
\begin{uniten}With the excess and remainder numbers of the left column, multiply correspondingly the use numbers of the right column: metal gets 1386000; silk gets 624000; stone gets 5898750; gourd gets 1430000; hide gets 411840; bamboo gets 17160; earth gets 17160; wood gets 171600. The figure is given below.\end{uniten}

\begin{unitzh}\unitlabel{shushu-007.27}併八位，得九百九十五萬六千五百一十分，為總。滿衍母八萬五千八百去之，不滿三千七百一十分，展為三丈七尺一寸，為基地深。\end{unitzh}
\begin{uniten}Add the 8 places, getting 9956510 fen, as the total. Remove full [multiples of] the extension denominator 85800; the unfilled [remainder] is 3710 fen, which, expanded, is 3 zhang 7 chi 1 cun: this is the depth of the foundation ground.\end{uniten}

\begin{unitzh}\unitlabel{shushu-007.28}推計土功\end{unitzh}
\begin{uniten}Reckoning the Earthwork Labor\end{uniten}

\begin{unitzh}\unitlabel{shushu-008.1}問：築堤，起四縣夫，分給里歩皆同。齊闊二丈，里法三百六十歩，歩法五尺八寸。人夫以物力差定：甲縣物力一十三萬八千六百貫，乙縣物力一十四萬六千三百貫，丙縣物力一十九萬二千五百貫，丁縣物力一十八萬四千八百貫。每力七百七十貫科一名，春程人功平方六十尺。先到縣先給。今甲、乙二縣俱畢，丙縣餘五十一丈，丁縣餘一十八丈，不及一日全功。欲知堤長及四縣夫所築各㡬何？\end{unitzh}
\begin{uniten}Question: To build a dike, laborers are raised from 4 counties, and the li (里) and bu (歩) allotted [to each] are all the same. The uniform width is 2 zhang (丈); the li rule (lifa) is 360 bu; the bu rule (bufa) is 5 chi (尺) 8 cun (寸). The laborers are fixed in graded proportion to property assets (wuli): County A's property assets are 138,600 guan (貫); County B's property assets are 146,300 guan; County C's property assets are 192,500 guan; County D's property assets are 184,800 guan. For every 770 guan of assets 1 man is levied; the spring quota (chuncheng) of a man's work is 60 square chi. The county that arrives first is allotted first. Now the two counties A and B have both finished; County C has 51 zhang remaining and County D has 18 zhang remaining, falling short of one full day's work. One wishes to know the length of the dike, and how much the laborers of the 4 counties each built?\end{uniten}

\begin{unitzh}\unitlabel{shushu-008.2}按：題意以四縣修堤總長相同，毎日所修之長不同，以各每日所修之長計總長，或適足，或有餘，以求總長也。但不正言其數，而設堤闊數、各縣物力數、一夫力數、一夫平方數，以取每日所修堤長數，故令人不能驟觧。\end{unitzh}
\begin{uniten}[Siku editors: The intent of the problem is this: the total length of dike repaired by the 4 counties is the same, while the length repaired each day differs; measuring the total length by the length each repairs per day, it either comes out exactly or leaves a surplus, and from this the total length is sought. But [the problem] does not state these numbers directly; instead it sets out the number for the width of the dike, the numbers for the property assets of each county, the number of assets per man and the number of square [chi] per man, from which to obtain the numbers for the length of dike repaired each day. Therefore it keeps one from understanding it at once.]\end{uniten}

\begin{unitzh}\unitlabel{shushu-008.3}答曰：堤長一十九里二百三十五歩五尺；甲縣夫築一千二十六丈{\zhsmall 乙、丙、丁同}；乙縣夫築一千七百六十八歩五尺六寸{\zhsmall 甲、丙、丁同}；丙縣夫築四里三百二十八歩五尺六寸{\zhsmall 甲、乙、丁同}；丁縣夫築{\zhsmall 同前三縣數}。\end{unitzh}
\begin{uniten}Answer: The length of the dike is 19 li 235 bu 5 chi. The laborers of County A built 1026 zhang [Note: B, C and D the same]; the laborers of County B built 1768 bu 5 chi 6 cun [Note: A, C and D the same]; the laborers of County C built 4 li 328 bu 5 chi 6 cun [Note: A, B and D the same]; the laborers of County D built [Note: the same as the numbers of the preceding 3 counties].\end{uniten}

\begin{unitzh}\unitlabel{shushu-008.4}按：四縣所築堤長等，則丈數、歩數、里數皆同。今以三数分載三縣下，而復註以與某縣同，殊混人目。\end{unitzh}
\begin{uniten}[Siku editors: The lengths of dike built by the 4 counties are equal, so the numbers in zhang, in bu and in li are all the same. Now the three numbers are recorded separately under three counties, and moreover annotated as being the same as such-and-such counties; this greatly confuses the reader's eye.]\end{uniten}

\begin{unitzh}\unitlabel{shushu-008.5}術曰：置各縣力，以程功程為實；以力率乗堤齊闊為法；除之，得各縣日築復數{\zhsmall 有分者通之，互乗之，得通數}。求總等，不約一位，約衆位，曰元數。連環求等，約竒，得定母。陸續求衍數、竒數、乗率、用數。以丙、丁縣不及數乘本用，併為總數。以定母相乘為衍母，滿母去總數，得各縣分給里歩積尺數。以縣數因之，為堤長。各以里法、歩法約之，為里歩。\end{unitzh}
\begin{uniten}Method: Set up the assets of each county; with the work quota (chenggong) "quota" [Translator's note: this second 程 is evidently an error for 乘 "multiply"] make the dividend (shi). With the assets rate (lilü) multiply the uniform width of the dike to make the divisor (fa). Divide, obtaining the daily building of each county as compound numbers (fushu) [Note: where there are fractions, convert them (tong) and cross-multiply them, obtaining converted numbers (tongshu)]. Seek the overall common measure (zongdeng); do not reduce one place, but reduce the many places; [the results] are called the original numbers (yuanshu). Seek the common measures in linked rings (lianhuan), reducing the odd, to obtain the fixed denominators (dingmu). In succession seek the extension numbers (yanshu), remainder numbers (jishu), multiplying rates (chenglü) and use numbers (yongshu). With the shortfall numbers of counties C and D multiply their own use [numbers], and add them to make the total number (zongshu). Multiply the fixed denominators together to make the extension denominator (yanmu); removing the denominator from the total number as often as it fills, obtain the number of accumulated chi of the li and bu allotted to each county. Multiply it by the number of counties to make the length of the dike. Reduce each by the li rule and the bu rule to make li and bu.\end{uniten}

\begin{unitzh}\unitlabel{shushu-008.6}草曰：置甲縣力一十三萬八千六百貫，乙縣力一十四萬六千三百貫，丙縣力一十九萬二千五百貫，丁縣力一十八萬四千八百貫，以程功六十尺徧乗之，皆以貫黙{\zhsmall 按：貫黙乃以一貫千文為法之名，與前官陌、市陌名相似}約之：甲得八百三十一萬六千尺，乙得八百七十七萬八千尺，丙得一千一百五十五萬尺，丁得一千一百八萬八千尺，各為實。次以力率七百七十貫乗堤齊闊二十尺，亦以貫黙約之，得一萬五千四百尺，為法。徧除諸各實：甲得五十四丈，乙得五十七丈，丙得七十五丈，丁得七十二丈，各為四縣衆夫每日築長率。按大衍術，命曰復數，列右行。\end{unitzh}
\begin{uniten}Working: Set up County A's assets, 138,600 guan, County B's assets, 146,300 guan, County C's assets, 192,500 guan, and County D's assets, 184,800 guan. Multiply them all by the work quota, 60 chi, and reduce them all by the silent guan (guanmo 貫黙) [Siku editors: guanmo 貫黙 is the name for taking 1 guan of 1000 wen (文) as the standard (fa); it is similar to the names official mo (guanmo 官陌) and market mo (shimo 市陌) above]: A obtains 8,316,000 chi, B obtains 8,778,000 chi, C obtains 11,550,000 chi, D obtains 11,088,000 chi; each is a dividend (shi). Next, with the assets rate, 770 guan, multiply the uniform width of the dike, 20 chi, and likewise reduce by the silent guan, obtaining 15,400 chi as the divisor (fa). Divide all the dividends throughout by it: A obtains 54 zhang, B obtains 57 zhang, C obtains 75 zhang, D obtains 72 zhang; each is the rate (lü) of the length built each day by the many laborers of [each of] the 4 counties. According to the Great Extension method (dayan shu), they are named compound numbers; list them in the right column.\end{uniten}

\begin{unitzh}\unitlabel{shushu-008.7}以復數{\zhsmall 按：前術以尾數在十以上者為復數，此數不合}求總等，得三大{\zhsmall 按：此條原本皆以丈為寸，於義無取，今皆改正}，以約三位多者，不約其少者：甲得五十四，乙得一十九，丙得二十五，丁得二十四，仍為元數。次以兩兩連環求等，各約之。\end{unitzh}
\begin{uniten}With the compound numbers [Siku editors: the earlier method takes those whose tail figure (weishu 尾數) is at the tens or above as compound numbers; these numbers do not accord with this] seek the overall common measure, obtaining 3 da [Translator's note: 大 "great" is evidently an error for 丈 zhang] [Siku editors: in this passage the original text everywhere takes zhang as cun, which makes no sense; now all have been corrected]. With it reduce the three places that are more, and do not reduce the lesser one: A obtains 54, B obtains 19, C obtains 25, D obtains 24; these then serve as the original numbers. Next seek the common measures pairwise in linked rings and reduce each by them.\end{uniten}

\begin{unitzh}\unitlabel{shushu-008.8}按：四縣每日築長數，皆以丈為单位，非位數也。但一等數可以度盡四數，必先求總等約之，然後可以為元數。即此可見總等法不獨用於通數、復數也。\end{unitzh}
\begin{uniten}[Siku editors: The numbers for the length built each day by the 4 counties all take the zhang as the unit; they are not place numbers (weishu 位數). But if one common measure can measure out the 4 numbers exactly, one must first seek the overall common measure and reduce them by it; only then can they serve as the original numbers. From this one can see that the method of the overall common measure is not used only for converted numbers and compound numbers.]\end{uniten}

\begin{unitzh}\unitlabel{shushu-008.9}先以丁、丙求等，又以丁、乙求等，皆得一，不約。次以丁、甲求等，得六，只約甲五十四，得九，不約丁。次以丙與乙求等，又以丙與甲九求等，皆得一，不約。後以乙與甲九求等，得一，不約。復驗甲九與丁二十四，猶可再約，又求等，得三，以約丁二十四，得八，復乘甲為二十七。\end{unitzh}
\begin{uniten}First seek the common measure of D and C, and again seek the common measure of D and B; both give 1, so do not reduce. Next seek the common measure of D and A, obtaining 6; reduce only A's 54, obtaining 9, and do not reduce D. Next seek the common measure of C and B, and again seek the common measure of C and A's 9; both give 1, so do not reduce. Afterwards seek the common measure of B and A's 9, obtaining 1, so do not reduce. On examining again, A's 9 and D's 24 can still be reduced once more: again seek the common measure, obtaining 3; with it reduce D's 24, obtaining 8, and in turn multiply A by it, making 27.\end{uniten}

\begin{unitzh}\unitlabel{shushu-008.10}次以定母四位相乗，求得一十萬二千六百，為衍母。各以定母約衍母：甲得三千八百，乙得五千四百，丙得四千一百四，丁得一萬二千八百二十五，為衍數。\end{unitzh}
\begin{uniten}Next multiply the 4 places of fixed denominators together, obtaining 102,600 as the extension denominator. Reduce the extension denominator by each fixed denominator: A obtains 3800, B obtains 5400, C obtains 4104, D obtains 12,825; these are the extension numbers.\end{uniten}

\begin{unitzh}\unitlabel{shushu-008.11}滿定母，各去衍數：甲不滿二十，乙不滿四，丙不滿四，丁不滿一，各為竒數。\end{unitzh}
\begin{uniten}Removing the fixed denominator from each extension number as often as it fills: for A what does not fill is 20, for B what does not fill is 4, for C what does not fill is 4, for D what does not fill is 1; each is the remainder number.\end{uniten}

\begin{unitzh}\unitlabel{shushu-008.12}以各定母與本竒數，用大衍求一術入之，各得乘率：甲得二十三，乙得五，丙得一十九，丁得一。\end{unitzh}
\begin{uniten}Take each fixed denominator together with its own remainder number and enter them into the Great Extension method of seeking one (dayan qiuyi shu); each obtains a multiplying rate: A obtains 23, B obtains 5, C obtains 19, D obtains 1.\end{uniten}

\begin{unitzh}\unitlabel{shushu-008.13}以右行乗率對乘寄左行衍數：甲得八萬七千四百，乙得二萬七千，丙得七萬七千九百七十六，丁得一萬二千八百二十五，各為用數。\end{unitzh}
\begin{uniten}With the multiplying rates of the right column multiply correspondingly the extension numbers lodged in the left column: A obtains 87,400, B obtains 27,000, C obtains 77,976, D obtains 12,825; each is a use number.\end{uniten}

\begin{unitzh}\unitlabel{shushu-008.14}次驗四縣所築有無不及零丈尺寸：今甲、乙俱畢，為無；丙餘五十一丈，丁餘一十八丈，為有。以丙、丁二縣餘丈各乗丙、丁二用數：其丙五十一丈乗丙用七萬七千九百七十六，得三百九十七萬六千七百七十六丈，為丙總；以丁餘一十八乗丁用一萬二千八百二十五，得二十三萬八百五十丈，為丁總。併二總，得四百二十萬七千六百二十六丈，為總數。亦以丈通衍母，得一十萬二千六百丈，仍為衍母。滿去總數，不滿一千二十六丈，為所求長率。以四縣因之，得四十一百四丈，為實。以歩法五尺八寸除之，得七千七十五歩五尺，為堤長歩。以里法三百六十歩約之，得一十九里二百三十五歩五尺，為堤通長。置長率一千二十六丈，以歩法約之，得一千七百六十八歩五尺六寸；又以里法約之，得四里三百二十八歩五尺六寸，為各縣所給道里歩尺數。\end{unitzh}
\begin{uniten}Next examine whether or not there is a shortfall of odd zhang, chi and cun in what the 4 counties built. Now A and B have both finished: this counts as "there is not". C has 51 zhang remaining and D has 18 zhang remaining: this counts as "there is". With the remaining zhang of the two counties C and D multiply respectively the two use numbers of C and D: thus C's 51 zhang multiplies C's use [number], 77,976, obtaining 3,976,776 zhang as C's total; with D's remainder, 18, multiply D's use [number], 12,825, obtaining 230,850 zhang as D's total. Add the two totals, obtaining 4,207,626 zhang as the total number. Likewise convert the extension denominator into zhang, obtaining 102,600 zhang, which still serves as the extension denominator. Remove it from the total number as often as it fills; what does not fill, 1026 zhang, is the length rate sought. Multiply it by the 4 counties, obtaining 4 ten 1 hundred 4 zhang [Translator's note: the text has 四十一百四, evidently for 四千一百四, 4104] as the dividend. Divide it by the bu rule, 5 chi 8 cun, obtaining 7075 bu 5 chi as the length of the dike in bu. Reduce it by the li rule, 360 bu, obtaining 19 li 235 bu 5 chi as the full length of the dike. Set up the length rate, 1026 zhang, and reduce it by the bu rule, obtaining 1768 bu 5 chi 6 cun; again reduce it by the li rule, obtaining 4 li 328 bu 5 chi 6 cun as the numbers of li, bu and chi of road allotted to each county.\end{uniten}

\begin{unitzh}\unitlabel{shushu-008.15}餘米推數\end{unitzh}
\begin{uniten}Deducing the amounts from the remaining rice\end{uniten}

\begin{unitzh}\unitlabel{shushu-009.1}問：有米鋪訴被盗去米一般三籮，皆適滿，不記細數。今左壁籮剰一合，中間籮剰一升四合，右壁籮剰一合。後獲賊，係甲、乙、丙三名。甲稱「當夜摸得馬杓，在左壁籮舀入布袋」；乙稱「踢着木履，在中籮舀入袋」；丙稱「摸得漆椀，在左邉籮舀入袋」。將歸食用，日乆不知數。索到三器：馬杓滿容一升九合，木履容一升七合，漆椀容一升二合。欲知所失米數，計贓結㫁，三盗各㡬何？答曰：共失米九石五斗六升三合。環甲米三石一斗九升二合，乙米三石一斗七升九合，丙米三石一斗九升二合。\end{unitzh}
\begin{uniten}Question: A rice shop reports that rice was stolen from it from three baskets alike (一般), all exactly full; it does not remember the detailed amount. Now in the basket by the left wall 1 ge (合) is left over; in the middle basket 1 sheng (升) 4 ge are left over; in the basket by the right wall 1 ge is left over. Later the thieves were caught: they are three persons, A, B and C. A states that on that night he groped about and got hold of a horse ladle, and at the basket by the left wall scooped [rice with it] into a cloth bag. B states that he kicked against a wooden clog, and at the middle basket scooped [with it] into a bag. C states that he groped about and got hold of a lacquered bowl, and at the basket on the left side scooped [with it] into a bag. They took it home and ate it; as many days have passed, they do not know the amount. The three vessels were obtained: the horse ladle holds 1 sheng 9 ge when full, the wooden clog holds 1 sheng 7 ge, the lacquered bowl holds 1 sheng 2 ge. One wishes to know the amount of rice lost, in order to reckon the stolen goods and settle the case: how much is it for each of the three thieves? Answer: In all 9 dan (石) 5 dou (斗) 6 sheng 3 ge of rice was lost. Huan (環) [Translator's note: unexplained here; apparently corrupt] A's rice 3 dan 1 dou 9 sheng 2 ge; B's rice 3 dan 1 dou 7 sheng 9 ge; C's rice 3 dan 1 dou 9 sheng 2 ge.\end{uniten}

\begin{unitzh}\unitlabel{shushu-009.2}術曰：以大衍求之。列三器所容為元數，連環求等，約為定母，以相乘為衍母。以定各約衍母，得衍數。各滿定母去之，得竒。以竒、定用大衍求得乘率，以乘衍數，得用數。次以各剰米乗用，併之為總，滿衍母去之，不滿為毎籮米。各以剰米減之，餘為甲、乙、丙盗米。併之為共失米。\end{unitzh}
\begin{uniten}Method: Solve it by the Great Extension (dayan 大衍). Set out what the three vessels hold as the original numbers (yuanshu 元數); seek the common divisors (deng 等) in a linked chain (lianhuan 連環) and reduce, to make the fixed denominators (dingmu 定母); multiply them together to make the extension denominator (yanmu 衍母). Divide the extension denominator by each of the fixed [denominators] to obtain the extension numbers (yanshu 衍數). From each remove the fixed denominator as many times as it fills, to obtain the remainders (qi 竒). With the remainders and the fixed [denominators] use the Great Extension to obtain the multiplier rates (chenglü 乘率); with them multiply the extension numbers to obtain the use numbers (yongshu 用數). Next, with each leftover rice multiply the use [number], and add them up as the total; remove from it the extension denominator as many times as it fills; what does not fill it is the rice of each basket. From it subtract each leftover rice; what remains is the rice stolen by A, B and C. Add them up for the rice lost in all.\end{uniten}

\begin{unitzh}\unitlabel{shushu-009.3}草曰：列三器所容一升九合、一升七合、一升二合為元數，連環求等，皆得一，不約。便以元數相乗，得三千八百七十六，為衍母。以各元數為定母，以定約衍母，得衍數：甲得二百四，乙得二百二十八，丙得三百二十三，各為衍數，列左行。以三定母甲十九、乙十七、丙十二列右行。具圖如後。\end{unitzh}
\begin{uniten}Working: Set out what the three vessels hold, 1 sheng 9 ge, 1 sheng 7 ge and 1 sheng 2 ge, as the original numbers. Seeking the common divisors in a linked chain, all give 1: do not reduce. Then directly multiply the original numbers together, obtaining 3876 as the extension denominator. Take each original number as a fixed denominator; divide the extension denominator by the fixed [denominators] to obtain the extension numbers: A obtains 204, B obtains 228, C obtains 323; each is an extension number; set them out in the left column. Set out the three fixed denominators, A 19, B 17, C 12, in the right column. The diagram is set out below.\end{uniten}

\begin{unitzh}\unitlabel{shushu-009.4}各滿定母去衍數，得竒數：甲得一十四，乙得七得七，丙得一十一。\end{unitzh}
\begin{uniten}From the extension numbers remove each fixed denominator as many times as it fills, to obtain the remainder numbers (qishu 竒數): A obtains 14; B obtains 7, obtains 7 [Translator's note: so repeated in the text]; C obtains 11.\end{uniten}

\begin{unitzh}\unitlabel{shushu-009.5}各以竒、定用大衍求一，各得乗率：甲得十五，乙得五，丙得一十一，各為乘率，列右行，對寄在行衍數。具圖如後。\end{unitzh}
\begin{uniten}With each remainder and fixed [denominator], use the Great Extension to seek one (qiuyi 求一), each obtaining a multiplier rate: A obtains 15, B obtains 5, C obtains 11; each is a multiplier rate. Set them out in the right column, facing the extension numbers lodged in the column. The diagram is set out below.\end{uniten}

\begin{unitzh}\unitlabel{shushu-009.6}以丙行對乗之，得用數：甲得三千六十，乙得一千一百四十，丙得三千五百五十三，列右行。具圖如後。\end{unitzh}
\begin{uniten}Multiply correspondingly with column C (丙行) [Translator's note: 丙 'C' is perhaps for 兩 'the two'], obtaining the use numbers: A obtains 3060, B obtains 1140, C obtains 3553; set them out in the right column. The diagram is set out below.\end{uniten}

\begin{unitzh}\unitlabel{shushu-009.7}既得用數，始驗問題三籮剰米，列左行，對三人所用。以兩行對乗之：甲得三千六十，乙得一萬五千九百六十，丙得三千五百五十三。\end{unitzh}
\begin{uniten}Having obtained the use numbers, only then examine the leftover rice of the three baskets in the problem; set it out in the left column, facing the use [numbers] of the three persons. Multiply the two columns correspondingly: A obtains 3060, B obtains 15960, C obtains 3553.\end{uniten}

\begin{unitzh}\unitlabel{shushu-009.8}併三數，得二萬二千五百七十三，為總數。滿衍母三千八百七十六去之，不滿三千一百九十三合，展為三石一斗九升三合，為三籮適滿細數。以左籮剰一合減之，餘三石一斗九升二合，為甲盗米，又為丙盗米。以中籮剰米一升四合減之，餘三石一斗七升九合，為乙盗米。併三人，共得九石五斗六升三合，為所失總米。合問。\end{unitzh}
\begin{uniten}Add the three numbers, obtaining 22573 as the total number. Remove from it the extension denominator 3876 as many times as it fills; what does not fill it is 3193 ge, which expanded is 3 dan 1 dou 9 sheng 3 ge: this is the detailed amount of the three baskets when exactly full. Subtract from it the 1 ge left over in the left basket; what remains, 3 dan 1 dou 9 sheng 2 ge, is the rice stolen by A, and is also the rice stolen by C. Subtract from it the rice left over in the middle basket, 1 sheng 4 ge; what remains, 3 dan 1 dou 7 sheng 9 ge, is the rice stolen by B. Adding the three persons, together one obtains 9 dan 5 dou 6 sheng 3 ge as the total rice lost. This agrees with the question.\end{uniten}

\begin{unitzh}\unitlabel{shushu-009.9}程行計地\end{unitzh}
\begin{uniten}Rates of Travel: Reckoning the Distance\end{uniten}

\begin{unitzh}\unitlabel{shushu-009.10}問：軍師獲㨗，當早㸃差急足三名，徃都下節節走報。具：甲於前數日申末到，乙後數日未正到，丙於今日辰未到。據供：甲日行三百里，乙日行二百四十里，丙日行一百八十里。問自軍前至都里數及三人各行日數㡬何？答曰：軍前至都三千三百里。甲行一十一日，乙行一十三日四時半，丙行一十八日二時。\end{unitzh}
\begin{uniten}Question: The army (軍師) won a victory, and that same morning appointed and dispatched three express couriers (jizu 急足) to go to the capital, running to report one after another (節節). It is recorded (具) that A arrived several days before, at the end (mo 末) of the shen (申) double-hour; B arrived several days after, at the middle (zheng 正) of the wei (未) double-hour; C arrived today at chen wei (辰未) [Translator's note: evidently for 辰末, the end of the chen double-hour]. According to their statements, A travels 300 li (里) a day, B travels 240 li a day, C travels 180 li a day. What are the number of li from the army front to the capital and the number of days each of the three persons traveled? Answer: From the army front to the capital is 3300 li. A traveled 11 days; B traveled 13 days 4 1/2 double-hours (shi 時); C traveled 18 days 2 double-hours.\end{uniten}

\begin{unitzh}\unitlabel{shushu-009.11}術曰：以大衍求之。置各行里，先求總等，存一約衆，得原里。次以連環求等，約竒復乘偶，得定母。以定相乘為衍母。滿定除衍，衍得乘率，以乘衍數，得用數。次置辰刻正末，乗各行里為實，以晝六時約之，得餘里。各乗用數，併為總，滿衍母去，得所求至都里。以各日行約之，得日辰刻數。\end{unitzh}
\begin{uniten}Method: Solve it by the Great Extension (dayan 大衍). Set out the li traveled by each; first seek the overall common divisor (zongdeng 總等); keep one and reduce the many, to obtain the original li (原里). Next seek the common divisors in a linked chain (lianhuan qiu deng 連環求等); reduce the odd (qi 竒) and multiply back the even (ou 偶), to obtain the fixed denominators (dingmu 定母). Multiply the fixed [denominators] together to make the extension denominator (yanmu 衍母). Remove the fixed [denominators] from the extension [numbers] as many times as they fill; by the Extension obtain the multiplier rates (chenglü 乘率); with them multiply the extension numbers (yanshu 衍數) to obtain the use numbers (yongshu 用數). Next set out the double-hour and ke (刻) (辰刻), middle or end (正末); multiply each one's li traveled by it as the dividend (shi 實); divide by the 6 double-hours of daytime to obtain the remaining li (yuli 餘里). Multiply each by its use number and add them up as the total; remove the extension denominator as many times as it fills, to obtain the sought li to the capital. Divide by each one's daily travel to obtain the numbers of days, double-hours and ke.\end{uniten}

\begin{unitzh}\unitlabel{shushu-009.12}草曰：置甲三百里，乙二百四十里，丙一百八十里。先求總等，得六十。只存甲三百勿約，乃約乙二百四十，得四；次約丙一百八十，得三，各為元數。\end{unitzh}
\begin{uniten}Working: Set out A 300 li, B 240 li, C 180 li. First seek the overall common divisor: it gives 60. Keep only A's 300 and do not reduce it; then reduce B's 240, obtaining 4; next reduce C's 180, obtaining 3; each is an original number (yuanshu 元數).\end{uniten}

\begin{unitzh}\unitlabel{shushu-009.13}連環求等：先以丙、乙求等，得一，不約。次以丙、甲求等，三。於術約竒不約偶，盖以等三約三，因得一為竒，盧無衍數，乃便徑先約甲三百為一百，復以等三兼丙三為九。既丙九為竒，甲百為偶，此即是約竒弗約偶。次以一四與甲百求等，得四，以四約一百，得二十五，為甲；復以四乗乙四，得一十六，為乙，各為定母。\end{unitzh}
\begin{uniten}Seek the common divisors in a linked chain: first seek the common divisor of C and B: it gives 1; do not reduce. Next seek the common divisor of C and A: 3. In the method one reduces the odd (qi) and does not reduce the even (ou); but reducing 3 by the common divisor 3 would thereby give 1 as the odd, and 盧 [Translator's note: evidently for 慮, 'fearing that'] there would be no extension number, so one then directly reduces A's 300 first, to 100, and further combines the common divisor 3 with C's 3 to make 9. Since C's 9 is the odd and A's 100 is the even, this is precisely 'reduce the odd and do not reduce the even'. Next seek the common divisor of 1 4 (一四) [Translator's note: perhaps for 乙四, B's 4] and A's 100: it gives 4. Reduce 100 by 4 to obtain 25 as A; further multiply B's 4 by 4 to obtain 16 as B; each is a fixed denominator.\end{uniten}

\begin{unitzh}\unitlabel{shushu-009.14}以定母相乗，得三千六百，為衍母。以各定約衍母，為衍數：甲得一百四十四，乙得二百二十五，丙得四百。\end{unitzh}
\begin{uniten}Multiply the fixed denominators together, obtaining 3600 as the extension denominator. Divide the extension denominator by each fixed [denominator] for the extension numbers: A obtains 144, B obtains 225, C obtains 400.\end{uniten}

\begin{unitzh}\unitlabel{shushu-009.15}術數各滿衍母去之，不滿為竒數：甲得一十九，乙得一，丙得四。\end{unitzh}
\begin{uniten}From each method number (術數) remove the extension denominator [Translator's note: 術數 and 衍母 here evidently stand for the extension numbers and the fixed denominators] as many times as it fills; what does not fill it is the remainder number: A obtains 19, B obtains 1, C obtains 4.\end{uniten}

\begin{unitzh}\unitlabel{shushu-009.16}以各竒數與定母用大衍入之，各得乘數：甲得四，乙得一，丙得七，各為乗率，列右行。\end{unitzh}
\begin{uniten}Take each remainder number with the fixed denominator and enter them into the Great Extension, each obtaining a multiplier number (chengshu 乘數): A obtains 4, B obtains 1, C obtains 7; each is a multiplier rate; set them out in the right column.\end{uniten}

\begin{unitzh}\unitlabel{shushu-009.17}以乗率對乘寄左行衍數：甲得五百七十六，乙得二百二十五，丙得二千八百，各為用數。\end{unitzh}
\begin{uniten}With the multiplier rates multiply correspondingly the extension numbers lodged in the left column: A obtains 576, B obtains 225, C obtains 2800; each is a use number.\end{uniten}

\begin{unitzh}\unitlabel{shushu-009.18}次置甲申末到者，其酉初為夜，此是甲以全日到，為無餘里。次置乙於未正到，乃於卯時數至未正，得四箇半時，以四半乗乙行二百四十里，得一千八十，為實，以畫六時約之，得一百八十里，為乙行不及全日之餘里。次置丙於辰未到，自卯初數至辰未，得二時，以因丙行一百八十里，得三百六十，為實，以六時除之，得六十里，為丙行不及全日之餘里。\end{unitzh}
\begin{uniten}Next set out A, who arrived at the end of shen: the beginning of you (酉) following it is night; this means A arrived in whole days, which counts as no remaining li. Next set out B's arrival at the middle of wei: counting from the mao (卯) double-hour to the middle of wei gives 4 1/2 double-hours; multiply B's travel of 240 li by 4 1/2, obtaining 1080 as the dividend; divide it by the 6 double-hours of daytime (畫) [Translator's note: 畫 'drawing' is evidently for 晝 'daytime'], obtaining 180 li as the remaining li of B's travel, not amounting to a whole day. Next set out C's arrival at chen wei: counting from the beginning of mao to chen wei gives 2 double-hours; with them multiply C's travel of 180 li, obtaining 360 as the dividend; divide it by the 6 double-hours, obtaining 60 li as the remaining li of C's travel, not amounting to a whole day.\end{uniten}

\begin{unitzh}\unitlabel{shushu-009.19}以乙餘一百八十乗乙用二百二十五，得四萬五百於中。以丙餘六十乗丙用二千八百，得十六萬八千，加中，共得二十萬八千五百，為總。滿衍母三千六百去之，不滿三千三百里，為軍前至都里。以甲三百除之，得一十一日；以乙二百四十除之，得一十三日四時半；以丙一百八十除之，得一十八日二時。合問。\end{unitzh}
\begin{uniten}With B's remaining 180 multiply B's use [number] 225, obtaining 40500, [placed] in the middle. With C's remaining 60 multiply C's use [number] 2800, obtaining 168000; add it to the middle, together obtaining 208500 as the total. Remove from it the extension denominator 3600 as many times as it fills; what does not fill it, 3300 li, is the li from the army front to the capital. Divide it by A's 300 to obtain 11 days; divide it by B's 240 to obtain 13 days 4 1/2 double-hours; divide it by C's 180 to obtain 18 days 2 double-hours. This agrees with the question.\end{uniten}

\begin{unitzh}\unitlabel{shushu-009.20}按：凡總等數必小於連環等數，若甚大，即為連環等數，此題數是也。故再約，即用求續等法，不然不能合也。\end{unitzh}
\begin{uniten}Note: In general the overall common divisor must be smaller than the linked-chain common divisors; if it is very large, it then becomes a linked-chain common divisor; the numbers of this problem are of this kind. Therefore one reduces again, that is, one uses the method of seeking the continued common divisor (qiu xu deng fa 求續等法); otherwise it cannot agree.\end{uniten}

\begin{unitzh}\unitlabel{shushu-009.21}程行相及\end{unitzh}
\begin{uniten}Rates of Travel: Catching Up\end{uniten}

\begin{unitzh}\unitlabel{shushu-010.1}問：有急足三名，甲日行三百里，乙日行二百五十里，丙日行二百里。先差丙徃他處下文字，旣兩日，又有文字，遣乙追付；已半日，復有文字，續令甲趕付乙。三人偶不相及，乃同時俱至彼所。先欲知乙果及丙、甲果及乙得日并里，次問彼處去此里數，各㡬何？\end{unitzh}
\begin{uniten}Question: There are 3 express couriers. A travels 300 li (里) a day, B travels 250 li a day, and C travels 200 li a day. First C was dispatched to go to another place to deliver documents. When 2 days had passed, there were more documents, and B was sent to pursue and hand them over. After half a day, there were documents once again, and A was then ordered to hurry and hand them to B. The three men happened not to catch up with one another, and all arrived at that place at the same time. First we wish to know the days and li obtained for B actually catching up with C and for A actually catching up with B; next we ask the number of li from that place to here. How much is each?\end{uniten}

\begin{unitzh}\unitlabel{shushu-010.2}按：題意謂三行遲疾不同，乙後丙兩日，甲後乙半日，問㡬日㡬里可以追及。又旣及之後，三人不能同行，及各至彼處之時刻，皆與各起程之時刻相同，盖言自此至彼所行皆為整日數也。\end{unitzh}
\begin{uniten}Note: The meaning of the problem is that the three travel at different speeds, slow and fast; B is 2 days behind C, and A is half a day behind B; it asks in how many days and over how many li they can catch up. Further, after they have caught up, the three cannot travel together, and the hour at which each arrives at that place is in every case the same as the hour at which each set out on the road: that is to say, what each travels from here to there is in every case a whole number of days.\end{uniten}

\begin{unitzh}\unitlabel{shushu-010.3}答曰：乙果追及丙，八日，行二千里；甲果追及乙，二日半，行七百五十里；彼處去此三千里。\end{unitzh}
\begin{uniten}Answer: B actually catches up with C in 8 days, traveling 2000 li; A actually catches up with B in 2 1/2 days, traveling 750 li; that place is 3000 li from here.\end{uniten}

\begin{unitzh}\unitlabel{shushu-010.4}術曰：以均輸求之，大衍入之。置乙己去日數，乗乙行里為實；以甲乙行里差為法，除之，得甲及乙日數辰刻；以乗甲行，得里。次置丙旣去日，乗丙行里為實；以丙乙行里差為法，除之，得乙及丙日數；以乗乙行，得里。然後置三人日行，求總等，約得原數；以連環求等，約得定母；以定相乗，得衍母；各定約衍，得衍數；滿定去衍，得竒；竒定大衍，得乗率；以乗寄衍，得用數。視甲及乙里為乙率，乙及丙里為丙率。以乙日行滿去乙率，不滿為乙餘；以丙日行滿去丙率，不滿為丙餘。以二餘各乗本用，併之為總，滿衍去之，不滿為彼去此里。\end{unitzh}
\begin{uniten}Method: Solve it by equitable transport (junshu), and bring it into the great extension (dayan). Set up the number of days B had already gone and multiply it by B's li of travel as the dividend (shi); take the difference between the li traveled by A and B as the divisor (fa); divide, giving the number of days, chen (辰) and ke (刻) for A to catch up with B; multiply A's travel by it, giving the li. Next set up the days C had already gone and multiply it by C's li of travel as the dividend; take the difference between the li traveled by C and B as the divisor; divide, giving the number of days for B to catch up with C; multiply B's travel by it, giving the li. Then set up the daily travel of the three men, seek the overall common divisor (zongdeng), and reduce, giving the original numbers (yuanshu 原數); by seeking the common divisor in linked rings (lianhuan qiudeng), reduce, giving the fixed denominators (dingmu); multiply the fixed [denominators] together, giving the extension denominator (yanmu); reduce the extension [denominator] by each fixed [denominator], giving the extension numbers (yanshu); remove the fixed [denominator] from the extension [number] as often as it is full, giving the remainder (qi); from the remainder and the fixed [denominator], by the great extension, obtain the multiplying rate (chenglü); multiply the lodged extension [numbers] by them, giving the use numbers (yongshu). Take the li of A catching up with B as B's rate (lü), and the li of B catching up with C as C's rate. Remove B's daily travel from B's rate as often as it is full; what does not fill is B's residue. Remove C's daily travel from C's rate as often as it is full; what does not fill is C's residue. Multiply each of the two residues by its own use [number] and add them to make the total; remove the extension [denominator] from it as often as it is full; what does not fill is the li from there to here.\end{uniten}

\begin{unitzh}\unitlabel{shushu-010.5}草曰：置乙己去日，乘乙日行二百五十里，得一百二十五里，為實。次置甲日行三百里，減去乙行二百五十里，餘五十里，為差，法除實，得二日五十刻，甲果及乙數；以乗甲行三百里，得七百五十，為甲及乙里數。次置丙既行二百，乗丙日二百里，得四百里，為實。次置乙行二百五十里，減丙行二百里，餘五十里，為差，法除實，得八百，為乙及丙日數；以乗乙行二百五十里，得二千里，為乙得及丙之里數。已上為先欲知果及數。次列甲乙丙三名日行，求總等，得五十，先約甲丙，存得甲六、乙二百五十、丙四。\end{unitzh}
\begin{uniten}Working: Set up the days B had already gone, multiply it by B's daily travel of 250 li, giving 125 li, as the dividend. Next set up A's daily travel of 300 li, subtract B's travel of 250 li, leaving 50 li, as the difference; the divisor divides the dividend, giving 2 days 50 ke, the number for A actually catching up with B; multiply A's travel of 300 li by it, giving 750, as the number of li for A to catch up with B. Next set up C's already traveled 200, multiply it by C's daily 200 li, giving 400 li, as the dividend. Next set up B's travel of 250 li, subtract C's travel of 200 li, leaving 50 li, as the difference; the divisor divides the dividend, giving 800, as the number of days for B to catch up with C [Translator's note: 八百 "800" here and 二百 "200" above are evidently errors for 八日 "8 days" and 二日 "2 days".]; multiply B's travel of 250 li by it, giving 2000 li, as the number of li in which B manages to catch up with C. The above gives the numbers for actually catching up that were first wished to be known. Next lay out the daily travel of the three, A, B and C, and seek the overall common divisor, getting 50; first reduce A and C, leaving A 6, B 250, C 4.\end{uniten}

\begin{unitzh}\unitlabel{shushu-010.6}以甲六、丙四求等，得二；以二約甲為三，復以二因丙為八。次將乙二百五十與丙八相約，得二；乃約乙為一百二十五，復以二因丙為十六。定得甲三、乙一百二十五、兩十六，為定母。\end{unitzh}
\begin{uniten}With A 6 and C 4 seek the common divisor, getting 2; reduce A by 2, making 3, and again multiply (yin) C by 2, making 8. Next reduce B 250 and C 8 against each other, getting 2; then reduce B, making 125, and again multiply C by 2, making 16. This definitively gives A 3, B 125, "two" 16 [Translator's note: 兩 "two" is evidently an error for 丙 "C".], as the fixed denominators.\end{uniten}

\begin{unitzh}\unitlabel{shushu-010.7}以定相乗，得六千，為衍母；以各定約衍母，得衍數：甲得二千，乙得四十八，丙得三百七十五。求竒數。\end{unitzh}
\begin{uniten}Multiplying the fixed [denominators] together gives 6000 as the extension denominator; reducing the extension denominator by each fixed [denominator] gives the extension numbers: A gets 2000, B gets 48, C gets 375. Seek the remainders (qishu).\end{uniten}

\begin{unitzh}\unitlabel{shushu-010.8}左上二千，以甲三去之，竒二；左中四十八，即為乙竒；左下三百七十五，丙十六去之，竒也。\end{unitzh}
\begin{uniten}Upper left, 2000: remove A's 3 from it, remainder 2. Middle left, 48: this is B's remainder. Lower left, 375, from which C's 16 is removed: the remainder.\end{uniten}

\begin{unitzh}\unitlabel{shushu-010.9}各以大衍求，得甲二、乙一百一十二、丙七，各為乘率。\end{unitzh}
\begin{uniten}Seeking each by the great extension gives A 2, B 112, C 7, each as the multiplying rate.\end{uniten}

\begin{unitzh}\unitlabel{shushu-010.10}以乗率對乗衍數，甲得四千，乙得五千三百七十六，丙得二千六百二十五，為泛用數。\end{unitzh}
\begin{uniten}Multiply the extension numbers correspondingly by the multiplying rates: A gets 4000, B gets 5376, C gets 2625, as the provisional use numbers (fanyongshu).\end{uniten}

\begin{unitzh}\unitlabel{shushu-010.11}併三泛，得二萬二千一，乃多衍母一倍，當半衍母六千，得三千，以消甲四千，餘一千；又消乙五千三百七十六，餘二千三百七十六；丙不消；各為定用數。\end{unitzh}
\begin{uniten}Adding the three provisional [numbers] gives 22001; this exceeds the extension denominator by one fold. One should halve the extension denominator 6000, giving 3000, and with it diminish A's 4000, leaving 1000; further diminish B's 5376, leaving 2376; C is not diminished. Each is a fixed use number (dingyongshu).\end{uniten}

\begin{unitzh}\unitlabel{shushu-010.12}既得用數，次視前草中甲及乙七百五十里為乙率，乙及丙一千里為丙率，各滿乙丙日行里去之。\end{unitzh}
\begin{uniten}Having obtained the use numbers, next take, from the working above, the 750 li of A catching up with B as B's rate, and the 1000 li of B catching up with C as C's rate; remove from each the li of the daily travel of B and C, respectively, as often as they are full.\end{uniten}

\begin{unitzh}\unitlabel{shushu-010.13}今乙丙二人所行，各皆適滿去之，無餘；雖稱同時俱至，乃各係全日所行。便以乙丙二人約六千里，得三千里，為彼去此里數。合問。\end{unitzh}
\begin{uniten}Now what the two men B and C traveled is in each case exactly removed when full, with no residue; although it is said that they all arrived at the same time, it is in each case what was traveled in whole days. So, with the two men B and C, reduce the 6000 li, giving 3000 li as the number of li from there to here. This agrees with the question.\end{uniten}

\begin{unitzh}\unitlabel{shushu-010.14}按：復數求元數，用總等法，尚屬未密。盖總等約後，有當連環求等者，有當即求續等者，其法不能定也。今少為變通：凡復數皆見十者，先以十為總等，徧約之{\zhsmall 百千萬同}，為元數；俟連環求等畢，復以總等十乗一數{\zhsmall 百千萬同}，然後再求續等，以得定數。爰依題數，具式於後。\end{unitzh}
\begin{uniten}Note: For seeking the primary numbers (yuanshu 元數) from repeated numbers (fushu), using the method of the overall common divisor is still not rigorous. For after reducing by the overall common divisor, there are cases in which one should seek the common divisor in linked rings, and cases in which one should immediately seek the continuing common divisor (xudeng); the method cannot be fixed. Now we make a slight adaptation: whenever the repeated numbers all show tens, first take 10 as the overall common divisor and reduce them all by it [Note: the same for hundreds, thousands and ten-thousands], making the primary numbers; when the seeking of the common divisor in linked rings is finished, multiply one number by the overall common divisor 10 [Note: the same for hundreds, thousands and ten-thousands], and after that seek the continuing common divisor again, so as to obtain the fixed numbers (dingshu). Accordingly, following the numbers of the problem, we set out the pattern below.\end{uniten}

\begin{unitzh}\unitlabel{shushu-010.15}法：列三數於上，以十為總等，徧約之，得甲三十、乙二十五、丙二十，即為元數。連環求等：以甲與乙約，得等數五{\zhsmall 竒}，約甲得六{\zhsmall 偶}；以甲與丙約，得等數二{\zhsmall 偶}，約甲得三{\zhsmall 竒}，為甲數徧約畢。次以乙與丙約，得等數五{\zhsmall 竒}，約丙得四{\zhsmall 偶}，為乙丙二數徧約畢。乃以總等十乗乙數，得二百五十。次求續等：以甲與乙與丙相約，俱無續等；以乙與丙約，得續等二{\zhsmall 偶}，約乙數得一百二十五{\zhsmall 竒}，復乗丙得八。則甲三、乙一百二十五、丙八，即為各定數也。以三定數連乗，得三千，為衍母，即所問彼處去此之里數。較舊術，算省而數亦確矣。\end{unitzh}
\begin{uniten}Rule: Lay out the three numbers on top, take 10 as the overall common divisor and reduce them all by it, giving A 30, B 25, C 20, which are the primary numbers. Seek the common divisor in linked rings: reducing A with B gives the common divisor (dengshu) 5 [Note: odd]; reducing A gives 6 [Note: even]; reducing A with C gives the common divisor 2 [Note: even]; reducing A gives 3 [Note: odd]; as A's number, its all-round reduction is finished. Next, reducing B with C gives the common divisor 5 [Note: odd]; reducing C gives 4 [Note: even]; for the two numbers of B and C the all-round reduction is finished. Then multiply B's number by the overall common divisor 10, giving 250. Next seek the continuing common divisor: reducing A with B and with C, in neither case is there a continuing common divisor; reducing B with C gives the continuing common divisor 2 [Note: even]; reducing B's number gives 125 [Note: odd]; multiplying C again gives 8. Then A 3, B 125, C 8 are the respective fixed numbers. Multiplying the three fixed numbers in succession gives 3000 as the extension denominator, which is the number of li from that place to here that is asked. Compared with the old method, the calculation is economical and the numbers are also exact.\end{uniten}

\begin{unitzh}\unitlabel{shushu-010.16}數學九章卷一下\end{unitzh}
\begin{uniten}Shuxue jiuzhang, juan 1, second part.\end{uniten}

\begin{unitzh}\unitlabel{shushu-010.17}子部，天文算法類，算書之屬，數學九章\end{unitzh}
\begin{uniten}Masters section (zibu), astronomy and mathematics category, subcategory of mathematical books, Shuxue jiuzhang.\end{uniten}

\section*{卷2上}\addcontentsline{toc}{section}{卷2上}

\begin{unitzh}\unitlabel{shushu-011.1}欽定四庫全書。《數學九章》卷二上。宋秦九韶撰。天時。\end{unitzh}
\begin{uniten}Imperially Commissioned Complete Library of the Four Treasuries (Qinding Siku quanshu). Mathematics in Nine Sections (Shuxue jiuzhang), juan 2, first part. Composed by Qin Jiushao of the Song. Heavenly Seasons (Tianshi).\end{uniten}

\begin{unitzh}\unitlabel{shushu-011.2}推氣治厯\end{unitzh}
\begin{uniten}Calculating the Qi to Regulate the Calendar (tuiqi zhili)\end{uniten}

\begin{unitzh}\unitlabel{shushu-011.3}問：太史測驗天道，慶元四年戊午嵗冬至三十九日九十二刻四十五分；紹定三年庚寅嵗冬至三十二日九十四刻一十二分。欲求中間嘉泰甲子嵗氣骨、嵗餘、斗分，各得幾何？按：紹定三年庚寅之冬至，實紹定四年辛夘之始。辛夘距戊午三十四年，積年為三十三。答曰：氣骨十一日三十八刻二十分八十一杪八十小分；餘嵗五日二十四刻二十九分三十杪三十小分；斗分空日二十四刻二十九分三十杪三十小分。\end{unitzh}
\begin{uniten}Question: The Grand Astrologer (taishi) observed and verified the Way of Heaven. The winter solstice of the year wuwu, the 4th year of Qingyuan, [was at] 39 days 92 ke (刻) 45 fen (分); the winter solstice of the year gengyin, the 3rd year of Shaoding, [was at] 32 days 94 ke 12 fen. One wishes to find, for the jiazi year of Jiatai in between, the qi-bone (qigu), the year remainder (suiyu) and the Dipper fraction (doufen): how much is each? Note: The winter solstice of the year gengyin, 3rd year of Shaoding, is in fact the beginning of the year xinmao, 4th year of Shaoding. Xinmao is 34 years from wuwu; the accumulated years (jinian) are 33. Answer: Qi-bone 11 days 38 ke 20 fen 81 miao (杪) 80 small fen (xiaofen 小分); year remainder 5 days 24 ke 29 fen 30 miao 30 small fen; Dipper fraction zero (kong) days 24 ke 29 fen 30 miao 30 small fen.\end{uniten}

\begin{unitzh}\unitlabel{shushu-011.4}術曰：先距前後年數為法。置前測曰刻分，减後測曰刻分，餘為率。{\zhsmall 不足减，則加紀䇿。}以紀䇿累加之，今及天道，合用五日以上數為實。以法除實，得嵗餘。去全日，餘為斗分。以所求中間年上距前測年數乘嵗餘，益入前測日刻分，滿紀䇿去之，餘為所求氣骨。\end{unitzh}
\begin{uniten}Method: First [take] the number of years separating the earlier and the later [observations] as the divisor (fa). Set up the days [Translator's note: the text writes 曰 "say" for 日 "days", here and in the next clause], ke and fen of the earlier observation and subtract it from the days, ke and fen of the later observation; the remainder is the rate (lü). [Note: If it is insufficient to be subtracted from, then add the cycle tally (jice).] Repeatedly add the cycle tally to it; now, to reach [agreement with] the Way of Heaven, it is fitting to use a number of 5 days or more as the dividend (shi). Divide the dividend by the divisor, obtaining the year remainder. Remove the whole days; the remainder is the Dipper fraction. Multiply the year remainder by the number of years from the sought intermediate year back up to the earlier observation, and add [this] into the days, ke and fen of the earlier observation; when it fills the cycle tally, remove it; the remainder is the sought qi-bone.\end{uniten}

\begin{unitzh}\unitlabel{shushu-011.5}草曰：置前測戊午嵗距前嵗庚寅嵗，得三十三，為法。置前測戊午嵗冬至三十九日{\zhsmall 日辰癸夘}九十二刻四十五分，減後測紹定三年庚寅嵗冬至三十二日{\zhsmall 日辰丙申}九十四刻一十二分。今後測者少，不及前測者以減，乃加紀法六十日於後測日内，得九十二日九十四日一十二分。然後用前測者減之，餘五十三日一刻六十七分，為率。按術，當以法三十三，餘率須使啇數必得五日以上，乃可。今率未得五日，乃兩度累加紀法一百二十入率内，共得一百七十三日一刻六十七分，為刻實。如法除之，得五日二十四刻二十九分三十杪三十小分{\zhsmall 不盡棄之}，為嵗餘。乃去全五日，得二十四刻二十九分三十杪三十小分，為斗分。\end{unitzh}
\begin{uniten}Working: Set up [the distance] from the earlier observation, the year wuwu, to the former year, the year gengyin, obtaining 33, as the divisor. Set up the winter solstice of the earlier observation, the year wuwu, 39 days [Note: day sign (richen) guimao] 92 ke 45 fen, and subtract it from the winter solstice of the later observation, the year gengyin, 3rd year of Shaoding, 32 days [Note: day sign bingshen] 94 ke 12 fen. Now the later observation is smaller and not enough for the earlier observation to be subtracted from it, so add the cycle divisor (jifa), 60 days, into the days of the later observation, obtaining 92 days 94 days [Translator's note: the second 日 "days" evidently stands for 刻 ke] 12 fen. Then subtract the earlier observation from it, leaving 53 days 1 ke 67 fen, as the rate. According to the method, one should use the divisor 33; the remaining rate must be made such that the quotient necessarily reaches 5 days or more, and only then is it acceptable. Now the rate does not yet reach 5 days, so twice cumulatively add the cycle divisor, 120, into the rate, obtaining altogether 173 days 1 ke 67 fen, as the ke-dividend (keshi). Divide it by the divisor, obtaining 5 days 24 ke 29 fen 30 miao 30 small fen [Note: what is not exhausted is discarded], as the year remainder. Then remove the whole 5 days, obtaining 24 ke 29 fen 30 miao 30 small fen, as the Dipper fraction.\end{uniten}

\begin{unitzh}\unitlabel{shushu-011.6}次推嘉泰甲子上距慶元戊午嵗，得六，以乘嵗餘五日二十刻刻二十九分三十杪三十小分，得三十一日四十五刻七十五分八十一杪八十小分。益入前測戊午嵗三十九日九十二刻四十五分，得七十一日三十八刻二十分八十一杪八十小分。滿紀法六十去之，餘一十一日三十八刻二十分八十一杪八十小分，為所求甲子氣骨之數。合問。\end{unitzh}
\begin{uniten}Next calculate [the years] from the jiazi [year] of Jiatai back up to the year wuwu of Qingyuan, obtaining 6; multiply by it the year remainder, 5 days 20 ke ke [Translator's note: 二十刻刻 evidently stands for 二十四刻, 24 ke] 29 fen 30 miao 30 small fen, obtaining 31 days 45 ke 75 fen 81 miao 80 small fen. Add it into the earlier observation, the year wuwu, 39 days 92 ke 45 fen, obtaining 71 days 38 ke 20 fen 81 miao 80 small fen. When it fills the cycle divisor, 60, remove it, leaving 11 days 38 ke 20 fen 81 miao 80 small fen, as the number of the sought jiazi qi-bone. This answers the question.\end{uniten}

\begin{unitzh}\unitlabel{shushu-011.7}按：氣骨者，年冬至時距甲子日子正初刻後之日分也；嵗餘者，嵗實去六甲子之餘日分也；斗分者，嵗實去三百六十五日之餘分也。此未知嵗實之法，故先以前後兩氣骨相減，餘數為實，以積年為法除之。嵗餘約五日餘，紀日六十，故實數内累加六十日，至啇得五日上而止，則實數為積嵗餘之數，以積年除之，得嵗餘日分。既得嵗餘，以甲子積年六乘之，得甲子積嵗餘，與前測氣骨相加，滿紀法去之，餘即甲子氣骨也。\end{unitzh}
\begin{uniten}Note: The qi-bone is the days and fen by which the time of the winter solstice of the year lies after the first ke of the exact zi [hour] (zizheng chuke) of a jiazi day. The year remainder is the days and fen remaining when six jiazi [cycles] are removed from the year dividend (suishi). The Dipper fraction is the fen remaining when 365 days are removed from the year dividend. This is a method for when the year dividend is not known; therefore one first subtracts the two qi-bones, earlier and later, from each other, takes the remaining number as the dividend, and divides it by the accumulated years taken as the divisor. The year remainder is about 5-odd days, and the days of a cycle are 60; therefore one repeatedly adds 60 days into the dividend, stopping when the quotient reaches 5 days or more; then the dividend is the number of the accumulated year remainders, and dividing it by the accumulated years gives the days and fen of the year remainder. Having obtained the year remainder, multiply it by 6, the accumulated years to jiazi, obtaining the accumulated year remainder of jiazi; add it to the qi-bone of the earlier observation; when it fills the cycle divisor, remove it; the remainder is the qi-bone of jiazi.\end{uniten}

\begin{unitzh}\unitlabel{shushu-011.8}治厯推閏\end{unitzh}
\begin{uniten}Regulating the Calendar and Calculating the Intercalation (zhili tuirun)\end{uniten}

\begin{unitzh}\unitlabel{shushu-011.9}問：開禧厯以嘉泰四年甲子嵗天正冬至為一十一日{\zhsmall 日辰乙亥}四十四刻六十一分五十四杪，十一月經朔一日{\zhsmall 日辰乙丑}七十五刻五十五分六十二杪。問閏骨、閏率各幾何？答曰：閏骨九日六十九刻五分九十一杪{\zhsmall 不盡一百六十九分杪之一百二十一}；閏骨率十六萬三千七百七十一。\end{unitzh}
\begin{uniten}Question: The Kaixi calendar takes the winter solstice of the Heavenly Standard (tianzheng) in the year jiazi, the 4th year of Jiatai, to be 11 days [Note: day sign yihai] 44 ke 61 fen 54 miao, and the regular new moon (jingshuo) of the 11th month to be 1 day [Note: day sign yichou] 75 ke 55 fen 62 miao. What are the intercalary bone (rungu) and the intercalary rate (runlü) each? Answer: Intercalary bone 9 days 69 ke 5 fen 91 miao [Note: not exhausted, 121/169 miao]; intercalary-bone rate (rungu lü) 163771.\end{uniten}

\begin{unitzh}\unitlabel{shushu-011.10}術曰：以日法各通氣、朔日刻分杪，各為氣骨、朔骨分。其氣骨分如約率而一，約盡者為可用。{\zhsmall 或收弃餘分在一刻以下者，亦可用。}然後與朔骨分相減，餘為閏骨率。以日法約之，為閏骨䇿。\end{unitzh}
\begin{uniten}Method: With the day divisor (rifa), convert the days, ke, fen and miao of the qi and of the new moon each, making respectively the qi-bone fen and the new-moon-bone (shuogu) fen. Divide the qi-bone fen by the reduction rate (yuelü); if it reduces exactly, it is usable. [Note: Or if the remaining fen that are rounded up or discarded are below 1 ke, it is also usable.] Then subtract it and the new-moon-bone fen from each other; the remainder is the intercalary-bone rate. Reduce it by the day divisor, making the intercalary-bone tally (rungu ce).\end{uniten}

\begin{unitzh}\unitlabel{shushu-011.11}草曰：置本厯曰法一萬六千九百，先通冬至一十一日四十四刻六十一分五十四杪，得一十九萬三千四百四十分二十六小分，為實。其歴約率係三千一百二十，以約之，得六十二，可用。其餘小分二十六，乃弃之，只用一十九萬三千四百四十為氣骨分。次置朔一日七十五刻五十五分六十二杪，以本厯日法一萬六千九百乘之，得二萬九千六百六十八分九十九杪七十八小分，將近一分，故於氣骨内所弃二十六小分借二十二小分，以補朔内，收上得二萬九千六百六十九，為朔。然後以朔骨分減氣骨分，餘有一十六萬三千七百七十一，為閏骨率。復以日法除之，得閏骨䇿九日六十九刻五分九十一杪，不盡一百二十一筭，直命之為一百六十九分杪之一百二十一。合問。\end{unitzh}
\begin{uniten}Working: Set up the day divisor [Translator's note: the text writes 曰 for 日] of this calendar, 16900. First convert the winter solstice, 11 days 44 ke 61 fen 54 miao, obtaining 193440 fen 26 small fen, as the dividend. The reduction rate of this calendar is 3120; reducing by it gives 62, which is usable. The remaining 26 small fen are then discarded, and only 193440 is used as the qi-bone fen. Next set up the new moon, 1 day 75 ke 55 fen 62 miao, and multiply it by the day divisor of this calendar, 16900, obtaining 29668 fen 99 miao 78 small fen, which is nearly 1 fen; therefore, of the 26 small fen discarded within the qi-bone, borrow 22 small fen to make up [the amount] within the new moon; rounding up gives 29669, as the new moon. Then subtract the new-moon-bone fen from the qi-bone fen, leaving 163771, as the intercalary-bone rate. Again divide it by the day divisor, obtaining the intercalary-bone tally, 9 days 69 ke 5 fen 91 miao, with 121 counting units (suan) not exhausted; designate it directly as 121/169 miao. This answers the question.\end{uniten}

\begin{unitzh}\unitlabel{shushu-011.12}按：此題若置冬至日分，内減經朔日分，餘九日六十九刻五分九十二杪，得閏骨䇿，此原草僅多一百六十九分杪之四十八。盡草中氣骨内棄小分二十六，朔骨分内進二十二，併之為一百六十九分杪之四十八。其不徑相減，而必用通分、約分、累乘、累除者，為向後推筭用耳。\end{unitzh}
\begin{uniten}Note: In this problem, if one sets up the days and fen of the winter solstice and subtracts from it the days and fen of the regular new moon, the remainder, 9 days 69 ke 5 fen 92 miao, gives the intercalary-bone tally; this is only 48/169 miao more than the original working. Altogether, in the working 26 small fen were discarded within the qi-bone and 22 were advanced within the new-moon-bone fen; combined, these make 48/169 miao. That it does not subtract directly but must use conversion to fen, reduction of fen, repeated multiplication and repeated division is merely for use in the later calculations.\end{uniten}

\begin{unitzh}\unitlabel{shushu-011.13}治厯演紀\end{unitzh}
\begin{uniten}Regulating the Calendar and Deriving the Epoch (zhili yanji)\end{uniten}

\begin{unitzh}\unitlabel{shushu-011.14}問：開禧歴積年七百八十四萬八千一百八十三，欲知推演之原，調日法，求朔餘、朔率、斗分、嵗率、嵗閏、入元嵗、入閏、朔定骨、閏泛骨、閏縮、紀率、氣元率、元閏、元數，及氣等率、因率、蔀率、朔等數、因數、蔀數、朔積年，二十三事各幾何？\end{unitzh}
\begin{uniten}Question: The accumulated years of the Kaixi calendar are 7848183. One wishes to know the origin of their derivation: adjusting the day divisor (tiao rifa), find the new-moon remainder (shuoyu), the new-moon rate (shuolü), the Dipper fraction, the year rate (suilü), the year intercalation (suirun), the entering-epoch years (ruyuan sui), the entering intercalation (rurun), the fixed new-moon bone (shuo dinggu), the approximate intercalary bone (run fangu), the intercalary contraction (runsuo), the cycle rate (jilü), the qi-epoch rate (qiyuan lü), the epoch intercalation (yuanrun), the epoch number (yuanshu), and the qi equal rate (qi denglü), the multiplying rate (yinlü), the bu-cycle rate (bulü), the new-moon equal number (shuo dengshu), the multiplying number (yinshu), the bu-cycle number (bushu) and the new-moon accumulated years (shuo jinian): how much is each of these 23 items?\end{uniten}

\begin{unitzh}\unitlabel{shushu-011.15}答曰：日法一萬六千九百；朔餘八千九百六十七；朔率四十九萬九千六十七；斗分四千一百八；嵗率六百一十七萬二千六百八；嵗閏一十八萬三千八百四；入元嵗九千一百八十；入閏四十七萬四千二百六十；朔定骨二萬九千六百六十九；閏泛骨一十六萬三千七百七十一；閏縮一十八萬八千五百七十八；紀率一百一萬四千；氣元率一萬九千五百；元閏三十七萬七千八百七十三；元數四百二；氣等率五十二；因率一百四十四；蔀率三百二十五；朔等數一；因數四十五萬七千九百九十九；蔀數四十九萬九千六十七；朔積年七百八十三萬九千；積年七百八十四萬八千一百八。\end{unitzh}
\begin{uniten}Answer: day divisor 16900; new-moon remainder 8967; new-moon rate 499067; Dipper fraction 4108; year rate 6172608; year intercalation 183804; entering-epoch years 9180; entering intercalation 474260; fixed new-moon bone 29669; approximate intercalary bone 163771; intercalary contraction 188578; cycle rate 1014000; qi-epoch rate 19500; epoch intercalation 377873; epoch number 402; qi equal rate 52; multiplying rate 144; bu-cycle rate 325; new-moon equal number 1; multiplying number 457999; bu-cycle number 499067; new-moon accumulated years 7839000; accumulated years 7848108.\end{uniten}

\begin{unitzh}\unitlabel{shushu-011.16}術曰：以厯法求之，大衍入之。調日法，如何承天術，用强弱母子互乘，得數併之，為朔餘。以二十九日通日法，增入朔餘，為朔率。又以日法乘前厯所測冬至氣刻分，收棄末位，為偶數，得斗分。與日法用大衍術入之，求等數、因率、蔀率。以紀乘等數，為約率。置所求氣定骨，如約率而一，得數以乘因率，滿蔀率去之，不滿以紀法乘之，為入元嵗。次置嵗日，以日法通之，併以斗定分，為嵗率。以十二月乘朔率，減嵗率，餘為嵗閏。以嵗閏乘入元嵗，滿朔率去之，不滿為入閏。與閏骨相減之，得差。{\zhsmall 或適足，便以入元嵗為積年，後術並不用；或差在刻分法半數以下者，亦以入元嵗為積年。}\end{unitzh}
\begin{uniten}Method: Seek it by the calendar method, entering it into the Great Extension (dayan). Adjust the day divisor as in the method of He Chengtian: use the strong and weak denominators and numerators (qiangruo muzi), cross-multiply them, and add the numbers obtained, making the new-moon remainder. Convert 29 days by the day divisor and add the new-moon remainder into it, making the new-moon rate. Further, multiply by the day divisor the ke and fen of the winter-solstice qi measured by the previous calendar, round up or discard the last place, making an even number, and obtain the Dipper fraction. With the day divisor, enter it into the Great Extension method (dayan shu) to find the equal number (dengshu), the multiplying rate and the bu-cycle rate. Multiply the equal number by the cycle (ji), making the reduction rate. Set up the sought fixed qi-bone (qi dinggu), divide it by the reduction rate, and multiply the number obtained by the multiplying rate; when it fills the bu-cycle rate, remove it; what does not fill it, multiply by the cycle divisor, making the entering-epoch years. Next set up the days of the year, convert them by the day divisor, and add the fixed Dipper fraction (dou dingfen) to them, making the year rate. Multiply the new-moon rate by 12 months and subtract it from the year rate; the remainder is the year intercalation. Multiply the entering-epoch years by the year intercalation; when it fills the new-moon rate, remove it; what does not fill it is the entering intercalation. Subtract it and the intercalary bone from each other to obtain the difference. [Note: If it is exactly sufficient, then take the entering-epoch years as the accumulated years, and the subsequent methods are not used at all; or if the difference is below half of the ke-fen divisor (kefen fa), also take the entering-epoch years as the accumulated years.]\end{uniten}

\begin{unitzh}\unitlabel{shushu-011.17}必在刻分法半數以上，却以閏泛骨併朔率，得數内减入閏，餘與朔率求閏縮。{\zhsmall 在朔率以下，便為閏縮；以上，用朔率減之，亦得。}以紀法乘日法，為紀率。以等數約之，為氣元率。以氣元乘嵗閏，滿朔率去之，不滿為元閏。虗置一億，減入元嵗，餘為實，元率除之，得乘限。乃以元閏與朔率用大衍入之，求得等數、因數、蔀數。以等約閏縮，得數以因數乘之，滿蔀數去之，不滿在乘限以下，以乘元率，為朔積年。併入元嵗，為演紀積年。又加成厯年。今人相乘演積年，其術如調日法，求朔餘、朔率，立\end{unitzh}
\begin{uniten}[If the difference] is definitely above half of the ke-fen divisor, then add the approximate intercalary bone to the new-moon rate, subtract the entering intercalation from the number obtained, and with the remainder and the new-moon rate find the intercalary contraction. [Note: If it is below the new-moon rate, it is then the intercalary contraction; if above, subtract the new-moon rate from it, and it is likewise obtained.] Multiply the day divisor by the cycle divisor, making the cycle rate. Reduce it by the equal number, making the qi-epoch rate. Multiply the year intercalation by the qi-epoch [rate]; when it fills the new-moon rate, remove it; what does not fill it is the epoch intercalation. Set up separately (xuzhi) 100000000 (yi), subtract the entering-epoch years, and take the remainder as the dividend; divide it by the epoch rate (yuanlü), obtaining the multiplication limit (chengxian). Then enter the epoch intercalation and the new-moon rate into the Great Extension, finding the equal number, the multiplying number and the bu-cycle number. Reduce the intercalary contraction by the equal [number]; multiply the number obtained by the multiplying number; when it fills the bu-cycle number, remove it; what does not fill it, being below the multiplication limit, multiply the epoch rate by it, making the new-moon accumulated years. Add the entering-epoch years, making the accumulated years of the epoch derivation (yanji jinian). Further add the years of the completed calendar. People today multiply to derive the accumulated years; their method is like adjusting the day divisor: find the new-moon remainder and the new-moon rate, set up [continues]\end{uniten}

\begin{unitzh}\unitlabel{shushu-012.1}斗分、嵗餘，求氣骨、朔骨、閏骨，及衍等數、約率、因率、蔀率，求入元嵗、嵗閏、入閏、元率、元閏，已上皆同此術。但其所以求朔積年之術，乃以閏骨減入閏，餘為之閏贏，却與閏縮、朔率列號甲、乙、丙、丁四位，除乘消減，謂之方程，乃求得元數，以乘元率，所得為之積年，加入元嵗，共為演紀嵗積年。所謂方程，正是大衍術{\zhsmall 今人少知}，非特置筭繁多，初無定法可傳，甚是惑悮後學，易失古人之術意。\end{unitzh}
\begin{uniten}[continued] the dou fraction (dou fen) and the year remainder (sui yu); finding the solar-term bone (qi gu), the new-moon bone (shuo gu) and the intercalation bone (run gu); deriving (yan) the equal number (deng shu), the reducing rate (yue lü), the multiplying rate (yin lü) and the bu-period rate (bu lü 蔀率); and finding the years of entry into the origin (ru yuan sui), the year intercalation (sui run), the entering intercalation (ru run), the origin rate (yuan lü) and the origin intercalation (yuan run): all the above is the same as in this method. But their method for finding the new-moon accumulated years (shuo ji nian) is to subtract the intercalation bone from the entering intercalation, the remainder being called the intercalary surplus (run ying); then, with the intercalary shortfall (run suo) and the new-moon rate (shuo lü), to set them out in four positions labelled jia, yi, bing and ding, and to divide, multiply, eliminate and subtract, which is called rectangular arrays (fangcheng); thus the origin number (yuan shu) is found, and with it the origin rate is multiplied; what is obtained is called the accumulated years (ji nian); adding the years of entry into the origin, the total is the accumulated years of the extended era (yan ji sui ji nian). What is called rectangular arrays is precisely the Great Extension method (dayan shu) [Note: people today seldom know this]; not only is the setting out of the counting rods (zhi suan) complicated and numerous, but from the start there is no fixed rule that can be handed down; it greatly confuses and misleads later students, and the intent of the methods of the ancients is easily lost.\end{uniten}

\begin{unitzh}\unitlabel{shushu-012.2}故今術不言閏贏，而曰入閏差者，蓋本將來可用入元嵗便為積年之意。故今止將元閏、朔率二項，以大衍先求等數、因數、蔀者，仍倣前前求入元嵗之術理：假閏骨如氣骨，以等數為約數，及求乘數、蔀數；以等約閏縮得，因乘數，滿蔀去之，不滿在限下，以乘元率，便為朔積年，亦加入元嵗，共為演紀積年。此術非惟止用乘除省便，又且於自然中取見積年，不惑不差矣。新術敢不用閏贏而求者，實知閏贏已存於入閏之中，但求朔積年之竒分與閏縮等，則自與入閏相合，必滿朔率所去故也。數理精微，不易窺識，窮年致志，感於夢寐，幸而得知，謹不敢隐。\end{unitzh}
\begin{uniten}Therefore the present method does not speak of the intercalary surplus but speaks of the difference of the entering intercalation (ru run cha), for it is based on the idea that in future the years of entry into the origin may be used directly as the accumulated years. Therefore now one takes only the two items, the origin intercalation and the new-moon rate, and by the Great Extension first finds the equal number, the multiplying number (yin shu) and the bu-period; this still imitates the principle of the preceding, earlier method for finding the years of entry into the origin: treat the intercalation bone like the solar-term bone, take the equal number as the reducing number (yue shu), and find the multiplier number (cheng shu) and the bu-period number (bu shu); reduce the intercalary shortfall by the equal (deng); what is obtained, multiply by the multiplier number; when it fills the bu-period, discard it; what does not fill, being below the limit, multiplies the origin rate, and this is then the new-moon accumulated years; likewise add the years of entry into the origin, the total being the accumulated years of the extended era (yan ji ji nian). This method not only uses nothing but multiplication and division and is economical and convenient, but moreover brings out the accumulated years from within what is natural (ziran), without confusion and without error. The new method dares to seek without using the intercalary surplus because it truly knows that the intercalary surplus is already contained within the entering intercalation: one only seeks an odd fraction (ji fen) of the new-moon accumulated years equal to the intercalary shortfall, and then it of itself combines with the entering intercalation, because they must fill the new-moon rate, which is discarded. The principles of numbers are subtle and fine, not easily perceived; I devoted my mind to them for years on end, was stirred by them in my dreams, and fortunately came to know this; respectfully, I dare not conceal it.\end{uniten}

\begin{unitzh}\unitlabel{shushu-012.3}草曰：本厯以何承天術調得一萬六千九百為日法，係三百三十九强、一十七弱。先以强數三百三十九乘强子二十六，得八千八百一十四於上；次以弱數一十七乘弱子九，得一百五十三，併上，共得八千九百六十七，為朔餘。次以日法通朔䇿二十九日，得四十九萬一百，增入朔餘，得四十九萬九千六十七，為朔率。又以日法乘統天厯所測毎嵗冬至周日下二十四刻三十一分，得四千一百八分三十九杪，為斗泛分。驗八分既偶，遂棄三十九杪，只以四千一百八分為斗定分，與日法以大衍術入之，□得五十二為等數，一百四十四為因率，三百二十五為蔀率。以甲子六十為紀法，乘等數，得三千一百二十，為約率。\end{unitzh}
\begin{uniten}Working: This calendar, by He Chengtian's method, adjusted and obtained 16,900 as the day divisor (ri fa); this consists of 339 strong (qiang) and 17 weak (ruo). First, with the strong number (qiang shu) 339 multiply the strong numerator (qiang zi) 26, obtaining 8,814 at the top; next, with the weak number (ruo shu) 17 multiply the weak numerator (ruo zi) 9, obtaining 153; combine with the above, obtaining in all 8,967 as the new-moon remainder (shuo yu). Next, with the day divisor convert (tong) the new-moon tally (shuo ce) of 29 days, obtaining 490,100; add in the new-moon remainder, obtaining 499,067 as the new-moon rate. Further, with the day divisor multiply the 24 ke (刻) 31 fen (分) below the cycle days (zhou ri) of the winter solstice of each year as measured by the Tongtian calendar, obtaining 4,108 fen 39 miao (杪) as the floating dou fraction (dou fan fen). Examining: since the 8 fen is even, discard the 39 miao and take only 4,108 fen as the fixed dou fraction (dou ding fen). With the day divisor enter it into the Great Extension method; [lost character] obtain 52 as the equal number, 144 as the multiplying rate and 325 as the bu-period rate. Take the jiazi 60 as the era divisor (ji fa), and with it multiply the equal number, obtaining 3,120 as the reducing rate.\end{uniten}

\begin{unitzh}\unitlabel{shushu-012.4}却置本厯上課所用嘉泰甲子嵗氣骨一十一日四十四刻六十一分五十一杪，以乘日法，得一十九萬三千四百四十分二十六杪，為氣泛骨。欲滿約率三千一百二十而一，故就近乃棄微杪，只以一十九萬三千四百四十為氣定骨。然後以約率三千一百二十除之，得六十二；以因率一百四十四乘之，得八千九百二十八；滿蔀率三百二十五去之，不滿一百五十三；以紀法六十乘之，得九千一百八十年，為入元嵗。次置嵗□三百六十五，以日法乘之，得六百一十六萬八千五百，併斗定分四千一百八，得六百一十七萬二千六百八，為嵗率。却以十二月乘朔率四十九萬九千六百七，得五百九十八萬八千八百四，率内減去此數，餘一十八萬三千八百四，為嵗閏。以嵗閏乘入元嵗九千一百八十，得一十億八千七百三十二萬七百二十，滿朔率去之，不滿四十七萬四千二百六十，為入閏。\end{unitzh}
\begin{uniten}Then set up the solar-term bone of the jiazi year of Jiatai used in the testing submitted (shang ke) for this calendar, 11 days 44 ke 61 fen 51 miao; with it multiply the day divisor, obtaining 193,440 fen 26 miao as the floating solar-term bone (qi fan gu). Since one wants it to be fully divisible by the reducing rate 3,120, therefore, approximating to the nearest, discard the wei (微) and miao, and take only 193,440 as the fixed solar-term bone (qi ding gu). Then divide it by the reducing rate 3,120, obtaining 62; multiply it by the multiplying rate 144, obtaining 8,928; when it fills the bu-period rate 325, discard it; what does not fill is 153; multiply it by the era divisor 60, obtaining 9,180 years as the years of entry into the origin. Next set up the year [lost character] 365; multiply it by the day divisor, obtaining 6,168,500; combine with the fixed dou fraction 4,108, obtaining 6,172,608 as the year rate (sui lü). Then with 12 months multiply the new-moon rate 499,607, obtaining 5,988,804; from the rate subtract this number; the remainder, 183,804, is the year intercalation. With the year intercalation multiply the years of entry into the origin 9,180, obtaining 1,087,320,720; when it fills the new-moon rate, discard it; what does not fill, 474,260, is the entering intercalation.\end{uniten}

\begin{unitzh}\unitlabel{shushu-012.5}次置本厯所用嘉泰甲子嵗天正十一月朔一日七十五刻五十五分六十二杪，以日法乘之，得二萬九千六百六十八分九千九百七十八杪，為朔泛骨。就近收杪為一分，共得二萬九千六百六十九，為朔定骨數。然後乃以朔定骨減氣骨一十九萬三千四百四十，餘一十六萬三千七百七十一，為閏泛骨。置日法，以二百歸除之，得八十四半，為半刻法。次以入閏數内減去閏泛骨，與入閏相課減之，餘三十一萬四百八十九{\zhsmall 此是閏贏}，為差。半刻法以上，乃以閏泛骨併朔率，共得六十六萬二千八百三十八，以入閏四十七萬四千二百六十減之，餘一十八萬八千五百七十八，在朔率下，便為閏縮。\end{unitzh}
\begin{uniten}Next set up the new moon of the tianzheng eleventh month of the jiazi year of Jiatai used in this calendar, 1 day 75 ke 55 fen 62 miao; multiply it by the day divisor, obtaining 29,668 fen 9,978 miao as the floating new-moon bone (shuo fan gu). Approximating to the nearest, gather the miao into 1 fen, obtaining in all 29,669 as the fixed new-moon bone (shuo ding gu) number. Then subtract the fixed new-moon bone from the solar-term bone 193,440; the remainder, 163,771, is the floating intercalation bone (run fan gu). Set up the day divisor and divide it by 200 by return division (gui chu), obtaining 84 1/2 as the half-ke divisor (ban ke fa). Next, from the entering-intercalation number subtract the floating intercalation bone, testing it against the entering intercalation by subtraction; the remainder, 310,489 [Note: this is the intercalary surplus], is the difference (cha). It is above the half-ke divisor; so combine the floating intercalation bone with the new-moon rate, obtaining in all 662,838; subtract from it the entering intercalation 474,260; the remainder, 188,578, being below the new-moon rate, is then the intercalary shortfall.\end{uniten}

\begin{unitzh}\unitlabel{shushu-012.6}次以紀䇿六十乘日法，得一百一萬四千，為紀率。以等數五十二歸除紀率，得一萬九千五百{\zhsmall 按：即六十乘三百二十五之數，為一蔀年數}，為氣元率。以氣元率乘嵗閏一十八萬三千八百四，得三十五億八千四百一十七萬八千，滿朔率去之，不滿三十七萬七千八百七十三，為元閏。次置一億{\zhsmall 按：此數似虗設，不過取一億之數為限耳。此所求過限，又將改率數以遷就之矣}，以入元嵗元千一百八十減之，餘九千九百九十九萬八百二十，為實；以元率一萬九千五百為法，除之，得五千一二十七，為乘元限數。\end{unitzh}
\begin{uniten}Next, with the era tally (ji ce) 60 multiply the day divisor, obtaining 1,014,000 as the era rate (ji lü). With the equal number 52 divide the era rate by return division, obtaining 19,500 [Siku editors: this is the number of 60 multiplied by 325, the number of years of one bu-period] as the solar-term origin rate (qi yuan lü). With the solar-term origin rate multiply the year intercalation 183,804, obtaining 3,584,178,000; when it fills the new-moon rate, discard it; what does not fill, 377,873, is the origin intercalation. Next set up 100,000,000 [Siku editors: this number seems to be set up arbitrarily; it merely takes the number 100,000,000 as the limit. If what is sought here exceeds the limit, the rate numbers will again be altered to accommodate it.]; subtract from it the years of entry into the origin, 元-thousand 180 [Translator's note: the text has 元 where a digit is expected]; the remainder, 99,990,820, is the dividend (shi); with the origin rate 19,500 as the divisor (fa), divide it, obtaining 5 thousand 1 27 [Translator's note: so the text, 五千一二十七; the numeral is garbled], as the multiplying-origin limit number (cheng yuan xian shu).\end{uniten}

\begin{unitzh}\unitlabel{shushu-012.7}乃以元閏三十七萬七千八百七十三餘，與朔率四十九萬九千六十七，用大衍術求之，得等數一，因等四十五萬七千九百九十九，蔀數四十九萬九千六十七。然後以等數一約閏縮，只得一十八萬八千五百七十八，以因數四十五萬七千九百九十九乘之，得八百六十三億六千八百五十三萬五千四百二十二，滿蔀數四十九萬九千六百七十去之，不滿四百二，在乘元限數以下，為可用。以乘元率一萬九千五百，得七百八十三萬九千年，為數積年；併入元嵗九千一百八十，共得七百八十四萬八千一百八十，為嘉泰四年甲子嵗積算。本厯係於丁夘嵗進呈，又加丁夘三年，共為七百八十四萬八千一百八十三年，為本厯積年。合具繪圖如後。\end{unitzh}
\begin{uniten}Then take the origin intercalation 377,873, the remainder, and the new-moon rate 499,067, and seek with the Great Extension method, obtaining the equal number 1, the multiplying equal (yin deng) 457,999 [Translator's note: so the text; elsewhere 因數, multiplying number], and the bu-period number 499,067. Then reduce the intercalary shortfall by the equal number 1, obtaining just 188,578; multiply it by the multiplying number 457,999, obtaining 86,368,535,422; when it fills the bu-period number 499,670, discard it; what does not fill, 402, is below the multiplying-origin limit number and is usable. With it multiply the origin rate 19,500, obtaining 7,839,000 years as the number (shu) accumulated years [Translator's note: so the text; the method has 朔積年]; combine with the years of entry into the origin 9,180, obtaining in all 7,848,180 as the accumulated count (ji suan) for the jiazi year, the 4th year of Jiatai. This calendar was presented to the throne in the dingmao year; adding a further 3 years for dingmao, the total is 7,848,183 years, as the accumulated years of this calendar. A drawn diagram is duly provided below.\end{uniten}

\begin{unitzh}\unitlabel{shushu-012.8}按：此術草内竒定相求，有等數，又有因數、蔀數之異。蓋等數即度盡定、竒兩數之數；因數為竒數之倍數，任倍定、竒二數相較，但得一等數，則竒之倍數即為因數；蔀數者，竒數最大之倍數也，任倍竒、定，至兩邊相等無較數，則竒數之倍數即謂之蔀數也。等數甚小者，因數不患其甚大，有蔀數以限之也。草中尚多訛舛，正之於後。\end{unitzh}
\begin{uniten}Note: In the working of this method, in the mutual seeking of the odd (ji) and the fixed (ding), there is the equal number, and further there is the distinction between the multiplying number and the bu-period number. For the equal number is the number that measures out exactly (du jin) both numbers, the fixed and the odd; the multiplying number is the multiple (bei shu) of the odd number: multiply the two numbers, the fixed and the odd, by arbitrary multiples and compare them; as soon as one equal number is obtained, the multiple of the odd is the multiplying number. The bu-period number is the largest multiple of the odd number: multiply the odd and the fixed by arbitrary multiples until the two sides are equal and there is no difference (jiao shu); then the multiple of the odd number is called the bu-period number. When the equal number is very small, there is no need to worry that the multiplying number is very large, since there is the bu-period number to limit it. In the working there are still many errors; they are corrected below.\end{uniten}

\begin{unitzh}\unitlabel{shushu-012.9}按：此題術草皆曰何承天調日法，而《宋書》所載何承天法並無甚率，且各用數亦與此不同。今細按其草，日法已有定數，所調者朔䇿餘分也。然從來朔䇿餘分，皆以實測之朔䇿分、嵗實分兩母子互乘相通即得，並無所謂調法。今所載强弱母子四數，大約已有朔䇿餘分與日法分相約而得，非别有所本，乃故設曲折以為竒也。\end{unitzh}
\begin{uniten}Note: The method and the working of this problem both speak of He Chengtian's adjusting the day divisor (tiao ri fa), but the method of He Chengtian recorded in the Song shu has no rates whatsoever, and moreover the numbers each uses also differ from these. Now, examining the working closely, the day divisor already has a fixed number, and what is adjusted is the remainder fraction of the new-moon tally. Yet the remainder fractions of the new-moon tally have always been obtained simply by cross-multiplying and mutually converting (tong) the denominators (mu) and numerators (zi) of the two, the observed new-moon-tally fraction and the year-actual fraction (sui shi fen); there has never been any so-called adjusting method. The four numbers recorded here, the strong and weak denominators and numerators, were in general obtained by mutually reducing the already existing new-moon-tally remainder fraction and the day-divisor fraction; they have no separate basis, but a roundabout course was deliberately set up to make it seem marvellous.\end{uniten}

\begin{unitzh}\unitlabel{shushu-012.10}試以朔餘分八千九百六十七分為第一條；置日法分，内減朔䇿餘分，餘七千九百三十三，為第二條。以此二數數取之：先置第一條，減第二條，餘一千零三十四，為第三條；七因第三條，以減第二條，餘六百九十五，為第四條；以第四條減第三條，餘三百三十九，為第五條；二因第五條，以減第四條，餘一十七，為第六條。是第五條即强母數，第六條即弱母數矣。\end{unitzh}
\begin{uniten}Let us try taking the new-moon remainder fraction, 8,967 fen, as the first item (tiao); set up the day-divisor fen and subtract from them the new-moon-tally remainder fen; the remainder, 7,933, is the second item. Count these two numbers off against each other (shu qu): first set up the first item and subtract the second item; the remainder, 1,034, is the third item; multiply the third item by 7 and subtract it from the second item; the remainder, 695, is the fourth item; subtract the fourth item from the third item; the remainder, 339, is the fifth item; multiply the fifth item by 2 and subtract it from the fourth item; the remainder, 17, is the sixth item. Thus the fifth item is the strong denominator number (qiang mu shu), and the sixth item is the weak denominator number (ruo mu shu).\end{uniten}

\begin{unitzh}\unitlabel{shushu-012.11}次用第五條、第六條轉求第一條，以取兩子數：置第六條於上，二因第五條加之，得第五條者二、第六條者一，共六百九十五，為第四條；以第四條加第五條，得第五條者三、第六條者一，共一千零三十四，為第三條；七因第三條，以加第四條，得第五條者二十三、第六條者八，共七千九百三十三，為第二條；以第二條加第三條，得第五條者二十六、第六條者九，共八千九百六十七，為第一條。是第五條倍數即强子數，第六條倍數即弱子數矣。\end{unitzh}
\begin{uniten}Next use the fifth item and the sixth item to work back to the first item, so as to obtain the two numerator numbers: set the sixth item at the top; multiply the fifth item by 2 and add it, obtaining 2 of the fifth item and 1 of the sixth item, together 695, as the fourth item; add the fourth item to the fifth item, obtaining 3 of the fifth item and 1 of the sixth item, together 1,034, as the third item; multiply the third item by 7 and add it to the fourth item, obtaining 23 of the fifth item and 8 of the sixth item, together 7,933, as the second item; add the second item to the third item, obtaining 26 of the fifth item and 9 of the sixth item, together 8,967, as the first item. Thus the multiple of the fifth item is the strong numerator number, and the multiple of the sixth item is the weak numerator number.\end{uniten}

\begin{unitzh}\unitlabel{shushu-012.12}至算式中以日法取强母數者，則又以第五條、第六條再約而得者。如以第五條為實，第六條為法，為法商之：初商得十，以法乘初商，得一百七十，為初商積，減實，餘一百六十九，恰與百分日法分之一等，故百除日法分為一數，又三乘之，得五百零七，為實；以四十九為法，初商得十，初商積為四百九十，減實，餘十七；乃以商十乘十七，得一百七十，與前初商積等，故加與前餘積等之百分日法分之一，得三百三十九，為五條之數，名之曰强數也。\end{unitzh}
\begin{uniten}As for taking the strong denominator number from the day divisor in the computation (suan shi), it is obtained by reducing again with the fifth item and the sixth item. For instance, take the fifth item as the dividend and the sixth item as the divisor; as divisor, estimate the quotient (shang): the first quotient (chu shang) is 10; multiply the first quotient by the divisor, obtaining 170 as the first-quotient product (chu shang ji); subtract it from the dividend; the remainder, 169, is exactly equal to 1/100 of the day-divisor fen; therefore divide the day-divisor fen by 100 to make one number, and further multiply it by 3, obtaining 507 as the dividend; with 49 as the divisor, the first quotient is 10 and the first-quotient product is 490; subtract it from the dividend; the remainder is 17. Then multiply 17 by the quotient 10, obtaining 170, equal to the earlier first-quotient product; therefore add the 1/100 of the day-divisor fen that is equal to the earlier remaining amount, obtaining 339 as the number of item 5; it is called the strong number.\end{uniten}

\begin{unitzh}\unitlabel{shushu-012.13}次以日法乘朔䇿日，得數併朔餘，為朔率。\end{unitzh}
\begin{uniten}Next, with the day divisor multiply the days of the new-moon tally; combine the number obtained with the new-moon remainder to make the new-moon rate.\end{uniten}

\begin{unitzh}\unitlabel{shushu-012.14}斗分見偶則棄，見竒則收為偶。\end{unitzh}
\begin{uniten}The dou fraction: if it appears even, discard; if it appears odd, gather it up to make it even.\end{uniten}

\begin{unitzh}\unitlabel{shushu-012.15}按：此係棄分以下數不用也。分為偶數，即用其數；分為竒數，則杪微進一分，併為偶數；如無杪微，即加一分。\end{unitzh}
\begin{uniten}Note: This means that the numbers below the fen are discarded and not used. If the fen are an even number, use that number; if the fen are an odd number, then advance the miao and wei into 1 fen, combining to make an even number; if there are no miao or wei, then add 1 fen.\end{uniten}

\begin{unitzh}\unitlabel{shushu-013.1}又按：算中用數，以日法分一萬六千九百分為主，斗分定為四千一百零八，為偶數，氣骨分亦定為偶數。其各時刻分皆由日分比例而得，故變時刻分為日分求之，無不合。\end{unitzh}
\begin{uniten}Further note: the numbers used in the computation take the day-divisor (ri fa) fen, 16900 fen (分), as the basis. The Dipper fen (dou fen 斗分) are fixed as 4108, which is an even number, and the solar-term bone (qi gu 氣骨) fen are also fixed as an even number. Their fen for each double-hour (shi) and ke are all obtained from the day fen by proportion. Therefore, if the double-hour and ke fen are converted into day fen and the computation is made, there is nothing that does not agree.\end{uniten}

\begin{unitzh}\unitlabel{shushu-013.2}子部·天文算法類·算書之屬·《數學九章》·卷二上。子部·天文算法類·算書之屬·《數學九章》·卷二上。\end{unitzh}
\begin{uniten}[Page header: Masters section (zibu) · Astronomy and mathematics (tianwen suanfa lei) · Mathematical books (suanshu zhi shu) · Shuxue jiuzhang · juan 2, part 1.] [Page header: Masters section · Astronomy and mathematics · Mathematical books · Shuxue jiuzhang · juan 2, part 1.]\end{uniten}

\begin{unitzh}\unitlabel{shushu-013.3}按：求入元嵗法，用斗分與日法分求等率、乘率。蓋以六十年之嵗實積分與紀法分相約，後以六十除紀法分，得日法分為定；以六十除嵗實積分，得斗分為竒。求得蔀數、乘數，皆與六十年之嵗實積分與紀法分所求者同，惟等數則為六十之一，故以六十乘之為乘分，以約氣骨分，然後以乘數乘之，滿蔀數去之，所得用數為六十年之周數，故以六十乘之始為年數。此立法之意也。然以六十年為周數，則六十年之間其氣骨數有合者則不可得，故所得年數較以嵗實分、紀分相求者為逺也。\end{unitzh}
\begin{uniten}Note: the rule for finding the years entering the epoch (ru yuan sui) uses the Dipper fen and the day-divisor fen to seek the equal rate (deng lü) and the multiplying rate (cheng lü). For it reduces the accumulated year-actual (sui shi) fen of 60 years and the cycle-divisor (ji fa) fen against each other, and afterwards divides the cycle-divisor fen by 60, obtaining the day-divisor fen as the fixed number (ding), and divides the accumulated year-actual fen by 60, obtaining the Dipper fen as the surplus (ji 竒). The bu number (bu shu 蔀數) and the multiplier (cheng shu) obtained are both the same as those sought from the accumulated year-actual fen of 60 years and the cycle-divisor fen; only the equal number (deng shu) is 1/60. Therefore it is multiplied by 60 to make the multiplied fen (cheng fen), by which the solar-term bone fen are reduced. After that, multiply by the multiplier and cast out the full bu numbers. The use number (yong shu) obtained is a number of 60-year rounds (zhou shu), so only when it is multiplied by 60 does it become the number of years. This is the intention behind setting up the rule. However, if 60 years are taken as the round, then those within the 60 years whose solar-term bone numbers agree cannot be obtained. Therefore the number of years obtained is farther than that sought by using the year-actual fen and the cycle fen with each other.\end{uniten}

\begin{unitzh}\unitlabel{shushu-013.4}子部·天文算法類·算書之屬·《數學九章》·卷二上。子部·天文算法類·算書之屬·《數學九章》·卷二上。子部·天文算法類·算書之屬·《數學九章》·卷二上。\end{unitzh}
\begin{uniten}[Page header: Masters section · Astronomy and mathematics · Mathematical books · Shuxue jiuzhang · juan 2, part 1.] [Page header: Masters section · Astronomy and mathematics · Mathematical books · Shuxue jiuzhang · juan 2, part 1.] [Page header: Masters section · Astronomy and mathematics · Mathematical books · Shuxue jiuzhang · juan 2, part 1.]\end{uniten}

\begin{unitzh}\unitlabel{shushu-013.5}按：氣元一萬九千五百，乃前蔀數三百二十五以六十乘之之數。蓋求入元嵗用六十倍者，故此仍用六十倍也。\end{unitzh}
\begin{uniten}Note: the solar-term epoch (qi yuan 氣元), 19500, is the number obtained by multiplying the previous bu number, 325, by 60. This is because 60-fold was used in finding the years entering the epoch; therefore here 60-fold is still used.\end{uniten}

\begin{unitzh}\unitlabel{shushu-013.6}又按：此皆用六十年嵗實分求得之數，與用一嵗實分求得之數同，蓋因積年數為六十度盡之數；若非六十度盡之數，則得數必逺也。今依其數另設一題，以明其法。\end{unitzh}
\begin{uniten}Further note: these are all numbers obtained using the year-actual fen of 60 years, and they are the same as the numbers obtained using the year-actual fen of one year. This is because the number of accumulated years is a number measured out exactly (du jin) by 60. If it were not a number measured out exactly by 60, the number obtained would certainly be farther. Now, following these numbers, a separate problem is set up to clarify the rule.\end{uniten}

\begin{unitzh}\unitlabel{shushu-013.7}設宋開禧厯：日法一萬六千九百分，嵗實分六百一十七萬二千六百零八法，紀率分一百零一萬四千分，朔率分四十九萬九千零六十七分；嘉泰甲子嵗天正冬至距甲子日子正後十九萬三千四百四十分{\zhsmall 古名氣骨}，十一月朔距甲子日子正後二萬九千六百六十九分{\zhsmall 古名朔骨}。問距歴元甲子子正初刻冬至朔之積年幾何？答曰：七百八十四萬八千一百八十年。\end{unitzh}
\begin{uniten}Suppose that the Kaixi calendar of the Song has: day divisor 16900 fen; year-actual fen 6172608 fa [Translator's note: 法 fa here is unclear; the parallel items end in 分 fen]; cycle-rate (ji lü) fen 1014000 fen; new-moon-rate (shuo lü) fen 499067 fen. In the jiazi year of Jiatai, the tianzheng (天正) winter solstice is 193440 fen after exact zi (zizheng 子正) of a jiazi day [Note: anciently named "solar-term bone"], and the new moon of the 11th month is 29669 fen after exact zi of a jiazi day [Note: anciently named "new-moon bone" (shuo gu)]. Question: how many are the accumulated years (ji nian) from the calendar epoch (li yuan), with the winter solstice and new moon at the first ke of exact zi on jiazi? Answer: 7848180 years.\end{uniten}

\begin{unitzh}\unitlabel{shushu-013.8}法：以紀率為紀定，紀率除嵗率{\zhsmall 即嵗實分}，餘八萬八千六百零八分，為紀竒。依大衍術求至竒一百零三，餘分六百二十四，餘分可以度盡上數，則命六百二十四為等數，一百零三為乘數。又求得竒一千六百二十五，無餘分，則命一千六百二十五為蔀數。乃以等數約氣骨分，得三百一十，以乘數乘之，得三萬一千九百三十，滿蔀數去之，餘一千零五十五，即專以氣骨分求得距厯年之積年數也。\end{unitzh}
\begin{uniten}Rule: take the cycle rate as the cycle fixed number; divide the year rate (sui lü) [Note: that is, the year-actual fen] by the cycle rate; the remainder, 88608 fen, is the cycle surplus. Seek by the Great Extension method (dayan shu) up to the surplus 103, with remaining fen 624; the remaining fen can measure out the number above exactly, so designate 624 as the equal number and 103 as the multiplier. Further, the surplus 1625 is obtained with no remaining fen, so designate 1625 as the bu number. Then reduce the solar-term bone fen by the equal number, obtaining 310; multiply it by the multiplier, obtaining 31930; cast out the full bu numbers, leaving 1055: this is the number of accumulated years from the calendar year, obtained by using the solar-term bone fen alone.\end{uniten}

\begin{unitzh}\unitlabel{shushu-013.9}舊法以斗分{\zhsmall 嵗餘分四千一百零八}為竒，日法為定，求得等數，一紀法六十乘之，以約氣骨，得數以乘數乘之，蔀數除之，餘數又以六千乘之為積年，名入元嵗。其術未密{\zhsmall 詳前}，故所得積年為九千一百八十，其數亦較逺也。\end{unitzh}
\begin{uniten}The old rule takes the Dipper fen [Note: the year-remainder fen, 4108] as the surplus and the day divisor as the fixed number, and finds the equal number; it multiplies this by one cycle divisor, 60, and reduces the solar-term bone by it; the number obtained is multiplied by the multiplier and divided by the bu number; the remainder is further multiplied by 6000 [Translator's note: 六千 6000 is evidently an error for 六十 60] to make the accumulated years, named the years entering the epoch. Its method is not precise [Note: details above], so the accumulated years obtained are 9180, and this number too is comparatively far.\end{uniten}

\begin{unitzh}\unitlabel{shushu-013.10}次以蔀數{\zhsmall 即嵗實紀法滿一會年數}乘嵗率，得一百億三千四十八萬八千，滿朔率去之，餘二十三萬九千四百三十四，舊名氣元閏，為朔竒，朔率為朔定。依前法求得等數一，乗數六千二百五十一，蔀率四十九萬九千零六十七。次以前所得積年乗嵗率，滿朔率去之，餘二十七萬五千二百二十四，為前朔距至前分數，舊名入元閏。以嘉泰甲子氣骨、朔骨相減，得十六萬三千七百七十一，為後朔距至前分數，舊名閏骨。\end{unitzh}
\begin{uniten}Next, multiply the year rate by the bu number [Note: that is, the number of years in which the year-actual and the cycle divisor complete one meeting (hui)], obtaining 10030488000; cast out the full new-moon rates, leaving 239434, formerly named "solar-term-epoch intercalary" (qi yuan run 氣元閏), as the new-moon surplus; the new-moon rate is the new-moon fixed number. By the previous rule obtain the equal number 1, the multiplier 6251, and the bu rate (bu lü) 499067. Next, multiply the year rate by the accumulated years obtained before; cast out the full new-moon rates, leaving 275224, which is the number of fen by which the earlier new moon is distant before the solstice, formerly named "entering-epoch intercalary" (ru yuan run 入元閏). Subtract the solar-term bone and the new-moon bone of the jiazi year of Jiatai from each other, obtaining 163771, which is the number of fen by which the later new moon is distant before the solstice, formerly named "intercalary bone" (run gu 閏骨).\end{uniten}

\begin{unitzh}\unitlabel{shushu-013.11}夫十一月朔常在冬至前退行，今前逺後近，是已退過一朔䇿，則於後閏骨内加一朔䇿，再減去入元閏，餘三十八萬七千六百一十四，為後朔前朔相差之分數，舊名閏縮。乃以等數約閏縮，仍得原數，以乘數乘之，滿蔀數去之，餘四千八百二十九，為會數。乃以一會年數{\zhsmall 即前蔀數}乗之，得七百八十四萬七千一百二十五，為朔積年，加入前積年，得七百八十四萬八千一百八十，為嘉泰甲子積年。\end{unitzh}
\begin{uniten}Now, the new moon of the 11th month always lies before the winter solstice, moving backward. Here the earlier is far and the later is near: this means it has already moved back past one new-moon tally (shuo ce 朔䇿). So add one new-moon tally into the later intercalary bone, then subtract the entering-epoch intercalary, leaving 387614, which is the number of fen by which the later new moon and the earlier new moon differ, formerly named "intercalary shrinkage" (run suo 閏縮). Then reduce the intercalary shrinkage by the equal number, still obtaining the original number; multiply it by the multiplier; cast out the full bu numbers, leaving 4829, which is the meeting number (hui shu). Then multiply it by the number of years of one meeting [Note: that is, the previous bu number], obtaining 7847125, which is the new-moon accumulated years; add in the previous accumulated years, obtaining 7848180, which is the accumulated years of the jiazi year of Jiatai.\end{uniten}

\begin{unitzh}\unitlabel{shushu-013.12}又法：仍按本法求之。先以嵗率、紀率、朔率求總等，朔率不盡，無總等，各率朔即為各元數。次連環求等，朔元不盡嵗元等數，等數為六百二十四，留嵗元不約，約紀元得一千六百二十五分，為紀泛定，嵗元、朔元即為嵗泛定、朔泛定。次求續等，紀泛定、嵗泛定等數為十三，約嵗泛定，乗紀泛定，得四十七萬四千八百一十六為嵗定，二萬一千一百二十五為紀定，朔泛定即朔定。\end{unitzh}
\begin{uniten}Another rule: still seek it according to the original rule. First seek the total equal (zong deng) of the year rate, the cycle rate and the new-moon rate: the new-moon rate is not exhausted (bu jin); there is no total equal; each rate, shuo [Translator's note: 朔 shuo here appears to be intrusive], is then the respective original number (yuan shu). Next seek the equals in linked rings (lianhuan qiu deng): the new-moon original number is not exhausted by the equal number of the year original number; the equal number is 624; keep the year original number unreduced and reduce the cycle original number, obtaining 1625 fen, as the cycle provisional fixed number (fan ding); the year original number and the new-moon original number are then the year provisional fixed number and the new-moon provisional fixed number. Next seek the continued equal (xu deng): the equal number of the cycle provisional fixed number and the year provisional fixed number is 13; reduce the year provisional fixed number by it and multiply the cycle provisional fixed number by it, obtaining 474816 as the year fixed number and 21125 as the cycle fixed number; the new-moon provisional fixed number is the new-moon fixed number.\end{uniten}

\begin{unitzh}\unitlabel{shushu-013.13}三定數連乗，得五○○五八八五五五四六九六○○○，為衍母。紀定、朔定相乗，得八一○九八三八七五，為嵗衍；嵗定、朔定相乗，得二三六九六四九九六六七二，為紀衍；嵗定、紀定相乗，得一○○三○四八八○○○，為朔衍。置各衍數，滿各定數去之，餘嵗竒四七二九六三，紀竒二○○四七，朔竒二三九四三四。\end{unitzh}
\begin{uniten}Multiply the three fixed numbers together, obtaining 5005885554696000, as the extension mother (yan mu 衍母). Multiply the cycle fixed number and the new-moon fixed number together, obtaining 810983875, as the year extension (sui yan); multiply the year fixed number and the new-moon fixed number together, obtaining 236964996672, as the cycle extension; multiply the year fixed number and the cycle fixed number together, obtaining 10030488000, as the new-moon extension. Set out each extension number and cast out the full respective fixed numbers; the remainders are: year surplus 472963, cycle surplus 20047, new-moon surplus 239434.\end{uniten}

\begin{unitzh}\unitlabel{shushu-013.14}次以各定、各竒求各乗數，得嵗乗數{\zhsmall 此題不用嵗乗數，求之以備其數}二四九六七，紀乗數二○○○八，朔乗數六二五一。以各乘數乗各衍數，得各泛用數：嵗泛用二○一九八九三二○七九四六二五，紀泛用四七四一一九五六五四四一三三七六，朔泛用六二七○○五八□四八八○○○。併三泛用，與衍母數等，則泛用即為定用。\end{unitzh}
\begin{uniten}Next, from each fixed number and each surplus seek each multiplier, obtaining: year multiplier [Note: this problem does not use the year multiplier; it is sought so as to complete its numbers] 24967, cycle multiplier 20008, new-moon multiplier 6251. Multiply each extension number by each multiplier, obtaining the provisional use numbers (fan yong shu): year provisional use 201989320794625, cycle provisional use 4741195654413376, new-moon provisional use 6270058[blank]488000. The three provisional uses added together equal the extension-mother number, so the provisional uses are the fixed uses (ding yong).\end{uniten}

\begin{unitzh}\unitlabel{shushu-013.15}乃以氣骨乗紀定用，得九一七一三六八八七一九六二八三四五五三四四○。置氣骨，減去朔骨，餘十六萬三千七百七十一，以乘朔定用，得一○二六八五三六七六七一○○二四八○○○。併數得九二七四○五四二三九六三三八三七○一四四○，滿衍母去之，餘四八四四三七三八六五三四四○，為實，以嵗實分為法除之，得七百八十四萬八千一百八十，即嘉泰甲子積年之數也。此法較前法數繁，然其理可互相發明。後復設一法，兼二法用之。\end{unitzh}
\begin{uniten}Then multiply the cycle fixed use by the solar-term bone, obtaining 9171368871962834553440. Set out the solar-term bone and subtract the new-moon bone, leaving 163771; multiply the new-moon fixed use by it, obtaining 10268536767100248000. Adding the numbers gives 927405423963383701440; cast out the full extension mothers, leaving 48443738653440, which is the dividend (shi); take the year-actual fen as the divisor (fa) and divide, obtaining 7848180: this is the number of accumulated years of the jiazi year of Jiatai. This rule is more cumbersome in its numbers than the previous rule, but their principles can illuminate each other. Below, another rule is set up, which uses the two rules combined.\end{uniten}

\begin{unitzh}\unitlabel{shushu-013.16}三法：先以嵗率、紀率求等數，得六百二十四，專約紀率，得一六二五分為紀元，嵗率即為嵗元。又求續等數，得十三，以約嵗元，得四七四八一六為嵗定；以乘紀元，得二一一二五為紀定。紀定、嵗定相乘，得一○○三○四八八○○○為衍母。以紀定二一一二五為嵗衍，以嵗定四七四八一六為紀衍。嵗衍小於嵗定，即為嵗竒；紀衍滿紀定去之，餘一○○六六為紀竒。\end{unitzh}
\begin{uniten}Third rule: first seek the equal number of the year rate and the cycle rate, obtaining 624; reduce only the cycle rate, obtaining 1625 fen as the cycle original number; the year rate is then the year original number. Further seek the continued equal number, obtaining 13; reduce the year original number by it, obtaining 474816 as the year fixed number; multiply the cycle original number by it, obtaining 21125 as the cycle fixed number. Multiply the cycle fixed number and the year fixed number together, obtaining 10030488000 as the extension mother. Take the cycle fixed number, 21125, as the year extension, and the year fixed number, 474816, as the cycle extension. The year extension is smaller than the year fixed number, so it is the year surplus; cast out the full cycle fixed numbers from the cycle extension, leaving 10066 as the cycle surplus.\end{uniten}

\begin{unitzh}\unitlabel{shushu-013.17}次以各定、各竒求各乘數，得嵗乗數二五四八六一，紀乗數九七八六。以各乗數乗各衍數，得各用數：嵗用數五三八三九三八六二五，紀用數四六四六五四九三七六。乃以氣骨乗紀用數，得八九八八二八五一一二九三四四○，滿衍母去之，餘六五一二一○一四四○為通積分，為實，以嵗實分為法除之，得一千零五十五，即專以氣骨求\end{unitzh}
\begin{uniten}Next, from each fixed number and each surplus seek each multiplier, obtaining: year multiplier 254861, cycle multiplier 9786. Multiply each extension number by each multiplier, obtaining the use numbers: year use number 5383938625, cycle use number 4646549376. Then multiply the cycle use number by the solar-term bone, obtaining 898828511293440; cast out the full extension mothers, leaving 6512101440 as the total accumulated fen (tong ji fen), which is the dividend; take the year-actual fen as the divisor and divide, obtaining 1055: that is, seeking by the solar-term bone alone [Translator's note: the sentence breaks off here at a page marker.]\end{uniten}

\begin{unitzh}\unitlabel{shushu-013.18}子部·天文算法類·算書之屬·《數學九章》·卷二上。\end{unitzh}
\begin{uniten}[Page header: Masters section · Astronomy and mathematics · Mathematical books · Shuxue jiuzhang · juan 2, part 1.]\end{uniten}

\begin{unitzh}\unitlabel{shushu-013.19}次以前衍母分{\zhsmall 即嵗率紀率一會一千六百二十五年之積分}一○○三○四八八○○○與朔率分求等數，得一，即以前衍衍母分為嵗紀元，亦即為嵗紀定；以朔率分四九九○六七為朔元，亦即為朔定。二定數相乗，得五○○五八八五五五四六九六○○○為衍母。以朔定為嵗紀衍，以嵗紀定為朔衍。嵗紀衍小於嵗紀定，即以嵗紀衍為嵗紀竒；明衍滿朔定去之，餘二三九四三四為朔竒。各以定、竒求乗數，得嵗紀乗數九九○四八五二四○三，朔乗數六二五一。以各乗數乗各衍數，得嵗紀泛用數四九四三一八四九七四二○八○○一，朔泛用數六二七○○五八○四八八○○○。併二泛數，與衍母等，則泛用數即為定用數。\end{unitzh}
\begin{uniten}Next, seek the equal number of the previous extension-mother fen [Note: that is, the accumulated fen of the 1625 years of one meeting of the year rate and the cycle rate], 10030488000, and the new-moon-rate fen, obtaining 1. Then take the previous extension, the extension-mother fen, as the year-cycle original number, which is also the year-cycle fixed number; take the new-moon-rate fen, 499067, as the new-moon original number, which is also the new-moon fixed number. Multiply the two fixed numbers together, obtaining 5005885554696000 as the extension mother. Take the new-moon fixed number as the year-cycle extension, and the year-cycle fixed number as the new-moon extension. The year-cycle extension is smaller than the year-cycle fixed number, so take the year-cycle extension as the year-cycle surplus; cast out the full new-moon fixed numbers from the ming [Translator's note: 明 ming is evidently an error for 朔 shuo] extension, leaving 239434 as the new-moon surplus. From each fixed number and surplus seek the multipliers, obtaining: year-cycle multiplier 9904852403, new-moon multiplier 6251. Multiply each extension number by each multiplier, obtaining: year-cycle provisional use number 4943184974208001, new-moon provisional use number 62700580488000. The two provisional numbers added together equal the extension mother, so the provisional use numbers are the fixed use numbers.\end{uniten}

\begin{unitzh}\unitlabel{shushu-013.20}乃置前通積分六五一二一○一四四○，滿朔率去之，餘二七五二二四，為入元第一千零五十五年之閏分。又置嘉泰甲子氣骨，減去朔骨，餘一六三七七一，為嘉泰甲子\end{unitzh}
\begin{uniten}Then set out the previous total accumulated fen, 6512101440; cast out the full new-moon rates, leaving 275224, which is the intercalary fen (run fen) of the 1055th year after entering the epoch. Further, set out the solar-term bone of the jiazi year of Jiatai and subtract the new-moon bone, leaving 163771, as the Jiatai jiazi [continues]\end{uniten}

\begin{unitzh}\unitlabel{shushu-014.1}之閏分。閏分每嵗漸加，今後數小於前數，是知已加過一朔率，乃於後閏分内加一朔率分，減去前閏分，得三八七六一四，為前後閏分差。以乗朔定用數，得二四三○三六二二八○五二七五六三二○○○，滿衍母去之，餘四八四三七二二六五五二○○○，為實。以嵗率為法收之，得七百八十四萬七千一百二十五，為後積年數。並前積年數，共得七百八十四萬八千一百八十年，為嘉泰甲子積年，與前數合。\end{unitzh}
\begin{uniten}[continued] intercalary fen (runfen). The intercalary fen increase gradually every year. Now the later number is smaller than the earlier number, so one knows that it has already increased past one lunation rate (shuolü). Therefore add the fen (分) of one lunation rate to the later intercalary fen and subtract the earlier intercalary fen, obtaining 387614, which is the difference between the earlier and later intercalary fen. Multiply the definite usable number (ding yongshu) of the lunation by it, obtaining 24303622805275632000. Discard the extended mother (yanmu) as often as the number fills it; the remainder, 48437226552000, is the dividend (shi). With the year rate (suilü) as divisor (fa), collect it (shou), obtaining 7847125, which is the later number of accumulated years. Combined with the earlier number of accumulated years, it gives altogether 7848180 years. This is the accumulated years (jinian) of the Jiatai jiazi year, and it agrees with the earlier number.\end{uniten}

\begin{unitzh}\unitlabel{shushu-014.2}按：右竒定相求，其上層竒一數，即大衍術中所謂立天元一也；其逐層數，即術中所謂遞互乗餘也；其下層竒得數，即術中所謂乗數也{\zhsmall 有等數者求蔀數}。古無筆算，舊式所載不詳，兼多重複舛偽之處。集中惟此問甚繁，故既設題以明其法，復備録加減乗除之數，以詳其算式，俾觀者易見焉。\end{unitzh}
\begin{uniten}Note: In the above mutual seeking with remainders (qi) and definite numbers (ding), the number "remainder 1" in the upper layer is what the Great Extension method (dayan shu) calls "setting up the celestial unknown (tianyuan)" [Translator's note: in the Great Extension procedure this is the unit 1 placed at the upper left of the counting board, not an unknown]. The numbers of the successive layers are what the method calls "successive mutual multiplication and remainders" (dihu cheng yu). The number obtained with the remainder in the bottom layer is what the method calls the multiplier (chengshu) [Note: where there is an equal number (dengshu), one finds the cycle number (bushu)]. In antiquity there was no brush-written calculation. What the old forms record is not detailed, and moreover there are many places of repetition and error. In the collection only this problem is very complicated. Therefore, having set out the problem to make its method clear, we have further recorded in full the numbers of the additions, subtractions, multiplications and divisions, to give its calculation forms in detail, so that those who look at it may see it easily.\end{uniten}

\begin{unitzh}\unitlabel{shushu-014.3}綴術推星\end{unitzh}
\begin{uniten}Calculating a Planet by the Linking Method (zhuishu)\end{uniten}

\begin{unitzh}\unitlabel{shushu-014.4}問：嵗星合伏，經一十六日九十分，行三度九十分，去曰一十三度，乃見。後順行一百一十三日，一十十七度八十三分，乃留。欲知合伏段、晨疾初段常度、初行率、末行率、平行率各幾何？\end{unitzh}
\begin{uniten}Question: The Year Star (suixing), in conjunction-occultation (hefu), passes through 16 days 90 fen and moves 3 du (度) 90 fen. At 13 du from the sun [Translator's note: the text has 曰 "say" for 日 "sun" here, and for 日 "day" in the following note] it then appears. Afterwards it moves direct (shunxing) for 113 days, "one ten ten seven" du 83 fen [Translator's note: so the text, 一十十七, a garbled number; the working below has 17 du 83 fen], and then becomes stationary (liu). We wish to know, for the conjunction-occultation stage (duan) and for the first stage of morning fast motion (chenji chuduan), how much each of the following is: the constant degrees (changdu), the initial motion rate (chuxing lü), the final motion rate (moxing lü) and the mean motion rate (pingxing lü).\end{uniten}

\begin{unitzh}\unitlabel{shushu-014.5}按：此以兩積曰之遞差積度求各行率也。蓋合伏初日其行最疾，以次漸遲，遲極則留。總其積度，略如遞減差分，故古法皆以其術歩之。\end{unitzh}
\begin{uniten}Note: This uses the accumulated degrees (jidu), with successive differences (dicha), of the two accumulated day counts to find each motion rate. For on the first day of conjunction-occultation its motion is fastest. In order it gradually becomes slower, and when the slowness reaches its limit it becomes stationary. Taken as a whole, its accumulated degrees are roughly like distribution with successively decreasing differences (dijian chafen); therefore the ancient methods all computed it (bu) by that procedure.\end{uniten}

\begin{unitzh}\unitlabel{shushu-014.6}答曰：合伏一十六日九十分，常度三度九十度，初行率二十三分九十七杪，平行率二十三分二杪，末行率二十二分七杪。晨疾初三十日，常度六度一十三分，初行率二十一分九十六杪，平行率二十分三十三杪，末行率一十八分七十杪。\end{unitzh}
\begin{uniten}Answer: Conjunction-occultation: 16 days 90 fen; constant degrees 3 du 90 du; initial motion rate 23 fen 97 miao (杪); mean motion rate 23 fen 2 miao; final motion rate 22 fen 7 miao. First stage of morning fast motion: 30 days; constant degrees 6 du 13 fen; initial motion rate 21 fen 96 miao; mean motion rate 20 fen 33 miao; final motion rate 18 fen 70 miao.\end{uniten}

\begin{unitzh}\unitlabel{shushu-014.7}術曰：以方程法求之。置見日，減一，餘半之，為見率。以伏日併見日，為初行法。以法半之，如見率，共為伏率。以伏日乘伏率，為伏差；以見日乘見率，為見差。以伏日乘見差於上，以見日乘伏差，減上，餘為法。以見日乘伏度，為泛；以伏日乘見度，減泛，餘為實。實滿法而一為度；不滿，退除為分杪，即得日差。\end{unitzh}
\begin{uniten}Method: Solve it by the rectangular arrays (fangcheng) method. Set out the appearance days (xianri), subtract 1, and halve the remainder: this is the appearance rate (xianlü). Combine the hidden days (furi) with the appearance days: this is the initial-motion divisor (chuxing fa). Halve the divisor; it and the appearance rate together make the hidden rate (fulü). Multiply the hidden rate by the hidden days: this is the hidden difference (fucha). Multiply the appearance rate by the appearance days: this is the appearance difference (xiancha). Multiply the appearance difference by the hidden days and place it at the top. Multiply the hidden difference by the appearance days and subtract it from the top; the remainder is the divisor. Multiply the hidden degrees (fudu) by the appearance days: this is the provisional (fan). Multiply the appearance degrees (xiandu) by the hidden days and subtract it from the provisional; the remainder is the dividend. The dividend, each time it fills the divisor, gives 1 du; what does not fill is divided with shifting down (tuichu) into fen and miao. This gives the daily difference (richa).\end{uniten}

\begin{unitzh}\unitlabel{shushu-014.8}按：此求逐日之遞差為日差也。術曰方程，非也。其所謂見數者，乃徒設一數，宛轉附會，使合於方程之行列也。如見日減一，折半，為見率；併伏見日，折半，為半總日。既以半總日加見率，先以伏日乘之，後以見日乗之；復置見率，先以見日乘之，後以復日乘之；相減，然後為法。豈非半總日不用加見率，但以伏日見日連乘之，即可為法乎？特多立名目，故為曲折顛倒，使人不易辦耳。今去其見率，另為歩算於後，以明其立法之本意焉。\end{unitzh}
\begin{uniten}Note: This finds the successive day-by-day difference, which is the daily difference. The method says "rectangular arrays"; that is wrong. What it calls the appearance number is merely a number set up to no purpose, twisted round and forced to fit the columns and rows (hanglie) of the rectangular arrays. For instance, the appearance days less 1, halved, make the appearance rate, and the hidden and appearance days combined and halved make the half total days. Having added the half total days to the appearance rate, one first multiplies by the hidden days and afterwards by the appearance days. One again sets out the appearance rate and first multiplies it by the appearance days and afterwards by the 復 days [Translator's note: 復 "again" evidently stands for 伏 "hidden"]. One subtracts the two from each other, and only then has the divisor. Is it not so that, without adding the appearance rate to the half total days, one need only multiply them successively by the hidden days and the appearance days, and that can serve as the divisor? It merely sets up many names and headings, and so deliberately makes it roundabout and topsy-turvy, so that people cannot easily make it out. Now we remove its appearance rate and make a separate computation (busuan) below, to make clear the original intention behind the establishment of the method.\end{uniten}

\begin{unitzh}\unitlabel{shushu-014.9}法：以合伏日除伏行度，得二十三分{\zhsmall 七六九二三}，為合伏日折中第八日四十五分一日之行度{\zhsmall 即第七日九十五分至第八日九十五分之行度}。以順行日除順行度，得十五分{\zhsmall 七七八七六一}，為順行日折中第五十六日五十分一日之行度。兩一日之行度相減，行七分{\zhsmall 一九八一六二}，為合伏第八日四十五分與順行第五十六日五十分兩一日之行度較，為實。併合伏順行兩日數而半之，得六十四日{\zhsmall 九五}，為合伏第八日四十五分至順行五十六日五十分之積日，為法。除之，得十一杪{\zhsmall 二三六五八五}，為一日遞差之數，即日差。若不先用除，則以兩日數與兩行度互乘，相減，為實；兩日數相乘，又併兩日數而半之，再乘，為法；得數亦同。\end{unitzh}
\begin{uniten}Rule: Divide the degrees of hidden motion by the conjunction-occultation days, obtaining 23 fen [Note: 76923] [Translator's note: small-print digits such as these continue the preceding number in further decimal places]. This is the motion in one day at the midpoint (zhezhong) of the conjunction-occultation days, the 8th day 45 fen [Note: that is, the motion from the 7th day 95 fen to the 8th day 95 fen]. Divide the degrees of direct motion by the days of direct motion, obtaining 15 fen [Note: 778761]. This is the motion in one day at the midpoint of the direct-motion days, the 56th day 50 fen. Subtract the two one-day motions from each other: the motion of 7 fen [Note: 198162] is the difference between the two one-day motions at the 8th day 45 fen of conjunction-occultation and at the 56th day 50 fen of direct motion; it is the dividend. Combine the two numbers of days, of conjunction-occultation and of direct motion, and halve, obtaining 64 days [Note: 95]. This is the accumulated days from the 8th day 45 fen of conjunction-occultation to 56 days 50 fen of direct motion; it is the divisor. Dividing, one obtains 11 miao [Note: 236585], which is the number of the successive difference for one day, that is, the daily difference. If one does not first use division, then cross-multiply the two numbers of days and the two motion degrees and subtract the products from each other to make the dividend. Multiply the two numbers of days together, further combine the two numbers of days and halve, and multiply by this again, to make the divisor. The result obtained is likewise the same.\end{uniten}

\begin{unitzh}\unitlabel{shushu-014.10}求初行率：置初法，減一，餘乘日差，為寄。以半初行法乘寄，得數，又加伏見度，共為初行實。以法退除之，得合伏日初行率。\end{unitzh}
\begin{uniten}To find the initial motion rate: set out the initial divisor, subtract 1, and multiply the remainder by the daily difference: this is the lodged number (ji). Multiply the lodged number by half the initial-motion divisor. To the number obtained further add the hidden and appearance degrees; together this is the initial-motion dividend (chuxing shi). Divide it by the divisor with shifting down, obtaining the initial motion rate for the conjunction-occultation days.\end{uniten}

\begin{unitzh}\unitlabel{shushu-014.11}按：此求合伏第一日最疾之行也。其法即遞減差分，有總數，有次數，有毎次差數，求初次最大之數也。初行法減一乘日差為寄者，合伏初日與順行末日兩行率之差也。半法乗寄，與積差等，故加共度為實，以共日為法除之，為合伏初日行率二十三分九十七杪也。\end{unitzh}
\begin{uniten}Note: This finds the fastest motion, that of the first day of conjunction-occultation. Its method is that of distribution with successively decreasing differences: there is a total number, there is a number of terms (cishu), there is the difference for each term, and one seeks the largest number, that of the first term. The initial-motion divisor less 1, multiplied by the daily difference, makes the lodged number: this is the difference between the two motion rates of the first day of conjunction-occultation and of the last day of direct motion. Half the divisor multiplied by the lodged number equals the accumulated difference (jicha). Therefore add the combined degrees to make the dividend, and divide by the combined days as divisor. This gives the first-day motion rate of conjunction-occultation, 23 fen 97 miao.\end{uniten}

\begin{unitzh}\unitlabel{shushu-014.12}求未行率：以段日乘日差，減初行率，餘為末行率。\end{unitzh}
\begin{uniten}To find the final (未) motion rate [Translator's note: 未 "not yet" evidently stands for 末 "final", here and in the rule for the mean motion rate]: multiply the daily difference by the days of the stage (duanri) and subtract from the initial motion rate; the remainder is the final motion rate.\end{uniten}

\begin{unitzh}\unitlabel{shushu-014.13}按：此求合伏末日之行率也。以段日乘日差，求合伏初末日兩行率之較也。既得初末日兩行率之較，以減初行率，即末行率也。\end{unitzh}
\begin{uniten}Note: This finds the motion rate of the last day of conjunction-occultation. Multiplying the daily difference by the days of the stage finds the difference between the two motion rates of the first and last days of conjunction-occultation. Having obtained the difference between the two motion rates of the first and last days, subtract it from the initial motion rate, and that is the final motion rate.\end{uniten}

\begin{unitzh}\unitlabel{shushu-014.14}求平行率：以初行率併未行率而半之，為平行率。\end{unitzh}
\begin{uniten}To find the mean motion rate: combine the initial motion rate with the final (未) motion rate and halve; this is the mean motion rate.\end{uniten}

\begin{unitzh}\unitlabel{shushu-014.15}按：此即均分合伏度為毎日之平行率也，與遞加遞減有首尾數求中數者同。應與伏日除伏度數同；不同者，本非遞差之數也。\end{unitzh}
\begin{uniten}Note: This divides the conjunction-occultation degrees evenly into the mean motion rate of each day. It is the same as cases of successive increase or successive decrease where the first and last numbers are known and the middle number is sought. It ought to be the same as the number from dividing the hidden degrees by the hidden days; that it is not the same is because these are originally not numbers with successive differences.\end{uniten}

\begin{unitzh}\unitlabel{shushu-014.16}求交段差：以各段常日下分數，減全日一百分，餘乗末日行率，為交段差。\end{unitzh}
\begin{uniten}To find the stage-junction difference (jiaoduan cha): subtract the number of fen below the days in the constant days (changri) of each stage from a whole day of 100 fen. Multiply the remainder by the motion rate of the last day; this is the stage-junction difference.\end{uniten}

\begin{unitzh}\unitlabel{shushu-014.17}按：此即各段日下分數不及一日所差之行分也。求之，以備後數加減。\end{unitzh}
\begin{uniten}Note: This is the motion in fen for the amount by which the fen below the days of each stage fall short of one day. It is found in preparation for the additions and subtractions of the later numbers.\end{uniten}

\begin{unitzh}\unitlabel{shushu-014.18}累減前段積度，以益後段積度，各為常度。\end{unitzh}
\begin{uniten}Cumulatively subtract it from the accumulated degrees of the preceding stage so as to add it to the accumulated degrees of the following stage; each is then the constant degrees.\end{uniten}

\begin{unitzh}\unitlabel{shushu-014.19}按：常度即各段積度也。求晨疾初段常度，見常中，專解於後。\end{unitzh}
\begin{uniten}Note: The constant degrees are the accumulated degrees of each stage. For finding the constant degrees of the first stage of morning fast motion, see under "constant" (chang zhong); it is explained specially below.\end{uniten}

\begin{unitzh}\unitlabel{shushu-014.20}草曰：以伏日隨伏度為右行，以見日隨見度為左行，以度對度，日對日，其度於上，日於中，空其下，列之。\end{unitzh}
\begin{uniten}Working: Take the hidden days together with the hidden degrees as the right column (hang), and the appearance days together with the appearance degrees as the left column, degrees facing degrees and days facing days. Set them out with the degrees at the top, the days in the middle, and the bottom left empty.\end{uniten}

\begin{unitzh}\unitlabel{shushu-014.21}置見日一百一十三，減一，餘一百一十二，以半之，得五十六，為見率。以伏日一十六日九十分，併見日一百一十三，得一百二十九日九十分，為初行法。\end{unitzh}
\begin{uniten}Set out the appearance days, 113; subtract 1, leaving 112; halve it, obtaining 56, which is the appearance rate. Take the hidden days, 16 days 90 fen, and combine with the appearance days, 113, obtaining 129 days 90 fen, which is the initial-motion divisor.\end{uniten}

\begin{unitzh}\unitlabel{shushu-014.22}以初行法半之，得六十四日九十五分，併見率五十六日，得一百二十日九十五分，為伏率。以初行法寄之，以伏率歸右下，以對見率，仍分左右兩行，為首圖。\end{unitzh}
\begin{uniten}Halve the initial-motion divisor, obtaining 64 days 95 fen; combine with the appearance rate, 56 days, obtaining 120 days 95 fen, which is the hidden rate. Lodge (ji) the initial-motion divisor. Put the hidden rate at the lower right, facing the appearance rate, still dividing them into left and right columns: this is the first diagram (shoutu).\end{uniten}

\begin{unitzh}\unitlabel{shushu-014.23}以首圖伏日一十六日九十分，乗伏率一百二十日九十五分，得二千四十四日五分五十杪，為伏差，於右下；以首圖見日一百一十三，乘見率五十六，得六十三百二十八日，為見差，於左下；乃成次圖。\end{unitzh}
\begin{uniten}With the hidden days of the first diagram, 16 days 90 fen, multiply the hidden rate, 120 days 95 fen, obtaining 2044 days 5 fen 50 miao, which is the hidden difference, at the lower right. With the appearance days of the first diagram, 113, multiply the appearance rate, 56, obtaining 63 hundred 28 days, which is the appearance difference, at the lower left. This then forms the second diagram (citu).\end{uniten}

\begin{unitzh}\unitlabel{shushu-014.24}凡方程之術，先欲得者存之，以未欲得者互偏乘兩行諸數。今驗此圖，先欲得日差，故存其左右之上下，以左右之中伏見日數，互偏乘兩行。乃以次圖右中伏日一十六日九十分，先偏乘左行，畢，左上得三百一度三十二分七十杪，左中得一千九百九日七十分，左下得一十萬六千九百四十三日二十分；又以次圖左中見日一百一十三，偏乗右行，畢，右上得四百四十日七十分，右中亦得一千九百九日七十分，右下得二十三萬九百七十八日二十一分五十杪。\end{unitzh}
\begin{uniten}In general, in the method of rectangular arrays, what one wishes to obtain first is kept, and with what one does not yet wish to obtain one mutually multiplies throughout (piancheng) all the numbers of the two columns. Now, examining this diagram, one first wishes to obtain the daily difference. Therefore keep the top and bottom of left and right, and with the hidden and appearance day numbers in the middle of left and right mutually multiply throughout the two columns. So, with the hidden days at the right middle of the second diagram, 16 days 90 fen, first multiply throughout the left column. When done, the left top gives 301 du 32 fen 70 miao, the left middle gives 1909 days 70 fen, and the left bottom gives 106943 days 20 fen. Again, with the appearance days at the left middle of the second diagram, 113, multiply throughout the right column. When done, the right top gives 440 days 70 fen, the right middle likewise gives 1909 days 70 fen, and the right bottom gives 230978 days 21 fen 50 miao.\end{uniten}

\begin{unitzh}\unitlabel{shushu-014.25}以兩行得變名泛積法，而成才圖。乃驗才圖，左上下皆少，用減右行，畢，行上餘一百三十九度三十七分三十杪，為日差實；右中空；右下得一十二萬四千三十五日一分五十杪，為日差法。今維圖法多實少，除得空度空分十一杪二十三小分六十五小杪，不盡十杪五十五小分三十九小杪五十二微分五十微杪，收為一小杪，為日定差一十一杪二十三小分六十六小秒。\end{unitzh}
\begin{uniten}With what the two columns have obtained, change their names to provisional, accumulated and divisor (fan, ji, fa), and so form the third diagram (caitu). Then, examining the third diagram, the left top and bottom are both smaller; use them to subtract from the right column. When done, the top of the column leaves 139 du 37 fen 30 miao, which is the dividend of the daily difference. The right middle is empty. The right bottom gives 124035 days 1 fen 50 miao, which is the divisor of the daily difference. Now, in the diagram the divisor is larger and the dividend smaller. Dividing gives 0 du 0 fen 11 miao 23 xiao fen (小分) 65 xiao miao (小杪), with an unexhausted remainder of 10 miao 55 xiao fen 39 xiao miao 52 wei fen (微分) 50 wei miao (微杪), which is collected as 1 xiao miao. This is the definite daily difference (ri dingcha), 11 miao 23 xiao fen 66 xiao miao.\end{uniten}

\begin{unitzh}\unitlabel{shushu-014.26}既得日差，乃行初行率。置法圖内初行法一百二十九日九十分，内減去一日，餘一百二十八日九十分，乗日差一十一杪二十三小分六十六小杪，得空度一十四分四十八杪三十九小分七十七小杪四十微分，為寄。次置{\zhsmall 按：此下脱三十一字，應作「次置寄，以半法乗之，得九度四十分七十三杪四十三小分三十二小杪一十三微杪」}。\end{unitzh}
\begin{uniten}Having obtained the daily difference, then proceed to the initial motion rate. Set out the initial-motion divisor in the divisor diagram, 129 days 90 fen; from it subtract 1 day, leaving 128 days 90 fen. Multiply by the daily difference, 11 miao 23 xiao fen 66 xiao miao, obtaining 0 du 14 fen 48 miao 39 xiao fen 77 xiao miao 40 wei fen, which is the lodged number. Next set out [Siku editors: below this, 31 characters have dropped out; it should read: "Next set out the lodged number, multiply it by half the divisor, obtaining 9 du 40 fen 73 miao 43 xiao fen 32 xiao miao 13 wei miao"].\end{uniten}

\begin{unitzh}\unitlabel{shushu-014.27}以得數加伏度三度九十分、見度一十七度八十三分，共得三十一度一十三分七十三杪四十三小分三十二小杪一十三微分，為初行實。如初行法一百二十九日九十分而一，乃行空度二十三分九十七杪，為伏合初日行率。餘三杪一十三分分三十二小杪一十三微分，棄之。\end{unitzh}
\begin{uniten}Add the number obtained to the hidden degrees, 3 du 90 fen, and the appearance degrees, 17 du 83 fen, obtaining altogether 31 du 13 fen 73 miao 43 xiao fen 32 xiao miao 13 wei fen, which is the initial-motion dividend. Divide by the initial-motion divisor, 129 days 90 fen; then it moves 0 du 23 fen 97 miao, which is the first-day motion rate of occultation-conjunction (fuhe). The remainder, 3 miao 13 fen fen 32 xiao miao 13 wei fen [Translator's note: "fen fen" (分分) is evidently corrupt], is discarded.\end{uniten}

\begin{unitzh}\unitlabel{shushu-014.28}求末行率：置合伏段日數一十六日九十分，乗日差一十一杪二十三小分六十六小杪，得一分八十九杪八十九小分八十五小杪四十微分，為得數。乃以得數減初行率二十三分九十七杪，餘二十二分七杪一十小分一十四小杪六十微分，為合伏末日行。\end{unitzh}
\begin{uniten}To find the final motion rate: set out the number of days of the conjunction-occultation stage, 16 days 90 fen, and multiply by the daily difference, 11 miao 23 xiao fen 66 xiao miao, obtaining 1 fen 89 miao 89 xiao fen 85 xiao miao 40 wei fen, which is the number obtained. Then subtract the number obtained from the initial motion rate, 23 fen 97 miao, leaving 22 fen 7 miao 10 xiao fen 14 xiao miao 60 wei fen, which is the motion of the last day of conjunction-occultation.\end{uniten}

\begin{unitzh}\unitlabel{shushu-014.29}求平行率：置初行率二十三分九十七杪，併末行率二十三分七杪，得四十六分四杪，以半之，得二十三分二杪，為平行率。\end{unitzh}
\begin{uniten}To find the mean motion rate: set out the initial motion rate, 23 fen 97 miao, and combine with the final motion rate, 23 fen 7 miao, obtaining 46 fen 4 miao. Halve it, obtaining 23 fen 2 miao, which is the mean motion rate.\end{uniten}

\begin{unitzh}\unitlabel{shushu-014.30}求交段差：置合伏日下，減全日一百分，餘一十分，乘末行率二十二分七杪，得二分二十杪七十小分，為交段差。\end{unitzh}
\begin{uniten}To find the stage-junction difference: set out the part below the days of conjunction-occultation and subtract it from a whole day of 100 fen, leaving 10 fen. Multiply by the final motion rate, 22 fen 7 miao, obtaining 2 fen 20 miao 70 xiao fen, which is the stage-junction difference.\end{uniten}

\begin{unitzh}\unitlabel{shushu-015.1}求晨疾初段常度：置合伏日一十六日九十分，乃收九十分作一日，通為一十七日；併舊厯所注晨疾初段常日三十，得四十七，為共日；乗合伏初行率二十三分九十七杪，得一十一度二十六分五十九杪，為寄上。\end{unitzh}
\begin{uniten}To find the constant degrees (changdu) of the first segment of morning fast (chenji chu duan): set up the days of conjunction-and-invisibility (hefu), 16 days 90 fen (分); then round up (shou) the 90 fen to make 1 day, making 17 days in all; add the constant days (changri) of the first segment of morning fast noted in the old calendar, 30, obtaining 47, which is the combined days (gongri); multiply by the initial rate of travel (chuxing lü) of conjunction-and-invisibility, 23 fen 97 miao (杪), obtaining 11 du (度) 26 fen 59 miao, which is reserved (ji) above.\end{uniten}

\begin{unitzh}\unitlabel{shushu-015.2}按：此有第一日行度，有逐日遞減之差，有前後各段日數，有前段積度，求後段積也。先以共日乗初行率者，以最疾為率之共積也；下求遞差以減之，故為寄。\end{unitzh}
\begin{uniten}Note: here there are the degrees travelled on the first day; there is the difference by which it decreases successively day by day; there are the numbers of days of each of the preceding and following segments; there are the accumulated degrees (jidu) of the preceding segment; and one seeks the accumulation of the following segment. First multiplying the initial rate of travel by the combined days gives the combined accumulation (gongji) taking the fastest [motion] as the rate; below, the successive difference (dicha) is sought in order to subtract it, hence this is reserved.\end{uniten}

\begin{unitzh}\unitlabel{shushu-015.3}乃副置共日四十七，減一，餘四十六，以半之，得二十三，乗副四十七，得一千八十一；以乗日差一十一杪二十三小分六十六小杪，得一度二十一分四十六杪七十六小分四十六小杪；以減上寄一十一度一十六分五十九杪，餘一十度五分一十二杪二十三小分五十四小杪，為合伏、晨疾初兩段共積度。\end{unitzh}
\begin{uniten}Then set up in duplicate (fuzhi) the combined days, 47; subtract 1, leaving 46; halve it, obtaining 23; multiply by the duplicate 47, obtaining 1081; with this multiply the daily difference (richa), 11 miao 23 xiaofen (小分) 66 xiaomiao (小杪), obtaining 1 du 21 fen 46 miao 76 xiaofen 46 xiaomiao; subtract this from the [number] reserved above, 11 du 16 fen 59 miao [Translator's note: above, 26 fen.], leaving 10 du 5 fen 12 miao 23 xiaofen 54 xiaomiao, which is the combined accumulated degrees of the 2 segments, conjunction-and-invisibility and first morning fast.\end{uniten}

\begin{unitzh}\unitlabel{shushu-015.4}按：此乃求積差，以減上數，得共日之積慶也。法應於共日内減一日，以乘日差，得數為共日數初末日行率之較；再以共日數乘之，得數折半，為積差。此先折半，次連乘，得積差，其理亦同。\end{unitzh}
\begin{uniten}Note: this is to find the accumulated difference (jicha), in order to subtract it from the number above and obtain the accumulated degrees of the combined days [Translator's note: the text writes 慶 for 度 'degrees'.]. The method should be: from the combined days subtract 1 day, and with this multiply the daily difference; the number obtained is the difference (jiao) between the rates of travel of the first and last days of the combined number of days; again multiply it by the combined number of days, and halve the number obtained, which is the accumulated difference. Here one first halves and then multiplies successively (liancheng) to obtain the accumulated difference; the principle is the same.\end{uniten}

\begin{unitzh}\unitlabel{shushu-015.5}置共積，内減合伏三度九十分，餘六度一十五分一十二杪二十三小分五十四小杪，為泛；次以交段差二分二十杪七十小分減泛，餘六度一十二分九十一杪五十三小分五十四小杪，為晨疾初段常度；注厯，乃收八杪五{\zhsmall 按應作四}十六小分四十六小杪，為全分常定度。\end{unitzh}
\begin{uniten}Set up the combined accumulation; from it subtract the 3 du 90 fen of conjunction-and-invisibility, leaving 6 du 15 fen 12 miao 23 xiaofen 54 xiaomiao, which is provisional (fan); next subtract the junction-of-segments difference (jiaoduan cha), 2 fen 20 miao 70 xiaofen, from the provisional [value], leaving 6 du 12 fen 91 miao 53 xiaofen 54 xiaomiao, which is the constant degrees of the first segment of morning fast; for noting in the calendar, then round up by 8 miao 5[Siku editors: should be 4]6 xiaofen 46 xiaomiao, making the constant fixed (ding) degrees in whole fen.\end{uniten}

\begin{unitzh}\unitlabel{shushu-015.6}按：此於共積内減去合伏段積，尚有合伏九十分不及一日所差之行分，即交段差，未減，故為泛數；再減交段差，為晨疾初段常泛度；再收為六度十三分，始為定常度也。\end{unitzh}
\begin{uniten}Note: here the accumulation of the conjunction-and-invisibility segment is subtracted from the combined accumulation; there still remain the fen of travel that differ because the 90 fen of conjunction-and-invisibility fall short of 1 day, that is, the junction-of-segments difference, not yet subtracted; hence it is a provisional number (fanshu). Further subtracting the junction-of-segments difference gives the constant provisional degrees of the first segment of morning fast; further rounding up to 6 du 13 fen, it only then becomes the fixed constant degrees.\end{uniten}

\begin{unitzh}\unitlabel{shushu-015.7}求晨疾初段初行率：以日差一十一杪二十三小分六十六小杪減合伏末行率二十二分七杪，餘二十一分九十六杪，為晨疾初段初行率，行泛收之為定者也。\end{unitzh}
\begin{uniten}To find the initial rate of travel of the first segment of morning fast: subtract the daily difference, 11 miao 23 xiaofen 66 xiaomiao, from the final rate of travel (moxing lü) of conjunction-and-invisibility, 22 fen 7 miao, leaving 21 fen 96 miao, which is the initial rate of travel of the first segment of morning fast; it is the provisional travel rounded up to make the fixed [value].\end{uniten}

\begin{unitzh}\unitlabel{shushu-015.8}按：此以合伏末日之次日為晨疾初段之初日也，故置合伏之末行率，減一日之差，即為晨疾初段之初行率。五杪餘收為六杪。凡寄零未收名泛數，已收名定數。下倣此。\end{unitzh}
\begin{uniten}Note: here the day after the last day of conjunction-and-invisibility is taken as the first day of the first segment of morning fast; hence one sets up the final rate of travel of conjunction-and-invisibility and subtracts the difference for 1 day, which is then the initial rate of travel of the first segment of morning fast. The 5 miao and more is rounded up to 6 miao. In general, when the reserved remainder (jiling) has not been rounded up, it is called a provisional number; when it has been rounded up, it is called a fixed number (dingshu). What follows imitates this.\end{uniten}

\begin{unitzh}\unitlabel{shushu-015.9}求晨疾初末行率：置晨疾初常日三十，減一，餘二十九日，乗日差一十一杪二十三小分六十六十小杪，得三分二十五杪八十六小分一十四小杪；以減晨疾初段初行率泛二十一分九十五杪七十六小分三十四小杪，餘一十八分六十九杪九十小分二十小杪，為晨疾初末行率。\end{unitzh}
\begin{uniten}To find the final rate of travel of first morning fast: set up the constant days of first morning fast, 30; subtract 1, leaving 29 days; multiply by the daily difference, 11 miao 23 xiaofen 66 xiaomiao [Translator's note: the text writes 六十六十小杪, with a superfluous 十.], obtaining 3 fen 25 miao 86 xiaofen 14 xiaomiao; subtract this from the provisional initial rate of travel of the first segment of morning fast, 21 fen 95 miao 76 xiaofen 34 xiaomiao, leaving 18 fen 69 miao 90 xiaofen 20 xiaomiao, which is the final rate of travel of first morning fast.\end{uniten}

\begin{unitzh}\unitlabel{shushu-015.10}按：此求晨疾初段末日之行率也。常日減一日，乗日差，得數為晨疾初段初末二日行率之較也，故減初行率，得末行率。\end{unitzh}
\begin{uniten}Note: this finds the rate of travel on the last day of the first segment of morning fast. The constant days less 1 day, multiplied by the daily difference, give a number which is the difference between the rates of travel of the 2 days, first and last, of the first segment of morning fast; hence subtracting it from the initial rate of travel gives the final rate of travel.\end{uniten}

\begin{unitzh}\unitlabel{shushu-015.11}求平行率：以晨疾初初行泛二十一分九十五杪七十六小分三十四小杪，併晨疾初末泛一十八分六十九杪九十小分二十小分，得四十分六十五杪六十六小分五十四小杪；以半之，得二十分三十二杪八十三小分二十七小杪，為晨疾初平行泛。乃以三泛收棄為之定。\end{unitzh}
\begin{uniten}To find the mean rate of travel (pingxing lü): take the provisional initial travel of first morning fast, 21 fen 95 miao 76 xiaofen 34 xiaomiao, and add the provisional final [travel] of first morning fast, 18 fen 69 miao 90 xiaofen 20 xiaofen, obtaining 40 fen 65 miao 66 xiaofen 54 xiaomiao; halve it, obtaining 20 fen 32 miao 83 xiaofen 27 xiaomiao, which is the provisional mean travel of first morning fast. Then the 3 provisional [values] are rounded up or discarded (qi) to make them fixed.\end{uniten}

\begin{unitzh}\unitlabel{shushu-015.12}按：此與求合伏平行率同。其言泛棄為定者，蓋截去杪下竒零，過半則收為一杪也，然語意欠明。\end{unitzh}
\begin{uniten}Note: this is the same as finding the mean rate of travel of conjunction-and-invisibility. Where it speaks of the provisional being discarded to make the fixed, it presumably means cutting off the odd remainder (qiling) below the miao, and if it exceeds half, rounding it up to 1 miao; but the sense of the wording is unclear.\end{uniten}

\begin{unitzh}\unitlabel{shushu-015.13}又按：五星行度遲疾差廻，非遞加遞減之數。術中僅以合伏與順見二段各取中數，至推遂日行度，仍用遞加遞減之法，故古法之疎，五星尤甚。原文語多隠晦，令悉為解之，可以見古今疎宻之所在焉。\end{unitzh}
\begin{uniten}Further note: the degrees travelled by the five planets—slow and fast, varying and turning back—are not numbers of successive addition and successive subtraction. In the method, only the 2 segments of conjunction-and-invisibility and direct appearance (shunjian) each take a mean number (zhongshu); when it comes to computing the degrees travelled day by day [Translator's note: the text writes 遂 for 逐.], the method of successive addition and successive subtraction is still used; hence the coarseness of the ancient methods is most severe for the five planets. The language of the original text is often obscure; let it all be explained, so that one can see where the coarseness and precision of ancient and modern lie.\end{uniten}

\begin{unitzh}\unitlabel{shushu-015.14}《數學九章》卷二上\end{unitzh}
\begin{uniten}Shuxue jiuzhang, juan 2, part 1.\end{uniten}

\begin{unitzh}\unitlabel{shushu-015.15}<子部、天文算法類、算書之屬、數學九章>\end{unitzh}
\begin{uniten}<Masters section, astronomy and mathematics category, mathematical books subcategory, Shuxue jiuzhang>\end{uniten}

\section*{卷2下}\addcontentsline{toc}{section}{卷2下}

\begin{unitzh}\unitlabel{shushu-016.1}欽定四庫全書\end{unitzh}
\begin{uniten}Imperially Commissioned Complete Library of the Four Treasuries (Siku quanshu)\end{uniten}

\begin{unitzh}\unitlabel{shushu-016.2}《數學九章》卷二下 宋 秦九韶 撰\end{unitzh}
\begin{uniten}Shuxue jiuzhang, juan 2, second part. Composed by Qin Jiushao of the Song.\end{uniten}

\begin{unitzh}\unitlabel{shushu-016.3}天時\end{unitzh}
\begin{uniten}Heavenly Seasons (tianshi)\end{uniten}

\begin{unitzh}\unitlabel{shushu-016.4}揆日究微\end{unitzh}
\begin{uniten}Gauging the Sun and Probing the Minute\end{uniten}

\begin{unitzh}\unitlabel{shushu-016.5}問：厯代測景，惟唐《大衍厯》最宻。本朝《崇天厯》陽城冬至景一丈二尺七寸一分五十杪，夏至景一尺四寸七分七十杪九，係與《大衍厯》同。今《開禧厯》臨安府冬至景一丈八寸二分二十五杪，夏至景九寸一分。欲求臨安府夏至後差㡬日而景與陽城夏至日等，較以《大衍厯》晷景所差尺寸各㡬何？答曰：大暑後五日午中景長一尺四寸八分八十五杪。\end{unitzh}
\begin{uniten}Question: Of the shadow measurements of the successive dynasties, only the Dayan calendar of the Tang is the most precise. In the Chongtian calendar of our dynasty, the winter-solstice shadow at Yangcheng is 1 zhang (丈) 2 chi (尺) 7 cun (寸) 1 fen (分) 50 miao (杪), and the summer-solstice shadow is 1 chi 4 cun 7 fen 70 miao 9 [Translator's note: 七十杪九 is evidently transposed; below, this shadow is given as 七十九杪, 79 miao]; these are the same as in the Dayan calendar. Now, in the Kaixi calendar, the winter-solstice shadow at Lin'an prefecture is 1 zhang 8 cun 2 fen 25 miao, and the summer-solstice shadow is 9 cun 1 fen. It is desired to find how many days after the summer solstice the shadow at Lin'an prefecture becomes equal to that of the summer-solstice day at Yangcheng, and, compared with the gnomon shadow (guijing) of the Dayan calendar, how many chi and cun the difference is in each case. Answer: 5 days after Major Heat (dashu), at noon, the shadow is 1 chi 4 cun 8 fen 85 miao long.\end{uniten}

\begin{unitzh}\unitlabel{shushu-016.6}按：舊本答數後有二圖，舛錯潦草，傳冩者失其真也。細考圖内所載之數，皆與今法頗合，知此悉當時實測所定，非同臆説也。因取其數，改正於後。\end{unitzh}
\begin{uniten}Note: In the old edition there were two diagrams after the numbers of the answer; they were erroneous and sloppy, and those who transmitted and copied them lost their true form. Examining closely the numbers recorded in the diagrams, all agree quite well with the present method, from which one knows that they were all fixed by actual measurement at the time and are not the same as conjecture. Therefore their numbers have been taken and set out, corrected, below.\end{uniten}

\begin{unitzh}\unitlabel{shushu-016.7}子部，天文算法類，算書之屬，《數學九章》，卷二下\end{unitzh}
\begin{uniten}Masters section (zibu), astronomy and mathematics category, mathematical books subcategory: Shuxue jiuzhang, juan 2, second part\end{uniten}

\begin{unitzh}\unitlabel{shushu-016.8}置臨安府所測冬至景一丈八寸二分二十五杪，以夏至景九寸一分減之，餘九尺九寸一分二十五杪，為景差，以為實。置象限度九十一度三十一分四十四杪，加一十一度二十五分二十七杪五十小分，命度為寸，得一百二寸五十六百七十一分五十杪，為法。以除前差實，得空寸九十六百六十四分四十杪，不盡棄之。自乗，得節泛數九千三百四十分，不盡棄之。\end{unitzh}
\begin{uniten}Set up the winter-solstice shadow measured at Lin'an prefecture, 1 zhang 8 cun 2 fen 25 miao; subtract from it the summer-solstice shadow, 9 cun 1 fen; the remainder, 9 chi 9 cun 1 fen 25 miao, is the shadow difference (jingcha); take it as the dividend (shi). Set up the degrees of the quadrant (xiangxian du), 91 du (度) 31 fen 44 miao; add 11 du 25 fen 27 miao 50 small fen (xiaofen); taking du as cun, one obtains 102 cun 5671 fen 50 miao, the divisor (fa). With it divide the preceding difference, the dividend, obtaining 0 cun 9664 fen 40 miao; the remainder is discarded. Multiply this by itself, obtaining the provisional node number (jie fanshu), 9340 fen; the remainder is discarded.\end{uniten}

\begin{unitzh}\unitlabel{shushu-016.9}先以小暑節乗率二十五乗節率九千三百四十分，得二十三寸三千五百分於上。次以臨安夏至九寸一分自乗，得八十二寸八千一百分，為夏至冪。乃以夏至冪加上，得一百六寸一千六百分，為小暑冪，以為實。以一寸為隅，開平方，得一尺三分，為臨安小暑節景。不盡七百分，即寸下七毫，棄之。\end{unitzh}
\begin{uniten}First, with the multiplying rate (chenglü) of the Minor Heat (xiaoshu) node, 25, multiply the node rate (jielü), 9340 fen, obtaining 23 cun 3500 fen, placed above. Next, multiply the Lin'an summer solstice, 9 cun 1 fen, by itself, obtaining 82 cun 8100 fen, the summer-solstice area (mi). Then add the summer-solstice area to that above, obtaining 106 cun 1600 fen, the Minor Heat area; take it as the dividend. With 1 cun as the corner (yu), extract the square root, obtaining 1 chi 3 fen, the shadow of the Minor Heat node at Lin'an. The remainder, 700 fen, that is, 7 hao (毫) below the cun, is discarded.\end{uniten}

\begin{unitzh}\unitlabel{shushu-016.10}按：各節氣影長，皆當時實測所定，本不待求。今所設求法，乃故為溟涬，使人不可觧也。細查其數，首以象限加十一度餘為法，以除影差，得數自乗為節率；四每節下又有乗率，以乗率、節率相乗，與夏至影冪相加，即為本節影冪。是知節率乃强取之數，蓋以此數先除各節影冪與夏至影冪之較，名為乗率，故以此與節率相乗，加夏至影冪，即各節影冪也。數家設術誤人，往往如此。\end{unitzh}
\begin{uniten}Note: The shadow lengths of the various solar terms (jieqi) were all fixed by actual measurement at the time and originally did not need to be found. The method of finding set up here deliberately makes things murky, so that people cannot understand it. Examining its numbers closely: first, the quadrant plus 11-odd du is taken as the divisor, by which the shadow difference is divided, and the number obtained, multiplied by itself, is the node rate; 4 [Translator's note: 四 "four" appears corrupt here] under each node there is moreover a multiplying rate; the multiplying rate and the node rate are multiplied together and added to the area of the summer-solstice shadow, which gives the shadow area of that node. From this one knows that the node rate is a forcibly chosen number: presumably this number was first used to divide the difference between each node's shadow area and the summer-solstice shadow area, and this was named the multiplying rate; therefore, multiplying this by the node rate and adding the summer-solstice shadow area gives each node's shadow area. The way mathematicians set up methods that mislead people is often like this.\end{uniten}

\begin{unitzh}\unitlabel{shushu-016.11}又以大暑乗率一百九乗節率九千三百四十分，得一百一寸八千六十分於頭位。仍以夏至冪八十二寸八千一百分加頭位一百一寸八千六十分，得一百八十四寸六千一百六十分，為大暑冪，以為實。以一寸為隅，開平方，得一尺三寸五分八十七杪，為大暑。不盡，棄之。\end{unitzh}
\begin{uniten}Again, with the Major Heat multiplying rate, 109, multiply the node rate, 9340 fen, obtaining 101 cun 8060 fen, placed in the head position (touwei). Likewise add the summer-solstice area, 82 cun 8100 fen, to the head position, 101 cun 8060 fen, obtaining 184 cun 6160 fen, the Major Heat area; take it as the dividend. With 1 cun as the corner, extract the square root, obtaining 1 chi 3 cun 5 fen 87 miao, which is Major Heat. The remainder is discarded.\end{uniten}

\begin{unitzh}\unitlabel{shushu-016.12}又置立秋乗率二百八十九，乗節率九千三百四十分，得二百六十九寸九千二百六十分於上。仍置夏至冪八十一寸八千一百分，加於上二百六十九寸九千二百六十分，得三百五十二寸七千三百六十八，為立秋冪。置立秋冪為實，以一寸為隅，開平方，得一尺八寸七分八十一杪，為立秋景。不盡，棄之。\end{unitzh}
\begin{uniten}Again set up the multiplying rate of the Beginning of Autumn (liqiu), 289, and multiply the node rate, 9340 fen, obtaining 269 cun 9260 fen, placed above. Likewise set up the summer-solstice area, 81 cun 8100 fen, and add it to the 269 cun 9260 fen above, obtaining 352 cun 7368, the Beginning of Autumn area. Set up the Beginning of Autumn area as the dividend; with 1 cun as the corner, extract the square root, obtaining 1 chi 8 cun 7 fen 81 miao, the Beginning of Autumn shadow. The remainder is discarded.\end{uniten}

\begin{unitzh}\unitlabel{shushu-016.13}乃驗陽城夏至景一尺四寸七分七十九杪，在大暑後、立秋前。乃置大暑一尺三寸五分八十七杪，併立秋景一尺八寸七分八十一杪，得三尺二寸三分六十八杪，以半之，得一尺六寸一分八十四杪，為大暑後九日景。又以九日景併大暑景，得二尺九寸七分七千一杪，以半之，得一尺四寸八分八十五杪半，為大暑後五日景。乃以大暑景得五日景，得二尺八寸四分七十二杪半，以半之，得一尺四寸二分三十六杪少，為大暑後三日景。又以五日景得三日景，得二尺九寸一分二十一杪大，以半之，得一尺四寸五分六十杪八十七小分五十小杪，為大暑後四日景。\end{unitzh}
\begin{uniten}Then examine: the summer-solstice shadow of Yangcheng, 1 chi 4 cun 7 fen 79 miao, lies after Major Heat and before the Beginning of Autumn. So set up Major Heat, 1 chi 3 cun 5 fen 87 miao, and combine (bing) it with the Beginning of Autumn shadow, 1 chi 8 cun 7 fen 81 miao, obtaining 3 chi 2 cun 3 fen 68 miao; halve it, obtaining 1 chi 6 cun 1 fen 84 miao, the shadow 9 days after Major Heat. Again combine the 9th-day shadow with the Major Heat shadow, obtaining 2 chi 9 cun 7 fen 7001 miao; halve it, obtaining 1 chi 4 cun 8 fen 85 1/2 miao, the shadow 5 days after Major Heat. Then with the Major Heat shadow "obtain" the 5th-day shadow [Translator's note: 得 "obtain", here and in the next sentence, evidently stands for 併 "combine"], obtaining 2 chi 8 cun 4 fen 72 1/2 miao; halve it, obtaining 1 chi 4 cun 2 fen 36 1/4 miao, the shadow 3 days after Major Heat. Again with the 5th-day shadow "obtain" the 3rd-day shadow, obtaining 2 chi 9 cun 1 fen 21 3/4 miao; halve it, obtaining 1 chi 4 cun 5 fen 60 miao 87 small fen 50 small miao (xiaomiao), the shadow 4 days after Major Heat.\end{uniten}

\begin{unitzh}\unitlabel{shushu-016.14}今驗陽城夏至景一尺四寸七分七十九杪，為入臨安府大暑後四日景一尺四寸五分六十杪太強。乃以四日減五日景，餘三分二十四杪太弱，為景差，以十二時除之，得二十七杪五小分二十小杪，為法。乃置陽城夏至景一尺四寸七分七十九抄，減臨安大暑後四日景一尺四寸五分六十杪八十七小分五十小杪，餘二分一十八杪一十二小分五十小杪，為實。後以法二十七杪五小分二十小杪除之，實如法而一，得商數八有餘。命大暑四日午後，數八辰，得大暑五日寅時，景與陽城夏至之日午景等。\end{unitzh}
\begin{uniten}Now examine: the summer-solstice shadow of Yangcheng, 1 chi 4 cun 7 fen 79 miao, enters (ru) the shadow 4 days after Major Heat at Lin'an prefecture, 1 chi 4 cun 5 fen 60 3/4 miao, slightly more (qiang). Then subtract the 4th day from the 5th-day shadow; the remainder, 3 fen 24 3/4 miao, slightly less (ruo), is the shadow difference; divide it by the 12 double-hours (shi 時), obtaining 27 miao 5 small fen 20 small miao, the divisor. Then set up the summer-solstice shadow of Yangcheng, 1 chi 4 cun 7 fen 79 miao, and subtract the Lin'an shadow 4 days after Major Heat, 1 chi 4 cun 5 fen 60 miao 87 small fen 50 small miao; the remainder, 2 fen 18 miao 12 small fen 50 small miao, is the dividend. Afterwards divide it by the divisor, 27 miao 5 small fen 20 small miao: dividing the dividend by the divisor, one obtains a quotient (shangshu) of 8 with something over. Naming from after wu (noon) on the 4th day of Major Heat and counting 8 chen (辰), one obtains the yin double-hour (yinshi) of the 5th day of Major Heat, when the shadow equals the noon shadow of the summer-solstice day at Yangcheng.\end{uniten}

\begin{unitzh}\unitlabel{shushu-016.15}求較以《大衍厯》晷景所差：乃置陽城大暑景長一尺九寸五分七十六杪，併陽城立秋景二尺五寸三分三十一杪，得四尺四寸九分七杪，以半之，得二尺二寸四分五十三杪半，為次暑後九日午中景。置九日景，復併大暑景一尺九寸五分七十六抄，得四尺二寸二十九杪半，以半之，得二尺一寸一十四杪太，為大暑後五日景。今驗《開禧厯》所推臨安府大暑後五日午中景一尺四寸八分八十五杪半，與陽城大暑後五日午中景尺一寸一十四杪太課之。乃以臨安府五日景減陽城五日景，差六寸一分二十九杪少。\end{unitzh}
\begin{uniten}To find the difference in comparison with the gnomon shadow of the Dayan calendar: set up the length of the Major Heat shadow at Yangcheng, 1 chi 9 cun 5 fen 76 miao, and combine it with the Beginning of Autumn shadow at Yangcheng, 2 chi 5 cun 3 fen 31 miao, obtaining 4 chi 4 cun 9 fen 7 miao; halve it, obtaining 2 chi 2 cun 4 fen 53 1/2 miao, the noon shadow 9 days after Next Heat (cishu). Set up the 9th-day shadow and again combine it with the Major Heat shadow, 1 chi 9 cun 5 fen 76 miao, obtaining 4 chi 2 cun 29 1/2 miao; halve it, obtaining 2 chi 1 cun 14 3/4 miao, the shadow 5 days after Major Heat. Now examine the noon shadow 5 days after Major Heat at Lin'an prefecture as reckoned by the Kaixi calendar, 1 chi 4 cun 8 fen 85 1/2 miao, and check (ke) it against the noon shadow 5 days after Major Heat at Yangcheng, chi 1 cun 14 3/4 miao. Then subtract the 5th-day shadow of Lin'an prefecture from the 5th-day shadow of Yangcheng; the difference is 6 cun 1 fen 29 1/4 miao.\end{uniten}

\begin{unitzh}\unitlabel{shushu-016.16}按：此法不過以臨安前後兩節氣影長比例一影長之日數時刻，復以所得節氣日數時刻比例一陽城影長，與之相較耳。題内引大行、崇天、開禧諸法名目，又稱其較同異差數，皆故為張皇之語。且影差逐日不同，皆以平派求之，法亦未宻也。\end{unitzh}
\begin{uniten}Note: This method does no more than use the shadow lengths of the two solar terms before and after, at Lin'an, to find by proportion (bili) the number of days and the hour for one shadow length, and then use the number of days and the hour of the solar term so obtained to find by proportion one shadow length at Yangcheng and compare it with that. In the problem, the citing of the names of the methods of the Daxing [Translator's note: 大行 is written for 大衍, Dayan], Chongtian and Kaixi, and the further talk of comparing their agreement and difference and the amounts of difference, are all words deliberately made grandiose. Moreover, the shadow differences are not the same from day to day, yet all are found by even apportionment (pingpai); the method, too, is not precise.\end{uniten}

\begin{unitzh}\unitlabel{shushu-016.17}草曰：置臨安府所測冬至景一丈八寸二分二十五抄，以夏至景九寸一分減之，餘九尺九寸一分二十五抄，為景差，以為實。置象度九十一度三十一分四十四杪，加一十一度二十五分二十七杪半，命度為寸，得一百二寸五千六百七十一分半，為法。除實，得空寸九千六百六十四分四十杪，以自乗之，得空寸九千三百四十分，為節率。先以小暑乗率二十五乗之，得二十三寸三千五百分於上。次以臨安夏至景九寸一分自乗，得八十二寸八千一百分，為夏至景冪。以加上，得一百六寸一千六百分，為小暑景冪。開平方，以一寸為隅開之，得一尺三分，為小暑景。\end{unitzh}
\begin{uniten}Working: Set up the winter-solstice shadow measured at Lin'an prefecture, 1 zhang 8 cun 2 fen 25 miao; subtract from it the summer-solstice shadow, 9 cun 1 fen; the remainder, 9 chi 9 cun 1 fen 25 miao, is the shadow difference; take it as the dividend. Set up the quadrant degrees (xiangdu), 91 du 31 fen 44 miao; add 11 du 25 fen 27 1/2 miao; taking du as cun, one obtains 102 cun 5671 1/2 fen, the divisor. Divide the dividend, obtaining 0 cun 9664 fen 40 miao; multiply it by itself, obtaining 0 cun 9340 fen, the node rate. First, with the Minor Heat multiplying rate, 25, multiply it, obtaining 23 cun 3500 fen, placed above. Next multiply the Lin'an summer-solstice shadow, 9 cun 1 fen, by itself, obtaining 82 cun 8100 fen, the area of the summer-solstice shadow. Add it to that above, obtaining 106 cun 1600 fen, the area of the Minor Heat shadow. Extract the square root: with 1 cun as the corner, extract it, obtaining 1 chi 3 fen, the Minor Heat shadow.\end{uniten}

\begin{unitzh}\unitlabel{shushu-016.18}又以大暑乗率一百九乗節率九千三百四十分，得一百一寸八千六十分於上，仍加夏至冪八十二寸八千一百分，得一百八十四寸六千一百六十分，為大暑冪，以為實。一寸寸為隅，開平方，得一尺三寸五分八十七杪，為大暑景。又置立秋乗率二百八十九，乗節率九千三百四十分，得二百六十九寸九千二百六十分於上，仍加夏至冪八十二寸八千一百分，共得三百五十二寸七千三百六十，為立秋冪，以為實。以一寸為隅，開平方，得一尺八寸七分八十一抄，為立秋景。\end{unitzh}
\begin{uniten}Again, with the Major Heat multiplying rate, 109, multiply the node rate, 9340 fen, obtaining 101 cun 8060 fen, placed above; likewise add the summer-solstice area, 82 cun 8100 fen, obtaining 184 cun 6160 fen, the Major Heat area; take it as the dividend. 1 cun cun as the corner, extract the square root, obtaining 1 chi 3 cun 5 fen 87 miao, the Major Heat shadow. Again set up the Beginning of Autumn multiplying rate, 289, and multiply the node rate, 9340 fen, obtaining 269 cun 9260 fen, placed above; likewise add the summer-solstice area, 82 cun 8100 fen, obtaining altogether 352 cun 7360, the Beginning of Autumn area; take it as the dividend. With 1 cun as the corner, extract the square root, obtaining 1 chi 8 cun 7 fen 81 miao, the Beginning of Autumn shadow.\end{uniten}

\begin{unitzh}\unitlabel{shushu-016.19}乃驗陽城夏至景一尺四寸七分七十九杪，在大暑後、立秋前。乃併大暑、立秋二景，半之，得一尺六寸一分八十四杪，為大暑後九日景。又併大暑景，半之，得一尺四寸八分八十五杪半，為大暑後五日午中景。又併大暑景，得數半之，得一尺四寸二分三十六杪少，為大暑後三日景。又併五日景一尺四寸八分八十五杪半，得數半之，得一尺四寸五分六十杪强，為大暑後四日景。驗得陽城夏至景入臨安大暑後四日，乃以四日景減五日景，餘三分二十四杪太弱，為差，以十二時除之，得二十七杪五小分二十小杪，為法。復除陽城景與本日景差二分一十八杪一十二小分五十小杪，得八，命外為在初五日寅時景等。\end{unitzh}
\begin{uniten}Then examine: the summer-solstice shadow of Yangcheng, 1 chi 4 cun 7 fen 79 miao, lies after Major Heat and before the Beginning of Autumn. So combine the two shadows of Major Heat and the Beginning of Autumn and halve, obtaining 1 chi 6 cun 1 fen 84 miao, the shadow 9 days after Major Heat. Again combine with the Major Heat shadow and halve, obtaining 1 chi 4 cun 8 fen 85 1/2 miao, the noon shadow 5 days after Major Heat. Again combine with the Major Heat shadow and halve the number obtained, obtaining 1 chi 4 cun 2 fen 36 1/4 miao, the shadow 3 days after Major Heat. Again combine with the 5th-day shadow, 1 chi 4 cun 8 fen 85 1/2 miao, and halve the number obtained, obtaining 1 chi 4 cun 5 fen 60 miao, slightly more, the shadow 4 days after Major Heat. Having examined and found that the summer-solstice shadow of Yangcheng enters the 4th day after Major Heat at Lin'an, subtract the 4th-day shadow from the 5th-day shadow; the remainder, 3 fen 24 3/4 miao, slightly less, is the difference; divide it by the 12 double-hours, obtaining 27 miao 5 small fen 20 small miao, the divisor. Again divide the difference between the Yangcheng shadow and the shadow of this day, 2 fen 18 miao 12 small fen 50 small miao, obtaining 8; naming it outward (ming wai), it is at the yin double-hour of the 5th day (chuwu ri) that the shadows are equal.\end{uniten}

\begin{unitzh}\unitlabel{shushu-016.20}求較以《大衍厯》晷景所差：乃置陽城大暑景長一尺九寸五分七十六杪，併陽城立秋景二尺五寸三分三十一杪，得四尺四寸九分七杪，以半之，得二尺二寸四分五十三杪半，為大暑後九日午中景。復併暑景一尺九寸五分七十六杪，得四尺二寸二十九杪半，以半之，得二尺一寸一十四杪太，為大暑後五日暑。以較今《開禧厯》當日景一尺四寸八分八十五杪半，差少六寸一分二十九杪少。合問。\end{unitzh}
\begin{uniten}To find the difference in comparison with the gnomon shadow of the Dayan calendar: set up the length of the Major Heat shadow at Yangcheng, 1 chi 9 cun 5 fen 76 miao, and combine it with the Beginning of Autumn shadow at Yangcheng, 2 chi 5 cun 3 fen 31 miao, obtaining 4 chi 4 cun 9 fen 7 miao; halve it, obtaining 2 chi 2 cun 4 fen 53 1/2 miao, the noon shadow 9 days after Major Heat. Again combine with the Heat shadow, 1 chi 9 cun 5 fen 76 miao, obtaining 4 chi 2 cun 29 1/2 miao; halve it, obtaining 2 chi 1 cun 14 3/4 miao, the heat (shu) 5 days after Major Heat. Comparing it with the shadow of that day in the present Kaixi calendar, 1 chi 4 cun 8 fen 85 1/2 miao, the difference: less by 6 cun 1 fen 29 1/4 miao. This answers the question.\end{uniten}

\begin{unitzh}\unitlabel{shushu-016.21}接：集中皆術在前，草次之，圖在後，此條之例不同。\end{unitzh}
\begin{uniten}Note [Translator's note: the text writes 接 for 按, "note"]: Throughout the collection, the method comes first, the working next, and the diagrams after; the format of this item is different.\end{uniten}

\begin{unitzh}\unitlabel{shushu-016.22}天池測雨\end{unitzh}
\begin{uniten}Measuring Rain with a Heavenly Pool (tianchi)\end{uniten}

\begin{unitzh}\unitlabel{shushu-017.1}問：今州郡多有天池盆，以測雨水。但知以盆中之水為得雨之數，不知器形不同，則受雨多少亦異，未可以所測便為平地得雨之數。假今盆口徑二尺八寸，底徑一尺二寸，深一尺八寸，接雨水深九寸，欲求平地雨䧏㡬何？答曰：平地雨䧏三寸。\end{unitzh}
\begin{uniten}Question: Nowadays the prefectures and commanderies mostly have heavenly-pool basins (tianchi pen) with which to measure rainwater. But they only know to take the water in the basin as the amount of rain received; they do not know that when the shapes of the vessels are not the same, the amount of rain received also differs, so that what is measured cannot simply be taken as the amount of rain received on level ground. Suppose now that a basin has a mouth diameter of 2 chi (尺) 8 cun (寸), a base diameter of 1 chi 2 cun and a depth of 1 chi 8 cun, and that the rainwater it has caught is 9 cun deep. It is desired to find how much rain fell on level ground. Answer: On level ground the rainfall is 3 cun.\end{uniten}

\begin{unitzh}\unitlabel{shushu-017.2}術曰：盆深乗底徑，為底率；二徑差乗水深，併底率，為靣率。以盆深為法，除靣率，得面徑。以二率相乗，又各自乗，三位併之，乗水深，為實。盆深乗口徑，以自之，又三因，為法。除之，得平地水深。\end{unitzh}
\begin{uniten}Method: Multiply the depth of the basin by the base diameter to make the base rate (di lü). Multiply the difference of the two diameters by the water depth and add the base rate to make the surface rate (mian lü). Take the depth of the basin as divisor (fa) and divide the surface rate, obtaining the surface diameter. Multiply the two rates together, and also multiply each by itself; add the three positions together and multiply by the water depth to make the dividend (shi). Multiply the depth of the basin by the mouth diameter, multiply it by itself, and further triple it (san yin) to make the divisor. Divide, obtaining the depth of water on level ground.\end{uniten}

\begin{unitzh}\unitlabel{shushu-017.3}草曰：以盆深及徑皆通為寸，盆深得一十八寸，底徑得一十二寸，相乗，得二百一十六寸，為底率。置口徑二十八寸，減底徑一十二寸，餘一十六寸，為差，以乗水深九寸，得一百四十四寸，併底率二百一十六寸，得三百六十寸，為靣率。以盆深一十八寸為法，除靣率，得二十寸，展為二尺，為水靣徑。以底率二百一十六寸乗靣率三百六十寸，得七萬七千七百六十寸於上；以底率二百一十六寸自乗，得四萬六千六百五十六寸，加上；又以靣率三百六十寸自乗，得一十二萬九千六百，併上，共得二十五萬四千一十六，以乗水深九寸，得二百二十八萬六千一百四十四寸，為實。以盆深一十八寸乗口徑二十八寸，得五百四寸，自乗，得二十五萬四千一十六寸，又三因，得七十六萬二千四十八寸，為法。除實，得三寸，為平地雨深。合問。\end{unitzh}
\begin{uniten}Working: Convert (tong) the depth of the basin and the diameters all into cun: the depth of the basin gives 18 cun and the base diameter gives 12 cun; multiplying them together gives 216 cun, which is the base rate. Set out the mouth diameter, 28 cun, and subtract the base diameter, 12 cun, leaving 16 cun, which is the difference; multiply it by the water depth, 9 cun, obtaining 144 cun; add the base rate, 216 cun, obtaining 360 cun, which is the surface rate. Take the depth of the basin, 18 cun, as divisor and divide the surface rate, obtaining 20 cun, which, converted (zhan), is 2 chi; this is the diameter of the water surface. Multiply the base rate, 216 cun, by the surface rate, 360 cun, obtaining 77760 cun, placed above. Multiply the base rate, 216 cun, by itself, obtaining 46656 cun, and add it to the above. Further multiply the surface rate, 360 cun, by itself, obtaining 129600; combining with the above, altogether one obtains 254016; multiply it by the water depth, 9 cun, obtaining 2286144 cun, which is the dividend. Multiply the depth of the basin, 18 cun, by the mouth diameter, 28 cun, obtaining 504 cun; multiplying it by itself gives 254016 cun; further tripling it gives 762048 cun, which is the divisor. Dividing the dividend gives 3 cun, which is the depth of rain on level ground. This agrees with the question.\end{uniten}

\begin{unitzh}\unitlabel{shushu-017.4}竹器驗雪\end{unitzh}
\begin{uniten}Testing Snow with a Bamboo Vessel\end{uniten}

\begin{unitzh}\unitlabel{shushu-017.5}問：以圓竹籮驗雪，籮口徑一尺六寸，深一尺七寸，底徑一尺二寸，雪䧏其中，髙一尺。籮體通風，受雪多，則平地少。欲知平地雪髙㡬何？\end{unitzh}
\begin{uniten}Question: A round bamboo basket (luo) is used to test snowfall. The mouth diameter of the basket is 1 chi 6 cun, its depth 1 chi 7 cun, and its base diameter 1 chi 2 cun. Snow falls into it to a height of 1 chi. The body of the basket lets the wind through; when it receives much snow, then level ground has less. It is desired to know how high the snow on level ground is.\end{uniten}

\begin{unitzh}\unitlabel{shushu-017.6}按：「籮體通風」一語，與算術不相涉。或籮口所降之雪歸於籮底，與前「天池測雨」題相同，然依上步算，平地雪深只七寸餘。今其數又不合，殆故為是語以誤人也。\end{unitzh}
\begin{uniten}Note: The phrase "the body of the basket lets the wind through" has nothing to do with the computation. Perhaps the snow falling at the mouth of the basket gathers at the bottom of the basket, which is the same as in the preceding problem "Measuring Rain with a Heavenly Pool"; but computing according to the steps above, the depth of snow on level ground is only 7-odd cun. Now, moreover, its number does not agree; probably the phrase was put in deliberately to mislead people.\end{uniten}

\begin{unitzh}\unitlabel{shushu-017.7}答曰：平地雪厚九寸三千四百二十九分之七百六十四。\end{unitzh}
\begin{uniten}Answer: The thickness of the snow on level ground is 9 764/3429 cun. [Translator's note: the working below gives 3439.]\end{uniten}

\begin{unitzh}\unitlabel{shushu-017.8}術曰：口徑減底徑，餘乗雪深，半之，自乗，為隅。以籮深冪乗雪深冪，併隅，又乗雪深冪，為實。隅實可約，約之，開連枝三乗方，得平地雪厚。\end{unitzh}
\begin{uniten}Method: Subtract the base diameter from the mouth diameter; multiply the remainder by the snow depth, halve it, and multiply it by itself to make the corner (yu). Multiply the basket-depth area (mi) by the snow-depth area, add the corner, and multiply again by the snow-depth area to make the dividend. When the corner and the dividend can be reduced, reduce (yue) them, extract the root of the linked-branch (lianzhi) thrice-multiplied square (san cheng fang) [Translator's note: i.e. a root of the fourth degree.], and obtain the thickness of the snow on level ground.\end{uniten}

\begin{unitzh}\unitlabel{shushu-017.9}草曰：列問數，各通為寸。置口徑一十六寸，減底徑一十二寸，餘四寸，乗雪深一十寸，得四十寸，以半之，得二十寸，自乗，得四百寸，為隅。以籮深一十七寸自乗，得二百八十九寸，為籮深冪。次置雪深一十寸自乗，得一百寸，為雪深冪，以乗籮深冪數，加隅，又乗深冪，得二百九十三萬寸，為實。隅實求等，得四百，俱約之，得七千三百二十五，為實；得一，為隅。開三乗方。歩法不可超，乃約實，置商九寸，與隅一相生，得九，為下亷；又與商相生，八十一寸，為上亷；又與商相生，得七百二十九，為從方。乃命上商除實，不盡七百六十四。已而復以商生隅，入二亷至方，陸續又生。畢，以方亷隅共併之，得三千四百三十九分寸之七百六十四，為平地雪厚九寸三千四百三十九分寸之七百六十四。合問。\end{unitzh}
\begin{uniten}Working: Lay out the numbers of the question, each converted into cun. Set out the mouth diameter, 16 cun, and subtract the base diameter, 12 cun, leaving 4 cun; multiply it by the snow depth, 10 cun, obtaining 40 cun; halve it, obtaining 20 cun; multiply it by itself, obtaining 400 cun, which is the corner. Multiply the basket depth, 17 cun, by itself, obtaining 289 cun, which is the basket-depth area. Next set out the snow depth, 10 cun, and multiply it by itself, obtaining 100 cun, which is the snow-depth area; multiply it by the number of the basket-depth area, add the corner, and multiply again by the depth area, obtaining 2930000 cun, which is the dividend. For the corner and the dividend seek the equal number (deng), obtaining 400; reduce both by it, obtaining 7325, which is the dividend, and obtaining 1, which is the corner. Extract the root of the thrice-multiplied square. The stepping of the divisor (bu fa) cannot leap (chao); so, gauging the dividend (yue shi), set the quotient (shang) at 9 cun; generating (sheng) it together with the corner, 1, gives 9, which is the lower edge (xia lian); generating again together with the quotient: 81 cun, which is the upper edge (shang lian); generating again together with the quotient gives 729, which is the accompanying square (cong fang). Then apply (ming) the quotient above to divide the dividend; there remains, unexhausted, 764. After that, again with the quotient generate the corner into the two edges and on to the square (fang), generating again in succession. When this is finished, add the square, edges and corner all together, obtaining 764/3439 cun; this makes the thickness of the snow on level ground 9 764/3439 cun. This agrees with the question.\end{uniten}

\begin{unitzh}\unitlabel{shushu-017.10}按：此法之意不可見，然以數考之，非通法也。設原題雪深為一寸，以口徑底徑較四寸乗雪深一寸，仍得四寸，半之，得二寸，自之，得四寸，為隅。以籮深一十七寸自之，得二百八十九寸，為籮深冪。雪深一寸自之，仍得一寸，為雪深冪。二深冪相乗，仍得二百八十九寸，併隅，得二百九十三寸，再以雪深冪乗之，仍得二百九十三寸，為實。隅實相約，得七十四寸二十五百分，為實；一為隅。開三乗方，得二寸又六千四百分寸之五千七百二十五。是平地雪反深於籮内矣。\end{unitzh}
\begin{uniten}Note: The meaning of this method cannot be seen; but examined by means of numbers, it is not a general method. Suppose the snow depth in the original problem were 1 cun. Multiply the difference (jiao) of the mouth diameter and the base diameter, 4 cun, by the snow depth, 1 cun: one still obtains 4 cun; halving it gives 2 cun; multiplying it by itself gives 4 cun, which is the corner. Multiply the basket depth, 17 cun, by itself, obtaining 289 cun, which is the basket-depth area. The snow depth, 1 cun, multiplied by itself still gives 1 cun, which is the snow-depth area. Multiplying the two depth areas together still gives 289 cun; adding the corner gives 293 cun; multiplying it again by the snow-depth area still gives 293 cun, which is the dividend. Reducing the corner and the dividend by each other gives 74 cun and 25 hundredths as the dividend, and 1 as the corner. Extracting the root of the thrice-multiplied square gives 2 cun and 5725/6400 cun. This means the snow on level ground would, on the contrary, be deeper than inside the basket.\end{uniten}

\begin{unitzh}\unitlabel{shushu-017.11}列問數，各通為十：口徑得一十六寸，深一十七寸，底徑一十二寸，籮中雪髙一十寸。乃以底徑減口徑，餘四寸，乗雪深一十寸，得四十寸，以中得數二十寸，自乗，得四百寸，為隅。以籮深一十七寸自乗，得二百八十九寸，為籮深冪。次置雪深一十寸自乗，得一百寸，為雪深冪。以雪深冪一百寸乗籮深冪二百八十九寸，得二萬八千九百寸，併隅四百寸，得二萬九千三百寸，為上。置頭位數二萬九千三百寸，又乗雪深冪一百寸，得二百九十三萬寸，為實，開三乗方。以隅實求等，得四百，俱為約之，得七千三百二十五，為實；一為隅，開之。\end{unitzh}
\begin{uniten}Lay out the numbers of the question, each converted into "ten" [Translator's note: 十 "ten" here evidently stands for 寸 "cun".]: the mouth diameter gives 16 cun, the depth 17 cun, the base diameter 12 cun, and the height of the snow in the basket 10 cun. Then subtract the base diameter from the mouth diameter, leaving 4 cun; multiply it by the snow depth, 10 cun, obtaining 40 cun; taking the middle (zhong), the number obtained is 20 cun; multiplying it by itself gives 400 cun, which is the corner. Multiply the basket depth, 17 cun, by itself, obtaining 289 cun, which is the basket-depth area. Next set out the snow depth, 10 cun, and multiply it by itself, obtaining 100 cun, which is the snow-depth area. Multiply the snow-depth area, 100 cun, by the basket-depth area, 289 cun, obtaining 28900 cun; add the corner, 400 cun, obtaining 29300 cun, which is placed above. Set out the number in the head position (tou wei), 29300 cun, and multiply it again by the snow-depth area, 100 cun, obtaining 2930000 cun, which is the dividend; extract the root of the thrice-multiplied square. With the corner and the dividend seek the equal number, obtaining 400; reduce both by it, obtaining 7325, which is the dividend; 1 is the corner; extract the root.\end{uniten}

\begin{unitzh}\unitlabel{shushu-017.12}歩法不可超，乃約實，置商九寸，與隅相生，得九，為下亷。下亷九又與商九相生，得八十一，為上亷。上亷又與商相生，得七百二十九，為從方。乃以從方七百二十九命上商九，除實七千三百二十五，訖，實餘七百六十四。既而後以商生隅，入下亷。下亷得一十八，又與啇九相生，入上亷。上亷得二百四十三，又與商相生，入方，得二千九百一十六。又以商九生隅一，入下亷一十八内，得二十七。又以商九生下亷二十七，入上亷二百四十三内，得四百八十六。又以商生隅，入下亷二十七内，得三十六，為求圖。乃以末圓方亷隅四者併之，得三千四百三十九，為母；以實餘七百六十四，為子。命為平地雪厚九寸三千四百三十九分寸之七百六十四。合問。\end{unitzh}
\begin{uniten}The stepping of the divisor cannot leap; so, gauging the dividend, set the quotient at 9 cun; generating it together with the corner gives 9, which is the lower edge. The lower edge, 9, is again generated together with the quotient, 9, giving 81, which is the upper edge. The upper edge is again generated together with the quotient, giving 729, which is the accompanying square. Then, with the accompanying square, 729, apply the quotient above, 9, and divide the dividend, 7325; when this is done, the remainder of the dividend is 764. After that, with the quotient generate the corner into the lower edge. The lower edge gives 18; again generate it together with the quotient, 9, into the upper edge. The upper edge gives 243; again generate it together with the quotient into the square, obtaining 2916. Again with the quotient, 9, generate the corner, 1, into the lower edge, 18, within which one obtains 27. Again with the quotient, 9, generate the lower edge, 27, into the upper edge, 243, within which one obtains 486. Again with the quotient generate the corner into the lower edge, 27, within which one obtains 36; this makes the "seeking diagram" (qiu tu). Then add together the four, the square, the edges and the corner, of the last circle (mo yuan), obtaining 3439, which is the denominator (mu); take the remainder of the dividend, 764, as the numerator (zi). [Translator's note: 求圖 and 末圓 appear garbled, perhaps for 末圖, "last diagram".] Express it as the thickness of the snow on level ground: 9 764/3439 cun. This agrees with the question.\end{uniten}

\begin{unitzh}\unitlabel{shushu-017.13}圓罌測雨\end{unitzh}
\begin{uniten}Measuring Rain with a Round Jar\end{uniten}

\begin{unitzh}\unitlabel{shushu-018.1}問：以圓罌接雨，口徑一尺五寸，腹徑二尺四寸，底徑八寸，深一尺六寸，並裏明。接得雨水深一尺二寸。圓法用宻率。問平地雨水深㡬何？\end{unitzh}
\begin{uniten}Question: A round jar (yuan ying) is used to catch rain. Its mouth diameter is 1 chi (尺) 5 cun (寸), its belly diameter 2 chi 4 cun, its base diameter 8 cun, and its depth 1 chi 6 cun, all inside clear measurements (li ming). The rainwater caught is 1 chi 2 cun deep. For the circle method (yuan fa), use the accurate rate (mi lü). What is the depth of the rainwater on level ground? [Translator's note: the working below uses 10 cun 5 fen for the mouth diameter.]\end{uniten}

\begin{unitzh}\unitlabel{shushu-018.2}按：此題問平地雨深，無闗圓法宻率，句贅。若求罌中雨積數，則當加此語。\end{unitzh}
\begin{uniten}Siku editors' note: This problem asks for the depth of rain on level ground, which has no bearing on the circle method and the accurate rate; the clause is superfluous. If the number for the volume of rain in the jar were sought, then this phrase ought to be added.\end{uniten}

\begin{unitzh}\unitlabel{shushu-018.3}答曰：平地雨深一尺八寸七萬四千八十八分寸之六萬四千四百八十三。\end{unitzh}
\begin{uniten}Answer: The depth of rain on level ground is 1 chi 8 64483/74088 cun.\end{uniten}

\begin{unitzh}\unitlabel{shushu-018.4}按：答數誤，改正見後。\end{unitzh}
\begin{uniten}Siku editors' note: The answer is wrong; the correction is given below.\end{uniten}

\begin{unitzh}\unitlabel{shushu-018.5}術曰：底徑與腹徑相乗，又各自乗，併之，乗半罌深，以一十一乗之，為下率；以四十二為下法，除得下積。以半罌深併雨深，減元罌深，餘為上深。以口徑減腹徑，餘乗上深，為次。以半罌深乗口徑，加次，為靣率。以半深除靣率，得水靣徑。以半深乗腹徑，為腹率。置靣率與腹率相乗，又各自乗，併之，以一十一乗之，為上率。以半深自乗為冪，以乗下法，為上法。上法除上率，得上積。半深冪乗下率，併上率，為總實。口徑冪乗上法，為總法。除實，得平地雨髙。\end{unitzh}
\begin{uniten}Method: Multiply the base diameter and the belly diameter together; also square each; combine them; multiply by half the jar depth; multiply by 11: this is the lower rate (xia lü). Take 42 as the lower divisor (xia fa); dividing gives the lower volume (xia ji). Add half the jar depth to the rain depth and subtract the original jar depth; the remainder is the upper depth (shang shen). Subtract the mouth diameter from the belly diameter; multiply the remainder by the upper depth: this is the secondary (ci). Multiply the mouth diameter by half the jar depth and add the secondary: this is the surface rate (mian lü). Divide the surface rate by the half depth to get the diameter of the water surface. Multiply the belly diameter by the half depth: this is the belly rate (fu lü). Set up the surface rate and the belly rate and multiply them together; also square each; combine them; multiply by 11: this is the upper rate (shang lü). Square the half depth to make the area (mi); multiply the lower divisor by it: this is the upper divisor (shang fa). Divide the upper rate by the upper divisor to get the upper volume (shang ji). Multiply the lower rate by the area of the half depth and add the upper rate: this is the total dividend (zong shi). Multiply the upper divisor by the area of the mouth diameter: this is the total divisor (zong fa). Dividing the dividend gives the height of rain on level ground.\end{uniten}

\begin{unitzh}\unitlabel{shushu-018.6}草曰：置底徑八寸，與腹徑二十四寸相乗，得一百九十二寸於上；又底徑八寸自乗，得六十四寸，加上；又腹徑二十四寸自乗，得五百七十六寸，併上，共得八百三十二寸。以乗半罌深八寸，得六千六百五十六寸。又以一十一乗之，得七萬三千二百一十六寸，為下率。{\zhsmall 按：此下法不合，皆為題中圖法句所誤。}置宻率法一十四，以所併三因之，得四十二，為下法。\end{unitzh}
\begin{uniten}Working: Set up the base diameter 8 cun and multiply it by the belly diameter 24 cun to get 192 cun in the upper position. Also square the base diameter 8 cun to get 64 cun and add it to the upper position. Also square the belly diameter 24 cun to get 576 cun and combine it with the upper position: altogether 832 cun. Multiply half the jar depth 8 cun by it to get 6656 cun. Again multiply it by 11 to get 73216 cun: this is the lower rate. [Siku editors: This lower divisor is not right; all of this was misled by the clause on the "diagram method" (tu fa) in the problem.] [Translator's note: 圖 "diagram" is evidently an error for 圓 "circle".] Set up the divisor of the accurate rate, 14; on account of what was combined, multiply it by 3 (san yin) to get 42: this is the lower divisor.\end{uniten}

\begin{unitzh}\unitlabel{shushu-018.7}以半深八寸併雨深一十二寸，得二十寸，以減元深一十六寸，餘四寸，為上深。以口徑一十寸五分減腹徑二十四寸，餘一十三寸五分，以乗上深四寸，得五十四寸，為次。以半罌深八寸乗口徑一十寸五分，得八十四寸，加次，共得一百三十八寸，為靣率。以半深八寸乗腹徑二十四寸，得一百九十二寸，為腹率。置靣率一百三十八寸，與腹率一百九十二寸相乗，得二萬六千四百九十六寸於上；又以靣率一百三十八寸自乗，得一萬九十四十四，加上；又以腹率一百九十二寸自乗，得三萬六千八百六十四，併上，共得八萬二千四百四寸。{\zhsmall 按：此條内落以上髙四寸乗之一層。}以一十一乗之，得九十萬六千四百四十四寸，為上率。\end{unitzh}
\begin{uniten}Add the half depth 8 cun to the rain depth 12 cun to get 20 cun; from it subtract the original depth 16 cun, leaving 4 cun: this is the upper depth. Subtract the mouth diameter 10 cun 5 fen (分) from the belly diameter 24 cun, leaving 13 cun 5 fen; multiply the upper depth 4 cun by it to get 54 cun: this is the secondary. Multiply the mouth diameter 10 cun 5 fen by half the jar depth 8 cun to get 84 cun; add the secondary, altogether getting 138 cun: this is the surface rate. Multiply the belly diameter 24 cun by the half depth 8 cun to get 192 cun: this is the belly rate. Set up the surface rate 138 cun and multiply it by the belly rate 192 cun to get 26496 cun in the upper position. Also square the surface rate 138 cun to get "1 ten-thousand 9 tens 4 tens 4" [Translator's note: the numeral 一萬九十四十四 is garbled, evidently for 一萬九千四十四, 19044.], and add it to the upper position. Also square the belly rate 192 cun to get 36864 and combine it with the upper position: altogether 82404 cun. [Siku editors: In this item the one step of multiplying it by the upper height 4 cun has dropped out.] Multiply it by 11 to get 906444 cun: this is the upper rate.\end{uniten}

\begin{unitzh}\unitlabel{shushu-018.8}以半深八寸自乗，得六十四寸，為半深冪；以乗下法四十二，得二千六百八十八，為上法。以半深冪六十四寸乗下率七萬三千二百一十六寸，得四百六十八萬五千八百二十四寸，併上率九十萬六千四百四十四，共得五百五十九萬二千二百六十八寸，為總實。以口徑一十寸五分自乗，得一百一十寸二分五厘，以乗上法二千六百八十八寸，得二十九萬六千三百五十二寸，為總法。除實，得一尺八寸，不盡二十五萬七千九百三十二，與法求等，得四，俱約之，為一尺八寸七萬四千八十八分寸之六萬四千四百八十三，為平地雨深。合問。\end{unitzh}
\begin{uniten}Square the half depth 8 cun to get 64 cun: this is the area (mi) of the half depth. Multiply the lower divisor 42 by it to get 2688: this is the upper divisor. Multiply the lower rate 73216 cun by the area of the half depth, 64 cun, to get 4685824 cun; add the upper rate 906444, altogether getting 5592268 cun: this is the total dividend. Square the mouth diameter 10 cun 5 fen to get 110 cun 2 fen 5 li (厘); multiply the upper divisor 2688 cun by it to get 296352 cun: this is the total divisor. Dividing the dividend gives 1 chi 8 cun, with 257932 not exhausted (bu jin). Seek the equal (qiu deng) of this and the divisor and obtain 4; reduce both by it, giving 1 chi 8 64483/74088 cun: this is the depth of rain on level ground. This matches the question.\end{uniten}

\begin{unitzh}\unitlabel{shushu-018.9}按：此法有二誤。法、實皆當用圓冪，或皆用方冪；今以圓冪率乗實，方冪率乗法，法實不同類，一誤也。罌内雨自腹徑截之，為雨圓䑓體，下髙八寸，上髙四寸；於下體併三冪，以髙乗之，於上體只併三冪，未以髙乗之，二誤也。有此二誤，故得平地雨深少三十五分之十七。今依本法改正於後。\end{unitzh}
\begin{uniten}Siku editors' note: This method has two errors. The divisor and the dividend should both use the circle area (yuan mi), or both use the square area (fang mi). Now the dividend is multiplied by the rate of the circle area and the divisor by the rate of the square area, so that divisor and dividend are not of the same kind: this is the first error. The rain inside the jar, cut off at the belly diameter, forms round-platform solids (yuan tai ti) of rain, the lower 8 cun high and the upper 4 cun high. For the lower solid the three areas (mi) were combined and multiplied by the height; for the upper solid the three areas were only combined and not multiplied by the height: this is the second error. Because of these two errors, the depth of rain on level ground obtained is short by 17/35. Now, following the original method, it is corrected below.\end{uniten}

\begin{unitzh}\unitlabel{shushu-018.10}法：以腹徑、底徑相乗，又各自乗，併三積，以半罌深八寸乗之，得六千六百五十六寸，為三倍方罌内腹下雨積。又以口徑、腹徑相減，餘一十三寸五分；以雨深減罌深，餘四寸；相乗，以半罌深除之，得六寸七分五厘，與口徑相加，得一十七寸二分五厘，為雨靣徑。與腹徑相乗，又各自乗，併三積，以雨上深四寸乗之，得五千一百五十寸二五，為方罌内三倍腹上雨積。併二雨積，得一萬一千八百零六寸二五，為方罌内三倍共雨積，為實。口徑自乗，三因，得三百三十寸七五，為法。除實，得三尺五寸又一千三百二十三分寸之九百二十，為平地雨深。\end{unitzh}
\begin{uniten}Rule: Multiply the belly diameter and the base diameter together, also square each, combine the three areas, and multiply by half the jar depth 8 cun to get 6656 cun: this is 3 times the volume of rain below the belly in a square jar (fang ying). Also subtract the mouth diameter and the belly diameter from each other, leaving 13 cun 5 fen; subtract the rain depth from the jar depth, leaving 4 cun; multiply these together and divide by half the jar depth to get 6 cun 7 fen 5 li; add this to the mouth diameter to get 17 cun 2 fen 5 li: this is the diameter of the rain surface. Multiply it and the belly diameter together, also square each, combine the three areas, and multiply by the upper rain depth 4 cun to get 5150.25 cun: this is, in the square jar, 3 times the volume of rain above the belly. Combine the two rain volumes to get 11806.25 cun: this is, in the square jar, 3 times the total rain volume; it is the dividend. Square the mouth diameter and multiply by 3 to get 330.75 cun: this is the divisor. Dividing the dividend gives 3 chi 5 cun and 920/1323 cun: this is the depth of rain on level ground.\end{uniten}

\begin{unitzh}\unitlabel{shushu-018.11}若不先用除，則以口徑、腹徑較與罌深、雨深較相乗之五十四寸，為雨靣徑口徑較加一半罌乗之數{\zhsmall 應以半罌除之，得雨徑較；今不除，即如雨徑較以半罌乗之。}，即為雨靣徑口徑較。此數既加一半罌乗，則諸數皆以半罌乗之，得口徑八十四寸，腹徑一百九十二寸。以口徑與雨靣徑口徑較相加，得雨靣徑一百三十八寸，與腹徑相乗，又各自乗，併三冪，以腹上雨深四寸乗之，得三十二萬九千六百一十六寸，為三倍上雨積。又以半罌深冪乗前三倍下雨積，得四十二萬五千九百八十四寸，為三倍下雨積。併二積，得七十五萬五千六百寸，為三倍共雨積，為實。以半罌深冪乗三因口徑冪，得二萬一千一百六十八寸，為法。除之，得數亦同。\end{unitzh}
\begin{uniten}If division is not used first, then the 54 cun from multiplying together the difference (jiao) of the mouth and belly diameters and the difference of the jar depth and the rain depth is the number of the difference of the rain-surface diameter and the mouth diameter with one multiplication by half the jar added [Note: It ought to be divided by half the jar to get the rain-diameter difference; now it is not divided, so it is just as if the rain-diameter difference had been multiplied by half the jar.]; this, then, serves as the difference of the rain-surface diameter and the mouth diameter. Since this number has had one multiplication by half the jar added, all the numbers are multiplied by half the jar, giving the mouth diameter 84 cun and the belly diameter 192 cun. Add the mouth diameter and the difference of the rain-surface diameter and the mouth diameter to get the rain-surface diameter 138 cun; multiply it and the belly diameter together, also square each, combine the three areas (mi), and multiply by the depth of rain above the belly, 4 cun, to get 329616 cun: this is 3 times the upper rain volume. Also multiply the earlier 3 times the lower rain volume by the area of half the jar depth to get 425984 cun: this is 3 times the lower rain volume. Combine the two volumes to get 755600 cun: this is 3 times the total rain volume; it is the dividend. Multiply 3 times the area of the mouth diameter by the area of half the jar depth to get 21168 cun: this is the divisor. Dividing, the number obtained is also the same.\end{uniten}

\begin{unitzh}\unitlabel{shushu-018.12}峻積驗雪\end{unitzh}
\begin{uniten}Verifying snow with a steep accumulation\end{uniten}

\begin{unitzh}\unitlabel{shushu-019.1}問：驗雪占年，墻髙一丈二尺，倚木去址五尺，梢與墻齊，木身積雪厚四寸。峻積薄，平積厚，欲知平地雪厚㡬何？答曰：平地雪厚一尺四分。\end{unitzh}
\begin{uniten}Question: Examining snow to prognosticate the year. A wall is 1 zhang (丈) 2 chi (尺) high. A leaning pole is 5 chi from the foot [of the wall], and its tip is level with the wall. The snow accumulated on the body of the pole is 4 cun (寸) thick. Snow accumulated on a steep [surface] is thin; snow accumulated on the level is thick. One wishes to know how thick the snow on level ground is. Answer: The snow on level ground is 1 chi 4 fen (分) thick.\end{uniten}

\begin{unitzh}\unitlabel{shushu-019.2}術曰：以少廣求之，連枝入之。以去址自乗為隅；以墻髙自乗，併隅於上；以雪厚自之，乗上為實。{\zhsmall 可約者約而開之。}開連枝平方，得地雪厚。\end{unitzh}
\begin{uniten}Method: Seek it by "small width" (shaoguang), entering it as "joined branches" (lianzhi). Multiply the distance from the foot by itself to make the corner (yu). Multiply the wall height by itself and add the corner, [placing the result] in the upper [position]. Multiply the snow thickness by itself, and multiply the upper [position] by it to make the dividend (shi). [Note: Where they can be reduced, reduce them and then extract.] Extract the joined-branch square root to obtain the thickness of the snow on the ground.\end{uniten}

\begin{unitzh}\unitlabel{shushu-019.3}草曰：以問數皆通為寸。置去址五十寸，自乗得二千五百，為隅；以墻髙一百二十寸自乗，得一萬四千四百寸，併隅，得一萬六千九百寸於上；以雪厚四寸自之，得一十六，乗上，得二十七萬四百寸，為實。開連枝平方：今隅、實可求等，得一百，俱約之，得二千七百四為實，得二十五為隅。開平方，得一十寸四分，展為一尺四分，為平地雪厚。合問。\end{unitzh}
\begin{uniten}Working: Convert all the numbers of the question into cun. Set up the distance from the foot, 50 cun. Multiplied by itself, it gives 2500, as the corner. Multiply the wall height, 120 cun, by itself to obtain 14400 cun. Adding the corner gives 16900 cun, [placed] in the upper [position]. Multiply the snow thickness, 4 cun, by itself to obtain 16. Multiplying the upper [position] by it gives 270400 cun, as the dividend. Extract the joined-branch square root. Here one can seek the equal [number] (qiu deng) for the corner and the dividend, obtaining 100. Reducing both by it gives 2704 as the dividend and 25 as the corner. Extracting the square root gives 10 cun 4 fen; converted (zhan), this is 1 chi 4 fen, as the thickness of the snow on level ground. This matches the question.\end{uniten}

\begin{unitzh}\unitlabel{shushu-019.4}按：此術理法皆確，然實用勾股；不曰勾股而曰少廣、曰連枝者，猶有所閉匿而不肯盡發也。試以圗明之：甲乙為牆上雪厚，即平地雪厚；乙丙為木上雪厚。甲乙丙勾股形與木倚牆所成勾股形同式：牆髙為大股，木為大弦，木去址為大勾；甲乙為小弦，甲丙為小股，乙丙為小勾。以牆髙大股自乗、木去址大勾自乗，併之為大弦冪，為實；以木上雪厚乙丙小勾冪乗之，以木去址大勾冪除之，得甲乙小弦冪；開平方，即為平地雪厚也。\end{unitzh}
\begin{uniten}Note: The reasoning and the procedure of this method are both correct, but in fact it uses base and altitude (gougu). That it does not say "base and altitude" but says "small width" and "joined branches" is as though there were something kept shut away and hidden that [he] was unwilling to disclose fully. Let us try to make this clear with a diagram. AB (甲乙) is the thickness of the snow on the wall, that is, the thickness of the snow on level ground; BC (乙丙) is the thickness of the snow on the pole. The base-altitude figure (gougu xing) ABC (甲乙丙) and the base-altitude figure formed by the pole leaning against the wall are of the same form (tongshi). The wall height is the large altitude, the pole is the large hypotenuse, and the pole's distance from the foot is the large base. AB is the small hypotenuse, AC (甲丙) is the small altitude, and BC is the small base. Multiply the wall height, the large altitude, by itself, and the pole's distance from the foot, the large base, by itself. Adding them makes the area (mi) of the large hypotenuse, as the dividend (shi). Multiply it by the area of the snow thickness on the pole, the small base BC. Divide it by the area of the pole's distance from the foot, the large base, to obtain the area of the small hypotenuse AB. Extracting the square root gives the thickness of the snow on level ground.\end{uniten}

\begin{unitzh}\unitlabel{shushu-019.5}數學九章卷二下\end{unitzh}
\begin{uniten}Shuxue jiuzhang (Mathematical Treatise in Nine Sections), juan 2, part 2\end{uniten}

\section*{卷3上}\addcontentsline{toc}{section}{卷3上}

\begin{unitzh}\unitlabel{shushu-020.1}欽定四庫全書\end{unitzh}
\begin{uniten}Imperially Commissioned Complete Library of the Four Treasuries (Qinding Siku quanshu)\end{uniten}

\begin{unitzh}\unitlabel{shushu-020.2}数學九章卷三上 宋 秦九韶 撰\end{unitzh}
\begin{uniten}Shuxue jiuzhang, juan 3, part 1 (shang). Composed by Qin Jiushao of the Song.\end{uniten}

\begin{unitzh}\unitlabel{shushu-020.3}田域\end{unitzh}
\begin{uniten}Fields and Domains (tianyu)\end{uniten}

\begin{unitzh}\unitlabel{shushu-020.4}按：此卷以方圓斜直冪積相求，即方田、少廣、勾股諸法；而術中累乘累除，錯綜變換，與常法迴然，其本則出於立天元一法。今擇其難解者，以立天元一法明之，皆不攻自破矣。\end{unitzh}
\begin{uniten}Note: In this juan, squares, circles, slants and straights, and areas (mi, ji) are sought from one another; these are the methods of Rectangular Fields (Fangtian), Lesser Breadth (Shaoguang) and Base and Altitude (Gougu). But the repeated multiplications and repeated divisions in the methods, intricately interwoven and transformed, are far removed from the ordinary methods; their basis, however, comes from the method of setting up the celestial unknown (tianyuan). Now we select those that are difficult to understand and clarify them by the method of setting up the celestial unknown; then all of them break down of themselves without being attacked.\end{uniten}

\begin{unitzh}\unitlabel{shushu-020.5}古池推元\end{unitzh}
\begin{uniten}Inferring the Original of an Ancient Pond\end{uniten}

\begin{unitzh}\unitlabel{shushu-020.6}問：有方中古圓池，堙圮，北餘一角，從外方隔斜至内圓邊七尺六寸。欲就古跡修之，欲求圓、方、方斜各幾何？荅曰：池圓徑三丈六尺六寸四百二十九分寸之四百一十二；方面三丈六尺六寸四百二十九分寸之四百一十二；方斜五丈一尺八寸四百二十九分寸之四百一十二。\end{unitzh}
\begin{uniten}Question: There is an ancient round pond within a square. It has silted up and collapsed; on the north one corner remains. From the outer square, diagonally across (gexie) to the edge of the inner circle, it is 7 chi (尺) 6 cun (寸). One wishes to repair it following the ancient traces, and wishes to find the circle, the square and the diagonal of the square: how much is each? Answer: The diameter of the pond's circle is 3 zhang (丈) 6 chi 6 412/429 cun; the side of the square is 3 zhang 6 chi 6 412/429 cun; the diagonal of the square (fangxie) is 5 zhang 1 chi 8 412/429 cun.\end{uniten}

\begin{unitzh}\unitlabel{shushu-020.7}術曰：以少廣求之，投胎術{\zhsmall 按即益積之名}入之。斜自乘，倍之，為實；倍斜，為益方；以半寸為從隅。開投胎平方，得徑，又為方面。以隅併之，共為方斜。\end{unitzh}
\begin{uniten}Method: Seek it by Lesser Breadth (shaoguang), entering it by the casting-into-the-womb method (toutai shu) [Siku editors: That is, a name for "augmenting the area" (yiji).]. Multiply the diagonal by itself and double it, making the dividend (shi); double the diagonal, making the augmenting square (yifang); take half a cun as the following corner (congyu). Extract the square root by casting into the womb, obtaining the diameter, which is also the side of the square. Add the corners to it; together they make the diagonal of the square.\end{uniten}

\begin{unitzh}\unitlabel{shushu-020.8}草曰：以斜七十六寸自乘，得五千七百七十六；倍之，得一萬一千五百五十二寸，為實。倍斜七十六寸，得一百五十二，為益方{\zhsmall 按有長方積，先求長，其長濶較名益方}。以半寸為從隅，開平方。置實一萬一千五百五十二於上，益方一百五十二於中，從隅五分於下。於下起歩，約得百。\end{unitzh}
\begin{uniten}Working: Multiply the diagonal, 76 cun, by itself, obtaining 5776; double it, obtaining 11,552 cun, which is the dividend. Double the diagonal, 76 cun, obtaining 152, which is the augmenting square [Siku editors: When there is the area of a rectangle and one first seeks the length, the difference (jiao) of its length and width is called the augmenting square.]. Take half a cun as the following corner and extract the square root. Place the dividend 11,552 at the top, the augmenting square 152 in the middle, and the following corner, 5 fen (分), at the bottom. Starting the steps (qibu) from the bottom, the estimate reaches the hundreds.\end{uniten}

\begin{unitzh}\unitlabel{shushu-020.9}古池圖\end{unitzh}
\begin{uniten}Diagram of the ancient pond.\end{uniten}

\begin{unitzh}\unitlabel{shushu-020.10}乃於實上商置三百寸。方再進，為一萬五千二百{\zhsmall 按再進者，以百乘之也}；隅五進為五千{\zhsmall 按隅五分以百再乗得五千}。以商隅相生，得一萬五千，為正方，以消益方一萬五千二百，其益方餘二百；以與商相生，得六百，投入實，得一萬二千一百五十二。又商隅相生，又得正方一萬五千，内消負方二百訖，餘一萬四千八百，為從方{\zhsmall 按倍正方減益方之數}。一退，為一千四百八十；以隅再退，為五十。\end{unitzh}
\begin{uniten}Then, above the dividend, place 300 cun as the quotient (shang). The square, advanced (jin) twice, becomes 15,200 [Siku editors: "Advancing twice" means multiplying it by 100.]; the corner 5, advanced, becomes 5000 [Siku editors: The corner, 5 fen, multiplied by 100 twice, gives 5000.]. Multiply the quotient and the corner together (xiangsheng), obtaining 15,000, which is the positive square (zhengfang); with it cancel (xiao) the augmenting square 15,200, and the augmenting square leaves 200. Multiply this with the quotient, obtaining 600, and cast it into the dividend, obtaining 12,152. Again multiply the quotient and the corner together, again obtaining the positive square 15,000; within it cancel the negative square (fufang) 200; when this is done, the remainder, 14,800, is the following square (congfang) [Siku editors: The number of twice the positive square minus the augmenting square.]. Moving it back (tui) one place makes 1480; moving the corner back two places makes 50.\end{uniten}

\begin{unitzh}\unitlabel{shushu-020.11}乃于上商之次續商置六十寸，與隅相生，増入正方，得一千七百八十。乃命續商除實，訖，實餘一千四百七十四。次以商生隅，増又正方，為二千八十。方一退，為二百八；隅再退，為五分。乃於續商之次又商置六寸，與隅相生，増入正方，為二百一十一。乃命商除實，訖，實不盡二百六寸，不開，為分子。\end{unitzh}
\begin{uniten}Then, next to the upper quotient, place 60 cun as the continued quotient (xushang); multiply it with the corner and add it into the positive square, obtaining 1780. Then use the continued quotient to divide the dividend; when this is done, the dividend leaves 1474. Next multiply the corner by the quotient (sheng) and add it again to the positive square, making 2080. The square, moved back one place, makes 208; the corner, moved back two places, makes 5 fen. Then, next to the continued quotient, again place 6 cun in the quotient; multiply it with the corner and add it into the positive square, making 211. Then use the quotient to divide the dividend; when this is done, 206 cun of the dividend remain unexhausted; they are not extracted and make the numerator (fenzi).\end{uniten}

\begin{unitzh}\unitlabel{shushu-020.12}乃以商生隅，増入正方，又併隅，共得二百一十四寸五分，為分母。以分母分子求等，得五分，為等數。皆以五分約其分母分子之數，為四百二十九分寸之四百一千二。通命之，得池圓徑及方面皆三丈六尺六寸四百二十九分寸之四百一十二。又倍隅斜七尺六十，得一丈五尺二寸，併徑三丈六尺六寸，共得五丈一尺八寸四百二十九分寸之四百一十二，為方斜。\end{unitzh}
\begin{uniten}Then multiply the corner by the quotient and add it into the positive square, and further add the corner, altogether obtaining 214 cun 5 fen, which is the denominator (fenmu). Seek the equal (qiu deng) between the denominator and the numerator, obtaining 5 fen as the equal number (dengshu). Reduce (yue) the numbers of the denominator and the numerator both by 5 fen, which makes "four hundred one thousand two"/429 cun [Translator's note: the numerator 四百一千二 is garbled; elsewhere it is 412.]. Stating them together (tongming), one obtains the diameter of the pond's circle and the side of the square, both 3 zhang 6 chi 6 412/429 cun. Further, double the corner diagonal, 7 chi 60, obtaining 1 zhang 5 chi 2 cun; add the diameter, 3 zhang 6 chi 6 cun, altogether obtaining 5 zhang 1 chi 8 412/429 cun, which is the diagonal of the square.\end{uniten}

\begin{unitzh}\unitlabel{shushu-020.13}按：此術以立天元一法明之。法：立天元一為池徑，即方邊；自之，得一平方，為方冪；倍之，得二平方，為斜冪，寄左。次倍斜至歩，加天元一，得一百五十二寸多一元，為方斜；自之，得二萬三千一百零四寸多三百零四元多一平方，亦為斜冪。與左相消，雨邊各減一平方，得二萬三千一百零四寸多三百零四元與一平方等。寸數為實，元數為較。\end{unitzh}
\begin{uniten}Note: This method is clarified by the method of setting up the celestial unknown. The method: set up the celestial unknown as the diameter of the pond, which is the side of the square. Multiplying it by itself gives 1 square power (pingfang), which is the area (mi) of the square; doubling it gives 2 square powers, which is the area of the diagonal; set it aside on the left (jizuo). Next, double the diagonal "reaching the bu" [Translator's note: 至歩 is unclear here, apparently corrupt.] and add the celestial unknown, obtaining 152 cun plus 1 yuan (unknown), which is the diagonal of the square; multiplying it by itself gives 23,104 cun plus 304 yuan plus 1 square power, which is also the area of the diagonal. Cancel it against the left (xiangxiao): subtracting 1 square power from each of the "rain" sides [Translator's note: 雨 "rain" is evidently written for 兩 "both".], one obtains 23,104 cun plus 304 yuan equal to 1 square power. The number of cun is the dividend; the number of yuan is the difference.\end{uniten}

\begin{unitzh}\unitlabel{shushu-020.14}或兩邊各半之，得一萬一千五百五十二寸多一百五十二元多半平方與一平方等，寸數為實，元數、半方數共為較。術中所用，葢次數也，然不如前數之便。至開方法，即有長方積、有長濶軗、帶縱先求長之法也。\end{unitzh}
\begin{uniten}Or, halving each of the two sides, one obtains 11,552 cun plus 152 yuan plus half a square power equal to 1 square power; the number of cun is the dividend, and the number of yuan and the number of half squares together make the difference. What is used in the method is probably the second set of numbers, but it is not as convenient as the former numbers. As for the root extraction, it is the method of having the area of a rectangle and having the 軗 of length and width [Translator's note: 軗 is evidently written for 較 "difference".], and first seeking the length with an attached extension (daizong).\end{uniten}

\begin{unitzh}\unitlabel{shushu-020.15}尖田求積\end{unitzh}
\begin{uniten}Finding the Area of a Pointed Field\end{uniten}

\begin{unitzh}\unitlabel{shushu-020.16}問：有兩尖田一叚，其尖長不等，兩大斜三十九歩，兩小斜二十五歩，中廣三十歩，欲知其積幾何？荅曰：曰積八百四十歩。\end{unitzh}
\begin{uniten}Question: There is one plot (duan) of a field pointed at both ends, whose points are unequal in length. The two large slants (xie) are 39 bu (歩), the two small slants 25 bu, and the middle breadth 30 bu. One wishes to know how much its area is. Answer: The area is 840 bu.\end{uniten}

\begin{unitzh}\unitlabel{shushu-020.17}術曰：以少廣求之，翻法入之。置半廣自乘，為半冪；與小斜冪相減，相乘，為小率。以半冪與大斜冪相減，相乘，為大率。以二率相減，餘自乘，為實。併二率，倍之，為從上亷。以一為益隅，開翻法三乘方，得積{\zhsmall 一位開盡者，不用翻法}。\end{unitzh}
\begin{uniten}Method: Seek it by Lesser Breadth, entering it by the turning-over method (fanfa). Set up half the breadth and multiply it by itself, making the half area (banmi); subtract it and the area of the small slant from each other and multiply them together, making the small rate (lü). Subtract the half area and the area of the large slant from each other and multiply them together, making the large rate. Subtract the two rates from each other and multiply the remainder by itself, making the dividend. Add the two rates and double the sum, making the following upper side (cong shanglian). Take 1 as the augmenting corner (yiyu), and extract the thrice-multiplied root (sancheng fang) by the turning-over method, obtaining the area [Note: Where it is completely extracted at one place, the turning-over method is not used.].\end{uniten}

\begin{unitzh}\unitlabel{shushu-020.18}草曰：置廣三十歩，以半之，得一十五；自乘，得二百二十五，為半冪。以小斜二十五歩自乘，得六百二十五，為小斜冪；與半冪相減，餘四百；與半冪二百二十五相乘，得九萬歩，為小率。置大斜三十九歩自乘，得一千五百二十一，為大斜冪；與半冪二百二十五相减，餘一千二百九十六；與半冪二百二十五相乘，得二十九萬一千六百，為大率。以小率九萬減大率，餘二十萬一千六百；自乘，得四百六億四千二百五十六萬，為實。以小率九萬併大率二十九萬一千六百，得三十八萬一千六百；倍之，得七十六萬三千二百，為從上亷{\zhsmall 按從上亷平方和數也}。\end{unitzh}
\begin{uniten}Working: Set up the breadth, 30 bu, and halve it, obtaining 15; multiply it by itself, obtaining 225, which is the half area. Multiply the small slant, 25 bu, by itself, obtaining 625, which is the area of the small slant; subtract it and the half area from each other, leaving 400; multiply this with the half area 225, obtaining 90,000 bu, which is the small rate. Set up the large slant, 39 bu, and multiply it by itself, obtaining 1521, which is the area of the large slant; subtract it and the half area 225 from each other, leaving 1296; multiply this with the half area 225, obtaining 291,600, which is the large rate. Subtract the small rate 90,000 from the large rate, leaving 201,600; multiply it by itself, obtaining 40,642,560,000, which is the dividend. Add the small rate 90,000 to the large rate 291,600, obtaining 381,600; double it, obtaining 763,200, which is the following upper side [Siku editors: The following upper side is the sum number of the square powers.].\end{uniten}

\begin{unitzh}\unitlabel{shushu-020.19}以一為益隅，開玲瓏翻法三乘方。歩法：乃以從亷超一位，益隅超三位，約商得十；今再超進，乃商置百。其從上亷為七十六億三千二百萬，其益隅為一億。約實，置商八百，為定商。以商生益隅，得八億，為益下亷；又以商生下亷，得六十四億，為益上亷；與從上亷七十六億三千二百萬相消，從上亷餘十二億三千二百萬；又與商相生，得九十八億五千六百萬，為從方；又與商相生，得七百八十八億四千八百萬，為正積；與元實四百六億四千二百五十六萬相消，正積餘三百八十二億五百四十四萬，為正實。\end{unitzh}
\begin{uniten}Take 1 as the augmenting corner and extract the thrice-multiplied root by the openwork (linglong) turning-over method. Stepping method (bufa): let the following side leap over (chao) one place and the augmenting corner leap over three places; the estimated quotient is in the tens. Now leap forward again, and then place the quotient in the hundreds. The following upper side becomes 7,632,000,000 and the augmenting corner becomes 100,000,000. Estimating from the dividend, place 800 in the quotient as the fixed quotient (dingshang). Multiply the augmenting corner by the quotient, obtaining 800,000,000, which is the augmenting lower side (yi xialian); again multiply the lower side by the quotient, obtaining 6,400,000,000, which is the augmenting upper side (yi shanglian); cancel it against the following upper side 7,632,000,000, and the following upper side leaves 1,232,000,000. Again multiply this with the quotient, obtaining 9,856,000,000, which is the following square; again multiply it with the quotient, obtaining 78,848,000,000, which is the positive area (zhengji); cancel it against the original dividend (yuanshi) 40,642,560,000, and the positive area leaves 38,205,440,000, which is the positive dividend (zhengshi).\end{uniten}

\begin{unitzh}\unitlabel{shushu-020.20}又以益隅一億與商相生，得八億，増入益下亷，為一十六億；又以益下亷與商相生，得一百二十八億，為益上亷；乃以益上亷與從上亷一十二億三千二百萬相消，餘一百一十五億六千八百萬，為益上亷；又與商相生，得九百二十五億四千四百萬，為益方；與從方九十八億五千六百萬相消，益餘八百二十六億八千八百萬，為益方。又以商生益隅一億，得八億，増入益下亷，得二十四億；又以商相生，得一百九十二億，入益上亷，得三百七億六千八百萬，為益上亷。又以商生益隅一億，得八億，入益下亷，得三十二億，畢。\end{unitzh}
\begin{uniten}Again multiply the augmenting corner, 100,000,000, with the quotient, obtaining 800,000,000; add it into the augmenting lower side, making 1,600,000,000. Again multiply the augmenting lower side with the quotient, obtaining 12,800,000,000, which is the augmenting upper side; then cancel the augmenting upper side against the following upper side, 1,232,000,000, leaving 11,568,000,000, which is the augmenting upper side. Again multiply it with the quotient, obtaining 92,544,000,000, which is the augmenting square; cancel it against the following square, 9,856,000,000; the augmenting leaves 82,688,000,000, which is the augmenting square. Again multiply the augmenting corner, 100,000,000, by the quotient, obtaining 800,000,000; add it into the augmenting lower side, obtaining 2,400,000,000. Again multiply with the quotient, obtaining 19,200,000,000; add it into the augmenting upper side, obtaining 30,768,000,000, which is the augmenting upper side. Again multiply the augmenting corner, 100,000,000, by the quotient, obtaining 800,000,000; add it into the augmenting lower side, obtaining 3,200,000,000. Done.\end{uniten}

\begin{unitzh}\unitlabel{shushu-020.21}其益方一退，為八十二億六千八百八十萬；益上亷再退，得三億七百六十八萬；益十亷三退，得三百二十萬；益隅四退，為一萬，畢。乃約正實，續置置四十歩，與益隅一萬相生，得四萬，入益下亷，為三百二十四萬；又與商相生，得一千二百九十六萬，入益上亷内，為三億二千六十四萬；又與商相生，得一十二億八千二百五十六萬，入從方内，為九十五億五千一百三十六萬。乃命上續啇四十除實，適盡。所得八百四十歩，為田積。今列求率開方圖於後。\end{unitzh}
\begin{uniten}The augmenting square, moved back one place, becomes 8,268,800,000; the augmenting upper side, moved back two places, gives 307,680,000; the augmenting "ten" side, moved back three places, gives 3,200,000 [Translator's note: 十 "ten" is evidently written for 下 "lower".]; the augmenting corner, moved back four places, becomes 10,000. Done. Then, estimating from the positive dividend, continue by placing 40 bu; multiply it with the augmenting corner 10,000, obtaining 40,000; add it into the augmenting lower side, making 3,240,000. Again multiply with the quotient, obtaining 12,960,000; add it into the augmenting upper side, making 320,640,000. Again multiply with the quotient, obtaining 1,282,560,000; add it into the following square, making 9,551,360,000. Then use the continued quotient 40 above to divide the dividend; it is exactly exhausted. The result, 840 bu, is the area of the field. Now the diagrams for finding the rates and extracting the root are set out below.\end{uniten}

\begin{unitzh}\unitlabel{shushu-020.22}按：此術以立天元一法明之。法：立天元一為尖積，即大小兩三角積和；自之，得一平方，為和自乘。以半廣冪減大斜冪，與餘積相乘，得二十九萬一千六百歩，為大三角積自乘；以坐廣冪減小斜冪，與餘數相乘，得九萬歩，為小三角積自乘。二自乘數併而倍之，内減去和自乘，得七十六萬三千二百歩少一平方，為較自乘；與和自乘再相乘，得七十六萬三千二百平方少一三乘方，寄左。次以大小兩三角積相減，餘二十萬零一千六百歩，為和較相乘數；自之，得四百零六億四千二百五十六萬歩，與左相等。則後歩數為實，前平方數為從上亷，三乘方數即益隅，與草中所取之數悉合。\end{unitzh}
\begin{uniten}Note: This method is clarified by the method of setting up the celestial unknown. The method: set up the celestial unknown as the pointed area, which is the sum (he) of the areas of the two triangles, large and small; multiplying it by itself gives 1 square power, which is the sum multiplied by itself. Subtract the area of the half breadth from the area of the large slant and multiply with the remaining area, obtaining 291,600 bu, which is the area of the large triangle multiplied by itself. Subtract the area of the "seated" breadth [Translator's note: 坐 "seated" is evidently written for 半 "half".] from the area of the small slant and multiply with the remaining number, obtaining 90,000 bu, which is the area of the small triangle multiplied by itself. Add the two self-multiplied numbers and double them; from this subtract the sum multiplied by itself, obtaining 763,200 bu less 1 square power, which is the difference multiplied by itself; multiply this again with the sum multiplied by itself, obtaining 763,200 square powers less 1 thrice-multiplied power; set it aside on the left. Next subtract the large and small triangle areas from each other; the remainder, 201,600 bu, is the number of the sum and the difference multiplied together; multiplying it by itself gives 40,642,560,000 bu, which is equal to the left. Then the latter number of bu is the dividend, the former number of square powers is the following upper side, and the number of thrice-multiplied powers is the augmenting corner; all agree with the numbers taken in the working.\end{uniten}

\begin{unitzh}\unitlabel{shushu-020.23}又按：此苦以小率九萬歩開平方，得三百歩，即小三角積；以大率二十九萬一千六百歩開平方，得五百四十歩，即大三角積；併之，得八百四十歩，即尖積。其法甚易，然必如此費算者，殆欲用立天元一法，不求分積，即得所問之總積也。\end{unitzh}
\begin{uniten}Further note: here [Translator's note: the text has 此苦 "this bitter"; 苦 is evidently corrupt.], extracting the square root of the small rate, 90,000 bu, gives 300 bu, which is the area of the small triangle; extracting the square root of the large rate, 291,600 bu, gives 540 bu, which is the area of the large triangle; adding them gives 840 bu, which is the pointed area. The method is very easy; yet that one must spend calculation like this is probably because one wishes to use the method of setting up the celestial unknown and, without seeking the partial areas, obtain directly the total area asked for.\end{uniten}

\begin{unitzh}\unitlabel{shushu-020.24}正負開三乘方圖\end{unitzh}
\begin{uniten}Diagram of Extracting the Thrice-Multiplied Root with Positive and Negative (zhengfu)\end{uniten}

\begin{unitzh}\unitlabel{shushu-020.25}術曰：商常為正，實常為負，從常為正，益常\end{unitzh}
\begin{uniten}Method: The quotient is always positive; the dividend is always negative; the following (cong) terms are always positive; the augmenting (yi) terms are always [Translator's note: the sentence breaks off here in the transcription.]\end{uniten}

\begin{unitzh}\unitlabel{shushu-020.26}《子部，天文算法類，算書之屬，數學九章，卷三上》《子部，天文算法類，算書之屬，數學九章，卷三上》《子部，天文算法類，算書之屬，數學九章，卷三上》《子部，天文算法類，算書之屬，數學九章，卷三上》\end{unitzh}
\begin{uniten}Masters section (zibu), Astronomy and mathematics category (tianwen suanfa lei), Mathematical books subcategory (suanshu zhi shu), Shuxue jiuzhang, juan 3, part 1. Masters section, Astronomy and mathematics category, Mathematical books subcategory, Shuxue jiuzhang, juan 3, part 1. Masters section, Astronomy and mathematics category, Mathematical books subcategory, Shuxue jiuzhang, juan 3, part 1. Masters section, Astronomy and mathematics category, Mathematical books subcategory, Shuxue jiuzhang, juan 3, part 1. [Translator's note: these four catalogue labels apparently mark page images in the transcription.]\end{uniten}

\begin{unitzh}\unitlabel{shushu-020.27}已上係開三乘方翻法圖，後篇效此。\end{unitzh}
\begin{uniten}The above are the diagrams of the turning-over method for extracting the thrice-multiplied root; the later sections follow this.\end{uniten}

\begin{unitzh}\unitlabel{shushu-020.28}三斜求積\end{unitzh}
\begin{uniten}Finding the Area from Three Slants\end{uniten}

\begin{unitzh}\unitlabel{shushu-021.1}問沙田一叚有三斜，其小斜一十三里，中斜一十四里，大斜一十五里，里法三百歩，欲知為田幾何？荅曰：田積三百一十五頃。\end{unitzh}
\begin{uniten}Question: A plot of sandy field has three obliques (xie): its small oblique is 13 li (里), its middle oblique 14 li, and its large oblique 15 li; the li standard (lifa) is 300 bu (歩). One wishes to know how much field it makes. Answer: The area of the field is 315 qing (頃).\end{uniten}

\begin{unitzh}\unitlabel{shushu-021.2}術曰：以少廣求之。以小斜冪併大斜冪，減中斜冪，餘半之，自乘于上；以小斜冪乘大斜冪，減上，餘四約之為實，一為從隅，開平方得積。\end{unitzh}
\begin{uniten}Method: Seek it by Lesser breadth (shaoguang). Combine the small-oblique area (mi) with the large-oblique area and subtract the middle-oblique area; halve the remainder, multiply it by itself, and place it above. With the small-oblique area multiply the large-oblique area and subtract the number above; reduce (yue) the remainder by 4 to make the dividend (shi); 1 is the joining corner (cong yu); extract the square root to obtain the area.\end{uniten}

\begin{unitzh}\unitlabel{shushu-021.3}草曰：以斜一十三里自乘，得一百六十九里，為小斜冪；以大斜一十五里自乘，得二百二十五里，為大斜冪；併小斜冪，得三百九十四里於上。以中斜一十四里自乘，得一百九十六里，為中斜冪；減上，餘一百九十八里；以半之，得九十九里；自乘，得九千八百一里於上。以小斜冪一百六十九乘大斜冪二百二十五，得三萬八千二十五；減上，餘二萬八千二百二十四；以四約之，得七千五十六里，為實。以為一隅，開平方。以隅超歩為一百，乃於實上商置八十，以商生隅，得八百，為從方；乃命上商除實，餘六百五十六。又以商生隅入方，得數退一位，為一百六十；隅退二位，為一。乃於實上續商四里，生隅入從方内，得一百六十四；乃命續商除實，適盡。所得八十四里為田積，其形長八十四，廣一里。以里法三百歩自乘，得九萬歩，乘八十四里，得七百五十六萬歩；以畝法二百四十除之，得三萬一千五百畆；又以頃法一百畆約之，得三百一十五頃。\end{unitzh}
\begin{uniten}Working: Multiply the oblique of 13 li by itself, obtaining 169 li, as the small-oblique area. Multiply the large oblique of 15 li by itself, obtaining 225 li, as the large-oblique area; combine it with the small-oblique area, obtaining 394 li, placed above. Multiply the middle oblique of 14 li by itself, obtaining 196 li, as the middle-oblique area; subtract it from the number above, remainder 198 li; halve it, obtaining 99 li; multiply it by itself, obtaining 9801 li, placed above. With the small-oblique area 169 multiply the large-oblique area 225, obtaining 38,025; subtract the number above, remainder 28,224; reduce it by 4, obtaining 7056 li, as the dividend. Take 1 as the corner and extract the square root. Move the corner by leaping steps (chaobu) to make it 100; then above the dividend place 80 as the quotient (shang); with the quotient generate from the corner (sheng yu), obtaining 800, as the joining square (congfang). Then command (ming) the upper quotient to divide the dividend: remainder 656. Again with the quotient generate from the corner and enter it into the square (fang); the number obtained, moved back one place, makes 160; the corner, moved back two places, makes 1. Then above the dividend continue the quotient with 4 li; generate from the corner and enter it into the joining square, obtaining 164. Then command the continued quotient to divide the dividend: it is exactly exhausted. The 84 li obtained is the area of the field; its shape is 84 long and 1 li wide. Multiply the li standard of 300 bu by itself, obtaining 90,000 bu; multiply the 84 li by it, obtaining 7,560,000 bu; divide this by the mu standard (mufa) of 240, obtaining 31,500 mu; further reduce it by the qing standard (qingfa) of 100 mu, obtaining 315 qing.\end{uniten}

\begin{unitzh}\unitlabel{shushu-021.4}按：此術以立天元一法明之。法：立天元一為三角積，倍之，得二元；自之，得四平方，為中長冪乘底冪，以大斜為底，寄之。又以小斜冪與大斜冪相加，内減中斜冪，得一百九十八里；半之，得九十九里，為小分底與底相乘長冪；自之，得九千八百零一里，為小分底冪乘底冪之數。又以小斜冪、大斜冪相乘，得三萬八千零二十五里，為小分底冪乘底冪、中長冪乘底冪各一；内減小分底冪乘底冪之數，餘二萬八千二百二十四里，為中長冪乘底冪之數，與寄數等。兩邊各以四約之，得七千零五十六里，與一平方等。里數為實，方數即從隅也。從二題同此。\end{unitzh}
\begin{uniten}Note: This method can be made clear with the method of setting up the celestial unknown (tianyuan). The method: set up the celestial unknown as the area of the triangle; double it, obtaining 2 yuan; multiply it by itself, obtaining 4 squares (pingfang), which is the middle-length area (zhongchang mi) multiplied by the base area (di mi), the large oblique being taken as the base; lodge it (ji). Again add the small-oblique area and the large-oblique area together, and from within subtract the middle-oblique area, obtaining 198 li; halve it, obtaining 99 li, which is the oblong area (chang mi) of the small part of the base (xiao fen di) multiplied with the base; multiply it by itself, obtaining 9801 li, which is the number of the area of the small part of the base multiplied by the base area. Again multiply the small-oblique area and the large-oblique area together, obtaining 38,025 li, which is one each of the area of the small part of the base multiplied by the base area and the middle-length area multiplied by the base area; from within subtract the number of the area of the small part of the base multiplied by the base area; the remainder, 28,224 li, is the number of the middle-length area multiplied by the base area, and is equal to the lodged number. Reduce both sides each by 4, obtaining 7056 li, equal to 1 square. The number of li is the dividend; the number of squares is precisely the joining corner. The following two problems are the same as this.\end{uniten}

\begin{unitzh}\unitlabel{shushu-021.5}斜蕩求積\end{unitzh}
\begin{uniten}Seeking the area of an oblique marsh\end{uniten}

\begin{unitzh}\unitlabel{shushu-021.6}問有蕩一所，正北濶一十七里，自南尖穿徑中長二十四里，東南斜二十里，東北斜一十五里，西斜二十六里，欲知畆積幾何？荅曰：蕩積一千九百一十一頃六十畆。\end{unitzh}
\begin{uniten}Question: There is a marsh (dang): due north it is 17 li wide (kuo); from the southern tip, cutting straight through, the middle length is 24 li; the south-eastern oblique is 20 li; the north-eastern oblique is 15 li; the western oblique is 26 li. One wishes to know how much its area in mu is. Answer: The area of the marsh is 1911 qing 60 mu.\end{uniten}

\begin{unitzh}\unitlabel{shushu-021.7}術曰：以少廣求之。置中長乘北濶，半之，為寄。以中長冪減西斜冪，餘為實，以一為隅，開平方，得數減北濶，餘自乘，併中長冪，共為内率。以小斜冪併率，减中斜冪，餘半之，自乘於上；以小斜冪乘率，減上，餘四約之為實，以一為隅，開平方，得數加寄，共為蕩積。\end{unitzh}
\begin{uniten}Method: Seek it by Lesser breadth. Set up the middle length, multiply it by the northern width, halve it, and take it as the lodged number. Subtract the middle-length area from the western-oblique area; the remainder is the dividend; take 1 as the corner and extract the square root; subtract the number obtained from the northern width; multiply the remainder by itself and combine it with the middle-length area; together they make the inner rate (nei lü). Combine the small-oblique area with the rate and subtract the middle-oblique area; halve the remainder, multiply it by itself, and place it above. With the small-oblique area multiply the rate and subtract the number above; reduce the remainder by 4 to make the dividend; take 1 as the corner and extract the square root; add the number obtained to the lodged number; together they make the area of the marsh.\end{uniten}

\begin{unitzh}\unitlabel{shushu-021.8}草曰：以中長二十四里乘北濶一十七里，得四百八，乃半之，得二百四里，為寄。以中長自乘，得五百七十六，為長冪；以西斜二十六里自乘，得六百七十六，為大斜冪；以減長冪，餘一百里，為實，開平方，得一十里；以減北濶數一十七里，餘七里，自乘，得四十九里，併長冪五百七十六，得六百二十五，為内率。次置東小斜一十五里自乘，得二百二十五，為小斜冪；又置東南中斜二十里自乘，得四百，為中冪；却以小斜冪併率，得八百五十，以減中冪四百，餘四百五十，乃半之，得二百二十五，自乘，得五萬六百二十五里于上。又以小斜冪二百二十五乘率六百二十五，得一十四萬六百二十五，減上，餘九萬里，以四約，得二萬二千五百，為實，開平方，得一百五十，併寄二百四里，得三百五十四里，為泛。以里法三百六十自乘，得一十二萬九千六百歩，乘泛，得四千五百八十七萬八千四百歩，以畆法二百四十歩約之，得一千九百一十一頃六十畆，為蕩積。\end{unitzh}
\begin{uniten}Working: With the middle length of 24 li multiply the northern width of 17 li, obtaining 408; then halve it, obtaining 204 li, as the lodged number. Multiply the middle length by itself, obtaining 576, as the length area; multiply the western oblique of 26 li by itself, obtaining 676, as the large-oblique area; subtract the length area from it, remainder 100 li, as the dividend; extract the square root, obtaining 10 li; subtract it from the northern-width number of 17 li, remainder 7 li; multiply it by itself, obtaining 49 li; combine it with the length area 576, obtaining 625, as the inner rate. Next set up the eastern small oblique of 15 li and multiply it by itself, obtaining 225, as the small-oblique area; again set up the south-eastern middle oblique of 20 li and multiply it by itself, obtaining 400, as the middle area; then combine the small-oblique area with the rate, obtaining 850; subtract the middle area 400 from it, remainder 450; then halve it, obtaining 225; multiply it by itself, obtaining 50,625 li, placed above. Again with the small-oblique area 225 multiply the rate 625, obtaining 140,625; subtract the number above, remainder 90,000 li; reduce by 4, obtaining 22,500, as the dividend; extract the square root, obtaining 150; combine it with the lodged 204 li, obtaining 354 li, as the provisional number (fan). Multiply the li standard of 360 by itself, obtaining 129,600 bu; multiply the provisional number by it, obtaining 45,878,400 bu; reduce it by the mu standard of 240 bu, obtaining 1911 qing 60 mu, as the area of the marsh.\end{uniten}

\begin{unitzh}\unitlabel{shushu-021.9}計地容民\end{unitzh}
\begin{uniten}Calculating the land to accommodate people\end{uniten}

\begin{unitzh}\unitlabel{shushu-021.10}問沙洲一叚，形如棹力，廣一千九百二十歩，從三十六百歩，大斜二千五百歩，小斜一千八百二十歩，以安集民，每户給一十五畆，欲知地積、容民幾何？荅曰：池積一百四十九頃九十五畆；容民九百九十九戸；餘地一十畆。\end{unitzh}
\begin{uniten}Question: A plot of sandbank is shaped like an "oar-strength" (棹力) [Translator's note: 棹力 is evidently corrupt, probably for 棹刀 "oar-blade".]; its width (guang) is 1920 bu, its length (cong) 36 hundred bu, its large oblique 2500 bu, and its small oblique 1820 bu. It is to be used to settle and gather people, each household being given 15 mu. One wishes to know how much the area of the land is and how many people it accommodates. Answer: The pond area is 149 qing 95 mu [Translator's note: 池 "pond" is evidently an error for 地 "land".]; it accommodates 999 households of people; remaining land, 10 mu.\end{uniten}

\begin{unitzh}\unitlabel{shushu-021.11}術曰：以少廣求之。置廣乘長，半之，為寄。以廣冪併從冪為中冪{\zhsmall 按：實大斜冪。}；以小斜冪併中冪減大斜冪{\zhsmall 按：實中斜冪。}，餘半之，自乘于上；以小斜冪棄中冪，減上，餘以四約之為實，以一為隅，開平方，得數加寄，共為積。以每户給數除積，得容民戸數。\end{unitzh}
\begin{uniten}Method: Seek it by Lesser breadth. Set up the width, multiply it by the length (chang), halve it, and take it as the lodged number. Combine the width area with the length area as the middle area [Siku editors: Note: in fact the large-oblique area.]; combine the small-oblique area with the middle area and subtract the large-oblique area [Siku editors: Note: in fact the middle-oblique area.]; halve the remainder, multiply it by itself, and place it above. With the small-oblique area "discard" the middle area [Translator's note: 棄 "discard" evidently stands for 乘 "multiply".] and subtract the number above; reduce the remainder by 4 to make the dividend; take 1 as the corner and extract the square root; add the number obtained to the lodged number; together they make the area. Divide the area by the number given to each household, obtaining the number of households of people accommodated.\end{uniten}

\begin{unitzh}\unitlabel{shushu-021.12}草曰：置廣一千九百二十歩，乘從三千六百歩，得六百九十一萬二千歩，乃半之，得三百四十五萬六千歩，為寄。以廣自乘，得三百六十八萬六千四百歩，為廣冪；又以從自乘，得一千二百九十六萬步，為從冪；併廣冪，得一千六百六十四萬六千四百步，為中冪。次以小斜一千八百二十歩自乘，得三百三十一萬二千四百歩，為小斜冪；又以大斜二千五百歩自乘，得六百二十五萬歩，為大斜冪。\end{unitzh}
\begin{uniten}Working: Set up the width of 1920 bu and multiply it by the length of 3600 bu, obtaining 6,912,000 bu; then halve it, obtaining 3,456,000 bu, as the lodged number. Multiply the width by itself, obtaining 3,686,400 bu, as the width area; again multiply the length by itself, obtaining 12,960,000 bu, as the length area; combine it with the width area, obtaining 16,646,400 bu, as the middle area. Next multiply the small oblique of 1820 bu by itself, obtaining 3,312,400 bu, as the small-oblique area; again multiply the large oblique of 2500 bu by itself, obtaining 6,250,000 bu, as the large-oblique area.\end{uniten}

\begin{unitzh}\unitlabel{shushu-021.13}却以小⿰冪併中冪，得一千九百九十五萬八千八百歩；以大斜冪減之，餘一千三百七十萬八千八百歩；乃半之，得六百八十五萬四千四百歩；自乘，得四十六萬九千八百二十七億九千九百三十六萬歩於上。次以小斜冪乘中冪，得五十五萬一千三百九十五億三千五百三十六萬歩；減上，餘八萬一千五百六十七億三千六百萬，為實；以四約之，得二萬三百九十一億八千四百萬，為實；以一為隅，開平方，得一十四萬二千八百歩；併寄三百四十五萬六千歩，共得三百五十九萬八千八百歩。以畆法二百四十歩除之，得一萬四千九百九十五畆；次以頃法一百畆約之，為一百四十九頃九十五畆，為地積。又為實，以每户所給一十五畆為法，除實，得九百九十九户，不盡一十畆，不及一戸所給數，以為餘地一十畆。\end{unitzh}
\begin{uniten}Then combine the small-⿰ area [Translator's note: ⿰ is a garbled character, evidently for 斜 "oblique".] with the middle area, obtaining 19,958,800 bu; subtract the large-oblique area from it, remainder 13,708,800 bu; then halve it, obtaining 6,854,400 bu; multiply it by itself, obtaining 46,982,799,360,000 bu, placed above. Next with the small-oblique area multiply the middle area, obtaining 55,139,535,360,000 bu; subtract the number above, remainder 8,156,736,000,000, as the dividend; reduce it by 4, obtaining 2,039,184,000,000, as the dividend; take 1 as the corner and extract the square root, obtaining 142,800 bu; combine it with the lodged 3,456,000 bu, together obtaining 3,598,800 bu. Divide it by the mu standard of 240 bu, obtaining 14,995 mu; next reduce it by the qing standard of 100 mu, making 149 qing 95 mu, as the area of the land. Again make it the dividend; take the 15 mu given to each household as the divisor (fa) and divide the dividend, obtaining 999 households, with 10 mu not exhausted; it does not reach the amount given to one household, and is taken as remaining land, 10 mu.\end{uniten}

\begin{unitzh}\unitlabel{shushu-021.14}蕉田求積\end{unitzh}
\begin{uniten}Seeking the area of a banana field (jiaotian)\end{uniten}

\begin{unitzh}\unitlabel{shushu-022.1}問：蕉葉田一叚，中長五百七十六歩，中廣三十四歩，不知其周，求積畆合幾何？荅曰：田積四十五畆一角{\zhsmall 按：六十歩為一角，葢四分畆之一也。}十一歩六萬三千七十分歩之五千二百一十三。\end{unitzh}
\begin{uniten}Question: One plot of banana-leaf field (jiaoye tian): its middle length is 576 bu (歩) and its middle breadth is 34 bu; its perimeter is not known. Find the area: how many mu (畆) does it come to? Answer: The area of the field is 45 mu 1 jiao (角) [Siku editors: 60 bu make 1 jiao; that is, 1/4 of a mu.] 11 5213/63070 bu.\end{uniten}

\begin{unitzh}\unitlabel{shushu-022.2}術曰：以長併廣，再自乗，又十乘之，為實；半廣、半長各自乘，所得相減，餘為從方；一為從隅；開平方，半之，得積。\end{unitzh}
\begin{uniten}Method: Add the length to the breadth, multiply it by itself twice, and further multiply it by 10, to make the dividend (shi). Halve the breadth and halve the length, multiply each by itself, and subtract the results from each other; the remainder is the following square (congfang); 1 is the following corner (congyu). Extract the square root, halve it, and obtain the area.\end{uniten}

\begin{unitzh}\unitlabel{shushu-022.3}草曰：以長五百七十六歩併廣三十四歩，得六百一十，兩度自乘{\zhsmall 按：即自乘再乘。}，得二億二千六百九十八萬一千歩；進一位，即是以十乘之，得二十二億六千九百八十一萬歩，定得此數，以為實。置長五百七十六，以半之，得二百八十八，自乘，得八萬二千九百四十四於上；又置廣三十四歩，以半之，得一十七，自乘，得二百八十九，減上，餘八萬二千六百五十五，為從方；以一為從隅，開平方，得二萬一千七百四十二歩，不盡一萬四百二十六歩。以商生隅入方，又併隅算，共得一十二萬六千一百四十，為母，與不盡及開方田積數皆半之，田積定得一萬八百七十一步六萬三千七十分歩之五千二百一十三。以畆法二百四十約之，得四十五畆一角一十一歩六萬三千七十分歩之五千二百一十三。\end{unitzh}
\begin{uniten}Working: Add the length, 576 bu, to the breadth, 34 bu, obtaining 610; multiply it by itself two times [Siku editors: that is, multiply by itself and multiply again], obtaining 226,981,000 bu; advance it one place, which is to multiply it by 10, obtaining 2,269,810,000 bu; definitively obtaining this number, take it as the dividend (shi). Set down the length, 576, halve it, obtaining 288; multiply it by itself, obtaining 82,944, in the upper position. Again set down the breadth, 34 bu, halve it, obtaining 17; multiply it by itself, obtaining 289; subtract it from the upper, leaving 82,655 as the following square (congfang). Take 1 as the following corner (congyu) and extract the square root, obtaining 21,742 bu, with 10,426 bu not exhausted. Use the quotient (shang) to generate from the corner and enter it into the square, and further add the corner counter (yusuan); together this gives 126,140 as the denominator. This, together with the unexhausted part and the field-area number from the root extraction, is all halved: the field area is fixed at 10,871 5213/63070 bu. Reduce it by the mu divisor, 240, obtaining 45 mu 1 jiao 11 5213/63070 bu.\end{uniten}

\begin{unitzh}\unitlabel{shushu-022.4}按：此術以長與廣相加，自乘再乘，又以十乘之，為長方積；以半長自乘、半廣自乘相減，為長濶較；求得闊，折半為田積，非法也。此題中廣甚小，故得數較古法多七百餘，較密法少二千七百餘。若設長為七百零七，廣為二百九十三，亦以此法求之：長廣相加，自之，再之，又十乘之，得一百億，為實；半長、半廣各自之，相減，得十萬零三千五百，為長闊較；求得闊，折半，得三萬零四百二十六歩餘，為田積。依宻法求之，實十四萬四千九百餘歩，所差甚逺，其術之不合顯然矣。葢數必三乘而後可以平方求之，今再乘之後僅以十進之，宜其不可用也。\end{unitzh}
\begin{uniten}Note: This method adds the length and the breadth, multiplies by itself and multiplies again, and further multiplies by 10, to make the area of a rectangle (changfang ji); it subtracts the half-length multiplied by itself and the half-breadth multiplied by itself from each other to make the difference of length and width (jiao); it finds the width and halves it to make the field area. This is not a correct method. In this problem the middle breadth is very small, so the number obtained is more than that of the ancient method (gufa) by over 700, and less than that of the accurate method (mifa) by over 2,700. If one sets the length as 707 and the breadth as 293 and seeks it likewise by this method: the length and the breadth are added, multiplied by itself and again, and further multiplied by 10, giving 10,000,000,000 as the dividend (shi); the half-length and the half-breadth are each multiplied by itself and subtracted, giving 103,500 as the difference of length and width; finding the width and halving it gives over 30,426 bu as the field area. Sought according to the accurate method, it is in fact over 144,900 bu. The discrepancy is very great; that the method does not fit is evident. For the number must be multiplied three times; only then can it be sought by the square root. Here, after multiplying a second time, it is merely advanced by ten: it is fitting that it cannot be used.\end{uniten}

\begin{unitzh}\unitlabel{shushu-022.5}漂田堆積\end{unitzh}
\begin{uniten}Washed-away field: piled-up area (duiji)\end{uniten}

\begin{unitzh}\unitlabel{shushu-022.6}問：三斜田被水衝去一隅，而成四不等直田之狀。元中斜一十六步，如多長；水直五歩，如少濶；殘小斜一十三歩，如弦；殘大斜二十歩，如元中斜之弦；横量徑一十二歩，如殘田之廣，又如元中斜之勾，亦是水直之股。欲求元積、殘積、水積、元大斜、元中斜、二水斜各幾何？荅曰：元積一百三十八歩一十一分歩之八；水積一十二歩一十一分歩一八{\zhsmall 按：應一十三歩一十一分歩之七}；殘積一百二十六歩；元大斜二十九歩一十一分歩之一；元小斜一十八歩一十一分歩之一{\zhsmall 按：應一十一分歩之十}；水大斜九歩一十一分歩之一；水小斜五歩一十一分歩之一{\zhsmall 按：應一十一分歩之一}。\end{unitzh}
\begin{uniten}Question: A three-oblique field (sanxie tian) has had one corner washed away by water, and has taken the shape of a four-unequal rectangular field (zhitian). The original middle oblique (zhongxie) is 16 bu, like the greater length; the water straight (shuizhi) is 5 bu, like the lesser width; the remaining small oblique is 13 bu, like the hypotenuse (xian); the remaining large oblique is 20 bu, like the hypotenuse of the original middle oblique; the transversely measured diameter (jing) is 12 bu, like the breadth of the remaining field, and also like the base (gou) of the original middle oblique; it is also the altitude (gu) of the water straight. One wishes to find the original area, the remaining area, the water area, the original large oblique, the original middle oblique, and the two water obliques: how much is each? Answer: Original area: 138 8/11 bu. Water area: 12 bu, "11 parts of a bu 1 8" [Translator's note: 一十一分歩一八 is garbled; 一 apparently stands for 之, i.e. 8/11 bu.] [Siku editors: it should be 13 7/11 bu.] Remaining area: 126 bu. Original large oblique: 29 1/11 bu. Original small oblique: 18 1/11 bu [Siku editors: it should be 10/11 bu.] Water large oblique: 9 1/11 bu. Water small oblique: 5 1/11 bu [Siku editors: it should be 1/11 bu.]\end{uniten}

\begin{unitzh}\unitlabel{shushu-022.7}術曰：以少廣求之，連枝入之，又勾股入之。置水直減中斜，餘為法。以中斜乘大殘，為大斜實；以法除實，得元大殘；以殘大斜減之，餘為水大斜。以法乘徑，又自之，為小斜隅。以水直冪併徑冪，為弦冪，又乘徑冪，又乘中斜冪，為小斜實；與隅可約，約之；閉連枝平方，得元小斜；以殘小斜減之，餘為水小斜；以水直乘之，為水實；倍水小母為法，除之，得水積{\zhsmall 按：此處法踈。}。以水直併中斜乘徑，為實；以二為法，除之，得殘積。以殘積併水積，共為元積。分者通之，重有者重通之。\end{unitzh}
\begin{uniten}Method: Seek it by "small breadth" (shaoguang), bring in "connected branches" (lianzhi), and further bring in base and altitude (gougu). Set down the water straight and subtract it from the middle oblique; the remainder is the divisor (fa). With the middle oblique multiply the large remnant (dacan) to make the large-oblique dividend (shi); divide the dividend by the divisor to obtain the original large remnant; subtract the remaining large oblique from it; the remainder is the water large oblique. With the divisor multiply the diameter, and further multiply it by itself, to make the small-oblique corner (yu). Add the water-straight area (mi) to the diameter area to make the hypotenuse area; further multiply by the diameter area, and further multiply by the middle-oblique area, to make the small-oblique dividend. If it and the corner can be reduced, reduce them; close the connected-branches square root [Translator's note: 閉 "close" evidently stands for 開 "open", i.e. extract.] to obtain the original small oblique; subtract the remaining small oblique from it; the remainder is the water small oblique; multiply it by the water straight to make the water dividend; double the water-small denominator to make the divisor; divide by it to obtain the water area [Siku editors: the method here is imprecise.]. Add the water straight to the middle oblique and multiply by the diameter to make the dividend; take 2 as the divisor; divide by it to obtain the remaining area. Add the remaining area to the water area; together they make the original area. Where there are fractions, convert them (tong); where they occur repeatedly, convert them repeatedly.\end{uniten}

\begin{unitzh}\unitlabel{shushu-022.8}草曰：以水直五減中斜一十六，餘一十一，為法。以中斜一十六乘大殘二十，得三百二十，為大斜實；以法除之，得二十九歩一十一分歩之一，為元大斜；内減殘大斜二十歩，餘九歩一十一分歩之一，為水大斜。以法一十一乘徑一十二{\zhsmall 按：乘徑可省。}，得一百三十二，自之，得一萬七十四百二十四，為小斜隅。以水直五自乘，得二十五，為水直冪；以徑一十二自之，得一百四十四，為徑冪；併水直冪，得一百六十九，為弦冪；以乘徑冪{\zhsmall 按：此乘徑冪亦可省，葢以此乘，復以此除，徒為多筭耳。}一百四十四，得二萬四千三百三十六於上；又以中斜一十六自乘，得二百五十六，為中斜冪；以乘上，得六百二十三萬一十六，為小斜實。開平方，與隅求等，得一百四十四，俱約之：實得四萬三千二百六十四，隅得一百二十一。\end{unitzh}
\begin{uniten}Working: Subtract the water straight, 5, from the middle oblique, 16, leaving 11 as the divisor (fa). With the middle oblique, 16, multiply the large remnant, 20, obtaining 320 as the large-oblique dividend (shi); divide it by the divisor, obtaining 29 1/11 bu as the original large oblique; from this subtract the remaining large oblique, 20 bu, leaving 9 1/11 bu as the water large oblique. With the divisor, 11, multiply the diameter, 12 [Siku editors: the multiplication by the diameter can be omitted.], obtaining 132; multiply it by itself, obtaining "1 wan 70 400 24" [Translator's note: 一萬七十四百二十四 is garbled; 十 "ten" apparently stands for 千 "thousand", i.e. 17,424.] as the small-oblique corner (yu). Multiply the water straight, 5, by itself, obtaining 25 as the water-straight area (mi); multiply the diameter, 12, by itself, obtaining 144 as the diameter area; add the water-straight area, obtaining 169 as the hypotenuse area; with it multiply the diameter area [Siku editors: this multiplication by the diameter area can also be omitted, for multiplying by it and then dividing by it again merely makes for more calculation.] 144, obtaining 24,336 in the upper position. Further multiply the middle oblique, 16, by itself, obtaining 256 as the middle-oblique area; with it multiply the upper, obtaining 6,230,016 as the small-oblique dividend. For extracting the square root, seek with the corner the equal number (deng), obtaining 144; reduce both by it: the dividend gives 43,264, and the corner gives 121.\end{uniten}

\begin{unitzh}\unitlabel{shushu-022.9}開方不盡，以連枝術入之：用隅一百二十一乘實四萬三千二百六十四，得五百二十三萬四千九百四十四，為定實；以一為定隅，開平方，得二千二百八十八，為實；以約隅一百二十一除之，得一十八歩，不盡一百一十一{\zhsmall 按：一百一十整。}，與法一百二十一俱以一十一約之，得一十一分歩之十，為元小斜。減殘小斜一十三歩，餘五歩一十一分歩之一{\zhsmall 按：十訛一，故下數誤。}，為水小斜；通歩内子，得五十六，以水直五歩乘之，得二百八十，為水實；倍水小母一十一，得二十二，為法，除之，得一十二歩，不盡一十六，與法俱以二約之，為一十二歩一十一分歩之八{\zhsmall 按：應一十二歩一十一分步之七。}水積。置中斜一十六，併水直五，得二十一，乗徑一十二，得二百五十二，以半之，得一百二十六，為殘積。以水併積，共得一百三十八歩一十一分歩之八，為元積{\zhsmall 按：應一百三十九步一十一分步之七。}。\end{unitzh}
\begin{uniten}The root extraction is not exhausted; bring in the connected-branches method (lianzhi shu): with the corner, 121, multiply the dividend, 43,264, obtaining 5,234,944 as the fixed dividend (dingshi); take 1 as the fixed corner (dingyu) and extract the square root, obtaining 2,288 as the dividend; divide it by the reduced corner, 121, obtaining 18 bu, with 111 not exhausted [Siku editors: exactly 110.]; this and the divisor, 121, are both reduced by 11, obtaining 10/11 bu, as the original small oblique. Subtract the remaining small oblique, 13 bu, leaving 5 1/11 bu [Siku editors: 十 "ten" has been corrupted into 一 "one"; therefore the numbers below are wrong.] as the water small oblique. Convert the bu and bring in the numerator (zi), obtaining 56; multiply it by the water straight, 5 bu, obtaining 280 as the water dividend; double the water-small denominator, 11, obtaining 22 as the divisor; divide by it, obtaining 12 bu, with 16 not exhausted; this and the divisor are both reduced by 2, making 12 8/11 bu [Siku editors: it should be 12 7/11 bu.] as the water area. Set down the middle oblique, 16, add the water straight, 5, obtaining 21; multiply by the diameter, 12, obtaining 252; halve it, obtaining 126 as the remaining area. Add the water to the area; together this gives 138 8/11 bu as the original area [Siku editors: it should be 139 7/11 bu.].\end{uniten}

\begin{unitzh}\unitlabel{shushu-022.10}《數學九章》卷三上\end{unitzh}
\begin{uniten}Shuxue jiuzhang (Mathematical Treatise in Nine Sections), juan 3, first part.\end{uniten}

\begin{unitzh}\unitlabel{shushu-022.11}子部，天文算法類，算書之屬，《數學九章》\end{unitzh}
\begin{uniten}Masters branch (zibu), category of astronomy and mathematics, subcategory of mathematical books: Shuxue jiuzhang.\end{uniten}

\section*{卷3下}\addcontentsline{toc}{section}{卷3下}

\begin{unitzh}\unitlabel{shushu-023.1}欽定四庫全書\end{unitzh}
\begin{uniten}Imperially Commissioned Complete Library of the Four Treasuries (Qinding Siku quanshu)\end{uniten}

\begin{unitzh}\unitlabel{shushu-023.2}數學九章卷三下 宋 秦九韶 撰\end{unitzh}
\begin{uniten}Shuxue jiuzhang (Mathematics in Nine Sections), juan 3, second part (xia). Composed by Qin Jiushao of the Song.\end{uniten}

\begin{unitzh}\unitlabel{shushu-023.3}田域\end{unitzh}
\begin{uniten}Fields and Territories (tianyu)\end{uniten}

\begin{unitzh}\unitlabel{shushu-023.4}環田三積\end{unitzh}
\begin{uniten}Three areas with a ring field (huantian)\end{uniten}

\begin{unitzh}\unitlabel{shushu-023.5}問：環田、大小圓田共三段。環田外周三十步，虗徑八步；大圓田徑一十步；小圓田周三十步。欲知三田積，及環田周、通、實徑，大圓周，小圓徑，各㡬何？答曰：環田積二十步二百三十六萬二千二百五十六分步之一百二十九萬八千二十五，通徑九步一十九分步之九，實徑一步一十九分步之九，内周二十五步一十七分步之五。大圓田積七十九步五十三分步之三，周三十一步二十一分步之十三。小圓田積七十一步二百八十六分步之四十三，徑九步一十九分步之九。\end{unitzh}
\begin{uniten}Question: A ring field and a large and a small circular field, three plots (duan) in all. The ring field's outer circumference is 30 bu (步) and its empty diameter (xujing) 8 bu; the large circular field's diameter is 10 bu; the small circular field's circumference is 30 bu. One wishes to know the areas of the three fields, and the ring field's circumference, through and solid diameters (tongjing, shijing), the large circle's circumference and the small circle's diameter: how much is each?
Answer: The ring field's area is 20 1298025/2362256 bu; the through diameter 9 9/19 bu; the solid diameter 1 9/19 bu; the inner circumference 25 5/17 bu.
The large circular field's area is 79 3/53 bu; the circumference 31 13/21 bu.
The small circular field's area is 71 43/286 bu; the diameter 9 9/19 bu.\end{uniten}

\begin{unitzh}\unitlabel{shushu-023.6}術曰：以方田及少廣率變求之。各置圓、環徑，自乗為冪，進位為實，以一為隅，開平方，得周。各置環、圓周，自乗為冪，退位為實，以一為隅，開平方，得徑。以周冪或徑冪乗各實，以一十六約之，為實，以一為隅，開平方，得圓積。置環周冪乘徑實，十六約之，為大率。置虚徑冪乗内周實，十六約之，為小率。以二率相減之餘，以自乗為實；併二率，倍之，為從上亷；一為益隅；開三乗方，得環積。置環周，自乗，退位為實，一為隅，開平方，得通徑。以虚徑減通徑，餘為實徑。其有開不盡者，約而命之。\end{unitzh}
\begin{uniten}Method: Find them by transforming the rates (lü) of Rectangular Fields (fangtian) and Diminishing Breadth (shaoguang). In each case set up the diameter of the circle or the ring, multiply it by itself to make the area (mi), advance it one place to make the dividend (shi), take 1 as the corner (yu) and extract the square root: this gives the circumference. In each case set up the circumference of the ring or the circle, multiply it by itself to make the area, retreat it one place to make the dividend, take 1 as the corner and extract the square root: this gives the diameter. Multiply each dividend by the circumference area or the diameter area, reduce (yue) it by 16 to make the dividend, take 1 as the corner and extract the square root: this gives the area of the circle. Set up the ring circumference area and multiply it by the diameter dividend; reduce it by 16: this makes the greater rate. Set up the empty-diameter area and multiply it by the inner-circumference dividend; reduce it by 16: this makes the lesser rate. Take the remainder of the two rates subtracted from each other and multiply it by itself to make the dividend; combine the two rates and double them to make the following upper side (cong shanglian); 1 is the augmenting corner (yi yu); extract the root of the thrice-multiplied power (sancheng fang): this gives the ring area. Set up the ring circumference, multiply it by itself and retreat it one place to make the dividend; 1 is the corner; extract the square root: this gives the through diameter. Subtract the empty diameter from the through diameter; the remainder is the solid diameter. Where the extraction is not exhausted, reduce it and name (ming) it.\end{uniten}

\begin{unitzh}\unitlabel{shushu-023.7}草曰：置大圓徑一十步，自乗得一百，為徑冪；進位得一千，為實；以一為隅，開平方，得三十一步，不盡三十九，為分子。乃以隅生方，又益隅，共得六十三，為分母。以分子與母求等，得三，俱以三約之，母子得二十一分步之一十三，為大圓周三十一步二十一分步之一十三。次以徑冪一百乗前實一千，得一十萬，以十六約之，得六千二百五十，為實；以一為隅，開平方，得七十九步，不盡九，為分子。乃以隅生方，又增隅，得一百五十九，為分母。以分子母求等，得三，俱以三約母子，得五十三分步之三，為大圓積七十九步五十三分步之三。\end{unitzh}
\begin{uniten}Working: Set up the large circle's diameter, 10 bu; multiply it by itself, getting 100 as the diameter area; advance it one place, getting 1000 as the dividend; take 1 as the corner and extract the square root, getting 31 bu, with 39 not exhausted, as the numerator. Then from the corner generate the square (fang), and further add the corner, getting 63 in all as the denominator. Seek the equal number (deng) of the numerator and the denominator, getting 3; reduce both by 3: the denominator and numerator give 13/21 bu. This makes the large circle's circumference, 31 13/21 bu. Next multiply the previous dividend, 1000, by the diameter area, 100, getting 100000; reduce it by 16, getting 6250 as the dividend; take 1 as the corner and extract the square root, getting 79 bu, with 9 not exhausted, as the numerator. Then from the corner generate the square, and further increase it by the corner, getting 159 as the denominator. Seek the equal number of the numerator and denominator, getting 3; reduce the denominator and numerator both by 3, getting 3/53 bu. This makes the large circle's area, 79 3/53 bu.\end{uniten}

\begin{unitzh}\unitlabel{shushu-023.8}次置小圓田周三十步，自乗得九百，為周冪；退位得九十，為徑實；以一為隅，開平方，得九步，不盡九。以隅生方，又益隅，得一十九步之九，為小圓徑九步一十九分步之九。次以周冪九百乗前實九十，得八萬一千，以十六約之，得五千六十二步五分，為實；以一為隅，開平方，得七十一步，有不盡數二十一步五分，為子。以隅生方，又益隅，得一百四十三，為分母。以分子母求等，得五分，俱約之，得二百八十六分步之四十三，為積。\end{unitzh}
\begin{uniten}Next set up the small circular field's circumference, 30 bu; multiply it by itself, getting 900 as the circumference area; retreat it one place, getting 90 as the diameter dividend; take 1 as the corner and extract the square root, getting 9 bu, with 9 not exhausted. From the corner generate the square, and further add the corner, getting 9 of 19 bu. This makes the small circle's diameter, 9 9/19 bu. Next multiply the previous dividend, 90, by the circumference area, 900, getting 81000; reduce it by 16, getting 5062 bu 5 fen (分) [Translator's note: fen here, and li (釐) below, are decimal places, tenths and hundredths.] as the dividend; take 1 as the corner and extract the square root, getting 71 bu, with an unexhausted number of 21 bu 5 fen as the numerator. From the corner generate the square, and further add the corner, getting 143 as the denominator. Seek the equal number of the numerator and denominator, getting 5 fen; reduce both by it, getting 43/286 bu. This makes the area.\end{uniten}

\begin{unitzh}\unitlabel{shushu-023.9}次置環田周三十步，自乗得九百，為周冪；退位得九十，為實；以一為隅，開平方，得九步，不盡九，為分子。以隅生方，併隅，得一十九，為分母，直命之，為環田通徑九步一十九分步之九。次以環周冪九百乗環實九十，得八萬一千，以十六約之，得五千六十二步五分，為大率。\end{unitzh}
\begin{uniten}Next set up the ring field's circumference, 30 bu; multiply it by itself, getting 900 as the circumference area; retreat it one place, getting 90 as the dividend; take 1 as the corner and extract the square root, getting 9 bu, with 9 not exhausted, as the numerator. From the corner generate the square and combine the corner, getting 19 as the denominator; name it directly: this makes the ring field's through diameter, 9 9/19 bu. Next multiply the ring dividend, 90, by the ring circumference area, 900, getting 81000; reduce it by 16, getting 5062 bu 5 fen as the greater rate.\end{uniten}

\begin{unitzh}\unitlabel{shushu-023.10}次置環田虗徑八步，自乗得六十四，為虗冪；進位得六百四十，為實；以一為隅，開平方，得二十五步，不盡一十五，為分子。以隅生方，又併隅，得五十一，為分母，與子求等，得三，俱約之，得一十七分步之五，為環田内周二十五步一十七分步之五。次以虗冪六十四乗周實六百四十，得四萬九百六十，以十六約之，得二千五百六十，為小率。\end{unitzh}
\begin{uniten}Next set up the ring field's empty diameter, 8 bu; multiply it by itself, getting 64 as the empty area; advance it one place, getting 640 as the dividend; take 1 as the corner and extract the square root, getting 25 bu, with 15 not exhausted, as the numerator. From the corner generate the square, and further combine the corner, getting 51 as the denominator; seek the equal number with the numerator, getting 3; reduce both by it, getting 5/17 bu. This makes the ring field's inner circumference, 25 5/17 bu. Next multiply the circumference dividend, 640, by the empty area, 64, getting 40960; reduce it by 16, getting 2560 as the lesser rate.\end{uniten}

\begin{unitzh}\unitlabel{shushu-023.11}以小率減大率，餘二千五百二步五分，自乗得六百二十六萬二千五百六步二分五釐，為實。以大小二率併之，得七千六百二十二步五分，倍之，得一萬五千二百四十五，為從上亷；以一為益隅；開玲瓏三乗方，得二十步，不盡三十二萬四千五百六步二分五釐，為分子。續啇無數，乃以益隅一、益下亷八十併之，得八十一，為減母。次以從上亷一萬二千八百四十五併從方五十七萬七千八百，得五十九萬六百四十五，以母八十一減之，餘五十九萬五百六十四，為分母。以分子求等，得二分五釐，俱約之，得二百三十六萬二千二百五十六分步之一百二十九萬八千二十五，為環田積二十步二百三十六萬二千二百五十六分步之一百二十九萬八千二十五。\end{unitzh}
\begin{uniten}Subtract the lesser rate from the greater rate; the remainder is 2502 bu 5 fen; multiply it by itself, getting 6262506 bu 2 fen 5 li as the dividend. Combine the two rates, greater and lesser, getting 7622 bu 5 fen; double it, getting 15245 as the following upper side; take 1 as the augmenting corner; extract the openwork (linglong) root of the thrice-multiplied power [Translator's note: here the positions of the square and of the lower side are empty.], getting 20 bu, with 324506 bu 2 fen 5 li not exhausted, as the numerator. Continuing the quotient (shang), there is no number; then take the augmenting corner, 1, and the augmenting lower side (xialian), 80, and combine them, getting 81 as the subtracting denominator (jianmu). Next combine the following upper side, 12845, with the following square (cong fang), 577800, getting 590645; subtract the denominator, 81, from it; the remainder, 590564, is the denominator. Seek the equal number with the numerator, getting 2 fen 5 li; reduce both by it, getting 1298025/2362256 bu. This makes the ring field's area, 20 1298025/2362256 bu.\end{uniten}

\begin{unitzh}\unitlabel{shushu-023.12}次置環田通徑九步一十九分步之九，以虚徑八步減之，餘一步一十九分步之九，為環田實徑。合問。\end{unitzh}
\begin{uniten}Next set up the ring field's through diameter, 9 9/19 bu; subtract the empty diameter, 8 bu, from it; the remainder, 1 9/19 bu, is the ring field's solid diameter. This fits the question.\end{uniten}

\begin{unitzh}\unitlabel{shushu-023.13}按：周徑相求，以進位、退位為實者，蓋以徑一周三有竒，徑一自之仍得一，周自之略與十等，故徑冪升一位為周冪，周冪降一位為徑冪，以省算，亦法之巧者。其徑求周，較宻率約大一百五十七分之一；周求徑，約小一百五十九分之一。然較古率則已宻矣。\end{unitzh}
\begin{uniten}Note: Where circumference and diameter are sought from each other by advancing or retreating one place to make the dividend, it is because, with the diameter 1, the circumference is 3 and an odd part; the diameter 1 multiplied by itself still gives 1, while the circumference multiplied by itself is roughly equal to 10. Therefore the diameter area raised one place makes the circumference area, and the circumference area lowered one place makes the diameter area, so as to save computation; this too is an ingenious point of method. When the circumference is sought from the diameter, compared with the precise rate (mi lü) it is larger by about 1/157; when the diameter is sought from the circumference, it is smaller by about 1/159. Yet compared with the ancient rate (gu lü) it is already precise.\end{uniten}

\begin{unitzh}\unitlabel{shushu-023.14}其周冪、徑冪相乗，十六約之，開平方得圓積者，蓋周徑相乗，四歸得圓積；徑自乗為方積，故四歸亦展為自乗十六之數約之，得四分徑之冪乗周冪之數，故開方得圓積。至求環積，與前求尖田積同，但彼立天元一為兩積之和，此立天元一為兩積之較耳。其式如左。\end{unitzh}
\begin{uniten}As for multiplying the circumference area and the diameter area together, reducing by 16 and extracting the square root to obtain the circle's area: it is because circumference and diameter multiplied together and divided by 4 (si gui) give the circle's area; the diameter multiplied by itself makes a square area, therefore the division by 4 is likewise expanded into the self-multiplied number 16, by which one reduces, obtaining the number of the area of a quarter of the diameter multiplied by the circumference area; therefore extracting the square root gives the circle's area. As for finding the ring area, it is the same as finding the area of the pointed field (jiantian) earlier, except that there the celestial unknown (tianyuan) was set up as the sum (he) of the two areas, and here the celestial unknown is set up as the difference (jiao) of the two areas. Its form is as on the left.\end{uniten}

\begin{unitzh}\unitlabel{shushu-023.15}法：立天元一為環田積，即内外兩圓積之較；自之，得一平方，為較自乗。以大小率{\zhsmall 即二圓積各自乗}併而倍之，得一萬五千二百四十五步，内減較自乗，得一萬五千二百四十五步少一平方，為和自乗。與較自乗再相乗，得一萬五千二百四十五平方少一三乗方，寄之。次以大小率相減，餘二千五百零二步五分，為和較相乗；再自之，得六百二十六萬二千五百零六步二分五釐，與寄數等，即為實。寄數内平方數即從上亷，三乗方數即益隅。\end{unitzh}
\begin{uniten}Rule: Set up the celestial unknown as the ring field's area, which is the difference of the two circle areas, inner and outer; multiply it by itself, getting 1 square power (pingfang), which is the difference multiplied by itself. Combine the greater and lesser rates [Note: that is, the two circle areas each multiplied by itself] and double them, getting 15245 bu; from this subtract the difference multiplied by itself, getting 15245 bu less 1 square power, which is the sum multiplied by itself. Multiply it together again with the difference multiplied by itself, getting 15245 square powers less 1 thrice-multiplied power; lodge it (ji). Next subtract the greater and lesser rates from each other; the remainder, 2502 bu 5 fen, is the sum and the difference multiplied together; multiply it again by itself, getting 6262506 bu 2 fen 5 li, which is equal to the lodged number: this is then the dividend. In the lodged number, the number of square powers is the following upper side, and the number of thrice-multiplied powers is the augmenting corner.\end{uniten}

\begin{unitzh}\unitlabel{shushu-023.16}此圗照問列位，以後照草運算。\end{unitzh}
\begin{uniten}This diagram sets out the positions according to the question; from here on, compute according to the working. [Translator's note: this and the following short paragraphs are captions to counting-board diagrams, which this transcription does not reproduce; "upper", "auxiliary", "middle", "next" and "lower" in them are positions on the board.]\end{uniten}

\begin{unitzh}\unitlabel{shushu-023.17}乃先置次大徑，以上副自乗得中，以中進位為次實，以一為下隅。\end{unitzh}
\begin{uniten}Then first set up, next, the large diameter; the upper (shang) and the auxiliary (fu) multiplied by themselves give the middle (zhong); advance the middle one place to make the next (ci) dividend; take 1 as the lower (xia) corner.\end{uniten}

\begin{unitzh}\unitlabel{shushu-023.18}凡九變，至此得大圓徑。次求大圓積。\end{unitzh}
\begin{uniten}Nine transformations (bian) in all; at this point one obtains the large circle's diameter. Next find the large circle's area.\end{uniten}

\begin{unitzh}\unitlabel{shushu-023.19}上副自乗得中，以次約之，得下為實。\end{unitzh}
\begin{uniten}The upper and the auxiliary multiplied by themselves give the middle; reduce it by the next, getting the lower as the dividend.\end{uniten}

\begin{unitzh}\unitlabel{shushu-023.20}凡十一變，至此得大圓積。次求小圓徑。\end{unitzh}
\begin{uniten}Eleven transformations in all; at this point one obtains the large circle's area. Next find the small circle's diameter.\end{uniten}

\begin{unitzh}\unitlabel{shushu-023.21}一十一變得小圓積。次求環田通徑。當求環田通徑，蓋環田之外周三十步與小圓田外周同，則不遇與前七變諸圗一理，兹不復繁。乃求實徑。\end{unitzh}
\begin{uniten}Eleven transformations give the small circle's area. Next find the ring field's through diameter. In finding the ring field's through diameter: since the ring field's outer circumference of 30 bu is the same as the small circular field's outer circumference, it then does not meet [Translator's note: 不遇, probably an error for 不過, "is nothing but"] — it is one principle with the preceding diagrams of seven transformations; here it is not repeated at length. Then find the solid diameter.\end{uniten}

\begin{unitzh}\unitlabel{shushu-023.22}凡九變得環田内周。次求環積。\end{unitzh}
\begin{uniten}Nine transformations in all give the ring field's inner circumference. Next find the ring area.\end{uniten}

\begin{unitzh}\unitlabel{shushu-023.23}<子部,天文算法類,算書之屬,數學九章,卷三下> <子部,天文算法類,算書之屬,數學九章,卷三下> <子部,天文算法類,算書之屬,數學九章,卷三下> <子部,天文算法類,算書之屬,數學九章,卷三下>\end{unitzh}
\begin{uniten}Masters Section (zibu), Category of Astronomy and Mathematics (tianwen suanfa lei), Subcategory of Mathematical Books (suanshu zhi shu), Shuxue jiuzhang, juan 3, second part.
Masters Section, Category of Astronomy and Mathematics, Subcategory of Mathematical Books, Shuxue jiuzhang, juan 3, second part.
Masters Section, Category of Astronomy and Mathematics, Subcategory of Mathematical Books, Shuxue jiuzhang, juan 3, second part.
Masters Section, Category of Astronomy and Mathematics, Subcategory of Mathematical Books, Shuxue jiuzhang, juan 3, second part.
[Translator's note: this Siku quanshu catalogue label, given four times in succession, is a marker inserted by the transcription, probably in place of pages of diagrams.]\end{uniten}

\begin{unitzh}\unitlabel{shushu-023.24}開三乗方，凡二十變，至此得環田積數。\end{unitzh}
\begin{uniten}Extracting the root of the thrice-multiplied power, twenty transformations in all; at this point one obtains the number of the ring field's area.\end{uniten}

\begin{unitzh}\unitlabel{shushu-023.25}求實徑，但以虚徑减通徑，餘一步一十九分步之九，為環田實徑。\end{unitzh}
\begin{uniten}To find the solid diameter, simply subtract the empty diameter from the through diameter; the remainder, 1 9/19 bu, is the ring field's solid diameter.\end{uniten}

\begin{unitzh}\unitlabel{shushu-023.26}均分梯田\end{unitzh}
\begin{uniten}Dividing a ladder field (titian) equally\end{uniten}

\begin{unitzh}\unitlabel{shushu-024.1}問：户業田一段，若梯之狀，南廣小，三十四步；北廣大，五十二步；正長一百五十步。合係兄弟三人均分其田，邉道各欲出入，其地形難分，經官乞分：南甲、乙，北内。欲知其田共積、各人合得田數及各段正長、大小廣㡬何？答曰：田共積二十六畆二百一十步。甲得八畆三角五十步：小廣三十四步，係元南廣；大廣四十步五萬八千七百九分步之五萬二千二百八十四，大約百分步之八十九；正長五十七步二千四十五分步之□百五十三，大約一百分步之四十一分。乙得八畆三角五十步：小廣同甲大廣；大廣四十六步八萬四千八百二十六億八千九百五十七萬二千六百五十一分步之六萬五千八百七十四億五千四百八十二萬五千二百八十三，計大率約百分步之七十七分半強；正長四十九步四億一千二百四十萬六千二百九分步之二千二十七萬六千三百一十九，大約百分步之四分九釐。丙得八畆三角五十步：小廣同乙大廣；大廣五十二步，係元北廣；正長四十三步八千四百三十三億七千九十萬一千九百五分步之四千四百八十八億八千六百二萬九千四十六，大約百分步之五十三分強。\end{unitzh}
\begin{uniten}Question: A household estate field, one plot, in the shape of a ladder (ti). The southern width is the smaller, 34 bu (步); the northern width is the larger, 52 bu; the straight length (zhengchang) is 150 bu. It falls to three brothers to divide the field equally, each wishing to go in and out by the side road. The shape of the land is hard to divide, so they went to the officials to request a division: south, A (jia), B (yi); north, nei [Translator's note: 内 "inside" is evidently an error for 丙 bing, C.]. One wishes to know the total area of the field, the amount of field each person should receive, and the straight length and the large and small widths of each section: how much is each?
Answer: The total area of the field is 26 mu (畆) 210 bu.
A receives 8 mu 3 jiao (角) 50 bu: small width 34 bu, which is the original southern width; large width 40 52,284/58,709 bu, roughly 89/100 bu; straight length 57 [lost digit]53/2045 bu, roughly 41/100 bu.
B receives 8 mu 3 jiao 50 bu: small width the same as A's large width; large width 46 6,587,454,825,283/8,482,689,572,651 bu, reckoned at the rough rate (dalü) about (77 1/2)/100 bu, a little more; straight length 49 20,276,319/412,406,209 bu, roughly 4 fen (分) 9 li (釐) of 100 parts of a bu.
C (bing) receives 8 mu 3 jiao 50 bu: small width the same as B's large width; large width 52 bu, which is the original northern width; straight length 43 448,886,029,046/843,370,901,905 bu, roughly 53/100 bu, a little more.\end{uniten}

\begin{unitzh}\unitlabel{shushu-024.2}術曰：以少廣及從法求之。併兩廣，乗長，得數，以分田人數約之，為通率；半之，為各積。以長乗各積，為共實；以長乗南廣，為甲從方；二廣差半之，為共隅。開連枝平方，得甲截長。以甲長除通率，得數減小廣，餘為甲廣，即為乙小廣。以元長乗乙小廣，為乙從方。置共隅、共實，開連枝平方，得乙截長。以乙長除通率，得數減乙小廣，餘為乙大廣，即為丙小廣。倂甲乙長，減元長，餘為丙長。以元大廣為丙大廣。各有分者通之。\end{unitzh}
\begin{uniten}Method: Seek it by Small Width (shaoguang) and the joined method (congfa). Add the two widths and multiply by the length; divide the obtained number (deshu) by the number of persons dividing the field, making the common rate (tonglü); halve it, making each one's area (geji). Multiply each one's area by the length, making the common dividend (gongshi); multiply the southern width by the length, making A's joined side (congfang); halve the difference of the two widths, making the common corner (gongyu). Extract the linked-branch square root (lianzhi pingfang), obtaining A's cut length (jiechang). Divide the common rate by A's length; from the obtained number subtract the small width; what remains is A's width, which is B's small width. Multiply B's small width by the original length, making B's joined side. Set up the common corner and the common dividend and extract the linked-branch square root, obtaining B's cut length. Divide the common rate by B's length; from the obtained number subtract B's small width; what remains is B's large width, which is C's small width. Add the lengths of A and B and subtract from the original length; what remains is C's length. Take the original large width as C's large width. Where there are parts in each, convert them (tong).\end{uniten}

\begin{unitzh}\unitlabel{shushu-024.3}草曰：置小廣三十四，併大廣五十二，得八十六，乗長一百五十，得一萬二千九百，為實。以兄弟三人約之，得四千三百，為通率；半之，得二十一百五十，為各積。以□法二百四十步約之，得八畆，不盡二百三十步；以角法六十步約之，得三角五十步。是三人各得八畆三角五十步。以元長一百五十步乗各積二千一百五十，得三十二萬二千五百，為共實。以長一百五乗小廣三十四，得五千一百，為甲從方。以小廣大廣，餘一十八，乃半之，得九，為共隅。開連枝平方{\zhsmall 開方革更不繁具}，得五十七步，不盡，三約為二千四十五分步之八百五十三，為甲截長。\end{unitzh}
\begin{uniten}Working: Set up the small width 34, add the large width 52, obtaining 86; multiply by the length 150, obtaining 12,900, as the dividend (shi). Divide it by the 3 brothers, obtaining 4300, as the common rate; halve it, obtaining 2150, as each one's area. Divide it by the [lost character] divisor (fa), 240 bu, obtaining 8 mu, remainder (bujin) 230 bu; divide this by the jiao divisor, 60 bu, obtaining 3 jiao 50 bu. Thus the 3 persons each receive 8 mu 3 jiao 50 bu. With the original length 150 bu multiply each one's area 2150, obtaining 322,500, as the common dividend. With the length 105 [Translator's note: 一百五, evidently for 150.] multiply the small width 34, obtaining 5100, as A's joined side. With the small width and the large width, there remains 18; then halve it, obtaining 9, as the common corner. Extract the linked-branch square root [Note: The root extraction ge (革) is not set out again in detail.] [Translator's note: 革 is evidently an error for 草 "working".], obtaining 57 bu, with a remainder; reduced by 3, it makes 853/2045 bu, as A's cut length.\end{uniten}

\begin{unitzh}\unitlabel{shushu-024.4}乃以分母二千四十五通全步，内子，共得一十一萬七千四百一十八，為法。又以分母乗通率四千三百，得八百七十九萬三千五百，為實。以法除之，得七十四步，不盡一十萬四千五百六十八，與法求等，得二，俱約之，為五萬八千七百九分步之五萬二千二百八十四。乃以小廣三十四步於所得全步七十四步内減之，餘四十步五萬八千七百九分步之五萬二千二百八十四，為甲大廣，即為乙小廣。\end{unitzh}
\begin{uniten}Then with the denominator (fenmu) 2045 convert the whole bu and add in the numerator (zi), obtaining in all 117,418, as the divisor. Again with the denominator multiply the common rate 4300, obtaining 8,793,500, as the dividend. Divide by the divisor, obtaining 74 bu, remainder 104,568; seeking the equal number (deng) with the divisor gives 2; reducing both by it makes 52,284/58,709 bu. Then from within the whole bu obtained, 74 bu, subtract the small width 34 bu; there remains 40 52,284/58,709 bu, as A's large width, which is B's small width.\end{uniten}

\begin{unitzh}\unitlabel{shushu-024.5}今次求乙長：乃以分母五萬八千七百九通乙小廣四十步，得二百三十四萬八千三百六十，内子五萬二千二百八十四，得二百四十萬六百四十四；又元長一百五十乗之，得三億六千九萬六千六百，為乙從方。又以分母五萬八千七百九通共實三十二萬二千五百，得一百八十九億三千三百六十五萬二千五百，為乙實。又以分母通共隅九，得五十二萬八千三百八十一，為乙從隅。開連枝平方{\zhsmall 更不立草}，得四十九步，不盡二千二十七萬六千三百十九；隅併方，共得四億一千二百四十萬六千三百九，為母。與不盡求等，單一，不可約，乃定為四十九步四億一千二百四十萬六千三百九分步之二千二十七萬六千三百一十九，為乙截長。\end{unitzh}
\begin{uniten}Now next, to seek B's length: with the denominator 58,709 convert B's small width 40 bu, obtaining 2,348,360; add in the numerator 52,284, obtaining 2,400,644; again multiply it by the original length 150, obtaining 360,096,600, as B's joined side. Again with the denominator 58,709 convert the common dividend 322,500, obtaining 18,933,652,500, as B's dividend. Again with the denominator convert the common corner 9, obtaining 528,381, as B's joined corner (congyu). Extract the linked-branch square root [Note: No further working is set up.], obtaining 49 bu, remainder 20,276,319; adding the corner to the side gives in all 412,406,309, as the denominator. Seeking the equal number with the remainder, it is unity (danyi) and cannot be reduced; so it is fixed as 49 20,276,319/412,406,309 bu, as B's cut length.\end{uniten}

\begin{unitzh}\unitlabel{shushu-024.6}以乙長母通全步，内子，得二百二億二千八百一十八萬五千四百六十，為法。以乙長七千七百三十三億四千七百一十二萬八千七百為實。以法除之，得八十七步，不盡一百三十四億九千四百九十九萬三千六百八十，與法求等，得一百四十，俱約之，為八十七步一億四千四百四十八萬七千三十九分步之九千六百三十九萬二千八百一十二，為得數。\end{unitzh}
\begin{uniten}With the denominator of B's length convert the whole bu and add in the numerator, obtaining 20,228,185,460, as the divisor. Take B's length 773,347,128,700 as the dividend. Divide by the divisor, obtaining 87 bu, remainder 13,494,993,680; seeking the equal number with the divisor gives 140; reducing both by it makes 87 96,392,812/144,487,039 bu, as the obtained number.\end{uniten}

\begin{unitzh}\unitlabel{shushu-024.7}乃以乙小廣母五萬八千七百九乗得數子九千六百三十九萬二千八百一十二，得五萬六千五百九十一億二千五百五十九萬九千七百八，為泛。却以得數母一億四千四百四十八萬七千三十九分乗乙小廣子五萬二千二百八十四，得七萬五千五百四十三億六千三十四萬七千七十六，為寄數於上。乃以小廣母五萬八千七百九乗得數母一億四千四百四十八萬七千三十九，得八萬四千八百二十六億八千九百五十七萬二千六百五十一。以寄減泛，今不及減，乃破全步一為分，併泛，得八十六步十四萬一千四百一十八億一千五百一十七萬二千三百五十九；減去小廣四十步及分，餘四十六步八萬四千八百二十六億八千九百五十七萬二千六百五十一分步之六萬五千八百七十四億五千四百八十二萬五千二百八十三，為乙大廣，亦丙小廣。\end{unitzh}
\begin{uniten}Then with the denominator of B's small width, 58,709, multiply the numerator of the obtained number, 96,392,812, obtaining 5,659,125,599,708, as the provisional (fan). In turn, with the denominator of the obtained number, 144,487,039 parts, multiply the numerator of B's small width, 52,284, obtaining 7,554,360,347,076, as the lodged number (jishu) above. Then with the denominator of the small width, 58,709, multiply the denominator of the obtained number, 144,487,039, obtaining 8,482,689,572,651. Subtract the lodged from the provisional: now it cannot be subtracted, so break 1 whole bu into parts and add them to the provisional, obtaining 86 bu 14,141,815,172,359; subtract the small width 40 bu and its parts; there remains 46 6,587,454,825,283/8,482,689,572,651 bu, as B's large width, which is also C's small width.\end{uniten}

\begin{unitzh}\unitlabel{shushu-024.8}求丙長：置甲長五十七步二千四十五分步之八百五十六，乙長四十九步四億一千二百四十萬六千三百九分步之二千二十七萬六千三百一十九。以甲乙分母互乗子，甲乙分母相乗，得甲正長五十七步八千四百三十三億七千九十萬一千九百五分步之三千五百三十億一千九百八十萬五百四億，乙正長四十九步八千四百三十三億七千九十萬一千九百五分步之四百一十四億六千五百七萬二千三百五十五。併甲乙長及分，共長一百六步三千九百四十四億八千四百八十七萬二千八百五十九分，用減元長一百五十步：先破一步，通分母，作八千四百三十三億七千九十萬一千九百五，減去甲乙長，長餘四十三步八千四百三十三億七千九十萬一千九百五分步之四千四百八十八億八千六百二萬九千四十六，為丙正長。\end{unitzh}
\begin{uniten}To seek C's length: set up A's length, 57 856/2045 bu, and B's length, 49 20,276,319/412,406,309 bu. With the denominators of A and B cross-multiply the numerators, and multiply the denominators of A and B together, obtaining A's straight length, 57 353,019,800,504/843,370,901,905 bu [Translator's note: the text adds a further 億 "hundred million" after this numerator.], and B's straight length, 49 41,465,072,355/843,370,901,905 bu. Add the lengths of A and B with their parts: the total length is 106 bu 394,484,872,859 parts; use it to subtract from the original length 150 bu: first break 1 bu and convert it with the denominator, making 843,370,901,905; subtract the lengths of A and B; the length remaining, 43 448,886,029,046/843,370,901,905 bu, is C's straight length.\end{uniten}

\begin{unitzh}\unitlabel{shushu-024.9}按：此術以立天元一法明之。法：立天元一為甲正長。南北廣差折半，得九，以乗天元，得九元；以共正長除之，得一百五十分天元之九，為甲之半廣差。與小廣相加，得三十四步多一百五十分元之九；再以天元乗之，得三十四元多一百五十分平方之九，即與每人分田二千一百五十步等。兩數各以分母一百五十乗之，得三十二萬二千五百步，與九平方多五千一百元等。步數為實，元數為從方，平方數為隅，得甲正長。求乙、丙長廣同此，但多一帶分，故其數較繁。\end{unitzh}
\begin{uniten}Note: This method is made clear by the method of setting up the celestial unknown (tianyuan). Method: set up the celestial unknown as A's straight length. Halve the difference of the southern and northern widths, obtaining 9; multiply the celestial unknown by it, obtaining 9 yuan (元); divide it by the total straight length, obtaining 9/150 of the celestial unknown, as A's half width-difference. Add it to the small width, obtaining 34 bu plus 9/150 yuan; multiply it again by the celestial unknown, obtaining 34 yuan plus 9/150 square (pingfang), which is then equal to the 2150 bu of field divided to each person. Multiply each of the two quantities by the denominator 150, obtaining 322,500 bu, equal to 9 squares plus 5100 yuan. The number of bu is the dividend, the number of yuan is the joined side, the number of squares is the corner (yu); this obtains A's straight length. Seeking the lengths and widths of B and C is the same as this, but there is one more attached fraction (daifen), hence the numbers are comparatively complicated.\end{uniten}

\begin{unitzh}\unitlabel{shushu-024.10}數學九章卷三下\end{unitzh}
\begin{uniten}Shuxue jiuzhang (Mathematics in Nine Sections), juan 3, part 2\end{uniten}

\section*{卷4上}\addcontentsline{toc}{section}{卷4上}

\begin{unitzh}\unitlabel{shushu-025.1}欽定四庫全書\end{unitzh}
\begin{uniten}Imperially Commissioned Complete Library of the Four Treasuries (Qinding Siku quanshu)\end{uniten}

\begin{unitzh}\unitlabel{shushu-025.2}數學九章卷四上 宋 秦九韶 撰\end{unitzh}
\begin{uniten}Shuxue jiuzhang (Mathematics in Nine Sections), juan 4, part 1. Compiled by Qin Jiushao of the Song.\end{uniten}

\begin{unitzh}\unitlabel{shushu-025.3}測望\end{unitzh}
\begin{uniten}Surveying (cewang)\end{uniten}

\begin{unitzh}\unitlabel{shushu-025.4}按：測望之法，見於晉劉徽《海島算經》，原名《重差》。其書一卷九題，法簡數宻。此卷本其法而擴充之，於古人之意實多所發明。然其中譌舛之處，較他卷尤甚，今悉為正之。至術有未合者，更設法以附其後焉。\end{unitzh}
\begin{uniten}Siku editors' note: The method of surveying is found in the Haidao suanjing (Sea Island Mathematical Manual) of Liu Hui of the Jin, originally named Chongcha (Double Differences). That book has 1 juan and 9 problems; its methods are concise and its numbers precise. This juan bases itself on that method and extends it, and truly brings much to light concerning the intentions of the ancients. However, the corrupt and erroneous places in it are even worse than in the other juan; we have now corrected them all. As for methods that do not agree, we have further devised methods and appended them after them.\end{uniten}

\begin{unitzh}\unitlabel{shushu-025.5}望山高逺\end{unitzh}
\begin{uniten}Viewing the height and distance of a mountain\end{uniten}

\begin{unitzh}\unitlabel{shushu-025.6}問：名山去城，不知高逺。城外平地有木一株，高二丈三尺，假為前表。乃立後表，與木齊高，相去一百六十四歩。先退前表三丈九寸，次退後表三丈一尺三寸，斜望山峯，各與其表之端參合。人目高五尺，里法三百六十歩，歩法五尺。欲知山高及逺各㡬何？\end{unitzh}
\begin{uniten}Question: A famous mountain is at a distance from a city; its height and distance are not known. On the level ground outside the city there is a tree, 2 zhang (丈) 3 chi (尺) high, which is taken provisionally as the front gnomon (biao). Then a rear gnomon is set up, of the same height as the tree, the two being 164 bu (歩) apart. First one retreats (tui) from the front gnomon 3 zhang 9 cun (寸), next one retreats from the rear gnomon 3 zhang 1 chi 3 cun, and sights obliquely at the mountain peak, in each case aligning it with the top of the gnomon. The person's eye is 5 chi high; the li (里) rule (lifa) is 360 bu; the bu rule (bufa) is 5 chi. One wishes to know how much the mountain's height and its distance each are.\end{uniten}

\begin{unitzh}\unitlabel{shushu-025.7}答曰：髙二十里半零三歩五分歩之三；逺二十七里三百二十八歩五百七十五分歩之六十七。\end{unitzh}
\begin{uniten}Answer: Height 20 1/2 li and 3 3/5 bu; distance 27 li 328 67/575 bu.\end{uniten}

\begin{unitzh}\unitlabel{shushu-025.8}按：術數誤，後入目距山係三十五里二百三十九歩一尺三寸，其故詳後。\end{unitzh}
\begin{uniten}Siku editors' note: The numbers of the method are wrong. The distance from the rear person's eye [Translator's note: the text has 入 "enter", evidently for 人 "person"] to the mountain is 35 li 239 bu 1 chi 3 cun; the reason is set out in detail below.\end{uniten}

\begin{unitzh}\unitlabel{shushu-025.9}術曰：以勾股求之，重差入之。置二退表相減，餘為高法。通表間，併法於上。以目高減表高，餘乗上為寔。寔如法而一，得山高。以法乗表高為逺法{\zhsmall 按：此條法誤，應以法乗表高與人目去地之較}。以退後表乗髙寔為逺寔，寔如法而一，得山逺。\end{unitzh}
\begin{uniten}Method: Solve it by base and altitude (gougu), entering it by double differences (chongcha). Set out the two retreats from the gnomons (tuibiao) and subtract one from the other; the remainder is the height divisor (gaofa). Convert the distance between the gnomons (biaojian) and add the divisor to it, in the upper position. Subtract the eye height from the gnomon height; multiply the upper by the remainder to make the dividend (shi). Divide the dividend by the divisor (fa); one obtains the mountain's height. Multiply the gnomon height by the divisor to make the distance divisor (yuanfa) [Siku editors: This rule is wrong; one should multiply the difference between the gnomon height and the height of the person's eye above the ground by the divisor]. Multiply the height dividend (gaoshi) by the retreat from the rear gnomon to make the distance dividend (yuanshi); divide the dividend by the divisor; one obtains the mountain's distance.\end{uniten}

\begin{unitzh}\unitlabel{shushu-025.10}望山髙逺圖{\zhsmall 按：舊圖畫山木在術前，今山改移於此。}\end{unitzh}
\begin{uniten}Diagram of viewing the height and distance of the mountain [Siku editors: The old diagram, drawing the mountain and the tree, was before the method; now the mountain has been moved here].\end{uniten}

\begin{unitzh}\unitlabel{shushu-025.11}草曰：置後退表三丈一尺三寸，減前退表三丈九寸，餘四寸，為高法。置表去木一百六十四歩，以歩法五十寸通，得八千二百寸，為表間。併法四寸，得八千二百四寸於上。以目髙五尺減表高二丈三尺，餘通之為一百八十寸，乗上，得一百四十七萬六千七百二十寸，為髙實。實如髙法四寸而一，得三十六萬九千一百八十寸，為積寸。次以歩法五十寸約之，得七千三百八十三寸五分歩之三。次以里法三百六十歩約之，得二十里一百八十三歩五分歩之三，為山髙{\zhsmall 按：此所得係人目上之山高，若加人目高，則多一歩}。\end{unitzh}
\begin{uniten}Working: Set out the rear retreat from the gnomon, 3 zhang 1 chi 3 cun; subtract the front retreat from the gnomon, 3 zhang 9 cun; the remainder, 4 cun, is the height divisor. Set out the distance of the gnomon from the tree, 164 bu; convert it by the bu rule, 50 cun, obtaining 8200 cun as the distance between the gnomons. Add the divisor, 4 cun, obtaining 8204 cun, in the upper position. Subtract the eye height, 5 chi, from the gnomon height, 2 zhang 3 chi; convert the remainder, making 180 cun; multiply the upper by it, obtaining 1,476,720 cun as the height dividend. Divide the dividend by the height divisor, 4 cun, obtaining 369,180 cun as the accumulated cun (jicun). Next reduce this by the bu rule, 50 cun, obtaining 7383 cun 3/5 bu [Translator's note: 寸 cun here evidently stands for 歩 bu]. Next reduce this by the li rule, 360 bu, obtaining 20 li 183 3/5 bu as the mountain's height [Siku editors: What is obtained here is the height of the mountain above the person's eye; if the height of the person's eye is added, it is 1 bu more].\end{uniten}

\begin{unitzh}\unitlabel{shushu-025.12}次以法四寸乘表高二丈三尺，得九百二十，為逺法{\zhsmall 按：誤同前}。以退後表三丈一尺三寸乗髙實一百四十七萬六千七百二十寸，得四億六千二百二十一萬三千三百六十寸，為逺實。實如逺法九百二十寸而一，得五十萬二千四百五寸二十三分寸之一十九，為積寸。乃以歩法五十寸乗逺法九百二十，得四萬六千寸為法，亦除逺實，得一萬四十八歩，不盡五千三百六十，與法求等，得八十，俱以約之，得五百七十五分歩之六十七。又以里法三百六十歩約，得二十七里三百二十八歩五百七十五分歩之六十七，為山去後表入立望處等。圖如後。\end{unitzh}
\begin{uniten}Next multiply the gnomon height, 2 zhang 3 chi, by the divisor, 4 cun, obtaining 920 as the distance divisor [Siku editors: The error is the same as before]. Multiply the height dividend, 1,476,720 cun, by the retreat from the rear gnomon, 3 zhang 1 chi 3 cun, obtaining 462,213,360 cun as the distance dividend. Divide the dividend by the distance divisor, 920 cun, obtaining 502,405 19/23 cun as the accumulated cun. Then multiply the distance divisor, 920, by the bu rule, 50 cun, obtaining 46,000 cun as the divisor, and with it likewise divide the distance dividend, obtaining 10,048 bu, with 5360 not exhausted. Seek the common measure (deng) of this and the divisor, obtaining 80; reduce both by it, obtaining 67/575 bu. Further reduce by the li rule, 360 bu, obtaining 27 li 328 67/575 bu as the distance of the mountain from the place at the rear gnomon where the person [Translator's note: the text has 入 "enter", evidently for 人 "person"] stands to sight, etc. The diagrams are as follows.\end{uniten}

\begin{unitzh}\unitlabel{shushu-025.13}按：術中求山髙法合；其求逺以表髙乗高法為逺法，則誤。葢本法應即以高法為逺法，以退後表乗表問并法為寔，即得後人目距山之逺。今以退後表乗髙寔為寔，而髙寔乃目高減表高乗表減併法之數，則逺法亦當以目髙減表髙乗髙法。今即以表髙乗之，則法數大，故得數小也。\end{unitzh}
\begin{uniten}Siku editors' note: In the method, the rule for finding the mountain's height is correct; but its finding of the distance by multiplying the height divisor by the gnomon height to make the distance divisor is wrong. For by the proper rule one should simply take the height divisor as the distance divisor, and multiply the distance between the gnomons plus the divisor [Translator's note: the text has 表問; 問 "ask" evidently for 間 "between"] by the retreat from the rear gnomon to make the dividend; then one obtains the distance from the rear person's eye to the mountain. Now [the method] multiplies the height dividend by the retreat from the rear gnomon to make the dividend, and the height dividend is the number obtained by multiplying the distance between the gnomons plus the divisor [Translator's note: the text has 表減; 減 "subtract" evidently for 間] by the gnomon height less the eye height; so the distance divisor too should be the height divisor multiplied by the gnomon height less the eye height. Now it is multiplied simply by the gnomon height; thus the divisor is large, and therefore the number obtained is small.\end{uniten}

\begin{unitzh}\unitlabel{shushu-025.14}乃以頭位八千二百四寸乗中一百八十寸，得一百四十七萬六千七百二十寸，為髙寔。\end{unitzh}
\begin{uniten}Then multiply the middle [position], 180 cun, by the head position (touwei), 8204 cun, obtaining 1,476,720 cun as the height dividend.\end{uniten}

\begin{unitzh}\unitlabel{shushu-025.15}子部，天文算法類，算書之屬，《數學九章》，卷四上\end{unitzh}
\begin{uniten}Masters section (zibu), Astronomy and Mathematics category, Mathematical Books subcategory; Shuxue jiuzhang, juan 4, part 1 [Translator's note: a classification heading of the Siku quanshu copy, reproduced at this point in the transcription].\end{uniten}

\begin{unitzh}\unitlabel{shushu-025.16}乃以歩法五十寸乗中位逺法九百二十寸，得下位四萬六千寸，為後圖中位歩寸法。\end{unitzh}
\begin{uniten}Then multiply the distance divisor in the middle position, 920 cun, by the bu rule, 50 cun, obtaining in the lower position 46,000 cun, which is the bu-cun divisor in the middle position of the following diagram.\end{uniten}

\begin{unitzh}\unitlabel{shushu-025.17}乃以中除上，得下位里數及零歩；其不盡寸與法求，得八十，俱約之，為歩分母子之數。\end{unitzh}
\begin{uniten}Then divide the upper by the middle, obtaining in the lower position the number of li and the odd bu; the cun not exhausted, sought together with the divisor, give 80; reduce both by it to make the numbers of the denominator and numerator of the fraction of a bu.\end{uniten}

\begin{unitzh}\unitlabel{shushu-025.18}臨臺測水\end{unitzh}
\begin{uniten}Measuring water from a terrace overlooking it\end{uniten}

\begin{unitzh}\unitlabel{shushu-025.19}問：臨水城臺，立髙三丈，其上架樓，其下址側脚濶二尺。謢岸排沙下橋，去址一丈二尺，外椿露土髙五尺，與址下平。遇水漲時，浸至址。今水退，不知多少。人從樓上欄杆腰中駕一竿出外，斜望水際，得四尺一寸五分，乃與竿端參合。人目髙五尺。欲知水退立深、涸岸斜長，自臺址至水際，各幾何？\end{unitzh}
\begin{uniten}Question: A city terrace overlooking the water stands 3 zhang high, with a tower built upon it; at its foundation below, the batter (cejiao, literally "side foot") is 2 chi wide. Protecting the bank and holding back the sand, the lower bridge (xiaqiao) [Translator's note: 橋 "bridge" is perhaps an error for 椿 "pile"] is 1 zhang 2 chi from the foundation; the outer piles stand out of the earth 5 chi high, level with the bottom of the foundation. When the water rises, it floods up to the foundation. Now the water has receded, and it is not known by how much. A person, from the waist of the railing on the tower, mounts a pole projecting outward and sights obliquely at the water's edge; he obtains 4 chi 1 cun 5 fen (分), and then [the sight] is aligned with the end of the pole. The person's eye is 5 chi high. One wishes to know the vertical depth of the water's recession and the slant length of the dried bank from the terrace foundation to the water's edge: how much is each?\end{uniten}

\begin{unitzh}\unitlabel{shushu-025.20}按：算題固不厭其難，然必簡而不漏，繁而不贅，始為合作。如此題本意，謂竿端與臺址上下懸直，則「側脚闊二尺」句已贅；又不明言人目距臺邊逺近，皆故為黯黮也。\end{unitzh}
\begin{uniten}Siku editors' note: A mathematical problem certainly does not shun difficulty, but only when it is concise without omissions and elaborate without superfluity is it a proper composition. As for this problem, its original intention is that the end of the pole and the terrace foundation hang in a vertical line above and below; then the phrase "the batter is 2 chi wide" is already superfluous. Moreover, it does not state clearly how far the person's eye is from the edge of the terrace. All this deliberately makes it obscure.\end{uniten}

\begin{unitzh}\unitlabel{shushu-025.21}術曰：以勾股變法兼少廣求之。求涸岸斜長：置出竿乘臺髙為叚；以去基乘叚為闊泛；以岸髙乘叚為淺泛；以目髙乘去基為約泛。三泛可約者約之，為定率；不可約，徑為率。以闊率自乘為闊冪；以淺率自乘為淺冪；併闊淺二冪，共為竣冪。復乘闊冪於上，以臺髙冪乘上為竣實。次以闊率乘淺率為寄。以臺髙數乘闊率，又乘約率，得數内減寄，餘自乘為竣隅。騐峻實、峻隅兩者，可約，求等約之，為峻定實、峻定隅。開同休連枝平方，得峻岸斜長。{\zhsmall 同休格：先以隅開平方，得數名同隅；以同隅乘定實，開之，得數為寔；以同隅為法除之，得竣斜。 按：此條誤，草中乃即以定實開平方，得數以同隅除之，為竣斜也。}求水退深：基岸髙冪乘竣定實為深實；以去岸冪併岸為冪，乗竣定為深隅，{\zhsmall 其深實、深□可約，約之，仍以同休格入之。}開連枝平方，得水退深。\end{unitzh}
\begin{uniten}Method: Solve it by the modified method of base and altitude (gougu bianfa) together with shaoguang ("Lesser Breadth"). To find the slant length of the dried bank: set out the projecting pole and multiply it by the terrace height to make the segment (duan). Multiply the segment by the distance from the base to make the provisional breadth (kuofan). Multiply the segment by the bank height to make the provisional shallowness (qianfan). Multiply the distance from the base by the eye height to make the provisional reducer (yuefan). If the three provisional [numbers] can be reduced, reduce them to make the fixed rates (dinglü); if they cannot be reduced, take them directly as the rates (lü). Multiply the breadth rate by itself to make the breadth area (mi); multiply the shallowness rate by itself to make the shallowness area; add the two areas of breadth and shallowness together to make the slope area (junmi) [Translator's note: here and in several places below the text has 竣 "complete" for 峻 "steep, slope"]. Multiply it again by the breadth area, in the upper position; multiply the upper by the area of the terrace height to make the slope dividend (junshi). Next multiply the shallowness rate by the breadth rate to make the lodged [number] (ji). Multiply the breadth rate by the number of the terrace height, and multiply further by the reducer rate; from the number obtained subtract the lodged [number]; multiply the remainder by itself to make the slope corner (junyu). Examine the two, the slope dividend and the slope corner: if they can be reduced, seek the common measure (deng) and reduce them, making the fixed slope dividend and the fixed slope corner. Extract the square root by the tongxiu connected-branches (tongxiu lianzhi) method [Translator's note: 休 is probably an error for 体, i.e. 體, "body"]; one obtains the slant length of the slope bank. [Note: The tongxiu pattern (ge): first extract the square root of the corner; the number obtained is called the same corner (tongyu). Multiply the fixed dividend by the same corner and extract its root; the number obtained is the dividend; divide it by the same corner as divisor; one obtains the slope's slant.] [Siku editors: This rule is wrong; in the working one directly extracts the square root of the fixed dividend and divides the number obtained by the same corner to make the slope's slant.] To find the depth of the water's recession: multiply the fixed slope dividend by the area of the base-bank height to make the depth dividend (shenshi); take the area of the distance from the bank, add the bank to make an area, and multiply it by the fixed slope to make the depth corner (shenyu) [Translator's note: these two sentences appear garbled]. [Note: If the depth dividend and the depth □ [character lost] can be reduced, reduce them, and still enter them by the tongxiu pattern.] Extract the square root by connected branches (lianzhi); one obtains the depth of the water's recession.\end{uniten}

\begin{unitzh}\unitlabel{shushu-025.22}臨䑓測水圖{\zhsmall 按：舊圖畫樓䑓，不畫正髙，在術前，今改正，移于此。}\end{unitzh}
\begin{uniten}Diagram of measuring water from the terrace overlooking it [Siku editors: The old diagram drew the tower and terrace but did not draw the true height, and it was before the method; now it has been corrected and moved here].\end{uniten}

\begin{unitzh}\unitlabel{shushu-025.23}草曰：以出竿四十一寸五分乘䑓髙三十尺，得一百二十四尺五寸，為段。以去\end{unitzh}
\begin{uniten}Working: Multiply the terrace height, 30 chi, by the projecting pole, 41 cun 5 fen, obtaining 124 chi 5 cun as the segment (duan). With the distance from [continues]\end{uniten}

\begin{unitzh}\unitlabel{shushu-026.1}址一十二尺乘叚，得一千四百九十四，以為闊泛。以䕶岸髙五尺乘段一百二十四尺五寸，得六百二十二尺五寸，為淺泛。以目髙五尺乗去址一十二尺，得六十尺，為約泛。以闊泛、淺、約泛三者求等，等得一尺五寸。皆以約，其闊泛得九百九十六尺為闊率，其淺泛得四百一十五尺為淺率，其約泛得四十尺為約率。\end{unitzh}
\begin{uniten}[continued] …the base, 12 chi (尺), multiply the segment (duan), obtaining 1,494, which is taken as the broad provisional (kuo fan). With the height of the protecting bank, 5 chi, multiply the segment, 124 chi 5 cun (寸), obtaining 622 chi 5 cun, as the shallow provisional (qian fan). With the eye height, 5 chi, multiply the distance from the base, 12 chi, obtaining 60 chi, as the reducing provisional (yue fan). With the three—the broad provisional, the shallow, and the reducing provisional—find the common measure (deng); the common measure obtained is 1 chi 5 cun. Reduce them all by it: the broad provisional gives 996 chi, as the broad rate (kuo lü); the shallow provisional gives 415 chi, as the shallow rate (qian lü); the reducing provisional gives 40 chi, as the reducing rate (yue lü).\end{uniten}

\begin{unitzh}\unitlabel{shushu-026.2}以闊率九百九十六自乗，得九十九萬二千一十六尺，為闊冪。以淺率四百一十五自乗，得一十七萬二千二百二十五尺，為淺冪。倂闊、淺二冪，得一百一十六萬四千二百四十一，為竣冪。以闊冪九十九萬二千一十六尺乘竣冪，得一萬一千五百四十九億四千五百六十九萬九千八百五十六尺於上。又以臺髙三十尺自乘，得九百尺，為臺髙冪，乗上，得一千三十九萬四千五百一十一億二千九百八十七萬四百尺，為竣實。\end{unitzh}
\begin{uniten}Multiply the broad rate, 996, by itself, obtaining 992,016 chi, as the broad area (kuo mi). Multiply the shallow rate, 415, by itself, obtaining 172,225 chi, as the shallow area (qian mi). Add the two areas, broad and shallow, obtaining 1,164,241, as the steep area (jun mi) [Translator's note: 竣 'complete' is written for 峻 'steep' here and below]. With the broad area, 992,016 chi, multiply the steep area, obtaining 1,154,945,699,856 chi, placed in the upper position. Again multiply the terrace height, 30 chi, by itself, obtaining 900 chi, as the terrace-height area; multiply the upper one by it, obtaining 1,039,451,129,870,400 chi, as the steep dividend (jun shi).\end{uniten}

\begin{unitzh}\unitlabel{shushu-026.3}次以闊率九百九十六乘淺率四百一十五，得四十一萬三千三百四十，為寄。以臺髙三十乘闊率九百九十六，得二萬九千八百八十，文乘約率四十，得一百一十九萬五千二百，内减寄，餘七十八萬一千八百六十尺，自乗，得六千一百一十三億五百五萬九千六百尺，為隅。以隅與竣實求等，得二千四百八十萬四百，俱以約之，得四千一百九十一萬二千六百七十六尺為竣定實，得二萬四千六百四十九為竣定隅。開同休連枝平方，得竣岸至水際斜長。\end{unitzh}
\begin{uniten}Next, with the broad rate, 996, multiply the shallow rate, 415, obtaining 413,340, as the lodged quantity (ji). With the terrace height, 30, multiply the broad rate, 996, obtaining 29,880; wen [Translator's note: 文 is evidently an error for 又 'again'] multiply the reducing rate, 40, obtaining 1,195,200; subtract the lodged quantity from it, leaving 781,860 chi; multiply this by itself, obtaining 611,305,059,600 chi, as the corner (yu). With the corner and the steep dividend find the common measure, obtaining 24,800,400; reduce both by it, obtaining 41,912,676 chi as the steep fixed dividend (jun ding shi), and obtaining 24,649 as the steep fixed corner (jun ding yu). Extract the tongxiu linked-branch (lianzhi) square root [Translator's note: 同休 'tongxiu' is probably 同体 (= 同體) 'same body'; 休 and 体 are nearly identical in form], obtaining the slant length from the steep bank to the water's edge.\end{uniten}

\begin{unitzh}\unitlabel{shushu-026.4}驗同休格：乃以定隅二萬四千六百四十九萬為實，先以一為隅，開平方，得一百五十七，為同休法；次以竣定實四千一百九十一萬二千六百七十六尺為實，亦以一為隅，開平方，得六千四百七十四尺，為同休實。實如同休法一百五十七而一，求等得一，俱以一各約之，其法與餘只得此數，乃直命之，得四丈一尺一百五十七分尺之三十七，為涸岸斜長至水際。\end{unitzh}
\begin{uniten}Verify by the tongxiu pattern (ge): then take the fixed corner, 24,649 wan [Translator's note: the text adds a superfluous 萬 'ten thousand'; the fixed corner above is 24,649], as the dividend; first, with 1 as the corner, extract the square root, obtaining 157, as the tongxiu divisor (fa); next take the steep fixed dividend, 41,912,676 chi, as the dividend, and also with 1 as the corner extract the square root, obtaining 6,474 chi, as the tongxiu dividend (shi). Divide the dividend by the tongxiu divisor, 157; find the common measure, obtaining 1; reduce each of them by 1: the divisor and the remainder yield only these numbers; so name it directly, obtaining 4 zhang (丈) 1 37/157 chi, as the slant length of the dry bank to the water's edge.\end{uniten}

\begin{unitzh}\unitlabel{shushu-026.5}求退水深：置岸髙五尺，自乗，得二十五，為岸髙冪，乗竣定實四千一百九十一萬二千六百七十六尺，得一十億四千七百八十一萬六千九百，為深泛。以去岸一十二尺自乗，得一百四十四尺，為去岸冪，併岸髙幕二十五尺，得一百六十九尺，以乗竣定隅二萬四千六百四十九，得四百一十六萬五千六百八十一，為隅泛。置二泛求等，得一百六十九，俱約二泛，得六百二十萬一百為定實，得二萬四千六百四十九為深定隅。開連枝平方，得水退立深。\end{unitzh}
\begin{uniten}To find the depth of the receded water: set up the bank height, 5 chi, and multiply it by itself, obtaining 25, as the bank-height area; with it multiply the steep fixed dividend, 41,912,676 chi, obtaining 1,047,816,900, as the depth provisional (shen fan). Multiply the distance from the bank, 12 chi, by itself, obtaining 144 chi, as the distance-from-bank area; add the bank-height area, 25 chi, obtaining 169 chi; with this multiply the steep fixed corner, 24,649, obtaining 4,165,681, as the corner provisional (yu fan). Set up the two provisionals and find the common measure, obtaining 169; reduce both provisionals by it, obtaining 6,200,100 as the fixed dividend (ding shi), and obtaining 24,649 as the depth fixed corner. Extract the linked-branch square root, obtaining the vertical depth of the water's recession.\end{uniten}

\begin{unitzh}\unitlabel{shushu-026.6}驗同休格：乃以深定隅二萬四千六百四十九為實，先以一為隅，開平方，得一百五十七，為同休法；次以深實實六百二十萬一百為實，亦以一為隅，開平方，得二千四百九十，為同休實。實如法一百五十七而一，得一十五尺，不盡一百三十五，與法求等，得一，俱以一各約，法與只得此數，乃直命之，得一丈五尺一百五十七分尺之一百三十五，為水退立深數也。\end{unitzh}
\begin{uniten}Verify by the tongxiu pattern: then take the depth fixed corner, 24,649, as the dividend; first, with 1 as the corner, extract the square root, obtaining 157, as the tongxiu divisor; next take the depth dividend—dividend 6,200,100—as the dividend, and also with 1 as the corner extract the square root, obtaining 2,490, as the tongxiu dividend. Divide the dividend by the divisor, 157, obtaining 15 chi, with 135 not exhausted; with the divisor find the common measure, obtaining 1; reduce each by 1; the divisor and — yield only these numbers; so name it directly, obtaining 1 zhang 5 135/157 chi, as the number of the vertical depth of the water's recession.\end{uniten}

\begin{unitzh}\unitlabel{shushu-026.7}按：此條術雖甚繁，理數皆極精宻，非兼通於勾股、通分之法者不能立也。但累乗累除，錯綜變換，皆未嘗明言，其不能無金鍼不度之疑。今繪圖以之，並條拆其乘除各數于後。\end{unitzh}
\begin{uniten}Note: although the method of this item is very complicated, its principles and numbers are all extremely refined and precise; one not versed in both the methods of base and altitude (gougu) and of bringing to a common denominator (tongfen) could not have set it up. But its repeated multiplications and repeated divisions, intricately interwoven and transformed, are nowhere explicitly stated; this cannot but raise the suspicion that the golden needle has not been passed on. Now we draw a diagram with it, and below we break down item by item the various numbers of its multiplications and divisions.\end{uniten}

\begin{unitzh}\unitlabel{shushu-026.8}如圖：甲乙為䑓正髙，乙丙為樁去䑓址，丙丁為岸髙，乙戊為臺址至水際，即為峻斜，己庚為人目髙，甲庚為出竿，戊癸為水面正深。\end{unitzh}
\begin{uniten}[Translator's note: the cyclical signs 甲乙丙丁戊己庚辛壬癸 are rendered A–J; several labels in this passage are garbled and are rendered as written.] As in the diagram: AB is the true height of the terrace; BC is the distance of the stake from the terrace base; CD is the bank height; BE, from the terrace base to the water's edge, is the steep slant; FG is the person's eye height; AG is the protruding pole; EJ is the true depth of the water surface.\end{uniten}

\begin{unitzh}\unitlabel{shushu-026.9}題有甲乙䑓髙、乙丙樁去址、丙丁去樁、甲庚出竿、己入庚目髙，求乙辛竣斜。自丁㸃與丙丁平行相等作丁辛線，自乙㸃與丙丁並行作丁，作乙辛線，自丁㸃與戊甲平行作丁壬㸃，得壬丁辛勾股形。内有乙丁辛勾股形一，與乙丙丁辛等；有乙丁壬三角形一，與甲乙戊形同式。法當以己庚小股乗庚辛大勾，以甲庚小勾除之，得壬辛大股；次以乙丁三小角形下斜邉乗甲乙戊形直邉，以乙辛减壬辛，餘壬乙為乙丁壬形直邉，為法除之，得乙戊為甲乙戊形下斜邉，即所求䑓址至水際之峻斜。\end{unitzh}
\begin{uniten}The problem gives AB, the terrace height; BC, the stake's distance from the base; CD, the distance from the stake; AG, the protruding pole; F into G, the eye height; and seeks BH, the steep slant. From point D, parallel and equal to CD, draw the line DH; from point B, parallel to CD, draw D, draw the line BH; from point D, parallel to EA, draw the point DI; this gives the base-altitude figure (gougu xing) IDH. Within it there is one base-altitude figure BDH, equal to BCDH; there is one triangle (sanjiao xing) BDI, of the same form (tongshi) as the figure ABE. The method should be: with the small altitude FG multiply the large base GH, and divide by the small base AG, obtaining the large altitude IH; next, with the lower slant side of the BD three-small-angle figure multiply the upright side of the figure ABE; subtract BH from IH, and the remainder IB, being the upright side of the figure BDI, is taken as the divisor; dividing, one obtains BE as the lower slant side of the figure ABE, which is the sought steep slant from the terrace base to the water's edge.\end{uniten}

\begin{unitzh}\unitlabel{shushu-026.10}其法只用丨除两次，甚屬易簡。即遇數不盡者，以通分御之，𠕂加一二次乗除可以。乃必增至十餘次多者，始欲窮數之變，就一題以為諸法之例，非徒為繁難也。試依術内逓次乗除之數，逐條細論之。\end{unitzh}
\begin{uniten}Its method uses only 1 (丨) division twice, which is very easy and simple. Even where one meets numbers that are not exhausted, one manages them by bringing to a common denominator, and adding one or two more multiplications and divisions would do. That it must be increased to as many as ten-odd times is because from the outset the intention was to exhaust the variations of numbers, taking one problem as the example for all methods; it is not merely for the sake of complexity and difficulty. Let us try, following the numbers of the successive multiplications and divisions in the method, to discuss them in detail item by item.\end{uniten}

\begin{unitzh}\unitlabel{shushu-026.11}出竿{\zhsmall 甲庚}乗髙䑓{\zhsmall 甲乙}為段，去址{\zhsmall 乙丙}乗段為闊率{\zhsmall 原名闊泛，約之為闊率，今即為闊率}，為去址乗䑓髙出竿長冪之數。闊率自乗為闊冪，即如去址冪乗䑓髙出竿長冪自乗之數，又即如去址冪乗䑓髙冪又乗出竿冪之數。\end{unitzh}
\begin{uniten}The protruding pole [Note: AG] multiplied by the high terrace [Note: AB] is the segment; the distance from the base [Note: BC] multiplied by the segment is the broad rate [Note: originally named the broad provisional; reduced, it was the broad rate; now it is directly the broad rate], which is the number of the distance from the base multiplied by the oblong area (chang mi) of the terrace height and the protruding pole. The broad rate multiplied by itself is the broad area, which is the same as the number of the distance-from-base area multiplied by the oblong area of the terrace height and the protruding pole multiplied by itself, and again the same as the number of the distance-from-base area multiplied by the terrace-height area and again multiplied by the protruding-pole area.\end{uniten}

\begin{unitzh}\unitlabel{shushu-026.12}岸髙{\zhsmall 丙丁}乗段為淺率{\zhsmall 原名淺泛，約之為淺，令即淺淺率}，為岸髙乗䑓髙出竿長冪之數。淺率自乗為淺冪，即如岸髙冪乗臺髙出竿長冪自乗之數，又即如岸髙冪乗䑓髙冪又乗出竿冪之數。\end{unitzh}
\begin{uniten}The bank height [Note: CD] multiplied by the segment is the shallow rate [Note: originally named the shallow provisional; reduced, it was the shallow; let it directly be the shallow shallow rate], which is the number of the bank height multiplied by the oblong area of the terrace height and the protruding pole. The shallow rate multiplied by itself is the shallow area, which is the same as the number of the bank-height area multiplied by the oblong area of the terrace height and the protruding pole multiplied by itself, and again the same as the number of the bank-height area multiplied by the terrace-height area and again multiplied by the protruding-pole area.\end{uniten}

\begin{unitzh}\unitlabel{shushu-026.13}併闊冪、淺冪為淺冪，為竣冪，□{\zhsmall 乙丁}冪乗臺髙出竿長冪自乗之數，又即如小斜冪乗䑓冪又乗出竿竿冪之數。\end{unitzh}
\begin{uniten}Adding the broad area and the shallow area is the shallow area, is the steep area; the number of the □ [Note: BD] area multiplied by the oblong area of the terrace height and the protruding pole multiplied by itself; and again the same as the number of the small slant (xiao xie) area multiplied by the terrace area and again multiplied by the protruding-pole pole area.\end{uniten}

\begin{unitzh}\unitlabel{shushu-026.14}闊冪、峻冪相乗為上數，即如小斜冪乗去址冪，又乗臺髙冪自乗，又乗出竿冪自乗之數。\end{unitzh}
\begin{uniten}The broad area and the steep area multiplied together are the upper number (shang shu), which is the same as the number of the small slant area multiplied by the distance-from-base area, again multiplied by the terrace-height area multiplied by itself, and again multiplied by the protruding-pole area multiplied by itself.\end{uniten}

\begin{unitzh}\unitlabel{shushu-026.15}闊率、淺率相乗為寄數，即土去址、岸髙相乗，又乗䑓髙冪，又乗出竿冪之數。\end{unitzh}
\begin{uniten}The broad rate and the shallow rate multiplied together are the lodged number (ji shu), which is, tu [Translator's note: 土 'earth' stands where 如 'the same as' is expected], the number of the distance from the base and the bank height multiplied together, again multiplied by the terrace-height area, and again multiplied by the protruding-pole area.\end{uniten}

\begin{unitzh}\unitlabel{shushu-026.16}去址、目髙{\zhsmall 己庚}相乗為約率，即如出竿乗壬辛。\end{unitzh}
\begin{uniten}The distance from the base and the eye height [Note: FG] multiplied together are the reducing rate, which is the same as the protruding pole multiplied by IH.\end{uniten}

\begin{unitzh}\unitlabel{shushu-026.17}臺髙乗闊率，即如去址乗臺髙冪又乗出竿之數；又以約率乗之，即如去址、壬辛相乗冪又乗臺髙冪又乗出竿冪之數。内减寄數，餘去址、壬乙相乗冪又乗臺髙冪又乗出竿冪之，再自乗之隅數，即如壬乙冪乗去址冪，又乗臺髙冪自乗，又乗出竿冪自乗之數。\end{unitzh}
\begin{uniten}The terrace height multiplied by the broad rate is the same as the number of the distance from the base multiplied by the terrace-height area and again multiplied by the protruding pole; multiplying it again by the reducing rate, it is the same as the number of the area of the distance from the base and IH multiplied together, again multiplied by the terrace-height area and again multiplied by the protruding-pole area. Subtracting the lodged number from it leaves the area of the distance from the base and IB multiplied together, again multiplied by the terrace-height area and again multiplied by the protruding-pole area; multiplying this by itself again, the corner number (yu shu) is the same as the number of the IB area multiplied by the distance-from-base area, again multiplied by the terrace-height area multiplied by itself, and again multiplied by the protruding-pole area multiplied by itself.\end{uniten}

\begin{unitzh}\unitlabel{shushu-026.18}上數、隅數内，去址冪、臺髙冪自乗、出竿冪自乗各數皆同，則用上數乗、隅數除，即如用小斜冪乗、壬乙冪除矣。以䑓髙冪乗上數，若以隅數除之，即得竣斜{\zhsmall 乙戊}冪。但數不能盡，故約之，帶隅數開平方，所謂連枝同休法也。至闊泛、淺用于乗數，約泛用于除數，故可两邉同約，又為省算也。求水立深同此。\end{unitzh}
\begin{uniten}In the upper number and the corner number, the respective numbers—the distance-from-base area, the terrace-height area multiplied by itself, and the protruding-pole area multiplied by itself—are all the same; so using the upper number to multiply and the corner number to divide is the same as using the small slant area to multiply and the IB area to divide. Multiply the upper number by the terrace-height area; if one divides it by the corner number, one obtains the steep slant [Note: BE] area. But the numbers cannot be exhausted; therefore one reduces them and extracts the square root carrying the corner number: this is the so-called linked-branch tongxiu method. As for the broad provisional and the shallow, they are used in the multiplying numbers, and the reducing provisional is used in the dividing numbers; therefore both sides can be reduced alike, which again saves computation. Finding the vertical depth of the water is the same as this.\end{uniten}

\begin{unitzh}\unitlabel{shushu-026.19}子部，天文算法類，算書之屬，《數學九章》，卷四上。子部，天文算法類，算書之屬，《數學九章》，卷四上。子部，天文算法類，算書之屬，《數學九章》，卷四上。子部，天文算法類，算書之屬，《數學九章》，卷四上。子部，天文算法類，算書之屬，《數學九章》，卷四上。\end{unitzh}
\begin{uniten}Masters section (zibu), Astronomy and Calculation Methods category, Mathematical Books subcategory, Shuxue jiuzhang, juan 4, upper part. Masters section, Astronomy and Calculation Methods category, Mathematical Books subcategory, Shuxue jiuzhang, juan 4, upper part. Masters section, Astronomy and Calculation Methods category, Mathematical Books subcategory, Shuxue jiuzhang, juan 4, upper part. Masters section, Astronomy and Calculation Methods category, Mathematical Books subcategory, Shuxue jiuzhang, juan 4, upper part. Masters section, Astronomy and Calculation Methods category, Mathematical Books subcategory, Shuxue jiuzhang, juan 4, upper part.\end{uniten}

\begin{unitzh}\unitlabel{shushu-026.20}陟岸測水\end{unitzh}
\begin{uniten}Ascending the Bank to Measure the Water\end{uniten}

\begin{unitzh}\unitlabel{shushu-027.1}問：行師遇水，須計篾䌫搭造浮橋。今垂繩量陟岸，髙山丈。人立其上，欲測水面六闊。以六尺為矩，平持，去目下五；今矩本抵頤，遥望水彼岸，與矩端、岸相合；又望水此岸沙際，入矩端三尺四寸。人目髙五尺。其水面闊幾何？答曰：水闊二十三丈四尺六寸。{\zhsmall 按：應二十三丈八尺。}\end{unitzh}
\begin{uniten}Question: An army on the march meets water and must calculate the bamboo cables for building a pontoon bridge. Now a rope is lowered to measure the high bank: its height is 山 zhang (丈) [Translator's note: 山 "mountain" evidently stands for 三 "three"; the working has 3 zhang]. A person stands on top of it and wishes to measure the 六 width of the water surface [Translator's note: 六 "six" is evidently intrusive]. Take 6 chi (尺) for the set-square (ju 矩) and hold it level, 5 below the eye; now the base of the set-square rests against the chin. Gazing from afar at the far bank of the water, the end of the set-square and the bank coincide; again gazing at the sandy edge of this bank of the water, it enters 3 chi 4 cun (寸) in from the end of the set-square. The person's eye is 5 chi high. How wide is the water surface? Answer: The water is 23 zhang 4 chi 6 cun wide. [Siku editors: It should be 23 zhang 8 chi.]\end{uniten}

\begin{unitzh}\unitlabel{shushu-027.2}術曰：以勾股重差求之。置短去目下寸為法；以人目併岸髙，減去法，{\zhsmall 按：減法誤。}餘乗人矩端為實；如法而一，得水闊。\end{unitzh}
\begin{uniten}Method: Find it by the double differences (chongcha) of base and altitude (gougu). Set up the distance of the 短 below the eye, cun, as the divisor (fa) [Translator's note: 短 "short" evidently stands for 矩 "set-square", and a 5 is missing before 寸 "cun"; below, 人 "person" stands for 入 "entering"]. Add the person's eye to the height of the bank and subtract the divisor [Siku editors: Subtracting the divisor is wrong.]; multiply the remainder by the person's set-square end (人矩端) as the dividend (shi); divide by the divisor; this gives the width of the water.\end{uniten}

\begin{unitzh}\unitlabel{shushu-027.3}陡岸測水圗{\zhsmall 按：舊圗畫岸、水、視緑，不能在術前，令改正，移於此。}\end{unitzh}
\begin{uniten}Diagram of measuring water from a steep bank [Siku editors: The old diagram drew the bank, the water and the sight-緑, and could not stand before the method; let it be corrected and moved here.] [Translator's note: 緑 "green" evidently stands for 綫 "line".]\end{uniten}

\begin{unitzh}\unitlabel{shushu-027.4}草曰：置矩本去目下五寸為法。以入目髙五尺併岸髙三丈，得三丈五尺，通為寸，得三百五十寸，减去五寸，餘三百四十五寸，乗沙際入矩端三十四寸，得一萬一千七百三十寸，為實。實如法五寸而一，得二千三百四十六寸，展為二十三丈四尺六寸，為水闊。合問。\end{unitzh}
\begin{uniten}Working: Set up the 5 cun by which the base of the set-square is below the eye as the divisor. Add the "entering eye" (入目) height of 5 chi to the bank height of 3 zhang, giving 3 zhang 5 chi; convert (tong) into cun, giving 350 cun; subtract 5 cun, leaving 345 cun; multiply by the 34 cun by which the sandy edge enters in from the end of the set-square, giving 11730 cun as the dividend. Divide the dividend by the divisor 5 cun, giving 2346 cun; spread out (zhan) as 23 zhang 4 chi 6 cun, this is the width of the water. This answers the question.\end{uniten}

\begin{unitzh}\unitlabel{shushu-027.5}按：測望諸線皆合于人目之一㸃，其髙正當自人目計之。今减去人目距矩，自矩下計之，不得其理矣。\end{unitzh}
\begin{uniten}Note: The lines of sight in surveying all meet at the single point of the person's eye, so the height should properly be reckoned from the person's eye. Here the distance from the person's eye to the set-square is subtracted and the height is reckoned from below the set-square; this fails to grasp the principle.\end{uniten}

\begin{unitzh}\unitlabel{shushu-027.6}表望方城\end{unitzh}
\begin{uniten}Sighting a square city with poles (biao)\end{uniten}

\begin{unitzh}\unitlabel{shushu-027.7}問：敵城不知廣逺。傍城南山原林間房之，林際有木二株，南北相去一百六十步，遥與城東方面参相直。乃于二水之東相對立两表，表間與木四方平，人目以繩維之。人自東後表向西行一十步，望城東北隅，入東前表一十五步；又望城東南隅，入東前表四十八歩強半歩。里法三百六十。欲知其方廣及相去幾何？答曰：城東廣各一十二里三百二十步，城去木九里三百二十歩。\end{unitzh}
\begin{uniten}Question: An enemy city, whose breadth and distance are not known. Near the city, in the woods of the hills and plains to the south, one 房 it [Translator's note: 房 "house" evidently stands for 望 "look at"]: at the edge of the woods there are two trees, 160 bu (步) apart north to south, which in the distance are in alignment with the eastern face of the city. Then, east of the two 水 [Translator's note: 水 "water" evidently stands for 木 "tree"], two poles are set up opposite them, the interval of the poles and the trees being square and level on all four sides; with the person's eye they are tied together by ropes. A person walks 10 bu westward from the eastern rear pole and sights the northeast corner of the city: it enters 15 bu in from the eastern front pole; again he sights the southeast corner of the city: it enters 48 bu and a strong half (qiangban) bu in from the eastern front pole. The li (里) standard (lifa) is 360. One wishes to know the breadth of the square and the distance apart: how much are they? Answer: The city's eastern breadth is 12 li 320 bu each; the city is 9 li 320 bu from the trees.\end{uniten}

\begin{unitzh}\unitlabel{shushu-027.8}按：答數皆誤。今推得城方廣各一十一里二百二十歩又三十一分步之二十，城東南隅至北木一里九十九歩又三十一分歩之一十二。\end{unitzh}
\begin{uniten}Note: The numbers of the answer are all wrong. Reckoning now gives: the square breadth of the city, each, 11 li 220 20/31 bu; from the southeast corner of the city to the northern tree, 1 li 99 12/31 bu.\end{uniten}

\begin{unitzh}\unitlabel{shushu-027.9}術曰：以勾股重差求之。置城東南隅景入表，减表間，餘乗表間，為城去木寔；以西方歩減城東北隅景入表，餘為法，{\zhsmall 按：此句法誤。}得城去木數。以城東北隅景入表減表間，餘乗表間，為廣實；實如前法而一，得城廣。{\zhsmall 按：此所得乃城東北隅至前木之逺，以為城廣數，誤也。}\end{unitzh}
\begin{uniten}Method: Find it by the double differences of base and altitude. Set up the image (jing 景) of the southeast corner of the city entering the pole and subtract it from the distance between the poles; multiply the remainder by the distance between the poles as the dividend for the city's distance from the trees. Subtract the westward bu (西方歩) from the image of the northeast corner of the city entering the pole; the remainder is the divisor [Siku editors: The divisor in this sentence is wrong.]; this gives the number for the city's distance from the trees. Subtract the image of the northeast corner of the city entering the pole from the distance between the poles and multiply the remainder by the distance between the poles as the dividend for the breadth; divide the dividend by the previous divisor; this gives the breadth of the city. [Siku editors: What is obtained here is the distance from the northeast corner of the city to the front tree; to take it as the number for the breadth of the city is wrong.]\end{uniten}

\begin{unitzh}\unitlabel{shushu-027.10}表望方城圖{\zhsmall 按：舊圖畫城南二木與城東面不城一直線，在術前，今改正，移于此。}\end{unitzh}
\begin{uniten}Diagram of sighting a square city with poles [Siku editors: The old diagram drew the two trees south of the city and the east face of the city as not 城 one straight line, and it stood before the method; now it has been corrected and moved here.] [Translator's note: 城 "city" here evidently stands for 成 "form".]\end{uniten}

\begin{unitzh}\unitlabel{shushu-027.11}草曰：以西行一十步減東北隅入表一十五步，餘五歩為法。以城東南隅景入表四十八步七分半減表間一百六十步，餘一百一十一歩二分半，乗表間一表六十歩，得一萬七千八百，為城去木寔；以法五歩除之，{\zhsmall 按：誤同上。}得三千五百六十歩；以里法三百六十約之，得九里三百二十步，為城去木里及步數。\end{unitzh}
\begin{uniten}Working: Subtract the westward walk of 10 bu from the 15 bu of the northeast corner entering the pole, leaving 5 bu as the divisor. Subtract the 48 bu 7 1/2 fen (分) of the image of the city's southeast corner entering the pole from the 160 bu of the distance between the poles, leaving 111 bu 2 1/2 fen; multiply by the distance between the poles, 1 表 ("pole") 60 bu [Translator's note: 表 "pole" here stands for 百 "hundred"], giving 17800 as the dividend for the city's distance from the trees; divide it by the divisor 5 bu [Siku editors: The error is the same as above.], giving 3560 bu; reduce it by the li standard 360, giving 9 li 320 bu as the li and bu of the city's distance from the trees.\end{uniten}

\begin{unitzh}\unitlabel{shushu-027.12}次置城東北隅京十表一十五步，减表間一百六十，餘一百四十五歩，乗表間一百六十，得二萬三千二百，為减實；以以前法五步除之，得四千六百四十歩；以里法三百六十約之，為一十二里三百二十歩，即城方廣里及歩數。{\zhsmall 按：誤亦同上。}合問。\end{unitzh}
\begin{uniten}Next set up the 15 bu of the city's northeast corner 京十表 ("capital ten pole") [Translator's note: 京十表 evidently stands for 景入表 "image entering the pole"] and subtract it from the distance between the poles, 160, leaving 145 bu; multiply by the distance between the poles, 160, giving 23200 as the "subtraction dividend" (减實); divide it by the previous divisor 5 bu, giving 4640 bu; reduce it by the li standard 360, giving 12 li 320 bu, which is the li and bu of the square breadth of the city [Siku editors: The error is also the same as above.]. This answers the question.\end{uniten}

\begin{unitzh}\unitlabel{shushu-027.13}按：此題之要，在二本與城東面成一直面，方城與表木方形各邉俱平。西行減城東南隅入表之較與表間成小勾股形，城東南隅入表間表間之表與城東南隅至前木成大勾股形，此二形同式，可以相比。故術草中第二求以城東北隅入表減表間之餘乗表間為實，以西行步減城東北隅入表之餘為法除實，是也。但所得為城東北隅至前木之逺，以為城廣，則誤矣。\end{unitzh}
\begin{uniten}Note: The key to this problem lies in the two trunks (二本) and the east face of the city forming one straight face, and in the sides of the square city and of the square of the poles and trees all being parallel. The difference of the westward walk subtracted from the southeast corner of the city entering the pole, together with the distance between the poles, forms a small base-and-altitude figure (gougu xing); the southeast corner of the city entering the pole-interval, the pole of the pole-interval, and the southeast corner of the city to the front tree form a large base-and-altitude figure [Translator's note: this clause is garbled, and the conclusion drawn below concerns the northeast corner]. These two figures are of the same form (tongshi) and can be put in proportion (xiangbi). Therefore, in the method and working, the second computation, taking the remainder of the northeast corner of the city entering the pole subtracted from the distance between the poles, multiplied by the distance between the poles, as the dividend, and taking the remainder of the westward bu subtracted from the northeast corner of the city entering the pole as the divisor to divide the dividend, is correct. But what is obtained is the distance from the northeast corner of the city to the front tree; to take it as the breadth of the city is wrong.\end{uniten}

\begin{unitzh}\unitlabel{shushu-027.14}又西行步減城東南隅入表之較與表間成小勾股形，城東南隅入表減表間之較與城東南隅至前木成大勾股形，此二形亦同式，可以相比。以城東南隅入表減表間之餘乗表間為實，應以西行步減城東南隅入表之為法除之，即得城東南隅之前木之逺。術草中以西行步減城東北隅入表之餘為法，故得數大七倍餘。既得城東面南北二隅距前木之逺，則相減為城廣可知矣。\end{unitzh}
\begin{uniten}Further, the difference of the westward bu subtracted from the southeast corner of the city entering the pole, together with the distance between the poles, forms a small base-and-altitude figure; the difference of the southeast corner of the city entering the pole subtracted from the distance between the poles, together with the southeast corner of the city to the front tree, forms a large base-and-altitude figure. These two figures are also of the same form and can be put in proportion. Taking the remainder of the southeast corner of the city entering the pole subtracted from the distance between the poles, multiplied by the distance between the poles, as the dividend, one should take the westward bu subtracted from the southeast corner of the city entering the pole as the divisor and divide by it; this gives the distance from the southeast corner of the city to the front tree. In the method and working, the remainder of the westward bu subtracted from the northeast corner of the city entering the pole is taken as the divisor; therefore the number obtained is more than seven times too large. Once the distances of the south and north corners of the east face of the city from the front tree have been obtained, it is clear that subtracting one from the other gives the breadth of the city.\end{uniten}

\begin{unitzh}\unitlabel{shushu-027.15}遥度城\end{unitzh}
\begin{uniten}Measuring a city from afar\end{uniten}

\begin{unitzh}\unitlabel{shushu-027.16}問：有圓城，不知周徑，四門中開。北外三里有喬木，出南門便折東門九里，乃見木。欲知城周徑各幾何？{\zhsmall 圓用故法。}答曰：徑九里，周二十七里。\end{unitzh}
\begin{uniten}Question: There is a round city whose circumference and diameter are not known; its four gates open at the middle. Outside, 3 li to the north, there is a tall tree; going out of the south gate and at once turning 東門 ("east gate") 9 li, one then sees the tree [Translator's note: 門 "gate" evidently stands for 行 "walk"]. One wishes to know the circumference and diameter of the city: how much is each? [Note: For the circle use the old method.] Answer: Diameter 9 li; circumference 27 li.\end{uniten}

\begin{unitzh}\unitlabel{shushu-027.17}術曰：以勾股夕桀求之。一為從隅；五因北外里，為從七亷；置北里冪，八因，為從五亷；以北里冪為正率，以東行冪為負率，二率差四乗北里，為益從從三亷；倍負率成五亷，為益上亷；以北里乗上亷，為實。開玲瓏九乗方，得數自乗為徑，以三因徑得周。\end{unitzh}
\begin{uniten}Method: Find it by base and altitude 夕桀 [Translator's note: 夕桀 is evidently corrupt and unintelligible]. 1 is the additive corner (cong yu 從隅); multiply the li outside to the north by 5 as the additive seventh edge (qi lian 七亷); set up the area (mi) of the north li and multiply it by 8 as the additive fifth edge; take the area of the north li as the positive rate (lü) and the area of the eastward walk as the negative rate; the difference of the two rates, multiplied by 4 and by the north li, is the subtractive (yi 益) 從從 ("additive additive") third edge; double the negative rate, 成 the fifth edge, as the subtractive upper edge (shang lian 上亷) [Translator's note: 從從 is evidently intrusive, and 成 "complete" stands for 乘 "multiply"]; multiply the upper edge by the north li as the dividend. Extract the root of the openwork (linglong 玲瓏) nine-times-multiplied power (jiu cheng fang 九乗方); multiply the number obtained by itself to give the diameter; multiply the diameter by 3 to obtain the circumference.\end{uniten}

\begin{unitzh}\unitlabel{shushu-027.18}遥度圓城{\zhsmall 圖。舊圖畫城，掛在術前，今刪改，移于此。}\end{unitzh}
\begin{uniten}Measuring a round city from afar [Note: Diagram. The old diagram drew the city and was hung before the method; now it has been cut and altered, and moved here.]\end{uniten}

\begin{unitzh}\unitlabel{shushu-027.19}草曰：以一為從隅。以五因北三里，得一十五里，為從七亷。以北三里自乗，得九里，為正率；以八因率，得七十二，為從五亷。以西行見里自乗，得八十一，為負率。以正率九減負率，餘七十二，位負差；以四因之，得二百八十八；以乗北三里，得八百六十四，係負差所乗者，為益三亷。倍負率八十一，得一百六十二，乗五㢘七十二，得一萬一千六百六十四，為益上亷。以北三里乗上亷，得三萬四千九百九十二，為實。\end{unitzh}
\begin{uniten}Working: Take 1 as the additive corner. Multiply the north 3 li by 5, giving 15 li, as the additive seventh edge. Multiply the north 3 li by itself, giving 9 li, as the positive rate; multiply the rate by 8, giving 72, as the additive fifth edge. Multiply the li walked west until sighting by itself, giving 81, as the negative rate. Subtract the positive rate 9 from the negative rate, leaving 72, which is placed as the negative difference; multiply it by 4, giving 288; multiply by the north 3 li, giving 864; since it is a product of the negative difference, it is the subtractive third edge. Double the negative rate 81, giving 162; multiply by the fifth edge 72, giving 11664, as the subtractive upper edge. Multiply the upper edge by the north 3 li, giving 34992, as the dividend.\end{uniten}

\begin{unitzh}\unitlabel{shushu-027.20}置實、亷、隅，玲瓏空隅位方亷，以約實。衆法不可超進，乃于實上定商三里。其隅與商相生，得三，為從下亷；又與商相生，入從七亷，共得二十四，為星亷；又與商相生，得七十二，為從六亷；又與商相生，入五□内，共得二百八十八；又與商相生，得八百六十四，為從四亷；又與商相生，得二千五百九十二，為正三亷，内消益三亷八百六十四訖，餘一千七百二十八，為從三亷；又與商相生，得五千一百八十四，為從二亷；又與商相生，得一萬五千五百八十二，為從上亷，内消益上亷一萬一千六百六十四訖，餘三千八百八十八，為正上亷；又與商相生，得一萬一千六百六十四，為從方。乃命上商三里，除實，適盡，得三里。以自乗之，得九里，為城圓徑之里數。又以故法圓率三因之，得二十七，為城周。\end{unitzh}
\begin{uniten}Set out the dividend, the edges and the corner in openwork, the 隅 ("corner") positions of the square and edges being empty, so as to estimate the dividend (yue shi). The various divisors (zhongfa) cannot be advanced by leaps (chaojin), so above the dividend the root (shang 商) is fixed at 3 li. The corner and the root generate each other (xiangsheng), giving 3, as the additive lower edge (xia lian 下亷); again it generates with the root and enters the additive seventh edge, together giving 24, as the "star" edge (星亷); again it generates with the root, giving 72, as the additive sixth edge; again it generates with the root and enters into the fifth [lost character], together giving 288; again it generates with the root, giving 864, as the additive fourth edge; again it generates with the root, giving 2592, as the positive third edge; within it the subtractive third edge 864 is cancelled, and when this is done 1728 remains, as the additive third edge; again it generates with the root, giving 5184, as the additive second edge; again it generates with the root, giving 15582, as the additive upper edge; within it the subtractive upper edge 11664 is cancelled, and when this is done 3888 remains, as the positive upper edge; again it generates with the root, giving 11664, as the additive square (fang 方). Then name (ming) the root above, 3 li, and eliminate the dividend: it is exactly exhausted, giving 3 li. Multiply it by itself, giving 9 li, as the number of li of the round city's diameter. Again multiply it by the circle rate of the old method, 3, giving 27, as the circumference of the city.\end{uniten}

\begin{unitzh}\unitlabel{shushu-027.21}以上求率圖，以後開方圖。實與益皆負，畫黑；商與從皆正，畫朱。\end{unitzh}
\begin{uniten}The above are diagrams for finding the rates; what follows are diagrams of root extraction (kaifang). The dividend and the subtractive terms are all negative and are drawn in black; the root and the additive terms are all positive and are drawn in vermilion.\end{uniten}

\begin{unitzh}\unitlabel{shushu-027.22}按：商、實、負字旁書已明，今皆用黒。\end{unitzh}
\begin{uniten}Note: The characters 商 (root), 實 (dividend) and 負 (negative) are written alongside, which already makes it clear; now black is used for all.\end{uniten}

\begin{unitzh}\unitlabel{shushu-027.23}子部，天文算法類，算書之屬，數學九章，卷四上\end{unitzh}
\begin{uniten}Masters section (zibu), astronomy and mathematics category, mathematical books subcategory, Shuxue jiuzhang, juan 4, first part [Translator's note: running page marker of the Siku quanshu edition.]\end{uniten}

\begin{unitzh}\unitlabel{shushu-027.24}按：凡勾股難題，用立天元一法取之，多至三乗方而至。元李冶《測圓海鏡》一百七十問，僅一題取至五乗方，猶自以為煩。此題非甚難者，乃取至九乗方，蓋未得其要也。細校術草中亷隅積實之數，與立天元一法自然相生者廻殊。且凡立天元一法，開方後未有不得所求之數者；今得數自乗始為所之數，尤于古人立法之意不合。爰另立取法，並歩算之式於後。\end{unitzh}
\begin{uniten}Note: In general, difficult base-and-altitude problems, when solved by the method of setting up the celestial unknown (tianyuan), go at most as far as the thrice-multiplied power (sancheng fang). In the Ceyuan haijing of Li Ye of the Yuan, of its 170 problems only one is taken as far as the five-times-multiplied power, and he himself still found it cumbersome. This problem is not a very difficult one, yet it is taken as far as the nine-times-multiplied power; probably its essentials were not grasped. On close collation, the numbers of the edges, corner and dividend (jishi 積實) in the method and working are entirely different from those that arise naturally from the method of setting up the celestial unknown. Moreover, with the method of setting up the celestial unknown, there has never been a case where the number sought was not obtained after the root extraction; here the number obtained only becomes the number aimed at (所之數) after being multiplied by itself, which especially disagrees with the intent of the ancients in establishing methods. Therefore a separate method of solution is set out below, together with the layout of the step-by-step computation (busuan).\end{uniten}

\begin{unitzh}\unitlabel{shushu-027.25}法：立天元一為圓城徑，加三里，得三里多一元，為大股；自之，得九里多六元多一平方，為大股冪。九里為大勾，自得之八十一里，為大勾冪。相併，得九十里多六元多一年方，為大弦冪。又以大股為小勾弦和，三里為小勾弦較，和較相乗，得九里多三元，為小股冪。二分天元之一為小勾，加小勾三里，得三里多二分元之一，為小弦；自得之九里多三元多四分平方之一，為小弦冪。\end{unitzh}
\begin{uniten}Rule: Set up the celestial unknown as the diameter of the round city; add 3 li, giving 3 li plus 1 unknown (yuan 元), as the large altitude (gu); multiply it by itself, giving 9 li plus 6 unknowns plus 1 square (pingfang 平方), as the area (mi) of the large altitude. 9 li is the large base (gou); multiplied by itself it gives 81 li, as the area of the large base. Adding them together gives 90 li plus 6 unknowns plus 1 "year-square" (年方), as the area of the large hypotenuse (xian). Further, take the large altitude as the sum (he) of the small base and hypotenuse, and 3 li as the difference (jiao) of the small base and hypotenuse; multiplying the sum and the difference together gives 9 li plus 3 unknowns, as the area of the small altitude. 1/2 of the celestial unknown is the small base; add 3 li to the small base, giving 3 li plus 1/2 unknown, as the small hypotenuse; multiplied by itself it gives 9 li plus 3 unknowns plus 1/4 square, as the area of the small hypotenuse.\end{uniten}

\begin{unitzh}\unitlabel{shushu-027.26}乃以小冪與大股冪相乗，得八十一里多八十一元多二十九平方又四分平方之一多四立方又二分之一多四分三乗方之一，寄之。又以大弦冪與小股冪相乗，得八百一十里多三百二十四元二十七平方三立方，與寄數等。両邉各减八十一里三百二十四元二十七平方三立方，得四分三乗方之一多一立方二平方又四分平方之一少二百四十三元，與七百二十九里等。各以四乗之，得一三乗方多六立方九平方九百七十二元，與二千九百一十六里等。乃以里數為實，以元數為益方，平方數為從上亷，立方数為從下亷，三乗方数為隅，開帯縱三乘方，得九里，為城徑。開方式附後。\end{unitzh}
\begin{uniten}Then multiply the small area (小冪) and the area of the large altitude together, giving 81 li plus 81 unknowns plus 29 1/4 squares plus 4 1/2 cubes (lifang 立方) plus 1/4 thrice-multiplied power; set this aside (ji). Again multiply the area of the large hypotenuse and the area of the small altitude together, giving 810 li plus 324 unknowns, 27 squares, 3 cubes, which is equal to the number set aside. From each of the two sides subtract 81 li, 324 unknowns, 27 squares, 3 cubes, giving 1/4 thrice-multiplied power plus 1 cube, 2 1/4 squares, minus 243 unknowns, equal to 729 li. Multiply each by 4, giving 1 thrice-multiplied power plus 6 cubes, 9 squares, 972 unknowns, equal to 2916 li. Then take the number of li as the dividend, the number of unknowns as the subtractive square (fang), the number of squares as the additive upper edge, the number of cubes as the additive lower edge, and the number of thrice-multiplied powers as the corner; extract the root of the thrice-multiplied power with attached lengths (daizong 帯縱), giving 9 li as the diameter of the city. The layout of the root extraction is appended below.\end{uniten}

\begin{unitzh}\unitlabel{shushu-027.27}法：列寔及方、亷、隅数，約商九里。乃以隅生商，得九，入下亷，得一十五；又以下亷生商，得一百三十五，入上亷，得一百四十四；又以商生上亷，得一\end{unitzh}
\begin{uniten}Rule: Lay out the dividend and the numbers of the square, edges and corner, and estimate the root at 9 li. Then let the corner generate (sheng) with the root, giving 9; put it into the lower edge, giving 15; again let the lower edge generate with the root, giving 135; put it into the upper edge, giving 144; again let the root generate with the upper edge, giving 1 [continues]\end{uniten}

\begin{unitzh}\unitlabel{shushu-028.1}千二百九十六，以消益方，得二百二十四，為從方。以商生從方，得二千九百一十六，減實恰盡。為開得三乘方為九里，即城徑也。\end{unitzh}
\begin{uniten}[continued] thousand 296; with it cancel (xiao) the negative square (yi fang), obtaining 224, which becomes the positive square (cong fang). [Translator's note: 1296 − 972 = 324, and 324 × 9 = 2916; 224 is evidently an error.] With the quotient (shang) generate (sheng) from the positive square, obtaining 2916; subtract it from the dividend (shi), and it is exactly exhausted. Thus what is obtained by extracting the three-multiplication power (sancheng fang) is 9 li (里), which is the diameter of the city.\end{uniten}

\begin{unitzh}\unitlabel{shushu-028.2}《數學九章》卷四上\end{unitzh}
\begin{uniten}Shuxue jiuzhang (Mathematics in Nine Sections), juan 4, part 1\end{uniten}

\begin{unitzh}\unitlabel{shushu-028.3}子部，天文算法類，算書之屬，《數學九章》\end{unitzh}
\begin{uniten}Masters section (zibu), astronomy and calculation methods category (tianwen suanfa lei), mathematical books subcategory (suanshu zhi shu), Shuxue jiuzhang. [Translator's note: Siku quanshu catalogue classification.]\end{uniten}

\section*{卷4下}\addcontentsline{toc}{section}{卷4下}

\begin{unitzh}\unitlabel{shushu-029.1}欽定四庫全書\end{unitzh}
\begin{uniten}Imperially Commissioned Complete Library of the Four Treasuries (Qinding Siku quanshu)\end{uniten}

\begin{unitzh}\unitlabel{shushu-029.2}《數學九章》卷四下 宋 秦九韶 撰\end{unitzh}
\begin{uniten}Shuxue jiuzhang (Nine Chapters on Mathematics), juan 4, part 2. Composed by Qin Jiushao of the Song.\end{uniten}

\begin{unitzh}\unitlabel{shushu-029.3}測望\end{unitzh}
\begin{uniten}Surveying by Sighting (cewang)\end{uniten}

\begin{unitzh}\unitlabel{shushu-029.4}望敵圓營\end{unitzh}
\begin{uniten}Sighting the Enemy's Circular Camp\end{uniten}

\begin{unitzh}\unitlabel{shushu-029.5}問：敵臨河為圎營，不知大小。自河南岸至其地七里，於其地立兩表，相去二步。其西表與敵營南北相直，人退西表一十二步，遥望東表，適與敵營圓邉參合。圓法用密率，里法三百六十。欲知其營周及徑各㡬何？答曰：營周六里一百二步七分步之六，徑二里。按：答數有誤，營周係六里七十六步又一萬一千九百二十一分步之四千九百四十八，徑係一里三百五十一步又一千七百零三分步之九百九十九。\end{unitzh}
\begin{uniten}Question: The enemy, beside a river, has made a circular camp, whose size is not known. From the south bank of the river to the site is 7 li (里). At the site two gnomons (biao) are set up, 2 bu (步) apart. The western gnomon is aligned south–north with the enemy camp. A person withdraws 12 bu from the western gnomon and, gazing from afar at the eastern gnomon, finds it exactly in line with the circular edge of the enemy camp. For the circle rule (yuanfa) use the accurate rate (milü); the li rule (lifa) is 360. One wishes to know what the circumference and the diameter of the camp each are. Answer: The circumference of the camp is 6 li 102 6/7 bu; the diameter is 2 li. Note: the numbers of the answer are in error. The circumference of the camp is 6 li 76 4948/11921 bu; the diameter is 1 li 351 999/1703 bu.\end{uniten}

\begin{unitzh}\unitlabel{shushu-029.6}術曰：以勾股夕桀求之。置表間自乗為勾冪，以退表自乗為股冪，併二冪為弦冪。置里通步，自之，乗勾冪為率，自乗為泛實{\zhsmall 按：此數當即為實，開𢃄從方。今不開平方，乃以此數自乗，並以此數升他數，開𢃄從三乗方，不合。}。半弦冪乗率為泛從上亷{\zhsmall 按：半弦冪即半弦自乗又倍之之數。}。以勾冪減股冪，餘四約之{\zhsmall 按：此即半勾半股各自乗相減之數。}，自乗為泛益隅{\zhsmall 三泛可約之為定。}。開連枝三乗玲瓏方，得營徑。以密率二千二乗七除為周。\end{unitzh}
\begin{uniten}Method: Seek it by the base-and-altitude (gougu) xijie (夕桀) [Translator's note: 夕桀 is obscure as written, probably corrupt.]. Set up the distance between the gnomons and multiply it by itself to make the area (mi) of the base (gou); multiply the withdrawal from the gnomon by itself to make the area of the altitude (gu); combine the two areas to make the area of the hypotenuse (xian). Set up the li, convert them into bu, multiply by itself, and multiply by the area of the base to make the rate (lü); multiply it by itself to make the provisional dividend (fan shi) [Siku editors: This number should directly be the dividend, and one should extract the square root with an accompanying following term (daicong fang). Now it does not extract the square root (pingfang), but multiplies this number by itself, and also raises the other numbers by this number, and extracts the thrice-multiplied root (sancheng fang) with an accompanying following term; this is not right.]. Halve the area of the hypotenuse and multiply by the rate to make the provisional following upper edge (fan cong shang lian) [Siku editors: Half the area of the hypotenuse is the number of the half-hypotenuse multiplied by itself and then doubled.]. Subtract the area of the base from the area of the altitude, reduce (yue) the remainder by 4 [Siku editors: This is the number obtained when the half-base and the half-altitude are each multiplied by themselves and subtracted from each other.], and multiply it by itself to make the provisional augmenting corner (fan yi yu) [Note: The three provisional numbers may be reduced to make the definitive (ding) ones.] [Translator's note: cong 從 "following" marks a positive coefficient, yi 益 "augmenting" a negative one.]. Extract the linked-branch (lianzhi) thrice-multiplied openwork (linglong) root, obtaining the camp diameter. Multiply by the accurate rate 2002 (二千二) and divide by 7 to make the circumference [Translator's note: evidently for 二十二, 22.].\end{uniten}

\begin{unitzh}\unitlabel{shushu-029.7}草曰：置表間二步，自乗得四，為勾冪；以退表一十二步，自乗得一百四十四，為股冪；以勾股二冪併之，得一百四十八，為弦冪。置七里，以里法三百六十步通之，得二千五百二十步，自乗得六百三十五萬四百，乗勾冪四，得二千五百四十萬一千六百，為率。以率自乗，得六百四十五萬二千四百一十二億八千一百五十六萬，為之實。乃半弦冪得七十四，乗率二千五百四十萬一千六百，得一十八億七千九百七十一萬八千四百，為泛從上亷。以勾冪四減股冪一百四十四，餘一百四十，以四約之，得三十五，以自乗得一千二百二十五，為泛益隅。置三泛求等，得一千二百二十五{\zhsmall 按：即泛益隅。}，俱以約之，得五千二百六十七億二千七百五十七萬七千六百為定實，一百五十三萬四千四百六十四為從上亷，一為定益隅。\end{unitzh}
\begin{uniten}Working: Set up the distance between the gnomons, 2 bu, and multiply it by itself, obtaining 4 as the area of the base; multiply the withdrawal from the gnomon, 12 bu, by itself, obtaining 144 as the area of the altitude; combine the two areas of base and altitude, obtaining 148 as the area of the hypotenuse. Set up 7 li and convert them by the li rule of 360 bu, obtaining 2520 bu; multiply by itself, obtaining 6,350,400; multiply by the area of the base, 4, obtaining 25,401,600 as the rate. Multiply the rate by itself, obtaining 645,241,281,560,000, making it the dividend. Then halve the area of the hypotenuse, obtaining 74, and multiply by the rate, 25,401,600, obtaining 1,879,718,400 as the provisional following upper edge. Subtract the area of the base, 4, from the area of the altitude, 144, leaving 140; reduce it by 4, obtaining 35; multiply it by itself, obtaining 1225 as the provisional augmenting corner. Set up the three provisional numbers and seek their common measure (deng), obtaining 1225 [Siku editors: This is precisely the provisional augmenting corner.]; reduce them all by it, obtaining 526,727,577,600 as the definitive dividend, 1,534,464 as the following upper edge, and 1 as the definitive augmenting corner.\end{uniten}

\begin{unitzh}\unitlabel{shushu-029.8}開玲瓏三乗方。乃以亷隅超二度約啇，置七百，上亷為一百五十三億四千四百六十四，益隅為一億。乃以上啇生隅，得七億，為益下亷；又以上啇生益亷，減從亷，餘一百四億四千四百六十四萬，為從上亷；又以上啇生從亷，得七百三十一億一千二百四十八萬，為從方。乃命上啇除實，實餘一百四十九億四千二十一萬七千六百。又以上啇生益隅入下亷，得一十四億，為益下亷；又以上啇生益亷，得益上亷，減從亷，餘六億四千四百六十四萬，為上亷；又以上啇生上亷入方，得七百七十六億二千四百九十六萬，為方；又以上啇生益隅入下亷，得二十一億；又以上啇生下亷為益上亷，減從亷，餘一百四十億五千五百三十六萬，為益上亷；又以上啇生益隅入下亷，得二十八億。\end{unitzh}
\begin{uniten}Extract the openwork thrice-multiplied root. Then, the edge and the corner leaping (chao) twice, estimate the quotient (yue shang): place 700; the upper edge becomes 15,300,004,464 [Translator's note: 萬 is evidently missing after 4464; the working proceeds from 15,344,640,000.], and the augmenting corner becomes 100,000,000. Then with the upper quotient generate (sheng) from the corner, obtaining 700,000,000 as the augmenting lower edge (xia lian); again with the upper quotient generate from the augmenting edge and subtract it from the following edge, leaving 10,444,640,000 as the following upper edge; again with the upper quotient generate from the following edge, obtaining 73,112,480,000 as the following square (cong fang). Then order the upper quotient to divide the dividend; the dividend leaves 14,940,217,600. Again with the upper quotient generate from the augmenting corner and enter (ru) it into the lower edge, obtaining 1,400,000,000 as the augmenting lower edge; again with the upper quotient generate from the augmenting edge, obtaining an augmenting upper edge, and subtract it from the following edge, leaving 644,640,000 as the upper edge; again with the upper quotient generate from the upper edge and enter it into the square, obtaining 77,624,960,000 as the square; again with the upper quotient generate from the augmenting corner and enter it into the lower edge, obtaining 2,100,000,000; again with the upper quotient generate from the lower edge to make an augmenting upper edge, and subtract it from the following edge, leaving 14,055,360,000 as the augmenting upper edge; again with the upper quotient generate from the augmenting corner and enter it into the lower edge, obtaining 2,800,000,000.\end{uniten}

\begin{unitzh}\unitlabel{shushu-029.9}諸法皆退：方一退，為七十七億六千二百四十九萬六千；益上亷再退，為一億四千五十五萬三千六百；益下亷三退，為二百八十萬；益隅四退，為一萬。乃於上啇之次續啇，置二十步。以續啇生隅入下亷，為二百八十二萬；又以續啇生下亷入上亷，為一億四千六百一十九萬三千六百；又以續啇生上亷減從方，餘七十四億七千一十萬八十八百。乃命續啇除實，適盡。所得七百二十步，以里法約之，得二萬，為營徑。次以密率二十二乗七百二十，得一萬五千八百四十，為實，以七除之，得二千二百六十二步七分步之六，以里法約之，得六里一百二步七分步之六，為營周。\end{unitzh}
\begin{uniten}All the divisors (fa) retreat (tui): the square retreats once, becoming 7,762,496,000; the augmenting upper edge retreats twice, becoming 140,553,600; the augmenting lower edge retreats three times, becoming 2,800,000; the augmenting corner retreats four times, becoming 10,000. Then, next after the upper quotient, continue the quotient (xu shang), placing 20 bu. With the continued quotient generate from the corner and enter it into the lower edge, making 2,820,000; again with the continued quotient generate from the lower edge and enter it into the upper edge, making 146,193,600; again with the continued quotient generate from the upper edge and subtract it from the following square, leaving 7,470,108,800. Then order the continued quotient to divide the dividend: it is exactly exhausted. What is obtained is 720 bu; reducing it by the li rule gives 20,000 (二萬) as the camp diameter [Translator's note: evidently for 二里, 2 li.]. Next, by the accurate rate, multiply 720 by 22, obtaining 15,840 as the dividend; divide it by 7, obtaining 2262 6/7 bu; reducing it by the li rule gives 6 li 102 6/7 bu as the camp circumference.\end{uniten}

\begin{unitzh}\unitlabel{shushu-029.10}子部·天文算法類·算書之屬·數學九章·卷四下 子部·天文算法類·算書之屬·數學九章·卷四下 子部·天文算法類·算書之屬·數學九章·卷四下 子部·天文算法類·算書之屬·數學九章·卷四下 子部·天文算法類·算書之屬·數學九章·卷四下\end{unitzh}
\begin{uniten}[Translator's note: catalogue label of the Siku quanshu edition, which the transcription repeats five times.]
Masters section (zibu) · Astronomy and calculation category · Mathematical books subcategory · Shuxue jiuzhang · Juan 4, part 2
Masters section (zibu) · Astronomy and calculation category · Mathematical books subcategory · Shuxue jiuzhang · Juan 4, part 2
Masters section (zibu) · Astronomy and calculation category · Mathematical books subcategory · Shuxue jiuzhang · Juan 4, part 2
Masters section (zibu) · Astronomy and calculation category · Mathematical books subcategory · Shuxue jiuzhang · Juan 4, part 2
Masters section (zibu) · Astronomy and calculation category · Mathematical books subcategory · Shuxue jiuzhang · Juan 4, part 2\end{uniten}

\begin{unitzh}\unitlabel{shushu-029.11}按：此題用平方可矣。術中所謂率者，即平方實也。乃復加自乗，開三乗方，徒為繁冗耳。且乗從亷用半弦自乗之倍數，乗隅數應用半股自乗之數；今用勾股冪較四分之一，即用半股冪半勾冪之較，比半股自乗數小一半勾自乗數，故得數較大。若轉求表間及退步，必與原數不合。試以相去七里為大勾弦和，營徑三里為倍大勾，相減得五里為大勾弦較，和相乗得較三十五為大股冪，大勾弦里自之仍得一為大勾冪。置小勾冪四步，以大股冪乗之，得百百四十步，以大勾冪除之，仍得一百四十步為小股冪，比原小股冪少四步，其術之疎可知矣。設用平方法如左。\end{unitzh}
\begin{uniten}Note: For this problem the square root (pingfang) suffices. What the method calls the rate is precisely the dividend of the square root. Yet it further multiplies it by itself and extracts the thrice-multiplied root, which merely makes it cumbersome. Moreover, as the multiplier for the following edge it uses twice the number of the half-hypotenuse multiplied by itself, and as the multiplier for the corner number it ought to use the number of the half-altitude multiplied by itself; now it uses 1/4 of the difference of the areas of the base and the altitude, that is, it uses the difference of the area of the half-altitude and the area of the half-base, which is smaller than the number of the half-altitude multiplied by itself by one number of the half-base multiplied by itself; therefore the number obtained is rather large. If one turns back to seek the distance between the gnomons and the withdrawal in bu, they will certainly not agree with the original numbers. As a trial, take the distance of 7 li as the sum of the large base and hypotenuse, and the camp diameter of 3 li as twice the large base; subtracting gives 5 li as the difference of the large base and hypotenuse; multiplying it together with the sum gives "difference" 35 as the area of the large altitude; the large base-hypotenuse li, multiplied by itself, still gives 1 as the area of the large base. Set up the area of the small base, 4 bu, and multiply it by the area of the large altitude, obtaining "hundred hundred forty" (百百四十) bu; divide it by the area of the large base, still obtaining 140 bu as the area of the small altitude, which is 4 bu less than the original area of the small altitude: the looseness of the method can thus be known. [Translator's note: this trial is garbled; evidently read "camp diameter 2 li", "the sum and difference multiplied together give 35", "the large base, 1 li, multiplied by itself", and "140 bu" for 百百四十.] Suppose the square-root rule is used, as follows.\end{uniten}

\begin{unitzh}\unitlabel{shushu-029.12}法：立天元一為營徑。相去七里，通為二千五百二十步，為大勾弦和。相減，得二千五百二十步少一元，為大勾弦較。和較相乗，得六百三十五萬零四百步二千五百二十元，為大少股冪。天元一半之，以減相去步，得二千五百二十步少二分天元之一，為大弦。自之，得六百三十五萬零四百步少二千五百二十元多四分平方之一，為大弦冪。退步十二為小股，自之，得一百四十四步，為小股冪。表間二步為小勾，自之，得四步，為小勾冪。相併，得一百四十八步，為小弦冪。\end{unitzh}
\begin{uniten}Rule: Set up the celestial unknown as the camp diameter. The distance of 7 li, converted, makes 2520 bu, as the sum of the large base and hypotenuse. Subtracting, one obtains 2520 bu less 1 yuan (元) as the difference of the large base and hypotenuse. Multiplying the sum and the difference together gives 6,350,400 bu 2520 yuan as the area of the large less altitude [Translator's note: 少 "less" is evidently misplaced; it belongs before 2520 yuan.]. Halve the celestial unknown and subtract it from the distance in bu, obtaining 2520 bu less 1/2 of the celestial unknown as the large hypotenuse. Multiplying it by itself gives 6,350,400 bu less 2520 yuan plus 1/4 of a square (pingfang) as the area of the large hypotenuse. The withdrawal in bu, 12, is the small altitude; multiplying it by itself gives 144 bu as the area of the small altitude. The distance between the gnomons, 2 bu, is the small base; multiplying it by itself gives 4 bu as the area of the small base. Combining them gives 148 bu as the area of the small hypotenuse.\end{uniten}

\begin{unitzh}\unitlabel{shushu-029.13}以小弦冪乗大股冪，得九億三千九百八十五萬九千二百步少三十七萬二千九百六十元，寄之。又以小股冪乗大弦冪，得九億一千四百四十五萬七千六百步少三十六萬二千八百八十元多三十六平方，與寄數為相等。兩邉各減九億一千四百四十五萬七千六百步，各加三十七萬二千九百六十元，得二千五百四十萬零一千六百步{\zhsmall 即術中率數}，與三十六平方多一萬零八十元等。三數求總等，得三十六。約步數，得七十萬零五千六百，為長方積，為實；約元數，得二百八十，為從方，為長闊較；約平方數，得一，為隅。用𢃄從平方法開，得闊七百一十一步又一千七百零三分步之九百九十九，為營徑步。\end{unitzh}
\begin{uniten}Multiply the area of the large altitude by the area of the small hypotenuse, obtaining 939,859,200 bu less 372,960 yuan; set it aside (ji). Again multiply the area of the large hypotenuse by the area of the small altitude, obtaining 914,457,600 bu less 362,880 yuan plus 36 squares, which is made equal (xiangdeng) to the set-aside number. From both sides subtract 914,457,600 bu and to each add 372,960 yuan, obtaining 25,401,600 bu [Note: This is the rate number in the method.], equal to 36 squares plus 10,080 yuan. For the three numbers seek the overall common measure (zongdeng), obtaining 36. Reducing the number of bu gives 705,600 as the area of the oblong (changfang), as the dividend; reducing the number of yuan gives 280 as the following square, as the difference of length and width; reducing the number of squares gives 1 as the corner. Extracting by the rule of the square root with an accompanying following term (daicong pingfang), one obtains the width, 711 999/1703 bu, as the camp diameter in bu.\end{uniten}

\begin{unitzh}\unitlabel{shushu-029.14}以密率用二十二乗之，徑七除之，得二千二百三十六步又一萬一千九百二十一分步之四千九百四十八，為營周步。各以里率収之，得營周六里七十六步又一萬一千九百二十一分步之四千九百四十八，營徑一里三百五十一步又一千七百零三分步之九百九十九。\end{unitzh}
\begin{uniten}By the accurate rate, multiply it by 22 and divide by the diameter 7, obtaining 2236 4948/11921 bu as the camp circumference in bu. Gather (shou) each by the li rate, obtaining the camp circumference 6 li 76 4948/11921 bu, and the camp diameter 1 li 351 999/1703 bu.\end{uniten}

\begin{unitzh}\unitlabel{shushu-029.15}還原之法：置營徑步數，通分納子，得一二一一八三二分，為倍大勾分數。又以分母通相去步，得四二九一五六□分，為大勾弦和分數。二數相減，餘三□七九七二八分，為大勾弦較分數。和較分數相乗，得一三二一六八三七四五九六八分，為大股冪分。倍大勾分折半，得六□五九一六，自之，得三六七一三四一九九□五六分，為大勾冪分。及以小勾冪四步乗大股冪分，得五二八六七三四九九八二七二□分，以大勾冪分除之，得一百四十四步為小股冪，與原數合。此猶用西表相去步也。若細較之，當用人目相去步，則營周當多十二步餘，營徑當多四步不足也。\end{unitzh}
\begin{uniten}Rule for returning to the origin (huanyuan): Set up the number of bu of the camp diameter, convert it into parts and add in the numerator (tongfen nazi), obtaining 1211832 parts as the number of parts of twice the large base. Again convert the distance in bu by the denominator, obtaining 429156□ parts as the number of parts of the sum of the large base and hypotenuse [Translator's note: here and below □ marks a blank in the text, evidently a lost 〇 (zero).]. Subtracting the two numbers leaves 3□79728 parts as the number of parts of the difference of the large base and hypotenuse. Multiplying together the numbers of parts of the sum and the difference gives 1321683745968 parts as the parts of the area of the large altitude. Halving the parts of twice the large base gives 6□5916; multiplying it by itself gives 367134199□56 parts as the parts of the area of the large base. Further, multiplying the parts of the area of the large altitude by the area of the small base, 4 bu, gives 5286734998272□ parts; dividing this by the parts of the area of the large base gives 144 bu as the area of the small altitude, agreeing with the original number. This still uses the distance in bu from the western gnomon. If one compares finely, one should use the distance in bu from the person's eye; then the camp circumference should be greater by 12-odd bu, and the camp diameter greater by not quite 4 bu.\end{uniten}

\begin{unitzh}\unitlabel{shushu-029.16}望知敵衆\end{unitzh}
\begin{uniten}Sighting to Know the Enemy's Numbers\end{uniten}

\begin{unitzh}\unitlabel{shushu-030.1}問：敵為圓營，在水北平沙，不知人數。諜稱彼營布卒，占地方八尺。我軍在水南山原，於下立表，髙八丈，與原山腰等平。自表端引䋲，虚量平至人足三十步。人立其處，望彼營北陵，與表端參合；又望營南陵，入表端八尺。人目髙四尺八寸。以圓密率入重差，求敵衆合得㡬何？答曰：敵衆八百四十九人。\end{unitzh}
\begin{uniten}Question: The enemy has made a circular camp on the level sands north of the river; the number of men is not known. Spies report that in that camp the soldiers are deployed, each occupying ground 8 chi (尺) square. Our army is on the mountain plateau south of the river. Below, a gnomon (biao) is set up, 8 zhang (丈) high, level with the mountainside of the plateau. From the top of the gnomon a cord is drawn and measured through the air horizontally to the man's feet: 30 bu (步). A man standing at that place sights the northern edge (ling) of that camp, and it coincides with the top of the gnomon; again he sights the southern edge of the camp, and it enters the gnomon 8 chi below its top. The man's eye is 4 chi 8 cun (寸) high. Using the close rate (milü) of the circle, brought into double differences (chongcha), find the enemy host: how many should it come to? Answer: The enemy host is 849 men.\end{uniten}

\begin{unitzh}\unitlabel{shushu-030.2}按數不合，應二百七十三人，其故詳草後。\end{unitzh}
\begin{uniten}Note (Siku editors): The number does not agree; it should be 273 men. The reason is given in detail after the working.\end{uniten}

\begin{unitzh}\unitlabel{shushu-030.3}術曰：以勾股乗之。置人退表步{\zhsmall 按此條誤，法應置表髙}，乗入表為實，以人目髙為法，除之，得徑。以密周率乗徑，得數為實；以密徑率因人立為法，約之，得外周人數，餘収為一。副置，加六，以乗副，得數為實，如一十二而一，餘亦収為全。\end{unitzh}
\begin{uniten}Method: Multiply it by base and altitude (gougu). Set up the bu by which the man steps back from the gnomon (tuibiao) [Siku editors: This item is mistaken; the method should set up the height of the gnomon], multiply by the entry into the gnomon (rubiao) to make the dividend (shi), take the height of the man's eye as the divisor (fa), and divide, obtaining the diameter. Multiply the diameter by the close circumference rate (mi zhoulü), and the resulting number is the dividend; with the close diameter rate (mi jinglü) multiply (yin) the standing room of a man to make the divisor; reduce (yue) by it, obtaining the number of men on the outer circumference; a remainder is collected (shou) as one. Set this up in duplicate (fu), add 6, and multiply the duplicate by it; the resulting number is the dividend; divide by 12; a remainder is likewise collected as a whole.\end{uniten}

\begin{unitzh}\unitlabel{shushu-030.4}望知敵衆圖{\zhsmall 按舊圖畵山水，在術前，今刪改移於此}\end{unitzh}
\begin{uniten}Diagram for "Sighting to know the enemy host" [Siku editors: The old diagram depicted mountains and water and stood before the method; it has now been pruned, revised and moved here.]\end{uniten}

\begin{unitzh}\unitlabel{shushu-030.5}草曰：置人立退表三十步，以步法五尺展為五十寸通之，得一千五百寸，乗入表八尺，得一十二萬寸，為實。以人目髙四十八寸為法，除之。以密率周法二十二乗徑二千五百，得五萬五千寸，為實。以密率徑法七因諜稱人立八尺，得五百六十，為法。以法五百六十寸約實五萬五千寸，得九十八人，為外周人數，不盡一百二十寸，棄之。副置外周九十八人，加六，得一百四人，乗副為實。\end{unitzh}
\begin{uniten}Working: Set up the 30 bu by which the man stands back from the gnomon; convert it with the bu-measure (bufa), 5 chi, expanded into 50 cun, obtaining 1500 cun; multiply by the entry into the gnomon, 8 chi, obtaining 120,000 cun as the dividend. Take the height of the man's eye, 48 cun, as the divisor and divide. With the circumference value (zhoufa) of the close rate, 22, multiply the diameter, 2500, obtaining 55,000 cun as the dividend. With the diameter value (jingfa) of the close rate, 7, multiply (yin) the 8 chi in which the spies report a man stands, obtaining 560 as the divisor. With the divisor, 560 cun, reduce the dividend, 55,000 cun, obtaining 98 men as the number of men on the outer circumference; the unexhausted remainder, 120 cun, is discarded. Set up the outer circumference, 98 men, in duplicate; add 6, obtaining 104 men; multiply the duplicate to make the dividend.\end{uniten}

\begin{unitzh}\unitlabel{shushu-030.6}按此術應置表髙加人目髙以入表，棄之，誤置退步以入表乗之，故人數差多二倍，盖思省偶未至耳。至求人數，先用密率，次用束箭法，亦未盡合題問。今依其數各步於後。\end{unitzh}
\begin{uniten}Note (Siku editors): This method ought to set up the height of the gnomon plus the height of the man's eye, with the entry into the gnomon; discarding this, it mistakenly set up the step-back and multiplied it by the entry into the gnomon; therefore the number of men is off, too many by two times. Presumably his reflection and checking simply happened not to reach this point. As for finding the number of men, first using the close rate and then using the method of the bundle of arrows (shujian fa) also does not fully accord with the question posed. Now each is calculated below according to its numbers.\end{uniten}

\begin{unitzh}\unitlabel{shushu-030.7}求營徑：置表髙八丈，加人目髙四尺八寸，得八丈四尺八寸，以入表八尺乗之，得六萬七千八百四十寸，以人目髙除之，得一十四丈一尺三寸又三分寸之一，為營徑。盖以表髙加人目髙為大股，營徑為大勾較，人目髙為小股，入表為小勾較。置大股，以小勾較乗之，以小股除之，即得大勾較也。\end{unitzh}
\begin{uniten}To find the diameter of the camp: Set up the height of the gnomon, 8 zhang; add the height of the man's eye, 4 chi 8 cun, obtaining 8 zhang 4 chi 8 cun; multiply by the entry into the gnomon, 8 chi, obtaining 67,840 cun; divide by the height of the man's eye, obtaining 14 zhang 1 chi 3 1/3 cun as the diameter of the camp. For the height of the gnomon plus the height of the man's eye is the large altitude (da gu); the diameter of the camp is the large base difference (da gou jiao); the height of the man's eye is the small altitude (xiao gu); the entry into the gnomon is the small base difference (xiao gou jiao). Set up the large altitude, multiply by the small base difference, divide by the small altitude, and one obtains the large base difference.\end{uniten}

\begin{unitzh}\unitlabel{shushu-030.8}求人數，用密率：置徑為實，倍每人占地八尺，得一十六尺，為法，除實，得八，為外層數；加最内一層，得九，為共層數；餘一十三尺三寸又三分寸之一，為最内徑。以最内徑與最外徑{\zhsmall 即營陘}相加，以九層乗之，折半，得六百九十六尺，為九層共徑數。以密周率二十二乗之，得一萬五千三百一十二尺，為實；以密率七乗每人占地八尺，得五十六尺，為法；除之，得二百七十三人，餘三尺又七分尺之三，棄之。此法應先求得内外圓周，再求九層共周數；今先求九層共徑數，然後變為圓周，其理一也。\end{unitzh}
\begin{uniten}To find the number of men, using the close rate: Set up the diameter as the dividend; double the 8 chi of ground each man occupies, obtaining 16 chi as the divisor; divide the dividend, obtaining 8 as the number of outer layers (ceng); add the one innermost layer, obtaining 9 as the total number of layers; the remainder, 13 chi 3 1/3 cun, is the innermost diameter. Add the innermost diameter and the outermost diameter [Note: that is, the camp diameter (營陘, sic)], multiply by the 9 layers, and halve, obtaining 696 chi as the sum of the diameters of the 9 layers. Multiply by the close circumference rate, 22, obtaining 15,312 chi as the dividend; with the close rate 7 multiply the 8 chi of ground each man occupies, obtaining 56 chi as the divisor; divide, obtaining 273 men; the remainder, 3 3/7 chi, is discarded. This method ought first to find the inner and outer circumferences and then find the sum of the circumferences of the 9 layers; now one first finds the sum of the diameters of the 9 layers and afterwards converts it into circumference; the principle is the same.\end{uniten}

\begin{unitzh}\unitlabel{shushu-030.9}望敵逺近\end{unitzh}
\begin{uniten}Sighting how far or near the enemy is\end{uniten}

\begin{unitzh}\unitlabel{shushu-030.10}問：敵軍處北山下原，不知相去逺近。乃於平地立一表，髙四尺。人退表九百步{\zhsmall 步法五尺}，遥望山原，適與表端參合。人目髙四尺八，欲之敵軍相去㡬何？答曰：一十二里半。\end{unitzh}
\begin{uniten}Question: An enemy army is stationed on the plain below the northern mountain, and it is not known how far or near it is. So on level ground a gnomon is set up, 4 chi high. A man steps back 900 bu from the gnomon [Note: the bu-measure is 5 chi] and from afar sights the mountain plain; it exactly coincides with the top of the gnomon. The man's eye is 4 chi 8 high; "wishing it" (欲之): how far away is the enemy army? [Translator's note: the text is evidently corrupt here; 欲之 presumably stands for 欲知 "one wishes to know".] Answer: 12 1/2 li (里).\end{uniten}

\begin{unitzh}\unitlabel{shushu-030.11}術曰：以勾股求之，重差入之。置人目髙，以表髙減之，餘為法；置退表，乗表髙，為實。實如法而一。\end{unitzh}
\begin{uniten}Method: Solve it by base and altitude, bringing it into double differences. Set up the height of the man's eye and subtract the height of the gnomon from it; the remainder is the divisor. Set up the step-back from the gnomon and multiply by the height of the gnomon to make the dividend. Divide the dividend by the divisor.\end{uniten}

\begin{unitzh}\unitlabel{shushu-030.12}草曰：置人目髙四尺八寸，減表髙四尺，餘八寸，為法。置退表九百步，以步五十寸通之，得四萬五千寸，乗表髙四十寸，得一百八十萬寸，為實。如法八寸而一，得二十二萬五千寸，以步法五十寸約之，得四千五百步，為相去步；以里法三百六十步約之，得一十二里半，為敵去表所。合問。\end{unitzh}
\begin{uniten}Working: Set up the height of the man's eye, 4 chi 8 cun; subtract the height of the gnomon, 4 chi; the remainder, 8 cun, is the divisor. Set up the step-back from the gnomon, 900 bu; convert it at 50 cun per bu, obtaining 45,000 cun; multiply by the height of the gnomon, 40 cun, obtaining 1,800,000 cun as the dividend. Divide by the divisor, 8 cun, obtaining 225,000 cun; reduce by the bu-measure, 50 cun, obtaining 4500 bu as the distance in bu; reduce by the li-measure (lifa), 360 bu, obtaining 12 1/2 li as the distance of the enemy from where the gnomon is. This agrees with the question.\end{uniten}

\begin{unitzh}\unitlabel{shushu-030.13}表望浮圖\end{unitzh}
\begin{uniten}Sighting a pagoda with a gnomon\end{uniten}

\begin{unitzh}\unitlabel{shushu-030.14}問：有浮圖欹側，欲換塔心木，不知其髙。去塔六丈有刹竿，亦不知其髙。竿本去地九尺二寸始釘鋦{\zhsmall 按鋦原本作鈎，今改正}，鋦一十四枚，枚長五寸，每鋦下股相去二尺五寸。就竿為表，人退竿三丈，遥望浮圖尖，適與竿端斜合；又望相輪之本，影鋦第七枚上股。人目入去地四尺八寸。心木放三尺為㔼{\zhsmall 按㔼原本作□，今改正}卯剪截。欲求塔髙、輪髙，合用塔心木長，各㡬何？答曰：塔髙一十一丈七尺；相輪髙三丈；塔身髙八丈七尺；竿髙四丈二尺二寸；塔心木九丈，内三尺為剪截穿鑿㔼卯。\end{unitzh}
\begin{uniten}Question: There is a pagoda (futu) that leans to one side; one wishes to replace its central post (taxinmu), but its height is not known. 6 zhang from the pagoda there is a temple mast (chagan), whose height is likewise not known. From the base of the mast, at 9 chi 2 cun above the ground, the nailing of staples (ju) begins [Siku editors: 鋦 "staple" was written 鈎 "hook" in the original edition; it has now been corrected]; there are 14 staples, each 5 cun long, and the lower legs of the staples are each 2 chi 5 cun apart. Using the mast as the gnomon, a man steps back 3 zhang from the mast and from afar sights the tip of the pagoda; it exactly coincides obliquely with the top of the mast. Again he sights the base of the finial rings (xianglun); it is "shadowed" (ying) at the upper leg of the 7th staple. The man's eye, "entering" (ru), is 4 chi 8 cun from the ground. [Translator's note: 影 and 入 are obscure here.] The central post allows 3 chi for the tenons [Siku editors: 㔼 was written □ in the original edition; it has now been corrected] and mortises and for trimming. One wishes to find the height of the pagoda, the height of the rings, and the length of central post that should be used: how much is each? Answer: The height of the pagoda is 11 zhang 7 chi; the height of the finial rings is 3 zhang; the height of the pagoda body (tashen) is 8 zhang 7 chi; the height of the mast is 4 zhang 2 chi 2 cun; the central post is 9 zhang, of which 3 chi is for trimming and for boring and chiselling the tenons and mortises.\end{uniten}

\begin{unitzh}\unitlabel{shushu-030.15}按塔髙、竿髙二數合，相輪髙、塔身髙、塔心木長三數俱誤。相輪髙四丈五尺，塔身髙七丈二尺，塔心木長七丈五尺，説詳草後。\end{unitzh}
\begin{uniten}Note (Siku editors): The two numbers for the height of the pagoda and the height of the mast are correct; the three numbers for the height of the finial rings, the height of the pagoda body and the length of the central post are all wrong. The height of the finial rings is 4 zhang 5 chi, the height of the pagoda body 7 zhang 2 chi, and the length of the central post 7 zhang 5 chi; the explanation is given in detail after the working.\end{uniten}

\begin{unitzh}\unitlabel{shushu-030.16}術曰：以勾股求之，重差入之。置鋦數，減一，餘乗鋦相去數，併一枚長數，加竿本，共為表竿髙。以退表為法。以人目髙減表竿髙，餘乗竿去塔，為實。實如法而一，得數，加表竿髙，共為塔髙。置相輪入本鋦數，減一，餘乗鋦相去，又乗竿去塔，為實，實如法而一，得相輪髙{\zhsmall 按未加入鋦人數，法誤}。以減塔髙，餘為塔身髙。以益㔼卯尺數，為塔心木長。\end{unitzh}
\begin{uniten}Method: Solve it by base and altitude, bringing it into double differences. Set up the number of staples and subtract 1; multiply the remainder by the distance between staples; combine with the length of one staple; add the base of the mast; together this makes the height of the gnomon-mast (biaogan). Take the step-back from the gnomon as the divisor. Subtract the height of the man's eye from the height of the gnomon-mast; multiply the remainder by the distance from the mast to the pagoda to make the dividend. Divide the dividend by the divisor; add the resulting number to the height of the gnomon-mast; together this makes the height of the pagoda. Set up the number of the staple into which the base of the finial rings enters (相輪入本鋦數); subtract 1; multiply the remainder by the distance between staples, and again multiply by the distance from the mast to the pagoda to make the dividend; divide the dividend by the divisor, obtaining the height of the finial rings [Siku editors: The number for the entry into the staples (入鋦人數, sic) has not been added; the method is wrong]. Subtract it from the height of the pagoda; the remainder is the height of the pagoda body. Increase it by the number of chi for the tenons and mortises to make the length of the central post.\end{uniten}

\begin{unitzh}\unitlabel{shushu-030.17}草曰：置鋦一十四枚，減一，餘一十三，以乗鋦相去二尺五寸，得三百二十五寸；併最上鋦一枚長五寸，得三百三十寸；又加竿本九尺二寸，共得四百二十二尺，為表竿髙。以人退表三丈，通為三百寸，為法。次以人目髙四尺八寸減表竿髙四百二十二寸，餘三百七千四寸，以乗竿去塔六丈，得二十二萬四千四百寸，為實。實如法三百兩一，得七百四十八寸，加表竿髙四百二十二寸，得一千一百七十寸，以十約之，為一十一丈七尺，為塔髙。\end{unitzh}
\begin{uniten}Working: Set up the 14 staples, subtract 1, remainder 13; multiply by the distance between staples, 2 chi 5 cun, obtaining 325 cun; combine with the length of the one topmost staple, 5 cun, obtaining 330 cun; further add the base of the mast, 9 chi 2 cun, obtaining altogether 422 chi [sic] as the height of the gnomon-mast. Take the man's step-back from the gnomon, 3 zhang, converted into 300 cun, as the divisor. Next subtract the height of the man's eye, 4 chi 8 cun, from the height of the gnomon-mast, 422 cun; the remainder is "3 hundred 7 thousand 4" (三百七千四) cun [Translator's note: evidently an error for 374]; multiply by the distance from the mast to the pagoda, 6 zhang, obtaining 224,400 cun as the dividend. Divide the dividend by the divisor, 300 [Translator's note: the text has 兩一 for the usual 而一], obtaining 748 cun; add the height of the gnomon-mast, 422 cun, obtaining 1170 cun; reduce by 10, making 11 zhang 7 chi as the height of the pagoda.\end{uniten}

\begin{unitzh}\unitlabel{shushu-030.18}置相輪本入第七鋦，減一，餘六，以乗鋦相去二尺五寸，得一百五十寸；又乗竿去塔六丈，得九萬寸，為實。實如前法三百寸而一，得三百寸，約為三丈，得相輪髙{\zhsmall 按不加入鋦即為相輪髙，誤}。以相輪髙三丈減塔髙一十一丈七尺，餘八丈七尺，為塔身髙{\zhsmall 按此數及塔心木數皆因上數而誤}。益三尺為剪截㔼卯，共得九丈，為塔心木長。合前問。\end{unitzh}
\begin{uniten}Set up the base of the finial rings entering the 7th staple; subtract 1, remainder 6; multiply by the distance between staples, 2 chi 5 cun, obtaining 150 cun; again multiply by the distance from the mast to the pagoda, 6 zhang, obtaining 90,000 cun as the dividend. Divide the dividend by the previous divisor, 300 cun, obtaining 300 cun, reduced to 3 zhang; this gives the height of the finial rings [Siku editors: To take this directly as the height of the finial rings without adding the entry into the staples is wrong]. Subtract the height of the finial rings, 3 zhang, from the height of the pagoda, 11 zhang 7 chi; the remainder, 8 zhang 7 chi, is the height of the pagoda body [Siku editors: This number and the number for the central post are both wrong because of the number above]. Increase it by 3 chi for trimming and the tenons and mortises, obtaining altogether 9 zhang as the length of the central post. This agrees with the question above.\end{uniten}

\begin{unitzh}\unitlabel{shushu-030.19}按此皆大小形同式相求法也。人目去塔為總勾；人目上塔尖髙、塔身皆為總股髙；相輪髙為總股較；人目去竿為分小勾；人目上竿髙及相輪本入竿髙俱為分小股；相輪本入竿為小股較；竿去塔為分大勾。術以竿去塔分大勾與小勾股乗除，得大股數，加小股為總股，故塔尖髙數合；以大勾與小勾、小股較乗除，得大股較數，即為總股較，故相輪髙數少一小股較一丈五尺也。塔身髙、塔心木長皆本此數加減而得，故誤數相等。\end{unitzh}
\begin{uniten}Note (Siku editors): These are all methods of finding one from another by large and small figures of the same form. The distance from the man's eye to the pagoda is the total base (zong gou); the height of the pagoda tip above the man's eye, and the pagoda body, are both total altitude (zong gu) heights; the height of the finial rings is the total altitude difference (zong gu jiao); the distance from the man's eye to the mast is the partial small base (fen xiao gou); the height of the mast above the man's eye and the height at which the base of the rings enters the mast are both partial small altitudes (fen xiao gu); the entry of the base of the rings into the mast is the small altitude difference (xiao gu jiao); the distance from the mast to the pagoda is the partial large base (fen da gou). The method multiplies and divides the partial large base, the distance from the mast to the pagoda, with the small base and altitude, obtaining the large altitude, and adds the small altitude to make the total altitude; therefore the number for the height of the pagoda tip is correct. It multiplies and divides the large base with the small base and the small altitude difference, obtaining the large altitude difference, and takes this directly as the total altitude difference; therefore the number for the height of the finial rings is short by one small altitude difference, 1 zhang 5 chi. The height of the pagoda body and the length of the central post are both obtained by addition and subtraction from this number; therefore the erroneous amounts are equal.\end{uniten}

\begin{unitzh}\unitlabel{shushu-030.20}數學九章卷四下\end{unitzh}
\begin{uniten}Shuxue jiuzhang (Mathematical Treatise in Nine Sections), juan 4, second part.\end{uniten}

\section*{卷5上}\addcontentsline{toc}{section}{卷5上}

\begin{unitzh}\unitlabel{shushu-031.1}欽定四庫全書\end{unitzh}
\begin{uniten}Qinding Siku quanshu (Imperially Commissioned Complete Library of the Four Treasuries)\end{uniten}

\begin{unitzh}\unitlabel{shushu-031.2}《數學九章》卷五上 宋 秦九韶 撰\end{unitzh}
\begin{uniten}Shuxue jiuzhang (Mathematical Treatise in Nine Sections), juan 5, part 1. Composed by Qin Jiushao of the Song.\end{uniten}

\begin{unitzh}\unitlabel{shushu-031.3}賦役\end{unitzh}
\begin{uniten}Taxes and Corvée (fuyi)\end{uniten}

\begin{unitzh}\unitlabel{shushu-031.4}按：賦役一章，即古《九章》中差分、粟米諸法。但從來算書設題之數，未有如此卷之多者。如首題答數至一百七十五條，毎條歩算之數至十餘位，而得數皆無不合，亦可謂能廣其用矣。且諸問皆足以明貴賤、廩税及交質、變易之同異，使公私各得其平，誠治民之要務也。但題語多晦，術、草中間有與所問不相應者，亦於各條下指出，俾觀者無惑焉。\end{unitzh}
\begin{uniten}[Siku editors: The chapter on Taxes and Corvée is the methods of Graded Distribution (chafen), Millet and Grain (sumi) and the like in the ancient Nine Chapters. But in the mathematical books hitherto, the numbers set in the problems have never been as many as in this juan. For example, the numbers of the answer to the first problem come to 175 items, the numbers reckoned for each item reach more than ten places, and the results obtained all agree without exception; it may indeed be said to have broadened their use. Moreover, the questions are all sufficient to make clear the sameness and difference of high and low prices, granary taxes, and exchanges of pledges and conversions, so that public and private each obtain what is fair; this is truly an essential matter in governing the people. But the wording of the problems is often obscure, and here and there in the methods and workings there are places that do not correspond to what is asked; these too are pointed out under the respective items, so that readers may not be perplexed.]\end{uniten}

\begin{unitzh}\unitlabel{shushu-031.5}復邑修賦\end{unitzh}
\begin{uniten}Restoring the County and Revising the Taxes\end{uniten}

\begin{unitzh}\unitlabel{shushu-031.6}問：有海坍縣地，今有復漲，歳乆鄉井再成，申諸剏邑。稱土排到六鄉，以附郭為甲，最逺為已，各有田九等，開具下項：\end{unitzh}
\begin{uniten}Question: There is the land of a county that collapsed into the sea; now it has risen again, after many years the villages have formed again, and a request was submitted to found a county. The land was measured and arranged into 6 townships (xiang), taking the one adjoining the city wall as Jia and the farthest as Ji [Translator's note: the text writes 已 for the stem 己, here and in 已鄉 below]; each has fields of 9 grades (deng), set out in the following items:\end{uniten}

\begin{unitzh}\unitlabel{shushu-031.7}甲鄉共計田十四萬一百九十三畝三角一十二步。上等：上田五千六百七十八畝一角四十八歩，中田四千八百九十二畝三十歩，下田六千六百二十一畝五十四歩。中等：上田八千二百二十五畝二十四歩，中田一萬三十五畝六歩，下田一萬六千五百三十畝。下等：上田二萬一千九十畝二十四歩，中田三萬二千六十畝三歩，下田三萬五千六十一畝三角三歩。\end{unitzh}
\begin{uniten}Township Jia: in total, fields 140193 mu (畝) 3 jiao (角) 12 bu (步). Upper grade: upper fields 5678 mu 1 jiao 48 bu, middle fields 4892 mu 30 bu, lower fields 6621 mu 54 bu. Middle grade: upper fields 8225 mu 24 bu, middle fields 10035 mu 6 bu, lower fields 16530 mu. Lower grade: upper fields 21090 mu 24 bu, middle fields 32060 mu 3 bu, lower fields 35061 mu 3 jiao 3 bu.\end{uniten}

\begin{unitzh}\unitlabel{shushu-031.8}乙鄉田共計八萬四千一十畝二角二歩。上等：上田六千七百八十九畝一角三十六歩，中田五千九百八十七畝二角，下田八千一十畝三角三歩。中等：上田七千五百四十一畝，中田九千一百二十一畝二角一十二歩，下田一萬九千六十歩。下等：上田一萬八千三十七畝一角六歩，中田九千四百五十六畝三角五十九歩，下田無。\end{unitzh}
\begin{uniten}The fields of Township Yi in total: 84010 mu 2 jiao 2 bu. Upper grade: upper fields 6789 mu 1 jiao 36 bu, middle fields 5987 mu 2 jiao, lower fields 8010 mu 3 jiao 3 bu. Middle grade: upper fields 7541 mu, middle fields 9121 mu 2 jiao 12 bu, lower fields 19060 bu [Translator's note: so in the text, with no figure in mu]. Lower grade: upper fields 18037 mu 1 jiao 6 bu, middle fields 9456 mu 3 jiao 59 bu, lower fields none.\end{uniten}

\begin{unitzh}\unitlabel{shushu-031.9}丙鄉田共計一十二萬九百三十五畝五十八歩五分。上等：上田四千八百六十八畝二角三歩，中田五千九百七十九畝三角六歩，下田六千八百八十八畝二角六歩。中等：上田七千九百八十四畝一角，中田一萬四千五十六畝一十二歩，下田二萬三千三百三十三畝一十二歩。下等：上田二萬七千七百五十五畝一十六歩五分，中田無，下田三萬七十畝三歩。\end{unitzh}
\begin{uniten}The fields of Township Bing in total: 120935 mu 58 bu 5 fen. Upper grade: upper fields 4868 mu 2 jiao 3 bu, middle fields 5979 mu 3 jiao 6 bu, lower fields 6888 mu 2 jiao 6 bu. Middle grade: upper fields 7984 mu 1 jiao, middle fields 14056 mu 12 bu, lower fields 23333 mu 12 bu. Lower grade: upper fields 27755 mu 16 bu 5 fen, middle fields none, lower fields 30070 mu 3 bu.\end{uniten}

\begin{unitzh}\unitlabel{shushu-031.10}丁郷田共計八萬九千六十六畝二歩三分。上等：上田一萬一千一百二畝一歩，中田九千八百七十六畝一角，下田八千七百六十五畝一角三十歩。中等：上田七千五百三十九畝三十四歩三分，中田無，下田一萬二千九百八十七畝四十二歩。下等：上田無，中田五千四百三十二畝一角六歩，下田三萬三千三百六十三畝三角九歩。\end{unitzh}
\begin{uniten}The fields of Township Ding in total: 89066 mu 2 bu 3 fen. Upper grade: upper fields 11102 mu 1 bu, middle fields 9876 mu 1 jiao, lower fields 8765 mu 1 jiao 30 bu. Middle grade: upper fields 7539 mu 34 bu 3 fen, middle fields none, lower fields 12987 mu 42 bu. Lower grade: upper fields none, middle fields 5432 mu 1 jiao 6 bu, lower fields 33363 mu 3 jiao 9 bu.\end{uniten}

\begin{unitzh}\unitlabel{shushu-031.11}戊鄉田共計二十萬四千四百七十四畝一角二十四歩四分。上等：上田二萬四千六百三十二畝三十九歩，中田一萬三千五百二十一畝二十七歩，下田九千九百八十八畝三角三歩。中等：上田八千八百七十七畝五十六歩四分，中田一萬一千三百三十三畝三角，下田二萬七千六十七畝。下等：上田一萬九千八百七十六畝三角六歩，中等田七萬九千一百三十五畝三角四十三歩，下田一萬四十一畝二角三十歩。\end{unitzh}
\begin{uniten}The fields of Township Wu in total: 204474 mu 1 jiao 24 bu 4 fen. Upper grade: upper fields 24632 mu 39 bu, middle fields 13521 mu 27 bu, lower fields 9988 mu 3 jiao 3 bu. Middle grade: upper fields 8877 mu 56 bu 4 fen, middle fields 11333 mu 3 jiao, lower fields 27067 mu. Lower grade: upper fields 19876 mu 3 jiao 6 bu, middle-grade fields 79135 mu 3 jiao 43 bu, lower fields 10041 mu 2 jiao 30 bu.\end{uniten}

\begin{unitzh}\unitlabel{shushu-031.12}已鄉田共計一十五萬八千四百六十畝三角一十八歩二分。上等：上田無，中田七千七百八十八畝三角五十一歩，下田無。中等：上田九千九百九十九畝二角三歩，中田一萬八百三十六畝五十六歩，下田無。下等：上田三萬二千八十九畝一角四十五歩六分，中田四萬三千六百七十八畝二角五十七歩，下田五萬四千六十七畝三角四十五歩六分。\end{unitzh}
\begin{uniten}The fields of Township Ji in total: 158460 mu 3 jiao 18 bu 2 fen. Upper grade: upper fields none, middle fields 7788 mu 3 jiao 51 bu, lower fields none. Middle grade: upper fields 9999 mu 2 jiao 3 bu, middle fields 10836 mu 56 bu, lower fields none. Lower grade: upper fields 32089 mu 1 jiao 45 bu 6 fen, middle fields 43678 mu 2 jiao 57 bu, lower fields 54067 mu 3 jiao 45 bu 6 fen.\end{uniten}

\begin{unitzh}\unitlabel{shushu-031.13}照得昨來本縣原科苗米一十萬三千五百六十七石八斗四升四合三勺，和買一萬三千四百九十八匹一丈七尺三寸七分六釐，夏税九千八百七十六匹三丈二尺六寸五分八釐。其六鄉田係三色：甲為上，乙、丙為次，丁、戊、己又為次。先令官物為三差，使上比中、中比下，皆十分外差一；次令各鄉九等，皆於十分内差一，抛科用合租額。其乙鄉田最肥，次丁，次甲，次丙，次己，次戊。欲知三色九等毎畝等則及共科數，各鄉幾何？\end{unitzh}
\begin{uniten}It is ascertained (zhaode) that formerly this county was originally levied: seedling rice (miaomi) 103567 dan (石) 8 dou 4 sheng 4 ge 3 shao; harmonious purchase (hemai) 13498 pi (匹) 1 zhang 7 chi 3 cun 7 fen 6 li (釐); summer tax (xiashui) 9876 pi 3 zhang 2 chi 6 cun 5 fen 8 li. The fields of the six townships belong to three classes (se): Jia is the upper, Yi and Bing are the next, and Ding, Wu and Ji are the next again. First let the government goods (guanwu) be made into three differences (cha), so that upper compared with middle and middle compared with lower all differ by 1 outside 10 parts; next let the nine grades of each township all differ by 1 within 10 parts; apportion the levy (paoke) so as to agree with the tax quota (zu'e). The fields of Township Yi are the most fertile, next Ding, next Jia, next Bing, next Ji, next Wu. It is desired to know the grade rate (dengze) per mu of the three classes and nine grades, and the amounts of the total levy (gongke): how much for each township?\end{uniten}

\begin{unitzh}\unitlabel{shushu-031.14}按：題内言「田有三色，甲為上」云云，又言「乙鄉田最肥」云云，語若相左。葢前言三色者，六鄉共科之等；後言田肥者，每畝所出之異也。若易「田係三色」為「科有三差」，則語意明矣。\end{unitzh}
\begin{uniten}[Siku editors: Within the problem it says 'the fields have three classes, Jia is the upper' and so on, and again it says 'the fields of Township Yi are the most fertile' and so on; the statements seem to contradict each other. For the three classes spoken of first are the grades of the total levy of the six townships, while the fertility of the fields spoken of afterwards is the difference in what each mu yields. If 'the fields belong to three classes' were changed to 'the levy has three differences', the meaning of the wording would be clear.]\end{uniten}

\begin{unitzh}\unitlabel{shushu-031.15}答曰：甲鄉上等上田：苗三斗二升三合二勺，和一尺六寸九分，税一尺二寸三分；中田：苗二斗九升九勺，和一尺五寸二分，税一尺一寸一分；下田：苗二斗五升八合六勺，和一尺三寸五分，税九寸九分。中等上田：苗二斗二升六合二勺，和一尺一寸八分，税八寸六分；中田：苗一斗九升三合九勺，和一尺一分，税七寸四分；下田：苗一斗六升一合六勺，和八寸四分，税六寸二分。下等上田：苗一斗二升九合三勺，和六寸七分，税四寸九分；中田：苗九升六合九勺，和五寸一分，税三寸七分；下田：苗六升四合六勺，和三寸四分，税二寸五分。\end{unitzh}
\begin{uniten}Answer: Township Jia, upper grade, upper fields: miao 3 dou 2 sheng 3 ge 2 shao, he 1 chi 6 cun 9 fen, shui 1 chi 2 cun 3 fen [Translator's note: miao, he and shui abbreviate the seedling rice, the harmonious purchase and the summer tax]; middle fields: miao 2 dou 9 sheng 9 shao, he 1 chi 5 cun 2 fen, shui 1 chi 1 cun 1 fen; lower fields: miao 2 dou 5 sheng 8 ge 6 shao, he 1 chi 3 cun 5 fen, shui 9 cun 9 fen. Middle grade, upper fields: miao 2 dou 2 sheng 6 ge 2 shao, he 1 chi 1 cun 8 fen, shui 8 cun 6 fen; middle fields: miao 1 dou 9 sheng 3 ge 9 shao, he 1 chi 1 fen, shui 7 cun 4 fen; lower fields: miao 1 dou 6 sheng 1 ge 6 shao, he 8 cun 4 fen, shui 6 cun 2 fen. Lower grade, upper fields: miao 1 dou 2 sheng 9 ge 3 shao, he 6 cun 7 fen, shui 4 cun 9 fen; middle fields: miao 9 sheng 6 ge 9 shao, he 5 cun 1 fen, shui 3 cun 7 fen; lower fields: miao 6 sheng 4 ge 6 shao, he 3 cun 4 fen, shui 2 cun 5 fen.\end{uniten}

\begin{unitzh}\unitlabel{shushu-031.16}乙鄉上等田：苗三斗六升三合三勺，和一尺八寸九分，税一尺三寸九分；中田：苗三斗二升六合九勺，和一尺七寸，税一尺二寸五分；下田：苗二斗九升六勺，和一尺五寸二分，税一尺一寸一分。中等上田：苗二斗五升四合三勺，和一尺三寸三分，税九寸七分；中田：苗二斗一升八合，和一尺一寸四分，税八寸三分；下田：苗一斗八升一合六勺，和九寸五分，税六寸九分。下等上田：苗一斗四升五合三勺，和七寸六分，税五寸五分；中田：苗一斗九合，和五寸七分，税四寸二分；下田：苗七升二合七勺，和三寸八分，税二寸八分。\end{unitzh}
\begin{uniten}Township Yi, upper-grade fields: miao 3 dou 6 sheng 3 ge 3 shao, he 1 chi 8 cun 9 fen, shui 1 chi 3 cun 9 fen; middle fields: miao 3 dou 2 sheng 6 ge 9 shao, he 1 chi 7 cun, shui 1 chi 2 cun 5 fen; lower fields: miao 2 dou 9 sheng 6 shao, he 1 chi 5 cun 2 fen, shui 1 chi 1 cun 1 fen. Middle grade, upper fields: miao 2 dou 5 sheng 4 ge 3 shao, he 1 chi 3 cun 3 fen, shui 9 cun 7 fen; middle fields: miao 2 dou 1 sheng 8 ge, he 1 chi 1 cun 4 fen, shui 8 cun 3 fen; lower fields: miao 1 dou 8 sheng 1 ge 6 shao, he 9 cun 5 fen, shui 6 cun 9 fen. Lower grade, upper fields: miao 1 dou 4 sheng 5 ge 3 shao, he 7 cun 6 fen, shui 5 cun 5 fen; middle fields: miao 1 dou 9 ge, he 5 cun 7 fen, shui 4 cun 2 fen; lower fields: miao 7 sheng 2 ge 7 shao, he 3 cun 8 fen, shui 2 cun 8 fen.\end{uniten}

\begin{unitzh}\unitlabel{shushu-031.17}丙鄉上等上田：苗三斗三合五勺，和一尺五寸八分，税一尺一寸六分；中田：苗二斗七升三合一勺，和一尺四寸二分，税一尺四分；下田：苗二斗四升二合八勺，和一尺二寸七分，税九寸三分。中等上田：苗二斗一升二合四勺，和一尺一寸一分，税八寸一分；中田：苗一斗八升二合一勺，和九寸五分，税六寸九分；下田：苗一斗五升一合七勺，和七寸九分，税五寸八分。下等上田：苗一斗二升一合四勺，和六寸三分，税四寸六分；中田：苗九升一合，和四寸七分，税三寸五分；下田：苗六升七勺，和三寸二分，税二寸三分。\end{unitzh}
\begin{uniten}Township Bing, upper grade, upper fields: miao 3 dou 3 ge 5 shao, he 1 chi 5 cun 8 fen, shui 1 chi 1 cun 6 fen; middle fields: miao 2 dou 7 sheng 3 ge 1 shao, he 1 chi 4 cun 2 fen, shui 1 chi 4 fen; lower fields: miao 2 dou 4 sheng 2 ge 8 shao, he 1 chi 2 cun 7 fen, shui 9 cun 3 fen. Middle grade, upper fields: miao 2 dou 1 sheng 2 ge 4 shao, he 1 chi 1 cun 1 fen, shui 8 cun 1 fen; middle fields: miao 1 dou 8 sheng 2 ge 1 shao, he 9 cun 5 fen, shui 6 cun 9 fen; lower fields: miao 1 dou 5 sheng 1 ge 7 shao, he 7 cun 9 fen, shui 5 cun 8 fen. Lower grade, upper fields: miao 1 dou 2 sheng 1 ge 4 shao, he 6 cun 3 fen, shui 4 cun 6 fen; middle fields: miao 9 sheng 1 ge, he 4 cun 7 fen, shui 3 cun 5 fen; lower fields: miao 6 sheng 7 shao, he 3 cun 2 fen, shui 2 cun 3 fen.\end{uniten}

\begin{unitzh}\unitlabel{shushu-031.18}丁鄉上等上田：苗三斗四升三合二勺，和一尺七寸九分，税一尺三寸一分；中田：苗三斗八合九勺，和一尺六寸一分，税一尺一寸八分；下田：苗二斗七升四合六勺，和一尺四寸三分，税一尺五分。中等上田：苗二斗四升二勺，和一尺二寸五分，税九寸二分；中田：苗二斗五合九勺，和一尺七分，税七寸九分；下田：苗一斗七升一合六勺，和八寸九分，税六寸五分。下等上田：苗一斗三升七合三勺，和七寸二分，税五寸二分；中田：苗一斗三合，和五寸四分，税三寸九分；下田：苗六升八合六勺，和三寸六分，税二寸六分。\end{unitzh}
\begin{uniten}Township Ding, upper grade, upper fields: miao 3 dou 4 sheng 3 ge 2 shao, he 1 chi 7 cun 9 fen, shui 1 chi 3 cun 1 fen; middle fields: miao 3 dou 8 ge 9 shao, he 1 chi 6 cun 1 fen, shui 1 chi 1 cun 8 fen; lower fields: miao 2 dou 7 sheng 4 ge 6 shao, he 1 chi 4 cun 3 fen, shui 1 chi 5 fen. Middle grade, upper fields: miao 2 dou 4 sheng 2 shao, he 1 chi 2 cun 5 fen, shui 9 cun 2 fen; middle fields: miao 2 dou 5 ge 9 shao, he 1 chi 7 fen, shui 7 cun 9 fen; lower fields: miao 1 dou 7 sheng 1 ge 6 shao, he 8 cun 9 fen, shui 6 cun 5 fen. Lower grade, upper fields: miao 1 dou 3 sheng 7 ge 3 shao, he 7 cun 2 fen, shui 5 cun 2 fen; middle fields: miao 1 dou 3 ge, he 5 cun 4 fen, shui 3 cun 9 fen; lower fields: miao 6 sheng 8 ge 6 shao, he 3 cun 6 fen, shui 2 cun 6 fen.\end{uniten}

\begin{unitzh}\unitlabel{shushu-031.19}戊鄉上等上田：苗一斗五升三合八勺，和八寸，税五寸九分；中田：苗一斗三升八合四勺，和七寸二分，税五寸三分；下田：苗一斗二升三合一勺，和六寸四分，税四寸七分。中等上田：苗一斗七合七勺，和五寸六分，税四寸一分；中田：苗九升二合三勺，和四寸八分，税三寸五分；下田：苗七升六合九勺，和四寸，税二寸九分。下等上田：苗六升一合五勺，和三寸二分，税二寸三分；中田：苗四升六合一勺，和二寸四分，税一寸八分；下田：苗三升八勺，和一寸六分，税一寸二分。\end{unitzh}
\begin{uniten}Township Wu, upper grade, upper fields: miao 1 dou 5 sheng 3 ge 8 shao, he 8 cun, shui 5 cun 9 fen; middle fields: miao 1 dou 3 sheng 8 ge 4 shao, he 7 cun 2 fen, shui 5 cun 3 fen; lower fields: miao 1 dou 2 sheng 3 ge 1 shao, he 6 cun 4 fen, shui 4 cun 7 fen. Middle grade, upper fields: miao 1 dou 7 ge 7 shao, he 5 cun 6 fen, shui 4 cun 1 fen; middle fields: miao 9 sheng 2 ge 3 shao, he 4 cun 8 fen, shui 3 cun 5 fen; lower fields: miao 7 sheng 6 ge 9 shao, he 4 cun, shui 2 cun 9 fen. Lower grade, upper fields: miao 6 sheng 1 ge 5 shao, he 3 cun 2 fen, shui 2 cun 3 fen; middle fields: miao 4 sheng 6 ge 1 shao, he 2 cun 4 fen, shui 1 cun 8 fen; lower fields: miao 3 sheng 8 shao, he 1 cun 6 fen, shui 1 cun 2 fen.\end{uniten}

\begin{unitzh}\unitlabel{shushu-031.20}已鄉上等上田：苗二斗八升二合二勺，和一尺四寸七分，税一尺八分；中田：苗二斗五升三合九勺，和一尺三寸二分，税九寸七分；下田：苗二斗二升五合七勺，和一尺一寸八分，税八寸六分。中等上田：苗一斗九升七合五勺，和一尺三分，税七寸五分；中田：苗一斗六升九合三勺，和八寸八分，税六寸五分；下田：苗一斗四升一合一勺，和七寸四分，税五寸四分。\end{unitzh}
\begin{uniten}Township Ji, upper grade, upper fields: miao 2 dou 8 sheng 2 ge 2 shao, he 1 chi 4 cun 7 fen, shui 1 chi 8 fen; middle fields: miao 2 dou 5 sheng 3 ge 9 shao, he 1 chi 3 cun 2 fen, shui 9 cun 7 fen; lower fields: miao 2 dou 2 sheng 5 ge 7 shao, he 1 chi 1 cun 8 fen, shui 8 cun 6 fen. Middle grade, upper fields: miao 1 dou 9 sheng 7 ge 5 shao, he 1 chi 3 fen, shui 7 cun 5 fen; middle fields: miao 1 dou 6 sheng 9 ge 3 shao, he 8 cun 8 fen, shui 6 cun 5 fen; lower fields: miao 1 dou 4 sheng 1 ge 1 shao, he 7 cun 4 fen, shui 5 cun 4 fen. [continues]\end{uniten}

\begin{unitzh}\unitlabel{shushu-032.1}下等上田：苗一斗一升二合九勺，和五寸九分，税四寸三分；中田：苗八升四合六勺，和四寸四分，税三寸二分；下田：苗五升六合四勺，和二寸九分，税二寸二分。\end{unitzh}
\begin{uniten}[continued] Lower grade, upper field: seedling (miao) 1 dou (斗) 1 sheng (升) 2 ge (合) 9 shao (勺); harmonious purchase (he) 5 cun (寸) 9 fen (分); tax 4 cun 3 fen. Middle field: seedling 8 sheng 4 ge 6 shao; harmonious purchase 4 cun 4 fen; tax 3 cun 2 fen. Lower field: seedling 5 sheng 6 ge 4 shao; harmonious purchase 2 cun 9 fen; tax 2 cun 2 fen.\end{uniten}

\begin{unitzh}\unitlabel{shushu-032.2}甲鄉：苗米一萬九千五百五十石二斗四升八合三勺，和買一萬一百九十二丈二尺六寸五分六釐，夏税七千四百五十七丈六尺八寸九分八釐。乙鄉：苗米一萬七千七百七十二石九斗五升三合，和買九千二百六十五丈六尺九寸六分，夏税六千七百七十九丈七尺一寸八分。丙鄉：苗米一萬七千七百七十二石九斗五升三合，和買九千二百六十五丈六尺九寸六分，夏税六千七百七十九丈七尺一寸八分。丁鄉：苗米一萬六千一百五十七石二斗三升，和買八千四百二十三丈三尺六寸，夏税六千一百六十三丈三尺八寸。戊鄉：苗米一萬六千一百五十七石二斗三升，和買八千四百二十三丈三尺六寸，夏税六千一百六十三丈三尺八寸。己鄉：苗米一萬六千一百五十七石二斗三升，和買八千四百二十三丈三尺六寸，夏税六千一百六十三丈三尺八寸。\end{unitzh}
\begin{uniten}Jia township (xiang): seedling rice (miaomi) 19550 dan (石) 2 dou 4 sheng 8 ge 3 shao; harmonious purchase (hemai) 10192 zhang (丈) 2 chi (尺) 6 cun 5 fen 6 li (釐); summer tax (xiashui) 7457 zhang 6 chi 8 cun 9 fen 8 li. Yi township: seedling rice 17772 dan 9 dou 5 sheng 3 ge; harmonious purchase 9265 zhang 6 chi 9 cun 6 fen; summer tax 6779 zhang 7 chi 1 cun 8 fen. Bing township: seedling rice 17772 dan 9 dou 5 sheng 3 ge; harmonious purchase 9265 zhang 6 chi 9 cun 6 fen; summer tax 6779 zhang 7 chi 1 cun 8 fen. Ding township: seedling rice 16157 dan 2 dou 3 sheng; harmonious purchase 8423 zhang 3 chi 6 cun; summer tax 6163 zhang 3 chi 8 cun. Wu township: seedling rice 16157 dan 2 dou 3 sheng; harmonious purchase 8423 zhang 3 chi 6 cun; summer tax 6163 zhang 3 chi 8 cun. Ji township: seedling rice 16157 dan 2 dou 3 sheng; harmonious purchase 8423 zhang 3 chi 6 cun; summer tax 6163 zhang 3 chi 8 cun.\end{uniten}

\begin{unitzh}\unitlabel{shushu-032.3}術：以衰分求之。先列本縣色位，自下錐行列之；又以鄉數對列而乗之，副併為法，以除官物數，得一分之率；以率數乗未併者，各得諸鄉之數。次列各鄉等位，自上等置十分，每以内分錐行九折之，至九等止；又各以畝歩乗之，副併為鄉法，以除諸各鄉所得官物數，所得為一分之率；以乗未併者，各得毎畝税色。\end{unitzh}
\begin{uniten}Method: Find it by graded distribution (cuifen). First lay out the category positions (sewei) of this county (xian), arranging them from the bottom in an awl-shaped column (zhuixing); then lay out the numbers of townships opposite them and multiply; add them in an auxiliary position (fu) to make the divisor (fa), and with it divide the amount of the official goods (guanwu), giving the rate (lü) for one part; multiply the ones not yet added by the rate, and each gives the amount for each of the townships. Next lay out the grade positions (dengwei) of each township: starting from the upper grade, set 10 parts (fen); each time, by the inner parts (neifen), reduce to nine-tenths (jiuzhe) in an awl-shaped column, stopping at the ninth grade; then multiply each by the mu (畝) and bu (歩), add them in an auxiliary position to make the township divisor (xiangfa), and with it divide the amount of official goods obtained by each of the townships; what is obtained is the rate for one part; multiply the ones not yet added by it, and each gives the tax categories (shuise) per mu.\end{uniten}

\begin{unitzh}\unitlabel{shushu-032.4}草曰：列本縣色位三，自下色例十分，中十一分，上比中身外加一；上得一百二十一，中得一百一十，下得一百，為上中下三率，列之右行。乃列甲一對上率，乙、丙共二對中率，丁、戊、己共三對下率，列左行。乃以左行列數各相對乗右行率數，其上得一百二十一，其十得二百二十，其下得三百。乃副置而併之，得六百四十一，為法。置元科苗米一萬三千五百六十七石八斗四升四合三勺，為寔；以法除之，得一百六十一石五斗七升二合三勺，為一分之率。以未併下率一百乗率，得一萬六千一百五十七石二斗三升，為丁、戊、己三鄉各科數；以於身下加一，得一萬七千七百七十二石九斗五升三合，為乙、丙二鄉各科數；又於身下加一，得一萬九千五百五十石二斗四升八合三勺，為甲合科數。\end{unitzh}
\begin{uniten}Working: Lay out the 3 category positions of this county. From the bottom, the category ratio (seli) is 10 parts; the middle, 11 parts; the upper, compared with the middle, add one outside the body (shenwai jia yi) [Translator's note: i.e., multiply by 1.1]; the upper gives 121, the middle gives 110, the lower gives 100; as the three rates of upper, middle and lower, lay them out in the right column. Then lay out Jia, 1, opposite the upper rate; Yi and Bing, together 2, opposite the middle rate; Ding, Wu and Ji, together 3, opposite the lower rate; lay them out in the left column. Then multiply with the numbers laid out in the left column each the rate number opposite it in the right column: the upper gives 121, the "ten" [Translator's note: 十 "ten" evidently stands for 中 "middle"] gives 220, the lower gives 300. Then set them in an auxiliary position and add them, giving 641, as the divisor. Set the originally assessed (yuanke) seedling rice, 13567 dan 8 dou 4 sheng 4 ge 3 shao, as the dividend (shi); divide it by the divisor, giving 161 dan 5 dou 7 sheng 2 ge 3 shao as the rate for one part. Multiply the rate by the not-yet-added lower rate 100, giving 16157 dan 2 dou 3 sheng as the assessed amount of each of the three townships Ding, Wu and Ji; with this, add one below the body (shenxia jia yi), giving 17772 dan 9 dou 5 sheng 3 ge as the assessed amount of each of the two townships Yi and Bing; again add one below the body, giving 19550 dan 2 dou 4 sheng 8 ge 3 shao as the amount to be assessed on Jia.\end{uniten}

\begin{unitzh}\unitlabel{shushu-032.5}求和買，亦用六百四十一為法。置元科和買一萬三千四百九十八匹，以匹法四丈通之，得五萬三千九百九十二丈；内零一丈七尺三寸七分六釐，得五萬三千九百九十三丈七尺三寸七分六釐，為寔。以法除之，得八十四丈二尺三寸三分六釐，為一分之率。亦以未併下率一百乗之，得八千四百二十三丈三尺六寸，為丁、戊、己三鄉各科數；次於身下加一，得九千二百六十五丈六尺九寸六分，為乙、丙二鄉各科數；又於身下加一，得一萬一百九十二丈二尺六寸五分六釐，為甲鄉合科數。\end{unitzh}
\begin{uniten}To find the harmonious purchase, also use 641 as the divisor. Set the originally assessed harmonious purchase, 13498 pi (匹); convert it (tong) by the bolt-rule (pifa), 4 zhang, giving 53992 zhang; add in the odd part (ling), 1 zhang 7 chi 3 cun 7 fen 6 li, giving 53993 zhang 7 chi 3 cun 7 fen 6 li as the dividend. Divide it by the divisor, giving 84 zhang 2 chi 3 cun 3 fen 6 li as the rate for one part. Also multiply it by the not-yet-added lower rate 100, giving 8423 zhang 3 chi 6 cun as the assessed amount of each of the three townships Ding, Wu and Ji; next add one below the body, giving 9265 zhang 6 chi 9 cun 6 fen as the assessed amount of each of the two townships Yi and Bing; again add one below the body, giving 10192 zhang 2 chi 6 cun 5 fen 6 li as the amount to be assessed on Jia township.\end{uniten}

\begin{unitzh}\unitlabel{shushu-032.6}求夏税，亦置六百四十一為法。置元科九千八百七十六匹，以匹法四丈通之，得三萬九千五百四丈；内零三丈二尺六寸五分八釐，得三萬九千五百七丈二尺六寸五分八釐，為寔。以法除之，得六十一丈六尺三寸三分八釐，為一分率。亦以未併下率一百乗之，得六千一百六十三丈三尺八寸，為丁、戊、己三鄉各科數；次於身下加一，得六千七百七十九丈七尺一寸八分，為乙、丙二鄉各科數；又於身下加一，得七千四百五十七丈六尺八寸九分八釐，為甲鄉數。\end{unitzh}
\begin{uniten}To find the summer tax, also set 641 as the divisor. Set the original assessment, 9876 pi; convert it by the bolt-rule, 4 zhang, giving 39504 zhang; add in the odd part, 3 zhang 2 chi 6 cun 5 fen 8 li, giving 39507 zhang 2 chi 6 cun 5 fen 8 li as the dividend. Divide it by the divisor, giving 61 zhang 6 chi 3 cun 3 fen 8 li as the one-part rate. Also multiply it by the not-yet-added lower rate 100, giving 6163 zhang 3 chi 8 cun as the assessed amount of each of the three townships Ding, Wu and Ji; next add one below the body, giving 6779 zhang 7 chi 1 cun 8 fen as the assessed amount of each of the two townships Yi and Bing; again add one below the body, giving 7457 zhang 6 chi 8 cun 9 fen 8 li as the amount for Jia township.\end{uniten}

\begin{unitzh}\unitlabel{shushu-032.7}次列九等：上以十，次九、八、七、六、五、四、三、二，各對乗六鄉九等田畝。其田畝下角歩，以畝法除之，得分、釐、毫、絲、忽，於畝上對乗之，各得率。\end{unitzh}
\begin{uniten}Next lay out the nine grades: the upper with 10, the next 9, 8, 7, 6, 5, 4, 3, 2; with each, multiply correspondingly the field mu of the nine grades of the six townships. As for the jiao (角) and bu below the mu of the fields, divide them by the mu-rule (mufa), giving fen, li, hao (毫), si (絲), hu (忽); with these on the mu, multiply correspondingly, and each gives its rate.\end{uniten}

\begin{unitzh}\unitlabel{shushu-032.8}按：題内云各鄉九等皆於十分内差一，則是逓次九折也；術内亦云「列各鄉等位，自上等置十分，每以内分錐行九折之，至九等止」，則其法當為逓次九折益明矣。乃草中則云「上以十，次九、八、七、六、五、四、三、二」，則是逓次減一，而非逓次九折矣。以此得答數，既與題問未合，而術與草亦不相應也。\end{unitzh}
\begin{uniten}Note: In the problem it says that the nine grades of each township all differ by one within ten parts; this, then, is successive reduction to nine-tenths. In the method it also says, "lay out the grade positions of each township: starting from the upper grade, set 10 parts; each time, by the inner parts, reduce to nine-tenths in an awl-shaped column, stopping at the ninth grade"; so it is all the clearer that the procedure should be successive reduction to nine-tenths. Yet in the working it says, "the upper with 10, the next 9, 8, 7, 6, 5, 4, 3, 2"; this is successive subtraction of one, and not successive reduction to nine-tenths. The answer figures obtained in this way not only fail to agree with the question of the problem; the method and the working also do not correspond to each other.\end{uniten}

\begin{unitzh}\unitlabel{shushu-032.9}甲：上等上田率五萬六千七百八十四分五釐，中田率四萬四千二十九分一釐二毫五絲，下田率五萬二千九百六十九分八釐；中等上田率五萬七千五百七十五分七釐，中田率六萬二百一十分一釐五毫，下田率八萬二千六百五十分；下等上田率八萬四千三百六十分四釐，中田率九萬六千一百八十分三毫七絲五忽，下田率七萬一百二十三分五釐二毫五絲。\end{unitzh}
\begin{uniten}Jia: upper grade, upper field: rate 56784 fen 5 li; middle field: rate 44029 fen 1 li 2 hao 5 si; lower field: rate 52969 fen 8 li. Middle grade, upper field: rate 57575 fen 7 li; middle field: rate 60210 fen 1 li 5 hao; lower field: rate 82650 fen. Lower grade, upper field: rate 84360 fen 4 li; middle field: rate 96180 fen 3 hao 7 si 5 hu; lower field: rate 70123 fen 5 li 2 hao 5 si.\end{uniten}

\begin{unitzh}\unitlabel{shushu-032.10}乙：上等上田率六萬七千八百九十四分，中田率五萬三千八百八十七分五釐，下田率六萬四千八十六分一釐；中等上田率五萬二千七百八十七分，中田率五萬四千七百二十九分三釐，下田率九萬五千三百三十分一釐二毫五絲；下等上田率七萬二千一百四十九分一釐，中田率二萬八千三百七十分九釐八毫八絲，下田無。\end{unitzh}
\begin{uniten}Yi: upper grade, upper field: rate 67894 fen; middle field: rate 53887 fen 5 li; lower field: rate 64086 fen 1 li. Middle grade, upper field: rate 52787 fen; middle field: rate 54729 fen 3 li; lower field: rate 95330 fen 1 li 2 hao 5 si. Lower grade, upper field: rate 72149 fen 1 li; middle field: rate 28370 fen 9 li 8 hao 8 si; lower field: none.\end{uniten}

\begin{unitzh}\unitlabel{shushu-032.11}丙：上等上田率四萬八千六百八十五分一釐二毫五絲，中田率五萬三千八百一十七分九釐七毫五絲，下田率五萬五千一百八分二釐；中等上田率五萬五千八百八十九分七釐五毫，中田率八萬四千三百三十六分三釐，下田率一十一萬六千六百六十五分二釐五毫；下等上田率一十一萬一千二十分二釐七毫五絲，中田無，下田率六萬一百四十分二毫五絲。\end{unitzh}
\begin{uniten}Bing: upper grade, upper field: rate 48685 fen 1 li 2 hao 5 si; middle field: rate 53817 fen 9 li 7 hao 5 si; lower field: rate 55108 fen 2 li. Middle grade, upper field: rate 55889 fen 7 li 5 hao; middle field: rate 84336 fen 3 li; lower field: rate 116665 fen 2 li 5 hao. Lower grade, upper field: rate 111020 fen 2 li 7 hao 5 si; middle field: none; lower field: rate 60140 fen 2 hao 5 si.\end{uniten}

\begin{unitzh}\unitlabel{shushu-032.12}丁：上等上田率一十一萬一千二十分五毫{\zhsmall 按：五毫數誤，應四釐一毫又三分毫之二。}，中田率八萬八千八百八十六分二釐五毫，下田率七萬一百二十三分；中等上田率五萬二千七百七十四分{\zhsmall 按：尚有四忽餘。}，中田無，下田率六萬四千九百三十五分八釐七毫五絲；下等上田無，中田率一萬六千二百九十六分八釐二毫五絲，下田率六萬六千七百二十七分五釐七毫五絲。\end{unitzh}
\begin{uniten}Ding: upper grade, upper field: rate 111020 fen 5 hao [Siku editors: the figure 5 hao is wrong; it should be 4 li 1 hao and 2/3 hao]; middle field: rate 88886 fen 2 li 5 hao; lower field: rate 70123 fen. Middle grade, upper field: rate 52774 fen [Siku editors: there is still more than 4 hu]; middle field: none; lower field: rate 64935 fen 8 li 7 hao 5 si. Lower grade, upper field: none; middle field: rate 16296 fen 8 li 2 hao 5 si; lower field: rate 66727 fen 5 li 7 hao 5 si.\end{uniten}

\begin{unitzh}\unitlabel{shushu-032.13}戊：上等上田率二十四萬六千三百二十一分六釐二毫五絲，中田率一十二萬一千六百九十分一毫二絲五忽，下田率七萬九千九百一十分一釐；中等上田率六萬二千一百四十分六釐四毫五絲，中田率六萬八千二分五釐，下田率一十三萬五千三百三十五分；下等上田率七萬九千五百七分一釐，中田率二十三萬七千四百七分七釐八毫七絲五忽，下田率二萬八十三分二釐五毫。\end{unitzh}
\begin{uniten}Wu: upper grade, upper field: rate 246321 fen 6 li 2 hao 5 si; middle field: rate 121690 fen 1 hao 2 si 5 hu; lower field: rate 79910 fen 1 li. Middle grade, upper field: rate 62140 fen 6 li 4 hao 5 si; middle field: rate 68002 fen 5 li; lower field: rate 135335 fen. Lower grade, upper field: rate 79507 fen 1 li; middle field: rate 237407 fen 7 li 8 hao 7 si 5 hu; lower field: rate 20083 fen 2 li 5 hao.\end{uniten}

\begin{unitzh}\unitlabel{shushu-032.14}已：上等上田無，中田率七萬一百分六釐六毫二絲五忽，下田無；中等上田率六萬九千九百九十六分五釐八毫七絲五忽，中田率六萬五千一十七分四釐一毫{\zhsmall 按：四釐整。}，下田無；下等上田率一十二萬八千三百五十七分七釐六毫，中田率一十三萬一千三十六分二釐一毫二絲五忽，下田率一十萬八千一百三十五分八釐八毫。\end{unitzh}
\begin{uniten}Ji [Translator's note: written 已 for 己]: upper grade, upper field: none; middle field: rate 70100 fen 6 li 6 hao 2 si 5 hu; lower field: none. Middle grade, upper field: rate 69996 fen 5 li 8 hao 7 si 5 hu; middle field: rate 65017 fen 4 li 1 hao [Siku editors: exactly 4 li]; lower field: none. Lower grade, upper field: rate 128357 fen 7 li 6 hao; middle field: rate 131036 fen 2 li 1 hao 2 si 5 hu; lower field: rate 108135 fen 8 li 8 hao.\end{uniten}

\begin{unitzh}\unitlabel{shushu-032.15}併六鄉之九率，為九鄉之法：甲法六十萬四千八百八十三分二釐三毫七絲五忽；乙法四十八萬九千二百三十四分一釐一毫三絲；丙法五十八萬五千六百六十二分九釐；丁法四十七萬七百六十三分五釐七毫五絲{\zhsmall 按：丁法差三釐餘，一等尾數差故也。}；戊法一百五萬三百九十八分二毫；己法五十七萬二千六百四十四分五釐一毫二絲五忽{\zhsmall 按：多一毫，五等分多一毫故也。}。以六鄉法各除諸鄉官物，得一分之率{\zhsmall 米至圭，帛至忽止；半巳上𭣣，巳下棄。}。\end{unitzh}
\begin{uniten}Add together the nine rates of the six townships to make the divisors of the nine townships [Translator's note: "nine townships" (九鄉) is evidently an error; there are six]: Jia's divisor 604883 fen 2 li 3 hao 7 si 5 hu; Yi's divisor 489234 fen 1 li 1 hao 3 si; Bing's divisor 585662 fen 9 li; Ding's divisor 470763 fen 5 li 7 hao 5 si [Siku editors: Ding's divisor is off by more than 3 li; this is because the final digits of the first grade are off]; Wu's divisor 1050398 fen 2 hao; Ji's divisor 572644 fen 5 li 1 hao 2 si 5 hu [Siku editors: 1 hao too much; this is because the parts of the fifth grade are 1 hao too many]. With the divisors of the six townships divide the official goods of each of the townships, giving the rate for one part [Note: rice down to the gui (圭), silk down to the hu, and stop; half and above is taken in, below that is discarded].\end{uniten}

\begin{unitzh}\unitlabel{shushu-032.16}甲鄉：苗米三升二合三勺二抄一撮，和買一寸六分八釐五毫四忽{\zhsmall 按：數多四忽。}，夏税一寸二分三釐二毫八絲{\zhsmall 按：係九絲。}。乙鄉：苗米三升六合三勺二抄八撮，和買一寸八分九釐三毫九絲二忽，夏税一寸三分八釐五毫八絲{\zhsmall 按：應七絲八忽。}。\end{unitzh}
\begin{uniten}Jia township: seedling rice 3 sheng 2 ge 3 shao 2 chao (抄) 1 cuo (撮); harmonious purchase 1 cun 6 fen 8 li 5 hao 4 hu [Siku editors: the figure is 4 hu too much]; summer tax 1 cun 2 fen 3 li 2 hao 8 si [Siku editors: it is 9 si]. Yi township: seedling rice 3 sheng 6 ge 3 shao 2 chao 8 cuo; harmonious purchase 1 cun 8 fen 9 li 3 hao 9 si 2 hu; summer tax 1 cun 3 fen 8 li 5 hao 8 si [Siku editors: it should be 7 si 8 hu]. [continues]\end{uniten}

\begin{unitzh}\unitlabel{shushu-033.1}丙鄉：苗米三升三勺四抄八撮{\zhsmall 按：應六撮七圭餘}，和買一寸五分八釐二毫一絲，夏税一寸一分五釐七毫五絲{\zhsmall 按：應六絲}。\end{unitzh}
\begin{uniten}Township Bing: seedling rice (miaomi) 3 sheng (升) 3 shao (勺) 4 chao (抄) 8 cuo (撮) [Siku editors: It should be 6 cuo 7-odd gui (圭).]; harmonized purchase (hemai) 1 cun (寸) 5 fen (分) 8 li (釐) 2 hao (毫) 1 si (絲); summer tax (xiashui) 1 cun 1 fen 5 li 7 hao 5 si [Siku editors: It should be 6 si.].\end{uniten}

\begin{unitzh}\unitlabel{shushu-033.2}丁鄉：苗米三升四合三勺二抄一撮四圭{\zhsmall 按：多四圭，分數少三釐餘故也}，和買一寸七分八釐九毫三絲，夏税一寸三分九毫二絲三忽。\end{unitzh}
\begin{uniten}Township Ding: seedling rice 3 sheng 4 ge (合) 3 shao 2 chao 1 cuo 4 gui [Siku editors: 4 gui too many; the reason is that the number of shares (fenshu) is short by 3-odd li.]; harmonized purchase 1 cun 7 fen 8 li 9 hao 3 si; summer tax 1 cun 3 fen 9 hao 2 si 3 hu (忽).\end{uniten}

\begin{unitzh}\unitlabel{shushu-033.3}戊鄉：苗米一升五合三勺八抄二撮，和買八分一毫九絲二忽，夏税五分八釐六毫七絲七忽。\end{unitzh}
\begin{uniten}Township Wu: seedling rice 1 sheng 5 ge 3 shao 8 chao 2 cuo; harmonized purchase 8 fen 1 hao 9 si 2 hu; summer tax 5 fen 8 li 6 hao 7 si 7 hu.\end{uniten}

\begin{unitzh}\unitlabel{shushu-033.4}己鄉：苗米二升八合二勺一抄五撮一圭，和買一寸四分七釐九絲六忽，夏税一寸七釐六毫三絲。\end{unitzh}
\begin{uniten}Township Ji: seedling rice 2 sheng 8 ge 2 shao 1 chao 5 cuo 1 gui; harmonized purchase 1 cun 4 fen 7 li 9 si 6 hu; summer tax 1 cun 7 li 6 hao 3 si.\end{uniten}

\begin{unitzh}\unitlabel{shushu-033.5}各用鄉錐行數九、八、七、六、五、四、三、二、一，各乗一分之率，為各鄉毎畝等則泛數。或自上等上田之則，以一分之率累減之，亦得。\end{unitzh}
\begin{uniten}In each case use the township awl-row numbers (zhuixing shu) 9, 8, 7, 6, 5, 4, 3, 2, 1, each multiplying the rate (lü) for 1 share (fen); this gives the provisional figures (fanshu) of each township's graded assessments (dengze) per mu (畝). Or, starting from the assessment (ze) of the upper-grade upper fields, successively subtract the rate for 1 share from it; this also gives them.\end{uniten}

\begin{unitzh}\unitlabel{shushu-033.6}甲鄉上等上田：苗米三斗二升三合二勺一抄，和買一尺六寸八分五釐四絲{\zhsmall 按：多四絲，因一分之率多四忽故也，下同此}，夏税一尺二寸三分二釐八毫{\zhsmall 按：少一毫，同上}；中田：苗米二斗九升八勺八抄九撮{\zhsmall 按：少一撮}，和買一尺五寸一分六釐五毫三絲六忽{\zhsmall 按：多三絲六忽}，夏税一尺一寸九釐五毫二絲{\zhsmall 按：少九絲}；下田：苗米二斗五升八合五勺六抄八撮，和買一尺三寸四分八釐三絲二忽{\zhsmall 按：多三絲二忽}，夏税九寸八分六釐二毫四絲{\zhsmall 按：少八絲}。\end{unitzh}
\begin{uniten}Township Jia, upper-grade upper fields: seedling rice 3 dou (斗) 2 sheng 3 ge 2 shao 1 chao; harmonized purchase 1 chi (尺) 6 cun 8 fen 5 li 4 si [Siku editors: 4 si too many, because the rate for 1 share has 4 hu too many; the same applies below.]; summer tax 1 chi 2 cun 3 fen 2 li 8 hao [Siku editors: 1 hao too few; same as above.]. Middle fields: seedling rice 2 dou 9 sheng 8 shao 8 chao 9 cuo [Siku editors: 1 cuo too few.]; harmonized purchase 1 chi 5 cun 1 fen 6 li 5 hao 3 si 6 hu [Siku editors: 3 si 6 hu too many.]; summer tax 1 chi 1 cun 9 li 5 hao 2 si [Siku editors: 9 si too few.]. Lower fields: seedling rice 2 dou 5 sheng 8 ge 5 shao 6 chao 8 cuo; harmonized purchase 1 chi 3 cun 4 fen 8 li 3 si 2 hu [Siku editors: 3 si 2 hu too many.]; summer tax 9 cun 8 fen 6 li 2 hao 4 si [Siku editors: 8 si too few.].\end{uniten}

\begin{unitzh}\unitlabel{shushu-033.7}中等上田：苗米二斗二升六合二勺四抄七撮，和買一尺一寸七分九釐五毫二絲八忽{\zhsmall 按：多二絲八忽}，夏税八寸六分二釐九毫六絲{\zhsmall 按：少七絲}；中田：苗米一斗九升三合九勺二抄六撮，和買一尺一寸一釐二絲四忽{\zhsmall 按：多二絲四忽}，夏税七寸三分九釐六毫八絲{\zhsmall 按：少六絲}；下田：苗米一斗六升一合六勺五撮，和買八寸四分一釐五毫二絲{\zhsmall 按：多二絲}，夏税六寸一分六釐四毫{\zhsmall 按：多五絲}。\end{unitzh}
\begin{uniten}Middle-grade upper fields: seedling rice 2 dou 2 sheng 6 ge 2 shao 4 chao 7 cuo; harmonized purchase 1 chi 1 cun 7 fen 9 li 5 hao 2 si 8 hu [Siku editors: 2 si 8 hu too many.]; summer tax 8 cun 6 fen 2 li 9 hao 6 si [Siku editors: 7 si too few.]. Middle fields: seedling rice 1 dou 9 sheng 3 ge 9 shao 2 chao 6 cuo; harmonized purchase 1 chi 1 cun 1 li 2 si 4 hu [Siku editors: 2 si 4 hu too many.]; summer tax 7 cun 3 fen 9 li 6 hao 8 si [Siku editors: 6 si too few.]. Lower fields: seedling rice 1 dou 6 sheng 1 ge 6 shao 5 cuo; harmonized purchase 8 cun 4 fen 1 li 5 hao 2 si [Siku editors: 2 si too many.]; summer tax 6 cun 1 fen 6 li 4 hao [Siku editors: 5 si too many.].\end{uniten}

\begin{unitzh}\unitlabel{shushu-033.8}下等上田：苗米一斗二升九合二勺八抄四撮，和買六寸七分四釐一絲六忽{\zhsmall 按：多一絲六忽}，夏税四寸九分三釐一毫二絲{\zhsmall 按：少四絲}；中田：苗米九升六合九勺六抄三撮，和買五寸五釐五毫一絲二忽{\zhsmall 按：多一絲二忽}，夏税三寸六分九釐八毫四絲{\zhsmall 按：少三絲}；下田：苗米六升四合六勺四抄二撮，和買三寸三分七釐八忽{\zhsmall 按：多八忽}，夏税二寸四分六釐五毫六絲{\zhsmall 按：少二絲}。\end{unitzh}
\begin{uniten}Lower-grade upper fields: seedling rice 1 dou 2 sheng 9 ge 2 shao 8 chao 4 cuo; harmonized purchase 6 cun 7 fen 4 li 1 si 6 hu [Siku editors: 1 si 6 hu too many.]; summer tax 4 cun 9 fen 3 li 1 hao 2 si [Siku editors: 4 si too few.]. Middle fields: seedling rice 9 sheng 6 ge 9 shao 6 chao 3 cuo; harmonized purchase 5 cun 5 li 5 hao 1 si 2 hu [Siku editors: 1 si 2 hu too many.]; summer tax 3 cun 6 fen 9 li 8 hao 4 si [Siku editors: 3 si too few.]. Lower fields: seedling rice 6 sheng 4 ge 6 shao 4 chao 2 cuo; harmonized purchase 3 cun 3 fen 7 li 8 hu [Siku editors: 8 hu too many.]; summer tax 2 cun 4 fen 6 li 5 hao 6 si [Siku editors: 2 si too few.].\end{uniten}

\begin{unitzh}\unitlabel{shushu-033.9}乙鄉上等上田：苗米三斗六升三合二勺八抄，和買一尺八寸九分三釐九毫二絲，夏税一尺三寸八分五釐八毫{\zhsmall 按：多二絲}；中田：苗米三斗二升六合九勺五抄二撮，和買一尺七寸四釐五毫二絲八忽，夏税一尺二寸四分七釐二毫二絲{\zhsmall 按：多一絲七忽}；下田：苗米二九升六勺二抄四撮，和買一尺五寸一分五釐一毫三絲六忽，夏税一尺一寸八釐六毫四絲{\zhsmall 按：多一絲四忽}。\end{unitzh}
\begin{uniten}Township Yi, upper-grade upper fields: seedling rice 3 dou 6 sheng 3 ge 2 shao 8 chao; harmonized purchase 1 chi 8 cun 9 fen 3 li 9 hao 2 si; summer tax 1 chi 3 cun 8 fen 5 li 8 hao [Siku editors: 2 si too many.]. Middle fields: seedling rice 3 dou 2 sheng 6 ge 9 shao 5 chao 2 cuo; harmonized purchase 1 chi 7 cun 4 li 5 hao 2 si 8 hu; summer tax 1 chi 2 cun 4 fen 7 li 2 hao 2 si [Siku editors: 1 si 7 hu too many.]. Lower fields: seedling rice 2 9 sheng 6 shao 2 chao 4 cuo [Translator's note: 二九升 has no unit after 二; evidently 斗 (dou) is lacking.]; harmonized purchase 1 chi 5 cun 1 fen 5 li 1 hao 3 si 6 hu; summer tax 1 chi 1 cun 8 li 6 hao 4 si [Siku editors: 1 si 4 hu too many.].\end{uniten}

\begin{unitzh}\unitlabel{shushu-033.10}中等上田：苗米二斗五升四合二勺九抄六撮，和買一尺三寸二分五釐七毫四絲四忽，夏税九寸七分六絲{\zhsmall 按：多一絲三忽}；中田：苗米二斗一升七合九勺六抄八撮，和買一尺一寸三分六釐三毫五絲二忽，夏税八寸三分一釐四毫八絲{\zhsmall 按：多一絲一忽}；下田：苗米一斗八升一合六勺四抄，和買九寸四分六釐九毫六絲，夏税六寸九分二釐九毫{\zhsmall 按：多一絲}。\end{unitzh}
\begin{uniten}Middle-grade upper fields: seedling rice 2 dou 5 sheng 4 ge 2 shao 9 chao 6 cuo; harmonized purchase 1 chi 3 cun 2 fen 5 li 7 hao 4 si 4 hu; summer tax 9 cun 7 fen 6 si [Siku editors: 1 si 3 hu too many.]. Middle fields: seedling rice 2 dou 1 sheng 7 ge 9 shao 6 chao 8 cuo; harmonized purchase 1 chi 1 cun 3 fen 6 li 3 hao 5 si 2 hu; summer tax 8 cun 3 fen 1 li 4 hao 8 si [Siku editors: 1 si 1 hu too many.]. Lower fields: seedling rice 1 dou 8 sheng 1 ge 6 shao 4 chao; harmonized purchase 9 cun 4 fen 6 li 9 hao 6 si; summer tax 6 cun 9 fen 2 li 9 hao [Siku editors: 1 si too many.].\end{uniten}

\begin{unitzh}\unitlabel{shushu-033.11}下等上田：苗米一斗四升五合三{\zhsmall 勺}一抄二撮，和買七寸五分七釐五毫六絲八忽，夏税五寸五分四釐三毫二絲{\zhsmall 按：多七忽}；中田：苗米一斗八合九勺八抄四撮，和買五寸六分八釐一毫七絲六忽，夏税四寸一分五釐七毫四絲{\zhsmall 按：多五忽}；下田：苗米七升二合六勺五抄六撮，和買三寸七分八釐七毫八絲四忽，夏税二寸七分七釐一毫六絲{\zhsmall 按：多四忽}。\end{unitzh}
\begin{uniten}Lower-grade upper fields: seedling rice 1 dou 4 sheng 5 ge 3 [Note: shao] 1 chao 2 cuo; harmonized purchase 7 cun 5 fen 7 li 5 hao 6 si 8 hu; summer tax 5 cun 5 fen 4 li 3 hao 2 si [Siku editors: 7 hu too many.]. Middle fields: seedling rice 1 dou 8 ge 9 shao 8 chao 4 cuo; harmonized purchase 5 cun 6 fen 8 li 1 hao 7 si 6 hu; summer tax 4 cun 1 fen 5 li 7 hao 4 si [Siku editors: 5 hu too many.]. Lower fields: seedling rice 7 sheng 2 ge 6 shao 5 chao 6 cuo; harmonized purchase 3 cun 7 fen 8 li 7 hao 8 si 4 hu; summer tax 2 cun 7 fen 7 li 1 hao 6 si [Siku editors: 4 hu too many.].\end{uniten}

\begin{unitzh}\unitlabel{shushu-034.1}丙鄉上等上田：苗米三斗三合四勺八抄{\zhsmall 按：多一抄三撮}，和買一尺五寸八分二釐一毫，夏税一尺一寸五分七釐五毫{\zhsmall 按：少一毫六絲}。中田：苗米二斗七升三合一勺三抄二撮{\zhsmall 按：多一抄一撮}，和買一尺四寸二分三釐八毫九絲，夏税一尺四分一釐七毫五絲{\zhsmall 按：少一毫}。下田：苗米二斗四升二合七勺八抄四撮{\zhsmall 按：多一抄}，和買一尺二寸六分五釐六毫八絲，夏税九寸二分六釐{\zhsmall 按：少九絲}。\end{unitzh}
\begin{uniten}Bing township (丙鄉), upper grade (shangdeng), upper fields (shangtian): seedling rice (miaomi) 3 dou (斗) 3 ge (合) 4 shao (勺) 8 chao (抄) [Siku editors: too much by 1 chao 3 cuo (撮)]; harmonized purchase (hemai) 1 chi (尺) 5 cun (寸) 8 fen (分) 2 li (釐) 1 hao (毫); summer tax (xiashui) 1 chi 1 cun 5 fen 7 li 5 hao [Siku editors: too little by 1 hao 6 si (絲)]. Middle fields: seedling rice 2 dou 7 sheng (升) 3 ge 1 shao 3 chao 2 cuo [Siku editors: too much by 1 chao 1 cuo]; harmonized purchase 1 chi 4 cun 2 fen 3 li 8 hao 9 si; summer tax 1 chi 4 fen 1 li 7 hao 5 si [Siku editors: too little by 1 hao]. Lower fields: seedling rice 2 dou 4 sheng 2 ge 7 shao 8 chao 4 cuo [Siku editors: too much by 1 chao]; harmonized purchase 1 chi 2 cun 6 fen 5 li 6 hao 8 si; summer tax 9 cun 2 fen 6 li [Siku editors: too little by 9 si].\end{uniten}

\begin{unitzh}\unitlabel{shushu-034.2}中等上田：苗米二斗一升二合四勺三抄四撮{\zhsmall 按：多七撮}，和買一尺一寸七釐四毫七絲，夏税八寸一分二毫五絲{\zhsmall 按：多八絲}。中田：苗米一斗八升二合八抄八撮{\zhsmall 按：多八撮}，和買九寸四分九釐二毫六絲，夏税六寸九分四釐五毫{\zhsmall 按：少七絲}。下田：苗米一斗五升一合七勺四抄{\zhsmall 按：多六撮}，和買七寸九分一釐五絲，夏税五寸七分八釐七毫五絲{\zhsmall 按：少六絲}。\end{unitzh}
\begin{uniten}Middle grade, upper fields: seedling rice 2 dou 1 sheng 2 ge 4 shao 3 chao 4 cuo [Siku editors: too much by 7 cuo]; harmonized purchase 1 chi 1 cun 7 li 4 hao 7 si; summer tax 8 cun 1 fen 2 hao 5 si [Siku editors: too much by 8 si]. Middle fields: seedling rice 1 dou 8 sheng 2 ge 8 chao 8 cuo [Siku editors: too much by 8 cuo]; harmonized purchase 9 cun 4 fen 9 li 2 hao 6 si; summer tax 6 cun 9 fen 4 li 5 hao [Siku editors: too little by 7 si]. Lower fields: seedling rice 1 dou 5 sheng 1 ge 7 shao 4 chao [Siku editors: too much by 6 cuo]; harmonized purchase 7 cun 9 fen 1 li 5 si; summer tax 5 cun 7 fen 8 li 7 hao 5 si [Siku editors: too little by 6 si].\end{uniten}

\begin{unitzh}\unitlabel{shushu-034.3}下等上田：苗米一斗二升一合三勺九抄二撮{\zhsmall 按：多六撮}，和買六寸三分二釐八毫四絲，夏税四寸六分三釐{\zhsmall 按：少四絲}。中田：苗米九斗一合四抄四撮{\zhsmall 按：多四撮}，和買四寸七分四釐六毫三絲，夏税三寸四分七釐二毫{\zhsmall 按：多八絲}。下田：苗米六升六勺九抄六撮{\zhsmall 按：多三撮}，和買三寸一分六釐四毫二絲，夏税二寸三分一釐五毫{\zhsmall 按：少二絲}。\end{unitzh}
\begin{uniten}Lower grade, upper fields: seedling rice 1 dou 2 sheng 1 ge 3 shao 9 chao 2 cuo [Siku editors: too much by 6 cuo]; harmonized purchase 6 cun 3 fen 2 li 8 hao 4 si; summer tax 4 cun 6 fen 3 li [Siku editors: too little by 4 si]. Middle fields: seedling rice 9 dou 1 ge 4 chao 4 cuo [Siku editors: too much by 4 cuo] [Translator's note: so in the text; the neighbouring figures point to 9 sheng (九升)]; harmonized purchase 4 cun 7 fen 4 li 6 hao 3 si; summer tax 3 cun 4 fen 7 li 2 hao [Siku editors: too much by 8 si]. Lower fields: seedling rice 6 sheng 6 shao 9 chao 6 cuo [Siku editors: too much by 3 cuo]; harmonized purchase 3 cun 1 fen 6 li 4 hao 2 si; summer tax 2 cun 3 fen 1 li 5 hao [Siku editors: too little by 2 si].\end{uniten}

\begin{unitzh}\unitlabel{shushu-034.4}丁鄉上等上田：苗米三斗四升三合二勺一抄四撮{\zhsmall 按：多四撮}，和買一尺七寸八分九釐三毫，夏税一尺三寸九釐二毫三絲。中田：苗米三斗八合八勺九抄二撮六圭{\zhsmall 按：多三撮餘}，和買一尺六寸一分三毫七絲，夏税一尺一寸七分八釐三毫七忽。下田：苗米二斗七升四合五勺七抄一撮二圭{\zhsmall 按：多三撮餘}，和買一尺四寸三分一釐四毫四絲，夏税一尺四分七釐三毫八絲四忽。\end{unitzh}
\begin{uniten}Ding township (丁鄉), upper grade, upper fields: seedling rice 3 dou 4 sheng 3 ge 2 shao 1 chao 4 cuo [Siku editors: too much by 4 cuo]; harmonized purchase 1 chi 7 cun 8 fen 9 li 3 hao; summer tax 1 chi 3 cun 9 li 2 hao 3 si. Middle fields: seedling rice 3 dou 8 ge 8 shao 9 chao 2 cuo 6 gui (圭) [Siku editors: too much by more than 3 cuo]; harmonized purchase 1 chi 6 cun 1 fen 3 hao 7 si; summer tax 1 chi 1 cun 7 fen 8 li 3 hao 7 hu (忽). Lower fields: seedling rice 2 dou 7 sheng 4 ge 5 shao 7 chao 1 cuo 2 gui [Siku editors: too much by more than 3 cuo]; harmonized purchase 1 chi 4 cun 3 fen 1 li 4 hao 4 si; summer tax 1 chi 4 fen 7 li 3 hao 8 si 4 hu.\end{uniten}

\begin{unitzh}\unitlabel{shushu-034.5}中等上田：苗米二斗四升二勺四抄九撮八圭{\zhsmall 按：多二撮八圭}，和買一尺二寸五分二釐五毫一絲，夏税九寸一分六釐四毫六絲一忽。中田：苗米二斗五合九勺二抄八撮四圭{\zhsmall 按：多二撮四圭}，和買一尺七分三釐五毫八絲，夏税七寸八分五釐五毫三絲八忽。下田：苗米一斗七升一合六勺七撮{\zhsmall 按：多二撮}，和買八寸九分四釐六毫五絲，夏税六寸五分四釐六毫一絲五忽。\end{unitzh}
\begin{uniten}Middle grade, upper fields: seedling rice 2 dou 4 sheng 2 shao 4 chao 9 cuo 8 gui [Siku editors: too much by 2 cuo 8 gui]; harmonized purchase 1 chi 2 cun 5 fen 2 li 5 hao 1 si; summer tax 9 cun 1 fen 6 li 4 hao 6 si 1 hu. Middle fields: seedling rice 2 dou 5 ge 9 shao 2 chao 8 cuo 4 gui [Siku editors: too much by 2 cuo 4 gui]; harmonized purchase 1 chi 7 fen 3 li 5 hao 8 si; summer tax 7 cun 8 fen 5 li 5 hao 3 si 8 hu. Lower fields: seedling rice 1 dou 7 sheng 1 ge 6 shao 7 cuo [Siku editors: too much by 2 cuo]; harmonized purchase 8 cun 9 fen 4 li 6 hao 5 si; summer tax 6 cun 5 fen 4 li 6 hao 1 si 5 hu.\end{uniten}

\begin{unitzh}\unitlabel{shushu-034.6}下等上田：苗米一斗三升七合二勺八抄五撮六圭{\zhsmall 按：多一撮六圭}，和買七寸一分五釐七毫二絲，夏税五寸二分三釐六毫九絲二忽。中田：苗米一斗二合九勺六抄四撮二圭{\zhsmall 按：多一撮二圭}，和買五寸三分六釐七毫九絲，夏税三寸九分二釐七毫六絲九忽。下田：苗米六升八合六勺四抄二撮八圭{\zhsmall 按：多八圭}，和買三寸五分七釐八毫六絲，夏税二寸六分一釐八毫四絲六忽。\end{unitzh}
\begin{uniten}Lower grade, upper fields: seedling rice 1 dou 3 sheng 7 ge 2 shao 8 chao 5 cuo 6 gui [Siku editors: too much by 1 cuo 6 gui]; harmonized purchase 7 cun 1 fen 5 li 7 hao 2 si; summer tax 5 cun 2 fen 3 li 6 hao 9 si 2 hu. Middle fields: seedling rice 1 dou 2 ge 9 shao 6 chao 4 cuo 2 gui [Siku editors: too much by 1 cuo 2 gui]; harmonized purchase 5 cun 3 fen 6 li 7 hao 9 si; summer tax 3 cun 9 fen 2 li 7 hao 6 si 9 hu. Lower fields: seedling rice 6 sheng 8 ge 6 shao 4 chao 2 cuo 8 gui [Siku editors: too much by 8 gui]; harmonized purchase 3 cun 5 fen 7 li 8 hao 6 si; summer tax 2 cun 6 fen 1 li 8 hao 4 si 6 hu.\end{uniten}

\begin{unitzh}\unitlabel{shushu-034.7}戊鄉上等上田：苗米一斗五升三合八勺二抄，和買八寸一釐九毫二絲，夏税五寸八分六釐七毫七絲。中田：苗米一斗三升八合四勺三抄八撮，和買七寸二分一釐七毫二絲八忽，夏税五寸二分八釐九絲三忽。下田：苗米一斗二升三合五抄六撮，和買六寸四分一釐五毫三絲六忽，夏税四寸六分九釐四毫一絲六忽。\end{unitzh}
\begin{uniten}Wu township (戊鄉), upper grade, upper fields: seedling rice 1 dou 5 sheng 3 ge 8 shao 2 chao; harmonized purchase 8 cun 1 li 9 hao 2 si; summer tax 5 cun 8 fen 6 li 7 hao 7 si. Middle fields: seedling rice 1 dou 3 sheng 8 ge 4 shao 3 chao 8 cuo; harmonized purchase 7 cun 2 fen 1 li 7 hao 2 si 8 hu; summer tax 5 cun 2 fen 8 li 9 si 3 hu. Lower fields: seedling rice 1 dou 2 sheng 3 ge 5 chao 6 cuo; harmonized purchase 6 cun 4 fen 1 li 5 hao 3 si 6 hu; summer tax 4 cun 6 fen 9 li 4 hao 1 si 6 hu.\end{uniten}

\begin{unitzh}\unitlabel{shushu-034.8}中等上田：苗米一斗七合六勺七抄四撮，和買五寸六分一釐三毫四絲四忽，夏税四寸一分七毫三絲九忽。中田：苗米九升二合二勺九抄二撮，和買四寸八分一釐一毫五絲二忽，夏税三寸五分二釐六絲二忽。下田：苗米七升六合九勺一抄，和買四寸九毫六絲，夏税二寸九分三釐三毫八絲五忽。\end{unitzh}
\begin{uniten}Middle grade, upper fields: seedling rice 1 dou 7 ge 6 shao 7 chao 4 cuo; harmonized purchase 5 cun 6 fen 1 li 3 hao 4 si 4 hu; summer tax 4 cun 1 fen 7 hao 3 si 9 hu. Middle fields: seedling rice 9 sheng 2 ge 2 shao 9 chao 2 cuo; harmonized purchase 4 cun 8 fen 1 li 1 hao 5 si 2 hu; summer tax 3 cun 5 fen 2 li 6 si 2 hu. Lower fields: seedling rice 7 sheng 6 ge 9 shao 1 chao; harmonized purchase 4 cun 9 hao 6 si; summer tax 2 cun 9 fen 3 li 3 hao 8 si 5 hu.\end{uniten}

\begin{unitzh}\unitlabel{shushu-034.9}下等上田：苗米六升一合五勺二抄八撮，和買三寸二分七毫六絲八忽，夏税二寸三分四釐七毫八忽。中田：苗米四升六合一勺四抄六撮，和買二寸四分五毫七絲六忽，夏税一寸七分六釐二絲一忽。下田：苗米三升七勺六抄四撮，和買一寸六分三毫八絲四忽，夏税一寸一分七釐一三毫五絲四忽。\end{unitzh}
\begin{uniten}Lower grade, upper fields: seedling rice 6 sheng 1 ge 5 shao 2 chao 8 cuo; harmonized purchase 3 cun 2 fen 7 hao 6 si 8 hu; summer tax 2 cun 3 fen 4 li 7 hao 8 hu. Middle fields: seedling rice 4 sheng 6 ge 1 shao 4 chao 6 cuo; harmonized purchase 2 cun 4 fen 5 hao 7 si 6 hu; summer tax 1 cun 7 fen 6 li 2 si 1 hu. Lower fields: seedling rice 3 sheng 7 shao 6 chao 4 cuo; harmonized purchase 1 cun 6 fen 3 hao 8 si 4 hu; summer tax 1 cun 1 fen 7 li 1 3 hao 5 si 4 hu [Translator's note: so in the text, 七釐一三毫; the 一 (1) appears superfluous].\end{uniten}

\begin{unitzh}\unitlabel{shushu-034.10}己鄉上等上田：苗米二斗八升二合一勺五抄一撮，和買一尺四寸七分九毫六絲，夏税一尺七分六釐三毫。中田：苗米二斗五升三合九勺三抄五撮九圭，和買一尺三寸二分三釐八毫六絲四忽，夏税九寸六分八釐六毫七絲。下田：苗米二斗二升五合七勺二抄八圭，和買一尺一寸七分六釐七毫六絲八忽，夏税八寸六分一釐四絲。\end{unitzh}
\begin{uniten}Ji township (己鄉), upper grade, upper fields: seedling rice 2 dou 8 sheng 2 ge 1 shao 5 chao 1 cuo; harmonized purchase 1 chi 4 cun 7 fen 9 hao 6 si; summer tax 1 chi 7 fen 6 li 3 hao. Middle fields: seedling rice 2 dou 5 sheng 3 ge 9 shao 3 chao 5 cuo 9 gui; harmonized purchase 1 chi 3 cun 2 fen 3 li 8 hao 6 si 4 hu; summer tax 9 cun 6 fen 8 li 6 hao 7 si. Lower fields: seedling rice 2 dou 2 sheng 5 ge 7 shao 2 chao 8 gui; harmonized purchase 1 chi 1 cun 7 fen 6 li 7 hao 6 si 8 hu; summer tax 8 cun 6 fen 1 li 4 si.\end{uniten}

\begin{unitzh}\unitlabel{shushu-034.11}中等上田：苗米一斗九升七合五勺五撮七圭，和買一尺二分九釐六毫七絲二忽，夏税七寸五分三釐四毫一絲。中田：苗米一斗六升九合二勺九抄六圭，和買八寸八分二釐五毫七絲六忽，夏税六寸四分五釐七毫八絲。下田：苗米一斗四升一合七抄五撮五圭，和買七寸三分五釐四毫八絲，夏税五寸三分八釐一毫五絲。\end{unitzh}
\begin{uniten}Middle grade, upper fields: seedling rice 1 dou 9 sheng 7 ge 5 shao 5 cuo 7 gui; harmonized purchase 1 chi 2 fen 9 li 6 hao 7 si 2 hu; summer tax 7 cun 5 fen 3 li 4 hao 1 si. Middle fields: seedling rice 1 dou 6 sheng 9 ge 2 shao 9 chao 6 gui; harmonized purchase 8 cun 8 fen 2 li 5 hao 7 si 6 hu; summer tax 6 cun 4 fen 5 li 7 hao 8 si. Lower fields: seedling rice 1 dou 4 sheng 1 ge 7 chao 5 cuo 5 gui; harmonized purchase 7 cun 3 fen 5 li 4 hao 8 si; summer tax 5 cun 3 fen 8 li 1 hao 5 si.\end{uniten}

\begin{unitzh}\unitlabel{shushu-034.12}下等上田：苗米一斗一升二合八勺六抄四圭，和買五寸八分八釐三毫八絲四忽，夏税四寸三分五毫二絲。中田：苗米八升四合六勺四抄五撮三圭，和買四寸四分一釐二毫八絲八忽，夏税三寸二分二釐八毫九絲。下田：苗米五升六合四勺三抄二圭，和買二寸九分四釐一毫九絲二忽，夏税二寸一分五釐二毫六絲。\end{unitzh}
\begin{uniten}Lower grade, upper fields: seedling rice 1 dou 1 sheng 2 ge 8 shao 6 chao 4 gui; harmonized purchase 5 cun 8 fen 8 li 3 hao 8 si 4 hu; summer tax 4 cun 3 fen 5 hao 2 si. Middle fields: seedling rice 8 sheng 4 ge 6 shao 4 chao 5 cuo 3 gui; harmonized purchase 4 cun 4 fen 1 li 2 hao 8 si 8 hu; summer tax 3 cun 2 fen 2 li 8 hao 9 si. Lower fields: seedling rice 5 sheng 6 ge 4 shao 3 chao 2 gui; harmonized purchase 2 cun 9 fen 4 li 1 hao 9 si 2 hu; summer tax 2 cun 1 fen 5 li 2 hao 6 si.\end{uniten}

\begin{unitzh}\unitlabel{shushu-034.13}已上田則，苗至勺、絹至分收歸，為等則定數，答在前。\end{unitzh}
\begin{uniten}The field rates (tianze) above, with the seedling rice down to the shao and the silk (juan) down to the fen, are gathered in (shougui) to make the fixed numbers (dingshu) of the grade rates (dengze); the answer is given before.\end{uniten}

\begin{unitzh}\unitlabel{shushu-034.14}按：右各條下有不合之數者，皆在寄零尾數，至答數收之，勺及分以下俱不用，故數無不合也。\end{unitzh}
\begin{uniten}Note: where, under the items on the right, there are numbers that do not agree, these all lie in the lodged-remainder tail digits (jiling weishu); when it comes to the figures of the answer, they are gathered in, and nothing below the shao and the fen is used at all; therefore none of the numbers disagree.\end{uniten}

\begin{unitzh}\unitlabel{shushu-034.15}圍田租畝\end{unitzh}
\begin{uniten}Rent per Mu (畝) on Enclosed Fields (weitian)\end{uniten}

\begin{unitzh}\unitlabel{shushu-035.1}問：有興復圍田已成，共計三千二十一頃五十一畝一十五歩，分三等：其上等毎畆起租六斗，中等四斗五升，下等四斗。中田多上田弱半，不及下田太半。欲知三色田畝及各租㡬何？\end{unitzh}
\begin{uniten}Question: There are restored enclosed fields (weitian), already completed, totalling 3021 qing (頃) 51 mu (畝) 15 bu (歩), divided into three grades. Of these, the upper grade is levied a rent (qizu) of 6 dou (斗) per mu, the middle grade 4 dou 5 sheng (升), and the lower grade 4 dou. The middle fields exceed the upper fields by a weak half (ruoban) and fall short of the lower fields by a greater half (taiban). It is desired to know the fields in mu of the three kinds and how much each rent is.\end{uniten}

\begin{unitzh}\unitlabel{shushu-035.2}按：題言「中田多上田弱半，不及下田太半」，盖以弱半為三分之一，太半為三分之二也。「多上田弱半」者，即比上田多三分之一也；「不及下田太半」者，即比下田少三分之二也。是上田加三分之一為中田，即下田三分之一也。轉言之，則中田為下田三分之一，上田為中田四分之三也。\end{unitzh}
\begin{uniten}Siku editors' note: The problem says, "The middle fields exceed the upper fields by a weak half and fall short of the lower fields by a greater half"; this presumably takes the weak half as 1/3 and the greater half as 2/3. "Exceed the upper fields by a weak half" means being more than the upper fields by 1/3; "fall short of the lower fields by a greater half" means being less than the lower fields by 2/3. Thus the upper fields with 1/3 added make the middle fields, which are 1/3 of the lower fields. Put the other way round, the middle fields are 1/3 of the lower fields, and the upper fields are 3/4 of the middle fields.\end{uniten}

\begin{unitzh}\unitlabel{shushu-035.3}答曰：上田四百七十七頃八畝一十五歩，米二萬八千六百二十四石八斗三升七合五勺；中田六百三十六頃一十畝三角，米二萬八千六百二十四石八斗三升七合五勺；下田一千九百八頃三十二畝一角，米七萬六千三百三十二石九斗。\end{unitzh}
\begin{uniten}Answer: Upper fields, 477 qing 8 mu 15 bu, rice 28,624 dan (石) 8 dou 3 sheng 7 ge (合) 5 shao (勺); middle fields, 636 qing 10 mu 3 jiao (角), rice 28,624 dan 8 dou 3 sheng 7 ge 5 shao; lower fields, 1908 qing 32 mu 1 jiao, rice 76,332 dan 9 dou.\end{uniten}

\begin{unitzh}\unitlabel{shushu-035.4}術曰：以衰分求之。列母子求田率，副併為法；以共田為寔；寔如法而一，得一分之率；以徧乗未併者，得三等田；各以起租乗之，各得米。\end{unitzh}
\begin{uniten}Method: Find it by proportional distribution (cuifen). Lay out the denominators and numerators (muzi) to find the field rates (lü); add them together in a duplicate position (fubing) to make the divisor (fa); take the total field as the dividend (shi). Divide the dividend by the divisor, obtaining the rate of one part. With it multiply in turn each of those not added together, obtaining the three grades of field. Multiply each by its levied rent, and each gives the rice.\end{uniten}

\begin{unitzh}\unitlabel{shushu-035.5}草曰：置弱半母四為中率，子三為上率；以太半子二减母三，餘一，以乗中率四，只得四，為中泛；又以餘一乗上率三，只得三，為上泛；次以太半母三乗中泛四，得一十二，為下泛；副併三泛，得一十九，為法。置田三千二十一頃五十一畆，以畝法二百四十通之，得七千二百五十一萬六千二百四十歩，内子一十五歩，得七千二百五十一萬六千二百五十五歩，為寔；以法一十九除之，得三百八十一萬六千六百四十五歩，為一分之數。\end{unitzh}
\begin{uniten}Working: Set up the denominator of the weak half, 4, as the middle rate, and the numerator, 3, as the upper rate. Subtract the numerator of the greater half, 2, from the denominator, 3; the remainder is 1. With it multiply the middle rate 4, obtaining just 4, as the middle provisional (fan). Again with the remainder 1 multiply the upper rate 3, obtaining just 3, as the upper provisional. Next with the denominator of the greater half, 3, multiply the middle provisional 4, obtaining 12, as the lower provisional. Add the three provisionals together in a duplicate position, obtaining 19, as the divisor. Set up the field, 3021 qing 51 mu; convert it (tong) by the mu divisor (mufa) 240, obtaining 72,516,240 bu; add in the numerator 15 bu, obtaining 72,516,255 bu, as the dividend. Divide it by the divisor 19, obtaining 3,816,645 bu, as the number of one part.\end{uniten}

\begin{unitzh}\unitlabel{shushu-035.6}以上泛三因之，得一千一百四十四萬九千九百三十五歩，為上積；又以中泛四因之一分數，得一千五百二十六萬六千五百八十歩，為中積；又以下泛一十二乗一分數，得四千五百七十九萬九千七百四十歩，為下積。其三積各以畝法二百四十約之為畆：其上田得四百七十七頃八畝一十五歩，其中田得六百三十六頃一十畝三角，其下積得一千九百八頃三十二畝一角。\end{unitzh}
\begin{uniten}With the upper provisional 3 multiply (yin) it, obtaining 11,449,935 bu, as the upper area. Again with the middle provisional 4 multiply the one-part number, obtaining 15,266,580 bu, as the middle area. Again with the lower provisional 12 multiply the one-part number, obtaining 45,799,740 bu, as the lower area. Reduce (yue) each of the three areas by the mu divisor 240 to make mu: the upper field gives 477 qing 8 mu 15 bu; the middle field gives 636 qing 10 mu 3 jiao; the lower area gives 1908 qing 32 mu 1 jiao.\end{uniten}

\begin{unitzh}\unitlabel{shushu-035.7}各以起租三積為三寔：其上積一千一百四十四萬九千九百三十五歩，乗上積六斗，得六百八十六萬九千九百六十一石，為寔；以畝法二百四十除之，得二萬八千六百二十四石八斗三升七合五勺，為上田租。其中積一千五百二十六萬六千五百八十歩，乗中田租四斗五升，得六百八十六萬九千九百六十一石，為寔；以畝法二百四十除之，得二萬八千六百二十四石八斗三升七合五勺。其下積四千五百七十九萬九千七石四十歩，乗下租四斗，得一千八百三十一萬九千八百九十六石，為寔；以畝法二百四十除之，得七萬六千三百三十二石九斗，為下田米。\end{unitzh}
\begin{uniten}Using each levied rent, make the three areas into three dividends. The upper area, 11,449,935 bu, multiplied by the upper area 6 dou [Translator's note: 上積 "upper area" is evidently an error for 上租 "upper rent"], gives 6,869,961 dan, as the dividend; divide it by the mu divisor 240, obtaining 28,624 dan 8 dou 3 sheng 7 ge 5 shao, as the upper-field rent. The middle area, 15,266,580 bu, multiplied by the middle-field rent 4 dou 5 sheng, gives 6,869,961 dan, as the dividend; divide it by the mu divisor 240, obtaining 28,624 dan 8 dou 3 sheng 7 ge 5 shao. The lower area, 4579 myriad 9 thousand 7 dan 40 bu [Translator's note: 石 "dan" is evidently an error for 百 "hundred"], multiplied by the lower rent 4 dou, gives 18,319,896 dan, as the dividend; divide it by the mu divisor 240, obtaining 76,332 dan 9 dou, as the lower-field rice.\end{uniten}

\begin{unitzh}\unitlabel{shushu-035.8}按：草内求各衰數，意謂以三分為上田衰數，加弱半三分之一，得四分，為中田衰數；三因中田衰數，得一十二分，為下田衰數；併之，得一十九分，為總衰數。其法本属顯明，但語中參入母子、副、泛等名目，其意反晦矣。\end{unitzh}
\begin{uniten}Siku editors' note: As for finding each proportional number (cuishu) in the working, the meaning is: take 3 parts as the proportional number of the upper field; add the weak half, 1/3, obtaining 4 parts, as the proportional number of the middle field; multiply the proportional number of the middle field by 3, obtaining 12 parts, as the proportional number of the lower field; add them together, obtaining 19 parts, as the total proportional number. The method is in itself plain and clear, but terms such as "denominator and numerator" (muzi), "duplicate" (fu) and "provisional" (fan) are mixed into the wording, so that its meaning is, on the contrary, obscured.\end{uniten}

\begin{unitzh}\unitlabel{shushu-035.9}《數學九章》卷五上\end{unitzh}
\begin{uniten}Shuxue jiuzhang, juan 5, part 1.\end{uniten}

\begin{unitzh}\unitlabel{shushu-035.10}子部·天文算法類·算書之屬·數學九章\end{unitzh}
\begin{uniten}Masters branch (zibu), category of astronomy and mathematical methods (tianwen suanfa lei), subdivision of mathematical books (suanshu zhi shu), Shuxue jiuzhang. [Translator's note: the Siku quanshu classification of the work.]\end{uniten}

\section*{卷5下}\addcontentsline{toc}{section}{卷5下}

\begin{unitzh}\unitlabel{shushu-036.1}欽定四庫全書\end{unitzh}
\begin{uniten}Imperially Approved Complete Library of the Four Treasuries (Qinding Siku quanshu)\end{uniten}

\begin{unitzh}\unitlabel{shushu-036.2}數學九章卷五下 宋 秦九韶 撰\end{unitzh}
\begin{uniten}Mathematical Treatise in Nine Sections (Shuxue jiuzhang), Juan 5, Part 2. Composed by Qin Jiushao of the Song.\end{uniten}

\begin{unitzh}\unitlabel{shushu-036.3}賦役\end{unitzh}
\begin{uniten}Taxes and Labour Services (fuyi)\end{uniten}

\begin{unitzh}\unitlabel{shushu-036.4}户税移割\end{unitzh}
\begin{uniten}Transfer and Division of Household Taxes\end{uniten}

\begin{unitzh}\unitlabel{shushu-036.5}問：某縣據甲稱：「本户田地原納苗三十五石七斗，和買本色一十一匹二丈二尺九寸四分八釐七毫五絲，折帛二十七匹二寸一分三釐七毫五絲，紬絹折帛八匹三丈九寸七寸三分九釐，紬絹本色二十匹二丈八尺九寸三分九釐。已将田四百七畆出與乙，五百十六畆出與丙訖，乞移割本户所出田上税賦，歸併乙、丙兩户送納。」㑹到乙原有田三百七十五畆，兩原有田四百六十三畆，並係本鄉本等。毎畆苗三升五合，税一尺一寸五分，物力一貫二百；本等田紬一尺三寸四分，物力九百。物力及三十二貫數，和買一匹，内三分本色，七分折帛；夏税七分本色，三分折帛；紬五分本色，五分折帛，併絹紬。欲知甲田地及甲、乙、丙分割合納苗米、和買、夏税折帛、本色、畸零各幾何？\end{unitzh}
\begin{uniten}Question: In a certain county, according to Jia's statement: 'The fields and land of this household originally paid a grain levy (miao) of 35 dan (石) 7 dou (斗); harmonious purchase (hemai): in kind (bense) 11 pi (匹) 2 zhang (丈) 2 chi (尺) 9 cun (寸) 4 fen (分) 8 li (釐) 7 hao (毫) 5 si (絲), commuted silk (zhebo) 27 pi 2 cun 1 fen 3 li 7 hao 5 si; pongee and silk (chou juan) commuted silk 8 pi 3 zhang 9 cun 7 cun 3 fen 9 li [Translator's note: '9 cun 7 cun' is evidently corrupt; the working has 9 chi 7 cun]; pongee and silk in kind 20 pi 2 zhang 8 chi 9 cun 3 fen 9 li. Having already given out 407 mu (畆) of fields to Yi and 516 mu to Bing, I request that the taxes and levies on the fields given out by this household be transferred and divided off (yige), to be merged into the two households of Yi and Bing for delivery and payment.' On checking, Yi originally had 375 mu of fields, and liang (兩) [Translator's note: 兩 'two' is evidently an error for 丙, Bing] originally had 463 mu of fields, all belonging to this township and this grade. Per mu: grain levy 3 sheng (升) 5 ge (合), tax 1 chi 1 cun 5 fen, property assessment (wuli) 1 guan (貫) 200; for this grade, field pongee 1 chi 3 cun 4 fen, property assessment 900. When the property assessment reaches the amount of 32 guan, the harmonious purchase is 1 pi, of which 3 parts are in kind and 7 parts commuted silk; summer tax (xiashui): 7 parts in kind, 3 parts commuted silk; pongee: 5 parts in kind, 5 parts commuted silk; silk and pongee combined. It is desired to know Jia's fields and land, and, as divided among Jia, Yi and Bing, how much each should pay in grain-levy rice (miaomi), harmonious purchase, summer-tax commuted silk, in kind, and odd remainder (jiling).\end{uniten}

\begin{unitzh}\unitlabel{shushu-036.6}按：此題科税分田地二項：田有科米、税絹、物力三色，地只有税紬、物力二色；絹與紬折色不同；物力、和買合田地總計之。又甲户有田者地，乙、丙二户有田無地。此等皆不可不明言者，題中殊未分晰。\end{unitzh}
\begin{uniten}Note: In this problem the assessed taxes (keshui) are divided into two items, fields and land: fields have three kinds, levied rice, tax silk and property assessment; land has only two kinds, tax pongee and property assessment; silk and pongee differ in their commutation (zhese); property assessment and harmonious purchase are reckoned on fields and land combined in total. Further, household Jia has fields zhe (者) land [Translator's note: 者 is evidently an error for 有, i.e., 'has fields and has land'], while the two households Yi and Bing have fields but no land. All these are matters that must be stated clearly, yet in the problem they are not distinguished clearly at all.\end{uniten}

\begin{unitzh}\unitlabel{shushu-036.7}答曰：甲原有田一千二十畆，地一十一畆二角四十八步。甲今有田九十七畆：苗米三石三斗九升五合；夏税折帛三丈三尺四寸六分五釐，本色一匹，畸零三丈八尺八分五釐；物力一百一十六貫四百文。地一十一畆二角四十八步：税紬折帛七尺八寸三分九釐，税紬畸零七尺八寸三分九釐；物力一十貫五百三十文。田地共物力一百二十六貫九百三十文；和買折帛二匹三丈一尺六分三釐七毫五絲，本色一疋，畸零七尺五寸九分八釐七毫五絲。乙今有田七百八十二畆：苗米二十七石三斗七升；夏税折帛六匹二丈九尺七寸九分，本色一十五匹，畸零二丈九尺五寸一分；物力九百三十八貫四百文；和買折帛二十匹二丈一尺一寸，本色八匹，畸零三丈一尺九寸。丙今有田九百七十九畆：苗米三十四石二斗六升五合；夏税折帛八匹一丈七尺七寸五分五釐，本色一十九匹，畸零二丈八尺九分五釐；物力一千一百七十四貫八百文；和買折帛二十五匹二丈七尺九寸五分，本色一十一匹，畸零五寸五分。\end{unitzh}
\begin{uniten}Answer: Jia originally had 1020 mu of fields, and 11 mu 2 jiao (角) 48 bu (步) of land. Jia now has 97 mu of fields: grain-levy rice 3 dan 3 dou 9 sheng 5 ge; summer-tax commuted silk 3 zhang 3 chi 4 cun 6 fen 5 li, in kind 1 pi, odd remainder 3 zhang 8 chi 8 fen 5 li; property assessment 116 guan 400 wen (文). Land 11 mu 2 jiao 48 bu: tax-pongee commuted silk 7 chi 8 cun 3 fen 9 li, tax-pongee odd remainder 7 chi 8 cun 3 fen 9 li; property assessment 10 guan 530 wen. Fields and land together: property assessment 126 guan 930 wen; harmonious-purchase commuted silk 2 pi 3 zhang 1 chi 6 fen 3 li 7 hao 5 si, in kind 1 pi (疋), odd remainder 7 chi 5 cun 9 fen 8 li 7 hao 5 si. Yi now has 782 mu of fields: grain-levy rice 27 dan 3 dou 7 sheng; summer-tax commuted silk 6 pi 2 zhang 9 chi 7 cun 9 fen, in kind 15 pi, odd remainder 2 zhang 9 chi 5 cun 1 fen; property assessment 938 guan 400 wen; harmonious-purchase commuted silk 20 pi 2 zhang 1 chi 1 cun, in kind 8 pi, odd remainder 3 zhang 1 chi 9 cun. Bing now has 979 mu of fields: grain-levy rice 34 dan 2 dou 6 sheng 5 ge; summer-tax commuted silk 8 pi 1 zhang 7 chi 7 cun 5 fen 5 li, in kind 19 pi, odd remainder 2 zhang 8 chi 9 fen 5 li; property assessment 1174 guan 800 wen; harmonious-purchase commuted silk 25 pi 2 zhang 7 chi 9 cun 5 fen, in kind 11 pi, odd remainder 5 cun 5 fen.\end{uniten}

\begin{unitzh}\unitlabel{shushu-036.8}術曰：以粟米及衰分求之。置甲原納米為實，以毎畆苗為法除之，得甲原有田；以乗毎畆税，得田絹。次以甲紬絹本色、折帛併之，内減田絹，餘為地紬；以毎紬約之，得甲原有地。次以甲出田併乙、丙原有田，各得三户今有田地。各以等則乗之，各得物力、苗、税。次以毎畆物力率約各户共物力，得和買。副之：三分因之，退位，為和本；七分因之，退位，為和折。夏税反其分而因之、退之。紬以半之，併入税，各得。\end{unitzh}
\begin{uniten}Method: Seek it by 'Millet and Rice' (sumi) and 'Distribution by Proportion' (cuifen). Set out the rice originally paid by Jia as the dividend (shi), and divide it with the grain levy per mu as the divisor (fa), obtaining Jia's original fields; with it multiply the tax per mu, obtaining the field silk. Next combine Jia's pongee-and-silk in kind and commuted silk, and subtract the field silk from it; the remainder is the land pongee; divide (yue) it by the per-unit pongee, obtaining Jia's original land. Next combine the fields given out by Jia with the original fields of Yi and Bing, obtaining for each of the three households the fields and land it now has. Multiply each by the grade rates (dengze), obtaining for each the property assessment, grain levy and tax. Next divide each household's total property assessment by the per-mu property-assessment rate (lü), obtaining the harmonious purchase. Duplicate it (fu): multiply (yin) it by 3 parts and move it back a place (tuiwei), which is the harmonious in-kind (heben); multiply it by 7 parts and move it back a place, which is the harmonious commuted (hezhe). For the summer tax, reverse the parts, then multiply and move back. Halve the pongee and combine it into the tax; each is obtained.\end{uniten}

\begin{unitzh}\unitlabel{shushu-036.9}草曰：置甲原納米三十五石七斗為實，以毎畆苗三升五合為法除之，得一千二十畆，為甲原有田。以乗毎畆税一尺一寸五分，得一百一十七丈三尺，為絹積尺。次以甲原納絹紬折帛八匹三丈九尺七寸三分九釐，併紬絹本色二十匹二丈八尺九寸三分九釐，得二十九匹二丈八尺六寸七分八釐。以匹法四丈通匹數，内零丈，得一百一十八丈八尺六寸七分八釐；内減田絹積丈尺一百一十七丈三尺，餘一丈五尺六寸七分八釐，為甲地紬。以毎畆紬一尺三寸四分納之，得一十一畆七分。其畆下七分倍之，得一百四十；身下加二，得一百六十八步；以六十步納之，得二角零四十八步。得一十一畆二角四十八步，為甲原有地。\end{unitzh}
\begin{uniten}Working: Set out the rice originally paid by Jia, 35 dan 7 dou, as the dividend (shi), and divide it with the grain levy per mu, 3 sheng 5 ge, as the divisor (fa), obtaining 1020 mu, which is Jia's original fields. With it multiply the tax per mu, 1 chi 1 cun 5 fen, obtaining 117 zhang 3 chi, which is the silk total in chi (jichi). Next take the silk-and-pongee commuted silk originally paid by Jia, 8 pi 3 zhang 9 chi 7 cun 3 fen 9 li, and combine it with the pongee-and-silk in kind, 20 pi 2 zhang 8 chi 9 cun 3 fen 9 li, obtaining 29 pi 2 zhang 8 chi 6 cun 7 fen 8 li. Convert (tong) the number of pi by the pi standard (pifa) of 4 zhang and add in (nei) the odd zhang, obtaining 118 zhang 8 chi 6 cun 7 fen 8 li; from this subtract the field-silk total in zhang and chi, 117 zhang 3 chi; the remainder, 1 zhang 5 chi 6 cun 7 fen 8 li, is Jia's land pongee. Take it in (na) by the pongee per mu, 1 chi 3 cun 4 fen [Translator's note: 納 'take in' is used here where the method has 約 'divide'], obtaining 11 mu 7 fen. Double the 7 fen below the mu, obtaining 140; 'add two below the body' (shenxia jia er) [Translator's note: i.e., multiply by 1.2], obtaining 168 bu; take it in by 60 bu, obtaining 2 jiao and an odd (ling) 48 bu. This gives 11 mu 2 jiao 48 bu, which is Jia's original land.\end{uniten}

\begin{unitzh}\unitlabel{shushu-036.10}乃以甲所出四百七畆及五百一十六畆，併得九百二十三畆，減原有一千二十畆，餘九十七畆，為甲今有田併原地{\zhsmall 按：三字誤，葢地一十一畆餘不在田一千二十畆之内也。}。次以甲出田四百七畆併乙原有田三百七十五畆，共得七百八十二畆，為乙令有田。次以甲所出田五百一十六畆併丙原有田四百六十三畆，共得九百七十九畆，為丙今有田。各列三户今有田地數於右行，副之。先以毎畆苗三升五合徧乗右行：甲得三石三斗九升五合，乙得二十七石三斗七升，兩得三十四石二斗六升五合，為本户苗米。\end{unitzh}
\begin{uniten}Then take what Jia gave out, 407 mu and 516 mu, which combined give 923 mu; subtract it from the original 1020 mu, leaving 97 mu, which is the fields Jia now has combined with the original land [Siku editors: These three characters are wrong; for the land of 11-odd mu is not within the 1020 mu of fields.]. Next combine the 407 mu of fields given out by Jia with Yi's original fields of 375 mu, together obtaining 782 mu, which is the fields Yi ling (令) has [Translator's note: 令 is evidently an error for 今 'now']. Next combine the 516 mu of fields given out by Jia with Bing's original fields of 463 mu, together obtaining 979 mu, which is the fields Bing now has. Set out the numbers of fields and land the three households now have in the right column (youxing), and duplicate it (fu). First multiply the right column throughout (biancheng) by the grain levy per mu, 3 sheng 5 ge: Jia obtains 3 dan 3 dou 9 sheng 5 ge, Yi obtains 27 dan 3 dou 7 sheng, liang (兩) [Translator's note: evidently an error for 丙, Bing] obtains 34 dan 2 dou 6 sheng 5 ge; these are the grain-levy rice of the respective households.\end{uniten}

\begin{unitzh}\unitlabel{shushu-036.11}次以毎畆税一尺一寸五分徧乗左副行：甲得一十一丈一尺五寸五分，乙得八十九丈九尺三寸，丙得一百一十二丈五尺八寸五分，各為丈積。各列二位，皆以三分因頭位，七分因下位，併退一位，即是自下三七折之。又各以四大約丈積成匹：其甲得三丈三尺四寸六分五釐，為夏税折帛；又得一匹三丈分尺八分五釐，為夏税本色。其一得六匹二丈九尺七寸九分，為夏税折帛；又得一十五匹二丈九尺五寸一分，為夏税本色。其丙得八匹一丈七尺七寸五分五釐，為夏税折帛；得一十九匹二丈分尺九分五釐，為夏税本色。\end{unitzh}
\begin{uniten}Next multiply the left auxiliary column (zuo fuxing) throughout by the tax per mu, 1 chi 1 cun 5 fen: Jia obtains 11 zhang 1 chi 5 cun 5 fen, Yi obtains 89 zhang 9 chi 3 cun, Bing obtains 112 zhang 5 chi 8 cun 5 fen; each is the total in zhang (zhangji). Set out each in two positions; in all cases multiply the head position by 3 parts and the lower position by 7 parts, and move both back one place: this is precisely taking them at three and seven (sanqi zhe) from below. Again divide each total in zhang by 4 da (大) [Translator's note: 大 is evidently an error for 丈, zhang] to make pi. Of these, Jia obtains 3 zhang 3 chi 4 cun 6 fen 5 li, which is the summer-tax commuted silk, and further obtains 1 pi 3 zhang fen chi 8 fen 5 li [Translator's note: garbled; the answer has 3 zhang 8 chi 8 fen 5 li], which is the summer tax in kind. 'One' (一) [Translator's note: evidently an error for 乙, Yi] obtains 6 pi 2 zhang 9 chi 7 cun 9 fen, which is the summer-tax commuted silk, and further obtains 15 pi 2 zhang 9 chi 5 cun 1 fen, which is the summer tax in kind. Bing obtains 8 pi 1 zhang 7 chi 7 cun 5 fen 5 li, which is the summer-tax commuted silk, and obtains 19 pi 2 zhang fen chi 9 fen 5 li [Translator's note: garbled; the answer has 2 zhang 8 chi 9 fen 5 li], which is the summer tax in kind.\end{uniten}

\begin{unitzh}\unitlabel{shushu-036.12}次以田力一貫二百徧乗右副行，併以地物力乗甲原地及納：甲得一百一十六貫四百，乙得九百三十八貫四百，丙得一千一百七十四貫八百，各為田物力。其甲又置一十一畆七分，以乗地紬一尺三寸四分，得一丈五尺六寸七分八釐，為甲地紬；半之，得七尺八寸三分九釐，各為折紬、本色。又以甲地乗物力九百，得一十貫五百三十，為甲地乗力；併田物力一百一十六貫四百，共得一百二十六貫九百三十文，為甲物力。\end{unitzh}
\begin{uniten}Next multiply the right auxiliary column throughout by the field assessment (tianli, 田力), 1 guan 200, and also multiply Jia's original land by the land property assessment and take it in (na): Jia obtains 116 guan 400, Yi obtains 938 guan 400, Bing obtains 1174 guan 800; each is the field property assessment. For Jia, further set out 11 mu 7 fen, and with it multiply the land pongee, 1 chi 3 cun 4 fen, obtaining 1 zhang 5 chi 6 cun 7 fen 8 li, which is Jia's land pongee; halve it, obtaining 7 chi 8 cun 3 fen 9 li, which are respectively the commuted pongee and the in kind. Further, with Jia's land multiply the property assessment 900, obtaining 10 guan 530, which is Jia's land 'multiplied assessment' (乗力) [Translator's note: 乗 'multiply' is evidently an error for 物, i.e., land property assessment]; combine it with the field property assessment, 116 guan 400, together obtaining 126 guan 930 wen, which is Jia's property assessment.\end{uniten}

\begin{unitzh}\unitlabel{shushu-036.13}列甲、乙、丙三户共物力，各為貫，皆以物力和買率三十二貫為法除之：甲得三匹九分六釐六毫五絲六忽二㣲五塵，乙得二十九匹三分二釐五毫，丙得三十六匹七分一釐二毫五絲，各為和買率。亦各列二位，各以七三折之：其頭位七折者為和買折帛，下位三折者為和買本色。折訖，其甲得二匹七分七釐六毫五絲九忽三㣲七塵五沙，為和買折帛率；又得一匹一分八釐九毫九絲六忽八㣲七塵五沙，為和買本色率。乙得二十匹五分二釐七毫五絲，為和買折帛率；又得八匹七分九釐七毫五絲，為和買本色率。丙得二十五匹六分九釐八毫七絲五忽，為和買折帛率；又得一十一匹一釐三毫七絲五忽，為和買本色率。\end{unitzh}
\begin{uniten}Set out the total property assessments of the three households Jia, Yi and Bing, each made into guan, and divide all of them with the property-assessment harmonious-purchase rate (lü) of 32 guan as the divisor (fa): Jia obtains 3 pi 9 fen 6 li 6 hao 5 si 6 hu (忽) 2 wei (㣲) 5 chen (塵), Yi obtains 29 pi 3 fen 2 li 5 hao, Bing obtains 36 pi 7 fen 1 li 2 hao 5 si; each is the harmonious-purchase rate. Likewise set out each in two positions, and take each at seven and three (qisan zhe): that in the head position, taken at seven, is the harmonious-purchase commuted silk; that in the lower position, taken at three, is the harmonious purchase in kind. When the taking is done, Jia obtains 2 pi 7 fen 7 li 6 hao 5 si 9 hu 3 wei 7 chen 5 sha (沙), which is the harmonious-purchase commuted-silk rate, and further obtains 1 pi 1 fen 8 li 9 hao 9 si 6 hu 8 wei 7 chen 5 sha, which is the harmonious-purchase in-kind rate. Yi obtains 20 pi 5 fen 2 li 7 hao 5 si, which is the harmonious-purchase commuted-silk rate, and further obtains 8 pi 7 fen 9 li 7 hao 5 si, which is the harmonious-purchase in-kind rate. Bing obtains 25 pi 6 fen 9 li 8 hao 7 si 5 hu, which is the harmonious-purchase commuted-silk rate, and further obtains 11 pi 1 li 3 hao 7 si 5 hu, which is the harmonious-purchase in-kind rate.\end{uniten}

\begin{unitzh}\unitlabel{shushu-036.14}各率除端匹外，乃以定下分毫皆以四因之，收為丈尺寸分：甲得和買折帛二匹三丈一尺六分三釐七毫五絲，本色一匹七尺五寸九分八釐七毫五絲；乙得和買折帛二十匹二丈一尺一寸，本色八匹三丈一尺九寸；丙得和買折帛二十五匹二丈七尺九寸五分，本色一十一匹五寸五分；為各户和買數及分本色下畸零數。合問。\end{unitzh}
\begin{uniten}For each rate, apart from the whole pi (duanpi), take the fen and hao below the fixed place (dingxia) and multiply them all by 4, gathering them (shou) into zhang, chi, cun and fen: Jia obtains harmonious-purchase commuted silk 2 pi 3 zhang 1 chi 6 fen 3 li 7 hao 5 si, in kind 1 pi 7 chi 5 cun 9 fen 8 li 7 hao 5 si; Yi obtains harmonious-purchase commuted silk 20 pi 2 zhang 1 chi 1 cun, in kind 8 pi 3 zhang 1 chi 9 cun; Bing obtains harmonious-purchase commuted silk 25 pi 2 zhang 7 chi 9 cun 5 fen, in kind 11 pi 5 cun 5 fen. These are the harmonious-purchase amounts of each household and the odd-remainder amounts below the in-kind portion. This answers the question.\end{uniten}

\begin{unitzh}\unitlabel{shushu-036.15}均科綿税\end{unitzh}
\begin{uniten}Equal Assessment of the Silk-Floss Tax\end{uniten}

\begin{unitzh}\unitlabel{shushu-037.1}問：縣科綿，有五等户，共一萬一千三十三户，共科綿八萬八千三百三十七兩六錢。上等一十二户，副等八十七户，中等四百六十四户，次等二千三十五户，下等八千四百三十五户。欲令上三等折半差，下二等此中等六四折差，科率求之。問各户納及各等幾何？\end{unitzh}
\begin{uniten}Question: A county levies silk floss (mian). There are households of five grades, 11033 households in all, and the floss levied in all is 88337 liang (兩) 6 qian (錢). Upper grade (shangdeng), 12 households; secondary grade (fudeng), 87 households; middle grade (zhongdeng), 464 households; lesser grade (cideng), 2035 households; lower grade (xiadeng), 8435 households. It is desired that the upper three grades have a reduced-by-half difference (zheban cha), and the lower two grades, this [Translator's note: 此 "this" is evidently an error for 比 "compared with"; cf. the note below] middle grade, a 6–4 reduced difference (liu si zhe cha); find them with assessment rates (kelü). What is the payment of each household, and of each grade?\end{uniten}

\begin{unitzh}\unitlabel{shushu-037.2}按：題言「下二等比中等六四折差」，是次等為中等六分之四，下等又為次等六分之四也。乃術草中皆以十分之四收之，是四分折差，而非四六折差矣。集中題與術不相應者多類此，豈成書之後未能重加校正歟？\end{unitzh}
\begin{uniten}Note: The problem says "the lower two grades, compared with the middle grade, have a 6–4 reduced difference"; this means that the lesser grade is 4/6 of the middle grade, and the lower grade in turn is 4/6 of the lesser grade. Yet in the method and the working both are taken at 4/10; this is a 4-fen reduced difference and not a 4–6 reduced difference. Cases in the collection where the problem and the method do not correspond are mostly of this kind; can it be that after the book was completed it could not be collated and corrected again?\end{uniten}

\begin{unitzh}\unitlabel{shushu-037.3}答曰：上等一户一百二十四兩，一十二户計一千四百八十八兩。副等一户六十二兩，八十七户計五千三百九十四兩。中等一户三十一兩，四百六十四户計一萬四千三百八十四兩。次等一户一十二兩四錢，二千三十五户計二萬五千二百三十四兩。下等一户四兩九錢六分，八千四百三十五户計四萬一千八百三十七兩六錢。\end{unitzh}
\begin{uniten}Answer: Upper grade: one household 124 liang; 12 households, in total 1488 liang. Secondary grade: one household 62 liang; 87 households, in total 5394 liang. Middle grade: one household 31 liang; 464 households, in total 14384 liang. Lesser grade: one household 12 liang 4 qian; 2035 households, in total 25234 liang. Lower grade: one household 4 liang 9 qian 6 fen (分); 8435 households, in total 41837 liang 6 qian.\end{uniten}

\begin{unitzh}\unitlabel{shushu-037.4}術曰：列五等户數，先以四折下等數加次等户，又以四折之加中等户數，却以半折之加二等户，又以半折之加上等户數，不折，便為法。除科綿，得上等一户之綿。復半之為副等，又半之為中等，又四折之為次等，又四折之為下等，各一户綿。却各以户數乗之，各得五等共出綿。具圖折之。\end{unitzh}
\begin{uniten}Method: Set out the numbers of households of the five grades. First reduce the number of the lower grade at 4 (si zhe) [Translator's note: i.e., take 4/10 of it] and add it to the households of the lesser grade; again reduce this at 4 and add it to the number of households of the middle grade; then reduce this by half and add it to the households of the second grade; again reduce this by half and add it to the number of households of the upper grade; without reducing, this is then the divisor (fa). Divide the levied floss by it, obtaining the floss of one household of the upper grade. Halve it again for the secondary grade; halve it again for the middle grade; again reduce it at 4 for the lesser grade; again reduce it at 4 for the lower grade: in each case the floss of one household. Then multiply each by its number of households, obtaining in each case the floss paid in all by the five grades. Set out a diagram and reduce them.\end{uniten}

\begin{unitzh}\unitlabel{shushu-037.5}上 副 中 次 下\end{unitzh}
\begin{uniten}Upper　Secondary　Middle　Lesser　Lower\end{uniten}

\begin{unitzh}\unitlabel{shushu-037.6}草曰：先置下等户八千四百三十五，以四折之，得三千三百七十四。乃以得數折入次户二千三十五内，得五千四百九户，訖。又以四折次數五千四百九户，得二千一百六十三户六分。乃以得次數併入中户四百六十四内，共得二千六百二十七户六分。次以五分折中數二千六百二十七户六分，得數。次以得數一千三百一十三户八分併副户八十七，得一千四百户八分。又以五折副數一千四百户八分，得七百户四分。既得數，乃以副數七百户四分併上户一十二户，得七百一十二户四分，不折，便以為法。為法。\end{unitzh}
\begin{uniten}Working: First set up the households of the lower grade, 8435; reduce them at 4, obtaining 3374. Then enter the obtained number by reduction (zheru) into the 2035 households of the lesser grade, obtaining 5409 households; done. Again reduce at 4 the lesser-grade number, 5409 households, obtaining 2163 households 6 fen [Translator's note: here fen are tenths of a household]. Then combine the obtained lesser-grade number into the 464 households of the middle grade, obtaining in all 2627 households 6 fen. Next reduce at 5 fen the middle-grade number, 2627 households 6 fen, obtaining a number. Next combine the obtained number, 1313 households 8 fen, with the 87 households of the secondary grade, obtaining 1400 households 8 fen. Again reduce at 5 the secondary-grade number, 1400 households 8 fen, obtaining 700 households 4 fen. Having obtained the number, then combine the secondary-grade number, 700 households 4 fen, with the 12 households of the upper grade, obtaining 712 households 4 fen; without reducing, take this as the divisor. As the divisor. [Translator's note: so repeated in the text.]\end{uniten}

\begin{unitzh}\unitlabel{shushu-037.7}累折至等，共得七百一十二户四分，以為總法。乃置科綿八萬八千三百三十七兩六錢為實，法法七百一十二户四分而一，得一百二十四兩，為上等一户所出綿。以五折之，得六十二兩，為副等一户所出綿。又五折之，得三十一兩，為副等一户所出綿。乃以四折之，得一十二兩四錢，為次等一户所出綿。又以四折之，得四兩九錢六分，為下等一户所出綿。乃以各等户數各乗一户所出，即各得每等共數之綿。\end{unitzh}
\begin{uniten}Reducing successively through the grades, one obtains in all 712 households 4 fen, which is taken as the total divisor (zongfa). Then set up the levied floss, 88337 liang 6 qian, as the dividend (shi); divide by the divisor, divisor [Translator's note: 法 "divisor" is written twice] 712 households 4 fen, obtaining 124 liang as the floss paid by one household of the upper grade. Reduce it at 5, obtaining 62 liang as the floss paid by one household of the secondary grade. Again reduce it at 5, obtaining 31 liang as the floss paid by one household of the secondary grade [Translator's note: so the text; the middle grade is evidently meant]. Then reduce it at 4, obtaining 12 liang 4 qian as the floss paid by one household of the lesser grade. Again reduce it at 4, obtaining 4 liang 9 qian 6 fen as the floss paid by one household of the lower grade. Then multiply, for each grade, what one household pays by the number of households of that grade; this gives in each case the total floss of each grade.\end{uniten}

\begin{unitzh}\unitlabel{shushu-037.8}均定勸分\end{unitzh}
\begin{uniten}Equitably Fixing the Exhortation to Share (quanfen)\end{uniten}

\begin{unitzh}\unitlabel{shushu-037.9}問：欲勸糶賑濟，據甲民户物力畆步排定，共計一百六十二户，作九等：上等三户，第二等五户，第三等七户，第四等八户，第五等十三户，第六等二十一户，第七等二十六户，第八等三十四户，第九等四十五户。今先觀諭，第一等上户願糶五千石，第九等户願糶二百石。欲知各等抛差石數，併數認米數，各幾何？答曰：認該米二十三萬七千六百石。上等一户米五千石，三户計一萬五千石。二等一户米四千四百石，五户計二萬二千石。三等一户米三千八百石，七户計二萬六千六百石。四等一户米三千二百石，八户計二萬五千六百石。五等一户米二千六百石，一十三户計三萬三千八百石。六等一户米二千石，二十一户計四萬二千石。七等一户米一千四百石，三十六户計三萬六千四百石。八等一户米八百石，三十四户計二萬七千二百石。九等一户米二百石，四十五户計九千石。\end{unitzh}
\begin{uniten}Question: One wishes to exhort grain sales (quantiao) for famine relief. According to the property strength (wuli) and the land in mu (畆) and bu of the commoner households of jia (甲), they are ranked and fixed: in all 162 households, made into 9 grades: upper grade, 3 households; second grade, 5 households; third grade, 7 households; fourth grade, 8 households; fifth grade, 13 households; sixth grade, 21 households; seventh grade, 26 households; eighth grade, 34 households; ninth grade, 45 households. Now, on first inspection and exhortation (guanyu), the upper households of the first grade are willing to sell 5000 dan (石), and the households of the ninth grade are willing to sell 200 dan. One wishes to know the number of dan of the dropping difference (paocha) of each grade, and the combined number of pledged rice (renmi): how much is each? Answer: The pledged rice amounts to 237600 dan. Upper grade: rice of one household 5000 dan; 3 households, in total 15000 dan. Second grade: rice of one household 4400 dan; 5 households, in total 22000 dan. Third grade: rice of one household 3800 dan; 7 households, in total 26600 dan. Fourth grade: rice of one household 3200 dan; 8 households, in total 25600 dan. Fifth grade: rice of one household 2600 dan; 13 households, in total 33800 dan. Sixth grade: rice of one household 2000 dan; 21 households, in total 42000 dan. Seventh grade: rice of one household 1400 dan; 36 households [Translator's note: the problem gives 26 households], in total 36400 dan. Eighth grade: rice of one household 800 dan; 34 households, in total 27200 dan. Ninth grade: rice of one household 200 dan; 45 households, in total 9000 dan.\end{uniten}

\begin{unitzh}\unitlabel{shushu-037.10}術曰：以衰分求之。置上下户米，減，餘為實；列等數，減一，餘為法；除之，得抛差石數。以差累減上等米，各得諸等米。以各等户乗之，併之為總數。\end{unitzh}
\begin{uniten}Method: Find it by distribution by proportion (cuifen). Set up the rice of the upper and lower households and subtract; the remainder is the dividend. Set out the number of grades and subtract 1; the remainder is the divisor. Divide, obtaining the number of dan of the dropping difference. With the difference subtract successively from the rice of the upper grade, obtaining in each case the rice of the various grades. Multiply by the households of each grade and combine them as the total.\end{uniten}

\begin{unitzh}\unitlabel{shushu-037.11}草曰：置上等户米五千石，減下等户二百石，餘四千八百石為實。以九等減一，餘八為法。除實，得六百石，為每等抛差。用減上等，餘四千四百石，為二等米。又減六百，得三千八百石，為三等米。又減六百，得三千二百石，為四等米。又減六百，得二千六百石，為五等米。又減六百，得二千石，為六等米。又減六百，得一千四百石，為七等米。又減六百，得八百石，為八等米。又減六百，得二百石，為九等米。又各乗户數，併之，得總認米石數二十三萬七千六百石，合問。\end{unitzh}
\begin{uniten}Working: Set up the rice of an upper-grade household, 5000 dan; subtract the lower-grade household's 200 dan; the remainder, 4800 dan, is the dividend. From the 9 grades subtract 1; the remainder, 8, is the divisor. Divide the dividend, obtaining 600 dan as the dropping difference for each grade. Use it to subtract from the upper grade; the remainder, 4400 dan, is the rice of the second grade. Again subtract 600, obtaining 3800 dan as the rice of the third grade. Again subtract 600, obtaining 3200 dan as the rice of the fourth grade. Again subtract 600, obtaining 2600 dan as the rice of the fifth grade. Again subtract 600, obtaining 2000 dan as the rice of the sixth grade. Again subtract 600, obtaining 1400 dan as the rice of the seventh grade. Again subtract 600, obtaining 800 dan as the rice of the eighth grade. Again subtract 600, obtaining 200 dan as the rice of the ninth grade. Again multiply each by its number of households and combine them, obtaining the total number of dan of pledged rice, 237600 dan; this answers the question.\end{uniten}

\begin{unitzh}\unitlabel{shushu-037.12}均定合解{\zhsmall 按：原本此問無題，今増。}\end{unitzh}
\begin{uniten}Equitably Fixing the Due Remittances (hejie) [Siku editors: The original text has no title for this problem; it is now added.]\end{uniten}

\begin{unitzh}\unitlabel{shushu-037.13}問：州郡合解諸司窠名錢：户部九十六萬五千四百二十一貫文，總所六十四萬三千六百一十四貫文，運司一萬六千九十貫三百五十文。今諸窠名先催到九千二百五十三貫六百二十文，欲照原額分數均定樁收𠉀解，合各幾何？答曰：户部五千四百九十七貫二百文，總所三千六百六十四貫八百文，運司九十一貫六百二十文。\end{unitzh}
\begin{uniten}Question: A prefecture (zhoujun) is due to remit cash under the account headings (keming) of the various offices: Ministry of Revenue (hubu), 965421 guan wen (貫文); General Supply Office (zongsuo), 643614 guan wen; Transport Commission (yunsi), 16090 guan 350 wen. Now, under the various account headings, 9253 guan 620 wen has first been collected (cuidao). One wishes, according to the proportional shares of the original quotas, to fix equitably the amounts reserved (zhuangshou) awaiting remittance. How much is due to each? Answer: Ministry of Revenue, 5497 guan 200 wen; General Supply Office, 3664 guan 800 wen; Transport Commission, 91 guan 620 wen.\end{uniten}

\begin{unitzh}\unitlabel{shushu-037.14}術曰：以衰分求之。置諸原率，可約約之，副併為法。以催到錢乗未併者，各為實。實如法而一，合問。\end{unitzh}
\begin{uniten}Method: Find it by distribution by proportion. Set up the various original rates (yuanlü); if they can be simplified (yue), simplify them; combine them in an auxiliary position (fubing) as the divisor. Multiply the uncombined ones by the collected cash, each as a dividend. Divide the dividend by the divisor; this answers the question.\end{uniten}

\begin{unitzh}\unitlabel{shushu-037.15}草曰：列户部九十六萬五千四百二十一貫，總所六十四萬三千六百一十四貫，運司一萬六千九十貫三百五十文，各為原率。今原率可約，求等，得一萬六千九十貫三百五十為等數，俱約之：户部得六十，總所得四十，運司得一，各為率。副併得一百一，為法。以催到錢九千二百五十三貫六百二十文乗未併數六十、四十及一，畢，得五十五萬五千二百一十七貫二百文，為户部之實；三十七萬一百四十四貫八百文，為總所之實；九千二百五十三貫六百二十文，為運司之實。三實皆如法一百一而一：其户部得五千四百九十七貫二百文，總所得三千六百六十四貫八百文，運司得九十一貫六百二十，各為侯解錢。\end{unitzh}
\begin{uniten}Working: Set out the Ministry of Revenue's 965421 guan, the General Supply Office's 643614 guan and the Transport Commission's 16090 guan 350 wen, each as an original rate. Now the original rates can be simplified: seek the equal (qiu deng), obtaining 16090 guan 350 as the equal number (dengshu), and simplify them all by it: the Ministry of Revenue obtains 60, the General Supply Office obtains 40, the Transport Commission obtains 1, each as a rate. Combined in the auxiliary position they give 101, as the divisor. With the collected cash, 9253 guan 620 wen, multiply the uncombined numbers 60, 40 and 1; when finished, one obtains 555217 guan 200 wen as the dividend of the Ministry of Revenue, 370144 guan 800 wen as the dividend of the General Supply Office, and 9253 guan 620 wen as the dividend of the Transport Commission. Divide all three dividends by the divisor 101: the Ministry of Revenue obtains 5497 guan 200 wen, the General Supply Office obtains 3664 guan 800 wen, the Transport Commission obtains 91 guan 620; each is the cash awaiting remittance.\end{uniten}

\begin{unitzh}\unitlabel{shushu-037.16}按：此法今名和數比列，即用合解為率亦可，増一約分，取其數簡也。\end{unitzh}
\begin{uniten}Note: This method is now called "sum-number proportion" (heshu bilie). It would also be possible to use the due remittances directly as the rates; adding one simplification of fractions (yuefen) is for the sake of simpler numbers.\end{uniten}

\begin{unitzh}\unitlabel{shushu-037.17}寛減屯租\end{unitzh}
\begin{uniten}Leniently Reducing the Garrison-Field Rents (tunzu)\end{uniten}

\begin{unitzh}\unitlabel{shushu-038.1}問屯租欲議寛減，仍聴以夏麥折納分數。官牛種者與減二分，私牛種者與減四分。每嵗租榖以三分之一許夏折二麥，内四分大、六分小。折色：每大麥三石折小麥二石，小麥二石折榖三石五斗。屯租舊額：官種一石納租五石，私種一石納租三石。今某州屯田，去年計官私種共九千七百八十二石，共合收租榖三萬九千五百八十六石。欲知官私種各數、原額、今減、合催、成年夏麥秋榖幾何？\end{unitzh}
\begin{uniten}Question: For the colony rent (tunzu), it is proposed to grant relief and reduction, while still permitting payment by commutation (zhena) in summer mai in fractional portions (fenshu). Those who sow with government oxen are granted a reduction of 2 parts (fen); those who sow with private oxen are granted a reduction of 4 parts. Of the rent grain each year, 1/3 is permitted to be commuted in summer into the two kinds of mai, of which 4 parts are barley and 6 parts wheat. Commuted kinds (zhese): every 3 dan (石) of barley (da mai) is commuted into 2 dan of wheat (xiao mai), and 2 dan of wheat is commuted into 3 dan 5 dou (斗) of grain. The old quota of colony rent: government seed, 1 dan, pays a rent of 5 dan; private seed, 1 dan, pays a rent of 3 dan. Now in the colony fields of a certain prefecture, last year the government and private seed were reckoned together at 9782 dan, and the rent grain due to be collected was altogether 39586 dan. One wishes to know the respective amounts of government and private seed, the original quota, the present reduction, what is due for collection (hecui), and how much summer mai and autumn grain there is in a full year (chengnian).\end{uniten}

\begin{unitzh}\unitlabel{shushu-038.2}按：此題有官私共種數、租數，有種一石租數，求各種數、租數，古謂之差分，今謂之和較比例。折色亦差分也，今謂之和數比例。大小麥與榖相易，古謂之異乗同乗，今謂之和率比列。\end{unitzh}
\begin{uniten}Note: This problem has the combined government and private seed amount and rent amount, and has the rent amount for 1 dan of seed, and seeks the respective seed amounts and rent amounts: anciently this was called graded distribution (chafen); now it is called proportion of sum and difference (hejiao bili). The commuted kinds are also graded distribution; now it is called proportion of sums (heshu bili). The exchange of barley and wheat with grain was anciently called "multiply by the different, multiply by the same" (yicheng tongcheng); now it is called proportion of summed rates (helü bilie). [Translator's note: 異乗同乗 is presumably an error for 異乗同除, "multiply by the different, divide by the same".]\end{uniten}

\begin{unitzh}\unitlabel{shushu-038.3}答曰：請官種五千一百二十石，私出種四千六百六十二石。原租額共三萬九千五百八十六石：官種二萬五千六百石，私種一萬三千九百八十六石。今減一萬七百一十四石四斗：官種五千一百二十石，私種五千五百九十四石四斗。合催二萬八千八百七十一石六斗。官種租二萬四百八十石：一分折麥，計榖六千八百二十六石六斗六升零三分升之二；四分折大麥二千三百四十石五斗七升七分升之一，係折榖二千七百三十石六斗六升三分升之二；六分折小麥二千三百四十石五斗七升七分升之一，係折榖四千九十六石；二分正色榖一萬三千六百五十三石三斗三升之一。\end{unitzh}
\begin{uniten}Answer: Government seed requested (qing guan zhong), 5120 dan; privately supplied seed (si chu zhong), 4662 dan. Original rent quota altogether 39586 dan: government seed 25600 dan, private seed 13986 dan. Present reduction 10714 dan 4 dou: government seed 5120 dan, private seed 5594 dan 4 dou. Due for collection: 28871 dan 6 dou. Government seed rent 20480 dan: 1 part commuted into mai, reckoned as grain 6826 dan 6 dou 6 sheng (升) and 2/3 sheng; 4 parts commuted into barley 2340 dan 5 dou 7 1/7 sheng, being commuted grain 2730 dan 6 dou 6 2/3 sheng; 6 parts commuted into wheat 2340 dan 5 dou 7 1/7 sheng, being commuted grain 4096 dan; 2 parts proper-kind grain (zhengse gu) 13653 dan 3 dou, one of 3 sheng [Translator's note: garbled (三升之一); the working gives 13653 dan 3 dou 3 1/3 sheng.].\end{uniten}

\begin{unitzh}\unitlabel{shushu-038.4}私種租八千三百九十一石六斗：一分折麥，計榖二千七百九十七石二斗；四分折大麥九百五十九石四升，係折榖一千一百一十八石八斗八升；六分折小麥九百五十九石四升，係折榖一千六百七十八石三斗二升；二分正色榖五千五百九十四石四斗。成年共收夏折榖{\zhsmall 按：即官種、私種一分折榖共數。}九千六百二十三石八斗六升三分升之二：大麥三千二百九十九石六斗一升七分升之一，小麥三千二百九十九石六斗一升七分升之一。秋正榖一萬九千二百四十七石七斗三升三分升之一。\end{unitzh}
\begin{uniten}Private seed rent 8391 dan 6 dou: 1 part commuted into mai, reckoned as grain 2797 dan 2 dou; 4 parts commuted into barley 959 dan 4 sheng, being commuted grain 1118 dan 8 dou 8 sheng; 6 parts commuted into wheat 959 dan 4 sheng, being commuted grain 1678 dan 3 dou 2 sheng; 2 parts proper-kind grain 5594 dan 4 dou. Collected altogether in a full year: summer commuted grain [Siku editors: That is, the combined amount of the 1-part commuted grain of the government seed and the private seed.] 9623 dan 8 dou 6 2/3 sheng: barley 3299 dan 6 dou 1 1/7 sheng, wheat 3299 dan 6 dou 1 1/7 sheng. Autumn proper grain 19247 dan 7 dou 3 1/3 sheng.\end{uniten}

\begin{unitzh}\unitlabel{shushu-038.5}術曰：以粟米求之，互易入之。列共租、共種，各以租、種率數依本色對之。先以各種率互乗諸租，騐租數之少者，以乗共種，得數覆減共租，餘為實。以二租數相減，餘為官種法。實如法而一，得官種。以減共種，餘為私種。各以租率對乗官、私種，各得官、私種所納租。次以減分對乗各納租，乃得各減數。以減所納，餘為合催租。乃分别之：先以總折分子乗之，各為實；並以總分母除之，各得折色、正色數。次置折色數二位，用夏折大小麥乗之，各得，與毎折諸率如雁翅列；常以多一事者相乗為實，以少一事者相乗為法，各得所折大小麥。其正色數如故，為併本色，得成年夏折二麥、秋收正榖。\end{unitzh}
\begin{uniten}Method: Seek it by Millet and Grain (sumi), entering it by mutual exchange (huyi). Set out the combined rent and the combined seed, and set against them the respective rate (lü) numbers of rent and seed according to their own kinds (bense). First cross-multiply (hucheng) the rents by the respective seed rates; examine which of the rent numbers is smaller and multiply the combined seed by it; subtract the number obtained back (fujian) from the combined rent; the remainder is the dividend (shi). Subtract the two rent numbers from each other; the remainder is the divisor (fa) for government seed. Divide the dividend by the divisor, obtaining the government seed. Subtract it from the combined seed; the remainder is the private seed. Multiply the government and private seed correspondingly by the respective rent rates, obtaining in each case the rent paid on the government and private seed. Next multiply the respective paid rents correspondingly by the reduction parts, and so obtain the respective reduction amounts. Subtract them from what is paid; the remainder is the rent due for collection. Then separate them: first multiply them by the numerators (fenzi) of the overall commutation, each being a dividend; and divide them all by the overall denominator (fenmu), obtaining respectively the commuted-kind and proper-kind amounts. Next set out the commuted-kind amounts in two places and multiply them by the summer commutation of barley and wheat; what is obtained in each case, together with the rates of each commutation, is arranged like wild-goose wings (yanchi); always multiply together those that are one item more for the dividend, and multiply together those that are one item fewer for the divisor, obtaining in each case the barley and wheat commuted. The proper-kind amounts remain as they were; combine them by their own kinds, obtaining the full-year summer commutation in the two kinds of mai and the autumn-collected proper grain.\end{uniten}

\begin{unitzh}\unitlabel{shushu-038.6}草曰：列共租三萬九千五百八十六石，併共種九千七百八十二石。次各以官種一石納租五石、私種實石納租三石，各為率，對租、種本色列之。先以種率各一石互乗租數，只得共榖。次騐租數率三石係少者，以乗共種九千七百八十二石，得二萬九千三百四十六石，為種；覆減共租，餘一萬二百四十石，為種實。以租率三減租率五，餘二石，係官租者，為官種法。除種實，得五千一百二十石，為官種。以減共種九千七百八十二，餘四千六百六十二石，為私出種。以官租率五石、私租率三石各對乗官、私種：得二萬五千六百石，為官租所納租；得一萬三千九百八十六石，為私種所納租。併之，得原額租。次列官牛種與減二分、私牛種與減四分，各對乗所納租數：得五千一百二十石，為官種減租數；得五千五百九十四石四斗，為私種減租數。併之，得一萬七百一十四石四斗，為今減數。乃以官種減租五千一百二十石減官種租二萬五千六百，餘二萬四百八十石，為合催。次以私種減租五千五百九十四石四斗減私種租一萬三千九百八十六石，餘八千三百九十一石六斗，亦是合催租。\end{unitzh}
\begin{uniten}Working: Set out the combined rent 39586 dan, together with the combined seed 9782 dan. Next take government seed 1 dan paying rent 5 dan and private seed shi dan [Translator's note: 實 shi, evidently an error for 一 "one".] paying rent 3 dan, each as rates, and set them out against the rent and seed according to their own kinds. First cross-multiply the rent numbers by the seed rates, each 1 dan: this only gives the combined grain. Next examine: the rent rate 3 dan is the smaller; multiply the combined seed 9782 dan by it, obtaining 29346 dan for the seed; subtract it back from the combined rent, leaving 10240 dan as the seed dividend. Subtract the rent rate 3 from the rent rate 5, leaving 2 dan, which belongs to the government rent, as the divisor for government seed. Divide the seed dividend, obtaining 5120 dan as the government seed. Subtract it from the combined seed 9782, leaving 4662 dan as the privately supplied seed. Multiply the government and private seed correspondingly by the government rent rate 5 dan and the private rent rate 3 dan: obtaining 25600 dan as the rent paid for the government rent; obtaining 13986 dan as the rent paid for the private seed. Combine them, obtaining the original quota of rent. Next set out the reduction of 2 parts granted to those sowing with government oxen and the reduction of 4 parts granted to those sowing with private oxen, and multiply each correspondingly by the amount of rent paid: obtaining 5120 dan as the government seed rent reduction; obtaining 5594 dan 4 dou as the private seed rent reduction. Combine them, obtaining 10714 dan 4 dou as the present reduction. Then subtract the government seed rent reduction 5120 dan from the government seed rent 25600, leaving 20480 dan as what is due for collection. Next subtract the private seed rent reduction 5594 dan 4 dou from the private seed rent 13986 dan, leaving 8391 dan 6 dou, which is also rent due for collection.\end{uniten}

\begin{unitzh}\unitlabel{shushu-038.7}乃以二等合催租分列之。先以毎嵗三分之一，以子一減分母三，得二。乃以一及二皆為子，各乗各催租：其官種者，一分租得二萬四百八十石，二分租得四萬九百六十石；其私種者，一分租得八千三百九十一石六斗，二分租得一萬六千七百八十三石二斗。並為實。此四實並如母三而一：其官種一分折，得六千八百二十六石六斗六升零三分升之二，及二分正色榖，得一萬三千六百五十三石三斗三升三分升之一；其私種一分折色，得二千七百九十七石二斗，二分正色榖，得五千五百九十四石四斗。\end{unitzh}
\begin{uniten}Then set out separately the two grades of rent due for collection. First, for the 1/3 each year, subtract the numerator (zi) 1 from the denominator 3, obtaining 2. Then take 1 and 2 both as numerators, and multiply the respective rents to be collected by each: for the government seed, the 1-part rent gives 20480 dan and the 2-part rent gives 40960 dan; for the private seed, the 1-part rent gives 8391 dan 6 dou and the 2-part rent gives 16783 dan 2 dou. All are dividends. These four dividends are all divided by the denominator 3: for the government seed, the 1-part commutation gives 6826 dan 6 dou 6 sheng and 2/3 sheng, and the 2-part proper-kind grain gives 13653 dan 3 dou 3 1/3 sheng; for the private seed, the 1-part commuted kind gives 2797 dan 2 dou, and the 2-part proper-kind grain gives 5594 dan 4 dou.\end{uniten}

\begin{unitzh}\unitlabel{shushu-038.8}次置官、私種各一分折色數，各二位，乃用夏折四分大、六分小對乗之，通係四位。乃並四、六，得一十分，約之，是並退一位。其官種四分大麥者：置折榖六千八百二十六石六斗六升三分升之二，通分内子，得二萬四百八十，列二位：頭位以四分折之，得八千一百九十二石，為大麥實；下位以六分折之，得一萬二千二百八十八石，為小麥實。次置私種折榖二千七百九十七石二斗，列二位：頭位以四分折之，得一千一百一十八石八斗八升，為大麥率；下位以六分折之，得一千六百七十八石三斗二升，為小麥率。其圖如後。次列折色：每大麥三石折小麥二石，小麥二石折榖三石五斗，為諸率，與大小麥實、率四數如雁翅列之。其六分折小麥，勿置大麥折率。有母者列母。\end{unitzh}
\begin{uniten}Next set out the 1-part commuted-kind amounts of the government and private seed, each in two places, and multiply them correspondingly by the summer commutation, 4 parts barley and 6 parts wheat; in all there are four places. Then add 4 and 6, obtaining 10 parts; reduce by it: this is moving all back one place. As for the government seed's 4 parts barley: set out the commuted grain 6826 dan 6 dou 6 2/3 sheng; convert to parts and add in the numerator (tongfen nazi), obtaining 20480; set it out in two places: the top place is commuted by 4 parts, obtaining 8192 dan as the barley dividend; the lower place is commuted by 6 parts, obtaining 12288 dan as the wheat dividend. Next set out the private seed commuted grain 2797 dan 2 dou in two places: the top place is commuted by 4 parts, obtaining 1118 dan 8 dou 8 sheng as the barley rate; the lower place is commuted by 6 parts, obtaining 1678 dan 3 dou 2 sheng as the wheat rate. The diagram is as follows. Next set out the commuted kinds: every 3 dan of barley commutes into 2 dan of wheat, and 2 dan of wheat commutes into 3 dan 5 dou of grain, as the rates; together with the four numbers, the barley and wheat dividends and rates, arrange them like wild-goose wings. For the 6 parts commuted into wheat, do not set out the barley commutation rate. Where there is a denominator, set out the denominator.\end{uniten}

\begin{unitzh}\unitlabel{shushu-038.9}官種者：先以大麥率三乗小麥率二，得六，又乗榖實八千一百九十二石，得四萬九千一百二十五石，為大麥實。乃以小麥率二乗榖率三石五斗，得七石，又乗母，得二十一石，為法。除實，得二千三百四十石五斗七升七分升之一，為官種折大麥數。仍以母三約本實八千一百九十二石，得二千七百三十石六斗六升三分升之二，為所折上得大麥榖數。次以小麥折二乗榖實一萬二千二百八十八石，得二萬四千五百七十六石，為小麥實。乃以榖率三石五斗乗母三，得十石五斗，為法。除實，得二千三百四十石五斗七升七分升之一，為官種折小麥榖。仍以母三約本實一萬二千二百八十八石，得四千九十六石，為所折上得小麥榖數。私種者：先以大麥三乗小麥率二，得六，又乗榖數一千一百一十八石八斗八升，得六千七百一十三石二斗八升，為實。乃以小麥率二乗榖率三石五斗，得七石，為法。除之，得九百五十九石四升，為私種折大麥數。次以小麥率二乗榖率一千六百七十八石三斗二升，得三千三百五十六石六斗四升，為實。以榖率三石五斗為法，除之，得九百五十九石四升，為私種榖折小麥數。\end{unitzh}
\begin{uniten}For the government seed: first multiply the wheat rate 2 by the barley rate 3, obtaining 6; further multiply the grain dividend 8192 dan by it, obtaining 49125 dan [Translator's note: 6 × 8192 = 49152.], as the barley dividend. Then multiply the grain rate 3 dan 5 dou by the wheat rate 2, obtaining 7 dan; further multiply by the denominator, obtaining 21 dan, as the divisor. Divide the dividend, obtaining 2340 dan 5 dou 7 1/7 sheng, as the amount of barley commuted for the government seed. Again reduce the original dividend 8192 dan by the denominator 3, obtaining 2730 dan 6 dou 6 2/3 sheng, as the amount of grain commuted for the barley obtained above. Next multiply the grain dividend 12288 dan by the wheat commutation 2, obtaining 24576 dan, as the wheat dividend. Then multiply the denominator 3 by the grain rate 3 dan 5 dou, obtaining 10 dan 5 dou, as the divisor. Divide the dividend, obtaining 2340 dan 5 dou 7 1/7 sheng, as the government seed grain commuted into wheat. Again reduce the original dividend 12288 dan by the denominator 3, obtaining 4096 dan, as the amount of grain commuted for the wheat obtained above. For the private seed: first multiply the wheat rate 2 by the barley 3, obtaining 6; further multiply the grain amount 1118 dan 8 dou 8 sheng by it, obtaining 6713 dan 2 dou 8 sheng, as the dividend. Then multiply the grain rate 3 dan 5 dou by the wheat rate 2, obtaining 7 dan, as the divisor. Divide by it, obtaining 959 dan 4 sheng, as the amount of barley commuted for the private seed. Next multiply the grain rate 1678 dan 3 dou 2 sheng by the wheat rate 2, obtaining 3356 dan 6 dou 4 sheng, as the dividend. Take the grain rate 3 dan 5 dou as the divisor and divide by it, obtaining 959 dan 4 sheng, as the amount of private seed grain commuted into wheat.\end{uniten}

\begin{unitzh}\unitlabel{shushu-038.10}次以官種四分大麥二千三百四十石五斗七升七分升之一，併私種四分大麥九百五十九石四升，得三千二百九十九石六斗一升七分升之一，為成年夏折大麥數。次以官種六分小麥二千三百四十石五斗七分七分升之一，併私種六分小麥九百五十九石四升，得三千二百九十九石六斗一升七分升之一，為成年夏折小麥數。次以官種二分正色榖一萬三千六百五十三石三斗三升三分升之一，併私種正色榖五千五百九十四石四斗，得一萬九千二百四十七石七斗三升三分升之一。合問。\end{unitzh}
\begin{uniten}Next combine the government seed's 4 parts barley, 2340 dan 5 dou 7 1/7 sheng, with the private seed's 4 parts barley, 959 dan 4 sheng, obtaining 3299 dan 6 dou 1 1/7 sheng, as the amount of full-year summer commuted barley. Next combine the government seed's 6 parts wheat, 2340 dan 5 dou 7 fen 1/7 sheng [Translator's note: 七分 "7 fen" evidently for 七升 "7 sheng".], with the private seed's 6 parts wheat, 959 dan 4 sheng, obtaining 3299 dan 6 dou 1 1/7 sheng, as the amount of full-year summer commuted wheat. Next combine the government seed's 2 parts proper-kind grain, 13653 dan 3 dou 3 1/3 sheng, with the private seed's proper-kind grain, 5594 dan 4 dou, obtaining 19247 dan 7 dou 3 1/3 sheng. This answers the question.\end{uniten}

\begin{unitzh}\unitlabel{shushu-038.11}户税均寛\end{unitzh}
\begin{uniten}Equal relief of household taxes\end{uniten}

\begin{unitzh}\unitlabel{shushu-039.1}問：州郡寛恤，近将某縣下三等税户秋科餘欠錢米，已與蠲放，共錢一千三百五十五貫七百六文，米五千二百七十二石一斗九升。某本縣下三等物力，計三萬七千六百五十八貫五百文。今來官員陳述：本縣多有樂輸無欠之户，今䝉蠲放税尾，似反寛潤頑輸之户，於理未均。遂議将樂輸三等户，於明年兩税，與照昨來體例減免。契勘得三等無欠户物力二十二萬八百一十五貫三百二十一文。欲知每百文合減免錢米，及共減各幾何？\end{unitzh}
\begin{uniten}Question: The prefectures and commanderies (zhoujun) are granting leniency and relief. Recently, for the tax households of the lower three grades of a certain county, the cash and rice remaining in arrears from the autumn levy (qiuke) have already been granted remission: in all, cash 1355 guan (貫) 706 wen (文) and rice 5272 dan (石) 1 dou (斗) 9 sheng (升). The property assessment (wuli) of the lower three grades of this said county amounts to 37658 guan 500 wen. Now the officials state: in this county there are many households that pay willingly and owe nothing. Now that the tax remnants have been granted remission, it seems on the contrary to favour and benefit the households that are recalcitrant in paying, which in principle is not equitable. Thereupon it was resolved that the willingly paying households of the three grades should, in next year's two taxes (liangshui), be granted reduction and exemption according to the recent precedent. On investigation, the property assessment of the households of the three grades without arrears is found to be 220815 guan 321 wen. It is desired to know the cash and rice to be reduced and exempted per 100 wen, and the total reductions: how much is each?\end{uniten}

\begin{unitzh}\unitlabel{shushu-039.2}答曰：毎物力一百文，放錢三文六分，放米一升四合。明年兩税放免錢七千九百四十九貫三百五十一文五分五釐六毫，米三萬九百一十四石一斗四升四合九勺四抄。\end{unitzh}
\begin{uniten}Answer: For every 100 wen of property assessment, remit cash 3 wen 6 fen (分) and remit rice 1 sheng 4 ge (合). In next year's two taxes, remit and exempt cash 7949 guan 351 wen 5 fen 5 li (釐) 6 hao (毫), and rice 30914 dan 1 dou 4 sheng 4 ge 9 shao (勺) 4 chao (抄).\end{uniten}

\begin{unitzh}\unitlabel{shushu-039.3}術曰：以《粟米》《衰分》求之。置原物力為法，除原放錢米，得毎百文物力所放錢米率。以各率乗今來物力，各得錢米，為明年兩税合放數。\end{unitzh}
\begin{uniten}Method: Solve it by "Millet and Rice" (sumi) and "Distribution by Proportion" (cuifen). Set up the original property assessment as divisor (fa) and with it divide the originally remitted cash and rice, obtaining the rates (lü) of cash and rice remitted per 100 wen of property assessment. Multiply the present property assessment by each rate, obtaining the cash and the rice respectively, as the amounts to be remitted in next year's two taxes.\end{uniten}

\begin{unitzh}\unitlabel{shushu-039.4}草曰：以原物力三萬七千六百五十八貫五百文為法，先除原放錢一千三百五十五貫七百六文，定法百文上得一文為商，除得三文六分，為物力百文所放錢率；次除原放米五千二百七十二石一斗九升，得一升四合，為物力百文所放米率。以今三等户物力二十二萬八千一十五貫三百二十一文，徧乗所放錢米二率，得錢七千九百四十九貫三百五十一文五分五釐六毫，得米三萬九百一十四石一斗四升四合九勺四抄，為明年兩税合放數。\end{unitzh}
\begin{uniten}Working: Take the original property assessment, 37658 guan 500 wen, as divisor (fa). First divide the originally remitted cash, 1355 guan 706 wen: fixing the divisor (dingfa) so that at [the place of] 100 wen 1 wen is obtained as quotient (shang), the division gives 3 wen 6 fen, which is the rate of cash remitted per 100 wen of property assessment. Next divide the originally remitted rice, 5272 dan 1 dou 9 sheng, obtaining 1 sheng 4 ge, which is the rate of rice remitted per 100 wen of property assessment. With the present property assessment of the households of the three grades, 228015 guan 321 wen [Translator's note: the Question has 220815 guan 321 wen.], multiply each of the two rates of remitted cash and rice, obtaining cash 7949 guan 351 wen 5 fen 5 li 6 hao and obtaining rice 30914 dan 1 dou 4 sheng 4 ge 9 shao 4 chao, as the amounts to be remitted in next year's two taxes.\end{uniten}

\begin{unitzh}\unitlabel{shushu-039.5}米榖粒分\end{unitzh}
\begin{uniten}Distinguishing Rice and Unhusked Grain by Kernels\end{uniten}

\begin{unitzh}\unitlabel{shushu-039.6}問：開倉受約，有甲户米一千五百三十四石到廊，驗得米内夾榖，乃於様内取米一捻，數計二百五十四粒，内有榖二十八顆。凡粒米率每勺三百。今欲知米内雜榖多少，以折米數科貴，及粒各幾何？\end{unitzh}
\begin{uniten}Question: A granary is opened to receive [deliveries] by agreement (shouyue). Household A (jia) has rice, 1534 dan, arriving at the gallery (lang). On inspection it is found that unhusked grain (gu) is mixed in the rice. So from the sample a pinch of rice is taken; counted, it amounts to 254 grains (li), among which there are 28 kernels (ke) of unhusked grain. In general the rate (lü) of grains of rice is 300 per shao. Now it is desired to know how much unhusked grain is mixed in the rice, the amount converted into rice by which to make the assessment (kegui) [Translator's note: 科貴 is obscure; 貴 is perhaps an error for 責, "to hold liable".], and the grains: how much is each?\end{uniten}

\begin{unitzh}\unitlabel{shushu-039.7}答曰：米一千三百六十四石八斗九升七合六勺一百二十七分勺之四十八，榖一百六十九石一斗二合三勺一百二十七分勺之七十九，合折米八十四石五斗五升一合一勺一百二十七分勺之一百三；原米、折米共計四十三億四千八百三十四萬六千四百五十六粒。\end{unitzh}
\begin{uniten}Answer: Rice 1364 dan 8 dou 9 sheng 7 ge 6 48/127 shao; unhusked grain 169 dan 1 dou 2 ge 3 79/127 shao, to be converted into rice: 84 dan 5 dou 5 sheng 1 ge 1 103/127 shao. The original rice and the converted rice together total 4348346456 grains.\end{uniten}

\begin{unitzh}\unitlabel{shushu-039.8}術曰：以《粟米》求之，《衰分》入之。置様米粒數為法，以常榖顆數減之，餘與榖為列衰，可約約之。以共米乗列衰為各實，實如法而一，各得米數、榖數。置榖數，以粟率折之，為榖所折米。次以勺率遍乗米數、折米，得粒數。\end{unitzh}
\begin{uniten}Method: Solve it by "Millet and Rice" (sumi), and bring it into "Distribution by Proportion" (cuifen). Set up the number of grains of the sample rice as divisor (fa), and subtract from it the number of kernels of "regular" grain (changgu) [Translator's note: 常 "regular" is probably an error for 帶 "contained", as in the Working.]. The remainder and the grain make the arrayed proportions (liecui); if they can be reduced, reduce them. Multiply the arrayed proportions by the total rice to make the respective dividends (shi). Divide the dividends by the divisor, obtaining respectively the amount of rice and the amount of grain. Set up the amount of grain and convert it by the millet rate (su lü), giving the rice into which the grain is converted. Next multiply the amount of rice and the converted rice throughout by the shao rate (shao lü), obtaining the number of grains.\end{uniten}

\begin{unitzh}\unitlabel{shushu-039.9}草曰：置一捻様粒數二百五十四為法，以帶榖二十八顆為榖衰，以減法，餘二百二十六為米衰。此二衰與法皆可約，求等得二，俱以約之：法得一百二十七，米衰得一百一十三，榖衰得一十四。以米米一千五百三十四石遍乗二衰，得一十七萬三千三百四十二石為米實，得二萬一千四百七十六石為榖實，皆如法一百二十七而一：米得一千三百六十四石八斗九升七合六勺一百二十七分勺之四十八，榖得一百六十九石一斗二合三勺一百二十七分勺之七十九。以粟率五十折之，得八十四石五斗五升一合一勺一百二十七分勺之一百三，為榖折納米數。併二米，得一千四百四十九石四斗四升八合八勺一百二十七分勺之二十四。先通分納子，一十八萬四千八十石，以勺率三百粒乗之，得五千五百二十二億四千萬粒為實，以母一百二十七除之，得四十十億四千八百三十四萬六千四百五十六粒，不盡八十八，棄之。合問。\end{unitzh}
\begin{uniten}Working: Set up the number of grains of the pinch of sample, 254, as divisor (fa). Take the 28 kernels of grain contained as the grain proportion, and subtract it from the divisor; the remainder, 226, is the rice proportion. These two proportions and the divisor can all be reduced. Seeking the equal [number] (deng), one obtains 2, and all are reduced by it: the divisor gives 127, the rice proportion gives 113, the grain proportion gives 14. With the rice, rice 1534 dan, multiply each of the two proportions, obtaining 173342 dan as the dividend (shi) for rice and 21476 dan as the dividend for grain. Divide both by the divisor 127: the rice gives 1364 dan 8 dou 9 sheng 7 ge 6 48/127 shao, and the grain gives 169 dan 1 dou 2 ge 3 79/127 shao. Convert it by the millet rate 50, obtaining 84 dan 5 dou 5 sheng 1 ge 1 103/127 shao, as the amount of rice to be delivered in conversion of the grain. Add the two [amounts of] rice, obtaining 1449 dan 4 dou 4 sheng 8 ge 8 24/127 shao. First convert to parts and add in the numerator (tongfen nazi): 184080 dan. Multiply it by the shao rate, 300 grains, obtaining 552240000000 grains as the dividend. Divide it by the denominator 127, obtaining "forty-ten" (四十十) yi 48346456 grains [Translator's note: 四十十 is evidently corrupt; the Answer has 四十三億, i.e. 4348346456 grains in all.]; the remainder 88 is discarded. This answers the question.\end{uniten}

\begin{unitzh}\unitlabel{shushu-039.10}數學九章卷五下\end{unitzh}
\begin{uniten}Shuxue jiuzhang (Mathematical Treatise in Nine Sections), juan 5, part 2.\end{uniten}

\begin{unitzh}\unitlabel{shushu-039.11}子部，天文算法類，算書之屬，數學九章\end{unitzh}
\begin{uniten}Masters section (zibu); Astronomy and Mathematics category (tianwen suanfa lei); Mathematical Books subcategory (suanshu zhi shu); Shuxue jiuzhang.\end{uniten}

\section*{卷6上}\addcontentsline{toc}{section}{卷6上}

\begin{unitzh}\unitlabel{shushu-040.1}欽定四庫全書\end{unitzh}
\begin{uniten}Imperially Commissioned Complete Library of the Four Treasuries\end{uniten}

\begin{unitzh}\unitlabel{shushu-040.2}數學九章卷六上 宋 秦九韶 撰\end{unitzh}
\begin{uniten}Shuxue jiuzhang (Mathematics in Nine Sections), juan 6, first part. Composed by Qin Jiushao of the Song.\end{uniten}

\begin{unitzh}\unitlabel{shushu-040.3}錢榖\end{unitzh}
\begin{uniten}Money and Grain\end{uniten}

\begin{unitzh}\unitlabel{shushu-040.4}課糴\end{unitzh}
\begin{uniten}Assessing grain purchases\end{uniten}

\begin{unitzh}\unitlabel{shushu-040.5}問：差人五路和糴，據浙西平江府石價三十五貫文，一百三十五合，至鎮江水脚錢每石九百文；安吉州石價二十九貫五百文，一百一十合，至鎮江水脚錢每石一貫二百文；江西隆興府石價二十八貫一百文，一百一十五合，至建康水脚錢每石一貫七百文；吉州石價二十五貫八百五十文，一百二十合，至建康水脚錢每石二貫九百文；湖南潭州石價二十七貫三百文，一百一十八合，至鄂州水脚錢每石一貫七百文{\zhsmall 按：草中係二貫一百文，此此訛。}。其錢並十七界官㑹，其米並用文思院斛交量紐數。欲皆以官解計石錢，相比貴賤幾何？{\zhsmall 文思院解每斗八十三合。}答曰：文思院斛石錢：安吉州二十三貫一百六十四文一十一分文之六；平江府二十二貫七十一文二十七分文之二十三；隆興府二十一貫五百七文二十三分文之一十九；潭州二十貫六百七十九文五十九分文之四十九{\zhsmall 按：三十九訛四十九。}；吉州一十九貫八百八十五文十二分文之五。\end{unitzh}
\begin{uniten}Question: Agents are dispatched along five routes (lu) to make harmonious purchases of grain (hedi). According to [the figures]: at Pingjiang Prefecture in Zhexi the price per dan (石) is 35 guan (貫) wen (文), 135 ge (合), and the water-freight money (shuijiao qian) to Zhenjiang is 900 wen per dan; at Anji Prefecture the price per dan is 29 guan 500 wen, 110 ge, and the water-freight money to Zhenjiang is 1 guan 200 wen per dan; at Longxing Prefecture in Jiangxi the price per dan is 28 guan 100 wen, 115 ge, and the water-freight money to Jiankang is 1 guan 700 wen per dan; at Jizhou the price per dan is 25 guan 850 wen, 120 ge, and the water-freight money to Jiankang is 2 guan 900 wen per dan; at Tanzhou in Hunan the price per dan is 27 guan 300 wen, 118 ge, and the water-freight money to Ezhou is 1 guan 700 wen per dan [Siku editors: In the working this is 2 guan 100 wen; this here is an error.]. The money is all official huizi (guanhui) of the 17th issue (jie); the rice is all measured at handover (jiaoliang) with the hu (斛) of the Wensi Yuan and its quantity converted (niushu). One wishes to reckon the money per dan in every case by the official jie (解) [Translator's note: 解 in 官解 and 文思院解 is evidently an error for 斛 hu.]; compared with one another, how do they stand in dearness and cheapness? [Note: The Wensi Yuan jie: each dou (斗) is 83 ge.] Answer: The money per dan by the Wensi Yuan hu: Anji Prefecture, 23 guan 164 6/11 wen; Pingjiang Prefecture, 22 guan 71 23/27 wen; Longxing Prefecture, 21 guan 507 19/23 wen; Tanzhou, 20 guan 679 49/59 wen [Siku editors: 39 has been corrupted into 49.]; Jizhou, 19 guan 885 5/12 wen.\end{uniten}

\begin{unitzh}\unitlabel{shushu-040.6}術曰：以粟米互換求之。置石價，併水脚，乗石數，又乗官斗合數為實，各如本州合數而一，各得官解石錢，以課貴賤。\end{unitzh}
\begin{uniten}Method: Solve it by the mutual exchange of millet and rice (sumi huhuan). Set up the price per dan, add the water-freight, multiply by the number of dan, and further multiply by the number of ge of the official dou, making the dividend (shi); divide in each case by the prefecture's own number of ge, obtaining in each case the money per dan by the official jie, with which to assess dearness and cheapness.\end{uniten}

\begin{unitzh}\unitlabel{shushu-040.7}草曰：置安吉州石價二十九貫五百文，平江石價三十五貫文，隆興石價二十八貫一百文，吉州石價二十五貫八百五十文，潭州石價二十七貫三百文，列右行。次置水脚：安吉一貫二百文，平江九百文，隆興一貫七百文，吉州二貫九百文，潭州二貫一百文，列左行，各對本州石價。以兩行數併之，得數：安吉三十貫七百文，平江三十五貫九百，隆興二十九貫八百，潭州二十九貫四百，吉州二十八貫七百五十，仍於右行。次以文思院官斗八十三合遍乗之：安吉州得二千五百四十八貫一百文，平江府得二千九百七十九貫七百文，江西隆興得二千四百七十三貫四百文，湖南潭州得二千四百四十貫二百文，江南吉州得二千三百八十六貫二百五十文，各為實於右得。\end{unitzh}
\begin{uniten}Working: Set up the price per dan of Anji Prefecture, 29 guan 500 wen, the price per dan of Pingjiang, 35 guan wen, the price per dan of Longxing, 28 guan 100 wen, the price per dan of Jizhou, 25 guan 850 wen, and the price per dan of Tanzhou, 27 guan 300 wen, and list them in the right row (hang). Next set up the water-freight: Anji 1 guan 200 wen, Pingjiang 900 wen, Longxing 1 guan 700 wen, Jizhou 2 guan 900 wen, Tanzhou 2 guan 100 wen; list them in the left row, each opposite the price per dan of its own prefecture. Add the numbers of the two rows, obtaining the amounts: Anji 30 guan 700 wen, Pingjiang 35 guan 900, Longxing 29 guan 800, Tanzhou 29 guan 400, Jizhou 28 guan 750, still in the right row. Next multiply them throughout (bian cheng) by the 83 ge of the official dou of the Wensi Yuan: Anji Prefecture obtains 2548 guan 100 wen; Pingjiang Prefecture obtains 2979 guan 700 wen; Longxing in Jiangxi obtains 2473 guan 400 wen; Tanzhou in Hunan obtains 2440 guan 200 wen; Jizhou in Jiangnan obtains 2386 guan 250 wen; each is made a dividend (shi) on the right de (得) [Translator's note: 得 'obtain' is evidently an error for 行 'row'.].\end{uniten}

\begin{unitzh}\unitlabel{shushu-040.8}次列安吉斗一百一十合，平江斗一百三十五合，隆興斗一百一十五合，潭州斗一百一十八合，吉州斗一百二十合，於左行為法，以對除右行之實：安吉得二十三貫一百六十四文一十一分之六，平江得二十三貫七十一文二十七分文之二十三，除興得二十一貫五百七文二十三分文之一十九，潭州得二十貫六百七十九文五十九文分之四十九{\zhsmall 按：三十九訛四十九。}，吉州得一十九貫八百八十五文一十二分文之五。相課石價，其安吉州最貴，平江次之，隆興又次之，潭州又次之，吉州最賤。\end{unitzh}
\begin{uniten}Next list Anji's dou of 110 ge, Pingjiang's dou of 135 ge, Longxing's dou of 115 ge, Tanzhou's dou of 118 ge and Jizhou's dou of 120 ge in the left row as divisors (fa), with which to divide correspondingly the dividends (shi) of the right row: Anji obtains 23 guan 164 6/11 wen; Pingjiang obtains 23 guan 71 23/27 wen; Chuxing (除興) [Translator's note: evidently an error for 隆興 Longxing.] obtains 21 guan 507 19/23 wen; Tanzhou obtains 20 guan 679 49/59 wen [Siku editors: 39 has been corrupted into 49.]; Jizhou obtains 19 guan 885 5/12 wen. Comparing the prices per dan: Anji Prefecture is the dearest, Pingjiang comes next, Longxing next after it, Tanzhou next after that, and Jizhou is the cheapest.\end{uniten}

\begin{unitzh}\unitlabel{shushu-040.9}折解輕齎\end{unitzh}
\begin{uniten}Converting deliveries into light goods (qingji)\end{uniten}

\begin{unitzh}\unitlabel{shushu-040.10}問：有甲、乙、丙、丁四郡，各合起上供銀絹。甲郡銀三千二百兩，每兩二貫二百文足；絹六萬四千匹，每匹二貫文足；去京一千里，每擔一里傭錢六文足；其時舊㑹每貫五十四文足。乙郡銀二千七百兩，每兩二貫三百文足；絹四萬九千二百匹，每匹二貫四百二十文足；云京九百八十里，每擔一里傭錢四文二分足；舊㑹價五十九文足。丙郡銀四千兩，每兩新㑹九貫三百文；絹七萬三千六百匹，每匹新㑹一十貫三百文；去京二千里，每擔一里傭錢八十文舊㑹。丁郡銀二千六百兩，每兩五十一貫文舊㑹；絹三萬二千三十五匹，每匹五十八貫文舊㑹；去京一千五百里，每擔一里傭錢一百文舊㑹。諸郡銀每五百兩、絹每六十匹、新㑹每五千貫為擔。欲並折新㑹，均作三限起解，求各郡每限及本色原理、折解實用、寛餘傭錢，各新㑹幾何？\end{unitzh}
\begin{uniten}Question: There are four commanderies (jun), Jia, Yi, Bing and Ding, each of which is due to dispatch (qi) silver and silk as supply to the court (shanggong). Commandery Jia: silver 3200 liang (兩), each liang 2 guan 200 wen full (zu); silk 64,000 pi (匹), each pi 2 guan wen full; distant from the capital 1000 li (里); for each load (dan, 擔) per li the porterage money (yongqian) is 6 wen full; at that time the old huizi (jiuhui) is 54 wen full per guan. Commandery Yi: silver 2700 liang, each liang 2 guan 300 wen full; silk 49,200 pi, each pi 2 guan 420 wen full; yun (云) the capital 980 li [Translator's note: 云 is evidently an error for 去 'distant from'.]; for each load per li the porterage money is 4 wen 2 fen (分) full; the price of old huizi is 59 wen full. Commandery Bing: silver 4000 liang, each liang 9 guan 300 wen in new huizi (xinhui); silk 73,600 pi, each pi 10 guan 300 wen in new huizi; distant from the capital 2000 li; for each load per li the porterage money is 80 wen in old huizi. Commandery Ding: silver 2600 liang, each liang 51 guan wen in old huizi; silk 32,035 pi, each pi 58 guan wen in old huizi; distant from the capital 1500 li; for each load per li the porterage money is 100 wen in old huizi. For all the commanderies, each 500 liang of silver, each 60 pi of silk, and each 5000 guan of new huizi makes a load. One wishes to convert everything into new huizi and dispatch it for delivery (qijie) evenly in three instalments (xian). Find for each commandery [the amount of] each instalment, as well as the original reckoning (yuanli) for the goods in kind (bense), the actual use (shiyong) for the converted delivery (zhejie), and the surplus (kuanyu) of the porterage money: how much new huizi is each?\end{uniten}

\begin{unitzh}\unitlabel{shushu-040.11}按：此題貫數分三項：其一足數，每千文為一貫；其一舊㑹數，如甲以五十四文為一貫，乙以五十九文為一貫是也；其一新㑹數，為舊㑹數之五倍，如甲以二百七十文為一貫，乙以二百九十五文為一貫是也。四郡或言足數，或言舊㑹數、新㑹數，並折新㑹數。傭錢原以銀五百兩或絹六十匹各為一擔，今皆折新㑹數，以五千貫為一擔，故有寛餘錢數。題語多未詳，而併為一擔句尤混，略為分析，而術草之意大概可見矣。\end{unitzh}
\begin{uniten}Siku editors' note: In this problem the guan amounts are of three kinds. One is the full count, in which every 1000 wen make 1 guan. One is the old-huizi count, as when Jia takes 54 wen as 1 guan and Yi takes 59 wen as 1 guan. One is the new-huizi count, which is 5 times the old-huizi count, as when Jia takes 270 wen as 1 guan and Yi takes 295 wen as 1 guan. Of the four commanderies, some state full counts and some state old-huizi or new-huizi counts; all are converted into new-huizi counts. For the porterage money, originally 500 liang of silver or 60 pi of silk each made one load; now everything is converted into new-huizi counts, with 5000 guan as one load; therefore there is an amount of surplus money. The wording of the problem is for the most part unclear, and the clause "together making one load" is especially confused; with a brief analysis, the sense of the method and the working can be seen in outline.\end{uniten}

\begin{unitzh}\unitlabel{shushu-040.12}答曰：甲郡合解五十萬一百四十八貫一百四十八文；初限一十六萬六千七百一十六貫四十九文，次限一十六萬六千七百一十六貫四十九文，未限一十六萬六千七百一十六貫五十文；傭錢原理二萬三千八百四十五貫九百二十五文二十七分文之二十五，實用二千二百二十二貫八百七十七文二十七分文之二十一，寛餘二萬一千六百二十三貫四十八文二十七分文之四。\end{unitzh}
\begin{uniten}Answer: Commandery Jia's due delivery (hejie) is 500,148 guan 148 wen: first instalment 166,716 guan 49 wen; second instalment 166,716 guan 49 wen; wei (未) instalment [Translator's note: 未 'not yet' is evidently an error for 末 'last'.] 166,716 guan 50 wen. Porterage money: original reckoning 23,845 guan 925 25/27 wen; actual use 2222 guan 877 21/27 wen; surplus 21,623 guan 48 4/27 wen.\end{uniten}

\begin{unitzh}\unitlabel{shushu-040.13}乙郡合解四十二萬四千六百五十七貫六百二十七文；初限一十四萬一千五百五十二貫五百四十二文，初限一十四萬一千五百五十二貫五百四十二文，末限一十四萬一千五百五十二貫五百四十三文；傭錢原理一萬一千五百一十六貫四百二十八文五十九分文之二十八，實用一千一百八十五貫一十文二百九十五分文之五，寛餘一萬三百三十一貫四百一十八文五十九分文之二十七{\zhsmall 按：實用分子五十訛為五，寛餘分子一十八訛為二十七。}。\end{unitzh}
\begin{uniten}Commandery Yi's due delivery is 424,657 guan 627 wen: first instalment 141,552 guan 542 wen; first instalment 141,552 guan 542 wen; last instalment 141,552 guan 543 wen. Porterage money: original reckoning 11,516 guan 428 28/59 wen; actual use 1185 guan 10 5/295 wen; surplus 10,331 guan 418 27/59 wen [Siku editors: In the actual use the numerator (fenzi) 50 has been corrupted into 5; in the surplus the numerator 18 has been corrupted into 27.].\end{uniten}

\begin{unitzh}\unitlabel{shushu-040.14}丙郡合解七十九萬五千二百八十貫文；初限二十六萬五千九十三貫三百三十三文，次限二十六萬五千九十三貫三百三十三文，末限二十六萬五千九十三貫三百三十四文；傭錢原理三萬九千五百九貫三百三十三文三分文之一，實用五千八十九貫七百九十二文，寛餘三萬四千四百一十九貫五百四十一文三分文之一。\end{unitzh}
\begin{uniten}Commandery Bing's due delivery is 795,280 guan wen: first instalment 265,093 guan 333 wen; second instalment 265,093 guan 333 wen; last instalment 265,093 guan 334 wen. Porterage money: original reckoning 39,509 guan 333 1/3 wen; actual use 5089 guan 792 wen; surplus 34,419 guan 541 1/3 wen.\end{uniten}

\begin{unitzh}\unitlabel{shushu-040.15}丁郡合解三十九萬八千一百二十六貫文；初限一十三萬二千七百八貫六百六十六文，次限一十三萬二千七百八貫六百六十六文，末限一十三萬二千七百八貫六百六十八文；傭錢原理一萬四千一百七十三貫五百文，實用二千三百八十八貫七百五十六文，寛餘一萬一千七百八十四貫七百四十四文{\zhsmall 按：傭錢原理六千訛為四千，寛餘三千訛為一千。}。\end{unitzh}
\begin{uniten}Commandery Ding's due delivery is 398,126 guan wen: first instalment 132,708 guan 666 wen; second instalment 132,708 guan 666 wen; last instalment 132,708 guan 668 wen. Porterage money: original reckoning 14,173 guan 500 wen; actual use 2388 guan 756 wen; surplus 11,784 guan 744 wen [Siku editors: In the original reckoning of the porterage money, 6000 has been corrupted into 4000; in the surplus, 3000 has been corrupted into 1000.].\end{uniten}

\begin{unitzh}\unitlabel{shushu-040.16}術曰：以均輸求之。置各郡銀絹，乗各價，併之，歸足，原展足為舊㑹；次以五約舊㑹為新㑹，各得合解錢；以限數除之，得每限錢，不盡，併歸末限。次置里數，乗每里傭價為率；以率乗原銀及原絹，各為傭實；以每擔銀絹率各為法，實如法而一，不滿者亦為擔，併之為原理傭錢。次以率乗合觧錢為實，乃以錢物毎擔率為法，實如法而一，各得實用傭錢；以減原理傭錢，餘為寛餘傭錢，合問。\end{unitzh}
\begin{uniten}Method: Solve it by equitable transport (junshu). Set up each commandery's silver and silk, multiply by their respective prices, and add them; divide (gui) the full [count], expanding (zhan) the original full [count] to make old huizi; next reduce (yue) the old huizi by 5 to make new huizi, obtaining in each case the money due for delivery. Divide it by the number of instalments, obtaining the money of each instalment; what does not divide exactly is all put into the last instalment. Next set up the number of li and multiply by the porterage price per li, making the rate (lü); multiply the original silver and the original silk by the rate, making in each case the porterage dividend (yong shi); take the per-load rates of silver and of silk respectively as the divisors (fa); divide the dividend by the divisor; what falls short of full is also taken as a load; add them, making the original-reckoning porterage money. Next multiply the money due for delivery by the rate, making the dividend (shi), and then take the per-load rate of the money goods (qianwu) as the divisor (fa); divide the dividend by the divisor, obtaining in each case the actual-use porterage money; subtract it from the original-reckoning porterage money; the remainder is the surplus porterage money, which fits the question.\end{uniten}

\begin{unitzh}\unitlabel{shushu-040.17}數圖\end{unitzh}
\begin{uniten}Number diagram (shutu)\end{uniten}

\begin{unitzh}\unitlabel{shushu-041.1}草曰：置各郡銀絹，乗各價。甲郡銀三千二百兩，乙郡銀二千七百兩，郡銀四千兩，丁郡銀二千六百兩貫，於右行；甲郡銀兩價二貫二百足，乙郡銀兩價二貫三百足，丙郡銀兩價九貫三百新㑹，丁郡銀兩價五十一貫舊㑹，於左行，對乗之：甲得七千四十貫足，乙得六千二百一十貫足，丙得三萬七千二百貫新㑹，丁得十三萬二千六百貫舊㑹。又列置各郡絹：甲六萬四千匹，乙四萬九千二百匹，丙七萬三千六百匹，丁三萬二千三十五匹，於右行；各郡絹匹價：甲二貫足，乙二貫四百二十足，丙新㑹十貫三百，丁五十八貫舊㑹，於左行，亦對乗之：甲得一十二萬八千貫足，乙得一十一萬九千六十四貫足，丙得七十五萬八千八十貫新㑹，丁得一百八十五萬八千三十貫舊㑹。\end{unitzh}
\begin{uniten}Working: Set up the silver and silk of each commandery and multiply them by the respective prices. Commandery A, silver 3,200 liang (兩); commandery B, silver 2,700 liang; commandery, silver 4,000 liang [Translator's note: the name 丙 (C) is wanting before 郡 "commandery"]; commandery D, silver 2,600 liang guan [Translator's note: 貫 guan is evidently out of place here]; in the right column. Commandery A's price of silver per liang, 2 guan (貫) 200, full-string (zu, 足); commandery B's price of silver per liang, 2 guan 300, full-string; commandery C's price of silver per liang, 9 guan 300, new huizi (xinhui); commandery D's price of silver per liang, 51 guan, old huizi (jiuhui); in the left column. Multiply them correspondingly (duicheng): A obtains 7,040 guan full-string, B obtains 6,210 guan full-string, C obtains 37,200 guan new huizi, D obtains 132,600 guan old huizi. Again lay out the silk of each commandery: A 64,000 pi (匹), B 49,200 pi, C 73,600 pi, D 32,035 pi, in the right column; the price per pi of silk of each commandery: A 2 guan full-string, B 2 guan 420 full-string, C new huizi 10 guan 300, D 58 guan old huizi, in the left column. Again multiply them correspondingly: A obtains 128,000 guan full-string, B obtains 119,064 guan full-string, C obtains 758,080 guan new huizi, D obtains 1,858,030 guan old huizi.\end{uniten}

\begin{unitzh}\unitlabel{shushu-041.2}乃併各郡銀絹價：甲共一十三萬五千四十貫足，乙共一十二萬五千二百七十四貫足，丙共七十九萬五千二百八十貫新㑹，丁共一百九十九萬六百三十貫舊㑹。甲以舊㑹價五十四文展足錢，得二百五十萬七百四十貫七百四十文；乙以舊㑹價五十九文展足錢，得二百一十二萬三千二百八十八貫一百三十六文；丙係新㑹；丁係舊㑹。今甲、乙、丁俱以五除之，皆為新㑹：甲得五十萬一百四十八貫一百四十八文，乙得四十二萬四千六百五十七貫六百二十七文，丙得七十九萬五千二百八十貫文，丁得三十九萬八千一百二十六貫，各為合解錢。\end{unitzh}
\begin{uniten}Then add together the prices of the silver and silk of each commandery: A in all 135,040 guan full-string, B in all 125,274 guan full-string, C in all 795,280 guan new huizi, D in all 1,990,630 guan old huizi. For A, expand (zhan) the full-string cash (zuqian) by the old-huizi price of 54 wen (文), obtaining 2,500,740 guan 740 wen; for B, expand the full-string cash by the old-huizi price of 59 wen, obtaining 2,123,288 guan 136 wen; C is in new huizi; D is in old huizi. Now divide A, B and D all by 5, and all become new huizi: A obtains 500,148 guan 148 wen, B obtains 424,657 guan 627 wen, C obtains 795,280 guan wen, D obtains 398,126 guan; each is the total remittance money (hejie qian).\end{uniten}

\begin{unitzh}\unitlabel{shushu-041.3}以限數三除之：甲得一十六萬六千七百一十六貫四十九文，為初限、次限數，不盡一文，増入次限數内，共得一十六萬六千七百一十六貫五十文，為末限數；乙得一十四萬一千五百五十二貫五百四十二文，為初限、次限數，不盡一文，増入，得一十四萬一千五百五十二貫五百四十三文，為末限數；丙得二十六萬五千九十三貫三百三十三文，為初限、次限數，不盡一文，増入，得二十六萬五千九十三貫三百三十四文，為末限數；丁舊一十三萬二千七百八貫六百六十六文，為初限、次限數，不盡二文，増入，得一十三萬二千七百八貫六百六十八文，為末限數。\end{unitzh}
\begin{uniten}Divide by the number of installments (xianshu), 3: A obtains 166,716 guan 49 wen as the amount of the first installment (chuxian) and of the second installment (cixian); the remainder (bujin), 1 wen, is added into the amount of the second installment, giving in all 166,716 guan 50 wen as the amount of the last installment (moxian). B obtains 141,552 guan 542 wen as the amount of the first and second installments; the remainder, 1 wen, is added in, giving 141,552 guan 543 wen as the amount of the last installment. C obtains 265,093 guan 333 wen as the amount of the first and second installments; the remainder, 1 wen, is added in, giving 265,093 guan 334 wen as the amount of the last installment. D, old, 132,708 guan 666 wen [Translator's note: 舊 "old" stands where the other entries have 得 "obtains"], as the amount of the first and second installments; the remainder, 2 wen, is added in, giving 132,708 guan 668 wen as the amount of the last installment.\end{uniten}

\begin{unitzh}\unitlabel{shushu-041.4}各以里數乗傭錢，各為率。置甲郡一千里，乙郡九百八十里，丙郡二千里，丁郡一千五百里，於右行；次置甲郡傭錢六文足，乙郡傭錢四文二分足，丙郡傭錢八十文舊㑹，丁郡傭錢一百舊㑹，於左行，與右行對乗之：甲得率六置足，乙得率四貫一百一十六足，丙得率一百六十貫舊，丁得率一百五十貫舊，於右行。\end{unitzh}
\begin{uniten}In each case multiply the hire money (yongqian) by the number of li (里), each giving the rate (lü). Set up commandery A 1,000 li, commandery B 980 li, commandery C 2,000 li, commandery D 1,500 li, in the right column; next set up commandery A's hire money 6 wen full-string, commandery B's hire money 4 wen 2 fen (分) full-string, commandery C's hire money 80 wen old huizi, commandery D's hire money 100 old huizi, in the left column, and multiply them correspondingly with the right column: A obtains the rate 6 zhi (置) full-string [Translator's note: 置 "set up" evidently stands for 貫 guan], B obtains the rate 4 guan 116 full-string, C obtains the rate 160 guan old, D obtains the rate 150 guan old, in the right column.\end{uniten}

\begin{unitzh}\unitlabel{shushu-041.5}以率乗原銀數，各為傭實。次置甲原銀三千二百兩，乙銀二千七百，丙銀四千，丁銀二千六百兩，於左行，與右行對乗之：甲得一萬九千二百貫，乙得一萬一千一百一十三貫二百文，丙得六十四萬貫舊，丁得三十九萬貫舊，皆銀傭，置於右行。次置甲、乙、丙、丁每擔率銀五百兩為法，遍除左行：甲得三十八貫四百足，乙得二十二貫二百二十六文四分足，丙得一千二百八十貫舊，丁得七百八十貫舊，為各郡銀傭錢，列實寄别行。\end{unitzh}
\begin{uniten}Multiply the original amounts of silver by the rates, each as the hire dividend (yongshi). Next set up A's original silver 3,200 liang, B's silver 2,700, C's silver 4,000, D's silver 2,600 liang, in the left column, and multiply them correspondingly with the right column: A obtains 19,200 guan, B obtains 11,113 guan 200 wen, C obtains 640,000 guan old, D obtains 390,000 guan old; all are silver hire; set them in the right column. Next set up the per-load (dan, 擔) rate of A, B, C and D, silver 500 liang, as divisor (fa), and divide throughout (bianchu) the left column: A obtains 38 guan 400 full-string, B obtains 22 guan 226 wen 4 fen full-string, C obtains 1,280 guan old, D obtains 780 guan old; these are the silver hire money of each commandery; lay them out as shi (實) and lodge them in a separate column.\end{uniten}

\begin{unitzh}\unitlabel{shushu-041.6}次置甲原絹六萬四千匹，乙絹四萬九千二百匹，丙絹七萬三千六百匹，丁絹三萬二千三十五匹，為左行，與右行各率對乗之：甲得三十八萬四千貫足，乙得二十萬二千五百七貫二百足，丙得一千一百七十七萬六千貫舊，丁得四百八十萬五千二百五十貫，各為絹傭實。次以四郡每擔絹率六十匹為法除之：甲得六千四百貫足，乙得三千三百七十五貫一百二十足，丙得一十九萬六千二百六十六貫六百六十六文三分文之二舊，丁得八萬八十七貫五百舊，為各郡絹傭錢。併入寄别行：甲得六千四百三十八貫四百足，乙得三千三百九十七貫三百四十六文四分足，丙得一十九萬七千五百四十六貫六百六十六文三分文之二舊，丁得八萬八百六十七貫五百舊，列右行。\end{unitzh}
\begin{uniten}Next set up A's original silk 64,000 pi, B's silk 49,200 pi, C's silk 73,600 pi, D's silk 32,035 pi, as the left column, and multiply them correspondingly by the respective rates of the right column: A obtains 384,000 guan full-string, B obtains 202,507 guan 200 full-string, C obtains 11,776,000 guan old, D obtains 4,805,250 guan; each is the silk hire dividend. Next take the per-load silk rate of the four commanderies, 60 pi, as divisor and divide: A obtains 6,400 guan full-string, B obtains 3,375 guan 120 full-string, C obtains 196,266 guan 666 2/3 wen old, D obtains 80,087 guan 500 old; these are the silk hire money of each commandery. Add them into the lodged separate column: A obtains 6,438 guan 400 full-string, B obtains 3,397 guan 346 wen 4 fen full-string, C obtains 197,546 guan 666 2/3 wen old, D obtains 80,867 guan 500 old; lay them out in the right column.\end{uniten}

\begin{unitzh}\unitlabel{shushu-041.7}其甲舊㑹價五十四文，五因之，得二百七十文足；乙舊㑹價五十九文，亦五因之，得二百九十五文；丙以五；丁亦以五：於左行。以對約右行，皆為新㑹：甲得二萬三千八百四十五貫九百二十五文二十七分文之二十五，乙得一萬一千五百一十六貫四百二十八文五十九分文之二十八，丙得三萬九千五百九貫三百三十三文三分文之一，丁得一萬六千一百七十三貫五百文，並新㑹，係四郡原傭價錢。\end{unitzh}
\begin{uniten}As for A's old-huizi price of 54 wen, multiply it by 5 (wu yin), obtaining 270 wen full-string; B's old-huizi price of 59 wen, likewise multiply it by 5, obtaining 295 wen; for C, 5; for D also 5: in the left column. Use them correspondingly to reduce (yue) the right column, and all become new huizi: A obtains 23,845 guan 925 25/27 wen, B obtains 11,516 guan 428 28/59 wen, C obtains 39,509 guan 333 1/3 wen, D obtains 16,173 guan 500 wen, all in new huizi; these are the original hire-price money of the four commanderies.\end{uniten}

\begin{unitzh}\unitlabel{shushu-041.8}次以原四郡率對乗四郡合解新㑹，各為實：其甲率六貫足，乗甲合解錢五十萬一百四十八貫一百四十八文，得三十億八十八萬八千八百八十貫；其乙率四貫一百一十六足，乗乙合解錢四十二萬四千六百五十七貫六百二十七文，得一十七億四千七百八十九萬七百九十二貫七百三十二文足；其丙率一百六十貫舊，乗丙合解錢七十九萬五千二百八十貫，得一千二百七十二億四千四百八十萬貫舊；其丁率一百五十貫舊，乗丁合解錢三十九萬八千一百二十六貫，得五百九十七億一千八百九十萬貫舊：各為實。\end{unitzh}
\begin{uniten}Next multiply the total remittances of the four commanderies in new huizi correspondingly by the original rates of the four commanderies, each as dividend (shi): A's rate, 6 guan full-string, multiplied by A's total remittance money, 500,148 guan 148 wen, gives 3,000,888,880 guan; B's rate, 4 guan 116 full-string, multiplied by B's total remittance money, 424,657 guan 627 wen, gives 1,747,890,792 guan 732 wen full-string; C's rate, 160 guan old, multiplied by C's total remittance money, 795,280 guan, gives 127,244,800,000 guan old; D's rate, 150 guan old, multiplied by D's total remittance money, 398,126 guan, gives 59,718,900,000 guan old: each is a dividend.\end{uniten}

\begin{unitzh}\unitlabel{shushu-041.9}乃以每擔率五千貫為法而一：甲得六百貫一百七十七文足，不盡三千八百八十貫文；乙得三百四十九貫五百七十八文足，不盡七百九十二貫七百三十二文；丙得二萬五千四百四十八貫九百六十文舊㑹；丁得一萬一千九百四十三貫七百八十文舊㑹：為各郡實用。\end{unitzh}
\begin{uniten}Then take the per-load rate, 5,000 guan, as divisor and divide: A obtains 600 guan 177 wen full-string, remainder 3,880 guan wen; B obtains 349 guan 578 wen full-string, remainder 792 guan 732 wen; C obtains 25,448 guan 960 wen old huizi; D obtains 11,943 guan 780 wen old huizi: these are the actual use (shiyong) of each commandery.\end{uniten}

\begin{unitzh}\unitlabel{shushu-041.10}甲以二百七十文約，乙以二百九十五文約，丙、丁皆五約，為新㑹：甲二千二百二十二貫八百七十七文，不盡二百一十文；乙一千一百八十五貫一十文，不盡五百；丙五千八十九貫七百九十二文；丁二千三百八十八貫七百五十六文。各減原理：甲餘二萬一千六百二十三貫四十八文，乙餘一萬三百三十一貫四百一十八文，丙餘三萬四千四百一十九貫五百四十一文，丁餘一萬一{\zhsmall 按二訛一}千七百八十四貫七百四十四文。合問。\end{unitzh}
\begin{uniten}Reduce A by 270 wen, reduce B by 295 wen, reduce C and D both by 5, making new huizi: A 2,222 guan 877 wen, remainder 210 wen; B 1,185 guan 10 wen, remainder 500; C 5,089 guan 792 wen; D 2,388 guan 756 wen. Subtract each from the original reckoning (yuanli): A leaves 21,623 guan 48 wen, B leaves 10,331 guan 418 wen, C leaves 34,419 guan 541 wen, D leaves 11[Siku editors: 二 (2) is miswritten as 一 (1).],784 guan 744 wen. This answers the question.\end{uniten}

\begin{unitzh}\unitlabel{shushu-041.11}<子部,天文算法類,算書之屬,數學九章,卷六上> <子部,天文算法類,算書之屬,數學九章,卷六上> <子部,天文算法類,算書之屬,數學九章,卷六上> <子部,天文算法類,算書之屬,數學九章,卷六上>\end{unitzh}
\begin{uniten}«Masters section (zibu), Astronomy and Mathematics category (tianwen suanfa lei), Mathematical Books subcategory (suanshu zhi shu), Shuxue jiuzhang, juan 6, part 1» «Masters section, Astronomy and Mathematics category, Mathematical Books subcategory, Shuxue jiuzhang, juan 6, part 1» «Masters section, Astronomy and Mathematics category, Mathematical Books subcategory, Shuxue jiuzhang, juan 6, part 1» «Masters section, Astronomy and Mathematics category, Mathematical Books subcategory, Shuxue jiuzhang, juan 6, part 1» [Translator's note: catalogue label of the Siku quanshu; in the transcription it apparently marks the place of a diagram.]\end{uniten}

\begin{unitzh}\unitlabel{shushu-041.12}今欲變右行足錢、舊㑹皆為新㑹，故以五遍乗：甲陌五十四，得二百七十；乙陌五十九，得二百九十五。\end{unitzh}
\begin{uniten}Now we wish to change the full-string cash and the old huizi of the right column all into new huizi; therefore multiply throughout by 5: A's string-rate (mo, 陌) 54 gives 270; B's string-rate 59 gives 295.\end{uniten}

\begin{unitzh}\unitlabel{shushu-041.13}<子部,天文算法類,算書之屬,數學九章,卷六上> <子部,天文算法類,算書之屬,數學九章,卷六上>\end{unitzh}
\begin{uniten}«Masters section, Astronomy and Mathematics category, Mathematical Books subcategory, Shuxue jiuzhang, juan 6, part 1» «Masters section, Astronomy and Mathematics category, Mathematical Books subcategory, Shuxue jiuzhang, juan 6, part 1»\end{uniten}

\begin{unitzh}\unitlabel{shushu-041.14}按：右數實用傭錢内第二層乙條下子數□訛□；原理傭錢内末層丁條左第二數丅訛□；寛餘傭錢内第二層乙條下子數□訛□，末層丁條左第二數川訛丨。餘數俱合。\end{unitzh}
\begin{uniten}Note: In the numbers at right: within the actual-use hire money, in the B entry (tiao) of the second tier (ceng), the lower numerator figure (zishu): □ is miswritten as □; within the original-reckoning hire money, in the D entry of the last tier, the second figure from the left: 丅 is miswritten as □; within the surplus (kuanyu) hire money, in the B entry of the second tier, the lower numerator figure: □ is miswritten as □, and in the D entry of the last tier, the second figure from the left: 川 is miswritten as 丨. The remaining numbers all agree. [Translator's note: 丅, 川 and 丨 are counting-rod numerals for 6, 3 and 1; □ stands for rod numerals not reproduced in the transcription.]\end{uniten}

\begin{unitzh}\unitlabel{shushu-041.15}《數學九章》卷六上\end{unitzh}
\begin{uniten}Shuxue jiuzhang (Mathematical Treatise in Nine Sections), juan 6, part 1.\end{uniten}

\section*{卷6下}\addcontentsline{toc}{section}{卷6下}

\begin{unitzh}\unitlabel{shushu-042.1}欽定四庫全書 數學九章卷六下 宋 秦九韶 撰\end{unitzh}
\begin{uniten}Imperially Approved Complete Library of the Four Treasuries (Qinding Siku quanshu). Mathematics in Nine Sections (Shuxue jiuzhang), juan 6, part 2. Composed by Qin Jiushao of the Song.\end{uniten}

\begin{unitzh}\unitlabel{shushu-042.2}錢穀\end{unitzh}
\begin{uniten}Money and Grain (qiangu)\end{uniten}

\begin{unitzh}\unitlabel{shushu-042.3}僦直推原\end{unitzh}
\begin{uniten}Inferring the Original from the Rent Charge (jiuzhi tuiyuan)\end{uniten}

\begin{unitzh}\unitlabel{shushu-042.4}問房廊數内一户，日納一百五十六文八分足。為準指揮：未曽經減者減三分，已曽經減三分者減二分，已曽經減二分者更減二分。今本户累經減者，欲知原額房錢幾何？答曰：原額三百五十文。\end{unitzh}
\begin{uniten}Question: Among the houses and galleries (fanglang), one household pays daily 156 wen (文) 8 fen (分), full-count (zu). In compliance with a directive (zhihui): those never yet reduced are reduced by 3 parts (fen); those already reduced by 3 parts are reduced by 2 parts; those already reduced by 2 parts are further reduced by 2 parts. Now this household has been reduced repeatedly; we wish to know how much the originally assessed (yuan'e) house rent (fangqian) was. Answer: The original assessment was 350 wen.\end{uniten}

\begin{unitzh}\unitlabel{shushu-042.5}術曰：以衰分求之。列一十分兩行，各三位；列減分，對減右行；以餘者相乘為法。以左行原列相乘，得納錢為實。實如法而一，得原額錢。\end{unitzh}
\begin{uniten}Method: Solve it by proportional distribution (cuifen). Set out 10 parts in two columns (hang), each of 3 places (wei); set out the reduction parts and subtract them correspondingly from the right column; multiply the remainders together to make the divisor (fa). Multiply together the entries originally set out in the left column; with what is obtained, the paid cash makes the dividend (shi). Divide the dividend by the divisor, obtaining the original assessed cash.\end{uniten}

\begin{unitzh}\unitlabel{shushu-042.6}草曰：列一十三位於左行，又列一十分三位於右行。其右上減去初減三分，右中減去次減二分，右下減去更減二分，右行餘七、八、八，以相乗，得四百四十八，為法。乃以左行三位一十分相乘，得一千，為乗率，以乘見日納錢一百五十六文八分，得一百五十六貫八百文，為實。實如法而一，得三百五十文，為本户原額戸錢。\end{unitzh}
\begin{uniten}Working: Set out 10, in 3 places, in the left column; further set out 10 parts, in 3 places, in the right column. From its right top subtract the first reduction of 3 parts; from the right middle subtract the next reduction of 2 parts; from the right bottom subtract the further reduction of 2 parts. The right column leaves 7, 8, 8; multiply these together, obtaining 448 as the divisor. Then multiply together the 10 parts in the 3 places of the left column, obtaining 1000 as the multiplying rate (chenglü); with it multiply the present daily paid cash, 156 wen 8 fen, obtaining 156 guan (貫) 800 wen as the dividend. Divide the dividend by the divisor, obtaining 350 wen as this household's original assessed household cash.\end{uniten}

\begin{unitzh}\unitlabel{shushu-042.7}推求本息\end{unitzh}
\begin{uniten}Deducing Principal and Interest (tuiqiu benxi)\end{uniten}

\begin{unitzh}\unitlabel{shushu-042.8}問三庫息例：萬貫以上一釐，千貫以上二釐五毫，百貫以上三釐。甲庫本四十九萬三千八百貫，乙庫本三十七萬三百貫，丙庫本二十四萬六千八百貫。今三庫共約到息錢二萬五千六百四十四貫二百文。其典率：甲反錐差，乙方錐差，丙蒺藜差。欲知原典三例本息各幾何？\end{unitzh}
\begin{uniten}Question: The interest rules (xili) of three treasuries (ku): for 10000 guan and above, 1 li (釐); for 1000 guan and above, 2 li 5 hao (毫); for 100 guan and above, 3 li. The principal of Treasury A (jia) is 493800 guan; the principal of Treasury B (yi) is 370300 guan; the principal of Treasury C (bing) is 246800 guan. Now the three treasuries together have brought in, by reckoning (yuedao), interest cash of 25644 guan 200 wen. Their pawning rates (dianlü): A, the inverted-cone difference (fanzhui cha); B, the square-cone difference (fangzhui cha); C, the caltrop difference (jili cha). We wish to know how much, in each case, the principal and interest originally pawned under the three rules were.\end{uniten}

\begin{unitzh}\unitlabel{shushu-042.9}按此即衰分題也。其差有反錐、方錐、蒺藜之名，蓋以一、二、三遞減如立錐為反錐，以一、四、九平方遞加為方錐，以一、三、六三數遞加為蒺藜，是必古有其名也。至以各差求各本，則因各本原依各差入之也。\end{unitzh}
\begin{uniten}Note: This is precisely a problem of proportional distribution. Its differences bear the names inverted cone, square cone and caltrop: presumably 1, 2, 3, decreasing successively like a standing cone, are taken as the inverted cone; 1, 4, 9, squares (pingfang) increasing successively, are taken as the square cone; 1, 3, 6, three numbers increasing successively, are taken as the caltrop. These must be names that existed in antiquity. As for using each difference to find each principal, that is because each principal was originally entered according to each difference.\end{uniten}

\begin{unitzh}\unitlabel{shushu-042.10}答曰：甲庫共納息九千五十三貫文，一釐息二千四百六十九貫文，二釐半息四千一百一十五貫文，三釐息二千四百六十九貫文；乙庫共納息一萬五十一貫文，一釐息二百六十四貫五百文，二釐半息二千六百四十五貫文，三釐息七千一百四十一貫五百文；丙庫共納息六千五百四十貫二百文，一釐息二百四十六貫八百文，二釐半息一千八百五十一貫文，三釐息四千四百四十二貫四百文。\end{unitzh}
\begin{uniten}Answer: Treasury A pays in combined interest of 9053 guan wen: interest at 1 li, 2469 guan wen; interest at 2 1/2 li, 4115 guan wen; interest at 3 li, 2469 guan wen. Treasury B pays in combined interest of 10051 guan wen: interest at 1 li, 264 guan 500 wen; interest at 2 1/2 li, 2645 guan wen; interest at 3 li, 7141 guan 500 wen. Treasury C pays in combined interest of 6540 guan 200 wen: interest at 1 li, 246 guan 800 wen; interest at 2 1/2 li, 1851 guan wen; interest at 3 li, 4442 guan 400 wen.\end{uniten}

\begin{unitzh}\unitlabel{shushu-042.11}術曰：置諸庫諸色之差，照釐率為三行，縱併之為約率，横命之為乗率。以約率各約自庫之本，各得，以遍乗未併乗率。然後各以釐率横乗之，次以縱併之，為各庫共息。\end{unitzh}
\begin{uniten}Method: Set out the differences of the various kinds for the various treasuries, according to the li rates (lilü), as three columns; sum them vertically to make the dividing rates (yuelü); designate them horizontally as the multiplying rates. With the dividing rates divide (yue) each treasury's own principal; with each result (de), multiply throughout the multiplying rates that have not been summed. Then multiply each horizontally by the li rates; next sum them vertically, making each treasury's combined interest.\end{uniten}

\begin{unitzh}\unitlabel{shushu-042.12}草曰：置甲庫反錐差，自下置三、二、一於右行；次置乙庫方錐差，自上置一、四、九於中行；次置丙庫蒺藜差，自上置一、三、六於左行，各為三庫上中下三等乘率。乃縱併甲差三、二、一，得六，為甲約率；縱併乙差一、四、九，得一十四，為乙約率；縱併丙差一、三、六，得一十，為丙約率。直命九位數，各為上中下乗率。乃先以約率各約自庫之本：乃以甲約率六約甲本四十九萬三千八百貫，得八萬二千三百貫，為甲得；次以乙約率一十四約乙本三十七萬三百貫，得二萬六千四百五十貫，為乙得；次以丙約率一十約丙本二十四萬六千八百貫，得二萬四千六百八十貫，為丙得。\end{unitzh}
\begin{uniten}Working: Set out Treasury A's inverted-cone difference: from the bottom, set 3, 2, 1 in the right column. Next set out Treasury B's square-cone difference: from the top, set 1, 4, 9 in the middle column. Next set out Treasury C's caltrop difference: from the top, set 1, 3, 6 in the left column. These are, respectively, the multiplying rates of the three grades, upper, middle and lower, of the three treasuries. Then sum vertically A's difference 3, 2, 1, obtaining 6 as A's dividing rate; sum vertically B's difference 1, 4, 9, obtaining 14 as B's dividing rate; sum vertically C's difference 1, 3, 6, obtaining 10 as C's dividing rate. Directly designate the numbers in the 9 places as the upper, middle and lower multiplying rates respectively. Then first, with the dividing rates, divide each treasury's own principal: thus with A's dividing rate 6 divide A's principal 493800 guan, obtaining 82300 guan as A's result; next with B's dividing rate 14 divide B's principal 370300 guan, obtaining 26450 guan as B's result; next with C's dividing rate 10 divide C's principal 246800 guan, obtaining 24680 guan as C's result.\end{uniten}

\begin{unitzh}\unitlabel{shushu-042.13}以各得乗未併乗率：其甲所得八萬二千三百貫，乘反錐乗率三、二、一，得二十四萬六千九百貫為上率，得一十六萬四千六百貫為中率，得八萬二千三百貫為下率；其乙所得二萬六千四百五十貫，以乗方錐差一、四、九，得二萬六千四百五十貫為上率，得一十萬五千八百貫為中率，得二十三萬八千五十貫為下率；其丙所得二萬四千六百八十貫，以乗蒺藜差一、三、六，得二萬四千六百八十貫，以乘蒺藜七萬四千四十貫為中率，得一十四萬八千八十貫為下率。\end{unitzh}
\begin{uniten}With each result multiply the multiplying rates that have not been summed. A's result, 82300 guan, multiplies the inverted-cone multiplying rates 3, 2, 1, obtaining 246900 guan as the upper rate, obtaining 164600 guan as the middle rate, obtaining 82300 guan as the lower rate. B's result, 26450 guan, is used to multiply the square-cone difference 1, 4, 9, obtaining 26450 guan as the upper rate, obtaining 105800 guan as the middle rate, obtaining 238050 guan as the lower rate. C's result, 24680 guan, is used to multiply the caltrop difference 1, 3, 6, obtaining 24680 guan, with which multiply the caltrop 74040 guan as the middle rate, obtaining 148080 guan as the lower rate. [Translator's note: the words "with which multiply the caltrop" appear garbled; the parallel cases have "as the upper rate; obtaining" in their place.]\end{uniten}

\begin{unitzh}\unitlabel{shushu-042.14}然後各以息釐數乘各庫三乗{\zhsmall 此是變文為庫}：其甲以一釐乗上率二十四萬六千九百貫，得二千四百六十九貫，為上息；以二釐五毫乘中率一十六萬四千六百貫，得四千一百一十五貫，為中息；以三釐乘下率八萬二千三百貫，得二千四百六十九貫，為下息。併上中下三息，得九千五十三貫文，為甲庫共息。其乙庫以一釐乗上率二萬六千四百五十貫，得二百六十四貫五百文，為上息；以二釐五毫乗中率一十萬五千八百貫，得二千六百四十五貫，為中息；以三釐乗下率二十三萬八千五十貫，得七千一百四十一貫五百，為下息。併上中下三息，得一萬五十一貫文，為乙庫共息。\end{unitzh}
\begin{uniten}Then, with the interest li numbers, multiply each treasury's three products (san cheng) [Note: here the wording is varied, using "treasury".]: for A, with 1 li multiply the upper rate 246900 guan, obtaining 2469 guan as the upper interest; with 2 li 5 hao multiply the middle rate 164600 guan, obtaining 4115 guan as the middle interest; with 3 li multiply the lower rate 82300 guan, obtaining 2469 guan as the lower interest. Sum the three interests, upper, middle and lower, obtaining 9053 guan wen as Treasury A's combined interest. For Treasury B, with 1 li multiply the upper rate 26450 guan, obtaining 264 guan 500 wen as the upper interest; with 2 li 5 hao multiply the middle rate 105800 guan, obtaining 2645 guan as the middle interest; with 3 li multiply the lower rate 238050 guan, obtaining 7141 guan 500 as the lower interest. Sum the three interests, upper, middle and lower, obtaining 10051 guan wen as Treasury B's combined interest.\end{uniten}

\begin{unitzh}\unitlabel{shushu-042.15}其丙庫以一釐乗上率二萬四千六百八十貫，得二百四十六貫八百，為上息；以二釐五毫乘中率七萬四千四十貫，得一千八百五十一貫，為中息；以三釐乗下率一十四萬八千八十貫，得四千四百四十二貫四百文，為下息。併上中下三息，得六千五百四十貫二百文，為丙庫共息。併三庫共息，得二萬五千六百四十四貫二百文，為總息。\end{unitzh}
\begin{uniten}For Treasury C, with 1 li multiply the upper rate 24680 guan, obtaining 246 guan 800 as the upper interest; with 2 li 5 hao multiply the middle rate 74040 guan, obtaining 1851 guan as the middle interest; with 3 li multiply the lower rate 148080 guan, obtaining 4442 guan 400 wen as the lower interest. Sum the three interests, upper, middle and lower, obtaining 6540 guan 200 wen as Treasury C's combined interest. Sum the combined interests of the three treasuries, obtaining 25644 guan 200 wen as the overall interest (zongxi).\end{uniten}

\begin{unitzh}\unitlabel{shushu-042.16}子部，天文算法類，算書之屬，數學九章，卷六下。子部，天文算法類，算書之屬，數學九章，卷六下。子部，天文算法類，算書之屬，數學九章，卷六下。\end{unitzh}
\begin{uniten}Masters Section (zibu), Astronomy and Mathematics Category, Mathematical Books Subcategory, Shuxue jiuzhang, juan 6, part 2. Masters Section, Astronomy and Mathematics Category, Mathematical Books Subcategory, Shuxue jiuzhang, juan 6, part 2. Masters Section, Astronomy and Mathematics Category, Mathematical Books Subcategory, Shuxue jiuzhang, juan 6, part 2.\end{uniten}

\begin{unitzh}\unitlabel{shushu-042.17}易牒知原{\zhsmall 按舊本此問無題，今増入}\end{unitzh}
\begin{uniten}Knowing the Original from an Exchange of Certificates (yidie zhiyuan) [Siku editors: In the old edition this problem has no title; it has now been added.]\end{uniten}

\begin{unitzh}\unitlabel{shushu-043.1}問：出度牒差人營運，毎三道易鹽一十三袋，鹽二袋易布八十四疋，布一十五疋易絹三疋半，絹六疋易銀七兩二錢。今趂到銀九千一百七十二兩八錢，欲知原闗度牒道數幾何？答曰：度牒一百八十道。\end{unitzh}
\begin{uniten}Question: Ordination certificates (dudie) are issued and men are sent to trade with them. Every 3 certificates are exchanged for 13 bags of salt; 2 bags of salt are exchanged for 84 pi (疋) of cloth; 15 pi of cloth are exchanged for 3 1/2 pi of silk; 6 pi of silk are exchanged for 7 liang (兩) 2 qian (錢) of silver. Now 9172 liang 8 qian of silver have been earned. One wishes to know how many ordination certificates were originally drawn. Answer: 180 ordination certificates.\end{uniten}

\begin{unitzh}\unitlabel{shushu-043.2}術曰：以粟米互乘易法求之。列各數，以本色相對如鴈翅，以多一事者相乗為實，以少一事者相乘為法，除之。\end{unitzh}
\begin{uniten}Method: Solve it by the "Millet and Grain" (sumi) method of exchange by cross-multiplication (hucheng). Lay out the numbers, setting each kind against its own kind (bense) like the wings of a wild goose (yanchi). Multiply together the ones having one item more to make the dividend (shi); multiply together the ones having one item fewer to make the divisor (fa); divide.\end{uniten}

\begin{unitzh}\unitlabel{shushu-043.3}草曰：先以度牒三道乗鹽二袋，得六；以乘布一十五，得九十；文乗絹六疋，得五百四十；乃乗銀九萬一千七百二十八錢，得四千九百五十三萬三千一百二十錢，為實。次以鹽一十三袋乗布八十四，得一千九十二；以乗絹三疋五分，得三千八百二十二；乃乘銀七兩二錢，得二十七萬五千一百八十四錢，為法。除實，得一百八十道，為原闗度牒。\end{unitzh}
\begin{uniten}Working: First, with the 3 ordination certificates multiply the 2 bags of salt, obtaining 6; with this multiply the cloth, 15, obtaining 90; wen [Translator's note: 文 is evidently an error for 又 "again".] multiply the silk, 6 pi, obtaining 540; then multiply the silver, 91,728 qian, obtaining 49,533,120 qian, as the dividend (shi). Next, with the salt, 13 bags, multiply the cloth, 84, obtaining 1092; with this multiply the silk, 3 pi 5 fen (分), obtaining 3822; then multiply the silver, 7 liang 2 qian, obtaining 275,184 qian, as the divisor (fa). Divide the dividend, obtaining 180 certificates: the ordination certificates originally drawn.\end{uniten}

\begin{unitzh}\unitlabel{shushu-043.4}粟米交易{\zhsmall 按：舊本此問無題，今増。}\end{unitzh}
\begin{uniten}Exchanging Millet and Grain [Siku editors: In the old edition this problem has no title; it has now been added.]\end{uniten}

\begin{unitzh}\unitlabel{shushu-043.5}問：菽三升易小麥二升，小麥一升五合易油麻八合，油麻一升二合易粳米一升八合。今將菽一十四石四斗欲易油麻，又將小麥二十一石六斗欲易粳米，問各幾何？答曰：油麻五石一斗二升；粳米一十七石二斗八升。\end{unitzh}
\begin{uniten}Question: 3 sheng (升) of beans are exchanged for 2 sheng of wheat; 1 sheng 5 ge (合) of wheat are exchanged for 8 ge of sesame; 1 sheng 2 ge of sesame are exchanged for 1 sheng 8 ge of non-glutinous rice. Now one takes 14 dan (石) 4 dou (斗) of beans and wishes to exchange them for sesame, and further takes 21 dan 6 dou of wheat and wishes to exchange them for non-glutinous rice. What is each? Answer: sesame, 5 dan 1 dou 2 sheng; non-glutinous rice, 17 dan 2 dou 8 sheng.\end{uniten}

\begin{unitzh}\unitlabel{shushu-043.6}術曰：以粟米換易求之。置原易率，本色對列如鴈翅，以多一事者相乘為實，以少一事者相乗為法，除之，各得。或問數不干其率者，不置。\end{unitzh}
\begin{uniten}Method: Solve it by "Millet and Grain" exchange (huanyi). Set out the original exchange rates (lü), arraying each kind against its own kind like the wings of a wild goose. Multiply together the ones having one item more to make the dividend (shi); multiply together the ones having one item fewer to make the divisor (fa); divide, and each is obtained. If the number in the question has nothing to do with certain rates, these are not set out.\end{uniten}

\begin{unitzh}\unitlabel{shushu-043.7}草曰：置四色六數，列六位率如鴈翅，皆化為合。先將菽一十四石四斗化作一萬四千四百合，乃對前二句率數四位如鴈翅，至欲易油麻止，共五事，為上圖。次將小麥二十一石六斗化為二萬一千六百合，乃對後兩句率四位鴈翅，至欲粳米止，共五事，為下圖。下圖。其上圖：以菽一萬四千四百合乘麥二十，得二十八萬八千；又乗油麻八合，得二百三十萬四千合，為油麻實。次以菽三十合乘麥一十五合，得四百五十合，為法。除之，得五千一百二十合，展為五石一斗二升，為油麻。其下圖：以小麥二萬一千六百合乗油麻八合，得一十七萬二千八百合；又乗粳米一十八合，得三百一十一萬四百合，為粳米實。以小麥一以五合乗油麻一十二合，得一百八十合，為法。除之，得一萬七千二百八十，展作一十七石二斗八升，為粳米。\end{unitzh}
\begin{uniten}Working: Set out the four kinds and six numbers, laying out the rates in six positions like the wings of a wild goose, all converted into ge. First convert the beans, 14 dan 4 dou, into 14,400 ge; then set this against the four positions of the rate numbers of the first two clauses, like the wings of a wild goose, up to the sesame desired in exchange, and stop: five items in all, making the upper diagram. Next convert the wheat, 21 dan 6 dou, into 21,600 ge; then set this against the four positions of the rates of the last two clauses, in wild-goose wings, up to the desired non-glutinous rice, and stop: five items in all, making the lower diagram. Lower diagram. In the upper diagram: with the beans, 14,400 ge, multiply the wheat, 20, obtaining 288,000; further multiply by the sesame, 8 ge, obtaining 2,304,000 ge, the dividend for the sesame. Next, with the beans, 30 ge, multiply the wheat, 15 ge, obtaining 450 ge, as the divisor. Divide, obtaining 5120 ge, which expanded (zhan) makes 5 dan 1 dou 2 sheng: the sesame. In the lower diagram: with the wheat, 21,600 ge, multiply the sesame, 8 ge, obtaining 172,800 ge; further multiply by the non-glutinous rice, 18 ge, obtaining 3,110,400 ge, the dividend for the non-glutinous rice. With the wheat, 1 yi 5 ge [Translator's note: 以 is evidently an error for 十, i.e. 15 ge.], multiply the sesame, 12 ge, obtaining 180 ge, as the divisor. Divide, obtaining 17,280, which expanded makes 17 dan 2 dou 8 sheng: the non-glutinous rice.\end{uniten}

\begin{unitzh}\unitlabel{shushu-043.8}計米易麴{\zhsmall 按：舊本此問無題，今增。}\end{unitzh}
\begin{uniten}Reckoning Rice and Exchanging for Leaven [Siku editors: In the old edition this problem has no title; it has now been added.]\end{uniten}

\begin{unitzh}\unitlabel{shushu-043.9}問：庫率粳榖七石出米三石，糯米一斗易小麥一斗七升，小麥五升踏麴二斤四兩，麴一百一十斤醞糯米一石三斗。今有糯穀一千七百五十九石三斗八升，欲出穀做米，易麥踏麴，還自醞餘穀之米，須令適足，各合幾何？答曰：共穀一千七百五十九石三斗八升；出穀九百二十四石，得米三百九十六石；易麥六百七十三石二斗；踏麴三萬二百九十四斤；餘榖八百三十五石三斗八升，醞米三百五十八石二升。\end{unitzh}
\begin{uniten}Question: The storehouse rates are: 7 dan of non-glutinous grain [Translator's note: 粳 "non-glutinous" here; elsewhere in this problem the grain is 糯 "glutinous".] yield 3 dan of rice; 1 dou of glutinous rice is exchanged for 1 dou 7 sheng of wheat; 5 sheng of wheat are trodden into 2 jin (斤) 4 liang of leaven (qu); 110 jin of leaven brew 1 dan 3 dou of glutinous rice. Now there are 1759 dan 3 dou 8 sheng of glutinous grain. One wishes to take out grain and make it into rice, exchange this for wheat and tread leaven, and in turn oneself brew the rice from the remaining grain; it must be made exactly sufficient. How much should each be? Answer: grain in all, 1759 dan 3 dou 8 sheng; grain taken out, 924 dan, yielding 396 dan of rice; wheat by exchange, 673 dan 2 dou; leaven trodden, 30,294 jin; remaining grain, 835 dan 3 dou 8 sheng, rice brewed, 358 dan 2 sheng.\end{uniten}

\begin{unitzh}\unitlabel{shushu-043.10}術曰：以粟米換易求之。置諸率，隨本色對列如鴈翅。有分者通之，異類者變之。以頭位者進乗之，以下位退乗之，得合數。有對者相乗之，無對者直命之，為諸率。併上下無對者，為法率{\zhsmall 諸率可約者，又約之}。以今有物徧乘諸率{\zhsmall 不乗法率}，各為實。諸實並如法而一，各得。其已變者，復互易乗除之，即得所求。\end{unitzh}
\begin{uniten}Method: Solve it by "Millet and Grain" exchange. Set out the various rates, arraying them against each other according to their own kinds, like the wings of a wild goose. Where there are fractions, make them common (tong); where kinds are different, transform them (bian). With the head position multiply advancing (jincheng); with the lower position multiply retreating (tuicheng); this gives the combined numbers (heshu). Those that have a counterpart (dui) are multiplied together; those that have no counterpart are taken directly as they are (zhiming); these are the various rates. Add the upper and lower ones that have no counterpart to make the divisor rate (falü) [Note: Where the various rates can be reduced, reduce them further.]. With the quantity of goods now on hand multiply the various rates throughout [Note: The divisor rate is not multiplied.], each making a dividend (shi). Divide all the dividends by the divisor, and each is obtained. Those that have been transformed are again exchanged (huyi) by multiplying and dividing; then what is sought is obtained.\end{uniten}

\begin{unitzh}\unitlabel{shushu-043.11}草曰：置糯穀七、出米三於右行上、副兩位；次置糯米一斗、麥一斗七升於副行副、中兩位；次置小麥五升、踏麥二斤四兩於次行中、次兩位；次置麥一百一十觔、醞米一石三斗於左行次、下兩位，隨本色對列如鴈翅。訖，乃驗次行二斤四兩是四分斤之一，以母四通次行兩位，以子一内次行次位，具中位得二，次位得九。又驗左行下位是糯米，是異類，於糯米合變為糯穀。乃以問中首句率穀七、米三變之：以七因米一石三斗，得九石一斗，於左下，為穀；却以米三因麴一百一十斤，為三百三十斤麴，於左行。得變圖數。\end{unitzh}
\begin{uniten}Working: Set the glutinous grain 7 and the rice yielded 3 in the top (shang) and deputy (fu) positions of the right column; next, set the glutinous rice 1 dou and the wheat 1 dou 7 sheng in the deputy and middle (zhong) positions of the deputy column; next, set the wheat 5 sheng and the trodden wheat [Translator's note: here and in the next clause 麥 "wheat" stands where 麴 "leaven" is expected.] 2 jin 4 liang in the middle and next (ci) positions of the next column; next, set the wheat 110 jin (觔) and the brewed rice 1 dan 3 dou in the next and bottom (xia) positions of the left column, arraying them against each other according to their own kinds like the wings of a wild goose. When this is done, examine: the 2 jin 4 liang of the next column is a case of 1/4 jin; with the denominator 4 make common (tong) the two positions of the next column, and put the numerator 1 into (na) the next position of the next column; altogether the middle position gets 2 and the next position gets 9. Further examine: the bottom position of the left column is glutinous rice, which is a different kind; the glutinous rice ought to be transformed into glutinous grain. Then transform it by the rate in the first clause of the question, grain 7 and rice 3: with 7 multiply (yin) the rice 1 dan 3 dou, obtaining 9 dan 1 dou, at the lower left, as grain; then with the rice 3 multiply the leaven 110 jin, making 330 jin of leaven, in the left column. This gives the numbers of the transformed diagram (biantu).\end{uniten}

\begin{unitzh}\unitlabel{shushu-043.12}以左行三百三十乗次行二，得六百六十；次以六百六十乘副行一十，得六千六百；次以六千六百乗右上七，得四萬六千二百：各於原位。却以右行副位三因副行一斗七升，得五斗一升；又以五斗一升乗次行九，得四百五十九；又以四百五十九乗左下九十一，得四萬一千七百六十九。列為合圖數。乃驗合圖四行：其副、中、次三位有對，以對相乗合之；其右上、左下無對者，直命之：皆為率。列右行上得四萬六千二百，為出糯穀率；副位得一萬九千八百，為得糯米率；中得三萬三千六百六十，為易得麥率；次得一萬五千一百四十七，為踏到麴率；下得四萬一千七百六十九，為餘下糯榖率。併上下率，共得八萬七千九百六十九，為法率。今六率共求等，得一，約之，只得原率，為率圖。\end{unitzh}
\begin{uniten}With the 330 of the left column multiply the 2 of the next column, obtaining 660; next, with the 660 multiply the 10 of the deputy column, obtaining 6600; next, with the 6600 multiply the 7 at the upper right, obtaining 46,200: each in its original position. Then with the 3 in the deputy position of the right column multiply the 1 dou 7 sheng of the deputy column, obtaining 5 dou 1 sheng; further, with the 5 dou 1 sheng multiply the 9 of the next column, obtaining 459; further, with the 459 multiply the 91 at the lower left, obtaining 41,769. Lay these out as the numbers of the combined diagram (hetu). Then examine the four columns of the combined diagram: their three positions deputy, middle and next have counterparts, and counterpart is multiplied by counterpart to combine them; the upper right and lower left, which have no counterparts, are taken directly: all are rates. Listed: the top of the right column gives 46,200, the rate for the glutinous grain taken out; the deputy position gives 19,800, the rate for the glutinous rice obtained; the middle gives 33,660, the rate for the wheat obtained by exchange; the next gives 15,147, the rate for the leaven trodden; the bottom gives 41,769, the rate for the remaining glutinous grain. Adding the top and bottom rates gives in all 87,969, the divisor rate. Now, seeking the common measure (deng) of the six rates together, one obtains 1; reducing by it, one obtains just the original rates, which make the rate diagram (lütu).\end{uniten}

\begin{unitzh}\unitlabel{shushu-043.13}始用今有穀一千七百五十九石三斗八升，皆化為升，徧乘五率，不乗法率：得八十一億二千八百三十三萬五千六百升，為出穀實；得三十四億八千三百五十七萬二千四百升，為糯米實；得五十九億二千二百七萬三千八十升，為易麥實；得二十六億六千四百九十三萬二千八百八十六，為踏麴實；得七十三億四千八百七十五萬四千三百二十二升，為餘穀實。其五實皆如法八萬七千九百六十九而一：得九百二十四石，為出穀；得三百九十六石，為做到糯米；得六百七十三石二斗，為易到小麥；得三萬二百九十四斤，為踏到麴；得八百三十五石三斗八升，為餘下穀。今將餘下穀變為米，乃以乗率三因餘穀八百三十五石三斗八升，得二千五百六石一斗四升，為實；以糯穀率七為法，除之，得三百五十八石二升，為醞米。\end{unitzh}
\begin{uniten}Only then take the grain now on hand, 1759 dan 3 dou 8 sheng, convert it all into sheng, and multiply the five rates throughout, not multiplying the divisor rate: obtaining 8,128,335,600 sheng, the dividend for the grain taken out; obtaining 3,483,572,400 sheng, the dividend for the glutinous rice; obtaining 5,922,073,080 sheng, the dividend for the wheat by exchange; obtaining 2,664,932,886, the dividend for the trodden leaven; obtaining 7,348,754,322 sheng, the dividend for the remaining grain. Divide each of the five dividends by the divisor, 87,969: obtaining 924 dan, the grain taken out; obtaining 396 dan, the glutinous rice made; obtaining 673 dan 2 dou, the wheat obtained by exchange; obtaining 30,294 jin, the leaven trodden; obtaining 835 dan 3 dou 8 sheng, the remaining grain. Now, to transform the remaining grain into rice: with the multiplying rate 3 multiply the remaining grain, 835 dan 3 dou 8 sheng, obtaining 2506 dan 1 dou 4 sheng, as the dividend; with the glutinous-grain rate 7 as the divisor, divide, obtaining 358 dan 2 sheng: the rice brewed.\end{uniten}

\begin{unitzh}\unitlabel{shushu-043.14}按：術中互乗、進乗、退乗、對乗，皆通分法也。張邱建云：「學者不患乗除之為難，而患通分之為難。」此術曲盡其妙，今各釋於後。\end{unitzh}
\begin{uniten}Note: The cross-multiplication, advancing multiplication, retreating multiplication and counterpart multiplication (duicheng) in the method are all methods of making fractions common (tongfen). Zhang Qiujian says: "Learners do not worry that multiplication and division are difficult; they worry that making fractions common is difficult." This method brings out its subtleties to the full; now each is explained below.\end{uniten}

\begin{unitzh}\unitlabel{shushu-043.15}第一圖互乗：以右上穀七乗左下米一十三，得九十一，應以右下米三除之，方得穀數；今不除，便如得穀數又以米三乗之矣。故以米三乗麴一十一，得三十三，與穀數九十一相當，仍同於麴一十一與米一十三相當也。\end{unitzh}
\begin{uniten}The first diagram, cross-multiplication: with the grain 7 at the upper right multiply the rice 13 at the lower left, obtaining 91; this ought to be divided by the rice 3 at the lower right, and only then is the number for grain obtained. Now, not dividing is just as if the number for grain had been obtained and then further multiplied by the rice 3. Therefore with the rice 3 multiply the leaven 11, obtaining 33, which corresponds to the grain number 91; this is still the same as the leaven 11 corresponding to the rice 13.\end{uniten}

\begin{unitzh}\unitlabel{shushu-043.16}第二圖：左進乗，自下而上；右退乗，自上而下。左四位連乗至穀，應以右上三位連乗之數除之，得踏麴三十三所用穀數；今不除，為寄右上三位連乗之數為分母。右四位連乘至穀，即如麴三十三所醞穀數，同寄右上三位連乗之數為分母也。併之，即如踏麴用穀、麴所醞穀總數，同寄右上三位連乗之分也。\end{unitzh}
\begin{uniten}The second diagram: on the left, advancing multiplication, from bottom to top; on the right, retreating multiplication, from top to bottom. The four positions on the left multiplied in succession (liancheng) up to the grain ought to be divided by the number from multiplying in succession the upper three positions on the right, to obtain the number of grain used in treading the leaven 33; now, not dividing amounts to lodging (ji) the number from multiplying in succession the upper three positions on the right as a denominator. The four positions on the right multiplied in succession up to the grain are then equivalent to the number of grain brewed by the leaven 33, likewise lodging the number from multiplying in succession the upper three positions on the right as the denominator. Adding them is then equivalent to the total number of the grain used in treading the leaven and the grain brewed by the leaven, likewise lodging the fraction (fen) from multiplying in succession the upper three positions on the right.\end{uniten}

\begin{unitzh}\unitlabel{shushu-043.17}第三圖對乗：左下一位麴數，以前寄分母右上三位連乗之數乗之，即同寄一分母也。左下二位相乗，應以右上第三位除之，得踏麴所用麥數；不除，為寄右上第三位麴數，又以右上二位相乘之數乗之，是應用麥數内同寄右上三位連乘之分母也。左下三位連乘，應以右中二位相乗之數除之，得易麥應用米數；不除，為寄右中二位相乗之數，又以右上第一位米數乗之，是應用米數内同寄右上三位連乗之分母也。右總穀、出穀、餘穀、易米、易麥、踏麴六數，皆同寄一分母，則用以乗除求得數，即與本數無異。故以題中總穀數乗寄分母各數，以寄分母總數除之，即得所求各數也。\end{unitzh}
\begin{uniten}The third diagram, counterpart multiplication: the leaven number in the one position at the lower left is multiplied by the previously lodged denominator, the number from multiplying in succession the upper three positions on the right; then it likewise lodges one and the same denominator. The two positions at the lower left multiplied together ought to be divided by the third position at the upper right, to obtain the number of wheat used in treading the leaven; not dividing amounts to lodging the leaven number of the third position at the upper right; and further multiplying it by the number from multiplying together the upper two positions on the right, this is the number of wheat that ought to be used, within which the denominator from multiplying in succession the upper three positions on the right is likewise lodged. The three positions at the lower left multiplied in succession ought to be divided by the number from multiplying together the two middle positions on the right, to obtain the number of rice that ought to be used in exchanging for wheat; not dividing amounts to lodging the number from multiplying together the two middle positions on the right; and further multiplying it by the rice number of the first position at the upper right, this is the number of rice that ought to be used, within which the denominator from multiplying in succession the upper three positions on the right is likewise lodged. The six numbers above, total grain, grain taken out, remaining grain, rice for exchange, wheat by exchange and trodden leaven, all lodge one and the same denominator; so the numbers obtained by using them in multiplication and division are no different from the original numbers. Therefore, with the total grain number in the problem multiply each of the lodged-denominator numbers, and divide by the lodged-denominator total; then each of the numbers sought is obtained.\end{uniten}

\begin{unitzh}\unitlabel{shushu-043.18}再，此術不獨法之巧，即圖式布置亦皆皆具精義，熟玩之，可以得其往來變通之故。原第三合圖仍斜排，為上圖；第四率圖即各寄寄母數直列，為下圖。今合圖改為正圖，列於右；各得寄母數並列於左。\end{unitzh}
\begin{uniten}Further: in this method it is not only that the method is ingenious; even the arrangement of the diagrams also all has refined meaning [Translator's note: 皆 "all" is written twice here, and 寄 "lodged" twice below.]; studying them thoroughly, one can grasp the reasons for their transformations back and forth. Originally the third, combined diagram was still arranged obliquely, as the upper diagram; the fourth, rate diagram, namely the numbers each with lodged denominator (jimu), was laid out straight, as the lower diagram. Now the combined diagram has been changed into an upright diagram, laid out on the right; the numbers each obtained with lodged denominator are laid out side by side on the left.\end{uniten}

\begin{unitzh}\unitlabel{shushu-043.19}𮅕囘運費\end{unitzh}
\begin{uniten}𮅕 Return-Transport Costs [Translator's note: 𮅕 is a rare character of uncertain meaning.]\end{uniten}

\begin{unitzh}\unitlabel{shushu-044.1}問：有江西水運米一十二萬三千四百石，原係鎮江交卸，計水程二千一百三十里，每石水脚錢一貫二百文。今截上件米就池州安頓，池州至鎮江八百八十里。欲收囘不該水脚錢幾何？答曰：收囘錢六萬一千一百七十八貫五百九十一文。\end{unitzh}
\begin{uniten}Question: There is Jiangxi rice transported by water, 123,400 dan (石), originally to be handed over and unloaded at Zhenjiang, the water route being reckoned at 2130 li (里); the water-freight money (shuijiao qian) for each dan is 1 guan (貫) 200 wen (文). Now the above rice is intercepted and stored at Chizhou; from Chizhou to Zhenjiang is 880 li. One wishes to recover the water-freight money that is not due: how much is it? Answer: The recovered money is 61,178 guan 591 wen.\end{uniten}

\begin{unitzh}\unitlabel{shushu-044.2}術曰：以粟米互易求之。置池州至鎮江里數，乗水脚錢，得數又乗運米，為實；以原至鎮江水程為法；除實，得收囘錢。\end{unitzh}
\begin{uniten}Method: Seek it by the mutual exchange (huyi) of 'Millet and rice' (sumi). Set up the number of li from Chizhou to Zhenjiang and multiply by the water-freight money; the number obtained is again multiplied by the transported rice, making the dividend (shi). Take the original water route to Zhenjiang as the divisor (fa). Divide the dividend, obtaining the recovered money.\end{uniten}

\begin{unitzh}\unitlabel{shushu-044.3}草曰：置池州至鎮江八百八十里，乗毎石水脚錢一貫二百，得一千五十六貫文；又乗運米一十二萬三千四百石，得一億三千三十一萬四百貫文，為實；以原至鎮江水程二千一百三十里為法；除實，得六萬一千一百七十八貫五百九十一文，為收囘錢。\end{unitzh}
\begin{uniten}Working: Set up Chizhou to Zhenjiang, 880 li; multiply by the water-freight money for each dan, 1 guan 200, obtaining 1056 guan wen; again multiply by the transported rice, 123,400 dan, obtaining 130,310,400 guan wen, as the dividend. Take the original water route to Zhenjiang, 2130 li, as the divisor. Divide the dividend, obtaining 61,178 guan 591 wen as the recovered money.\end{uniten}

\begin{unitzh}\unitlabel{shushu-044.4}三合均價{\zhsmall 按舊本此問無題，今増。}\end{unitzh}
\begin{uniten}Averaging the Price of Three Combined [Siku editors: In the old edition this problem has no title; it is now added.]\end{uniten}

\begin{unitzh}\unitlabel{shushu-044.5}問：庫有三色金共五千兩，内八分金一千二百五十兩，兩價四百貫文；七分五釐金一千六百兩，兩價三百七十五貫文；八分五釐金二千一百五十兩，兩價四百二十五貫文。並欲煉為足色，每兩工食藥炭錢三貫文，耗金九百七十二兩五錢。欲知色分及兩價各幾何？答曰：色十分，兩價五百三貫七百二十四文五百三十七分文之二百一十二。\end{unitzh}
\begin{uniten}Question: A treasury has gold of three colours (se), together 5000 liang (兩). [Translator's note: the colour (se) of gold is its fineness, counted in tenths (fen) and hundredths (li).] Of this, gold of 8 fen (分) is 1250 liang, price per liang 400 guan wen; gold of 7 fen 5 li (釐) is 1600 liang, price per liang 375 guan wen; gold of 8 fen 5 li is 2150 liang, price per liang 425 guan wen. All are to be refined to full colour; for each liang the money for labour, food, drugs and charcoal is 3 guan wen, and the gold lost is 972 liang 5 qian (錢). One wishes to know the colour in fen and the price per liang: how much is each? Answer: Colour 10 fen; price per liang 503 guan 724 212/537 wen.\end{uniten}

\begin{unitzh}\unitlabel{shushu-044.6}術曰：以方田及粟米求之。置共數，以耗減之，餘為法；以三色分數各乗兩數，併之，為色分實；以三色價數各乗兩數，為寄；以工藥價乗共金，併寄，共為價實。二實皆如法而一，即各得。\end{unitzh}
\begin{uniten}Method: Seek it by 'Rectangular fields' (fangtian) and 'Millet and rice'. Set up the total and subtract the loss from it; the remainder is the divisor. With the fen numbers of the three colours multiply the respective numbers of liang and combine them, making the dividend for the colour fen. With the prices of the three colours multiply the respective numbers of liang, making the deposit (ji). With the labour-and-drug price multiply the total gold, and combine it with the deposit; together this makes the dividend for the price. Divide both dividends by the divisor; then each is obtained.\end{uniten}

\begin{unitzh}\unitlabel{shushu-044.7}草曰：置共金五千兩，減耗九百七十二兩五錢外，餘四千二十七兩五錢，為法。次置一千二百五十兩，乘八分，得一萬分於上；置一千六百兩，以七分五釐乗之，得一萬二千分，加上；置二千一百五十兩，乗八分五釐，得一萬八千二百七十五分，又加上；共得四萬二百七十五分，為分實。次置一千二百五十兩，乗四百貫，得五十萬貫，為寄；次置一千六百兩，乘價三百七十五貫，得六十萬貫，加寄；次置二千一百五十兩，乗價四百二十五貫，得九十一萬三千七百五十貫，又加寄；次置共金五千兩，乗工藥錢三貫，得一萬五千貫，又加寄；共得二百二萬八千七百五十貫，為貫實。二實並如法四千二十七兩五錢而一，得其色一十分；其價每兩得五百三貫七百二十四文，不盡一貫五百九十，與法求等，得七十五，俱約之，為五百三十七文之二百一十二。\end{unitzh}
\begin{uniten}Working: Set up the total gold, 5000 liang; subtract the loss, 972 liang 5 qian; beyond that there remain 4027 liang 5 qian, as the divisor. Next set up 1250 liang, multiply by 8 fen, obtaining 10,000 fen, on top. Set up 1600 liang, multiply it by 7 fen 5 li, obtaining 12,000 fen; add to the top. Set up 2150 liang, multiply by 8 fen 5 li, obtaining 18,275 fen; again add to the top. Together this gives 40,275 fen as the fen dividend. Next set up 1250 liang, multiply by 400 guan, obtaining 500,000 guan, as the deposit. Next set up 1600 liang, multiply by the price, 375 guan, obtaining 600,000 guan; add to the deposit. Next set up 2150 liang, multiply by the price, 425 guan, obtaining 913,750 guan; again add to the deposit. Next set up the total gold, 5000 liang, multiply by the labour-and-drug money, 3 guan, obtaining 15,000 guan; again add to the deposit. Together this gives 2,028,750 guan as the guan dividend. Divide both dividends by the divisor, 4027 liang 5 qian, obtaining its colour, 10 fen; its price per liang is obtained as 503 guan 724 wen, with 1 guan 590 not exhausted (bujin). Seek the equal (deng) with the divisor, obtaining 75; reduce (yue) both by it, making 212/537 wen.\end{uniten}

\begin{unitzh}\unitlabel{shushu-044.8}子部·天文算法類·算書之屬·數學九章·卷六下\end{unitzh}
\begin{uniten}Masters section (zibu) · Category of astronomy and calculation methods (tianwen suanfa lei) · Group of calculation books (suanshu zhi shu) · Shuxue jiuzhang · juan 6, lower part\end{uniten}

\begin{unitzh}\unitlabel{shushu-044.9}子部·天文算法類·算書之屬·數學九章·卷六下\end{unitzh}
\begin{uniten}Masters section · Category of astronomy and calculation methods · Group of calculation books · Shuxue jiuzhang · juan 6, lower part\end{uniten}

\begin{unitzh}\unitlabel{shushu-044.10}其色實餘盡，得一十分，為金色；其價除得五百三貫七百二十四文，為十分金毎兩價；不盡一貫五百九十文，與法求等，得七半，俱以約之，為五百三十七分文之二百一十二。\end{unitzh}
\begin{uniten}Its colour dividend, the remainder being exhausted, gives 10 fen as the colour of the gold. Its price, divided, gives 503 guan 724 wen as the price per liang of 10-fen gold; not exhausted, 1 guan 590 wen: seek the equal with the divisor, obtaining 7 1/2; reduce both by it, making 212/537 wen.\end{uniten}

\begin{unitzh}\unitlabel{shushu-044.11}《數學九章》卷六下\end{unitzh}
\begin{uniten}Shuxue jiuzhang, juan 6, lower part\end{uniten}

\begin{unitzh}\unitlabel{shushu-044.12}子部·天文算法類·算書之屬·數學九章\end{unitzh}
\begin{uniten}Masters section · Category of astronomy and calculation methods · Group of calculation books · Shuxue jiuzhang\end{uniten}

\section*{卷7上}\addcontentsline{toc}{section}{卷7上}

\begin{unitzh}\unitlabel{shushu-045.1}欽定四庫全書\end{unitzh}
\begin{uniten}Qinding Siku quanshu (Imperially Approved Complete Library of the Four Treasuries)\end{uniten}

\begin{unitzh}\unitlabel{shushu-045.2}《數學九章》卷七上 宋 秦九韶 撰\end{unitzh}
\begin{uniten}Shuxue jiuzhang (Mathematics in Nine Sections), juan 7, first part; composed by Qin Jiushao of the Song\end{uniten}

\begin{unitzh}\unitlabel{shushu-045.3}營建\end{unitzh}
\begin{uniten}Construction (yingjian)\end{uniten}

\begin{unitzh}\unitlabel{shushu-045.4}計作清臺\end{unitzh}
\begin{uniten}Reckoning the building of an observatory terrace (qingtai)\end{uniten}

\begin{unitzh}\unitlabel{shushu-045.5}問：剏築清臺一所，正高一十二丈，上廣五丈，袤七丈，下廣一十五丈，袤一十七丈，其袤當東西，廣當南北。秋程人日行六十里，里法三百六十步。钁土、鍬土每工各二百尺，築土每工九十尺。每擔土壤一尺三寸，徃來一百六十步，内四十步上下棚道。築高至少半，其棚道三當平道五；至中半，三當七；至太半，二當五。踟蹰之間，十加一；載輸之間，二十步，定一返。今甲、乙、丙三縣差夫：甲縣附郭，税力一十三萬三千八百六十六；乙縣去臺所一百二十里，税力二十三萬七千九百八十四；丙縣去臺所一百八十里，税力三十一萬二千三百五十四。俱以道里逺近、税力多少均科之。\end{unitzh}
\begin{uniten}Question: An observatory terrace is newly built, one structure. Its vertical height is 12 zhang (丈); at the top the breadth (guang) is 5 zhang and the length (mao) 7 zhang; at the bottom the breadth is 15 zhang and the length 17 zhang. Its length lies east–west and its breadth north–south. By the autumn schedule (qiucheng) a person travels 60 li (里) a day; the li standard (lifa) is 360 bu (步). Mattocking earth and shoveling earth are each 200 chi (尺) per workday (gong); ramming earth is 90 chi per workday. Each load of earth is 1 chi 3 cun (寸) of loose earth (rang); going and coming is 160 bu, of which 40 bu is up and down the scaffold path (pengdao). When the rammed height reaches 1/3 (shaoban), on the scaffold path 3 equal 5 on the level path (pingdao); when it reaches 1/2 (zhongban), 3 equal 7; when it reaches 2/3 (taiban), 2 equal 5. For the interval of dawdling (chichu), add 1 in 10; for the interval of loading and delivering (zaishu), 20 bu; fix one round trip (ding yi fan). Now three counties, A, B and C (jia, yi, bing), are to dispatch laborers (chaifu). County A adjoins the city wall (fuguo), with a tax strength (shuili) of 133,866; County B is 120 li from the terrace site, with a tax strength of 237,984; County C is 180 li from the terrace site, with a tax strength of 312,354. In all cases apportion it evenly (junke) according to the distance of the road and the amount of the tax strength.\end{uniten}

\begin{unitzh}\unitlabel{shushu-045.6}臺下舖石脚七層。先用甎砌臺身，次用甎疊砌轉道，周圍五帶，並濶六尺。須令南北二平道、東西三峻道相間。始自臺之艮隅，於東外道向南順升，由巽隅以西左轉，週迴厯北，復東，再升東裏道，至巽隅，乃登臺頂。其東裏道艮隅與北平道兩隅及西道乾隅之高，皆以強半；其西道坤隅與南道兩隅、東外道異隅之高，皆以五分之二。峻道毎級履高六寸。其東裏道級數取弱半，東外道級數取五分之二，西道級數取強半。石長五尺，濶二尺，厚五寸；甎長一尺二寸，濶六寸，厚二寸五分。欲知土積、定一返步、毎工人到土及總用功、各縣起夫、甎石、峻平道髙長、級數、踏蹤各幾何？\end{unitzh}
\begin{uniten}Below the terrace lay a stone footing (shijiao) of 7 layers. First build the body of the terrace with bricks; next, with bricks laid in courses, build the turning path (zhuandao) around it in 5 bands, all 6 chi wide. The 2 level paths on the south and north and the 3 steep paths (jundao) on the east and west must alternate. Starting from the northeast (gen) corner of the terrace, on the east outer path one ascends steadily toward the south; from the southeast (xun) corner one turns left westward, circles round through the north, returns to the east, ascends again by the east inner path to the southeast corner, and then mounts the top of the terrace. The heights of the northeast corner of the east inner path, of both corners of the north level path and of the northwest (qian) corner of the west path are all at 3/4 (qiangban); the heights of the southwest (kun) corner of the west path, of both corners of the south path and of the yi (異) corner of the east outer path are all at 2/5 [Translator's note: 異 is evidently an error for 巽, the southeast corner]. On the steep paths each step rises (lügao) 6 cun. For the east inner path the number of steps takes 1/4 (ruoban); for the east outer path the number of steps takes 2/5; for the west path the number of steps takes 3/4. The stones are 5 chi long, 2 chi wide and 5 cun thick; the bricks are 1 chi 2 cun long, 6 cun wide and 2 cun 5 fen (分) thick. One wishes to know the volume of earth, the bu of the fixed round trip, the earth delivered (daotu) by each worker per workday and the total labor used, the laborers raised by each county, the bricks and stones, and the heights and lengths, numbers of steps and tread lengths (tazong) of the steep and level paths: how much is each?\end{uniten}

\begin{unitzh}\unitlabel{shushu-045.7}答曰：土積一百五十四萬尺。定一返二百步一十八分步之五。{\zhsmall 按：此條起共十一條數皆誤。}毎工人到土一百四十尺七百二十一分尺之一百四十八。總用工四萬五千六百八十六工。甲縣差一萬七千一百三十二工，乙縣差一萬五千二百二十九工，丙縣差一萬三千三百二十五工。石四千三百一十七片。甎一百一十四萬二千二十四片。東裏道峻，艮隅高九丈，巽隅高一十一丈九尺四寸，級五十踏，踏蹤二尺二寸九分。東外道峻，艮隅高六寸，巽隅高四丈八尺，級八十踏，踏蹤二尺一寸五分。西峻道坤隅高四丈八尺，乾隅高九丈，級七十踏，踏蹤二尺一寸四分寸之九。南平道高四丈八尺。北平道九丈。\end{unitzh}
\begin{uniten}Answer: The volume of earth is 1,540,000 chi. The fixed round trip is 200 5/18 bu. [Siku editors: From this item on, the numbers of altogether 11 items are all erroneous.] The earth delivered by each worker per workday is 140 148/721 chi. The total labor used is 45,686 workdays. County A is assigned 17,132 workdays, County B is assigned 15,229 workdays, County C is assigned 13,325 workdays. Stone: 4317 pieces (pian). Bricks: 1,142,024 pieces. The east inner path, steep: the northeast corner is 9 zhang high, the southeast corner is 11 zhang 9 chi 4 cun high; steps: 50 treads; tread length 2 chi 2 cun 9 fen. The east outer path, steep: the northeast corner is 6 cun high, the southeast corner is 4 zhang 8 chi high; steps: 80 treads; tread length 2 chi 1 cun 5 fen. The west steep path: the southwest corner is 4 zhang 8 chi high, the northwest corner is 9 zhang high; steps: 70 treads; tread length 2 chi 1 9/4 cun. The south level path is 4 zhang 8 chi high. The north level path: 9 zhang.\end{uniten}

\begin{unitzh}\unitlabel{shushu-045.8}按：右答數中，惟第一條土積、未二條南北平道高三數無誤，餘數俱有誤處。其所以不合之故，具詳草後。\end{unitzh}
\begin{uniten}Note: Among the answer numbers above, only the first item, the volume of earth, and the last two items [Translator's note: the text has 未, evidently for 末 "last"], the heights of the south and north level paths, three numbers in all, are without error; all the remaining numbers have points of error. The reasons why they do not agree are given in detail after the working.\end{uniten}

\begin{unitzh}\unitlabel{shushu-045.9}術曰：以啇功求之，均輸入之。倍臺上袤，加下袤，乗上廣，為寄；倍下袤，加上袤，乗下廣，併寄，乗高，為土率；如六而一，得堅積。以築工尺為法，除堅積，得築工。以穿率乗堅積為實，以堅率乗钁鍬工尺，半之，為鍬法，除實，得钁鍬共工。以壤率因堅積，如堅率而一，為壤積。\end{unitzh}
\begin{uniten}Method: Seek it by Assessing Works (shanggong) and enter it by Fair Transport (junshu). Double the upper length of the terrace, add the lower length, multiply by the upper breadth: this is the reserve (ji). Double the lower length, add the upper length, multiply by the lower breadth, add the reserve, multiply by the height: this is the earth rate (tulü). Divide by 6, obtaining the firm volume (jianji). Take the chi of a ramming workday as divisor (fa) and divide the firm volume, obtaining the ramming labor (zhugong). Multiply the firm volume by the dug rate (chuanlü) as dividend (shi); multiply the chi of a mattock-and-spade workday by the firm rate (jianlü) and halve it, making the spade divisor (qiaofa); divide the dividend, obtaining the combined labor of mattock and spade. Multiply (yin) the firm volume by the loose rate (ranglü) and divide by the firm rate: this is the loose volume (rangji).\end{uniten}

\begin{unitzh}\unitlabel{shushu-045.10}求負土者：先列全分及等至高諸母子，以母互乗諸子，為寄左行；以諸母相乗，為寄母。次列棚道全分及所當鮮母衍子，以鮮母互乗衍子，為左行；以鮮母相乗，以乗寄，得數，又以列位乗之，為總母。以左右兩行諸子對乗之，併之，為總子。其總母求等約之，為定母子。以定子乘棚道為次，以定母乗平道，加次，义以踟蹰之數身下加之；又以載輸步乗定母，併次，為統數。以定母除統，得定一返步，亦為列土法。{\zhsmall 有分復通為法。}置程里通步，乗擔土尺，有歩母則又以步母乗之，為到實。實如法而一，得每工人到土，亦為壤法。以除壤積，得負土工。併前鍬、築二工，為總用工。\end{unitzh}
\begin{uniten}To find the earth-carrying (futu): first set out the whole (quanfen) and the denominators and numerators of the grades reached in height; with the denominators cross-multiply (hucheng) the numerators, making the reserved left column; multiply the denominators together, making the reserve denominator (jimu). Next set out the whole of the scaffold path and the scant denominators (xianmu) and extended numerators (yanzi) of what it equals; with the scant denominators cross-multiply the extended numerators, making the left column; multiply the scant denominators together and with this multiply the reserve; multiply the number obtained again by the listed positions (liewei): this is the total denominator (zongmu). Multiply correspondingly (duicheng) the numerators of the two columns, left and right, and add them: this is the total numerator (zongzi). For the total denominator seek the common measure (qiu deng) and reduce (yue) by it, making the fixed denominator and numerator (dingmu, dingzi). With the fixed numerator multiply the scaffold path: this is the secondary (ci). With the fixed denominator multiply the level path and add the secondary; again [Translator's note: the text has 义, evidently for 又 "again"], with the number for dawdling, add to it below its body (shenxia jia); again with the bu for loading and delivering multiply the fixed denominator and add the secondary: this is the comprehensive number (tongshu). With the fixed denominator divide the comprehensive, obtaining the bu of the fixed round trip; this is also the divisor for setting out earth (lietu fa). [Note: If there is a fraction, convert (tong) it again to make the divisor.] Set out the li of the day's journey (chengli), convert them to bu and multiply by the chi of earth per load; if there is a bu denominator (bumu), then multiply also by the bu denominator: this is the delivery dividend (daoshi). Divide the dividend by the divisor, obtaining the earth delivered by each worker per workday; this is also the loose divisor (rangfa). With it divide the loose volume, obtaining the earth-carrying labor. Add the two labors from before, of spade and of ramming: this is the total labor used.\end{uniten}

\begin{unitzh}\unitlabel{shushu-045.11}以各縣日程數約税力，各得力率，副併為科法。以共用工乗未併者，各為科實。實如法而一，各得縣夫。\end{unitzh}
\begin{uniten}With the numbers of days' journey (richeng) of the respective counties reduce the tax strengths, obtaining in each case the strength rate (lilü); add them together in a side position (fubing), making the apportionment divisor (kefa). With the combined labor used multiply the ones not added, each making an apportionment dividend (keshi). Divide the dividend by the divisor, obtaining in each case the laborers of the county.\end{uniten}

\begin{unitzh}\unitlabel{shushu-045.12}求甎者：倍轉道濶，遍加臺上下廣袤，變名上下濶長。以甎厚加臺高，為臺直。次列甎石長濶厚，各相乗，為甎石積法。通廣袤如法，乃倍上長，加下長，乗上濶，為寄；次倍下長，加上長，乗濶，併寄，共乗直，得數減土率，除如六而一，為泛。\end{unitzh}
\begin{uniten}To find the bricks: double the width of the turning path and add it throughout to the upper and lower breadths and lengths of the terrace; change their names to upper and lower width (kuo) and length (chang). Add the thickness of a brick to the height of the terrace: this is the terrace vertical (taizhi). Next set out the length, width and thickness of the brick and of the stone and multiply each together: these are the volume divisors (jifa) of brick and stone. Convert the breadths and lengths according to the rule; then double the upper length, add the lower length and multiply by the upper width: this is the reserve. Next double the lower length, add the upper length, multiply by the width and add the reserve; multiply all together by the vertical; from the number obtained subtract the earth rate and divide by 6: this is the provisional (fan).\end{uniten}

\begin{unitzh}\unitlabel{shushu-045.13}置南北道高子，各乗臺高為實，如各母而一，得五道諸隅高。以北道高減臺，餘為上停高；以南道高減北道高，餘為中停高；命南道高為下停高。以履寸除諸高，得級數。以上長減下長，餘半之，為勾。以勾乗南道高為實，如臺高而一，得底率。以底率減下長，餘為底股。以外道級數約底股，得外道踏蹤。又以底率減底股，餘為中股。以勾乗中停高為實，如臺高而一，得中率。以中率減中股，餘為上股。以西道級除上股，得西道踏蹤。以勾乗上停高為實，如臺高而一，得中率。以中率減中股，餘為上股。以西道級除上股，得西道踏級。以勾乗上停高為實，如臺高而一，得上率。以上率併中率，共減上股，餘為實，如裏道級而一，得裏道踏蹤。\end{unitzh}
\begin{uniten}Set out the numerators of the heights of the south and north paths, multiply each by the terrace height as dividend, and divide by the respective denominators, obtaining the heights of all the corners of the five paths. Subtract the north path height from the terrace; the remainder is the upper-stage height (shangting gao). Subtract the south path height from the north path height; the remainder is the middle-stage height. Designate the south path height as the lower-stage height. Divide the heights by the cun of the step rise, obtaining the numbers of steps. Subtract the upper length from the lower length and halve the remainder: this is the base (gou). Multiply the south path height by the base as dividend and divide by the terrace height, obtaining the bottom rate (dilü). Subtract the bottom rate from the lower length; the remainder is the bottom altitude (digu). Reduce the bottom altitude by the number of steps of the outer path, obtaining the tread length of the outer path. Again subtract the bottom rate from the bottom altitude; the remainder is the middle altitude (zhonggu). Multiply the middle-stage height by the base as dividend and divide by the terrace height, obtaining the middle rate (zhonglü). Subtract the middle rate from the middle altitude; the remainder is the upper altitude (shanggu). Divide the upper altitude by the steps of the west path, obtaining the tread length of the west path. Multiply the upper-stage height by the base as dividend and divide by the terrace height, obtaining the middle rate. Subtract the middle rate from the middle altitude; the remainder is the upper altitude. Divide the upper altitude by the steps of the west path, obtaining the tread steps (taji) of the west path. [Translator's note: these three sentences evidently repeat the three preceding ones.] Multiply the upper-stage height by the base as dividend and divide by the terrace height, obtaining the upper rate (shanglü). Add the upper rate and the middle rate and subtract them together from the upper altitude; the remainder is the dividend; divide by the steps of the inner path, obtaining the tread length of the inner path.\end{uniten}

\begin{unitzh}\unitlabel{shushu-045.14}次以道濶併下長為補，以南北道高併臺高乗補，為需。次上廣減下廣，餘為址。以址乗南道高為實，以臺高除之，得數減下廣，餘為南道長。以南長併下廣，乗南道高，加需。又以址乗北道高為實，以臺高除之，得數減下廣，餘為北道長。以北道長併下廣，乗北道高，又加需，共乗半道濶，得數併泛，為共率。以基脚層數乗石版厚，為基高。次倍道濶，併下濶，乗下長，為基率。次以下廣乗下袤，減基率，餘乗基高，為石率。以石率減共率，餘為甎率。以甎積法除甎率，得甎數；以石級積法除石率，得石版數。\end{unitzh}
\begin{uniten}Next add the path width to the lower length: this is the supplement (bu 補). Add the heights of the south and north paths to the terrace height and multiply the supplement by it: this is the requisite (xu). Next subtract the upper breadth from the lower breadth; the remainder is the site (zhi). Multiply the south path height by the site as dividend and divide it by the terrace height; subtract the number obtained from the lower breadth; the remainder is the length of the south path. Add the south length to the lower breadth, multiply by the south path height, and add the requisite. Again multiply the north path height by the site as dividend and divide it by the terrace height; subtract the number obtained from the lower breadth; the remainder is the length of the north path. Add the north path length to the lower breadth, multiply by the north path height, and again add the requisite; multiply all together by half the path width; add the number obtained to the provisional: this is the joint rate (gonglü). Multiply the thickness of the stone slab (shiban) by the number of layers of the foundation footing (jijiao): this is the foundation height (jigao). Next double the path width, add the lower width, and multiply by the lower length: this is the foundation rate (jilü). Next multiply the lower breadth by the lower length (mao), subtract it from the foundation rate, and multiply the remainder by the foundation height: this is the stone rate (shilü). Subtract the stone rate from the joint rate; the remainder is the brick rate (zhuanlü). Divide the brick rate by the brick volume divisor, obtaining the number of bricks; divide the stone rate by the stone-step volume divisor, obtaining the number of stone slabs.\end{uniten}

\begin{unitzh}\unitlabel{shushu-045.15}草曰：倍上袤七丈，得一十四；加下袤一十七丈，得三十四丈；加上袤七丈，得四十一；乗下廣一十五丈，得六百一十五；併寄，得七百七十；乗高一十二丈，得九千二百四十丈，為土率。以六除率，一千五百四十丈。以千尺通之{\zhsmall 按：舊本誤為十八通之，今改正。}，為一百五十四萬尺，堅積。以築工九十尺除之，得一萬七千一百一十一工九分工之一，為築工。次以穿率四因堅積，得六百一十六萬尺，為實；以堅率三因钁鍬工二百尺，得六百尺，為法；除實，得一萬二百六十六工三分工之二，為穿工。次以壤率五因堅積，得七百七十萬尺，為壤積，以三為母。具圖如後。\end{unitzh}
\begin{uniten}Working: Double the upper length, 7 zhang, obtaining 14; add the lower length, 17 zhang, obtaining 34 zhang [Translator's note: steps are evidently missing here; 14 plus 17 does not give 34]; add the upper length, 7 zhang, obtaining 41; multiply by the lower breadth, 15 zhang, obtaining 615; add the reserve, obtaining 770; multiply by the height, 12 zhang, obtaining 9240 zhang: this is the earth rate. Divide the rate by 6: 1540 zhang. Convert it by 1000 chi [Siku editors: The old edition erroneously reads "convert it by 18"; now corrected], making 1,540,000 chi, the firm volume. Divide it by the 90 chi of a ramming workday, obtaining 17,111 1/9 workdays: this is the ramming labor. Next multiply the firm volume by the dug rate 4, obtaining 6,160,000 chi as dividend; multiply the 200 chi of a mattock-and-spade workday by the firm rate 3, obtaining 600 chi as divisor; divide the dividend, obtaining 10,266 2/3 workdays: this is the digging labor (chuangong). Next multiply the firm volume by the loose rate 5, obtaining 7,700,000 chi as the loose volume, with 3 as denominator. The diagram is set out below.\end{uniten}

\begin{unitzh}\unitlabel{shushu-045.16}乃以右行母互乗左行子：左上得一十二，副位得九，次得六，下得一十八。乃變右行名為寄子。以諸母相乗，得一十八，為寄母。具圖如後。\end{unitzh}
\begin{uniten}Then with the denominators of the right column cross-multiply the numerators of the left column: the left top gives 12, the second position (fuwei) gives 9, the next gives 6, the bottom gives 18. Then change the name of the right column to reserve numerators (jizi). Multiply the denominators together, obtaining 18, as the reserve denominator. The diagram is set out below.\end{uniten}

\begin{unitzh}\unitlabel{shushu-045.17}次以棚道全分及所當鮮母衍子：三當五，及三當七，并二當五。\end{unitzh}
\begin{uniten}Next take the whole of the scaffold path and the scant denominators and extended numerators of what it equals: 3 equals 5, and 3 equals 7, together with 2 equals 5.\end{uniten}

\begin{unitzh}\unitlabel{shushu-045.18}乃以右行鮮母互乗左行衍子：上得四十五，副得四十二，次得三十，下得一十八，為右行。以鮮母相乗，得一十八。乃對寄左圖列之。\end{unitzh}
\begin{uniten}Then with the scant denominators of the right column cross-multiply the extended numerators of the left column: the top gives 45, the second gives 42, the next gives 30, the bottom gives 18; this is the right column. Multiply the scant denominators together, obtaining 18. Then set them out in a diagram opposite the reserve on the left.\end{uniten}

\begin{unitzh}\unitlabel{shushu-045.19}乃以左右兩行母子對乗之：上得五百四十，副得三百七十八，次得一百八十，下得三百二十四，母得三百二十四。\end{unitzh}
\begin{uniten}Then multiply correspondingly the denominators and numerators of the two columns, left and right: the top gives 540, the second gives 378, the next gives 180, the bottom gives 324, and the denominator gives 324.\end{uniten}

\begin{unitzh}\unitlabel{shushu-045.20}今以平分術入之，併四子，得一千四百二十二，為總子；以列位四乗乗母三百二十四，得一千二百九十六，為總母。\end{unitzh}
\begin{uniten}Now enter it by the method of equal division (pingfen shu): add the 4 numerators, obtaining 1422, as the total numerator; with the listed positions, 4, multiply the multiplied denominator 324, obtaining 1296, as the total denominator.\end{uniten}

\begin{unitzh}\unitlabel{shushu-046.1}乃以總母、子求等，得一十八，俱約之：總子得七十九，為定子；總母得七十二，為定母。以定子七十九乗棚道四十步，得三千一百六十，於次。以棚道一四十步減徃來一百六十步，餘一百二十，為平道，以乗定七十二，得八千六百四十，加次，共得一萬一千八百。又以踟蹰十加一，放身下加一，得一萬二千九百八十，仍於次。又以載二十歩乗定母七十二，得一千四百四十，併次，得一萬四千四百二十，為統數。以定母七十二除統數，得二百步，不盡約為一十八分步之五，為定一返步。乃復通分内子，得三千六百五，為到法。\end{unitzh}
\begin{uniten}Then find the common measure (qiudeng) of the total denominator and numerator, obtaining 18, and reduce both by it: the total numerator gives 79, which is the definite numerator (dingzi); the total denominator gives 72, which is the definite denominator (dingmu). Multiply the scaffold road (pengdao), 40 bu (步), by the definite numerator 79, obtaining 3160, [placed] in the next position (ci). Subtract the scaffold road 1 40 bu [Translator's note: the 一 here appears superfluous.] from the going and coming (wanglai), 160 bu; the remainder, 120, is the level road (pingdao); use it to multiply the definite 72, obtaining 8640; add the next position, altogether obtaining 11,800. Further, for lingering (chichu), add 1 per 10: placing below itself (shen), add 1, obtaining 12,980, again [placed] in the next position. Further, multiply the definite denominator 72 by the loading (zai), 20 bu, obtaining 1440; combine with the next position, obtaining 14,420, which is the comprehensive number (tongshu). Divide the comprehensive number by the definite denominator 72, obtaining 200 bu; the remainder, reduced, makes 5/18 bu; this is the definite bu of one return trip (ding yi fan bu). Then again convert to fractional parts and add in the numerator (tongfen nazi), obtaining 3605, which is the arrival divisor (daofa).\end{uniten}

\begin{unitzh}\unitlabel{shushu-046.2}乃置程里六十，以三百六十通之，得二萬一千六百步，乗擔上一尺三寸，得二萬八千八十尺，以到母一十八乗之，得五十萬五千四百四十尺，為到實。實如法除，得一百四十尺，不盡七百四十，與法求等，得五，約之，為七百二十一分尺之一百四十八，為到土。復通分内子，得一十萬一千八十八，又以壤母三因，得三十萬二千二百六十四尺，為壤法。次以到土母七百二十一乗壤積七百七十萬尺，得五十五億五千一百七十萬尺，為壤實。實如法而一，得一萬八千三百六工，不盡一十四萬九千二百一十六，與法求等，得三十二，俱約之，為一萬八千三百六工九千四百七十七分工之四千六百六十三，為擔土工。\end{unitzh}
\begin{uniten}Then set up the stage distance (chengli), 60 li (里); convert it with 360, obtaining 21,600 bu; multiply by the load on the carrying pole (danshang), 1 chi (尺) 3 cun (寸), obtaining 28,080 chi; multiply it by the arrival denominator (daomu) 18, obtaining 505,440 chi, which is the arrival dividend (daoshi). Divide the dividend (shi) by the divisor (fa), obtaining 140 chi, with a remainder of 740; find the common measure with the divisor, obtaining 5; reduce by it, making 148/721 chi; this is the arrived earth (daotu). Again convert to fractional parts and add in the numerator, obtaining 101,088; further multiply by the loose-earth denominator (rangmu), 3, obtaining 302,264 chi, which is the loose-earth divisor (rangfa). Next multiply the loose-earth volume (rangji), 7,700,000 chi, by the denominator of the arrived earth, 721, obtaining 5,551,700,000 chi, which is the loose-earth dividend (rangshi). Divide the dividend by the divisor, obtaining 18,306 gong (工, units of labor), with a remainder of 149,216; find the common measure with the divisor, obtaining 32; reduce both by it, making 18,306 4663/9477 gong; this is the labor for carrying earth (dantu gong).\end{uniten}

\begin{unitzh}\unitlabel{shushu-046.3}次列前土工圖：築一萬七千一百一十一工九分工之一，及穿工一萬二百六十六工三分工之二，各通分内子：築率得一十五萬四千，穿率得三萬八百，擔率一億七千三百四十九萬六百二十五。按術當以通率圖諸母互乗諸率，令騐擔母九千四百七十七可用築母九約，亦可用穿母三約，故從省：以築母九約擔母九千四百七十七，得一千五十三，為築率乗數；又以穿母三約擔母九千四百七十七，得三千一百五十九，為穿率乗數。各以乗數乗本率，名曰就母圖。乃以九千四百七十七變名曰就母。先以築率一十五萬四千乗乗數一千五十三，得一億六千二百一十六萬二千，為築分；次以穿率三萬八百乗乗率三千一百五十九，得九千七百二十九萬七千二百，為穿分；就以擔率一億七千三百四十九萬六百二十五為擔分。併三分，共得四億三千二百九十四萬九千八百二十五，為總工分實。以就母九千四百七十七除之，得四萬九千六百八十五工九千四百七十七分工之三千八十，為總用功。具圖如後。\end{unitzh}
\begin{uniten}Next lay out the diagram of the foregoing earth labors: pounding (zhu), 17,111 1/9 gong, and digging (chuan) labor, 10,266 2/3 gong; convert each to fractional parts and add in the numerator: the pounding rate (lü) gives 154,000, the digging rate gives 30,800, and the carrying (dan) rate is 173,490,625. According to the method, one ought to cross-multiply all the rates by all the denominators of the diagram of converted rates (tonglü tu); on verifying, the carrying denominator 9477 can be reduced by the pounding denominator 9, and can also be reduced by the digging denominator 3; therefore we follow the shorter way: reduce the carrying denominator 9477 by the pounding denominator 9, obtaining 1053, which is the multiplier (chengshu) of the pounding rate; further reduce the carrying denominator 9477 by the digging denominator 3, obtaining 3159, which is the multiplier of the digging rate. Multiply each original rate by its multiplier; this is called the diagram of the adopted denominator (jiumu tu). Then 9477 changes its name and is called the adopted denominator (jiumu). First multiply the multiplier 1053 by the pounding rate 154,000, obtaining 162,162,000, which is the pounding parts (zhufen); next multiply the multiplying rate 3159 by the digging rate 30,800, obtaining 97,297,200, which is the digging parts; directly take the carrying rate 173,490,625 as the carrying parts. Combine the three parts, altogether obtaining 432,949,825, which is the dividend of the total labor parts. Divide it by the adopted denominator 9477, obtaining 49,685 3080/9477 gong, which is the total work used. The diagram is set out below.\end{uniten}

\begin{unitzh}\unitlabel{shushu-046.4}置各縣日程約税力，得力率，副併為科法工。置甲縣力一十三萬三千八百六十六，以一日程為之，只得此，為甲率；又置乙縣力二十三萬七千九百八十四，以二日約，得一十一萬八千九百九十二，為乙率；又置丙縣力三十一萬二千三百五十四，以三日約，得一十萬四千一百一十八，為内率。\end{unitzh}
\begin{uniten}Set up the day schedules (richeng) of the counties and divide the tax strength (shuili) by them, obtaining the strength rates (lilü); add them in an auxiliary position to make the allotment divisor (kefa) of the labor. Set up the strength of county A (jia), 133,866; with its schedule of 1 day it gives just this, which is the rate of A. Further set up the strength of county B (yi), 237,984; divide by 2 days, obtaining 118,992, which is the rate of B. Further set up the strength of county C (bing), 312,354; divide by 3 days, obtaining 104,118, which is the rate of nei [Translator's note: 内 is evidently an error for 丙, C.].\end{uniten}

\begin{unitzh}\unitlabel{shushu-046.5}列三，求等，得一萬四千八百七十四，俱約之：甲得九，乙得八，丙得七，各為定率；副併，得二十四。具圖如後。\end{unitzh}
\begin{uniten}List the three and find the common measure, obtaining 14,874; reduce all of them by it: A gives 9, B gives 8, C gives 7; each is a definite rate (dinglü); added in an auxiliary position they give 24. The diagram is set out below.\end{uniten}

\begin{unitzh}\unitlabel{shushu-046.6}乃以總分四億三千二百九十四萬九千八百二十五遍乗三縣定率，為各實；以就母九千四百七十七乗併率二十四，為科法。甲得三十八億九千六百五十四萬八千四百二十五，乙得三十四億六千三百五十九萬八千六百，丙得三十億三千六十四萬八千七百七十五，各為實；法得二十二萬七千四百四十八，為科法，除各實。具圖如後。\end{unitzh}
\begin{uniten}Then multiply the definite rates of the three counties one by one by the total parts, 432,949,825, to make the respective dividends; multiply the combined rate 24 by the adopted denominator 9477 to make the allotment divisor. A gives 3,896,548,425, B gives 3,463,598,600, C gives 3,030,648,775; each is a dividend; the divisor obtained, 227,448, is the allotment divisor, by which each dividend is divided. The diagram is set out below.\end{uniten}

\begin{unitzh}\unitlabel{shushu-046.7}乃以科法除各實：甲得一萬七千一百三十一工，不盡一十三萬六千七百三十七，為甲縣工；乙得一萬五千二百二十八，不盡二萬四百五十六，為乙縣工；丙得一萬三千三百二十四，不盡一十三萬一千六百二十三，為丙縣工。諸縣不盡皆輩為一工：甲合科一萬七千一百三十二工，乙合科一萬五千二百二十九工，丙合科一萬三千三百二十五工。\end{unitzh}
\begin{uniten}Then divide each dividend by the allotment divisor: A gives 17,131 gong, with a remainder of 136,737, which is the labor of county A; B gives 15,228, with a remainder of 20,456, which is the labor of county B; C gives 13,324, with a remainder of 131,623, which is the labor of county C. The remainders of all the counties are all counted (bei) as 1 gong: A should be allotted 17,132 gong, B should be allotted 15,229 gong, C should be allotted 13,325 gong.\end{uniten}

\begin{unitzh}\unitlabel{shushu-046.8}按：此以上所求擔土定一返數誤，其到土、工數及總用工數、三縣合科工數皆誤。盖題言築臺至少半、至中半、至太半，當自平地至三分之一，三分之一至二分之一，二分之一至三分之二，三分之二至臺頂，共四段，其分數自下而上逐層數，應取兩分數之較。草中所列諸分子，下設一數，上即列三分數，其段數既不確，又即用各分子全數，故求得棚道七十二僅當平道七十九，與題中所言二當五、三當七、當五之數顯然不合矣。今推歩改正於後。\end{unitzh}
\begin{uniten}Note: in the above, the number found for the definite one return trip of carrying earth is wrong; its arrived earth and number of gong, the number of total labor used, and the numbers of labor allotted to the three counties are all wrong. For the problem says that the terrace is built up to the lesser half (shaoban), up to the middle half (zhongban), and up to the greater half (taiban): this should be from the level ground to 1/3, from 1/3 to 1/2, from 1/2 to 2/3, and from 2/3 to the top of the terrace, altogether 4 sections; and their fractions, taken layer by layer from the bottom upward, should be the differences of two fractions. Of the fractions listed in the working, one number was set below and three fractional numbers were listed directly above; not only was the number of sections incorrect, but the full numbers of the fractions were also used directly; therefore it was found that 72 of scaffold road correspond to only 79 of level road, which obviously does not agree with the numbers stated in the problem, 2 corresponding to 5, 3 corresponding to 7, corresponding to 5 [Translator's note: a numeral seems to be missing before this last 當五.]. Now we work it out (tuibu) and correct it below.\end{uniten}

\begin{unitzh}\unitlabel{shushu-046.9}法：先求分子較。以三之一為第一段分數。以三之一與二之一分母對乗，得六；分母互乗分子，三之一得六之二，二之一得六之三，相減，餘六之一，得第二段分數。以二之一與三之二分母對乗，得六；分母互乗分子，二之一得三，三之二得四，相減，餘六之一，為第三段分數。以三分之二與全臺三分之三相減，餘三之一，為第四段分數。然後列四段母數於右，子數於左，以各分母連乗，又互乗分子，得三百二十四為總分母，第一段一百零八，第二段五十四，第三段五十四，第四段一百零八，為各分子。求等數，得五十四，徧約之：總母得六；分子第一段得二，第二段得一，第三段得一，第四段得二，為各段分母子。\end{unitzh}
\begin{uniten}Rule: first find the differences of the fractions. Take 1/3 as the fraction of the first section. Multiply together (duicheng) the denominators of 1/3 and 1/2, obtaining 6; cross-multiply (hucheng) the numerators by the denominators: 1/3 gives 2/6, 1/2 gives 3/6; subtract them from each other, and the remainder, 1/6, gives the fraction of the second section. Multiply together the denominators of 1/2 and 2/3, obtaining 6; cross-multiply the numerators by the denominators: 1/2 gives 3, 2/3 gives 4; subtract them from each other, and the remainder, 1/6, is the fraction of the third section. Subtract 2/3 and the whole terrace, 3/3, from each other; the remainder, 1/3, is the fraction of the fourth section. Then list the denominators of the 4 sections on the right and the numerators on the left; multiply all the denominators successively (liancheng), and also cross-multiply the numerators, obtaining 324 as the total denominator, and first section 108, second section 54, third section 54, fourth section 108 as the respective numerators. Find the common measure (dengshu), obtaining 54, and reduce all by it: the total denominator gives 6; of the numerators, the first section gives 2, the second section gives 1, the third section gives 1, the fourth section gives 2; these are the denominator and numerators of the sections.\end{uniten}

\begin{unitzh}\unitlabel{shushu-046.10}次列棚道一當平道一為下第一段，三當五為第二段，三當七為第三段，二當五為第四段。以棚道數為分母，以平道數為分子，分母連乗，又互乗分子，得一十八為總分母，第一段得一十八，第二段得三十，第三段得四十二，第四段得四十五，為各分子。求等數，得三，徧約之：總母得六；分子第一段得六，第二段得十，第三段得一十四，第四段得一十五。\end{unitzh}
\begin{uniten}Next list scaffold road 1 corresponding to level road 1 as the lowest, first section; 3 corresponding to 5 as the second section; 3 corresponding to 7 as the third section; 2 corresponding to 5 as the fourth section. Take the numbers of scaffold road as denominators and the numbers of level road as numerators; multiply the denominators successively, and also cross-multiply the numerators, obtaining 18 as the total denominator, and the first section gives 18, the second section gives 30, the third section gives 42, the fourth section gives 45, as the respective numerators. Find the common measure, obtaining 3, and reduce all by it: the total denominator gives 6; of the numerators, the first section gives 6, the second section gives 10, the third section gives 14, the fourth section gives 15.\end{uniten}

\begin{unitzh}\unitlabel{shushu-046.11}次列前後兩分母分子{\zhsmall 已約之數}各對乗，得總母三十六；分子第一段得一十二，第二段得一十，第三段得一十四，第四段得三十。併諸子，得六十六，為總子。又求等，得六，各約之，得定母六，為棚道數；定子一十一，為平道數。此即合全臺數總計，凡棚道六當平道一十一也。\end{unitzh}
\begin{uniten}Next list the earlier and later two [sets of] denominators and numerators [Note: the numbers already reduced], and multiply each pair together, obtaining the total denominator 36; of the numerators, the first section gives 12, the second section gives 10, the third section gives 14, the fourth section gives 30. Combine all the numerators, obtaining 66, which is the total numerator. Again find the common measure, obtaining 6, and reduce each by it, obtaining the definite denominator 6, which is the number of scaffold road, and the definite numerator 11, which is the number of level road. This is the combined total for the whole terrace: altogether, 6 of scaffold road correspond to 11 of level road.\end{uniten}

\begin{unitzh}\unitlabel{shushu-046.12}乃以定子一十一乗棚道四十步，得四百四十步，為次。以棚道四十步減一百六十步，餘一百二十步，為平道，以定母六乗之，得七百二十步，加次，共得一千一百六十步。以踟蹰加一於身，得一千二百七十六步。又以載輸二十步乗定母六，得一百二十步，义加次，共得一千三百九十六歩，為統數。以定母六除之，得二百三十二步又三分步之二，為定一返步。乃復通分内子，得六百九十八，為到法。\end{unitzh}
\begin{uniten}Then multiply the scaffold road, 40 bu, by the definite numerator 11, obtaining 440 bu, as the next position. Subtract the scaffold road, 40 bu, from 160 bu; the remainder, 120 bu, is the level road; multiply it by the definite denominator 6, obtaining 720 bu; add the next position, altogether obtaining 1160 bu. For lingering, add 1 to itself (shen), obtaining 1276 bu. Further, multiply the definite denominator 6 by the loading and delivering (zaishu), 20 bu, obtaining 120 bu; again add the next position, altogether obtaining 1396 bu, which is the comprehensive number. Divide it by the definite denominator 6, obtaining 232 bu and 2/3 bu, which is the definite bu of one return trip. Then again convert to fractional parts and add in the numerator, obtaining 698, which is the arrival divisor.\end{uniten}

\begin{unitzh}\unitlabel{shushu-046.13}乃置程里六十，以三百六十步通之，得二萬一千六百步，乗上一尺三寸，得二萬八千零八十尺，以到法分母三乗之，得八萬四千二百四十尺，為到實。以到法除之，得一百二十尺，不盡四百八十，約之為三百四十九分尺之二百四十，舊名到土復。到土復者，一工每日所擔之堅土數也。通分内子，得四萬二千一百二十，以壤母三因之，得一十二萬六千□百六十尺，為壤法。次以到土復母乗壤積七百七十萬尺，得二十六億八千七百三十萬尺，為壤實。實如法而一，得二萬一千二百六十七工，不盡一千八百八十，約之得三千一百五十九分工之四十七，為擔土工。\end{unitzh}
\begin{uniten}Then set up the stage distance, 60 li; convert it with 360 bu, obtaining 21,600 bu; multiply by the above 1 chi 3 cun, obtaining 28,080 chi; multiply it by the denominator of the arrival divisor, 3, obtaining 84,240 chi, which is the arrival dividend. Divide it by the arrival divisor, obtaining 120 chi, with a remainder of 480; reduced, this makes 240/349 chi; this was formerly named daotufu ("arrived earth again"). The daotufu is the amount of hard earth carried by 1 gong each day. Convert to fractional parts and add in the numerator, obtaining 42,120; multiply it by the loose-earth denominator 3, obtaining 126 thousand [lost] hundred and 60 chi, which is the loose-earth divisor. Next multiply the loose-earth volume, 7,700,000 chi, by the denominator of the daotufu, obtaining 2,687,300,000 chi, which is the loose-earth dividend. Divide the dividend by the divisor, obtaining 21,267 gong, with a remainder of 1880; reduced, this gives 47/3159 gong; this is the labor for carrying earth.\end{uniten}

\begin{unitzh}\unitlabel{shushu-046.14}次列築工、穿工、擔工三數及分母分子，互乗以齊之。\end{unitzh}
\begin{uniten}Next list the three numbers of pounding labor, digging labor and carrying labor, together with their denominators and numerators, and cross-multiply to make them uniform (qi).\end{uniten}

\begin{unitzh}\unitlabel{shushu-047.1}列三工數，各通分内子，得築率一十五萬四千，穿率三萬零八百，擔率六千七百一十八萬二千五百。驗擔母數，築、穿二母數皆可度盡，用就母法：以築母除擔母，得三百五十一，為築東率；以穿母除擔母，得一千零五十三，為穿乗率。各以乗率乗本率，得築分五千四百零五萬四千，穿分三千二百四十三萬二千四百，擔率六千七百一十萬二千五百即為擔分。同以擔分母為總分母。乃併三分數，得一億五千三百六十六萬八千九百分，為實；以總分母為法除之，得四萬八千六百四十四工，不盡二千五百零四，約之不變，即命為三千一百五十九分二之二千五百零四，為總用工數。\end{unitzh}
\begin{uniten}Set out the three numbers of man-days (gong); for each, reduce to parts and add in the numerator (tongfen nazi), obtaining the ramming rate (lü) 154,000, the digging rate 30,800 and the carrying rate 67,182,500. Examining the carrying denominator, the two denominators of ramming and digging can both measure it exactly (dujin), so use the method of the adopted denominator (jiumu fa): divide the carrying denominator by the ramming denominator, obtaining 351, as the ramming 'east' rate [Translator's note: 東 'east' is evidently a miscopy of 乗 'multiplying'.]; divide the carrying denominator by the digging denominator, obtaining 1053, as the digging multiplying rate (chenglü). Multiply each one's own rate by its multiplying rate, obtaining the ramming parts 54,054,000 and the digging parts 32,432,400; the carrying rate, 67,102,500, is then the carrying parts. Take the carrying denominator in common as the total denominator. Then add the three numbers of parts, obtaining 153,668,900 parts, as the dividend (shi); take the total denominator as the divisor (fa) and divide by it, obtaining 48,644 man-days, with 2504 left over; reduced, it does not change, so it is named 2504/3159 of 二 [Translator's note: 二 'two' is evidently a miscopy of 工 'man-day'.], as the total number of man-days used.\end{uniten}

\begin{unitzh}\unitlabel{shushu-047.2}次置三縣力役，以里數遠近通之。先以乙、丙二里數求等，得六十，各約之，乙得二，丙得三。以乙二徧乗甲、丙，以丙三徧乗甲、乙，得甲率八十萬零三千一百六十九，乙率七十一萬三千九百五十二，丙率六十二萬四千七百零八。三率求等，得八萬九千二百四十四，徧約之，甲得，乙得八，丙得七，各為率。併率，得二十四。乃以工總分一億五千三百六十六萬八千九百分乗三縣各率，得甲實一十三億八千三百零二萬零一百，乙實一十二億二千九百三十五萬一千二百，丙實一十億零七千五百六十八萬二千三百。以就母三千一百五十九乗併率二十四，得七萬五千八百一十六，為科法。除各實，不盡者各進一工，得甲縣合科一萬八千二百四十二工，乙縣合科一萬六千二百一十五工，丙縣合科一萬四千一百八十九工。\end{unitzh}
\begin{uniten}Next set out the labor service (liyi) of the three counties and convert (tong) them by the li (里) numbers, far and near. First find the common measure (deng) of the two li-numbers of B (yi) and C (bing), obtaining 60; reduce each by it: B gives 2 and C gives 3. Multiply A (jia) and C throughout by B's 2, and multiply A and B throughout by C's 3, obtaining A's rate 803,169, B's rate 713,952 and C's rate 624,708. Find the common measure of the three rates, obtaining 89,244; reduce them all by it: A gives [Translator's note: no number is given for A.], B gives 8, C gives 7; each is taken as a rate. Add the rates, obtaining 24. Then multiply each rate of the three counties by the total parts of man-days, 153,668,900 parts, obtaining A's dividend 1,383,020,100, B's dividend 1,229,351,200 and C's dividend 1,075,682,300. Multiply the sum of the rates, 24, by the adopted denominator 3159, obtaining 75,816, as the assessment divisor (kefa). Divide each dividend; where it does not divide exactly, in each case advance by 1 man-day, obtaining: county A is to be assessed 18,242 man-days, county B is to be assessed 16,215 man-days, county C is to be assessed 14,189 man-days.\end{uniten}

\begin{unitzh}\unitlabel{shushu-047.3}求甎者，倍轉道並濶六尺，得一丈二尺，遍加臺上下廣袤，變各為上下濶長；以甎厚六寸加臺高。\end{unitzh}
\begin{uniten}To find the bricks: double the turning ramps (zhuandao), all (bing) 6 chi (尺) wide, obtaining 1 zhang (丈) 2 chi; add it to all of the platform's upper and lower breadth (guang) and extent (mao), changing each into the upper and lower width (kuo) and length (chang); add the brick thickness of 6 cun (寸) to the platform height.\end{uniten}

\begin{unitzh}\unitlabel{shushu-047.4}先以甎長一尺二寸、濶六寸、厚二寸五分相乗之，得一百八十寸，為甎積法。次列石版長五尺、濶二尺、厚五寸，相乗，得五千寸，為石積法。具圖如後。\end{unitzh}
\begin{uniten}First multiply together the brick's length 1 chi 2 cun, width 6 cun and thickness 2 cun 5 fen (分), obtaining 180 cun, as the brick-volume divisor (zhuanji fa). Next set out the stone slab's length 5 chi, width 2 chi and thickness 5 cun; multiply them together, obtaining 5000 cun, as the stone-volume divisor (shiji fa). The diagram is set out as follows.\end{uniten}

\begin{unitzh}\unitlabel{shushu-047.5}騐得諸法皆寸。乃以變各圖上下長、濶、直，按術求率：倍上長八百二十寸，得一千六百四十寸，加下長一千八百二十寸，得三千四百六十，乗上濶六百二十，得二百一十四萬五千二百，為寄。次倍下長一千八百二十，得三千六百四十，加上長八百二十，得四千四百六十，乗下濶一千六百二十，得七百二十二萬五千二百，併寄，得九百三十七萬四百，乗臺直一千二百六寸，得一百一十三億七十萬二千四百寸，仍為寄。乃騐土積圖，土率九千二百四十丈，以一百萬寸通之，得九十二億四千萬寸，以減寄，餘二十億六千七十萬二千四百寸，如六而一，得三億四千三百四十五萬四百寸，為泛。\end{unitzh}
\begin{uniten}Checking, the divisors obtained are all in cun. Then, with the upper and lower lengths and widths and the upright (zhi, 直) of the diagram of the changed [measures], seek the rate according to the method: double the upper length 820 cun, obtaining 1640 cun; add the lower length 1820 cun, obtaining 3460; multiply by the upper width 620, obtaining 2,145,200, as the deposit (ji). Next double the lower length 1820, obtaining 3640; add the upper length 820, obtaining 4460; multiply by the lower width 1620, obtaining 7,225,200; add the deposit, obtaining 9,370,400; multiply by the platform's upright 1206 cun, obtaining 11,300,702,400 cun, which again is taken as the deposit. Then check the earth-volume diagram: the earth rate is 9240 zhang; convert it by 1,000,000 cun, obtaining 9,240,000,000 cun; subtract it from the deposit, leaving 2,060,702,400 cun; divide by 6, obtaining 343,450,400 cun, as the provisional amount (fan).\end{uniten}

\begin{unitzh}\unitlabel{shushu-047.6}次置南道高五分之二，北道高强半，係四分之三，及臺原高一千二百寸。具圖如後。\end{unitzh}
\begin{uniten}Next set out the south ramp's height, 2/5; the north ramp's height, 'a strong half' (qiangban), which is 3/4; and the platform's original height, 1200 cun. The diagram is set out as follows.\end{uniten}

\begin{unitzh}\unitlabel{shushu-047.7}乃以南道高子二乗原臺高，得二千四百寸，為南寔；以南道高子三乗原臺高，得三千六百寸，為北寔。各如本母而一，得四百八十寸，約為四丈八尺，為南道兩隅，又為東外道巽隅高，又為西道坤隅高；所得九百寸，約為九丈，為北道兩隅高，又為西道乾隅高，又為東裏道艮隅高。具圖如後。\end{unitzh}
\begin{uniten}Then multiply the original platform height by the numerator (zi) 2 of the south ramp height, obtaining 2400 cun, as the south dividend; multiply the original platform height by the numerator 3 of the south [Translator's note: 南 'south' is evidently an error for 北 'north'.] ramp height, obtaining 3600 cun, as the north dividend. Divide each by its own denominator, obtaining 480 cun, which reduces to 4 zhang 8 chi, as the two corners of the south ramp, also as the height of the xun (巽) corner of the east outer ramp [Translator's note: the trigram names denote directions: xun southeast, kun southwest, qian northwest, gen northeast.], and also as the height of the kun (坤) corner of the west ramp; the 900 cun obtained, which reduces to 9 zhang, is the height of the two corners of the north ramp, also the height of the qian (乾) corner of the west ramp, and also the height of the gen (艮) corner of the east inner ramp. The diagram is set out as follows.\end{uniten}

\begin{unitzh}\unitlabel{shushu-047.8}以北道高九丈減臺高一十二丈，餘三丈，為上停高；以南道高四丈八尺減北道高九丈，餘四丈二尺，為中停髙；命南道高四丈八尺為下停高。\end{unitzh}
\begin{uniten}Subtract the north ramp height 9 zhang from the platform height 12 zhang, leaving 3 zhang, as the height of the upper stage (ting); subtract the south ramp height 4 zhang 8 chi from the north ramp height 9 zhang, leaving 4 zhang 2 chi, as the height of the middle stage; designate the south ramp height 4 zhang 8 chi as the height of the lower stage.\end{uniten}

\begin{unitzh}\unitlabel{shushu-047.9}三停高皆如履級寸而一，得五十，為東裏道級數；得七十，為西道級數；得八十，為東外道級數。次以上長八百二十寸減下長一千八百二十寸，餘一千寸，以半之，得五百寸，為勾。以乗南道高四百八十寸，得二十四萬寸，為實；如臺高一千二百寸而一，半之{\zhsmall 按上下長較已半得數，又半之，誤。}，得一百寸，為底率。以率減下長一千八百二十，餘一千七百二十寸，為底股。以外道級八十約之，得二尺一寸五分，為外道踏蹤。\end{unitzh}
\begin{uniten}Divide each of the three stage heights by the cun of a step (lüji), obtaining 50, as the number of steps of the east inner ramp; obtaining 70, as the number of steps of the west ramp; obtaining 80, as the number of steps of the east outer ramp. Next subtract the upper length 820 cun from the lower length 1820 cun, leaving 1000 cun; halve it, obtaining 500 cun, as the base (gou). Multiply the south ramp height 480 cun by it, obtaining 240,000 cun, as the dividend; divide by the platform height 1200 cun, and halve it [Siku editors: the difference of the upper and lower lengths was already halved to obtain the number; halving it again is an error.], obtaining 100 cun, as the bottom rate (dilü). Subtract the rate from the lower length 1820, leaving 1720 cun, as the bottom altitude (digu). Reduce it by the 80 steps of the outer ramp, obtaining 2 chi 1 cun 5 fen, as the tread (tazong) of the outer ramp.\end{uniten}

\begin{unitzh}\unitlabel{shushu-047.10}又以底率一百寸減底股一千七百二十寸，餘一千六百二十寸，為中股。乃以勾五百寸乗中停高四丈二尺，得二十一萬寸，為實；如臺高一千二百寸而一，得一百七十五寸，為中率。以中率減中股一千六百二十寸，餘一千四百四十五寸，為上股。以西道級七十除上股一千四百四十五寸，得二尺一寸四分寸之九，為西道踏蹤。又以勾五百寸乗上停高三百寸，得一十五萬寸，為實；如臺高一千二百寸而一，得一百二十五寸，為上率。併中率一百七十五寸，得三百，減上股一千四百四十五寸，餘一千一百四十五寸，為實；如裏道級五十而一，得二尺二寸九分，為裏道踏蹤。各得，具圖以見如後。\end{unitzh}
\begin{uniten}Again subtract the bottom rate 100 cun from the bottom altitude 1720 cun, leaving 1620 cun, as the middle altitude. Then multiply the middle stage height 4 zhang 2 chi by the base 500 cun, obtaining 210,000 cun, as the dividend; divide by the platform height 1200 cun, obtaining 175 cun, as the middle rate. Subtract the middle rate from the middle altitude 1620 cun, leaving 1445 cun, as the upper altitude. Divide the upper altitude 1445 cun by the 70 steps of the west ramp, obtaining 2 chi 1 cun 9/4 cun, as the tread of the west ramp. Again multiply the upper stage height 300 cun by the base 500 cun, obtaining 150,000 cun, as the dividend; divide by the platform height 1200 cun, obtaining 125 cun, as the upper rate. Add the middle rate 175 cun, obtaining 300; subtract it from the upper altitude 1445 cun, leaving 1145 cun, as the dividend; divide by the 50 steps of the inner ramp, obtaining 2 chi 2 cun 9 fen, as the tread of the inner ramp. Each having been obtained, the diagram is set out to show them, as follows.\end{uniten}

\begin{unitzh}\unitlabel{shushu-047.11}次以道濶六尺併下長一千八百二十寸，得一千八百八十寸，為補。以南北道高併臺高，共得二千五百八十寸，乗補，得四百八十五萬四百寸，為需。{\zhsmall 按此二數不合，其法太疎故也。}次以上廣五百寸減下廣一千五百寸，餘一千寸，為址，乗南道高四百八十寸，得四十八萬寸，為實；以臺高一千二百寸除之，得四百寸，為減率。以減下廣一千五百寸，餘一千一百寸，為南道長。併下廣一千五百寸，得二千六百，乗南道高四百八十寸，得一百二十四萬八千寸，加需。又以址一千寸乗北道高九百寸，得九十萬寸，為實；亦如臺高一千二百而一，得七百五十。以減下廣一千五百寸，餘七百五十寸，為北道長。併下廣一千五百，得二千二百五十寸，乗北道九百寸，得二百二萬五千寸，又加需，共得八百一十二萬三千四百。以半道濶三尺乗之，得二億四千三百七十萬二千寸，併泛三億四千三百四十五萬四百，得五億八千七百一十五萬二千四百寸，為共率。\end{unitzh}
\begin{uniten}Next add the ramp width 6 chi to the lower length 1820 cun, obtaining 1880 cun, as the supplement (bu). Add the heights of the south and north ramps to the platform height, obtaining in all 2580 cun; multiply by the supplement, obtaining 4,850,400 cun, as the requirement (xu). [Siku editors: these two numbers do not fit; it is because the method is too crude.] Next subtract the upper breadth 500 cun from the lower breadth 1500 cun, leaving 1000 cun, as the site (zhi, 址); multiply by the south ramp height 480 cun, obtaining 480,000 cun, as the dividend; divide it by the platform height 1200 cun, obtaining 400 cun, as the subtraction rate (jianlü). Subtract it from the lower breadth 1500 cun, leaving 1100 cun, as the south ramp length. Add the lower breadth 1500 cun, obtaining 2600; multiply by the south ramp height 480 cun, obtaining 1,248,000 cun; add it to the requirement. Again multiply the north ramp height 900 cun by the site 1000 cun, obtaining 900,000 cun, as the dividend; likewise divide by the platform height 1200, obtaining 750. Subtract it from the lower breadth 1500 cun, leaving 750 cun, as the north ramp length. Add the lower breadth 1500, obtaining 2250 cun; multiply by the north ramp 900 cun, obtaining 2,025,000 cun; again add it to the requirement, obtaining in all 8,123,400. Multiply it by half the ramp width, 3 chi, obtaining 243,702,000 cun; add the provisional amount 343,450,400, obtaining 587,152,400 cun, as the combined rate (gonglü).\end{uniten}

\begin{unitzh}\unitlabel{shushu-047.12}次以基脚七層乗石版厚五寸，得三十五寸，為基高。次倍道濶，得一百二十，併下濶一千六百二十寸，得一千七百四十，乗下長一千八百二十，得二百一十六萬六千八百，為基率。次以下廣一千五百寸乗下袤一千七百寸，得二百五十五萬，減基率，餘六十一萬六千八百寸，乗基高三十五寸，得二千一百五十八萬八千寸，為石率。以石率減共率，餘五億六千五百五十六萬四千四百寸，為甎率。以甎積法一百八十寸除之，得一{\zhsmall 按三訛一}百一十四萬二千二十四片九分片之四。乃以石積法五千寸除石率二千一百五十八萬八千寸，得四千三百一十七片，為石版；不盡三十寸，棄之不輩。合問。\end{unitzh}
\begin{uniten}Next multiply the stone-slab thickness 5 cun by the 7 courses of the foundation footing (jijiao), obtaining 35 cun, as the foundation height. Next double the ramp width, obtaining 120; add the lower width 1620 cun, obtaining 1740; multiply by the lower length 1820, obtaining 2,166,800, as the foundation rate. Next multiply the lower breadth 1500 cun by the lower extent 1700 cun, obtaining 2,550,000; subtract it from the foundation rate, leaving 616,800 cun; multiply by the foundation height 35 cun, obtaining 21,588,000 cun, as the stone rate. Subtract the stone rate from the combined rate, leaving 565,564,400 cun, as the brick rate. Divide it by the brick-volume divisor 180 cun, obtaining 1,142,024 4/9 pieces [Siku editors: 三 'three' has been miswritten as 一 'one'.]. Then divide the stone rate 21,588,000 cun by the stone-volume divisor 5000 cun, obtaining 4317 pieces, as the stone slabs; the remainder of 30 cun is discarded, 'not bei' (不輩) [Translator's note: 不輩 is unintelligible here and evidently corrupt.]. This matches the question.\end{uniten}

\begin{unitzh}\unitlabel{shushu-048.1}按：右草自求甎砌臺身積至級數，皆無誤；求踏蹤數有誤；至求補數、需數，其法更多未合。如題内云：先用甎砌臺身，次用甎砌轉道，周圍五帶，並濶六尺，南北二平道，東西三峻道。如自臺下艮隅南升至巽，西折行至坤，北升至乾，東折行至艮，復南升至巽，乃登臺頂。是南北二平道至臺址，皆如一面平堆形；西峻道至臺址、東裏峻道至臺址，皆上如一面三角斜堆，下如一面平堆形；東外峻道至臺址，如一面三角斜堆形。法當按各形取積。\end{unitzh}
\begin{uniten}Note: In the working above, from finding the volume of the brick-built terrace body (taishen) up to the numbers of steps, everything is without error. Finding the tread-lengths (tazong) contains errors. As for finding the supplement number (bu) and the requirement number (xu), even more of the method fails to agree. As the problem states: first build the terrace body with bricks; next build with bricks the winding roads (zhuandao), going round in 5 bands, all 6 chi (尺) wide, 2 level roads (pingdao) on the south and north and 3 steep roads (jundao) on the east and west. That is, from the gen corner at the foot of the terrace one ascends southward to xun, turns west and walks to kun, ascends northward to qian, turns east and walks to gen, again ascends southward to xun, and then mounts the top of the terrace. [Translator's note: gen, xun, kun and qian are trigram names for northeast, southeast, southwest and northwest.] Hence the 2 level roads on the south and north, down to the terrace site (taizhi), are each like the shape of a one-faced level pile (pingdui). The west steep road down to the terrace site and the east inner steep road down to the terrace site are each, above, like a one-faced triangular sloping pile (sanjiao xiedui) and, below, like the shape of a one-faced level pile. The east outer steep road down to the terrace site is like the shape of a one-faced triangular sloping pile. The method ought to take the volume according to each shape.\end{uniten}

\begin{unitzh}\unitlabel{shushu-048.2}草中取二平道，皆如法；至取峻道，乃以道濶併下長為補，以南北道高、臺高三數併乗之為需，然後以半道濶乗之為積，其意殊不可解；然以數考之，有差至三之一者。又臺身外四面甎砌各加六尺，轉道東面加二層，三面各加一層，是臺基長一十九丈四尺，濶一十八丈；草中以下濶為一千七百四十寸，差少六尺；下長仍用一千八百二十寸，差少一丈二尺：皆疎漏之誤。\end{unitzh}
\begin{uniten}In the working, the 2 level roads are taken entirely according to the method. But when it comes to taking the steep roads, it adds the road width to the lower length to make the supplement. It multiplies this by the sum of the 3 numbers, the heights of the south and north roads and the terrace height, to make the requirement, and afterwards multiplies by half the road width to make the volume. Its meaning is quite unintelligible; and, checked by the numbers, there are discrepancies reaching as much as 1/3. Moreover, outside the terrace body the brickwork on each of the 4 sides adds 6 chi. The winding roads add 2 layers on the east side and 1 layer on each of the other 3 sides, so the terrace base (taiji) is 19 zhang (丈) 4 chi long and 18 zhang wide. In the working, the lower width is taken as 1740 cun (寸), short by 6 chi, and the lower length still uses 1820 cun, short by 1 zhang 2 chi: all are errors of carelessness and omission.\end{uniten}

\begin{unitzh}\unitlabel{shushu-048.3}今自求踏蹤至末，皆另步之。法：以臺上長減下長，餘一千寸，半之，得五百寸，為勾率。以乗南道高四百八十寸，得二十四萬寸，為實；如臺高一千二百寸而一，得二百寸，為底率。以減下長，得一千六百二十寸，為底股。以外道級八十約之，得二尺零二分半，為東外道踏蹤。又以底率減底股，餘一千四百二十寸，為中底。以勾率五百乗中停高四百二十寸，得二十一萬寸，如臺高寸而一，得一百七十五寸，為中率。以減中底，餘一千二百四十五寸，為中股。以西道級七十約之，得一尺七寸七分又七分分之六，為西道踏蹤。以中率減中股，餘一千零七十寸，為上底。以勾率五百乗上停高三百寸，得一十五萬寸，如臺高而一，得一百二十五寸，為上率。以減上底，餘九百四十五寸，為上股。以裏道五十約之，得一尺八寸九分，為裏道踏蹤。次以道濶加下長，與下停高相乗，得九十萬二千四百寸，為倍東外峻道三角形立面。中底與中停高相乗，得五十九萬六千四百寸，為倍\end{unitzh}
\begin{uniten}Now, from finding the tread-lengths to the end, everything is worked out anew. The method: subtract the upper length of the terrace from the lower length; the remainder is 1000 cun. Halve it, getting 500 cun, as the base rate (goulü). Multiply the south road height, 480 cun, by it, getting 240000 cun, as the dividend (shi). Divide by the terrace height, 1200 cun, getting 200 cun, as the bottom rate (dilü). Subtract it from the lower length, getting 1620 cun, as the bottom altitude (digu). Divide this by the outer road's 80 steps, getting 2 chi and 2 1/2 fen (分), as the tread-length of the east outer road. Further, subtract the bottom rate from the bottom altitude; the remainder, 1420 cun, is the middle bottom (zhongdi). Multiply the base rate 500 by the middle-stage (zhongting) height, 420 cun, getting 210000 cun. Divide by the terrace height in cun, getting 175 cun, as the middle rate. Subtract it from the middle bottom; the remainder, 1245 cun, is the middle altitude. Divide this by the west road's 70 steps, getting 1 chi 7 cun 7 fen and 6/7 of a fen, as the tread-length of the west road. Subtract the middle rate from the middle altitude; the remainder, 1070 cun, is the upper bottom. Multiply the base rate 500 by the upper-stage height, 300 cun, getting 150000 cun. Divide by the terrace height, getting 125 cun, as the upper rate. Subtract it from the upper bottom; the remainder, 945 cun, is the upper altitude. Divide this by the inner road's 50, getting 1 chi 8 cun 9 fen, as the tread-length of the inner road. Next, add the road width to the lower length and multiply by the lower-stage height, getting 902400 cun, as double the triangular-shaped vertical face (limian) of the east outer steep road. Multiply the middle bottom and the middle-stage height together, getting 596400 cun, as double the [Translator's note: the sentence breaks off here; it evidently continues at the start of unit 5 ("upper triangular-shaped vertical face of the west steep road"). Units 4 and 5 appear to be transposed in this text.]\end{uniten}

\begin{unitzh}\unitlabel{shushu-048.4}上下外濶相減，餘一千寸，以北道高乗之，得九十萬，以臺高寸數除之，得七百五十寸，以減北下外濶，得九百九十寸，為北道上濶；與北下外濶相加，以北道乗之，得二百四十五萬七千寸，為倍北道至臺址平堆形立面。併倍二面數，得三百九十九萬三千寸。又併前五數之共數，得九百九十六萬九千寸。以半道濶三十寸乗之，得二億九千九百零七萬寸，為共道峻道積。又併前砌臺身積三億四千三百四十五萬零四百寸，得總積六億四千二百五十二萬零四百寸，為共率。次二因道濶，與下長相加，得一千九百四十寸，與南道下外濶一千八百寸相乗，得三百四十九萬二千寸。又以原下廣一千五百寸、下袤一千七百寸相乗，得二百五十五萬寸，以減外長濶相乗，餘九十四萬二千寸，以石基三十五寸乗之，得三千二百九十七萬寸，為\end{unitzh}
\begin{uniten}upper and lower outer widths are subtracted from each other, leaving 1000 cun. Multiply by the north road height, getting 900000. Divide by the number of cun of the terrace height, getting 750 cun. Subtract it from the north lower outer width, getting 990 cun, as the upper width of the north road. Add it to the north lower outer width and multiply by the north road, getting 2457000 cun, as double the level-pile-shaped vertical face of the north road down to the terrace site. Add together the numbers of the 2 doubled faces, getting 3993000 cun. Further, add the combined number of the previous 5 numbers, getting 9969000 cun. Multiply by half the road width, 30 cun, getting 299070000 cun, as the combined volume of the roads and steep roads. Further, add the previous volume of the brick-built terrace body, 343450400 cun, getting the total volume, 642520400 cun, as the combined rate (gonglü). Next, multiply the road width by 2 and add it to the lower length, getting 1940 cun. Multiply by the lower outer width of the south road, 1800 cun, getting 3492000 cun. Further, multiply together the original lower breadth (guang), 1500 cun, and lower extent (mao), 1700 cun, getting 2550000 cun. Subtract it from the product of the outer length and width, leaving 942000 cun. Multiply by the stone footing (shiji), 35 cun, getting 32970000 cun, as the [Translator's note: continued by "stone rate" at the start of unit 6.]\end{uniten}

\begin{unitzh}\unitlabel{shushu-048.5}西峻道上三角形立面。中底、下底併，與下停高相乗，得一百五十五萬五千二百寸，為倍西峻道下平堆形立面。上底、上停高相乗，得三十二萬一千寸，為倍東裏道上三角形立面。上底、下底併，與中停高、下停高併相乗，得二百六十萬零一千寸，為倍東裏道下平堆形立面。倂五數，共得五百九十七萬六千寸，比舊草所用需數大一百一十二萬五千六百寸。次求南北平道立面：先以上下濶俱加三道濶，得上外濶八百寸，下外濶一千八百寸。上下外濶相減，餘一千寸，以南道高乗之，得四十八萬寸，以臺高寸除之，得四百寸，為減率。以減下外濶，得一千四百寸，為南道上濶；與下外濶相併，以南道高乗之，得一百五十三萬六千寸，為倍南道至臺址平堆形立面。又以上下濶各加倍道濶，得北上外濶七百四十寸，北下外濶一千七百四十寸。北\end{unitzh}
\begin{uniten}upper triangular-shaped vertical face of the west steep road. Add the middle bottom and the lower bottom, and multiply by the lower-stage height, getting 1555200 cun, as double the lower level-pile-shaped vertical face of the west steep road. Multiply the upper bottom and the upper-stage height together, getting 321000 cun, as double the upper triangular-shaped vertical face of the east inner road. Add the upper bottom and the lower bottom, and multiply by the sum of the middle-stage height and the lower-stage height, getting 2601000 cun, as double the lower level-pile-shaped vertical face of the east inner road. Add the 5 numbers, getting in all 5976000 cun, which is larger than the requirement number used in the old working by 1125600 cun. Next, find the vertical faces of the south and north level roads. First add 3 road widths to both the upper and the lower widths, getting the upper outer width, 800 cun, and the lower outer width, 1800 cun. The upper and lower outer widths are subtracted from each other, leaving 1000 cun. Multiply by the south road height, getting 480000 cun. Divide by the terrace height in cun, getting 400 cun, as the subtraction rate (jianlü). Subtract it from the lower outer width, getting 1400 cun, as the upper width of the south road. Add it to the lower outer width and multiply by the south road height, getting 1536000 cun, as double the level-pile-shaped vertical face of the south road down to the terrace site. Further, add double the road width to each of the upper and lower widths, getting the north upper outer width, 740 cun, and the north lower outer width, 1740 cun. The north [Translator's note: the sentence evidently continues at the start of unit 4.]\end{uniten}

\begin{unitzh}\unitlabel{shushu-048.6}石率。以石率減共率，餘六億零九百五十五萬零四百寸，為甎率。以甎長、濶、厚連乗，得一百八十寸，為甎法；以石長、濶、厚連乗，得五千寸，為石法。以甎法除甎率，得三百三十八萬六千三百九十一片九分片之一，為共用甎數；以石法除石率，得六千五百九十四片，為共用石數。\end{unitzh}
\begin{uniten}stone rate (shilü). Subtract the stone rate from the combined rate; the remainder, 609550400 cun, is the brick rate (zhuanlü). Multiply successively the length, width and thickness of a brick, getting 180 cun, as the brick divisor (fa). Multiply successively the length, width and thickness of a stone, getting 5000 cun, as the stone divisor. Divide the brick rate by the brick divisor, getting 3386391 1/9 pieces, as the number of bricks used in all. Divide the stone rate by the stone divisor, getting 6594 pieces, as the number of stones used in all.\end{uniten}

\begin{unitzh}\unitlabel{shushu-048.7}堂隍程築\end{unitzh}
\begin{uniten}Ramming for a hall by work quota\end{uniten}

\begin{unitzh}\unitlabel{shushu-048.8}問：有營造地基，長二十一丈，濶一十七丈。先令七人築堅三丈，計功二日。今涓吉立木有日，欲限三日築了，每日合收杵手幾何？答曰：日收五百五十五工三分工之一。\end{unitzh}
\begin{uniten}Question: There is a construction foundation, 21 zhang long and 17 zhang wide. Previously 7 men were set to ram 3 zhang firm, and the work was reckoned at 2 days. Now an auspicious day has been chosen and the day for raising the timbers is fixed; it is desired to complete the ramming within a limit of 3 days. How many rammers (chushou) ought to be taken on each day? Answer: Each day take on 555 1/3 workers (gong).\end{uniten}

\begin{unitzh}\unitlabel{shushu-048.9}術曰：以長乗濶，又乗原日、原人，為實；以限日乗築丈數，為法；除之，得人夫。\end{unitzh}
\begin{uniten}Method: Multiply the length by the width, and again by the original days and the original men, as the dividend (shi). Multiply the limited days by the number of zhang rammed, as the divisor (fa). Divide, getting the laborers (renfu).\end{uniten}

\begin{unitzh}\unitlabel{shushu-048.10}草曰：以長二十一丈乗濶一十七丈，得三百五十七丈；又乗原二日，得七百一十四；又乗原七人，得四千九百九十八，為工實。以限三日乗原築三丈，得九，為法。除實，得五百五十五工，不盡三，與法俱三約之，為三分工之一。為日收五百五十五工三分工之一。合問。\end{unitzh}
\begin{uniten}Working: Multiply the length, 21 zhang, by the width, 17 zhang, getting 357 zhang. Again multiply by the original 2 days, getting 714. Again multiply by the original 7 men, getting 4998, as the labor dividend (gongshi). Multiply the limited 3 days by the original 3 zhang rammed, getting 9, as the divisor. Divide the dividend, getting 555 workers. The remainder 3, reduced together with the divisor by 3, makes 1/3 of a worker. This is 555 1/3 workers taken on each day. This answers the question.\end{uniten}

\begin{unitzh}\unitlabel{shushu-048.11}砌甎計積\end{unitzh}
\begin{uniten}Laying bricks and reckoning the volume\end{uniten}

\begin{unitzh}\unitlabel{shushu-049.1}問有交到六門甎一十五垜，每垜高五尺，濶八尺，長一丈；其甎毎片長八寸，濶四寸，厚一寸。欲砌地面使用：堂屋三間，各深三丈，共濶五丈二尺；書院六間，各深一丈五尺，各濶一丈二尺；後閤四間，各深一丈三尺，内二間濶一丈，次二間濶一丈五尺；亭子地面之長濶幾何？答曰：見有一十八萬七千五百片，今用一萬六千四百六片四分片之一，外餘一十七萬一千九十三片四分片之三。\end{unitzh}
\begin{uniten}Question: There have been delivered 15 stacks (duo 垜) of six-gate bricks (liumen zhuan 六門甎). Each stack is 5 chi (尺) high, 8 chi wide and 1 zhang (丈) long. Each piece (pian 片) of the bricks is 8 cun (寸) long, 4 cun wide and 1 cun thick. It is desired to use them to pave floor surfaces: a hall building (tangwu) of 3 bays (jian 間), each 3 zhang deep, together 5 zhang 2 chi wide; a study (shuyuan) of 6 bays, each 1 zhang 5 chi deep and each 1 zhang 2 chi wide; rear chambers (houge) of 4 bays, each 1 zhang 3 chi deep, of which 2 bays are 1 zhang wide and the next 2 bays are 1 zhang 5 chi wide. What are the length and width of the floor surfaces of the pavilions (tingzi)? [Translator's note: no data for the pavilions are given here. The working squares a pavilion site of 1 zhang 4 chi and multiplies by 10 sites, and the answer gives numbers of bricks, not dimensions.] Answer: On hand there are 187,500 pieces; now used, 16,406 1/4 pieces; left over besides, 171,093 3/4 pieces.\end{uniten}

\begin{unitzh}\unitlabel{shushu-049.2}術曰：以少廣求之。置各地面深濶相乗，以間數若所數乗之，共為實；甎長濶相乗為甎平法，除，得今用甎數。次以甎垜高、長、濶相乗為實，却以甎法乗厚，得數為甎積法，除之，得毎垜甎數。次以垜數乗之，得見有甎，以減今用甎，得餘甎。\end{unitzh}
\begin{uniten}Method: Seek it by means of small width (shaoguang). Set up the depth and width of each floor surface and multiply them together. Multiply by the number of bays or the number of sites (suo 所); together they make the dividend (shi). Multiply the length and width of a brick together to make the brick-plane divisor (zhuanping fa); divide, obtaining the number of bricks now used. Next multiply the height, length and width of a brick stack together to make the dividend. Then multiply the brick divisor by the thickness; the number obtained makes the brick-volume divisor (zhuanji fa). Divide by it, obtaining the number of bricks in each stack. Next multiply by the number of stacks, obtaining the bricks on hand; with this subtract the bricks now used, obtaining the remaining bricks.\end{uniten}

\begin{unitzh}\unitlabel{shushu-049.3}草曰：置堂濶五丈二尺，乗深三丈，得一十五萬六千寸於上。又置書院深一丈五尺，乗濶一丈二尺，得一萬八千寸，又以六間乗之，得一十萬八千寸，加上，共得二十六萬四千寸，併上。又置後閤濶一丈，併濶一丈五尺，得二丈五尺，又以各二間乗之，得五百寸，以乗各深一丈三尺，得六萬五千寸，加上，得三十二萬九千寸，併上。次置亭基一丈四尺，自乗，得一萬九千六百寸，以一十所乗之，得一十九萬六千寸，又併上，共得五十二萬五千寸，為實。以甎長八寸乗濶四寸，得三十二寸，為甎平法，除之，得一萬六千四百六片四分片之一，為共用甎。次置毎垜高五尺，乗濶八尺，得四千寸，又乗長一丈，得四十萬寸，為毎垜實；却以甎平法三十二寸乗厚一寸，只得三十二寸，為甎積法，除之，得一萬二千五百片，又以十五垜乗之，得一十八萬七千五百片，為見有甎；内減今用甎，餘有一十七萬一千九十三片四分片之三，為外餘瓦數。合問。\end{unitzh}
\begin{uniten}Working: Set up the hall's width, 5 zhang 2 chi, and multiply by the depth, 3 zhang, obtaining 156,000 cun, at the top. Further set up the study's depth, 1 zhang 5 chi, and multiply by the width, 1 zhang 2 chi, obtaining 18,000 cun. Further multiply by 6 bays, obtaining 108,000 cun. Add to the top, obtaining altogether 264,000 cun, combined at the top. Further set up the rear chambers' width, 1 zhang, and combine it with the width 1 zhang 5 chi, obtaining 2 zhang 5 chi. Further multiply by 2 bays each, obtaining 500 cun. With this multiply the depth, 1 zhang 3 chi each, obtaining 65,000 cun. Add to the top, obtaining 329,000 cun, combined at the top. Next set up the pavilion site, 1 zhang 4 chi, and multiply it by itself, obtaining 19,600 cun. Multiply by 10 sites, obtaining 196,000 cun. Further combine at the top, obtaining altogether 525,000 cun as the dividend. Multiply the brick's length, 8 cun, by its width, 4 cun, obtaining 32 cun as the brick-plane divisor. Divide by it, obtaining 16,406 1/4 pieces as the bricks used in all. Next set up the height of each stack, 5 chi, and multiply by the width, 8 chi, obtaining 4,000 cun. Further multiply by the length, 1 zhang, obtaining 400,000 cun as the dividend for each stack. Then multiply the brick-plane divisor, 32 cun, by the thickness, 1 cun, obtaining just 32 cun as the brick-volume divisor. Divide by it, obtaining 12,500 pieces. Further multiply by 15 stacks, obtaining 187,500 pieces as the bricks on hand. From this subtract the bricks now used, leaving 171,093 3/4 pieces as the number of tiles (wa 瓦) left over besides. This answers the question.\end{uniten}

\begin{unitzh}\unitlabel{shushu-049.4}《數學九章》卷七上\end{unitzh}
\begin{uniten}Shuxue jiuzhang, juan 7, first part.\end{uniten}

\begin{unitzh}\unitlabel{shushu-049.5}子部，天文算法類，算書之屬，《數學九章》\end{unitzh}
\begin{uniten}Masters section (zibu), category of astronomy and mathematical methods (tianwen suanfa lei), class of mathematical books (suanshu zhi shu), Shuxue jiuzhang.\end{uniten}

\section*{卷7下}\addcontentsline{toc}{section}{卷7下}

\begin{unitzh}\unitlabel{shushu-050.1}《欽定四庫全書》 《數學九章》卷七下 宋·秦九韶撰 營廷\end{unitzh}
\begin{uniten}Qinding Siku quanshu (Imperially Approved Complete Library of the Four Treasuries). Shuxue jiuzhang (Mathematics in Nine Sections), juan 7, second part. Composed by Qin Jiushao of the Song. Ying ting [Translator's note: evidently for 營建, "Construction".]\end{uniten}

\begin{unitzh}\unitlabel{shushu-050.2}計定城築\end{unitzh}
\begin{uniten}Calculating and Settling the Building of a City Wall\end{uniten}

\begin{unitzh}\unitlabel{shushu-050.3}問：郡築一城，圍長一千五百一十丈；外築羊馬墻、𨳩壕，長與城同。城身髙三丈，靣闊三丈，下闊七丈五尺。羊馬墻髙一丈，靣闊五尺，下闊一丈。𨳩壕靣闊三十丈，下闊二十五丈。女頭鵲䑓共髙五尺五寸，共闊三尺六寸，共長一丈。鵲臺長一丈，髙五寸，闊五尺四寸。座子長一丈，髙二尺二寸五分，闊三尺六寸。肩子髙一尺二寸五分，闊三尺六寸，長八尺四寸。㡌子髙一尺五寸，闊三尺六寸，長六尺六寸。箭窻三眼，各闊六寸，長七寸五分，外眼比内眼斜低三寸。\end{unitzh}
\begin{uniten}Question: A prefecture (jun) builds a city wall whose circumference is 1510 zhang (丈) long; outside it, it builds a sheep-and-horse wall (yangma qiang) and digs a moat, their length being the same as the city wall's. The body of the city wall is 3 zhang high, 3 zhang wide at the top face (mian), and 7 zhang 5 chi (尺) wide at the bottom. The sheep-and-horse wall is 1 zhang high, 5 chi wide at the top face, and 1 zhang wide at the bottom. The dug moat is 30 zhang wide at the top face and 25 zhang wide at the bottom. The battlement (nütou) and the magpie platform (quetai) are together 5 chi 5 cun (寸) high, together 3 chi 6 cun wide, and together 1 zhang long. The magpie platform is 1 zhang long, 5 cun high, and 5 chi 4 cun wide. The seat (zuozi) is 1 zhang long, 2 chi 2 cun 5 fen (分) high, and 3 chi 6 cun wide. The shoulder (jianzi) is 1 chi 2 cun 5 fen high, 3 chi 6 cun wide, and 8 chi 4 cun long. The cap (maozi) is 1 chi 5 cun high, 3 chi 6 cun wide, and 6 chi 6 cun long. The arrow windows (jianchuang), 3 openings (yan), are each 6 cun wide and 7 cun 5 fen long; the outer opening, compared with the inner opening, is lower on the slant by 3 cun.\end{uniten}

\begin{unitzh}\unitlabel{shushu-050.4}取土用穿四堅三為率。週廻石版鋪城脚三層，每片長五尺，闊二尺，厚五寸。通身用甎包砌：下一丈九幅{\zhsmall 按：九幅即九層。}，中一丈七幅，上一丈五幅。甎每片長一尺二寸，闊六寸，厚二寸五分。䕶嶮牆髙三尺，闊一尺二寸{\zhsmall 按：此句贅。盖下既論幅，則不用闊，下當云每幅之闊為甎之長數。}，下脚髙一尺五寸，鋪甎三幅；上一尺五寸，鋪甎二幅。\end{unitzh}
\begin{uniten}For taking the earth, dug (chuan) 4 to compacted (jian) 3 is used as the rate (lü). All around the circuit, stone slabs are laid at the foot of the city wall in 3 layers (ceng); each slab is 5 chi long, 2 chi wide and 5 cun thick. The whole body is encased in laid bricks: the lower 1 zhang in 9 strips (fu) [Siku editors: 9 strips means 9 layers], the middle 1 zhang in 7 strips, the upper 1 zhang in 5 strips. Each brick is 1 chi 2 cun long, 6 cun wide and 2 cun 5 fen thick. The danger-guarding wall (huxian qiang) is 3 chi high and 1 chi 2 cun wide [Siku editors: This phrase is superfluous. For since below it is reckoned in strips, the width is not used; below it ought to say that the width of each strip is the length of a brick]; its lower foot, 1 chi 5 cun high, is laid with 3 strips of bricks; the upper 1 chi 5 cun is laid with 2 strips of bricks.\end{uniten}

\begin{unitzh}\unitlabel{shushu-050.5}每長一丈，用木物料：永定柱二十條，長三丈五尺，徑一尺，每條栽埋工七分，串鑿工三分；爬頭拽後木共八十條，長二丈，徑七寸，每條作工三分，串鑿工二分；搏子木二百條，長一丈，徑三寸，每條作工二分，搬扛工二分；絍橛二千個，每個長一尺，方一寸，每個工七毫；絍索二千條，長一丈，徑五分，每條工九毫；石版一十片，匠一工，搬一工，每片灰一十觔；搬灰千觔用一工；甎匠每工砌七百片；石灰每甎一觔；蘆蓆一百五十領；青茅五百束；絲竿筀竹五十條；笍子水竹一十把，每把二尺圍；钁手、鍬手、擔土、杵手每工各六十尺；火頭一名受六十工；部押壕寨一名管一百二十工；每工日支新會一百文，米二升五合。欲知城牆堅積、壕積、壕深，共用木、竹、橛、索、甎、石、灰、蘆、茅、人工、錢、米數各㡬何？\end{unitzh}
\begin{uniten}For each 1 zhang of length, the timber and materials used are: permanently fixed posts (yongding zhu), 20 pieces, 3 zhang 5 chi long and 1 chi in diameter, each piece requiring 7 fen of a workday (gong) for planting and burying and 3 fen for piercing and chiselling [Translator's note: in the labour figures, fen and hao are decimal fractions of a workday, 1/10 and 1/1000]; crawling-head and dragging-back timbers (patou yehou mu), together 80 pieces, 2 zhang long and 7 cun in diameter, each piece requiring 3 fen of a workday for making and 2 fen for piercing and chiselling; clasp timbers (bozi mu), 200 pieces, 1 zhang long and 3 cun in diameter, each piece requiring 2 fen of a workday for making and 2 fen for carrying and shouldering; tying pegs (renjue), 2000, each 1 chi long and 1 cun square, each requiring 7 hao (毫) of a workday; tying ropes (rensuo), 2000 pieces, 1 zhang long and 5 fen in diameter, each requiring 9 hao of a workday; stone slabs, 10 pieces, masons 1 workday, carrying 1 workday, and for each slab 10 jin (觔) of lime; carrying 1000 jin of lime uses 1 workday; a brick mason lays 700 pieces per workday; lime, 1 jin per brick; reed mats, 150 sheets (ling); green thatch, 500 bundles; silk-pole (sigan) gui bamboo (guizhu), 50 pieces; ruizi water bamboo (shuizhu), 10 bunches, each bunch 2 chi in circumference; mattock-men (jueshou), spade-men (qiaoshou), earth-carriers (dantu) and pestle-men (chushou), each 60 chi per workday; one cook (huotou) takes on 60 workdays; one supervising moat-and-stockade officer (buya haozhai) manages 120 workdays; for each workday there is paid daily 100 wen (文) of new hui (xinhui) and 2 sheng (升) 5 ge (合) of rice. One wishes to know the compacted volumes of the city wall and the wall, the volume of the moat, the depth of the moat, and the total quantities used of timber, bamboo, pegs, ropes, bricks, stone, lime, reeds, thatch, labour, money and rice: how much is each?\end{uniten}

\begin{unitzh}\unitlabel{shushu-050.6}按：題意掘土為壕以築城，城身及羊馬牆身共積即壕身積，語中未詳。羊馬牆及壕周遶城外，當長于城周，題中未載距城尺數。城牆用甎包砌，當計三靣，題中只計一靣。皆屬疎漏。䕶險牆應以甎長為闊，題中復言闊之尺數；柱、木、䋲、橛等徑、方、長短，術草不用其數，題中亦皆𨳩載，未免冗繁。然古啇功之略猶可見焉，故于草中就其所問之意而改正之。\end{unitzh}
\begin{uniten}Note: The intent of the problem is that earth is dug to make the moat in order to build the city wall; the combined volume of the body of the city wall and the body of the sheep-and-horse wall is the volume of the body of the moat; the wording does not make this clear. The sheep-and-horse wall and the moat encircle the outside of the city wall and ought to be longer than the city wall's circumference; the problem does not record their distance from the city wall in chi. The city wall is encased in laid bricks, and three faces ought to be reckoned; the problem reckons only one face. All of these are oversights. The danger-guarding wall ought to take the brick's length as its width, yet the problem also states the width in chi; and as for the diameters, squares and lengths of the posts, timbers, ropes, pegs and so on, the method and the working do not use these numbers, yet the problem also sets them all out; this cannot avoid being superfluous and cumbersome. Nevertheless, an outline of the ancient assessing of works (shanggong) can still be seen in it; therefore, in the working, we correct it in accordance with the intent of what is asked.\end{uniten}

\begin{unitzh}\unitlabel{shushu-050.7}答曰：城積二千三百七十八萬二千五百尺，堅積；牆積一百一十三萬二千五百尺，堅積；壕積三千三百二十二萬尺，穿積；壕深八丈；永定柱三萬二百條，每條長三丈五尺，徑一尺；爬頭拽後木一十二萬八百條，每條長二丈，徑七寸；搏子木三十萬二千條，毎條長一丈，徑三寸；絍橛子三百二萬個，每個長一尺，方一寸；絍索三百二萬條，每條長一丈，徑五分；蘆蓆二十二萬六千五百領；青茅七十五萬五千束，每束六尺圍；筀竹七萬五千五百竿；水竹一萬五千一百把，毎把二寸圍；石版一萬五千一百片；城甎一千二百八十三萬三千四百九十片；石灰一千二百九十八萬四千四百九十觔；用工二百萬三千七百七十工；新會二十萬三百七十七貫文；支米五萬九十四石二斗五升。\end{unitzh}
\begin{uniten}Answer: Volume of the city wall, 23782500 chi, compacted volume; volume of the wall, 1132500 chi, compacted volume; volume of the moat, 33220000 chi, dug volume; depth of the moat, 8 zhang; permanently fixed posts, 30200 pieces, each 3 zhang 5 chi long and 1 chi in diameter; crawling-head and dragging-back timbers, 120800 pieces, each 2 zhang long and 7 cun in diameter; clasp timbers, 302000 pieces, each 1 zhang long and 3 cun in diameter; tying pegs, 3020000, each 1 chi long and 1 cun square; tying ropes, 3020000 pieces, each 1 zhang long and 5 fen in diameter; reed mats, 226500 sheets; green thatch, 755000 bundles, each bundle 6 chi in circumference; gui bamboo, 75500 poles; water bamboo, 15100 bunches, each bunch 2 cun in circumference; stone slabs, 15100 pieces; city-wall bricks, 12833490 pieces; lime, 12984490 jin; labour used, 2003770 workdays; new hui, 200377 guan (貫) wen; rice paid out, 50094 dan (石) 2 dou (斗) 5 sheng.\end{uniten}

\begin{unitzh}\unitlabel{shushu-050.8}按：右答數石版、城甎、石灰三數俱誤，餘數亦多偶合，具詳草後。\end{unitzh}
\begin{uniten}Note: Of the answers above, the three numbers for the stone slabs, the city-wall bricks and the lime are all wrong; many of the remaining numbers also agree only by chance. Details are given after the working.\end{uniten}

\begin{unitzh}\unitlabel{shushu-050.9}術曰：以啇功求之。置城及牆上下廣，各併之，乘髙，進位，半之，各得每丈積率；併之，為共率。先以每功尺除之，又以諸色工各數乘之，為土功丈率。次置柱、木、橛、索，乘其每條叚功，得各共功。次置城方一丈，自之，乘用甎總幅數，為實；以甎長乘厚，為則法；除實，得城身用甎。次置鵲臺、座子、肩子、㡌子各髙、闊、長相乘，為寄，併之于上。次箭窻眼髙依差寸求斜深虚積，減寄，餘為女頭甎實；以則法乘甎闊，為甎積法；除之，得女頭臺用磚。\end{unitzh}
\begin{uniten}Method: Seek it by assessing works (shanggong). Set up the upper and lower widths of the city wall and of the wall; add each pair; multiply by the height; advance a place; halve it; in each case one obtains the volume rate (lü) per zhang. Add them, making the combined rate. First divide it by the chi per unit of work (gong); then multiply by the respective numbers of the various kinds of labour, making the earthwork rate per zhang. Next set up the posts, timbers, pegs and ropes, and multiply them by the work of each item per piece, obtaining the respective total work. Next set up the city wall's square of 1 zhang, multiply it by itself, and multiply by the total number of strips of bricks used, making the dividend (shi); multiply the brick's length by its thickness, making the standard divisor (zefa); divide the dividend, obtaining the bricks used for the body of the city wall. Next set up the height, width and length of each of the magpie platform, the seat, the shoulder and the cap, multiply them together, making deposits (ji), and add them together above. Next, from the height of the arrow-window openings, according to the difference in cun, seek the empty volume of the slanting depth; subtract it from the deposit; the remainder makes the dividend for the battlement bricks; multiply the standard divisor by the brick's width, making the brick-volume divisor; divide by it, obtaining the bricks used for the battlement platform.\end{uniten}

\begin{unitzh}\unitlabel{shushu-050.10}又置䕶險牆髙，以丈乘而半之，又乘上下幅共數，為實；以磚闊、厚相乘，為法；除之，得䕶險牆用甎。併三項用甎，為都實；以毎功片為法，除之，得甎匠功。以每丈用石叚數求石匠工；以搬每丈石求搬石工。以片用灰數乘都甎，得甎用灰；以每丈石版數乘片用灰，得石用灰；併之，為甎石共灰；以每工搬灰數除之，得搬灰工。併諸作工，為實；以火頭、壕寨每管人數各為法；除之，得各數；又併之，為都工。然後以城圍通長偏乘諸項毎丈率積、灰、各工料，得共數。\end{unitzh}
\begin{uniten}Again set up the height of the danger-guarding wall, multiply by the zhang and halve it, and again multiply by the combined number of upper and lower strips, making the dividend; multiply the brick's width and thickness together, making the divisor (fa); divide by it, obtaining the bricks used for the danger-guarding wall. Add the bricks used in the three items, making the total dividend (dushi); take the pieces per unit of work as divisor and divide by it, obtaining the brick masons' work. With the number of stone pieces used per zhang, seek the stone masons' workdays; with the carrying of the stone per zhang, seek the stone-carrying workdays. Multiply the total bricks by the amount of lime used per piece, obtaining the lime used for the bricks; multiply the lime used per piece by the number of stone slabs per zhang, obtaining the lime used for the stone; add them, making the combined lime for bricks and stone; divide by the amount of lime carried per workday, obtaining the lime-carrying workdays. Add the workdays of all the tasks, making the dividend; take the numbers of people managed by each cook and each moat-and-stockade officer respectively as divisors; divide, obtaining the respective numbers; again add them, making the total workdays. After that, with the full length of the city wall's circumference, multiply throughout all the items — the per-zhang rate volumes, the lime, and each of the labour and materials — obtaining the totals.\end{uniten}

\begin{unitzh}\unitlabel{shushu-050.11}女頭圍{\zhsmall 按：舊圖在術前，其尺數多未合，今改正，移于此。}\end{unitzh}
\begin{uniten}Battlement circumference [Translator's note: 圍 "circumference" is evidently for 圖 "diagram".] [Siku editors: The old diagram was placed before the method, and many of its measurements in chi did not agree; it has now been corrected and moved here.]\end{uniten}

\begin{unitzh}\unitlabel{shushu-050.12}草曰：置城上廣三丈，併下廣七丈五尺，得一十丈五尺；乘髙三丈，得三千一百五十尺；進位{\zhsmall 按：即十尺乘之。}，得三萬一千五百\end{unitzh}
\begin{uniten}Working: Set up the city wall's upper width, 3 zhang; add the lower width, 7 zhang 5 chi, obtaining 10 zhang 5 chi; multiply by the height, 3 zhang, obtaining 3150 chi; advance a place [Siku editors: that is, multiply it by 10 chi], obtaining 31500 [continues]\end{uniten}

\begin{unitzh}\unitlabel{shushu-051.1}尺，半之，得一萬五千七百五十尺，為每丈城積率。置羊馬牆闊五尺，併下闊一丈，得一十五尺，乘髙一丈，得一百五十尺，進位，得一千五百尺，以半之，得七百五十尺，為羊馬牆每丈積率。併城牆二率，得一萬六千五百尺，為共率，以為實。以钁、鍬、擔工、杵手各六十尺為法，除實，得二百七十五工，以四色因之，得一千一百工，為钁、鍬、擔土、杵手工。\end{unitzh}
\begin{uniten}[continued] chi (尺); halve it, obtaining 15750 chi, as the rate (lü) of city-wall volume per zhang (丈). Set up the width of the sheep-and-horse wall (yangma qiang), 5 chi; add the lower width, 1 zhang, obtaining 15 chi; multiply by the height, 1 zhang, obtaining 150 chi; advance a place (jinwei), obtaining 1500 chi; halve it, obtaining 750 chi, as the rate of volume per zhang of the sheep-and-horse wall. Add the 2 rates of the city wall and the wall, obtaining 16500 chi, as the combined rate, and take it as the dividend (shi). Take the mattock (jue), spade (qiao), carrying labor (dan gong) and rammer-hands (chushou), 60 chi each, as the divisor (fa); dividing the dividend, obtain 275 gong (工, workdays); multiply it by the 4 kinds (se), obtaining 1100 gong, as the labor of mattock, spade, carrying earth and rammer-hands.\end{uniten}

\begin{unitzh}\unitlabel{shushu-051.2}置永定柱二十條，乘每條栽埋工七分，得一十四工；又乘串鑿工三分，得六工；計二十工，為永定柱工。置爬頭拽後木八十條，乘作工三分，得二十四工；又乘串鑿工二分，得一十六工；計四十工，為爬頭拽後木工。置搏子木二百條，乘作工二分，得四十工；又乘搬扛工二分，得四十工；計八十工，為搏子木工。置絍撅二千箇，乘作工七毫，得一十四工；絍索二千條，乘作工九毫，得一十八工；計三十二工，為橛索共工。\end{unitzh}
\begin{uniten}Set up the 20 permanent posts (yongding zhu); multiply by the planting-and-burying labor (zaimai gong) per post, 7 fen (分), obtaining 14 gong; again multiply by the threading-and-chiselling labor (chuanzao gong), 3 fen, obtaining 6 gong; in all 20 gong, as the labor for the permanent posts. Set up the 80 climbing-head drag-back timbers (patou yehou mu); multiply by the making labor (zuo gong), 3 fen, obtaining 24 gong; again multiply by the threading-and-chiselling labor, 2 fen, obtaining 16 gong; in all 40 gong, as the labor for the climbing-head drag-back timbers. Set up the 200 clasp timbers (bozi mu); multiply by the making labor, 2 fen, obtaining 40 gong; again multiply by the carrying-and-shouldering labor (bankang gong), 2 fen, obtaining 40 gong; in all 80 gong, as the labor for the clasp timbers. Set up the 2000 threading stakes (renjue); multiply by the making labor, 7 hao (毫), obtaining 14 gong; the 2000 threading ropes (rensuo), multiplied by the making labor, 9 hao, give 18 gong; in all 32 gong, as the combined labor for stakes and ropes.\end{uniten}

\begin{unitzh}\unitlabel{shushu-051.3}乃以城牆女頭甎積求甎匠工。置城身方一丈自乘，一百尺于上；次置下九幅、中七幅、上五幅，併之，得二十一幅，乘上，得二千一百尺。甎長有寸，以寸通之，為二十一萬寸，為實。{\zhsmall 按：此條誤，應減去三層石脚積，餘為甎實，方合。}以甎長一十二寸乘厚二寸五分，得三十寸，為則法{\zhsmall 按：則法即甎立長靣數，葢論幅數，則不用求積也。}，除實，得七千片，為城身甎數。\end{unitzh}
\begin{uniten}Then, from the brick volume of the city wall and the parapets (nütou), find the labor of the brickmasons. Set up the wall body (chengshen), 1 zhang square, multiplied by itself: 100 chi, placed above; next set up the lower 9 widths (fu), the middle 7 widths and the upper 5 widths; add them, obtaining 21 widths; multiply the above, obtaining 2100 chi. The brick length has cun (寸), so convert it into cun, making 210000 cun, as dividend. [Siku editors: This item is wrong: one should subtract the volume of the 3 layers of stone footing and take the remainder as the brick dividend; only then does it agree.] Multiply the brick length, 12 cun, by the thickness, 2 cun 5 fen, obtaining 30 cun, as the standard divisor (zefa) [Siku editors: The standard divisor is the number for the upright long face of a brick; for since one reckons by the number of widths, one does not need to find a volume]; dividing the dividend, obtain 7000 pieces (pian), as the number of bricks for the wall body.\end{uniten}

\begin{unitzh}\unitlabel{shushu-051.4}又置鵲臺髙五寸，乘闊五尺四寸，得二百七十寸；又乘長一丈，得二萬七千寸，寄上。又置座子髙二尺二寸五分，乘闊三尺六寸，得八百一十寸；又乘長一丈，得八萬一千寸，加寄。又置肩子髙一尺二寸五分，乘闊三尺六寸，得四百五十寸；又乘長八尺四寸，得三萬七千八百寸，又加寄。又置㡌子髙一尺五寸，乘闊三尺六寸，又乘長六尺六寸，得三萬五千六百四十寸，又加寄，共得一十八萬一千四百四十寸，共為寄。\end{unitzh}
\begin{uniten}Again set up the height of the magpie platform (quetai), 5 cun; multiply by the width, 5 chi 4 cun, obtaining 270 cun; again multiply by the length, 1 zhang, obtaining 27000 cun; lodge it above. Again set up the height of the seat (zuozi), 2 chi 2 cun 5 fen; multiply by the width, 3 chi 6 cun, obtaining 810 cun; again multiply by the length, 1 zhang, obtaining 81000 cun; add it to the lodged amount. Again set up the height of the shoulder (jianzi), 1 chi 2 cun 5 fen; multiply by the width, 3 chi 6 cun, obtaining 450 cun; again multiply by the length, 8 chi 4 cun, obtaining 37800 cun; again add it to the lodged amount. Again set up the height of the cap (maozi), 1 chi 5 cun; multiply by the width, 3 chi 6 cun, and again by the length, 6 chi 6 cun, obtaining 35640 cun; again add it to the lodged amount; altogether obtaining 181440 cun, all as the lodged amount.\end{uniten}

\begin{unitzh}\unitlabel{shushu-051.5}其箭窻内外眼雖差三寸，于斜深虚積将盈補虧，與直深等。以窻闊六寸乘長七寸五分，得四十五寸；又乘座闊三尺六寸，得一千六百二十寸，為窻積。{\zhsmall 按：箭眼三，當三因之，此用一眼數，誤。}以減寄，餘一十七萬九千八百二十寸，為實。置甎則法三十寸，乘甎闊六寸，得一百八十寸，為甎則法，除實，得九百九十九片，為女頭鵲臺共甎。\end{unitzh}
\begin{uniten}As for the arrow windows (jianchuang), although their inner and outer openings (yan) differ by 3 cun, in the empty volume of the slanting depth the excess is taken to make up the deficiency, so that it is equal to a straight depth. Multiply the window width, 6 cun, by the length, 7 cun 5 fen, obtaining 45 cun; again multiply by the width of the seat, 3 chi 6 cun, obtaining 1620 cun, as the window volume. [Siku editors: There are 3 arrow openings, so one should multiply by 3; here the figure for 1 opening is used, which is wrong.] Subtract it from the lodged amount; the remainder, 179820 cun, is the dividend. Set up the brick standard divisor, 30 cun; multiply by the brick width, 6 cun, obtaining 180 cun, as the brick standard divisor; dividing the dividend, obtain 999 pieces, as the bricks for the parapets and the magpie platform together.\end{uniten}

\begin{unitzh}\unitlabel{shushu-051.6}又置䕶險牆髙三尺，乘每丈，得三千寸，以牆法當半折之，得一千五百寸；又乘上下五幅，得七千五百寸，為實。次以甎厚二寸五分乘闊六寸，得一十五寸，為甎法，除實，得五百片，為䕶險牆甎。次併三項甎八千四百九十九片，為毎丈用甎都實。以每工七百片為法，除實，得一十二工七百分工之九十九，為甎工。\end{unitzh}
\begin{uniten}Again set up the height of the danger-guarding wall (huxian qiang), 3 chi; multiply by each zhang, obtaining 3000 cun; by the rule for walls it should be reduced by half, obtaining 1500 cun; again multiply by the upper and lower 5 widths, obtaining 7500 cun, as dividend. Next multiply the brick thickness, 2 cun 5 fen, by the width, 6 cun, obtaining 15 cun, as the brick divisor; dividing the dividend, obtain 500 pieces, as the bricks for the danger-guarding wall. Next add the bricks of the 3 items, 8499 pieces, as the total dividend (dushi) of bricks used per zhang. Take 700 pieces per gong as divisor; dividing the dividend, obtain 12 gong and 99/700 of a gong, as the brick labor.\end{uniten}

\begin{unitzh}\unitlabel{shushu-051.7}每丈用石版十片{\zhsmall 按：此條誤，應用一十六片二分。}，計一工；搬石十片，計一工。甎每片用灰一觔，命都甎，即甎用灰之數。又置每丈用石版一十片，每片用灰一十觔，相乘之，得一百觔，為石版用灰；併甎用灰八千四百九十九，得八千五百九十九觔，為甎石用灰數，為實。以毎工搬一千觔為法，除之，得八工一千分工之四百九十九，為搬灰工。\end{unitzh}
\begin{uniten}Per zhang 10 pieces of stone slab (shiban) are used [Siku editors: This item is wrong; 16 pieces 2 fen should be used], reckoned as 1 gong; carrying 10 pieces of stone is reckoned as 1 gong. Each brick uses 1 jin (觔) of lime; naming it by the total of bricks gives the amount of lime used for the bricks. Again set up the 10 pieces of stone slab used per zhang, each piece using 10 jin of lime; multiplying them together, obtain 100 jin, as the lime used for the stone slabs; add the lime used for the bricks, 8499, obtaining 8599 jin, as the amount of lime used for bricks and stone, as dividend. Take 1000 jin carried per gong as divisor; dividing, obtain 8 gong and 499/1000 of a gong, as the labor for carrying lime.\end{uniten}

\begin{unitzh}\unitlabel{shushu-051.8}并石匠搬石二工，通前列土工一千一百、定柱工二十、爬頭拽後木工四十、搏子木工八十、橛索工三十二、甎工一十二工七百分工之九十九{\zhsmall 按：此數誤。}、搬灰八工千分工之四百九十九{\zhsmall 按：此數誤。}、石匠搬石共二工{\zhsmall 按：此數誤。}，併諸作工，餘分不同者，合分術入之，共得一千二百九十四工七千分工之一千四百八十九。{\zhsmall 按：甎工差多一工，石工差少一工，故相併之工數無差；分子則差少四千八百五十二分。}\end{unitzh}
\begin{uniten}Combine the 2 gong of the stonemasons and the stone carrying; together with those listed before — earth labor 1100, fixed-post labor 20, climbing-head drag-back timber labor 40, clasp-timber labor 80, stake-and-rope labor 32, brick labor 12 gong and 99/700 of a gong [Siku editors: This number is wrong], lime carrying 8 gong and 499/1000 of a gong [Siku editors: This number is wrong], stonemasons and stone carrying together 2 gong [Siku editors: This number is wrong] — add the labor of all the works; where the remaining fractions are not the same, bring them in by the method of combining fractions (hefen shu); altogether obtain 1294 gong and 1489/7000 of a gong. [Siku editors: The brick labor is too large by 1 gong and the stone labor too small by 1 gong, so the number of gong when added together has no discrepancy; the numerator, however, is too small by 4852 parts.]\end{uniten}

\begin{unitzh}\unitlabel{shushu-051.9}通分内子，得一百二十九萬五千四百八十九，為衆工實。置火頭每管六十人，分母乘之，得六萬，為法{\zhsmall 按：數誤，分母七千乘六十，應得四十二萬。}，除都工實，得火頭二十一人六萬分之三萬五千四百八十九。壕寨每部一百二十人，就倍火頭法六萬為十二萬，亦除衆工實，得壕寨十人十萬分之九萬五千四百八十九。\end{unitzh}
\begin{uniten}Converting by the denominator and including the numerator (tongfen nazi), obtain 1295489, as the dividend of all the labor. Set up the 60 persons managed by each cook (huotou); multiply by the denominator, obtaining 60000, as divisor [Siku editors: The number is wrong: the denominator 7000 multiplied by 60 should give 420000]; dividing the dividend of the total labor, obtain for the cooks 21 persons and 35489/60000. Each moat-stockade overseer (haozhai) commands 120 persons; accordingly double the cooks' divisor, 60000, making 120000; again dividing the dividend of all the labor, obtain for the moat-stockade overseers 10 persons and 95489/100000.\end{uniten}

\begin{unitzh}\unitlabel{shushu-051.10}列兩餘分及前諸作工餘七千分之一千四百八十九，三項，以合分術入之，得一工，不盡五十萬四十億分之三十萬二千三百六十九億四千萬分；求等，又約之，為八十四萬分之五十萬三千九百四十九分。乃又併之，共得一千三百二十六工，其餘分大約百分中之五十九，在半以上，收為一工，共定得一千三百二十七工{\zhsmall 按：此數偶合。}，為每丈都工。\end{unitzh}
\begin{uniten}Set out the 2 remaining fractions and the remainder of the labor of all the works before, 1489/7000: 3 items; bring them in by the method of combining fractions, obtaining 1 gong, with an unexhausted remainder of 30236940000000/50004000000000 parts [Translator's note: the reduction that follows implies 五十萬四千億, 50400000000000]; find the common measure (qiu deng) and reduce it again, making 503949/840000 parts. Then add them again, altogether obtaining 1326 gong; the remaining fraction is roughly 59 parts in 100; being above half, it is taken as 1 gong; altogether it is settled at 1327 gong [Siku editors: This number agrees by chance], as the total labor per zhang.\end{uniten}

\begin{unitzh}\unitlabel{shushu-051.11}然後以城通長徧乘乘諸項。置城長一千五百一十丈，乘城率一萬五千七百五十尺，得二千三百七十八萬二千五百尺，為城堅積。又以城長乘牆率七百五十尺，得一百一十三萬二千五百尺，為牆堅積。併城牆二積，得二千四百九十一萬五千尺；又以墟率四因之，得九千九百六十六萬尺，為實；以堅率三約，得三千三百二十二萬尺，為壕積，以為實。以壕闊三十丈併下闊二十五丈，得五十五丈，以半之，得二百七十五尺，乘壕長一千五百一十丈，得四十一萬五千二百五十尺，為壕法，除實，得八丈，為壕深。\end{unitzh}
\begin{uniten}After that, by the total length of the wall multiply throughout, multiplying the various items. Set up the length of the wall, 1510 zhang; multiply by the city-wall rate, 15750 chi, obtaining 23782500 chi, as the firm volume (jianji) of the city wall. Again multiply the length of the wall by the wall rate, 750 chi, obtaining 1132500 chi, as the firm volume of the wall. Add the 2 volumes of the city wall and the wall, obtaining 24915000 chi; again multiply it by the hollow-ground rate (xu lü), 4, obtaining 99660000 chi, as dividend; reduce by the firm-earth rate (jian lü), 3, obtaining 33220000 chi, as the moat volume, and take it as dividend. Add the moat width, 30 zhang, and the lower width, 25 zhang, obtaining 55 zhang; halve it, obtaining 275 chi; multiply by the moat length, 1510 zhang, obtaining 415250 chi, as the moat divisor; dividing the dividend, obtain 8 zhang, as the moat depth.\end{uniten}

\begin{unitzh}\unitlabel{shushu-051.12}求工料共數：如術，以城通長徧乘丈率工、永定柱、爬頭拽後木、搏子木、橛子木、絍索、蘆蓆、筀竹、水竹、青茅、城甎、石版、石灰，各得。以共工乘日支錢米，得共錢米。更不立革。\end{unitzh}
\begin{uniten}To find the total amounts of labor and materials: as in the method, by the total length of the wall multiply throughout the per-zhang rates of labor, permanent posts, climbing-head drag-back timbers, clasp timbers, stake timbers (juezi mu), threading ropes, reed mats (luxi), gui bamboo (guizhu), water bamboo (shuizhu), green thatch (qingmao), wall bricks, stone slabs and lime; each is obtained. Multiply the total labor by the daily allowance of cash and rice, obtaining the total cash and rice. No further ge (革) is set up [Translator's note: 革 evidently for 草, "working"].\end{uniten}

\begin{unitzh}\unitlabel{shushu-051.13}按：草中法數有不合者，逐條改正如左。\end{unitzh}
\begin{uniten}Note: Where the methods and numbers in the working do not agree, they are corrected item by item as follows.\end{uniten}

\begin{unitzh}\unitlabel{shushu-051.14}一、題言城脚週廻鋪石版三層，通身用甎包砌。應于先求得一丈自乘幅數再乘之二十一萬寸數内，減去三層石脚之一萬三千五百寸，餘一十九萬六千五百寸，為甎實；以甎則三十寸除之，得六千五百五十觔，為甎數。草中即以全積為甎實，故得數差多四百五十片。\end{unitzh}
\begin{uniten}Item: The problem says that around the foot of the wall stone slabs are laid in 3 layers all round, and that the whole body is faced with bricks. One should, from the number 210000 cun obtained before by multiplying 1 zhang by itself and again by the number of widths, subtract the 13500 cun of the 3 layers of stone footing; the remainder, 196500 cun, is the brick dividend; dividing it by the brick standard, 30 cun, obtain 6550 jin, as the number of bricks. In the working the whole volume was taken directly as the brick dividend; hence the number obtained is too large by 450 pieces.\end{uniten}

\begin{unitzh}\unitlabel{shushu-051.15}一、女頭鵲臺全積一十八萬一千四百四十寸，一箭窻虚積一千六百二十寸，三眼應三因之，得四千八百六十寸，以減全積，餘一十七萬六千五百八十寸，為實；以一甎積一百八十寸除之，得九百八十一片，為甎數。草中只減一眼虚積，故得數差多一十八片。\end{unitzh}
\begin{uniten}Item: The whole volume of the parapets and the magpie platform is 181440 cun; the empty volume of 1 arrow window is 1620 cun; for the 3 openings one should multiply it by 3, obtaining 4860 cun; subtract it from the whole volume; the remainder, 176580 cun, is the dividend; dividing it by the volume of 1 brick, 180 cun, obtain 981 pieces, as the number of bricks. In the working only the empty volume of 1 opening was subtracted; hence the number obtained is too large by 18 pieces.\end{uniten}

\begin{unitzh}\unitlabel{shushu-051.16}一、併城牆、鵲臺、䕶險牆三甎數，得共用甎八千零三十一片，草中差多四百六十八片。\end{unitzh}
\begin{uniten}Item: Adding the 3 numbers of bricks for the city wall, the magpie platform and the danger-guarding wall, obtain a total of 8031 pieces of brick used; in the working it is too large by 468 pieces.\end{uniten}

\begin{unitzh}\unitlabel{shushu-051.17}一、題言城脚鋪石版三層，應以甎闊六寸乘九幅，得五十四寸，以一丈乘之，得五千四百寸，三因之，得一萬六千二百寸，為實；以石之長闊相乘，得一千寸，為法，除之，得一十六片二分。草中並無求法，但云毎丈用石版十片，差少六片二分。\end{unitzh}
\begin{uniten}Item: The problem says that at the foot of the wall stone slabs are laid in 3 layers. One should multiply the brick width, 6 cun, by the 9 widths, obtaining 54 cun; multiply it by 1 zhang, obtaining 5400 cun; multiply by 3, obtaining 16200 cun, as dividend; multiply together the length and width of a stone, obtaining 1000 cun, as divisor; dividing, obtain 16 pieces 2 fen. In the working there is no method for finding this at all; it merely says that 10 pieces of stone slab are used per zhang, which is too small by 6 pieces 2 fen.\end{uniten}

\begin{unitzh}\unitlabel{shushu-051.18}一、石每片用灰十觔，甎每片用灰一觔，毎丈石用灰一百六十二觔，甎用灰八千零三十一觔，二共用灰八千一百九十三觔。草中石用灰差少六十二觔，甎用灰差多四百六十八觔，二共用灰差多四百零六觔。\end{unitzh}
\begin{uniten}Item: Each piece of stone uses 10 jin of lime, and each brick uses 1 jin of lime; per zhang the stone uses 162 jin of lime and the bricks use 8031 jin of lime; the 2 together use 8193 jin of lime. In the working, the lime used for the stone is too small by 62 jin, the lime used for the bricks is too large by 468 jin, and the lime used for the 2 together is too large by 406 jin.\end{uniten}

\begin{unitzh}\unitlabel{shushu-051.19}一、石十片二工，甎七百片一工，灰千觔一工。每丈應用石工三工二分四釐，甎工一十一工七百分工之三百三十一，灰工八工千分工之一百九十三；併之，又併土工一千一百工、定柱工二十工、爬頭木工四十工、搏子木工八十工、橛索工三十二工，共得一千二百九十四工七千分之六千三百四十一。草中分子為一千四百八十九，差少四千八百五十二分。\end{unitzh}
\begin{uniten}Item: Stone, 10 pieces, 2 gong; bricks, 700 pieces, 1 gong; lime, 1000 jin, 1 gong. Per zhang one should use stone labor of 3 gong 2 fen 4 li (釐), brick labor of 11 gong and 331/700 of a gong, and lime labor of 8 gong and 193/1000 of a gong; add them, and further add the earth labor 1100 gong, fixed-post labor 20 gong, climbing-head timber labor 40 gong, clasp-timber labor 80 gong and stake-and-rope labor 32 gong; altogether obtain 1294 gong and 6341/7000. In the working the numerator is 1489, which is too small by 4852 parts.\end{uniten}

\begin{unitzh}\unitlabel{shushu-051.20}一、求大頭工，應以共工分母通工數，内子，為實；以分母通所受六十人，為法；除實，得火頭二十一工又一十四萬分工之八萬一千四百四十七。草中誤以千分通之，故分母子數不合。求壕寨，應仍用前實，以分母七千通所管一百二十人，為法，除實，得壕寨十工又二十八萬分之二十二萬一千四百四十七。草中同前誤，分母子數不合。\end{unitzh}
\begin{uniten}Item: To find the labor of the datou (大頭) [Translator's note: evidently for 火頭], one should convert the number of gong by the denominator of the total labor and include the numerator, as dividend; convert the 60 persons received by the denominator, as divisor; dividing the dividend, obtain for the cooks 21 gong and further 81447/140000 of a gong. In the working it was wrongly converted with thousandths; hence the numbers of the denominator and numerator do not agree. To find the moat-stockade overseers, one should still use the dividend above, and convert the 120 persons managed by the denominator 7000, as divisor; dividing the dividend, obtain for the moat-stockade overseers 10 gong and further 221447/280000. In the working there is the same error as above; the numbers of the denominator and numerator do not agree.\end{uniten}

\begin{unitzh}\unitlabel{shushu-052.1}一、前共工、火頭、壕寨三分數通而併之，得二工又二十八萬分之七萬七千九百八十一；然後總併之，得一千三百二十七工又二十八萬分工之七萬七千九百八十一。其分子不及分母十之三，棄之，即以一千三百二十七工為定數。草中分子僅進一工，餘數及十之六，又進一工，故所得定工數亦同。\end{unitzh}
\begin{uniten}Item: The fractions of the three, namely the preceding total labour, the cooks (huotou) and the moat-and-stockade officers (haozhai), are brought to a common denominator and added, giving 2 gong (工) and 77981/280000. Then, adding everything together, one obtains 1327 gong and 77981/280000 gong. Its numerator does not reach 3/10 of the denominator, so it is discarded, and 1327 gong is taken as the fixed number. In the Working, the numerators advanced only 1 gong, and the remaining number reached 6/10 and advanced 1 more gong; therefore the fixed number of gong obtained is also the same.\end{uniten}

\begin{unitzh}\unitlabel{shushu-052.2}一、共用石版二萬四千四百六十二觔，答數内差少九千三百六十二觔。一、共用甎一千二百一十二萬六千八百一十觔，答數内差多七萬零六百六十八觔。一、共用灰一千二百三十七萬一千四百三十觔，答數内差多六十一萬三千六十觔。右十餘條，皆依舊草正其舛訛；至立法之疎密，未暇論也。\end{unitzh}
\begin{uniten}Item: In all, stone slabs used: 24462 jin (觔); the answer is short by 9362 jin. Item: In all, bricks used: 12126810 jin; the answer is over by 70668 jin. Item: In all, lime used: 12371430 jin; the answer is over by 613060 jin. The ten-odd items above all correct its errors following the old Working; as for the looseness or rigour of the method as set up, there has been no leisure to discuss it.\end{uniten}

\begin{unitzh}\unitlabel{shushu-052.3}樓櫓功料\end{unitzh}
\begin{uniten}Labour and Materials for Watchtowers (loulu)\end{uniten}

\begin{unitzh}\unitlabel{shushu-052.4}問：築城合葢樓櫓六十處，每處一十間。䕶險髙四尺，長三丈，厚隨甎長。卧牛木一十一條，長一丈六尺，徑一尺一寸；搭腦木一十一條，長二丈，徑一尺；㸔壕柱一十一條，長一丈六尺，徑一尺二寸；副壕柱一十一條，長一丈五尺，徑一尺二寸；掛甲柱一十一條，長一丈三尺，徑一尺一寸；虎蹲柱一十一條，長七尺五寸，徑一尺；仰艎板木四十五條，長一丈，徑一尺二寸；平靣板木三十五條，長一丈，徑一尺二寸；串掛枋木七十三條，長五尺，徑一尺。仰板四八甎結砌三層，計六千片，每片用灰半觔。共用紙觔一百觔。牆甎長一尺二寸，闊六寸，厚二寸半。中板瓦七千五百片。一尺釘八箇，八寸釘二百七十箇，五寸釘一百箇，四寸釘五十箇，丁環二十箇。用工三百九十六人。欲知共用工料各幾何？{\zhsmall 按：䕶險牆每甎一片用灰一觔，題内缺。}\end{unitzh}
\begin{uniten}Question: In building a city wall, watchtowers are to be built at 60 places, each place with 10 bays (jian). The protective parapet (huxian) is 4 chi (尺) high and 3 zhang (丈) long; its thickness follows the length of a brick. Lying-ox timbers (woniu mu): 11 pieces, 1 zhang 6 chi long, 1 chi 1 cun (寸) in diameter; head-rest timbers (danao mu): 11 pieces, 2 zhang long, 1 chi in diameter; moat-watching posts (kanhao zhu): 11 pieces, 1 zhang 6 chi long, 1 chi 2 cun in diameter; auxiliary moat posts (fuhao zhu): 11 pieces, 1 zhang 5 chi long, 1 chi 2 cun in diameter; armour-hanging posts (guajia zhu): 11 pieces, 1 zhang 3 chi long, 1 chi 1 cun in diameter; crouching-tiger posts (hudun zhu): 11 pieces, 7 chi 5 cun long, 1 chi in diameter; upward-facing boat-plank timbers (yanghuang banmu): 45 pieces, 1 zhang long, 1 chi 2 cun in diameter; flat-surface plank timbers (pingmian banmu): 35 pieces, 1 zhang long, 1 chi 2 cun in diameter; linking hanging tie-beam timbers (chuangua fangmu): 73 pieces, 5 chi long, 1 chi in diameter. Upward-facing planks (yangban): four-eight bricks (siba zhuan) bonded and laid in 3 courses, totalling 6000 pieces, each piece using 1/2 jin of lime. In all, 100 jin of paper fibre (zhijin) is used. The wall bricks (qiang zhuan) are 1 chi 2 cun long, 6 cun wide and 2 1/2 cun thick. Medium flat tiles (zhong banwa): 7500 pieces. 1-chi nails: 8; 8-cun nails: 270; 5-cun nails: 100; 4-cun nails: 50; nail-rings (dinghuan): 20. Labour used: 396 men. One wishes to know how much of each item of labour and materials is used in all. [Siku editors: For the protective parapet wall, each brick uses 1 jin of lime; this is missing in the problem.]\end{uniten}

\begin{unitzh}\unitlabel{shushu-052.5}答曰：卧牛木六百六十條，搭腦木六百六十條，㸔壕柱六百六十條，副壕柱六百六十條，挂甲柱六百六十條，虎蹲柱六百六十條，串挂枋四千三百八十條，仰板木二千七百條，平板木二千一百條，城甎四萬八千片，四八甎三十六萬片，石灰二十萬八千觔，紙觔六千觔，中板瓦四十五萬片，丁環一千二百箇，一尺釘四百八十箇，八寸釘一萬六千二百箇，五寸釘六千箇，四寸釘三千箇，用工二萬三千七百六十人。\end{unitzh}
\begin{uniten}Answer: lying-ox timbers, 660 pieces; head-rest timbers, 660 pieces; moat-watching posts, 660 pieces; auxiliary moat posts, 660 pieces; armour-hanging posts, 660 pieces; crouching-tiger posts, 660 pieces; linking hanging tie-beams, 4380 pieces; upward-facing plank timbers (yangban mu), 2700 pieces; flat plank timbers (pingban mu), 2100 pieces; city-wall bricks (cheng zhuan), 48000 pieces; four-eight bricks, 360000 pieces; lime, 208000 jin; paper fibre, 6000 jin; medium flat tiles, 450000 pieces; nail-rings, 1200; 1-chi nails, 480; 8-cun nails, 16200; 5-cun nails, 6000; 4-cun nails, 3000; labour used, 23760 men.\end{uniten}

\begin{unitzh}\unitlabel{shushu-052.6}術曰：以啇功求之。置牆髙乘長，得寸為實；以甎闊乘厚為法，除之，得用甎及用灰。以處數並乘諸工料，得總用工料。\end{unitzh}
\begin{uniten}Method: Seek it by Assessing Works (shanggong). Set up the wall height, multiply by the length, obtaining cun, as the dividend (shi); multiply the brick width by the thickness as the divisor (fa); divide, obtaining the bricks used and the lime used. Multiply all the items of labour and materials alike by the number of places, obtaining the total labour and materials used.\end{uniten}

\begin{unitzh}\unitlabel{shushu-052.7}草曰：置牆髙四尺，通為四十寸；置長三丈，通為三百寸；相乘，得一萬二千寸為實。以甎闊六寸乗厚二寸五分，得一十五寸為法，除之，得八百片為牆甎，又為灰。并四八甎灰三千觔，共灰三千八百觔。乃以六十處徧乗總用工料：卧牛木、搭腦木、㸔壕柱、副壕柱、挂甲柱、虎蹲柱、仰板木、平板木、串挂枋木、四八甎、城甎、石灰、紙觔、中板瓦、一尺釘、八寸釘、五寸釘、四寸釘、釘環、工數，各項總數在前。\end{unitzh}
\begin{uniten}Working: Set up the wall height, 4 chi, converted to 40 cun; set up the length, 3 zhang, converted to 300 cun; multiply them together, obtaining 12000 cun as the dividend. Multiply the brick width, 6 cun, by the thickness, 2 cun 5 fen (分), obtaining 15 cun as the divisor; divide, obtaining 800 pieces as the wall bricks, and also as the lime. Together with the lime for the four-eight bricks, 3000 jin, the lime is in all 3800 jin. Then multiply throughout by the 60 places, for the total labour and materials used: lying-ox timbers, head-rest timbers, moat-watching posts, auxiliary moat posts, armour-hanging posts, crouching-tiger posts, upward-facing plank timbers, flat plank timbers, linking hanging tie-beam timbers, four-eight bricks, city-wall bricks, lime, paper fibre, medium flat tiles, 1-chi nails, 8-cun nails, 5-cun nails, 4-cun nails, nail-rings and the number of labourers; the totals for each item are given above.\end{uniten}

\begin{unitzh}\unitlabel{shushu-052.8}計造石壩\end{unitzh}
\begin{uniten}Computing the Building of a Stone Dam (shiba)\end{uniten}

\begin{unitzh}\unitlabel{shushu-052.9}問：創石壩一座，長三十丈，水深四丈二尺，令靣闊三丈。石版每片長五尺，闊二尺，厚五寸，用灰一十觔。每層髙二尺，差闊一尺。石匠每工九片；搬扛五片，用工四人；搬灰兼用，每工一百一十觔。火頭每名管六十人，部押每名管一百二十人。所用石湏依原叚，不許鑿動。欲知壩下闊及用石并灰、共工各幾何？\end{unitzh}
\begin{uniten}Question: A stone dam is to be newly built, 30 zhang long; the water is 4 zhang 2 chi deep; let the top width be 3 zhang. Each stone slab is 5 chi long, 2 chi wide and 5 cun thick, and uses 10 jin of lime. Each layer is 2 chi high, and the width differs by 1 chi. A stonemason does 9 slabs per gong; carrying 5 slabs on poles uses the labour of 4 men; carrying the lime is employed in addition, at 110 jin per gong. Each cook is in charge of 60 men; each supervisor (buya) is in charge of 120 men. The stone used must keep to its original pieces; it is not permitted to chisel it. One wishes to know: how much are the lower width of the dam, the stone and the lime used, and the total labour, each?\end{uniten}

\begin{unitzh}\unitlabel{shushu-052.10}答曰：壩下闊五丈，石版一十萬八千片，石灰一百萬八千觔，用夫一十萬三千五百二十八工十一分工之八。\end{unitzh}
\begin{uniten}Answer: the lower width of the dam is 5 zhang; stone slabs, 108000 pieces; lime, 1008000 jin; labourers used, 103528 8/11 gong.\end{uniten}

\begin{unitzh}\unitlabel{shushu-052.11}術曰：以啇功求之，招法入之。置層髙尺數，乗靣闊及長，為初率；次以差闊尺數乗髙，又乗長，為次率；却以石版長、闊、厚相乗為法，以除二率，各得石版，為上積及次積。置深，以層髙尺數約之，得層數，對二積列之一行，各添一撥天地數，各以累乗對約之，得乗率，以對上次積，併之為石版。以每觔用灰乗石為灰數。以匠工片數約版，得石匠。以搬夫數乗石版為實，以扛片數為法，除之，得人數。以搬用灰數除灰，得人數。併諸工，以火頭管數約之，為火頭；半之，為部。\end{unitzh}
\begin{uniten}Method: Seek it by Assessing Works, bringing in the recruitment method (zhaofa). Set up the layer height in chi, multiply it by the top width and by the length, as the initial rate (chu lü); next multiply the height by the width difference in chi, and again by the length, as the next rate (ci lü); then multiply together the length, width and thickness of a stone slab as the divisor, and divide the two rates by it, each giving stone slabs, as the upper volume (shang ji) and the next volume (ci ji). Set up the depth, divide (yue) it by the layer height in chi, obtaining the number of layers; set it out in one row against the two volumes; to each add one set of heaven-and-earth numbers (tiandi shu); multiply each cumulatively (leicheng) and divide correspondingly, obtaining the multiplying rates (cheng lü); use them against the corresponding upper and next volumes, and add these together as the stone slabs. Multiply the stone by the lime used per jin, as the amount of lime. [Translator's note: 每觔 "per jin" presumably stands for 每片 "per slab"; cf. the Working.] Divide the slabs by the number of slabs per mason's gong, obtaining the stonemasons. Multiply the stone slabs by the number of carrying labourers as the dividend; take the number of slabs carried on poles as the divisor; divide, obtaining the number of men. Divide the lime by the amount of lime carried, obtaining the number of men. Add all the labour; divide it by the number a cook is in charge of, as the cooks; halve this, as the supervisors (bu).\end{uniten}

\begin{unitzh}\unitlabel{shushu-052.12}草曰：置層髙二尺、靣闊三丈，相乘，得六十尺；又乘長三十，又得一萬八千尺，為初率。次以差闊一尺乘髙二尺，又乘長三十丈，得六百尺，為次率。却以石版長五尺、闊二尺、厚五寸相乘，得五尺為法，以除二率，得三千六百片為初積，一百二十片為次積，列右行。置深四丈二尺，以每層二尺約之，得二十一層，乘初積三千六百片，得七萬五千六百片于上。次置二十一層，減一，餘二十，以乘二十一層，得四百二十，半之，得二百一十，乗次積一百二十觔，得二萬五千二百片，加入上，共得一十萬八百片，為石版數。次置二十一層，減去一，餘二十，以差闊一尺乗之，得二丈，併上闊三丈，共得五丈，為下闊之數。\end{unitzh}
\begin{uniten}Working: Set up the layer height, 2 chi, and the top width, 3 zhang; multiply them together, obtaining 60 chi; multiply again by the length, 30, again obtaining 18000 chi, as the initial rate. Next multiply the height, 2 chi, by the width difference, 1 chi, and again by the length, 30 zhang, obtaining 600 chi, as the next rate. Then multiply together the stone slab's length 5 chi, width 2 chi and thickness 5 cun, obtaining 5 chi as the divisor; divide the two rates by it, obtaining 3600 slabs as the initial volume (chu ji) and 120 slabs as the next volume; set them out in the right column. Set up the depth, 4 zhang 2 chi; divide it by the 2 chi of each layer, obtaining 21 layers; multiply the initial volume, 3600 slabs, by it, obtaining 75600 slabs in the upper position. Next set up the 21 layers, subtract 1, leaving 20; multiply the 21 layers by it, obtaining 420; halve it, obtaining 210; multiply the next volume, 120 jin, by it, obtaining 25200 slabs; add this into the upper position, obtaining in all 100800 slabs, as the number of stone slabs. Next set up the 21 layers, subtract 1, leaving 20; multiply it by the width difference, 1 chi, obtaining 2 zhang; add the upper width, 3 zhang, obtaining in all 5 zhang, as the lower width.\end{uniten}

\begin{unitzh}\unitlabel{shushu-052.13}又置石一十萬八百片，以每工九約，得一萬一千二百。置石版數，以搬扛四人乗之，得四十萬三千二百為實，以五片為法，除之，得八萬六百四十工。又置石版數，以每片用灰一十觔乗之，得一百萬八千觔為灰實，以毎人擔用一百一十觔約之，得九千一百六十三工一十一分工之七，為灰工，于上。又併石版工一萬一千二百，及併搬扛工八萬六百四十工，加上，共得一十萬一千三工一十一分工之七。\end{unitzh}
\begin{uniten}Again set up the stone, 100800 slabs; divide by 9 per gong, obtaining 11200. Set up the number of stone slabs; multiply it by the 4 men for carrying on poles, obtaining 403200 as the dividend; take 5 slabs as the divisor; divide, obtaining 80640 gong. Again set up the number of stone slabs; multiply it by the 10 jin of lime used per slab, obtaining 1008000 jin as the lime dividend; divide it by the 110 jin each man carries on a shoulder-pole, obtaining 9163 7/11 gong, as the lime labour, in the upper position. Then add the stone-slab labour, 11200, and also add the carrying labour, 80640 gong, into the upper position: in all 101003 7/11 gong.\end{uniten}

\begin{unitzh}\unitlabel{shushu-052.14}通分内子，得一百一十一萬一千四十為實。置火頭每名管六十工，以乗分母一十一，得六百六十為法，除實，得一千六百八十三工，不盡二百六十，與法約之，為三十三分工之一十三，為火頭工數。半之，得八百四十一工三十三分工之二十三，為部押壕寨數。今衆工下十一分工之七，以母十一除火頭分母三十三，得三，以三因衆工下母子，為一十萬一千三工三十分工之二十一。併三項母子，得一十萬三千五百二十九工，分子五十七，滿母三十三，收一工，餘二十四，與母各三約，為十一分工之八，為共用一十萬三千五百二十八工十一分工之八。\end{unitzh}
\begin{uniten}Convert to parts and add in the numerator, obtaining 1111040 as the dividend. Set up the 60 gong that each cook is in charge of; multiply it by the denominator, 11, obtaining 660 as the divisor; divide the dividend, obtaining 1683 gong, with a remainder of 260; reducing this with the divisor gives 13/33 gong; this is the number of gong of the cooks. Halve it, obtaining 841 23/33 gong, as the number of the supervising moat-and-stockade officers (buya haozhai). Now, under the total labour there is 7/11 gong; divide the cooks' denominator, 33, by the denominator 11, obtaining 3; multiply (yin) by 3 the denominator and numerator under the total labour, giving 101003 21/30 gong. Add the denominators and numerators of the three items, obtaining 103529 gong, with numerators 57; filling the denominator 33, collect 1 gong; the remaining 24 and the denominator are each reduced by 3, giving 8/11 gong; so in all 103528 8/11 gong are used.\end{uniten}

\begin{unitzh}\unitlabel{shushu-052.15}<子部,天文算法類,算書之屬,數學九章,卷七下> <子部,天文算法類,算書之屬,數學九章,卷七下> <子部,天文算法類,算書之屬,數學九章,卷七下>\end{unitzh}
\begin{uniten}‹Masters (zi) division, Astronomy and Mathematics category, Mathematical Books subcategory, Shuxue jiuzhang, juan 7, part 2› ‹Masters (zi) division, Astronomy and Mathematics category, Mathematical Books subcategory, Shuxue jiuzhang, juan 7, part 2› ‹Masters (zi) division, Astronomy and Mathematics category, Mathematical Books subcategory, Shuxue jiuzhang, juan 7, part 2›\end{uniten}

\begin{unitzh}\unitlabel{shushu-052.16}築埂均功\end{unitzh}
\begin{uniten}Building an Embankment (geng) with Equalized Labour\end{uniten}

\begin{unitzh}\unitlabel{shushu-053.1}問曰：縣共興築圩埂，長三十六里半。甲縣出二千七百八十人，乙縣出一千九百九十人，丙縣出一千六百三十人，丁縣出一千三百二十人。其甲縣先差到一千五百四十四，丙縣先差到九百六十五夫。欲知各合賦役埂長幾何？{\zhsmall 里法三百六十歩。}築曰：甲先到人築二千六百二十八歩{\zhsmall 計七里一百八歩}；丙先到人築一千六百四十二歩半{\zhsmall 計四里二百二歩半}。\end{unitzh}
\begin{uniten}Question: It says the counties jointly undertake to build a polder dike (weigeng) 36 1/2 li (里) long. [Translator's note: 曰 'says' is probably an error for 四 'four'.] County A provides 2780 men, County B provides 1990 men, County C provides 1630 men, and County D provides 1320 men. Of these, County A has dispatched 1544 who arrived first, and County C has dispatched 965 laborers (fu) who arrived first. One wishes to know what length of dike each ought to be assigned as corvée. [Note: The li-rate (lifa) is 360 bu (歩).] Answer: [Translator's note: 築 'build' stands here for 答 'answer'.] The first-arrived men of A build 2628 bu [Note: amounting to 7 li 108 bu]; the first-arrived men of C build 1642 1/2 bu [Note: amounting to 4 li 202 1/2 bu].\end{uniten}

\begin{unitzh}\unitlabel{shushu-053.2}術曰：以啇功求之。置里通歩作尺，為積率；併諸縣人數，為均法。法與率可約者約之。以科率各乗先到人為實，皆法而一，各得先築里歩，為先賦埂長。續到人合賦功，凖此求之。\end{unitzh}
\begin{uniten}Method: Seek it by 'assessing works' (shanggong). Set up the li, convert them into bu and make them into chi (尺), as the total rate (jilü); sum the numbers of men of all the counties, as the equalizing divisor (junfa). Where the divisor and the rate can be reduced, reduce them. With the levy rate (kelü) multiply the first-arrived men of each, as the dividends (shi); divide all by the divisor (fa); each gives the li and bu built first, as the length of dike assigned first. The work that the men arriving later ought to be assigned is sought according to this.\end{uniten}

\begin{unitzh}\unitlabel{shushu-053.3}草曰：置三十六里五分，以里法三百六十歩通之，得一萬三千一百四十歩；又以歩五尺乗之，得六萬五千七百尺，為長率。併甲、乙、丙、丁縣四縣合科人，得七千七百二十，為均法。合法與率，可求等，得二十，以約之，率得三千二百八十五，法得三百八十六。以長率三千二百八十五尺乗甲縣先到人一千五百四十四夫，得五百七萬二千四十人，為甲實；實如法三百八十六而一，得一萬三千一百四十尺；以歩五尺約之，得二千六百二十八歩，為甲縣先到人所築積歩。又以長率三千二百八十五乗丙縣先到人九百六十五夫，得三百一十七萬二十五尺，為丙實；實如法三百八十六而一，得八千二百一十二尺五寸；以歩法五尺約之，得一千六百四十二歩半，為丙縣先到人所築歩。\end{unitzh}
\begin{uniten}Working: Set up 36 li and 5 fen (tenths); convert it by the li-rate of 360 bu, giving 13,140 bu; again multiply it by 5 chi per bu, giving 65,700 chi, as the length rate (changlü). Sum the men due for levy of the four counties A, B, C and D, giving 7720, as the equalizing divisor. Combining the divisor and the rate, their equal (deng) can be sought, giving 20; reducing them by it, the rate gives 3285 and the divisor gives 386. With the length rate 3285 chi multiply the first-arrived men of County A, 1544 laborers, giving 5,072,040 men, as A's dividend; dividing the dividend by the divisor 386 gives 13,140 chi; reducing it by 5 chi per bu gives 2628 bu, as the total bu built by the first-arrived men of County A. Again, with the length rate 3285 multiply the first-arrived men of County C, 965 laborers, giving 3,170,025 chi, as C's dividend; dividing the dividend by the divisor 386 gives 8212 chi 5 cun (寸); reducing it by the bu-rate of 5 chi gives 1642 1/2 bu, as the bu built by the first-arrived men of County C.\end{uniten}

\begin{unitzh}\unitlabel{shushu-053.4}計浚河渠\end{unitzh}
\begin{uniten}Calculating the Dredging of a Canal\end{uniten}

\begin{unitzh}\unitlabel{shushu-053.5}問𨳩通運河，就土築堤。令面廣六丈，底廣四丈，上流深八尺，下流深一丈六尺，長四十八里。其堤下廣二丈四尺，上廣一丈八尺，與河等，未知髙。以墟四堅三為率。秋程入功，每名自𨳩、運、築，墟堅共積常六十尺。築堤至半為棚道，取土上下，工減五分之一。限一月畢。欲知河積及堤積尺、共用工并日役工數及堤髙各幾何？答曰：河積六千二百二十萬八千尺；堤積四千六百六十五萬六千尺；堤髙二丈一尺七分尺之三；共用工二十四萬四千九百四十四；日役工八千一百六十四五分工之四。\end{unitzh}
\begin{uniten}Question: A transport canal (yunhe) is opened, and with its earth an embankment is built. Let the surface width be 6 zhang (丈), the bottom width 4 zhang, the upstream depth 8 chi, the downstream depth 1 zhang 6 chi, and the length 48 li. The lower width of the embankment is 2 zhang 4 chi and its upper width 1 zhang 8 chi, equal to the canal; its height is not known. Take loose earth (xu) 4 and compacted earth (jian) 3 as the rates (lü). By the autumn quota (qiucheng) of work put in, each man's own digging, carrying and building of loose and compacted earth together amounts to a constant (chang) volume of 60 chi. When the building of the embankment reaches half, a scaffold way (pengdao) is made; in fetching earth up and down, the labor is reduced by 1/5. The limit for finishing is one month. One wishes to know the volume of the canal and the volume of the embankment in chi, the total labor (gong) used together with the amount of labor employed per day, and the height of the embankment: how much is each? Answer: Volume of the canal 62,208,000 chi; volume of the embankment 46,656,000 chi; height of the embankment 2 zhang 1 3/7 chi; total labor used 244,944; labor employed per day 8164 4/5 gong.\end{uniten}

\begin{unitzh}\unitlabel{shushu-053.6}按：共用工、日用工二數俱誤。若以共工分工積，則每工各得四百四十四尺餘，其不合明矣。辨詳草後。\end{unitzh}
\begin{uniten}Note: The two numbers, the total labor used and the labor used per day, are both wrong. If the total labor is used to divide the volume of work, then each gong gets 444-odd chi; that they do not agree is clear. The discussion is detailed after the working.\end{uniten}

\begin{unitzh}\unitlabel{shushu-053.7}術曰：以啇功求之。併河上下廣于上，併河上下流深乗之，又以長乗為實，以四為法除，得河積。以堅率乗河積為實，以墟率為法除，得堤積。併堤上下廣，乗堤長，半之為法，除堤積，得堤髙。併河堤二積，以棚道母半之，副置；以棚道減工子乗之，以棚道減工母除之，得數以併其副，共為寄。以子減母，餘乗常尺，為増子；以母乗常尺，為増分。併増分、増子，乗寄，倍為用工實。以増分乗増子，又乗限月日，為法，除實，得用人工數。\end{unitzh}
\begin{uniten}Method: Seek it by 'assessing works' (shanggong). Sum the upper and lower widths of the canal, placed at the top; sum the upstream and downstream depths of the canal and multiply it; again multiply by the length, as the dividend (shi); with 4 as the divisor (fa) divide, giving the volume of the canal. With the compacted rate multiply the volume of the canal, as the dividend; with the loose rate as the divisor divide, giving the volume of the embankment. Sum the upper and lower widths of the embankment, multiply by the length of the embankment, and halve it, as the divisor; divide the volume of the embankment, giving the height of the embankment. Sum the two volumes of canal and embankment, halve it by the denominator of the scaffold way, and set up a copy; multiply it by the numerator of the scaffold-way labor reduction and divide it by the denominator of the scaffold-way labor reduction; add the number obtained to its copy, together as the lodged quantity (ji). Subtract the numerator from the denominator; with the remainder multiply the constant chi (changchi), as the augmented numerator (zengzi); with the denominator multiply the constant chi, as the augmented parts (zengfen). Sum the augmented parts and the augmented numerator, multiply the lodged quantity, and double it, as the dividend for the labor used. With the augmented parts multiply the augmented numerator, and further multiply by the days of the month limit, as the divisor; dividing the dividend gives the number of man-labor employed.\end{uniten}

\begin{unitzh}\unitlabel{shushu-053.8}草曰：置河上廣六丈，併底廣四丈，通之，折半，得五十尺，于上。又置河上流深八尺，併下流深一丈六尺，併之，折半，得一十二尺，以乗上數，得六百尺，為次。置長四十八里，以尺里法二千一百六十通之，得一十萬三千六百八十尺，得堤河長；以乗次，得六千二百二十萬八千尺，為河積。以堅率三因河積，得一億八千六百六十二萬四千尺，為實；以穿率四為法除之，得四千六百六十五萬六千尺，為堤積。置上廣一丈八尺、下廣二丈四尺，併之，為四十二尺，以乗堤長一十萬三千六百八十尺，得四百三十五萬四千五百六十尺，以半之，得二百一十七萬七千二百八十尺，為法；除堤積，得二十一尺，為堤髙；不盡九十三萬三千一百二十，與法求等，得三十一萬一千四十，俱約之，為七分尺之三。\end{unitzh}
\begin{uniten}Working: Set up the upper width of the canal, 6 zhang, add the bottom width, 4 zhang, convert, and halve, giving 50 chi, at the top. Again set up the upstream depth of the canal, 8 chi, and add the downstream depth, 1 zhang 6 chi; adding them and halving gives 12 chi; with it multiply the number at the top, giving 600 chi, as the next (ci). Set up the length, 48 li; convert it by the chi-li rate of 2160, giving 103,680 chi, which gives the length of the embankment and canal; with it multiply the next, giving 62,208,000 chi, as the volume of the canal. With the compacted rate 3 multiply (yin) the volume of the canal, giving 186,624,000 chi, as the dividend; with the excavated rate (chuanlü) 4 as the divisor divide it, giving 46,656,000 chi, as the volume of the embankment. Set up the upper width, 1 zhang 8 chi, and the lower width, 2 zhang 4 chi, and add them, making 42 chi; with it multiply the length of the embankment, 103,680 chi, giving 4,354,560 chi; halve it, giving 2,177,280 chi, as the divisor; divide the volume of the embankment, giving 21 chi, as the height of the embankment; the remainder, 933,120, with the divisor seek the equal (deng), giving 311,040; reducing both by it gives 3/7 chi.\end{uniten}

\begin{unitzh}\unitlabel{shushu-053.9}次置河積六千二百二十萬八十尺，併堤積四千六百六十五萬六千尺，得一億八百八十六萬四千尺。以棚道築至半是二除之，得五千四百四十三萬二千尺，副之。先以減工之子一{\zhsmall 按此以下法皆誤}乗之，只得此數，為實；乃後以減工母五為法除之，得一千八十八萬六千四百尺；併副五千四百四十三萬二千尺，共得六千五百三十一萬八千四百尺，為寄。\end{unitzh}
\begin{uniten}Next set up the volume of the canal, 62,200,080 chi, and add the volume of the embankment, 46,656,000 chi, giving 108,864,000 chi. Because the scaffold way is built at half, which is 2, divide it by this, giving 54,432,000 chi, and make a copy of it. First, with the numerator of the labor reduction, 1, [Siku editors: From here on the methods are all wrong.] multiply it, which only gives this number, as the dividend; then afterwards, with the denominator of the labor reduction, 5, as the divisor, divide it, giving 10,886,400 chi; add the copy, 54,432,000 chi, together giving 65,318,400 chi, as the lodged quantity.\end{uniten}

\begin{unitzh}\unitlabel{shushu-053.10}以折減工五分之一，以子一減母五，餘四，以乗常尺六十，得二百四十尺，為増子；以母五乗常尺六十，得三百尺，為増分。以二増併之，得五百四十，乗寄，得三百五十二億七千一百九十三萬六千尺；以半之，得一百七十億三千五百九十六萬八千尺，為用工實{\zhsmall 按此下又有脱誤}，得八千一百六十四，為每日人工數；不盡一百七十二萬八千，與法求等四十三萬二千，俱約之，為五分工之四。得每日用工八千一百六十四工五分工之四。復通分内子，得四萬七百二十四，以三十日乗之，得一百二十二萬一千七百二十，為寔；仍以母五約之，得二十四萬四千三百四十四工，為共用工。合問。\end{unitzh}
\begin{uniten}Taking the labor reduction of 1/5: subtract the numerator 1 from the denominator 5, leaving 4; with it multiply the constant chi, 60, giving 240 chi, as the augmented numerator; with the denominator 5 multiply the constant chi, 60, giving 300 chi, as the augmented parts. Add the two augmented quantities, giving 540; multiply the lodged quantity, giving 35,271,936,000 chi; halve it, giving 17,035,968,000 chi, as the dividend for the labor used [Siku editors: Below this there are again omissions and errors.], giving 8164, as the number of man-labor per day; the remainder, 1,728,000, with the divisor seek the equal, 432,000; reducing both gives 4/5 gong. This gives the labor used per day, 8164 4/5 gong. Again, converting into parts and adding in the numerator gives 40,724; multiplying it by 30 days gives 1,221,720, as the dividend; then reducing it by the denominator 5 gives 244,344 gong, as the total labor used. This agrees with the question.\end{uniten}

\begin{unitzh}\unitlabel{shushu-053.11}按：草中求堤積，至宻至㨗，誠數家之要法也。至減工子母乗除而下，則法與數皆有誤焉。葢題言棚道減工五分之一，是棚道為平道四分之五也。四為分母，五為分子。應以分子五乗上下積五千四百四十三萬二千尺，得二億七千二百一十六萬尺；以分母四乗副半積，得二億一千七百七十二萬八千尺；併之，得四億八千九百八十八萬八千尺，為實。以分母四乗常尺六十，得二百四十尺，為法，除實，得二百零四萬一千二百工，為共工數。以一月三十日除之，得六萬八千零四十工，為每日工數。\end{unitzh}
\begin{uniten}Note: In the working, the finding of the volume of the embankment is extremely precise and extremely quick; it is truly an essential method of the mathematicians. From the multiplication and division by the numerator and denominator of the labor reduction downward, both the method and the numbers have errors. For the problem says that the scaffold way reduces labor by 1/5; this means that the scaffold way is 5/4 of the level way (pingdao). 4 is the denominator, 5 is the numerator. One should, with the numerator 5, multiply the upper-and-lower volume, 54,432,000 chi, giving 272,160,000 chi; with the denominator 4 multiply the copied half volume, giving 217,728,000 chi; add them, giving 489,888,000 chi, as the dividend. With the denominator 4 multiply the constant chi, 60, giving 240 chi, as the divisor; dividing the dividend gives 2,041,200 gong, as the total labor. Dividing it by one month of 30 days gives 68,040 gong, as the labor per day.\end{uniten}

\begin{unitzh}\unitlabel{shushu-053.12}或置分子五乗上半積之得數于上，又併分子母得九，乗常尺六十，得五百四十尺，乗上數，得一千四百六十九億六千六百四十萬尺，為實。以分母四乗六十尺，得二百四十尺，為增母；以分子五乗六十尺，得三百尺，為増子；増母子相乗，得七萬二千尺，為法，除實，得共工數，亦與前同。此特不用副半積數，然不若前法之省。草中以五為分母，以一為分子，母子既以顛倒，而又以餘分為分子；後雖易一為四，而母子之名未正，故其中累乗累除之數漫無可據，而所差甚逺也。\end{unitzh}
\begin{uniten}Or: set up at the top the number obtained by multiplying the upper half volume by the numerator 5; again add the numerator and denominator, giving 9; multiply the constant chi, 60, giving 540 chi; multiply the number at the top, giving 146,966,400,000 chi, as the dividend. With the denominator 4 multiply 60 chi, giving 240 chi, as the augmented denominator (zengmu); with the numerator 5 multiply 60 chi, giving 300 chi, as the augmented numerator; multiply the augmented denominator and numerator together, giving 72,000 chi, as the divisor; dividing the dividend gives the total labor, again the same as before. This merely does not use the copied half-volume number, but it is not as economical as the former method. In the working, 5 is taken as the denominator and 1 as the numerator: the denominator and numerator having thus been inverted, moreover the remaining parts were taken as the numerator; and although afterwards 1 was changed into 4, the names of denominator and numerator were still not set right. Therefore the numbers of the repeated multiplications and repeated divisions in it are vague and have nothing to rely on, and the discrepancy is very great.\end{uniten}

\begin{unitzh}\unitlabel{shushu-053.13}圍田先計\end{unitzh}
\begin{uniten}Advance Calculation for a Polder Field (weitian)\end{uniten}

\begin{unitzh}\unitlabel{shushu-054.1}問：有草蕩一所，廣三里，縱一百一十八里，夏日水深二尺五寸，與溪靣等平。溪闊一十三丈，流長一百三十五里入湖。冬日水深一尺，欲趁此時圍𥚃成田。于蕩中順縱𨳩大港一條，磬折通溪；順廣𨳩小港二十四條，其深同。其小港闊比大港六分之一，大港深比大港靣三分之一，大小港底各不及靣一尺。取土為埂，髙一丈，上廣六尺，下廣一丈二尺。蕩縱當溪，其岸髙廣倍其埂數。上下流各立斗門一所，湏令田内止容水八寸，遏餘水復溪入湖。里法三百六十歩，歩法五尺。欲知田積、埂土積、大小港底靣深闊、冬夏積水、田港容水、過水、溪靣泛髙幾何？\end{unitzh}
\begin{uniten}Question: There is a grassy marsh (cao dang), one tract, 3 li (里) wide and 118 li long. In summer its water is 2 chi (尺) 5 cun (寸) deep, level with the surface of a creek. The creek is 13 zhang (丈) wide and flows a length of 135 li into a lake. In winter the water is 1 chi deep, and one wishes to take advantage of this time to enclose it and make fields. In the marsh, along the length, open one large channel (gang), bent like a chime-stone (qingzhe) to connect with the creek; along the width, open 24 small channels; their depths are the same. The width of the small channels, compared with the large channel, is 1/6; the depth of the large channel, compared with the surface of the large channel, is 1/3; the bottoms of the large and small channels each fall short of the surface by 1 chi. Take the earth to make embankments (geng), 1 zhang high, 6 chi wide at the top and 1 zhang 2 chi wide at the bottom. Where the length of the marsh faces the creek, the height and width of its bank are double the figures of the embankment. At the upper and lower reaches set up one sluice gate (doumen) each; it is required that only 8 cun of water be held within the fields, and the remaining water is shut out, returning to the creek and entering the lake. The li standard (lifa) is 360 bu (歩); the bu standard (bufa) is 5 chi. One wishes to know the area of the fields; the volume of embankment earth; the bottoms, surfaces, depths and widths of the large and small channels; the water accumulated in winter and in summer; the water held by fields and channels; the passing water (guo shui); and the rise (fan gao) of the creek surface: how much are they?\end{uniten}

\begin{unitzh}\unitlabel{shushu-054.2}按：題意掘土為港，即以其土四邉為埂，當溪者髙闊倍之，餘三邉等，語皆未詳。\end{unitzh}
\begin{uniten}Note: By the intent of the problem, earth is dug to make the channels, and that earth is then used to make embankments on the four sides, the one facing the creek being doubled in height and width and the remaining three sides equal; but the wording makes none of this explicit.\end{uniten}

\begin{unitzh}\unitlabel{shushu-054.3}答曰：田積一千八百六十六頃八畆二十四歩。埂土積九百六十五億五千二百萬立方寸。大港靣闊六丈一尺七寸，底闊六丈七寸，深六尺八寸。小港靣闊一丈二寸六分寸之五，底闊九尺二寸六分寸之五，深六尺八寸。夏積水二萬八千六百七十四億立方寸。冬積水一萬一千四百六十九億六千萬立方寸。田容水九千七十二億六千九百一十二萬立方寸。港容水九百六十五億五千二百萬立方寸{\zhsmall 港上者在田内}。遏出水一萬八千六百三十五億七千八百八十八萬立方寸。溪靣泛髙一尺三寸一十三萬一千六百二十五分寸之一萬四千四百一十一。\end{unitzh}
\begin{uniten}Answer: Area of the fields, 1866 qing (頃) 8 mu (畆) 24 bu. Volume of embankment earth, 96,552,000,000 cubic cun (lifang cun). Large channel: surface width 6 zhang 1 chi 7 cun, bottom width 6 zhang 7 cun, depth 6 chi 8 cun. Small channel: surface width 1 zhang 2 5/6 cun, bottom width 9 chi 2 5/6 cun, depth 6 chi 8 cun. Summer accumulated water, 2,867,400,000,000 cubic cun. Winter accumulated water, 1,146,960,000,000 cubic cun. Water held by the fields, 907,269,120,000 cubic cun. Water held by the channels, 96,552,000,000 cubic cun [Note: what is above the channels lies within the fields]. Water shut out (e chu shui), 1,863,578,880,000 cubic cun. Rise of the creek surface, 1 chi 3 14411/131625 cun.\end{uniten}

\begin{unitzh}\unitlabel{shushu-054.4}按：題言大港深比大港靣三分之一，答數中大港靣六丈一尺七寸，深六尺八寸，是九之一不足，而非三之一矣。此數與題不合，以下俱誤。\end{unitzh}
\begin{uniten}Note: The problem says that the depth of the large channel, compared with the surface of the large channel, is 1/3; in the answer, the surface of the large channel is 6 zhang 1 chi 7 cun and the depth 6 chi 8 cun, which falls short of 1/9 and is not 1/3. These figures do not agree with the problem, and everything below is wrong.\end{uniten}

\begin{unitzh}\unitlabel{shushu-054.5}子部，天文算法類，算書之屬，《數學九章》，卷七下\end{unitzh}
\begin{uniten}Masters Section (zibu), Astronomy and Mathematics Category (tianwen suanfa lei), Subcategory of Mathematical Books (suanshu zhi shu), Shuxue jiuzhang, juan 7, part 2. [Translator's note: Siku quanshu classification header.]\end{uniten}

\begin{unitzh}\unitlabel{shushu-054.6}草曰：先通歩法為五十寸，通三百六十歩，得一萬八千寸，為里法。以里法通蕩廣三里，得五萬四千，為廣率。又通蕩縱一百一十八里，得二百一十二萬四千，為縱率。以縱併廣率，得二百一十七萬八千，為和。以縱率乗廣率，得一千一百四十六億九千六百萬，為寄。三因縱率二百一十二萬四千，得六百三十七萬二千于上；倍和二百一十七萬八千，得四百三十五萬六千，加上，得一千七十二萬八千，為叚。次以埂上廣六尺併下廣一丈二尺，得一十八尺，乗半埂髙五十寸，得九千寸，又乗叚一千七十二萬八千，得九百六十五億五千二百萬，為土積，亦為港容水。\end{unitzh}
\begin{uniten}Working: First convert the bu standard into 50 cun; with it convert 360 bu, obtaining 18,000 cun as the li standard. With the li standard convert the width of the marsh, 3 li, obtaining 54,000 as the width rate (guang lü). Again convert the length of the marsh, 118 li, obtaining 2,124,000 as the length rate (zong lü). Add the length to the width rate, obtaining 2,178,000 as the sum (he). Multiply the width rate by the length rate, obtaining 114,696,000,000 as the reserve (ji). Multiply the length rate 2,124,000 by 3, obtaining 6,372,000 in the upper position; double the sum 2,178,000, obtaining 4,356,000; add it to the upper position, obtaining 10,728,000 as the segment (duan). Next add the upper width of the embankment, 6 chi, to the lower width, 1 zhang 2 chi, obtaining 18 chi; multiply by half the height of the embankment, 50 cun, obtaining 9000 cun; again multiply by the segment 10,728,000, obtaining 96,552,000,000 as the volume of earth, which is also the water held by the channels.\end{uniten}

\begin{unitzh}\unitlabel{shushu-054.7}以港闊母六因土積，得五千七百九十三億一千二百萬，為實。以闊子一乗小港二十四條，又乗廣率五萬四千，得一百二十九萬六千，為泛。以闊母六因縱率二百一十二萬四千，得一千二百七十四萬四千，併泛，得一千四百四萬，為隅。𨳩平方{\zhsmall 按：大港闊六分差一尺，小港闊一分差一尺，故此法不免有差}，得二百三寸，不盡七百三十七萬六千四百，收為所得一寸，乃得二百四寸，為堢。以深子一乗之，以深母三除，得六尺八寸，為大小港等深。次以深母三因堢二百四寸，得六百一十二寸，為實，如深子一而一，得六丈一尺二寸，為中。以不及一尺半之，得五寸，加中，得六丈一尺七寸{\zhsmall 按：如此則不成分數矣}，為大港靣闊。如母六而一，得一丈二寸六分寸之五，為小港靣。以不及一尺各減大小港靣，得六丈七寸，為大港底；得九尺二寸六分寸之五，為小港底。\end{unitzh}
\begin{uniten}Multiply the volume of earth by the channel-width denominator 6, obtaining 579,312,000,000 as the dividend (shi). Multiply the 24 small channels by the width numerator 1, and again multiply by the width rate 54,000, obtaining 1,296,000 as the provisional (fan). Multiply the length rate 2,124,000 by the width denominator 6, obtaining 12,744,000; add the provisional, obtaining 14,040,000 as the corner (yu). Extract the square root [Siku editors: The width of the large channel is 6 parts with a difference of 1 chi, and the width of the small channel is 1 part with a difference of 1 chi; therefore this method cannot avoid having an error], obtaining 203 cun, with a remainder of 7,376,400; take it in as 1 cun of the result, thus obtaining 204 cun as the bulwark (bao). Multiply it by the depth numerator 1 and divide by the depth denominator 3, obtaining 6 chi 8 cun as the equal depth of the large and small channels. Next multiply the bulwark 204 cun by the depth denominator 3, obtaining 612 cun as the dividend; divide by the depth numerator 1, obtaining 6 zhang 1 chi 2 cun as the middle (zhong). Halve the shortfall of 1 chi, obtaining 5 cun; add it to the middle, obtaining 6 zhang 1 chi 7 cun [Siku editors: If it is done like this, it does not make the fraction] as the surface width of the large channel. Divide by the denominator 6, obtaining 1 zhang 2 5/6 cun as the surface of the small channel. Subtract the shortfall of 1 chi from each of the surfaces of the large and small channels, obtaining 6 zhang 7 cun as the bottom of the large channel, and obtaining 9 chi 2 5/6 cun as the bottom of the small channel.\end{uniten}

\begin{unitzh}\unitlabel{shushu-054.8}次以埂下廣一丈二尺乗叚一千七十二萬八千寸，得一十二億八千七百三十六萬，為址。以大港靣六丈一尺七寸乗隅一千四百四萬，得八十六億六千二百六十八萬，為實；以闊母六除之，得一十四億四千三百七十八萬，為港平。以併址一千二億八千七百三十六萬，得二十七億三千一百一十四萬，減寄一千一百四十六億九千六百萬，餘一千一百一十九億六千四百八十六萬，為田積寸。以歩法五十寸自乗，得二千五百，除積寸，得四千四百七十八萬五千九百四十四歩，為田積歩。以畆法二百四十歩約之，得一千八百六十六頃八畆，不盡二十四歩，為田積。\end{unitzh}
\begin{uniten}Next multiply the segment 10,728,000 cun by the lower width of the embankment, 1 zhang 2 chi, obtaining 1,287,360,000 as the footing (zhi). Multiply the corner 14,040,000 by the surface of the large channel, 6 zhang 1 chi 7 cun, obtaining 8,662,680,000 as the dividend; divide it by the width denominator 6, obtaining 1,443,780,000 as the channel plane (gangping). Add to it the footing 100,287,360,000 [Translator's note: so in the text; the footing above is 1,287,360,000], obtaining 2,731,140,000; subtract it from the reserve 114,696,000,000, leaving 111,964,860,000 as the field area in cun. Multiply the bu standard 50 cun by itself, obtaining 2500; with it divide the area in cun, obtaining 44,785,944 bu as the field area in bu. Reduce it by the mu standard (mufa) 240 bu, obtaining 1866 qing 8 mu, with 24 bu remaining, as the area of the fields.\end{uniten}

\begin{unitzh}\unitlabel{shushu-054.9}以址一十二億八千七百三十六萬減寄一千一百四十六億九千六百萬，餘一千一百三十四億八百六十四萬，乗令容水八寸，得九千七十二億六千九百一十二萬，為田容水。次以夏水深二尺五寸乗寄一千一百四十六億九千六百萬，得二萬八千六百七十四億寸，為夏積水。次以冬水深一尺乗寄，得一萬一千四百六十九億六千萬寸，為冬積水。乃以田容水九千七十二億六千九百一十二萬併港容水九百六十五億五千二百萬，得一萬三十八億二千一百一十二萬，減夏積水二萬八千六百七十四億寸，餘一萬八千六百三十五億七千八百八十八萬，為遏出水。\end{unitzh}
\begin{uniten}Subtract the footing 1,287,360,000 from the reserve 114,696,000,000, leaving 113,408,640,000; multiply by the prescribed held water of 8 cun, obtaining 907,269,120,000 as the water held by the fields. Next multiply the reserve 114,696,000,000 by the summer water depth 2 chi 5 cun, obtaining 2,867,400,000,000 cun as the summer accumulated water. Next multiply the reserve by the winter water depth 1 chi, obtaining 1,146,960,000,000 cun as the winter accumulated water. Then add the water held by the fields, 907,269,120,000, to the water held by the channels, 96,552,000,000, obtaining 1,003,821,120,000; subtract it from the summer accumulated water 2,867,400,000,000 cun, leaving 1,863,578,880,000 as the water shut out.\end{uniten}

\begin{unitzh}\unitlabel{shushu-054.10}當以八節乗之，嵗日三百六十除之，為實；今從省，先以八節約嵗日三百六十，得四十五，為率。次以里法一萬八千寸通流長一百三十五里，得二百四十三萬，又乗溪闊一十三丈，得三十一億五千九百萬，以乗除率四十五，得一千四百二十一億五千五百萬，為法。除遏出水一萬八千六百三十五億七千八百八十八萬，得一尺三寸，為溪靣泛髙。不盡一百五十五億六千三百八十八萬，與法一千四百二十一億五千五百萬求等，得一百八萬，俱以約之，為一十三萬一千六百二十五分寸之一萬四千四百一十一，為泛髙寸下分母之數。合問。\end{unitzh}
\begin{uniten}It should be multiplied by the 8 seasonal nodes (ba jie) and divided by the 360 days of the year to make the dividend; now, for brevity, first reduce the 360 days of the year by the 8 nodes, obtaining 45 as the rate (lü). Next with the li standard 18,000 cun convert the length of the flow, 135 li, obtaining 2,430,000; again multiply by the width of the creek, 13 zhang, obtaining 3,159,000,000; with it multiply the dividing rate 45, obtaining 142,155,000,000 as the divisor (fa). Divide the water shut out, 1,863,578,880,000, obtaining 1 chi 3 cun as the rise of the creek surface. For the remainder 15,563,880,000 and the divisor 142,155,000,000 seek the equal (deng), obtaining 1,080,000; reduce both by it, giving 14411/131625 cun, as the figures of the fraction below the cun of the rise. This agrees with the question.\end{uniten}

\begin{unitzh}\unitlabel{shushu-054.11}按：草中自求大小港闊深以後，既與題問不合，且法多疎漏。今以立天元一術推明求大小港闊深之法于後。至田積、水積等，不過冪積、體積相較，初無深義，可無論也。\end{unitzh}
\begin{uniten}Note: In the working, from the finding of the widths and depths of the large and small channels onward, it does not agree with the question of the problem, and moreover the method has many careless omissions. Now the method for finding the widths and depths of the large and small channels is worked out and made clear below by the technique of setting up the celestial unknown (tianyuan). As for the area of the fields, the volumes of water and the like, they are no more than comparisons of areas (mi ji) and volumes (ti ji); they have from the outset no deep meaning and need not be discussed.\end{uniten}

\begin{unitzh}\unitlabel{shushu-054.12}法：立天元一為一分，六因之，得六元，為大港闊；減十寸，得六元少十寸，為大港底闊。併之，得十二元少十寸；以半深一元乗之，得十二平方少十元；以縱率因之，得二千五百四十八萬八千平方少二千一百二十四萬元。又以天元一為小港闊；減十寸，得一元少十分，為小港底闊。併之，得二元少十寸；以半深一元因之，得二平方少十元；以二十四廣率因之，得二百五十九萬二千平方少一千二百九十六萬元。併二數，得二千八百零八萬平方少三千四百二十萬元，與土埂共積等。而邉各以一萬除之，得二十八百零八平方少三千四百二十元，與九百六十五萬五千二百寸等。\end{unitzh}
\begin{uniten}Rule: Set up the celestial unknown as 1 part (fen); multiply it by 6, obtaining 6 yuan (元) as the width of the large channel; subtract 10 cun, obtaining 6 yuan less 10 cun as the bottom width of the large channel. Add them, obtaining 12 yuan less 10 cun; multiply by half the depth, 1 yuan, obtaining 12 squares (pingfang) less 10 yuan; multiply by the length rate, obtaining 25,488,000 squares less 21,240,000 yuan. Again take the celestial unknown 1 as the width of the small channel; subtract 10 cun, obtaining 1 yuan less 10 fen [Translator's note: so in the text; 10 cun is subtracted] as the bottom width of the small channel. Add them, obtaining 2 yuan less 10 cun; multiply by half the depth, 1 yuan, obtaining 2 squares less 10 yuan; multiply by 24 width rates, obtaining 2,592,000 squares less 12,960,000 yuan. Add the two quantities, obtaining 28,080,000 squares less 34,200,000 yuan, which is equal to the total volume of the earth embankments. Then divide each side by 10,000, obtaining 2808 squares less 3420 yuan, equal to 9,655,200 cun.\end{uniten}

\begin{unitzh}\unitlabel{shushu-054.13}三數又求等，得三十六，徧約之，得七十八平方少九十五元，與二十六萬八千二百寸等。乃以寸數為實，以元數為縱，以方數為隅，𨳩帶縱平方，得方邉為一分數，五尺九寸{\zhsmall 二五○五三二九}，即小港闊；減一尺，得四尺九寸{\zhsmall 二五○五三二九}，為小港底闊。六因一分數，得三丈五尺五寸{\zhsmall 五三○一九七三}，為大港闊；減一尺，得三丈四尺五寸{\zhsmall 五○三一九七三}，為大港底闊。二因一分數，得一丈一尺八寸{\zhsmall 五○一○六五八}，為同深。以此轉求共港容水數，乃與共土埂原積數相合。\end{unitzh}
\begin{uniten}For the three numbers again seek the equal (deng), obtaining 36; reduce them all by it, obtaining 78 squares less 95 yuan, equal to 268,200 cun. Then take the number of cun as the dividend, the number of yuan as the length (zong), and the number of squares as the corner; extract the square root with attached length (dai zong pingfang), obtaining the side of the square as the value of 1 part, 5 chi 9 cun [Note: 2505329] [Translator's note: the small-print digits give further decimal places], which is the width of the small channel; subtract 1 chi, obtaining 4 chi 9 cun [Note: 2505329] as the bottom width of the small channel. Multiply the value of 1 part by 6, obtaining 3 zhang 5 chi 5 cun [Note: 5301973] as the width of the large channel; subtract 1 chi, obtaining 3 zhang 4 chi 5 cun [Note: 5031973] as the bottom width of the large channel. Multiply the value of 1 part by 2, obtaining 1 zhang 1 chi 8 cun [Note: 5010658] as the common depth. Using these, find in turn the total amount of water held by the channels; it then agrees with the original total volume of the earth embankments.\end{uniten}

\begin{unitzh}\unitlabel{shushu-054.14}《數學九章》巻七下\end{unitzh}
\begin{uniten}Shuxue jiuzhang, juan 7, part 2.\end{uniten}

\begin{unitzh}\unitlabel{shushu-054.15}子部，天文算法類，算書之屬，《數學九章》\end{unitzh}
\begin{uniten}Masters Section (zibu), Astronomy and Mathematics Category, Subcategory of Mathematical Books, Shuxue jiuzhang.\end{uniten}

\section*{卷8上}\addcontentsline{toc}{section}{卷8上}

\begin{unitzh}\unitlabel{shushu-055.1}欽定四庫全書\end{unitzh}
\begin{uniten}Imperially Commissioned Complete Library of the Four Treasuries (Qinding Siku quanshu)\end{uniten}

\begin{unitzh}\unitlabel{shushu-055.2}《數學九章》巻八上 宋 秦九韶 撰\end{unitzh}
\begin{uniten}Mathematical Treatise in Nine Sections (Shuxue jiuzhang), juan 8, part 1. Composed by Qin Jiushao of the Song.\end{uniten}

\begin{unitzh}\unitlabel{shushu-055.3}軍旅\end{unitzh}
\begin{uniten}Military Affairs (junlü)\end{uniten}

\begin{unitzh}\unitlabel{shushu-055.4}計立方營\end{unitzh}
\begin{uniten}Reckoning the setting up of a square camp\end{uniten}

\begin{unitzh}\unitlabel{shushu-055.5}問：一軍三将，将三十三隊，隊一百二十五人。遇暮立營，人占立地方八尺，須令隊間容隊，師居中央。欲知營方幾何？答曰：方營一百七十一丈，隊方九丈。\end{unitzh}
\begin{uniten}Question: One army has 3 generals; each general has 33 companies (dui); each company has 125 men. At nightfall they set up camp; each man occupies standing ground 8 chi (尺) square. It must be arranged that the space between companies holds a company, and that the commander occupies the center. One wishes to know how large the camp square is. Answer: The square camp is 171 zhang (丈); the company square is 9 zhang.\end{uniten}

\begin{unitzh}\unitlabel{shushu-055.6}術曰：以少廣求之。置人占方冪，乗每隊人，為隊實；以一為隅，開平方，所得為隊方圖{\zhsmall 或開不盡，就為全盡}。次置隊數，乗将數，又四因之，増三，為寔；以二為從隅，開平方，得率，以乗隊方面，為營方面{\zhsmall 開不盡為全數}。\end{unitzh}
\begin{uniten}Method: Seek it by lesser breadth (shaoguang). Set up the square area (mi) occupied by one man and multiply it by the men of each company, to make the company dividend (shi); take 1 as the corner (yu) and extract the square root; what is obtained makes the company square figure (tu) [Note: if the extraction does not come out exactly, then make it complete]. Next set up the number of companies, multiply it by the number of generals, further multiply it by 4, and add 3, to make the dividend; take 2 as the following corner (congyu), extract the square root, and obtain the rate (lü); with it multiply the side of the company square, to make the side of the camp square [Note: if the extraction does not come out exactly, make it a whole number].\end{uniten}

\begin{unitzh}\unitlabel{shushu-055.7}按：舊圖各隊四眼内，每人作一小圓為識，今去之。總圖内各隊仍畫四眼，今只以一小方為一隊。舊總圖太大，難於撿閲，今收入半頁内。又按：總圖内係百隊，筭内只有九十九隊，圖中應虚一隊，舊本未詳。\end{unitzh}
\begin{uniten}Note: In the old diagram, within the four eyes of each company, each man was drawn as a small circle as a mark; now these have been removed. In the overall diagram each company was still drawn with four eyes; now only one small square is used for one company. The old overall diagram was too large and difficult to consult; now it has been fitted within half a page. Further note: the overall diagram is made up of 100 companies, while in the calculation there are only 99 companies, so in the diagram one company should be left empty; the old edition did not make this clear.\end{uniten}

\begin{unitzh}\unitlabel{shushu-055.8}草曰：置人占八尺，自乗得六十四尺，為人占方冪。以乗每隊一百二十五人，得八千尺，為寔。以一為隅，開平方。歩法常超一位，今隅超一度，至寔之百下。約寔，置啇八十尺。以啇八十生隅一，得八十，為方。乃命上啇除寔，訖，實餘一千六百。次以啇生隅入方，得一百六十。畢，方一退，隅再退之。復於上啇之次續啇九尺，乃以續啇九生隅一入方，得一百六十九。乃命續啇除寔，訖，得八十九尺，不盡七十九尺，就為九十尺，得隊方面。\end{unitzh}
\begin{uniten}Working: Set up the 8 chi occupied by a man; multiplying it by itself gives 64 chi, as the square area occupied by a man. With it multiply the 125 men of each company, obtaining 8000 chi, as the dividend. Take 1 as the corner and extract the square root. The step method (bufa) always skips (chao) one place; now the corner skips once, reaching below the hundreds of the dividend. Estimating (yue) the dividend, set up the quotient (shang) 80 chi. With the quotient 80 generate (sheng) the corner 1, obtaining 80, as the side (fang). Then call upon the upper quotient to eliminate (chu) the dividend; when this is done, the dividend has 1600 remaining. Next, with the quotient generate the corner and enter (ru) it into the side, obtaining 160. When this is finished, the side retreats (tui) once and the corner retreats twice. Again, after the upper quotient, continue the quotient with 9 chi; then with the continued quotient 9 generate the corner 1 and enter it into the side, obtaining 169. Then call upon the continued quotient to eliminate the dividend; when this is done, one obtains 89 chi, with 79 chi not exhausted; make it 90 chi; this gives the side of the company square.\end{uniten}

\begin{unitzh}\unitlabel{shushu-055.9}次置三十三隊，乗三将，得九十九；又四因，得三百九十六；増三，得三百九十九，為寔。以二為從方，一為從隅，開平方。歩法以從方進一位，至寔之十下；隅隅超一位，至寔之百下。乃約寔，置啇一十，生隅一入方，得一十二。乃命上啇除寔，訖，餘二百七十九。又以啇一十生隅入方，得二十二。畢，方一退，隅再退之。續于寔上啇九隊，以續啇九生隅入方，得三十一。乃命續啇除寔，適盡，得一十九。乗隊方面九十，得一千七百十尺，展為營方一百七十一丈，合問。\end{unitzh}
\begin{uniten}Next set up the 33 companies and multiply by the 3 generals, obtaining 99; again multiply by 4, obtaining 396; add 3, obtaining 399, as the dividend. Take 2 as the following side (congfang) and 1 as the following corner, and extract the square root. By the step method, the following side advances (jin) one place, to below the tens of the dividend; the corner, the corner skips one place, to below the hundreds of the dividend. Then, estimating the dividend, set up the quotient 10; generate the corner 1 and enter it into the side, obtaining 12. Then call upon the upper quotient to eliminate the dividend; when this is done, the remainder is 279. Again, with the quotient 10 generate the corner and enter it into the side, obtaining 22. When this is finished, the side retreats once and the corner retreats twice. Continuing above the dividend, take the quotient 9 companies; with the continued quotient 9 generate the corner and enter it into the side, obtaining 31. Then call upon the continued quotient to eliminate the dividend; it is exactly exhausted, giving 19. Multiplying the side of the company square, 90, gives 1710 chi; expanded (zhan), this makes the camp square 171 zhang, which agrees with the question.\end{uniten}

\begin{unitzh}\unitlabel{shushu-055.10}子部，天文算法類，算書之屬，《數學九章》，卷八上。子部，天文算法類，算書之屬，《數學九章》，卷八上。\end{unitzh}
\begin{uniten}Masters section (zibu), category of astronomy and mathematics, subcategory of mathematical books, Shuxue jiuzhang, juan 8, part 1. Masters section, category of astronomy and mathematics, subcategory of mathematical books, Shuxue jiuzhang, juan 8, part 1. [Translator's note: catalogue tags of the transcription.]\end{uniten}

\begin{unitzh}\unitlabel{shushu-055.11}方變鋭陣\end{unitzh}
\begin{uniten}Changing a square into a pointed formation\end{uniten}

\begin{unitzh}\unitlabel{shushu-055.12}問：歩兵五軍，軍一萬二千五百人，作方陣，人立地方八尺。欲變為前後鋭陣，陣後闊令多原方面半倍，陣間仍容騎路五丈以上，順鋭形出入。求方陣面、鋭陣長及前後鋭陣各布兵幾何？答曰：方面二百丈，方面布兵二百五十人；鋭後廣二百丈，鋭廣列共三百六十二人；鋭通正長三百丈；騎路二條，各闊五丈三尺；内鋭陣廣一百四十五丈六尺，列一百八十二人，長一百四十五丈六尺，計布兵一萬六千六百五十三人；外鋭兩廣各七十二丈，列九十人，計布兵四萬五千八百四十七人。\end{unitzh}
\begin{uniten}Question: Infantry of 5 armies, each army 12500 men, form a square formation (fangzhen); each man stands on ground 8 chi square. One wishes to change it into front and rear pointed formations (ruizhen): the rear breadth of the formation is to exceed the original square side by half (banbei); between the formations there must still be room for cavalry roads (qilu) of 5 zhang or more, for going out and coming in along the pointed shape. Find the side of the square formation, the length of the pointed formation, and how many soldiers are deployed in each of the front and rear pointed formations. Answer: The square side is 200 zhang; the soldiers deployed (bubing) along the square side are 250 men. The rear breadth of the pointed formation is 200 zhang; along the breadth of the pointed formation 362 men in all are ranked. The overall straight length of the pointed formation is 300 zhang. There are 2 cavalry roads, each 5 zhang 3 chi wide. The inner pointed formation is 145 zhang 6 chi in breadth, ranked with 182 men, and 145 zhang 6 chi in length, with 16653 soldiers deployed in all. The outer pointed formation has two breadths, each 72 zhang, ranked with 90 men, with 45847 soldiers deployed in all.\end{uniten}

\begin{unitzh}\unitlabel{shushu-055.13}按：鋭陣數，惟内鋭數合；外鋭通廣丈數及布兵數皆不合。詳見草後。\end{unitzh}
\begin{uniten}Note: Of the numbers of the pointed formation, only those of the inner pointed formation agree; the overall breadth in zhang of the outer pointed formation and its number of deployed soldiers all disagree. For details see after the working.\end{uniten}

\begin{unitzh}\unitlabel{shushu-055.14}術曰：以少廣求之。置兵，開平方，得方面人數{\zhsmall 開不盡，方為補隊}；以人立尺數乗之，為原方面。置原方面，以欲多數加之，為鋭後闊，亦為通長。倍馬路減之，餘為實；以人立尺約為闊布兵，不盡半歸馬路。以四約闊布兵，得外鋭一邉人。倍一邉人，併不歸，為内鋭長闊人數。副置，減一，餘乗其副，得數半之，為内鋭布兵。以減總兵，餘為外鋭布兵。\end{unitzh}
\begin{uniten}Method: Seek it by lesser breadth. Set up the soldiers and extract the square root, obtaining the number of men along the square side [Note: if the extraction does not come out exactly, this then makes the supplementary company (budui)]; multiply it by the number of chi each man stands on, to make the original square side. Set up the original square side and add to it the desired excess, to make the rear breadth of the pointed formation, which is also the overall length. Double the horse road (malu) and subtract it; the remainder is the dividend. Divide (yue) by the chi each man stands on, to make the soldiers deployed along the breadth; what is not exhausted is halved and returned to the horse roads. Divide the soldiers deployed along the breadth by 4, obtaining the men on one side of the outer pointed formation. Double the men on one side and add the undivided part (bugui), to make the number of men along the length and breadth of the inner pointed formation. Set this up in duplicate (fuzhi), subtract 1, multiply its duplicate by the remainder, and halve the number obtained, to make the soldiers deployed in the inner pointed formation. Subtract it from the total soldiers; the remainder is the soldiers deployed in the outer pointed formation.\end{uniten}

\begin{unitzh}\unitlabel{shushu-055.15}計立六萬二千五百人。鋭後廣、通長各三百丈。内鋭立一萬六千六百五十三人。外鋭立四萬五千八百四十七人。{\zhsmall 按：舊圖式不細，且在題後，今改正，移于此。}\end{unitzh}
\begin{uniten}In all, 62500 men stand. The rear breadth and the overall length of the pointed formation are each 300 zhang. In the inner pointed formation stand 16653 men. In the outer pointed formation stand 45847 men. [Siku editors: The old diagram was not detailed in form and was placed after the problem; now it has been corrected and moved here.]\end{uniten}

\begin{unitzh}\unitlabel{shushu-055.16}草曰：置一軍一萬二千五百，以五軍因之，得總兵六萬二千五百人，為寔。開平方，得二百五十人；以人立八尺乗之，得方面二百丈。置二百丈，加半倍一百丈，得三百丈，為鋭陣後闊，亦為鋭陣道長。先倍騎路五丈，得一十丈，以減後闊二百丈，餘二百九十丈，為實。以人立八尺約之，得三百六十二，為鋭後闊布兵，不盡四尺；以半之，得二尺，軰歸騎路，作五丈二尺。以四約鋭後闊布兵三百六十二人，得九十人，為外鋭一邉人。倍一邉九十，得一百八十，併不盡二人，共得一百八十二人，為内鋭廣布兵數，亦為長布兵。副置，加一，得一百八十三，乗副一百八十二，得三萬三千三百六；以半之，得一萬六千六百五十三人，為内鋭陣布兵。以減總兵六萬二千五百，餘四萬五千八百四十七人，為外鋭兵。\end{unitzh}
\begin{uniten}Working: Set up the 12500 of one army and multiply it by the 5 armies, obtaining the total soldiers, 62500 men, as the dividend. Extract the square root, obtaining 250 men; multiply by the 8 chi each man stands on, obtaining the square side, 200 zhang. Set up 200 zhang and add the half, 100 zhang, obtaining 300 zhang, as the rear breadth of the pointed formation, which is also the road length (daochang) of the pointed formation. First double the cavalry road of 5 zhang, obtaining 10 zhang, and subtract it from the rear breadth, 200 zhang; the remainder, 290 zhang, is the dividend. Divide it by the 8 chi each man stands on, obtaining 362, as the soldiers deployed along the rear breadth of the pointed formation, with 4 chi not exhausted; halve this, obtaining 2 chi, which go each to the cavalry roads, making 5 zhang 2 chi. Divide the 362 men deployed along the rear breadth of the pointed formation by 4, obtaining 90 men, as the men on one side of the outer pointed formation. Double the 90 of one side, obtaining 180; add the 2 men not exhausted, obtaining in all 182 men, as the number of soldiers deployed along the breadth of the inner pointed formation, which is also the soldiers deployed along the length. Set it up in duplicate and add 1, obtaining 183; multiply by the duplicate, 182, obtaining 33306; halve it, obtaining 16653 men, as the soldiers deployed in the inner pointed formation. Subtract it from the total soldiers, 62500; the remainder, 45847 men, is the soldiers of the outer pointed formation.\end{uniten}

\begin{unitzh}\unitlabel{shushu-055.17}按：草中以内鋭陣兵數减前方陣兵數，餘為外鋭陣兵數，非是，盖無以知兩總數為相等也。試以數明之：依束箭法，以總闊求得總三角數七萬零五百；以鋭闊求得内三角數一萬六千六百五十三；又以每人八尺除兩騎路濶十丈零四尺，得一十三人，與内鋭闊相加，得闊二百，求得内外間三角數二萬零一百。置總三角數，減内外間三角數，加内三角數，得六萬七千零五十三，與前方陣兵數相較，多四千五百五十三，安得謂之等乎？今另設歩法于後。\end{unitzh}
\begin{uniten}Note: In the working, the number of soldiers of the inner pointed formation is subtracted from the number of soldiers of the former square formation, and the remainder is taken as the number of soldiers of the outer pointed formation. This is not right, for there is no way of knowing that the two totals are equal. Let us try to make this clear with numbers. Following the method of bundled arrows (shujian fa), from the total breadth one finds the total triangular number (sanjiao shu) 70500; from the pointed breadth one finds the inner triangular number 16653; further, dividing the width of the two cavalry roads, 10 zhang 4 chi, by 8 chi per man gives 13 men; adding this to the inner pointed breadth gives a breadth of 200, from which one finds the triangular number between inner and outer, 20100. Set up the total triangular number, subtract the triangular number between inner and outer, and add the inner triangular number, obtaining 67053; compared with the number of soldiers of the former square formation, it is greater by 4553. How can they be called equal? Now a separate step method is set out below.\end{uniten}

\begin{unitzh}\unitlabel{shushu-055.18}法：設騎路之闊當二十人。先以總三角數與前方陣數相减，得今多八千人。乃倍騎路闊人數，得四十人，為截騎路上小三角之闊，求得小三角數八百二十，以减今多數，餘七千一百八十，為寔；以四十為法除之，得一百七十九人，為内鋭闊，餘二十人，依術内不盡者為補隊兵。次置總闊，減去内闊，餘一百九十六人；再減併騎闊四十人，餘一百五十六人；半之，得七十八人，為外後闊。是内鋭闊長皆為一百七十九人，外鋭長為三百七十五人，後兩闊共一百五十六人，騎路闊二十人。乃以内鋭闊求得内三角數一萬六千一百一十人；以内闊併兩騎闊，得二百一十九人，為闊，求得内外間三角數二萬四千零九十人。未置總三角數，内減去内外間三角數，餘四萬六千四百一十人；加内三角數，得六萬二千五百二十人；再減補隊兵二十人，得六萬二千五百人，與方陣總兵原數脗合。\end{unitzh}
\begin{uniten}Rule: Suppose the width of a cavalry road corresponds to 20 men. First take the difference between the total triangular number and the number of the former square formation, obtaining the present excess of 8000 men. Then double the number of men in the width of a cavalry road, obtaining 40 men, as the breadth of the small triangle cut off above the cavalry roads; find the small triangular number, 820, and subtract it from the present excess; the remainder, 7180, is the dividend. Take 40 as the divisor (fa) and divide, obtaining 179 men, as the inner pointed breadth, with 20 men remaining; in accordance with the method, what is not exhausted makes the soldiers of the supplementary company. Next set up the total breadth and subtract the inner breadth, leaving 196 men; further subtract the combined cavalry width, 40 men, leaving 156 men; halve this, obtaining 78 men, as the outer rear breadth. Thus the breadth and length of the inner pointed formation are both 179 men; the length of the outer pointed formation is 375 men; the two rear breadths together are 156 men; the cavalry road width is 20 men. Then from the inner pointed breadth find the inner triangular number, 16110 men; combining the inner breadth with the two cavalry widths gives 219 men, as the breadth, from which one finds the triangular number between inner and outer, 24090 men. Not yet set up the total triangular number [Translator's note: 未 'not yet' evidently stands for 末 'finally'.]; from it subtract the triangular number between inner and outer, leaving 46410 men; add the inner triangular number, obtaining 62520 men; further subtract the 20 soldiers of the supplementary company, obtaining 62500 men, which tallies exactly with the original number of the total soldiers of the square formation.\end{uniten}

\begin{unitzh}\unitlabel{shushu-055.19}子部，天文算法類，算書之屬，《數學九章》，卷八上。\end{unitzh}
\begin{uniten}Masters section, category of astronomy and mathematics, subcategory of mathematical books, Shuxue jiuzhang, juan 8, part 1.\end{uniten}

\begin{unitzh}\unitlabel{shushu-055.20}圓營敷布\end{unitzh}
\begin{uniten}Spreading out a circular camp\end{uniten}

\begin{unitzh}\unitlabel{shushu-056.1}問：周制一軍，欲布圓營九重，毎卒立圓邉六尺，重間相去比立尺數倍之。于内摘差兵四分之一出竒，不可縮營示弱，須令仍用原營布滿餘兵。欲知原營内外周及立人數，並出竒後每卒數、立尺數、外周人數各幾何？\end{unitzh}
\begin{uniten}Question: One army (jun) of the Zhou system wishes to deploy a circular camp in 9 layers; each soldier stands at 6 chi (尺) along the circle edge (yuanbian), and the distance between the layers is double the number of chi stood. From within, 1/4 of the soldiers are picked out and dispatched to go out as a surprise force (chuqi). The camp may not be shrunk to show weakness, so the original camp must still be used, filled out with the rest of the soldiers. One wishes to know the inner and outer perimeters of the original camp and the numbers of men standing, and also, after the surprise force has gone out, the number per soldier, the number of chi stood, and the number of men on the outer perimeter: how much is each?\end{uniten}

\begin{unitzh}\unitlabel{shushu-056.2}答曰：周制一軍一萬二千五百人，出竒三千一百二十五人。原内周八百四丈，立一千三百四十人；原外周八百六十一丈六尺，立一千四百三十六人。出竒後，原外周立一千八十九人，原内周立一千一十六人，内外周人立七尺九寸一分。\end{unitzh}
\begin{uniten}Answer: One army of the Zhou system is 12500 men; the surprise force sent out is 3125 men. The original inner perimeter is 804 zhang (丈), on which stand 1340 men; the original outer perimeter is 861 zhang 6 chi, on which stand 1436 men. After the surprise force has gone out, 1089 men stand on the original outer perimeter and 1016 men stand on the original inner perimeter; the men on the inner and outer perimeters stand at 7 chi 9 cun (寸) 1 fen (分).\end{uniten}

\begin{unitzh}\unitlabel{shushu-056.3}術曰：以啇功求之。置重數，減一，餘為叚；以叚乗圓差為衰；以衰乗重數為率。求圓周：以率減兵，餘如重數而一，得内周人數，不滿為餘兵。以人立圓邉乗内周人，得内周尺。倍衰乗圓邉為泛。以泛併内周尺，得外周尺，為實，如圓邉而一，得外周人。求出竒後：以率加存兵，如重數而一，得外周人，不滿為餘兵。以外周人約原外周尺，得後立尺。以後立尺約原内周，得内周人。\end{unitzh}
\begin{uniten}Method: Solve it by assessing works (shanggong). Set out the number of layers and subtract 1; the remainder is the segments (duan). Multiply the circular difference (yuancha) by the segments to make the gradation (cui). Multiply the number of layers by the gradation to make the rate (lü). To find the circular perimeters: subtract the rate from the soldiers; divide the remainder by the number of layers to obtain the number of men on the inner perimeter; what does not fill up is the surplus soldiers. Multiply the inner-perimeter men by the circle edge each man stands on to obtain the chi of the inner perimeter. Double the gradation and multiply by the circle edge to make the provisional quantity (fan). Add the provisional quantity to the chi of the inner perimeter to obtain the chi of the outer perimeter as the dividend (shi); divide by the circle edge to obtain the outer-perimeter men. To find the figures after the surprise force goes out: add the rate to the retained soldiers and divide by the number of layers to obtain the outer-perimeter men; what does not fill up is the surplus soldiers. Divide the chi of the original outer perimeter by the outer-perimeter men to obtain the later standing chi. Divide the original inner perimeter by the later standing chi to obtain the inner-perimeter men.\end{uniten}

\begin{unitzh}\unitlabel{shushu-056.4}按：求圓陣草中用圓束法。圓束實六等邉形，非圖形也。盖圓形重數相距等，則弧邉上相距不等；弧邉上相距等，則重數相距不等。惟圓束可並取相等，故用其法。至次陣減人數，不減營周尺數，則各重周上相距不能相等，故草中又以尺數求内周人數。然未免與圓束逐層相差數不合，亦僅取其大畧也。又舊用二圖，各㸃為圓周九重；今用一圖，㸃為六等邉形三重，惟取易見，則二圖九重，其理一也。\end{unitzh}
\begin{uniten}Note: In the working for finding the circular formation, the method of the circular bundle (yuanshu) is used. The circular bundle is in fact an equilateral hexagon, not a 'diagram' shape [Translator's note: 圖 'diagram' is evidently an error for 圓 'circle'.]. For in a circular shape, if the distances between the layers are equal, then the distances along the arc edges are not equal; if the distances along the arc edges are equal, then the distances between the layers are not equal. Only the circular bundle can take both as equal together; therefore its method is used. As for the next formation, the number of men is reduced but the number of chi of the camp perimeters is not reduced, so the distances along the perimeters of the various layers cannot be equal; therefore in the working the number of men on the inner perimeter is further found from the number of chi. Yet this cannot avoid disagreeing with the layer-by-layer differences of the circular bundle; it, too, only takes the rough approximation. Further, formerly two diagrams were used, each dotted as circular perimeters in 9 layers; now one diagram is used, dotted as an equilateral hexagon in 3 layers, only to make it easy to see; so for the two diagrams with 9 layers the principle is one and the same.\end{uniten}

\begin{unitzh}\unitlabel{shushu-056.5}草曰：置九重，減一，餘八，為叚；以乗圓束差六，得四十八，為衰{\zhsmall 按：圓束每層差不。今内外重數相距倍於人立相距，則每層差一十二，為倍差。常法：重數減一，與半差相乗為衰；今倍差，故即與差數相乗為衰也。}；九重得四百三十二，為率。\end{unitzh}
\begin{uniten}Working: Set out the 9 layers and subtract 1; the remainder 8 is the segments. Multiply the circular-bundle difference 6 by it to obtain 48 as the gradation [Siku editors: In the circular bundle each layer's difference is 不 [Translator's note: 不 'not' is evidently corrupt; a number, presumably 6, is expected.]. Now the distance between the inner and outer layers is double the distance at which the men stand, so each layer differs by 12, which is a doubled difference. By the usual rule, the number of layers minus 1 is multiplied with half the difference to make the gradation; now the difference is doubled, so it is multiplied directly with the number of the difference to make the gradation.]; by the 9 layers, obtain 432 as the rate.\end{uniten}

\begin{unitzh}\unitlabel{shushu-056.6}求原周：以率四百三十二減周制一軍一萬二千五百，餘一萬二千六十八，為實；如重數九而一，得一千三百四十人，為内周人數；不滿八人，為餘兵。\end{unitzh}
\begin{uniten}To find the original perimeters: subtract the rate, 432, from the Zhou-system army of 12500; the remainder 12068 is the dividend. Divide by the number of layers, 9, to obtain 1340 men as the number of men on the inner perimeter; the 8 men that do not fill up are the surplus soldiers.\end{uniten}

\begin{unitzh}\unitlabel{shushu-056.7}次以人立圓邉六尺乗内周人一千三百四十，得八千四十尺，收作八百四丈，為内周尺數。\end{unitzh}
\begin{uniten}Next multiply the inner-perimeter men, 1340, by the circle edge each man stands on, 6 chi, to obtain 8040 chi, which, converted, is 804 zhang, as the number of chi of the inner perimeter.\end{uniten}

\begin{unitzh}\unitlabel{shushu-056.8}倍衰四十八，得九十六，乗圓邉六尺，得五百七十六尺，為泛。\end{unitzh}
\begin{uniten}Double the gradation 48 to obtain 96; multiply by the circle edge, 6 chi, to obtain 576 chi as the provisional quantity.\end{uniten}

\begin{unitzh}\unitlabel{shushu-056.9}以泛五百七十六尺併内周八千四十尺，得八千六百一十六尺，為外周尺。\end{unitzh}
\begin{uniten}Add the provisional quantity, 576 chi, to the inner perimeter, 8040 chi, to obtain 8616 chi as the chi of the outer perimeter.\end{uniten}

\begin{unitzh}\unitlabel{shushu-056.10}以外周尺八千六百一十六為實，如圓邉六尺而一，得一千四百三十六人，為外周人數。\end{unitzh}
\begin{uniten}Take the chi of the outer perimeter, 8616, as the dividend; divide by the circle edge, 6 chi, to obtain 1436 men as the number of men on the outer perimeter.\end{uniten}

\begin{unitzh}\unitlabel{shushu-056.11}求出竒後：以竒母四約一萬二千五百，得三千一百二十五，為竒兵；以減總軍，餘九千三百七十五，為存兵。次以率四百三十二加之，得九千八百七十，為實；加重數九而一，得一千八十九，為外周人，不盡六。\end{unitzh}
\begin{uniten}To find the figures after the surprise force goes out: divide 12500 by the surprise denominator (qimu), 4, to obtain 3125 as the surprise soldiers; subtract it from the total army, the remainder 9375 being the retained soldiers. Next add the rate, 432, to it to obtain 9870 as the dividend; "add" the number of layers, 9, and make one, obtaining 1089 as the outer-perimeter men, with 6 not exhausted. [Translator's note: 9375 + 432 = 9807 = 9 × 1089 + 6; 加 "add" evidently stands for 如, i.e. "divide by" the number of layers.]\end{uniten}

\begin{unitzh}\unitlabel{shushu-056.12}次以原外周八千六百一十六尺為實，以外周人一千八十九約之，得七尺九寸一分，不盡二尺一分；與法求等，得三，俱約之，為分下三百六十三分之六十七。\end{unitzh}
\begin{uniten}Next take the original outer perimeter, 8616 chi, as the dividend and divide it by the outer-perimeter men, 1089, to obtain 7 chi 9 cun 1 fen, with 2 chi 1 fen not exhausted. Seek the common measure (deng) of this and the divisor (fa), obtaining 3; reduce both by it, giving 67/363 below the fen.\end{uniten}

\begin{unitzh}\unitlabel{shushu-056.13}置原内周八千四十尺為實，以後立尺七尺九寸一分約之，得一千一十六，為内周人數；不盡三尺四寸四分，為寛地。\end{unitzh}
\begin{uniten}Set out the original inner perimeter, 8040 chi, as the dividend and divide it by the later standing chi, 7 chi 9 cun 1 fen, to obtain 1016 as the number of men on the inner perimeter; the 3 chi 4 cun 4 fen not exhausted is spare ground (kuandi).\end{uniten}

\begin{unitzh}\unitlabel{shushu-056.14}本術所求内外周之人數既定，不拘竒出竒入，皆以六人為重差，或累差加減，各得諸重圍數；或併九重人，課總軍存兵。\end{unitzh}
\begin{uniten}Once the numbers of men on the inner and outer perimeters sought by this method have been fixed, regardless of whether the surprise force goes out or comes in, in every case take 6 men as the layer difference (zhongcha), or add or subtract the accumulated differences (leicha), to obtain the numbers for each of the layered circuits; or sum the men of the 9 layers to check the total army and the retained soldiers.\end{uniten}

\begin{unitzh}\unitlabel{shushu-056.15}計布圓陣\end{unitzh}
\begin{uniten}Calculating the Deployment of a Circular Formation\end{uniten}

\begin{unitzh}\unitlabel{shushu-056.16}問：歩率二千六百人為圖陣，人立圓九尺，形如車輻，無麗布陣。陣重間倍人立圓邉尺數，須合内徑七十二丈，圓法用周三徑一之率。欲知陣重幾數，及内外周、通徑，並所立人數各幾何？\end{unitzh}
\begin{uniten}Question: An infantry contingent of 2600 men forms a 'diagram' formation; each man stands at 9 chi on the circle; its shape is like the spokes of a cart wheel, the formation deployed in wuli (無麗) fashion [Translator's note: 圖 'diagram' here and in 圖邉 below is evidently for 圓 'circle'; 無麗 is presumably for 魚麗, the 'fish-scale' formation.]. The interval between the layers of the formation is double the number of chi of circle edge each man stands on; it must fit an inner diameter of 72 zhang; the circle rule uses the rate of circumference 3, diameter 1. One wishes to know how many layers the formation has, the inner and outer perimeters and the through-diameter (tongjing), and the numbers of men standing on them: how much is each?\end{uniten}

\begin{unitzh}\unitlabel{shushu-056.17}答曰：内周二百一十六丈，立二百四十人；外周三百二丈四尺，立三百三十六人；通徑一百丈八尺；陣計九重{\zhsmall 不盡八人}。\end{unitzh}
\begin{uniten}Answer: The inner perimeter is 216 zhang, on which stand 240 men; the outer perimeter is 302 zhang 4 chi, on which stand 336 men; the through-diameter is 100 zhang 8 chi; the formation counts 9 layers [Note: 8 men not exhausted].\end{uniten}

\begin{unitzh}\unitlabel{shushu-056.18}術曰：以啇功求之。以圓率因内徑為内周，以人立尺約之，為内周人數。乃以圓求差率為隅。次置内周人，減隅，餘約從方。列兵數為寛，開平方，得重數，不盡為餘兵。置重數，減一，餘四因，又乗圓邉尺數，併内徑，共為通徑。以圓率因通徑，得外周。\end{unitzh}
\begin{uniten}Method: Solve it by assessing works (shanggong). Multiply (yin) the inner diameter by the circle rate to make the inner perimeter; divide it by the chi each man stands, making the number of men on the inner perimeter. Then, from the circle, find the difference rate (chalü) to be the corner (yu). Next set out the inner-perimeter men and subtract the corner; the remainder "divide" (約) the joining square (congfang). Set out the number of soldiers as the "breadth" (寛) [Translator's note: 約 "divide" and 寛 "breadth" evidently stand for 為 "make" and 實 "dividend"; cf. the working below.]; extract the square root to obtain the number of layers; what is not exhausted is the surplus soldiers. Set out the number of layers, subtract 1, multiply the remainder by 4, and further multiply by the number of chi of the circle edge; add the inner diameter, the total making the through-diameter. Multiply the through-diameter by the circle rate to obtain the outer perimeter.\end{uniten}

\begin{unitzh}\unitlabel{shushu-056.19}按：舊本有圖，前題同，今删去。\end{unitzh}
\begin{uniten}Note: The old copy has a diagram, the same as in the preceding problem; it has now been deleted.\end{uniten}

\begin{unitzh}\unitlabel{shushu-056.20}草曰：以圓率三因内徑七十二丈，得二千一百六十尺，為内周。以圓邉九尺約内周，得二百四十，為内周人數。乃以圓束差六為從隅。次置内周二百四十人，減隅，餘二百三十四，為從方。列兵二千六百為寔，開平方。歩法：從方進一位，隅法超一位；今方隅皆不可超進，乃于寔約啇，置九重。以啇生隅六，得五十四，増入從方内，共得二百八十八。乃命上啇九重，除寔訖，寔餘八人，為餘兵。副置九重，減一，餘八，以四因之{\zhsmall 按：九重八間，徑両端應二因之；間倍于立歩，又應二因之；今合為四因。}，得三十二；又乗圖邉九尺，得二百八十八尺；併内徑七百二十尺，得一千八尺，為通徑。又以圓率三因通徑，得三千二十四尺，為外周。次以圓邉九尺為法，除外周尺數，得三百三十六人，為外周人數。合問。\end{unitzh}
\begin{uniten}Working: Multiply the inner diameter, 72 zhang, by 3, the circle rate, obtaining 2160 chi as the inner perimeter. Divide the inner perimeter by the circle edge, 9 chi, obtaining 240 as the number of men on the inner perimeter. Then take the circular-bundle difference, 6, as the joining corner (congyu). Next set out the 240 inner-perimeter men and subtract the corner; the remainder 234 is the joining square. Set out the soldiers, 2600, as the dividend and extract the square root. By the stepping rule (bufa), the joining square advances one place and the corner divisor (yufa) leaps one place; now neither the square nor the corner can leap or advance, so at the dividend estimate the quotient (shang), placing 9 layers. With the quotient generate the corner 6, obtaining 54; add it into the joining square, together obtaining 288. Then with the upper quotient, 9 layers, divide the dividend; when this is done, the dividend has 8 men left, as the surplus soldiers. Set out separately the 9 layers, subtract 1, remainder 8; multiply it by 4 [Siku editors: 9 layers have 8 intervals; the diameter has two ends, so one should multiply by 2; the interval is double the standing step, so one should again multiply by 2; now these are combined into multiplying by 4.], obtaining 32; further multiply by the 'diagram' edge, 9 chi, obtaining 288 chi; add the inner diameter, 720 chi, obtaining 1008 chi as the through-diameter. Further multiply the through-diameter by the circle rate 3, obtaining 3024 chi as the outer perimeter. Next take the circle edge, 9 chi, as the divisor and divide the number of chi of the outer perimeter, obtaining 336 men as the number of men on the outer perimeter. This agrees with the question.\end{uniten}

\begin{unitzh}\unitlabel{shushu-056.21}按：圓束環積有内周求重數法：置積為寔，圓束差半之為從隅，又以半差減内周，餘為從方，開平方，得重數。此圓束環積每層為倍差，故即以圓束差為從隅，減内周為從方也。又按：周三徑一，正與六邉形相合，故人數、尺俱無竒零也。\end{unitzh}
\begin{uniten}Note: The rule for finding the number of layers of a circular-bundle ring area (huanji) when the inner perimeter is given: set out the area as the dividend; halve the circular-bundle difference to make the joining corner; further subtract the half difference from the inner perimeter, the remainder making the joining square; extract the square root to obtain the number of layers. In this circular-bundle ring area each layer has a doubled difference; therefore the circular-bundle difference itself is taken as the joining corner, and is subtracted from the inner perimeter to make the joining square. Further note: circumference 3, diameter 1 agrees exactly with the equilateral hexagon; therefore the numbers of men and the chi are all without odd fractions.\end{uniten}

\begin{unitzh}\unitlabel{shushu-056.22}數學九章卷八上\end{unitzh}
\begin{uniten}Mathematical Treatise in Nine Sections (Shuxue jiuzhang), juan 8, part 1\end{uniten}

\section*{卷8下}\addcontentsline{toc}{section}{卷8下}

\begin{unitzh}\unitlabel{shushu-057.1}欽定四庫全書。數學九章卷八下。宋·秦九韶撰。\end{unitzh}
\begin{uniten}Imperially Commissioned Complete Library of the Four Treasuries (Qinding Siku quanshu). Mathematical Treatise in Nine Sections (Shuxue jiuzhang), juan 8, part 2. Composed by Qin Jiushao of the Song.\end{uniten}

\begin{unitzh}\unitlabel{shushu-057.2}軍旅\end{unitzh}
\begin{uniten}Military Affairs\end{uniten}

\begin{unitzh}\unitlabel{shushu-057.3}軍器功程\end{unitzh}
\begin{uniten}Work Quotas for Military Equipment\end{uniten}

\begin{unitzh}\unitlabel{shushu-057.4}問：今欲造弓、刀各一萬副，箭一百萬隻。據工程：七人九日造弓八張，八人六日造刀五副，三人二日造箭一百五十隻。作院見管弓作二百人，刀作五百四十人，箭作二百七十六人。欲知畢日㡬何？答曰：造弓一萬張，三百九十三日四分日之三；造刀一萬副，一百七十七日九分日之七；造箭一百萬隻，一百四十四日二百七分日之一百八十二。{\zhsmall 按：六十九分之六十四，訛二百七分之一百八十二。}\end{unitzh}
\begin{uniten}Question: Now one wishes to make bows and swords, 10,000 sets of each, and 1,000,000 arrows. According to the work quotas (gongcheng): 7 people in 9 days make 8 bows; 8 people in 6 days make 5 sets of swords; 3 people in 2 days make 150 arrows. The workshop (zuoyuan) currently has under its management 200 bow-workers, 540 sword-workers and 276 arrow-workers. One wishes to know: in how many days will it be finished? Answer: Making 10,000 bows, 393 3/4 days; making 10,000 sets of swords, 177 7/9 days; making 1,000,000 arrows, 144 182/207 days. [Siku editors: 64/69 has been corrupted into 182/207.]\end{uniten}

\begin{unitzh}\unitlabel{shushu-057.5}術曰：以粟米求之，互換入之。置各功程原人率於右行，置原日數於中行，置欲求數為左行，以三行對之，為各實，列右行。置原物數於中行，置見𬋩人為左行，以左行乘中行，各為法。以對除右行，各得日數。\end{unitzh}
\begin{uniten}Method: Solve it by Millet and Grain (sumi), entering it by interchange (huhuan). Set the original rates (lü) of people of each work quota in the right column; set the original numbers of days in the middle column; set the numbers sought as the left column; take the three columns correspondingly (dui) to make the respective dividends (shi), arranged in the right column. Set the original numbers of articles in the middle column; set the people currently under management as the left column; multiply the middle column by the left column, each to make a divisor (fa). Divide the right column correspondingly; each gives the number of days.\end{uniten}

\begin{unitzh}\unitlabel{shushu-057.6}草曰：置原造弓七人、造刀八人、造箭三人於右行；次置造弓九日、造刀六日、造箭二日，列中行；又置於造弓一萬、欲造刀一萬、欲造箭一百萬，列左行；以三行對乗。\end{unitzh}
\begin{uniten}Working: Set the original 7 people for making bows, 8 people for making swords and 3 people for making arrows in the right column; next set 9 days for making bows, 6 days for making swords and 2 days for making arrows, arranged as the middle column; further set, for making bows, 10,000, for the swords one wishes to make, 10,000, and for the arrows one wishes to make, 1,000,000, arranged as the left column; multiply the three columns correspondingly (dui cheng).\end{uniten}

\begin{unitzh}\unitlabel{shushu-057.7}次列原造弓八張、刀五副、箭一百五十隻放中；又列見管弓作二百人、刀作五百四十人、箭作二百七十六人於左行。\end{unitzh}
\begin{uniten}Next arrange the original 8 bows made, 5 sets of swords and 150 arrows, putting them in the middle; further arrange the currently managed 200 bow-workers, 540 sword-workers and 276 arrow-workers in the left column.\end{uniten}

\begin{unitzh}\unitlabel{shushu-057.8}以兩行對乗之，上得一千六百，中得二千七百，下得四萬一千四百，各為法。\end{unitzh}
\begin{uniten}Multiply the two columns correspondingly: the top gives 1,600, the middle gives 2,700, the bottom gives 41,400; each is a divisor.\end{uniten}

\begin{unitzh}\unitlabel{shushu-057.9}先以上法一千六百除寄右行弓日實六十三萬日，得三百九十三日，為造弓一萬張日數。\end{unitzh}
\begin{uniten}First, with the upper divisor 1,600, divide the bow-days dividend of 630,000 days lodged (ji) in the right column, obtaining 393 days as the number of days for making 10,000 bows.\end{uniten}

\begin{unitzh}\unitlabel{shushu-057.10}不盡一千二百，與法求等，得四百，俱約之，為四分日之三。\end{unitzh}
\begin{uniten}The remainder (bu jin) 1,200: seek the equal (qiu deng) with the divisor, obtaining 400; reduce both by it, giving 3/4 day.\end{uniten}

\begin{unitzh}\unitlabel{shushu-057.11}次以中得二千七百除寄右行刀日實四十八萬日，得一百七十七日，為造刀一萬副日數。\end{unitzh}
\begin{uniten}Next, with what the middle gave, 2,700, divide the sword-days dividend of 480,000 days lodged in the right column, obtaining 177 days as the number of days for making 10,000 sets of swords.\end{uniten}

\begin{unitzh}\unitlabel{shushu-057.12}不盡二千一百，與法求等，得三百，俱約之，為九分日之七。\end{unitzh}
\begin{uniten}The remainder 2,100: seek the equal with the divisor, obtaining 300; reduce both by it, giving 7/9 day.\end{uniten}

\begin{unitzh}\unitlabel{shushu-057.13}次以下法四萬□□百除寄右行箭日實六百萬日，得一百四十四日，為造箭一百萬隻日數。\end{unitzh}
\begin{uniten}Next, with the lower divisor 40,000 [blank] [blank] hundred, divide the arrow-days dividend of 6,000,000 days lodged in the right column, obtaining 144 days as the number of days for making 1,000,000 arrows.\end{uniten}

\begin{unitzh}\unitlabel{shushu-057.14}不盡三萬六千四百日，與法四萬一千四百求等，得二百，俱以約之，得二百七分日之一百八十二，為造箭日分。合問。{\zhsmall 按：不盡數三萬八千四百，誤為三萬六千四百；分母六十九，訛二百零七；分子六十四，訛一百八十二。}\end{unitzh}
\begin{uniten}The remainder of 36,400 days: seek the equal with the divisor 41,400, obtaining 200; reduce both by it, obtaining 182/207 day as the fraction of a day for making the arrows. This answers the question. [Siku editors: The remainder 38,400 has been mistakenly given as 36,400; the denominator (fenmu) 69 has been corrupted into 207; the numerator (fenzi) 64 has been corrupted into 182.]\end{uniten}

\begin{unitzh}\unitlabel{shushu-057.15}計造軍衣\end{unitzh}
\begin{uniten}Calculating the Making of Military Clothing\end{uniten}

\begin{unitzh}\unitlabel{shushu-057.16}問：庫有布、綿、絮三色，計料欲製軍衣。其布，六人八疋，少一百六十疋；七人九疋，剰五百六十疋。其綿，八人一百五兩，剰一萬六千五百兩；九人一百七十兩，剰一萬四千四百兩。其絮，四人一十三斤，少六千八百四斤；五人一十四斤，適足。欲知軍士及布、綿、絮各幾何？答曰：兵士一萬五千一百二十人，布二萬疋，綿三十萬兩，絮四萬二十三百二十六斤。\end{unitzh}
\begin{uniten}Question: A storehouse has three kinds of goods, cloth, silk floss (mian) and wadding (xu); reckoning the materials, it wishes to make military clothing. Of the cloth: at 6 people to 8 pi (疋), it is short by 160 pi; at 7 people to 9 pi, 560 pi are left over. Of the silk floss: at 8 people to 105 liang (兩) [Translator's note: the working below has 150.], 16,500 liang are left over; at 9 people to 170 liang, 14,400 liang are left over. Of the wadding: at 4 people to 13 jin (斤), it is short by 6,804 jin; at 5 people to 14 jin, it is exactly sufficient (shizu). One wishes to know how many army men there are, and how much cloth, silk floss and wadding each. Answer: soldiers, 15,120 people; cloth, 20,000 pi; silk floss, 300,000 liang; wadding, forty thousand twenty-three hundred twenty-six jin.\end{uniten}

\begin{unitzh}\unitlabel{shushu-057.17}術曰：以盈拙求之。置人數於左右之中，置所給物各於其上，置盈拙各於其下，令維乘之。先以人數互數乘其所給率，相減，餘為法。次以人數相乘為寄。後以盈拙互乘其上未減者，是謂未乘。騐其下係一盈一拙，以上下皆併之為物，其上併之實，其下併之乘寄為兵寔，如法而一，各得。騐其係兩盈或兩拙者，以上下皆相減之，其上減之餘為物寔，其下減之餘乘寄為兵寔，二寔皆如法而一，各得。騐其或一盈一足或一拙一足者，其適足乃以空互乘其上未減者，去之，只以所用盈拙數互乘其上為物寔，以盈或拙一數乘寄為兵寔，皆如法而一，各得。\end{unitzh}
\begin{uniten}Method: Solve it by excess and deficit (ying zhuo) [Translator's note: 拙 'clumsy' here, like 朏 in the working below, evidently stands for 朒 'deficit'.]. Set the numbers of people in the middle on the left and right; set the goods allotted, each above them; set the excess and deficit, each below them; and cross-multiply them (wei cheng). First, with the numbers of people, mutually multiply (hu shu cheng) the rates (lü) allotted; subtract them from each other; the remainder is the divisor. Next multiply the numbers of people together to make the lodged number (ji). Afterwards mutually multiply (hu cheng) the excess and deficit with those above that have not been subtracted; this is called 'not-yet multiplying' (wei cheng) [Translator's note: 未 evidently stands for 維.]. If on checking (yan) those below are one excess and one deficit, add together both the upper and the lower, for the goods: the sum of the upper, the dividend; the sum of the lower, multiplied by the lodged number, is the dividend for the soldiers; divide by the divisor, and each is obtained. If on checking they are two excesses or two deficits, subtract both the upper and the lower from each other: the remainder of the upper subtraction is the dividend for the goods; the remainder of the lower subtraction, multiplied by the lodged number, is the dividend for the soldiers; divide both dividends by the divisor, and each is obtained. If on checking it is either one excess and one sufficient or one deficit and one sufficient, then for the exactly sufficient one use the empty (kong) to mutually multiply the unsubtracted one above it, and discard this; only mutually multiply the number of excess or deficit that is used with the one above to make the dividend for the goods; multiply the single number of excess or deficit by the lodged number to make the dividend for the soldiers; divide both by the divisor, and each is obtained.\end{uniten}

\begin{unitzh}\unitlabel{shushu-057.18}求布草曰：置布於六人左中，八疋於左朏上，一百六十四疋於左下；置七人於右中，九疋於右上，盈百五六十於布下。先以左右之中六七互乘左右之上，訖，左上得五十五，右上得五十四，以相減之，餘二為法。次以左右中六七相乘，得四十二，為寄於中。次以左下虧一百六十乘右上未减五十四，得八十六百四十；又以右下盈五百六十乘左上未减五十六，得三萬一千三百六十。驗得左右之下係一盈一朏，當併之。以三萬一千三百六十併右上八千六百四十，得四萬，為布實。次以左下朏一百六十併左下盈五百六十，得七百二十，乗寄四十二，得三萬二百四十，為兵實。二實皆如法二而一，得二萬疋為布，得一萬五千一百二十為兵。\end{unitzh}
\begin{uniten}Working for finding the cloth: Set the cloth at 6 people in the left middle, 8 pi in the left deficit (fei) top, 164 pi in the left bottom; set 7 people in the right middle, 9 pi in the right top, and the excess, a hundred five sixty, below the cloth [Translator's note: this sentence is garbled in the text.]. First, with the middles on the left and right, 6 and 7, mutually multiply the tops on the left and right; when done, the left top gives 55 and the right top gives 54; subtract them from each other, and the remainder 2 is the divisor. Next multiply the middles on the left and right, 6 and 7, together, obtaining 42, as the lodged number in the middle. Next take the left-bottom shortfall (kui) 160 to multiply the not-yet-subtracted right top 54, obtaining eighty-six hundred forty; further take the right-bottom excess 560 to multiply the not-yet-subtracted left top 56, obtaining 31,360. On checking it is found that the bottoms on the left and right are one excess and one deficit; they should be added. Add 31,360 to the right top 8,640, obtaining 40,000 as the dividend for the cloth. Next add the left-bottom deficit 160 to the left-bottom excess 560, obtaining 720; multiply by the lodged 42, obtaining 30,240 as the dividend for the soldiers. Divide both dividends by the divisor 2: obtaining 20,000 pi as the cloth, and obtaining 15,120 as the soldiers.\end{uniten}

\begin{unitzh}\unitlabel{shushu-057.19}求布圖。<子部,天文算法類,算書之屬,數學九章,卷八下>\end{unitzh}
\begin{uniten}Diagram for finding the cloth. <Masters branch, Astronomy and Mathematics category, Books on Calculation subcategory, Shuxue jiuzhang, juan 8, part 2>\end{uniten}

\begin{unitzh}\unitlabel{shushu-057.20}求綿草曰：置八人於左中，綿一百五十兩於左上，餘一萬六千五百兩於左下；次置九人於右中，一百七十兩於右上，餘一萬四千四百兩於右下。以左右中八九五乘各上，訖，左上得一千三百五十，右上得一千三百六十，相減，餘一十為法。次以中八九相乘，得七萬二，為寄於中。次以左下一萬六千五百乘右上一千三百六萬，得二千二百四十四萬；却以右下一萬四千四百乘左上一千三百五十，得一千九百四十四萬。驗其下係兩盈，當相減之。其右上餘三百萬為綿實；其左右之下亦相減之，餘二千一百，乘寄七十二，得二十五萬一千二百，為兵實。二實皆如法一十而一，綿得三十萬兩，兵得一萬五千一百二十人。\end{unitzh}
\begin{uniten}Working for finding the silk floss: Set 8 people in the left middle, the silk floss 150 liang in the left top, the leftover (yu) 16,500 liang in the left bottom; next set 9 people in the right middle, 170 liang in the right top, the leftover 14,400 liang in the right bottom. With the middles on the left and right, 8 and 9, 'five'-multiply (wu cheng) each top [Translator's note: 五 evidently stands for 互 'mutually'.]; when done, the left top gives 1,350 and the right top gives 1,360; subtract them from each other, and the remainder 10 is the divisor. Next multiply the middles 8 and 9 together, obtaining 70,002, as the lodged number in the middle. Next take the left bottom 16,500 to multiply the right top 13,060,000, obtaining 22,440,000; then take the right bottom 14,400 to multiply the left top 1,350, obtaining 19,440,000. On checking, those below are two excesses; they should be subtracted from each other. The remainder at the right top, 3,000,000, is the dividend for the silk floss; the bottoms on the left and right are likewise subtracted from each other, the remainder 2,100, multiplied by the lodged 72, giving 251,200 as the dividend for the soldiers. Divide both dividends by the divisor 10: the silk floss obtained is 300,000 liang, the soldiers obtained are 15,120 people.\end{uniten}

\begin{unitzh}\unitlabel{shushu-057.21}求綿圖。<子部,天文算法類,算書之屬,數學九章,卷八下>\end{unitzh}
\begin{uniten}Diagram for finding the silk floss. <Masters branch, Astronomy and Mathematics category, Books on Calculation subcategory, Shuxue jiuzhang, juan 8, part 2>\end{uniten}

\begin{unitzh}\unitlabel{shushu-057.22}求絮草曰：置四人於左中，一十三觔於左上，少六千八百四觔於左下；又置五人於右中，一十四觔於右上，適足為空於右下。以左右之中四五互乘其上，訖，左上得六十五，右上得五□六，利減，餘九為法。以中四五相乘，得二十，為寄於中。先以左下六千八百四互乘右上五十六，得三十八萬一千四百二十；却以右適足之空乘左上六十五，亦為空，乃去之。只以右上三十八萬一千二十四觔為絮實；只以左下六千八百四乘寄二千人，得一十三萬六千八十，為兵實。二實皆如法九而一，具絮得四萬二千三百三十六觔，其兵得一萬五千一百二十人。合問。\end{unitzh}
\begin{uniten}Working for finding the wadding: Set 4 people in the left middle, 13 jin (觔) in the left top, the shortfall of 6,804 jin in the left bottom; further set 5 people in the right middle, 14 jin in the right top, and the exactly sufficient, making an empty (kong), in the right bottom. With the middles on the left and right, 4 and 5, mutually multiply those above them; when done, the left top gives 65 and the right top gives 5[blank]6; subtracting them (li jian) [Translator's note: 利 evidently stands for 相 'mutually'.], the remainder 9 is the divisor. Multiply the middles 4 and 5 together, obtaining 20, as the lodged number in the middle. First take the left bottom 6,804 to mutually multiply the right top 56, obtaining 381,420; then take the empty of the exactly sufficient on the right to multiply the left top 65, which likewise gives empty; so discard it. Only take the right top, 381,024 jin, as the dividend for the wadding; only take the left bottom 6,804 to multiply the lodged 2,000 people, obtaining 136,080 as the dividend for the soldiers. Divide both dividends by the divisor 9: the wadding, in full (ju), obtained is 42,336 jin, and the soldiers obtained are 15,120 people. This answers the question.\end{uniten}

\begin{unitzh}\unitlabel{shushu-057.23}求絮圖。<子部,天文算法類,算書之屬,數學九章,卷八下>\end{unitzh}
\begin{uniten}Diagram for finding the wadding. <Masters branch, Astronomy and Mathematics category, Books on Calculation subcategory, Shuxue jiuzhang, juan 8, part 2>\end{uniten}

\begin{unitzh}\unitlabel{shushu-057.24}已上布、綿、絮三項，求人兵數皆同。今仍於各圖立算求之，以合本術。\end{unitzh}
\begin{uniten}For the three items above, cloth, silk floss and wadding, the numbers of soldiers (ren bing) sought are all the same. Now we nevertheless set up the computation in each diagram to find them, so as to accord with the method proper.\end{uniten}

\begin{unitzh}\unitlabel{shushu-057.25}先計軍程\end{unitzh}
\begin{uniten}Reckoning the Army's Journey in Advance\end{uniten}

\begin{unitzh}\unitlabel{shushu-058.1}問：一軍三將，三十隊，隊七十五人。每將分左右傔，作九行，爬頭拽行，每日六十里。明日路狹，以軍拽行，至晚。欲知先宿程里數幾何？答曰：六里二百四十歩。\end{unitzh}
\begin{uniten}Question: An army has 3 generals (jiang) and 30 squads (dui), each squad 75 men. Each general divides them into left and right attendants (qian 傔), making 9 files (hang). They march in a drawn-out column with the head crawling (patou yexing), 60 li (里) each day. The next day the road is narrow, and the army marches in a drawn-out column until evening. One wishes to know: how many li is the advance lodging stage (xian sucheng)? Answer: 6 li 240 bu (歩).\end{uniten}

\begin{unitzh}\unitlabel{shushu-058.2}術曰：以均輸求之。置行數為法，以單數一行乘日程為寔，寔如法兩一，得宿程里歩。\end{unitzh}
\begin{uniten}Method: Solve it by equitable transport (junshu). Set up the number of files as the divisor (fa). With the single number (danshu), 1 file, multiply the day's stage (richeng) as the dividend (shi). Divide the dividend by the divisor [Translator's note: the text writes 兩一 where 而一 is usual], obtaining the lodging stage in li and bu.\end{uniten}

\begin{unitzh}\unitlabel{shushu-058.3}草曰：置行數九為法，以單傔數一行用乘六十里為寔，寔如法而一，得六十，不盡六里。以里法三百六十歩通之，得二千一百六十歩，又為寔，仍如法九而一，得二百四十歩，為六里二百四十歩宿程。\end{unitzh}
\begin{uniten}Working: Set up the number of files, 9, as the divisor. With the number of single attendants, 1 file, multiply 60 li as the dividend. Divide the dividend by the divisor, obtaining 60 [Translator's note: evidently an error for 6], with 6 li unexhausted. Convert this (tong) by the li standard (lifa) of 360 bu, obtaining 2,160 bu. Again take it as the dividend and again divide by the divisor 9, obtaining 240 bu. This makes 6 li 240 bu as the lodging stage.\end{uniten}

\begin{unitzh}\unitlabel{shushu-058.4}按：此草易見，舊有算式，今刪。\end{unitzh}
\begin{uniten}Note: This working is easy to see. Formerly there was a computation layout (suanshi); it has now been deleted.\end{uniten}

\begin{unitzh}\unitlabel{shushu-058.5}移運均勞{\zhsmall 分郡縣郷科均}\end{unitzh}
\begin{uniten}Equalizing the labor of transport (yiyun junlao) [Note: dividing among commandery, counties and cantons, and levying equally]\end{uniten}

\begin{unitzh}\unitlabel{shushu-058.6}問：今起夫移運縣餉於某郡交納，合起一萬二千夫。甲州有三縣：上縣力五十七萬三千二百五十九貫五百文，至輸所九百二十五里；中縣力五十萬四千九百八十三貫七百八十文，至輸所六百五十二里；下縣力四十九萬八千七百六十貫九百五十文，至輸所四百六十五里。乙軍倚郭一縣五郷：仁郷力一百二萬八千三百七十一貫九百八十文，至輸所七百六里；義鄉力一□□□□□□□□□□□六百文，至輸所七百九十五里；禮郷力一十萬八千四百六十三貫五十文，至輸所七百九十里；智鄉力八萬四千千二百三十六貫二百八十五文，至輸所七百十九里；信鄉力九千三百四十五貫一百六十六文，至輸所八百四里。欲知以物力多寡、道里逺近均運之，令費勞等，各合科夫幾何？答曰：甲州上縣差二千四百三十夫，中縣差三千三十七夫，下縣差四千二百六夫；乙軍郭縣仁鄉七百一十三夫，義鄉五百八十九夫，禮鄉五百三十八夫，智鄉四百四十一夫，信郷四十六夫。\end{unitzh}
\begin{uniten}Question: Now laborers (fu) are raised to transport county provisions to a certain commandery (jun 郡) for delivery; in all, 12,000 laborers are to be raised. Prefecture A (Jiazhou 甲州) has 3 counties (xian). The upper county's strength (li 力) is 573,259 guan (貫) 500 wen (文), and to the delivery place (shusuo) it is 925 li. The middle county's strength is 504,983 guan 780 wen, and to the delivery place it is 652 li. The lower county's strength is 498,760 guan 950 wen, and to the delivery place it is 465 li. Military prefecture B (Yijun 乙軍) has 1 county attached to its city wall (yiguo), with 5 cantons (xiang). The strength of Ren (仁) canton is 1,028,371 guan 980 wen, and to the delivery place it is 706 li. The strength of Yi (義) canton is 1 [blank of 11 characters] 600 wen, and to the delivery place it is 795 li. The strength of Li (禮) canton is 108,463 guan 50 wen, and to the delivery place it is 790 li. The strength of Zhi (智) canton is 84,236 guan 285 wen [Translator's note: the text writes 八萬四千千二百三十六, with 千 'thousand' doubled], and to the delivery place it is 719 li. The strength of Xin (信) canton is 9,345 guan 166 wen, and to the delivery place it is 804 li. One wishes to know: if the transport is equalized according to the greater or lesser property strength (wuli) and the farther or nearer road distance, so that expense and labor are equal, how many laborers should each be levied (ke)? Answer: Prefecture A: the upper county is assigned (chai) 2,430 laborers; the middle county is assigned 3,037 laborers; the lower county is assigned 4,206 laborers. Military prefecture B, the city county (guoxian): Ren canton 713 laborers; Yi canton 589 laborers; Li canton 538 laborers; Zhi canton 441 laborers; Xin canton 46 laborers.\end{uniten}

\begin{unitzh}\unitlabel{shushu-058.7}術曰：以均輸求之。置各縣及鄉力，皆如里而一，不盡者約之，復通分内子互乘之，或就母遷退之，各得變力。可約約之，為定力。副併為法。以合起夫徧乘未併定力，各得為實。並如前法而一，各得夫。其餘分軰之。\end{unitzh}
\begin{uniten}Method: Solve it by equitable transport (junshu). Set up the strength of each county and canton and divide each by its li. Where there is an unexhausted remainder, reduce it. Then convert into parts and add in the numerators (tongfen nazi) and multiply crosswise (hucheng), or shift them back according to the denominators (jiu mu qiantui). Each gives the transformed strength (bianli). Where they can be reduced, reduce them, making the fixed strengths (dingli). Add these together in an auxiliary position (fubing) as the divisor. With the total laborers to be raised, multiply over all the fixed strengths not yet added together; each gives a dividend. Divide all by the preceding divisor; each gives its laborers. Apportion the remaining fractions by rank (fenbei).\end{uniten}

\begin{unitzh}\unitlabel{shushu-058.8}草曰：置甲州三縣及乙軍五鄉物力、里數，作八行列之，具圖于後。\end{unitzh}
\begin{uniten}Working: Set up the property strengths and li distances of the 3 counties of Prefecture A and the 5 cantons of Military prefecture B, arranging them in 8 rows (hang); the diagram is provided below.\end{uniten}

\begin{unitzh}\unitlabel{shushu-058.9}子部·天文算法類·算書之屬·《數學九章》·卷八下\end{unitzh}
\begin{uniten}Masters section (zibu) · Astronomy and mathematics category · Mathematical books subcategory · Shuxue jiuzhang · juan 8, part 2 [Translator's note: a Siku quanshu catalogue label, which in the transcription apparently marks the place of the diagram image.]\end{uniten}

\begin{unitzh}\unitlabel{shushu-058.10}置上縣力五十七萬三千二百五十九貫五百文，如九百二十五里而一，得力六十一萬九千七百四十。置中縣五十萬四千九百八十三貫七百八十文，如六百五十二里而一，得力七十七萬四千五百一十五。置下縣四十九萬八千七百六十貫九百五十，如四百六十五里而一，得力一百七萬二千六百四，不盡九十文，與法求等，得十五，約之，得三十一分之六。\end{unitzh}
\begin{uniten}Set up the upper county's strength, 573,259 guan 500 wen, and divide by 925 li, obtaining a strength of 619,740. Set up the middle county's 504,983 guan 780 wen and divide by 652 li, obtaining a strength of 774,515. Set up the lower county's 498,760 guan 950 and divide by 465 li, obtaining a strength of 1,072,604, with 90 wen unexhausted. Seek the equal number (deng) of this and the divisor, obtaining 15. Reduce by it, obtaining 6/31.\end{uniten}

\begin{unitzh}\unitlabel{shushu-058.11}置仁郷一十二萬八千三百七十一貫九百八十，如七百里而一，得力一十八萬一千八百三十。置義鄉鄉一十一萬九千四百七十二貫六百文，如七百九十五里而一，得力一十五萬二百八十。置禮鄉一十萬八千四百六十三貫五十文，如七百九十而一，得力一十三萬七千二百九十五。置智鄉八萬四千二百三十六貫二百八十五文，如七百四十九里而一，得力一十一萬二千六百四十五。置信鄉九千三百四十五貫一百六十文，如八百四里而一，得力一萬一千二百六十三，不盡二百六十八文，與法求等，得二百六十八，約為三分之一。\end{unitzh}
\begin{uniten}Set up Ren canton's 128,371 guan 980 and divide by 700 li, obtaining a strength of 181,830. Set up Yi canton's [Translator's note: 鄉 'canton' is written twice] 119,472 guan 600 wen and divide by 795 li, obtaining a strength of 150,280. Set up Li canton's 108,463 guan 50 wen and divide by 790, obtaining a strength of 137,295. Set up Zhi canton's 84,236 guan 285 wen and divide by 749 li, obtaining a strength of 112,645. Set up Xin canton's 9,345 guan 160 wen and divide by 804 li, obtaining a strength of 11,263, with 268 wen unexhausted. Seek the equal number of this and the divisor, obtaining 268. Reducing gives 1/3.\end{uniten}

\begin{unitzh}\unitlabel{shushu-058.12}其下縣、信鄉二處帶母子者，各以母互徧乘八處所得，畢，二處各内本子。上得五千七百六十三萬五千八百二十，中得七千二百二萬九千八百九十五，下得九千九百七十五萬二千一百九十，仁得一千六百九十一萬一百九十，義得一千三百九十七萬六千四十，禮得一千二百七十六萬八千四百三十五，智得一千四十五萬九千二百四十五，信得一百八萬九百七十。已上為三郷五縣變力率。\end{unitzh}
\begin{uniten}As for the two places that carry denominators and numerators (mu zi), the lower county and Xin canton: with each denominator, multiply crosswise over all that was obtained for the 8 places. When this is finished, the two places each add in their own numerators. Upper obtains 57,635,820; Middle obtains 72,029,895; Lower obtains 99,752,190; Ren obtains 16,910,190; Yi obtains 13,976,040; Li obtains 12,768,435; Zhi obtains 10,459,245; Xin obtains 1,080,970. The above are the transformed-strength rates (lü) of the 3 cantons and 5 counties.\end{uniten}

\begin{unitzh}\unitlabel{shushu-058.13}可約者，復求等約之，求得五，故俱以五約之。上得一千一百五十二萬七千一百六十四，中得一千一百四十萬五千九百七十九，下得一千九百九十五萬四百三十八，仁得三十八萬二千五十八{\zhsmall 按：三十八訛五十八。}，義得二百七十九萬五千二百八，禮得二百五十五萬三千六百七十八，智得二百九萬一千八百四十九，信得二十一萬六千一百九十四。已上並為定力。\end{unitzh}
\begin{uniten}For those that can be reduced, again seek the equal number and reduce them. The number found is 5, so reduce all of them by 5. Upper obtains 11,527,164; Middle obtains 11,405,979; Lower obtains 19,950,438; Ren obtains 382,058 [Siku editors: Note: 38 has been miswritten as 58.]; Yi obtains 2,795,208; Li obtains 2,553,678; Zhi obtains 2,091,849; Xin obtains 216,194. The above together are the fixed strengths.\end{uniten}

\begin{unitzh}\unitlabel{shushu-058.14}副併八處定力，得五千六百九十二萬二千五百七十七，為法{\zhsmall 按：法多二十，仁鄉定力多二十故也。}。以合起一萬二千夫遍乘定力訖，上得一千三百八十三億二千五百九十六萬八千為實，中得一千七百二十八億七千一百四十七萬八千為中實，下得二千三百九十四億五百二十五萬六千為下實，{\zhsmall 按}得四百五億八千四百六十九萬六千{\zhsmall 按：多二十四萬，定力多二十故也。}為仁實，義得三百三十五億四千二百四十九萬六千為義寔，禮得三百六億四千四百二十萬萬四千為禮寔，智得三百五十一億二百一十八萬八千為智寔，信得二十五億九千四百三十二萬八千為信寔。已上八寔皆如前法而一，餘分軰之，各得夫數，合。\end{unitzh}
\begin{uniten}Add together in an auxiliary position the fixed strengths of the 8 places, obtaining 56,922,577, as the divisor [Siku editors: Note: the divisor is 20 too many; this is because Ren canton's fixed strength is 20 too many.]. With the total of 12,000 laborers to be raised, multiply over all the fixed strengths. When this is done, Upper obtains 138,325,968,000 as the dividend; Middle obtains 172,871,478,000 as the middle dividend; Lower obtains 239,405,256,000 as the lower dividend; [Siku editors: Note.] [Translator's note: the small print here is only 按 'note', with no further text.] obtains 40,584,696,000 [Siku editors: Note: it is 240,000 too many; this is because the fixed strength is 20 too many.] as the Ren dividend; Yi obtains 33,542,496,000 as the Yi dividend; Li obtains 30,644,204,000 [Translator's note: the text writes 三百六億四千四百二十萬萬四千, with 萬 'ten thousand' doubled] as the Li dividend; Zhi obtains 35,102,188,000 as the Zhi dividend; Xin obtains 2,594,328,000 as the Xin dividend. Divide each of the above 8 dividends by the preceding divisor and apportion the remaining fractions by rank. Each obtains its number of laborers, which agrees.\end{uniten}

\begin{unitzh}\unitlabel{shushu-058.15}均敷徑役\end{unitzh}
\begin{uniten}Equitably apportioning path service (junfu jingyi)\end{uniten}

\begin{unitzh}\unitlabel{shushu-059.1}問：軍戍差坐烽、擺舖，切慮差徑不均。今諸軍共合差一千二百六十人，契勘諸軍見管：前軍六千一百七十人，右軍四千九百三十六人，中軍七千四百四人，左軍三千七百二人，後軍二千四六十八人。各軍合差幾何？答曰：前軍差三百一十五人，右軍差二百五十二人，中軍差三百七十八人，□軍差一百八十九人，後軍差一百二十六人。\end{unitzh}
\begin{uniten}Question: Garrison troops are assigned to be stationed at the beacons and the relay posts (baipu), and there is earnest concern that the assignments (chai jing) may be unequal [Translator's note: 徑 'path' evidently stands for 徭 'labour service'.]. Now the armies together should assign 1260 men. On checking, the armies currently have under command (xianguan): the Front Army 6170 men, the Right Army 4936 men, the Central Army 7404 men, the Left Army 3702 men, the Rear Army 2 thousand 4 68 men [Translator's note: the text reads 二千四六十八, lacking 百 'hundred' after 四.]. How many should each army assign?
Answer: The Front Army assigns 315 men; the Right Army assigns 252 men; the Central Army assigns 378 men; the [blank] Army assigns 189 men; the Rear Army assigns 126 men.\end{uniten}

\begin{unitzh}\unitlabel{shushu-059.2}術曰：以均輸求之。置各軍見管人，騐可求約等，以約之為衰，副併為法。以共合差數乘列衰，各為寔。寔如法而一，各得。\end{unitzh}
\begin{uniten}Method: Solve it by equitable transport (junshu). Set up the men currently under command in each army; test whether a reducing common measure (yue deng) can be found, and reduce them by it to make the proportions (shuai); add them together on the side (fubing) as the divisor (fa). Multiply the listed proportions by the total number to be assigned; each is a dividend (shi). Divide the dividends by the divisor; each is obtained.\end{uniten}

\begin{unitzh}\unitlabel{shushu-059.3}草曰：置諸軍見管，求等，得一千二百三十四。俱以約見管，前軍得五，右軍得四，中軍得六，左軍得三，後軍二，列為各衰。副併諸衰，得二十，為法。以共差一千得二百六十人乘諸衰，前軍得六千三百，右軍得五千四十，中軍得七千五百六十，左軍得三千七百八十，後軍得二千五百二十，各為寔。皆如法二十而一，前軍合差三百一十五人，右軍合差二百五十二人，中軍合差三百七十二人，左軍合差一百八十九人，後軍合差一百二十六人。\end{unitzh}
\begin{uniten}Working: Set up what the armies currently have under command and find the common measure (deng), obtaining 1234. Reduce all those under command by it: the Front Army obtains 5, the Right Army obtains 4, the Central Army obtains 6, the Left Army obtains 3, the Rear Army 2; arrange them as the respective proportions. Add all the proportions together on the side, obtaining 20, as the divisor. Multiply all the proportions by the total assignment, 1 thousand obtain 260 men [Translator's note: the text reads 一千得二百六十; 得 'obtain' stands intrusively within the number.]: the Front Army obtains 6300, the Right Army obtains 5040, the Central Army obtains 7560, the Left Army obtains 3780, the Rear Army obtains 2520; each is a dividend. Divide all of them by the divisor 20: the Front Army should assign 315 men, the Right Army should assign 252 men, the Central Army should assign 372 men, the Left Army should assign 189 men, the Rear Army should assign 126 men.\end{uniten}

\begin{unitzh}\unitlabel{shushu-059.4}《數學九章》卷八下\end{unitzh}
\begin{uniten}Shuxue jiuzhang (Mathematics in Nine Sections), juan 8, part 2.\end{uniten}

\begin{unitzh}\unitlabel{shushu-059.5}子部·天文算法類·算書之屬·《數學九章》\end{unitzh}
\begin{uniten}Masters branch (zibu), astronomy and calculation category (tianwen suanfa lei), class of mathematical books (suanshu zhi shu), Shuxue jiuzhang.\end{uniten}

\section*{卷9上}\addcontentsline{toc}{section}{卷9上}

\begin{unitzh}\unitlabel{shushu-060.1}欽定四庫全書\end{unitzh}
\begin{uniten}Imperially Authorized Complete Library of the Four Treasuries (Qinding Siku quanshu)\end{uniten}

\begin{unitzh}\unitlabel{shushu-060.2}數學九章卷九上 宋 秦九韶 撰\end{unitzh}
\begin{uniten}Mathematical Treatise in Nine Sections (Shuxue jiuzhang), juan 9, part 1. Compiled by Qin Jiushao of the Song.\end{uniten}

\begin{unitzh}\unitlabel{shushu-060.3}市易\end{unitzh}
\begin{uniten}Market Exchange (shiyi)\end{uniten}

\begin{unitzh}\unitlabel{shushu-060.4}積木計餘\end{unitzh}
\begin{uniten}Stacked Timber: Reckoning the Remainder\end{uniten}

\begin{unitzh}\unitlabel{shushu-060.5}問：原管杉木一尖垜，偶不記數。從上取用至中間，見存九條為靣。問原木及見存各幾何？荅曰：原木一百五十三條，見存木一百一十七條。\end{unitzh}
\begin{uniten}Question: Originally one pointed stack (jianduo) of fir logs was in custody; it happened that the number was not recorded. They were taken for use from the top down to the middle, and now 9 logs remain, forming the face. What are the original timber and the timber now remaining, each? Answer: Original timber, 153 logs; timber now remaining, 117 logs.\end{uniten}

\begin{unitzh}\unitlabel{shushu-060.6}術曰：商功求之，堆積入之。倍中靣，副置，減一，以乘其副，得數半之，為原木。副置上層，減一，以乘其副，得數半之，用減原木，餘為見存。{\zhsmall 其非中一層數者，各以自地上至面層數立數求之。}\end{unitzh}
\begin{uniten}Method: Seek it by construction works (shanggong); enter it by pile volumes (duiji). Double the middle face, set it up in duplicate (fuzhi), subtract 1, and multiply its duplicate by it; halve the number obtained: this is the original timber. Set up the upper layer in duplicate, subtract 1, and multiply its duplicate by it; halve the number obtained, and use it to subtract from the original timber; the remainder is what now remains. [Note: Where it is not the number of the middle layer, in each case set up the number using the number of layers from the ground up to the face layer, and seek it.]\end{uniten}

\begin{unitzh}\unitlabel{shushu-060.7}草曰：倍中靣九條，得一十八。副置，減一，餘一十七，以乘副一十八，得三百六條，以半之，得一百五十三條，為原木之數。副置中靣九條，減一，餘八，以乘副九，得七十二，以半之，得三十六，以減原木一百五十三，餘一百一十七條，為見存木數。合問。\end{unitzh}
\begin{uniten}Working: Double the middle face, 9 logs, giving 18. Set it up in duplicate and subtract 1, leaving 17; multiply the duplicate 18 by it, giving 306 logs; halve it, giving 153 logs: this is the number of the original timber. Set up the middle face, 9 logs, in duplicate; subtract 1, leaving 8; multiply the duplicate 9 by it, giving 72; halve it, giving 36; subtract it from the original timber 153, leaving 117 logs: this is the number of timber now remaining. This agrees with the question.\end{uniten}

\begin{unitzh}\unitlabel{shushu-060.8}按：此即一靣平堆形。中層為九，上下必各有八層，共十七層，即原尖堆形；上八層，即用過尖堆形。其義甚明。舊餘木圖今刪。\end{unitzh}
\begin{uniten}Note: This is precisely the shape of a flat pile with one face. The middle layer being 9, above and below there must be 8 layers each, 17 layers in all: this is the shape of the original pointed pile; the upper 8 layers are the shape of the pointed pile that has been used. Its meaning is very clear. The former diagram of the remaining timber is now deleted.\end{uniten}

\begin{unitzh}\unitlabel{shushu-060.9}竹圍蘆束\end{unitzh}
\begin{uniten}Bamboo Circumferences and Reed Bundles\end{uniten}

\begin{unitzh}\unitlabel{shushu-060.10}問：受給場交收竹二千三百七十四把，内筀竹一千一百五十一把，每把外圍三十六竿；水竹一千二百二十三把，每把外圍四十二竿；蘆三千六十五束，每束圍五尺。其蘆原様五尺五寸，今納到圍小。合凖原蘆幾束？及水、筀竹各幾何？荅曰：筀竹一十四萬六千一百七十七竿，水竹二十萬六千六百八十七竿，合凖原蘆二千五百三十三束{\zhsmall 一百二十一分束之七}。\end{unitzh}
\begin{uniten}Question: A receiving-and-issuing yard (shougei chang) took delivery of 2374 bunches (ba) of bamboo: among them, gui bamboo (guizhu), 1151 bunches, each bunch with an outer circumference of 36 stalks (gan); water bamboo (shuizhu), 1223 bunches, each bunch with an outer circumference of 42 stalks; reeds, 3065 bundles (shu), each bundle 5 chi (尺) in circumference. The original specification for these reeds is 5 chi 5 cun (寸); those now delivered are smaller in circumference. To how many bundles of the original reeds are they equivalent? And how many are each of the water and gui bamboo? Answer: Gui bamboo, 146,177 stalks; water bamboo, 206,687 stalks; equivalent to the original reeds, 2533 bundles [Note: 7/121 bundle].\end{uniten}

\begin{unitzh}\unitlabel{shushu-060.11}術曰：以方田及圓率求之。置原束差，併竹外圍竹數，以乘外圍，又乘把數，為竹實；倍圓束差，為竹法，除之，各得二竹數，皆以把數為心加入，各得竹條數。置蘆圍尺數自乘，以乘蘆束，為蘆實；以蘆原尺數自乘，為蘆法，除實，得所凖蘆束數。\end{unitzh}
\begin{uniten}Method: Seek it by rectangular fields (fangtian) and the circle rate (yuanlü). Set up the original bundle difference (yuanshu cha), add the number of stalks in the outer circumference of the bamboo, multiply by the outer circumference, and multiply again by the number of bunches, as the bamboo dividend (shi); double the round-bundle difference (yuanshu cha) as the bamboo divisor (fa); divide, and each gives the number for the two bamboos; in each case add in the number of bunches as the centres, and each gives the number of bamboo stems. Set up the number of chi of the reed circumference, multiply it by itself, and multiply by the reed bundles, as the reed dividend; multiply the original number of chi of the reeds by itself, as the reed divisor; divide the dividend, giving the number of reed bundles to which they are equivalent.\end{uniten}

\begin{unitzh}\unitlabel{shushu-060.12}草曰：置圍束差六，併筀竹外圍三十六竿，得四十二竿，以乘外圍三十六竿，得一千五百一十二竿，又乘筀竹一千一百五十一把，得一百七十四萬三百一十二竿，為筀竹實；倍圓束差六，得一十二，為竹法，除實，得一十四萬五千二十六竿，以把數一千一百五十一併之，得一十四萬六千一百七十七竿，為筀竹。又置原差六，併水竹外圍四十二竿，得四十八竿，以乘水外圍四十二竿，得二千一十六竿，又乘水竹一千二百二十三把，得二百四十六萬五千五百六十八竿，為水竹實；亦以竹法一十二除之，得二十萬五千四百六十四竿，以水竹把數一千二百二十三併之，得二十萬六千六百八十七竿，為水竹數。次置蘆圍五尺，通為五十寸，以自乘，得二千五百寸，又乘蘆束數三千六十五，得七百六十六萬二千五百寸，為蘆實；以原様蘆圍五尺五寸，亦通為五十五寸，以自乗，得三千二十五寸，為蘆法，除實，得二千五百三十三束，不盡一百七十五寸，與法求等，得二十五，俱以約之，得一百二十一分束之七，為蘆二千五百三十三束一百二十一分束之七。合問。\end{unitzh}
\begin{uniten}Working: Set up the circumference-bundle difference (weishu cha) 6, add the outer circumference of the gui bamboo, 36 stalks, giving 42 stalks; multiply by the outer circumference, 36 stalks, giving 1512 stalks; multiply again by the gui bamboo, 1151 bunches, giving 1,740,312 stalks, as the gui bamboo dividend. Double the round-bundle difference 6, giving 12, as the bamboo divisor; divide the dividend, giving 145,026 stalks; add to it the number of bunches, 1151, giving 146,177 stalks: this is the gui bamboo. Again set up the original difference (yuan cha) 6, add the outer circumference of the water bamboo, 42 stalks, giving 48 stalks; multiply by the water [bamboo] outer circumference, 42 stalks, giving 2016 stalks; multiply again by the water bamboo, 1223 bunches, giving 2,465,568 stalks, as the water bamboo dividend. Likewise divide it by the bamboo divisor 12, giving 205,464 stalks; add to it the number of water bamboo bunches, 1223, giving 206,687 stalks: this is the number of water bamboo. Next set up the reed circumference 5 chi, convert it into 50 cun, multiply it by itself, giving 2500 cun; multiply again by the number of reed bundles, 3065, giving 7,662,500 cun, as the reed dividend. Take the original-specification reed circumference 5 chi 5 cun, likewise convert it into 55 cun, multiply it by itself, giving 3025 cun, as the reed divisor; divide the dividend, giving 2533 bundles, with 175 cun not exhausted (bujin); seek the common measure (deng) with the divisor, giving 25; reduce both by it, giving 7/121 bundle: this is the reeds, 2533 7/121 bundles. This agrees with the question.\end{uniten}

\begin{unitzh}\unitlabel{shushu-060.13}寄倉知總{\zhsmall 按：舊本此問無題，今増。}\end{unitzh}
\begin{uniten}Lodging in Granaries: Knowing the Total [Siku editors: In the old edition this problem has no title; it is now added.]\end{uniten}

\begin{unitzh}\unitlabel{shushu-060.14}問：和糴米運，借倉權頓，計五十厫，毎厫濶一丈五尺，深三丈，米高一丈二尺；又借寺屋四十間，内二十五間濶一丈二尺，深二丈五尺，米高一丈；内一十五間各濶一丈三尺，深三丈，米高一丈二尺。欲知寺屋及倉容米共計幾何？答曰：共計米一十六萬六千八十石。倉五十厫，米一十萬八千石；寺屋四十間，米五萬八千八十石。\end{unitzh}
\begin{uniten}Question: Rice of a government purchase (hedi) is being shipped; granaries are borrowed to lodge it temporarily, 50 storehouses (ao) in all, each storehouse 1 zhang (丈) 5 chi wide and 3 zhang deep, the rice 1 zhang 2 chi high. Also 40 rooms (jian) of temple buildings are borrowed: of these, 25 rooms are 1 zhang 2 chi wide and 2 zhang 5 chi deep, the rice 1 zhang high; of these, 15 rooms are each 1 zhang 3 chi wide and 3 zhang deep, the rice 1 zhang 2 chi high. One wishes to know how much rice the temple buildings and the granaries hold in total. Answer: In total, rice 166,080 dan (石). The granary, 50 storehouses: rice 108,000 dan; the temple buildings, 40 rooms: rice 58,080 dan.\end{uniten}

\begin{unitzh}\unitlabel{shushu-060.15}術曰：啇功求之。置厫并屋深、濶、米高相乗，併之為實，如斛法而一。\end{unitzh}
\begin{uniten}Method: Seek it by construction works. Set up the depth, width and rice height of the storehouses and of the rooms, multiply them together, add them up as the dividend (shi), and divide by the hu (斛) divisor (hufa).\end{uniten}

\begin{unitzh}\unitlabel{shushu-060.16}草曰：先以厫深三丈，通為三十尺，乗濶十五尺，得四百五十尺，又乗高一十二尺，得五千四百尺，以乗五十厫，得二十七萬尺，為實；以斛法二尺五寸除之，得一十萬八千石，為倉五寸厫共容米。次置寺屋深二十五尺，乗濶一十二尺，得三百尺，乗米高一十尺，得三千尺，以二十五間乗之，得七萬五千尺，於上。次置深三十尺，乗濶一十三尺，得三百九十尺，又乗米高十十二尺，得四千六百八十尺，以乗一十五間，得七萬二百尺，加上，共得一十四萬五千二百尺，以斛法二尺五寸除之，得五萬八千八十石，為寺屋四十間共容米。以併厫米，共得一十六萬六千八十石，為共和糴到米。\end{unitzh}
\begin{uniten}Working: First take the storehouse depth 3 zhang, convert it into 30 chi, multiply by the width 15 chi, giving 450 chi; multiply again by the height 12 chi, giving 5400 chi; multiply by the 50 storehouses, giving 270,000 chi, as the dividend. Divide it by the hu divisor 2 chi 5 cun, giving 108,000 dan: this is the rice held altogether by the granary's 5 cun [Translator's note: 五寸 "5 cun" is evidently an error for 五十 "50"] storehouses. Next set up the temple-building depth 25 chi, multiply by the width 12 chi, giving 300 chi; multiply by the rice height 10 chi, giving 3000 chi; multiply it by the 25 rooms, giving 75,000 chi, placed above. Next set up the depth 30 chi, multiply by the width 13 chi, giving 390 chi; multiply again by the rice height of "ten ten two" chi [Translator's note: 十十二 so in the text; evidently for 一十二, 12], giving 4680 chi; multiply by the 15 rooms, giving 70,200 chi; add it to the above, giving 145,200 chi in all; divide it by the hu divisor 2 chi 5 cun, giving 58,080 dan: this is the rice held altogether by the 40 rooms of temple buildings. Add it to the storehouse rice, giving 166,080 dan in all: this is the total government-purchased rice that has arrived.\end{uniten}

\begin{unitzh}\unitlabel{shushu-060.17}方圓同積{\zhsmall 按：舊本此問無題，今増。}\end{unitzh}
\begin{uniten}Square and Round of Equal Volume [Siku editors: In the old edition this problem has no title; it is now added.]\end{uniten}

\begin{unitzh}\unitlabel{shushu-060.18}問：有圓囤米二十五箇，内有大囤一十二箇，上徑一丈，下徑九尺，高一丈二尺；小囤一十三箇，上徑九尺，下徑八尺，高一丈。今出租斗一隻，口方九寸六分，底方七寸，正深四寸，並裹明凖尺。令凖數造五斗方斛及圓斛各二隻，湏令二斛口徑、正深大小不同。各得多少？及囤積米幾何？答曰：方斛一隻，口方六寸四分，底方一尺二寸，深一尺五寸九分二厘；又一隻，口方一尺，底方一尺二寸，深一尺一寸四分五厘。圓斛一隻，口徑一尺二寸七分，底徑一尺二寸，深一尺二寸一分四厘；又一隻，口徑一尺三寸，底徑一尺二寸，深一尺一寸八分五厘。囤米計八千六十七石四升七合四勺一抄八撮。\end{unitzh}
\begin{uniten}Question: There are round bins (dun) of rice, 25 in number; among them there are 12 large bins, upper diameter 1 zhang, lower diameter 9 chi, height 1 zhang 2 chi; and 13 small bins, upper diameter 9 chi, lower diameter 8 chi, height 1 zhang. Now one rent dou (zudou) is issued: mouth square 9 cun 6 fen (分), bottom square 7 cun, true depth 4 cun, all clear inside measures by the standard chi [Translator's note: 裹 "wrap" is presumably for 裏 "inside"]. Let square hu and round hu of 5 dou be made according to the standard number, 2 of each; it is required that the two hu differ in the size of their mouth diameters and true depths. How much does each come to? And how much rice is stored in the bins? Answer: Square hu, one: mouth square 6 cun 4 fen, bottom square 1 chi 2 cun, depth 1 chi 5 cun 9 fen 2 li (厘); another one: mouth square 1 chi, bottom square 1 chi 2 cun, depth 1 chi 1 cun 4 fen 5 li. Round hu, one: mouth diameter 1 chi 2 cun 7 fen, bottom diameter 1 chi 2 cun, depth 1 chi 2 cun 1 fen 4 li; another one: mouth diameter 1 chi 3 cun, bottom diameter 1 chi 2 cun, depth 1 chi 1 cun 8 fen 5 li. The bin rice amounts to 8067 dan 4 sheng (升) 7 ge (合) 4 shao (勺) 1 chao (抄) 8 cuo (撮).\end{uniten}

\begin{unitzh}\unitlabel{shushu-060.19}按：共囤米數誤，應得二千零一十六石七斗六升一合八勺五抄。草内少一四歸，故差多三倍。\end{unitzh}
\begin{uniten}Note: The number of the total bin rice is wrong; it should be 2016 dan 7 dou 6 sheng 1 ge 8 shao 5 chao. In the working one division by 4 (sigui) is missing; therefore the discrepancy is 3 times too much.\end{uniten}

\begin{unitzh}\unitlabel{shushu-060.20}術曰：以啇功及少廣求之。置出斗上下方相乗之，又各自乗，併之，乗深，又以五斗乗之，為積于上。求方斛：先自如意立數為斛深，又如意立數為底方，置深為從隅，以底方乗隅為從方，又以底乗從方為減率，以減上積，餘為實，開連枝平方，得方斛口方。不盡，以所得數為基，增損求之，以口底方相乗，又各自求，併之為法，除前上積，得深。餘分収棄之。求圖斛：置四數，以因前積為寄，如意立數為斛深，别如意立數為底徑，以三因深為從隅，以底徑乗隅為從方，以底徑乗從方為減率，以減寄，餘為實，開連枝平方，得口徑。不盡，以所得數為基，如意求差，以口底徑相乗，又各自乗，併之為法，除寄，得深。餘分収棄之。求囤米：置各囤上徑、下徑相乗，又各自乗，併之，乗高，又乗囤數，所得之數為積{\zhsmall 囤有大小，以類}，併之為共積，如四而一為實，以斛求之，得米。\end{unitzh}
\begin{uniten}Method: Seek it by construction works and diminishing breadth (shaoguang). Set up the upper and lower squares of the issued dou and multiply them together; also multiply each by itself; add them, multiply by the depth, and again multiply it by 5 dou, as the volume, placed above. To find the square hu: first set up a number of one's own choosing as the depth of the hu, and again set up a number of one's choosing as the bottom square; set up the depth as the accompanying corner (congyu); multiply the corner by the bottom square as the accompanying square (congfang); and again multiply the accompanying square by the bottom as the subtraction rate (jianlü); subtract it from the upper volume, and the remainder is the dividend (shi); extract the connected-branch square root (lianzhi pingfang), giving the mouth square of the square hu. If it is not exhausted, take the number obtained as the base and seek it by increasing or decreasing; multiply the mouth and bottom squares together, and also seek each by itself [Translator's note: 求 "seek" is evidently for 乗 "multiply"]; add them as the divisor (fa); divide the previous upper volume, giving the depth. The remaining fractions: take them in or discard them. To find the diagram hu [Translator's note: 圖 "diagram" is evidently for 圓 "round"]: set up the number 4 and multiply the previous volume by it, as the lodged number (ji); set up a number of one's choosing as the depth of the hu; separately set up a number of one's choosing as the bottom diameter; multiply the depth by 3 as the accompanying corner; multiply the corner by the bottom diameter as the accompanying square; multiply the accompanying square by the bottom diameter as the subtraction rate; subtract it from the lodged number, and the remainder is the dividend; extract the connected-branch square root, giving the mouth diameter. If it is not exhausted, take the number obtained as the base and seek the difference at will; multiply the mouth and bottom diameters together, and also multiply each by itself; add them as the divisor; divide the lodged number, giving the depth. The remaining fractions: take them in or discard them. To find the bin rice: set up the upper diameter and lower diameter of each bin and multiply them together; also multiply each by itself; add them, multiply by the height, and multiply again by the number of bins; the number obtained is the volume [Note: the bins are large and small; by kind]; add them as the total volume; divide by 4, as the dividend; seek it by the hu, giving the rice.\end{uniten}

\begin{unitzh}\unitlabel{shushu-061.1}草曰：置出租斗口方九寸六分，與底方七寸相乗，得六十七寸二分於上。又於口方九寸六分自乗之，得九十二寸一分六厘，加上。又以底方七寸自乗，得四十九寸，又加上，共得二百八寸三分六厘。乗深四寸，得八百三十三寸四分四厘。又以五斗乗之，得四千一百六十七寸二分，為三叚斛積于上。\end{unitzh}
\begin{uniten}Working: Set up the mouth square (koufang) of the rent-paying dou (chuzu dou 出租斗), 9 cun (寸) 6 fen (分), and multiply it together with the bottom square (difang), 7 cun, obtaining 67 cun 2 fen, [placed] above. Again take the mouth square, 9 cun 6 fen, and multiply it by itself, obtaining 92 cun 1 fen 6 li (厘); add it to the above. Again multiply the bottom square, 7 cun, by itself, obtaining 49 cun; again add it to the above; altogether obtaining 208 cun 3 fen 6 li. Multiply by the depth, 4 cun, obtaining 833 cun 4 fen 4 li. Again multiply this by 5 dou, obtaining 4167 cun 2 fen, as 3 portions (duan) of the hu (斛) volume, [placed] above.\end{uniten}

\begin{unitzh}\unitlabel{shushu-061.2}求方斛：如意立一尺六寸為斛深，又加意立一尺二寸為斛底。以深一十六寸為從隅，以底一十二寸乗隅，得一百九十二寸，為從方。又以底一十二寸乗從方一百九十二寸，得二千三百四寸，為減率。以減上積四千一百六十七寸二分，餘一千八百六十三寸二分，為實。開連枝平方，得六寸三分五厘，為基。其積不及一寸一分六厘，係有虧數，其基數為米可用，湏合損益基數。今益作六寸四分為口方，以原立一尺二寸為底方。以口方乗底方，得七十六寸八分以上；又以口方六寸四分自乗，得四十寸九分六厘；又以底方十二寸自乗，得一百四十四寸；併以上，共得二百六十一寸七分六厘，為法。以除前積四千一百六十七寸二分，得一尺五寸九分二厘，為方斛深。其積不及一厘九毫二絲，収為閏。\end{unitzh}
\begin{uniten}To find the square hu: at will set up 1 chi (尺) 6 cun as the depth of the hu; again deliberately set up 1 chi 2 cun as the bottom of the hu. Take the depth, 16 cun, as the following corner (congyu); multiply the corner by the bottom, 12 cun, obtaining 192 cun, as the following square (congfang). Again multiply the following square, 192 cun, by the bottom, 12 cun, obtaining 2304 cun, as the subtracting rate (jianlü). Subtract it from the volume above, 4167 cun 2 fen, leaving 1863 cun 2 fen as the dividend (shi). Extract the connected-branch square root (kai lianzhi pingfang), obtaining 6 cun 3 fen 5 li as the base (ji). Its volume falls short by 1 cun 1 fen 6 li; this involves a deficient number (kuishu); its base number can be used for rice [Translator's note: 米 "rice" is probably an error for 未 "not yet".]; one must decrease or increase the base number. Now increase it, making 6 cun 4 fen the mouth square, and take the originally set 1 chi 2 cun as the bottom square. Multiply the bottom square by the mouth square, obtaining 76 cun 8 fen, [placed] above; again multiply the mouth square, 6 cun 4 fen, by itself, obtaining 40 cun 9 fen 6 li; again multiply the bottom square, 12 cun, by itself, obtaining 144 cun; add them to the above, altogether obtaining 261 cun 7 fen 6 li as the divisor (fa). With it divide the previous volume, 4167 cun 2 fen, obtaining 1 chi 5 cun 9 fen 2 li as the depth of the square hu. Its volume falls short by 1 li 9 hao (毫) 2 si (絲); gather it in as surplus (run).\end{uniten}

\begin{unitzh}\unitlabel{shushu-061.3}又累増至一十寸為口方，仍以一十二寸為底方。乃以口方一十寸乗底方一十二寸，得一百二十寸於上；又口方自乗，得一百寸，加上；又以底方自乗，得一百四十四寸，又加上，共得三百六十四寸，為法。亦除前實積四千一百六十七寸二分，得一十一寸四分五厘，為方斛深。其積不及六分，収為閏。此是求出両等斛數，在人擇而用之。\end{unitzh}
\begin{uniten}Again, increase step by step up to 10 cun as the mouth square, still taking 12 cun as the bottom square. Then multiply the bottom square, 12 cun, by the mouth square, 10 cun, obtaining 120 cun, [placed] above; again multiply the mouth square by itself, obtaining 100 cun; add it to the above; again multiply the bottom square by itself, obtaining 144 cun; again add it to the above; altogether obtaining 364 cun as the divisor. Likewise divide the previous dividend volume, 4167 cun 2 fen, obtaining 11 cun 4 fen 5 li as the depth of the square hu. Its volume falls short by 6 fen; gather it in as surplus. This is finding two kinds of hu numbers; it is up to people to choose and use them.\end{uniten}

\begin{unitzh}\unitlabel{shushu-061.4}求圓斛：置四數，以因前積四千一百六十七寸二分，得一萬六千六百六十八寸八分，為寄。如意立一尺二寸為圓斛深，又如意立一尺為底徑。以三因深，得三十六寸，為從隅；以底一十寸乗隅，得三百六十寸，為從方；又以底一十寸乗從方，得三千六百寸，為减率。以減寄一萬六千六百六十八寸八分，餘一萬三千六十八寸八分，為實。開連枝平方，得一尺四寸七分，為基。其實不及二寸四分四厘，収為閏。\end{unitzh}
\begin{uniten}To find the round hu: set up the number 4 and with it multiply (yin) the previous volume, 4167 cun 2 fen, obtaining 16668 cun 8 fen, as the lodged [value] (ji). At will set up 1 chi 2 cun as the depth of the round hu; again at will set up 1 chi as the bottom diameter (dijing). Multiply the depth by 3, obtaining 36 cun, as the following corner; multiply the corner by the bottom, 10 cun, obtaining 360 cun, as the following square; again multiply the following square by the bottom, 10 cun, obtaining 3600 cun, as the subtracting rate. Subtract it from the lodged [value], 16668 cun 8 fen, leaving 13068 cun 8 fen as the dividend. Extract the connected-branch square root, obtaining 1 chi 4 cun 7 fen as the base. Its dividend falls short by 2 cun 4 fen 4 li; gather it in as surplus.\end{uniten}

\begin{unitzh}\unitlabel{shushu-061.5}次以原立底徑一尺併基一尺四寸七分，得二尺四寸七分，只減七分為差，以餘三尺四寸以半之，得一尺二寸，為底徑。以差七分併底徑，得一尺二寸七分，為口徑。始以口徑一尺二寸七分乗底徑一尺二寸，得一百五十二寸四分於上；次以口徑自乗，得一百六十一寸二分九厘，加上；又以底徑自乗，得一百四十四寸，又加上，共得四百五十七寸六分九厘。以三因之，得一千三百七十三寸七厘，為法。除前寄一萬六千六百六十八寸八分，得一尺二寸一分四厘，為圓斛正深。其實不及二毫六絲九忽八微，収為閏。\end{unitzh}
\begin{uniten}Next, add the originally set bottom diameter, 1 chi, and the base, 1 chi 4 cun 7 fen, obtaining 2 chi 4 cun 7 fen; subtract only 7 fen as the difference (cha); take the remainder, 3 chi 4 cun, and halve it, obtaining 1 chi 2 cun as the bottom diameter. Add the difference, 7 fen, to the bottom diameter, obtaining 1 chi 2 cun 7 fen as the mouth diameter (koujing). First multiply the bottom diameter, 1 chi 2 cun, by the mouth diameter, 1 chi 2 cun 7 fen, obtaining 152 cun 4 fen, [placed] above; next multiply the mouth diameter by itself, obtaining 161 cun 2 fen 9 li; add it to the above; again multiply the bottom diameter by itself, obtaining 144 cun; again add it to the above; altogether obtaining 457 cun 6 fen 9 li. Multiply it by 3, obtaining 1373 cun 7 li, as the divisor. Divide the previous lodged [value], 16668 cun 8 fen, obtaining 1 chi 2 cun 1 fen 4 li as the true depth (zhengshen) of the round hu. Its dividend falls short by 2 hao 6 si 9 hu (忽) 8 wei (微); gather it in as surplus.\end{uniten}

\begin{unitzh}\unitlabel{shushu-061.6}又以基一尺四寸七分增三分，得一尺五寸，併底徑一尺，得二尺五寸，減一寸為差，餘二尺四寸，以半之，得一尺二寸，為底徑。以差一寸併底徑一尺二寸，得一尺三寸，為口徑。始以口徑一十三寸乗底徑一尺二寸，得一百五十六寸於上；又以口徑一十三寸自乗，得一百六十九寸，加上；又以底徑一十二寸自乗，得一百四十四寸，又加上，共得四百六十九寸。以三因之，得一千四百七寸，為法。除前寄一萬六千六百六十八寸八分，得一尺一寸八分四厘七毫，為圓斛深。寄餘七厘一毫，却収深七毫作一厘，通得一尺一寸八分五厘，為圓斛深。此是求出両等圓斛，在人擇而用之。\end{unitzh}
\begin{uniten}Again, increase the base, 1 chi 4 cun 7 fen, by 3 fen, obtaining 1 chi 5 cun; add the bottom diameter, 1 chi, obtaining 2 chi 5 cun; subtract 1 cun as the difference, leaving 2 chi 4 cun; halve it, obtaining 1 chi 2 cun as the bottom diameter. Add the difference, 1 cun, to the bottom diameter, 1 chi 2 cun, obtaining 1 chi 3 cun as the mouth diameter. First multiply the bottom diameter, 1 chi 2 cun, by the mouth diameter, 13 cun, obtaining 156 cun, [placed] above; again multiply the mouth diameter, 13 cun, by itself, obtaining 169 cun; add it to the above; again multiply the bottom diameter, 12 cun, by itself, obtaining 144 cun; again add it to the above; altogether obtaining 469 cun. Multiply it by 3, obtaining 1407 cun, as the divisor. Divide the previous lodged [value], 16668 cun 8 fen, obtaining 1 chi 1 cun 8 fen 4 li 7 hao as the depth of the round hu. The lodged [value] leaves a remainder of 7 li 1 hao; then gather in the 7 hao of the depth, making it 1 li; in all obtaining 1 chi 1 cun 8 fen 5 li as the depth of the round hu. This is finding two kinds of round hu; it is up to people to choose and use them.\end{uniten}

\begin{unitzh}\unitlabel{shushu-061.7}求囤米：置大囤上徑一丈，通為百寸，乗下徑九十寸，得九千於上；又以上徑自乗，得一萬寸，加上；又以下徑九十寸自乗，得八千一百寸，加上，共得二萬七千一百寸。乗高一百二十寸，得三百二十五萬二千寸。又乗大囤一十二個，得三千九百二萬四千寸，為寄。\end{unitzh}
\begin{uniten}To find the rice in the bins (dun): set up the upper diameter of a large bin, 1 zhang (丈), converted (tong) into 100 cun, and multiply it by the lower diameter, 90 cun, obtaining 9000, [placed] above; again multiply the upper diameter by itself, obtaining 10000 cun; add it to the above; again multiply the lower diameter, 90 cun, by itself, obtaining 8100 cun; add it to the above; altogether obtaining 27100 cun. Multiply by the height, 120 cun, obtaining 3252000 cun. Again multiply by the 12 large bins, obtaining 39024000 cun, as the lodged [value].\end{uniten}

\begin{unitzh}\unitlabel{shushu-061.8}次置小囤上徑九十寸、下徑八十寸，相乗，得七千二百寸於次；又上徑自乗，得八千一百寸，加次；又下徑自乗，得六千四百寸，加次，共得二萬一千七百寸。又乗高一百寸，得二百一十七萬寸。又乗小囤一十三箇，得二千八百二十一萬寸。併寄，共得六千七百二十三萬四千寸，為實{\zhsmall 按：應四歸為實，草中遺漏，故得數誤。}。倍前斛積四千一百六十七寸二分為法，除之，得八千六十七石四升七合四勺一抄四撮。\end{unitzh}
\begin{uniten}Next set up the upper diameter of a small bin, 90 cun, and the lower diameter, 80 cun; multiply them together, obtaining 7200 cun, [placed] in the next [position]; again multiply the upper diameter by itself, obtaining 8100 cun; add it to the next [position]; again multiply the lower diameter by itself, obtaining 6400 cun; add it to the next [position]; altogether obtaining 21700 cun. Again multiply by the height, 100 cun, obtaining 2170000 cun. Again multiply by the 13 small bins, obtaining 28210000 cun. Add the lodged [value], altogether obtaining 67234000 cun as the dividend [Siku editors: It should be divided by 4 (sigui) to make the dividend; in the working this was omitted, therefore the number obtained is wrong.]. Double the previous hu volume, 4167 cun 2 fen, as the divisor; divide by it, obtaining 8067 dan (石) 4 sheng (升) 7 ge (合) 4 shao (勺) 1 chao (抄) 4 cuo (撮).\end{uniten}

\begin{unitzh}\unitlabel{shushu-061.9}按：求方斛積法，以上下徑相乗，又各自乗，併，而以深再乗，三除之，得積。求徑深法：三因積。有二徑，以二徑相乗、各自乗數併而除之，得深。有一深一徑，以深除積，得數内減徑自來，餘為實，徑為縱，開帶縱平方，得又一徑。或徑自乗，深再乗，减積除為實，以深為縱隅，深徑相乗為縱方，開連枝平方，得又一徑。\end{unitzh}
\begin{uniten}Note: The method for finding the volume of the square hu: multiply the upper and lower diameters together, and also multiply each by itself; add them, and multiply again by the depth; divide by 3, obtaining the volume. The method for finding diameters and depth: multiply the volume by 3. If there are two diameters, add the numbers of the two diameters multiplied together and of each multiplied by itself, and divide by this, obtaining the depth. If there is one depth and one diameter, divide the volume by the depth; from the number obtained subtract the diameter "self-come" (zilai) [Translator's note: 自來 is probably an error for 自乗, "multiplied by itself".]; the remainder is the dividend, and the diameter is the lengthwise [term] (zong); extract the square root with attached lengthwise [term] (kai daizong pingfang), obtaining the other diameter. Or multiply the diameter by itself and multiply again by the depth, subtract it from the volume, removing (chu) it, as the dividend; take the depth as the lengthwise corner (zongyu), and the depth and diameter multiplied together as the lengthwise square (zongfang); extract the connected-branch square root, obtaining the other diameter.\end{uniten}

\begin{unitzh}\unitlabel{shushu-061.10}草中求斗積不加三除，為三倍斗積；五因之，為三倍斛積。故設正深、底徑二數，開連枝平方，以求口徑；既得口徑，復設二徑數，以求正深也。求圓斛徑，即如求圓外切方邊，當以方圓冪率變圓積為方積，故四因前積，而以三因正深代三除也。求圓囤米尺積，先用方斛求積法，次變方為圓。方斛法用三除，變圓法用三因四除，合之則三因四除，合之則三除並可省，惟以四除之，卽得圓囤積。術中詳言之，而草中歩算遺漏四除，故得數誤為四倍也。\end{unitzh}
\begin{uniten}In the working, in finding the dou volume, the division by 3 is not applied, making 3 times the dou volume; multiplying it by 5 makes 3 times the hu volume. Therefore one posits the two numbers, the true depth and the bottom diameter, and extracts the connected-branch square root in order to find the mouth diameter; having obtained the mouth diameter, one again posits the numbers of the two diameters in order to find the true depth. Finding the diameter of the round hu is just like finding the side of the square circumscribed about a circle: one should use the rates (lü) of the areas (mi) of square and circle to change the round volume into a square volume; therefore one multiplies the previous volume by 4, and uses multiplication of the true depth by 3 in place of the division by 3. To find the volume in chi (chiji) of the rice in round bins, first use the method for finding the volume of the square hu, and next change the square into a round. The square-hu method uses division by 3, and the method of changing to round uses multiplication by 3 and division by 4; combining them, then multiplication by 3 and division by 4; combining them, then the division by 3 can likewise be omitted [Translator's note: the text appears garbled here.]; only divide by 4, and one obtains the volume of the round bins. The method (shu) states this in detail, but in the working the reckoning (busuan) omitted the division by 4; therefore the number obtained is erroneously fourfold.\end{uniten}

\begin{unitzh}\unitlabel{shushu-061.11}子部·天文算法類·算書之屬·《數學九章》·卷九上 子部·天文算法類·算書之屬·《數學九章》·卷九上 子部·天文算法類·算書之屬·《數學九章》·卷九上 子部·天文算法類·算書之屬·《數學九章》·卷九上 子部·天文算法類·算書之屬·《數學九章》·卷九上 子部·天文算法類·算書之屬·《數學九章》·卷九上 子部·天文算法類·算書之屬·《數學九章》·卷九上 子部·天文算法類·算書之屬·《數學九章》·卷九上 子部·天文算法類·算書之屬·《數學九章》·卷九上 子部·天文算法類·算書之屬·《數學九章》·卷九上 子部·天文算法類·算書之屬·《數學九章》·卷九上\end{unitzh}
\begin{uniten}Masters section (zibu), astronomy and mathematics category (tianwen suanfa lei), mathematical books subcategory (suanshu zhi shu), Shuxue jiuzhang, juan 9, part 1
Masters section, astronomy and mathematics category, mathematical books subcategory, Shuxue jiuzhang, juan 9, part 1
Masters section, astronomy and mathematics category, mathematical books subcategory, Shuxue jiuzhang, juan 9, part 1
Masters section, astronomy and mathematics category, mathematical books subcategory, Shuxue jiuzhang, juan 9, part 1
Masters section, astronomy and mathematics category, mathematical books subcategory, Shuxue jiuzhang, juan 9, part 1
Masters section, astronomy and mathematics category, mathematical books subcategory, Shuxue jiuzhang, juan 9, part 1
Masters section, astronomy and mathematics category, mathematical books subcategory, Shuxue jiuzhang, juan 9, part 1
Masters section, astronomy and mathematics category, mathematical books subcategory, Shuxue jiuzhang, juan 9, part 1
Masters section, astronomy and mathematics category, mathematical books subcategory, Shuxue jiuzhang, juan 9, part 1
Masters section, astronomy and mathematics category, mathematical books subcategory, Shuxue jiuzhang, juan 9, part 1
Masters section, astronomy and mathematics category, mathematical books subcategory, Shuxue jiuzhang, juan 9, part 1\end{uniten}

\begin{unitzh}\unitlabel{shushu-061.12}《數學九章》卷九上\end{unitzh}
\begin{uniten}Shuxue jiuzhang, juan 9, part 1\end{uniten}

\section*{卷9下}\addcontentsline{toc}{section}{卷9下}

\begin{unitzh}\unitlabel{shushu-062.1}欽定四庫全書。《數學九章》卷九下，宋·秦九韶撰。市易。\end{unitzh}
\begin{uniten}Imperially Commissioned Complete Library of the Four Treasuries (Qinding Siku quanshu). Shuxue jiuzhang (Mathematics in Nine Sections), juan 9, part 2; composed by Qin Jiushao of the Song. Market Exchange (shiyi).\end{uniten}

\begin{unitzh}\unitlabel{shushu-062.2}推求典本\end{unitzh}
\begin{uniten}Finding the principal of a pawn\end{uniten}

\begin{unitzh}\unitlabel{shushu-062.3}問：典庫今年二月二十九日，有人取觧一號。主家聽得當事共筭本息一百六十貫八百三十二文，稱係前嵗頭臘月半觧去，月息二分二釐。欲知原本幾何？答曰：本一百二十貫文。\end{unitzh}
\begin{uniten}Question: At a pawnshop (dianku), on the 29th day of the 2nd month of this year, someone redeemed one numbered pledge. The proprietor heard that the person in charge had reckoned principal and interest together at 160 guan (貫) 832 wen (文), stating that it had been pawned in the middle of the first la month (toula yue) of the year before last, at a monthly interest of 2 fen (分) 2 li (釐). It is desired to know how much the original principal was. Answer: The principal is 120 guan wen.\end{uniten}

\begin{unitzh}\unitlabel{shushu-062.4}術曰：以粟米求之。置積日，乘息分數，增三百，為法；以三百乘共錢，為寔。寔如法而一，得本。\end{unitzh}
\begin{uniten}Method: Seek it by Millet and Grain (sumi). Set up the accumulated days (jiri), multiply by the interest in fen (xi fenshu), add 300, and make this the divisor (fa); multiply the combined money by 300 and make this the dividend (shi). Divide the dividend by the divisor, obtaining the principal.\end{uniten}

\begin{unitzh}\unitlabel{shushu-062.5}草曰：置前年頭臘月半，係四十五日，并去年三百六十，又加今年五十九日，共得四百六十四日，為積日。乘息二分二釐，得一百二十文八釐；增三百文，得四百二文八釐，為法。以三百文乘共錢一百六十貫八百三十二文，得四萬八千二百四十九貫六百文，為寔。如法而一，得一百二十貫文，為原本。\end{unitzh}
\begin{uniten}Working: Set up the middle of the first la month of the year before last, which comes to 45 days; combine it with last year's 360, and further add this year's 59 days: altogether 464 days are obtained, which are the accumulated days. Multiply by the interest 2 fen 2 li, obtaining 120 wen 8 li [Translator's note: 464 × 2 fen 2 li gives 102 wen 8 li, the value the next step requires.]; add 300 wen, obtaining 402 wen 8 li, which is the divisor (fa). Multiply the combined money 160 guan 832 wen by 300 wen, obtaining 48,249 guan 600 wen, which is the dividend (shi). Divide by the divisor, obtaining 120 guan wen, which is the original principal.\end{uniten}

\begin{unitzh}\unitlabel{shushu-062.6}按：題意謂前年十一月半典錢，去年未取，今年二月二十九日取去，按日計利，共合本息錢一十六千八百三十二文，月息每千二十二文，問本錢若干。法當以積日乗息二十二文，以三十日除之，得一千文之共息；加一千文，為一千文之本息共数。然後置今本息共數，以一千乗之，一千之本息共數除之，即得今本錢數。草中以十文為本，以二分二厘為息，不以三十除之為共息，而以三十乗十文，得三百為本，其理一也。\end{unitzh}
\begin{uniten}Note: The meaning of the problem is that money was raised on pawn in the middle of the 11th month of the year before last, was not redeemed last year, and was redeemed on the 29th day of the 2nd month of this year; interest is reckoned by the day, and principal and interest together come to 16,832 wen [Translator's note: the problem gives 160 guan 832 wen.]; the monthly interest is 22 wen per thousand; what is the principal? The rule should be: multiply the accumulated days by the interest of 22 wen and divide by 30 days, obtaining the total interest on 1000 wen; add 1000 wen, making the combined sum of principal and interest for 1000 wen. Then set up the present combined sum of principal and interest, multiply it by 1000, and divide by the combined sum of principal and interest for 1000; this gives the present principal. In the working, 10 wen is taken as the principal and 2 fen 2 li as the interest; it does not divide by 30 to make the total interest, but multiplies 10 wen by 30, obtaining 300 as the principal. The principle is the same.\end{uniten}

\begin{unitzh}\unitlabel{shushu-062.7}子部·天文算法類·算書之屬·《數學九章》·卷九下。子部·天文算法類·算書之屬·《數學九章》·卷九下。\end{unitzh}
\begin{uniten}$\langle$Masters section, Astronomy and Mathematics category, Mathematical Books subcategory, Shuxue jiuzhang, juan 9, part 2$\rangle$ $\langle$Masters section, Astronomy and Mathematics category, Mathematical Books subcategory, Shuxue jiuzhang, juan 9, part 2$\rangle$\end{uniten}

\begin{unitzh}\unitlabel{shushu-062.8}推知糴數\end{unitzh}
\begin{uniten}Inferring the amount of grain purchased\end{uniten}

\begin{unitzh}\unitlabel{shushu-062.9}問：和糴三百萬貫，求石數。聞每石牙錢三十文，糴埸量米折支牙人所得；毎石出牽錢八百，牙人量米四石六斗八合，折與牽頭。欲知米數、石價、牙錢、米、牽錢各幾何？\end{unitzh}
\begin{uniten}Question: A harmonious purchase of grain (hedi) of 3,000,000 guan; find the number of dan (石). It is heard that the brokerage money (yaqian) is 30 wen per dan, and that the grain-purchasing yard measures out rice to pay, by conversion, what the broker (yaren) has earned; for each dan 800 of leading money (qianqian) is paid out, and the broker measures out 4 dan 6 dou (斗) 8 ge (合) of rice and gives it, by conversion, to the leader (qiantou). It is desired to know how much each of these is: the amount of rice, the price per dan, the brokerage money, the rice, and the leading money.\end{uniten}

\begin{unitzh}\unitlabel{shushu-062.10}按：題意買米共用錢三十億，每石牙錢三十文，共牙錢折米給之；共折米每石牽錢八百文，共牽錢牙人又折米四石六斗零八合給之。求各數。\end{unitzh}
\begin{uniten}Note: The meaning of the problem is: in buying the rice, 3,000,000,000 of money is used in all; the brokerage money is 30 wen per dan, and the total brokerage money is given in converted rice; on the total converted rice the leading money is 800 wen per dan, and for the total leading money the broker in turn gives 4 dan 6 dou and 8 ge of converted rice. Find each quantity.\end{uniten}

\begin{unitzh}\unitlabel{shushu-062.11}答曰：糴到米一十二萬石，石價二十五貫文，牙錢三千六百貫文，折米一百四十四石，牽錢一百一十五貫二百文。\end{unitzh}
\begin{uniten}Answer: Rice purchased, 120,000 dan; price per dan, 25 guan wen; brokerage money, 3,600 guan wen; converted rice, 144 dan; leading money, 115 guan 200 wen.\end{uniten}

\begin{unitzh}\unitlabel{shushu-062.12}術曰：以商功求之，率變入之。置糴本、牙錢、牽錢相乘，為寔；以牽米為隅，開連枝立方，得石價。以價除本，得糴到米。以牙錢乘米，得縂牙錢；以價除之，得牙米。以牽錢乘牙米，得共牽錢。\end{unitzh}
\begin{uniten}Method: Seek it by Construction Works (shanggong), and enter it with the rates (lü) transformed. Set up the purchasing capital, the brokerage money and the leading money, multiply them together, and make this the dividend (shi); take the leading rice as the corner (yu), and extract the connected-branches cube root (kai lianzhi lifang), obtaining the price per dan. Divide the capital by the price, obtaining the rice purchased. Multiply the rice by the brokerage money, obtaining the total brokerage money; divide it by the price, obtaining the brokerage rice. Multiply the brokerage rice by the leading money, obtaining the total leading money.\end{uniten}

\begin{unitzh}\unitlabel{shushu-062.13}草曰：置糴本三百萬貫，乘牙錢三十文，得九千萬貫文；乘牽錢八百文，得七百二十億萬貫{\zhsmall 按：七百二十億貫，萬字誤。後餘寔内萬字同此。}，為價寔。置牽米四石六斗八合於寔數零文之下，為立方從隅，起歩。歩法常超二位，每超一度，商進之。今隅凡超四度，當於寔上約定首商二十貫。乃以商生隅超四石六斗八合，得九十二貫一百六十文，乃以為亷。又以商生亷，得一百八十四萬三千二百貫，為方。乃以方命上商二十貫除寔，訖，寔餘三百五十一億三千六百萬貫。復以商生隅四石六斗八合入亷，得一百八十四貫三百二十文；又以商生亷加入方内，得五百五十二萬九千六百貫，為方法；復以商又生隅四石六斗八合加入亷，得二百七十六貫四百八十文，為亷法。其方法一退，亷法二退，從隅三退。乃於首商之次約寔，續商五貫。以續商生隅四石六斗八合入亷，得二百九十九貫五百二十文；又以續商生亷入方，得七百二萬七千二百貫。乃命續商五貫除寔，適盡。\end{unitzh}
\begin{uniten}Working: Set up the purchasing capital 3,000,000 guan, multiply by the brokerage money 30 wen, obtaining 90,000,000 guan wen; multiply by the leading money 800 wen, obtaining 720,000,000,000,000 guan [Siku editors: 72,000,000,000 guan; the character wan (萬) is an error. The character wan in the remaining dividend later is likewise.], which is the dividend (shi) for the price. Place the leading rice, 4 dan 6 dou 8 ge, beneath the odd wen (lingwen) of the dividend, as the accompanying corner (congyu) of the cube, and start stepping (qibu). The stepping rule always leaps over two places; for each leap, the quotient advances. Now the corner leaps four times in all, so above the dividend one should estimate and fix the first quotient as 20 guan. Then with the quotient generate (sheng) the leaped corner, 4 dan 6 dou 8 ge, obtaining 92 guan 160 wen, and take this as the edge (lian). Again with the quotient generate the edge, obtaining 1,843,200 guan, as the square (fang). Then with the square call (ming) the upper quotient 20 guan and divide the dividend; when this is done, the dividend has a remainder of 35,136,000,000 guan. Again with the quotient generate the corner, 4 dan 6 dou 8 ge, and enter it into the edge, obtaining 184 guan 320 wen; again with the quotient generate the edge and add it into the square, obtaining 5,529,600 guan, as the square divisor (fangfa); again with the quotient once more generate the corner, 4 dan 6 dou 8 ge, and add it into the edge, obtaining 276 guan 480 wen, as the edge divisor (lianfa). The square divisor retreats (tui) once, the edge divisor retreats twice, and the accompanying corner retreats three times. Then, next to the first quotient, estimate from the dividend and continue the quotient with 5 guan. With the continued quotient generate the corner, 4 dan 6 dou 8 ge, and enter it into the edge, obtaining 299 guan 520 wen; again with the continued quotient generate the edge and enter it into the square, obtaining 7,027,200 guan. Then call the continued quotient 5 guan and divide the dividend: it is exactly exhausted.\end{uniten}

\begin{unitzh}\unitlabel{shushu-062.14}所得二十五貫，為毎石米價，以為法；以糴本三百萬貫為寔。寔如法而一，得一十二萬石，為糴到米數。以米數乘牙錢三十，得三千六百貫，為牙錢。以石價二十五貫除牙錢三千六百貫文，得一百四十四石，為糴埸量米折牙錢。以牽錢八百乘牙米一百四十四石，得一百一十五貫二百文，為牽頭得牙人所與牽錢之數。以石價二十五貫文約牽錢一百一十五貫二百文，得四石六斗八合，為牽米折錢。合問。\end{unitzh}
\begin{uniten}The 25 guan obtained is the price of rice per dan; take it as the divisor (fa), and take the purchasing capital 3,000,000 guan as the dividend (shi). Divide the dividend by the divisor, obtaining 120,000 dan, which is the amount of rice purchased. Multiply the amount of rice by the brokerage money 30, obtaining 3,600 guan, which is the brokerage money. Divide the brokerage money 3,600 guan wen by the price per dan 25 guan, obtaining 144 dan, which is the rice measured out by the grain-purchasing yard in conversion of the brokerage money. Multiply the brokerage rice 144 dan by the leading money 800, obtaining 115 guan 200 wen, which is the amount of leading money given by the broker that the leader obtains. Reduce (yue) the leading money 115 guan 200 wen by the price per dan 25 guan wen, obtaining 4 dan 6 dou 8 ge, which is the leading rice in conversion of the money. This agrees with the question.\end{uniten}

\begin{unitzh}\unitlabel{shushu-062.15}按：此術之意，以立天元一觧之。法：立天元一為米每石之價，以折牽米四石六斗零八合乘之，得四元{\zhsmall 六○八}，為共牽錢。應以每石牽錢八百除之，寄分母代除，得八百分元之四分{\zhsmall 六○八}，為共折給牙人米數。又以天元一乘之，得八百分平方之四分{\zhsmall 六○八}，為共牙錢。應以每石牙錢三十除之，寄分母代除，得三十分又八百分平方之四分{\zhsmall 六○八}，為共米數。又以天元一乗之，得三十分又八百分立方之四分{\zhsmall 六○八}，為共米價，與三百萬貫相等。兩邉各以三十乗之，又以八百乗之，得四立方{\zhsmall 六○八}，與七百二十億貫等。以立方數除兩邉，得一立方，與一百五十六億二千五百萬貫等。置貫數為寔，開立方，得二十五貫，為米每石價數。草中不除，開連枝立方，為除之不能盡者，便於用也。\end{unitzh}
\begin{uniten}Note: The intent of this method is to solve it by setting up the celestial unknown. Rule: set up the celestial unknown as the price of rice per dan; multiply it by the converted leading rice, 4 dan 6 dou and 8 ge, obtaining 4 yuan [608] [Translator's note: the small-print digits are the decimal places, i.e. 4.608.], which is the total leading money. It ought to be divided by the leading money per dan, 800; lodging the denominator (fenmu) in place of dividing, one obtains 4 parts [608] out of 800 parts of yuan, which is the total amount of rice converted and given to the broker. Again multiply it by the celestial unknown, obtaining 4 parts [608] out of 800 parts of the square (pingfang), which is the total brokerage money. It ought to be divided by the brokerage money per dan, 30; lodging the denominator in place of dividing, one obtains 4 parts [608] out of 30 parts and further 800 parts of the square, which is the total amount of rice. Again multiply it by the celestial unknown, obtaining 4 parts [608] out of 30 parts and further 800 parts of the cube (lifang), which is the total price of the rice and is equal to 3,000,000 guan. Multiply each side by 30, and further multiply by 800, obtaining 4 cubes [608], equal to 72,000,000,000 guan. Divide both sides by the number of cubes, obtaining 1 cube, equal to 15,625,000,000 guan. Set up the number of guan as the dividend (shi) and extract the cube root, obtaining 25 guan, which is the price of rice per dan. In the working one does not divide but extracts the connected-branches cube root: for cases where the division cannot be carried out to the end, this is convenient to use.\end{uniten}

\begin{unitzh}\unitlabel{shushu-062.16}子部·天文算法類·算書之屬·《數學九章》·卷九下。子部·天文算法類·算書之屬·《數學九章》·卷九下。子部·天文算法類·算書之屬·《數學九章》·卷九下。子部·天文算法類·算書之屬·《數學九章》·卷九下。\end{unitzh}
\begin{uniten}$\langle$Masters section, Astronomy and Mathematics category, Mathematical Books subcategory, Shuxue jiuzhang, juan 9, part 2$\rangle$ $\langle$Masters section, Astronomy and Mathematics category, Mathematical Books subcategory, Shuxue jiuzhang, juan 9, part 2$\rangle$ $\langle$Masters section, Astronomy and Mathematics category, Mathematical Books subcategory, Shuxue jiuzhang, juan 9, part 2$\rangle$ $\langle$Masters section, Astronomy and Mathematics category, Mathematical Books subcategory, Shuxue jiuzhang, juan 9, part 2$\rangle$\end{uniten}

\begin{unitzh}\unitlabel{shushu-062.17}累收庫本\end{unitzh}
\begin{uniten}Accumulated receipts of treasury capital\end{uniten}

\begin{unitzh}\unitlabel{shushu-063.1}問：有庫本錢五十萬貫，月息六厘半。今令掌事每月帶本約息共還一十萬，欲知幾何月而納足，并末復畸錢多少？答曰：本息納足共七箇月，末後一月畸錢二萬四千七百六貫二百七十九文三分四厘八毫四絲六忽七微七沙三莽一輕二清五烟。\end{unitzh}
\begin{uniten}Question: There is treasury capital money of 500000 guan (貫), with a monthly interest of 6 1/2 li (厘). Now the person in charge (zhangshi) is made to repay each month, bringing the principal together with the interest, 100000 in all. One wishes to know in how many months it is paid in full, and also how much odd money (jiqian) is paid back at the end. Answer: Principal and interest are paid in full in 7 months in all; the odd money of the last month is 24706 guan 279 wen (文) 3 fen (分) 4 li 8 hao (毫) 4 si (絲) 6 hu (忽) 7 wei (微) 7 sha (沙) 3 mang (莽) 1 qing (輕) 2 qing (清) 5 yan (烟).\end{uniten}

\begin{unitzh}\unitlabel{shushu-063.2}術曰：以盈朏變法求之。置原本，以息數退位乗，歸本位，每出共納，累得月數，以末後不及數為足月錢數。\end{unitzh}
\begin{uniten}Method: Solve it by the transformed method (bianfa) of yingfei (盈朏) [Translator's note: probably for 盈朒, "excess and deficit"]. Set up the original principal; multiply it by the interest figure with the place shifted back (tuiwei), and return it to the principal's place; each time take out the joint payment; accumulating, obtain the number of months; take the final amount that falls short as the money of the month of full payment.\end{uniten}

\begin{unitzh}\unitlabel{shushu-063.3}草曰：置本五十萬貫，以六釐五毫乘之，入共本内，得五十三萬二千五百貫文。内减初月一十萬，餘四十三萬二千五百貫文。又以六釐五毫乘之，餘本内得四十六萬六百一十二貫五百文。又减次月一十萬貫，餘三十六萬六百一十二貫五百文。又以六釐五毫乘之，入餘本内，得三十八萬四千五十二貫三百一十二文五分。又减第三月錢一十萬貫文，餘二十八萬四千五十二貫三百一十二文五分。又以六釐五毫乘之，得三十萬二千五百一十五貫七百一十二文八分一釐二毫五絲。内减第四月錢一十萬貫，餘二十萬二千五百一十五貫七百一十二文二分一釐二毫五絲。\end{unitzh}
\begin{uniten}Working: Set up the principal, 500000 guan; multiply it by 6 li (釐) 5 hao and add it into the principal, obtaining together 532500 guan wen. From this subtract the first month's 100000; remainder 432500 guan wen. Again multiply it by 6 li 5 hao; within the remaining principal obtain 460612 guan 500 wen. Again subtract the next month's 100000 guan; remainder 360612 guan 500 wen. Again multiply it by 6 li 5 hao and add it into the remaining principal, obtaining 384052 guan 312 wen 5 fen. Again subtract the third month's money, 100000 guan wen; remainder 284052 guan 312 wen 5 fen. Again multiply it by 6 li 5 hao, obtaining 302515 guan 712 wen 8 fen 1 li 2 hao 5 si. From this subtract the fourth month's money, 100000 guan; remainder 202515 guan 712 wen 2 fen 1 li 2 hao 5 si.\end{uniten}

\begin{unitzh}\unitlabel{shushu-063.4}又以六釐五毫乘之，入餘本内，得二十一萬五千六百七十九貫二百三十四文一分四釐五毫三絲一忽二微五塵。内减第五月錢一十萬貫，餘一十一萬五千六百七十九貫二百三十四文一分四釐五毫三絲一忽二微五塵。又以六釐五毫乘之，入餘本内，得一十二萬三千一百九十八貫三百八十四文三分六釐四毫七絲五忽七微八塵一沙二𣺌五莾。减第六月錢一十萬貫，餘二萬三千一百九十八貫三百八十四文三分六釐四毫七絲五忽七微八塵一沙二𣺌五莾。又以六釐五毫乘之，入餘本内，得二萬四千七百六貫二百七十九文三分四釐八毫四絲六忽七微無塵七沙無𣺌三莾一輕二清五烟，為第七月納足本息畸錢。\end{unitzh}
\begin{uniten}Again multiply it by 6 li 5 hao and add it into the remaining principal, obtaining 215679 guan 234 wen 1 fen 4 li 5 hao 3 si 1 hu 2 wei 5 chen (塵). From this subtract the fifth month's money, 100000 guan; remainder 115679 guan 234 wen 1 fen 4 li 5 hao 3 si 1 hu 2 wei 5 chen. Again multiply it by 6 li 5 hao and add it into the remaining principal, obtaining 123198 guan 384 wen 3 fen 6 li 4 hao 7 si 5 hu 7 wei 8 chen 1 sha 2 miao (𣺌) 5 mang. Subtract the sixth month's money, 100000 guan; remainder 23198 guan 384 wen 3 fen 6 li 4 hao 7 si 5 hu 7 wei 8 chen 1 sha 2 miao 5 mang. Again multiply it by 6 li 5 hao and add it into the remaining principal, obtaining 24706 guan 279 wen 3 fen 4 li 8 hao 4 si 6 hu 7 wei, no chen, 7 sha, no miao, 3 mang 1 qing (輕) 2 qing (清) 5 yan; this is the odd money with which principal and interest are paid in full in the seventh month.\end{uniten}

\begin{unitzh}\unitlabel{shushu-063.5}均貨推本\end{unitzh}
\begin{uniten}Distributing Goods by Inferring the Capitals (jun huo tui ben)\end{uniten}

\begin{unitzh}\unitlabel{shushu-063.6}問：海舶赴務抽畢，除納主家貨物外，有沉香五千八十八兩，胡椒一萬四百三十包{\zhsmall 包四十觔}，象牙二百一十二合{\zhsmall 以大小為合，觔兩俱等}，係甲、乙、丙、丁四人合本博到。縁昨來凑本，互有假借，甲分到官供稱：甲本金二百兩，鹽四袋鈔一十道；乙本銀八百兩，鹽三袋鈔八十八道；丙本銀一千六百七十兩，度蝶一十五道；丁本度蝶五十二道，金五十八兩八銖。已上共估值四十二萬四千貫。甲借乙鈔，乙借丙銀，丙借丁度蝶，丁借甲金。今合撥各借物歸原主名下為率，均分上件貨物。欲知金、銀、袋鹽、度牒原價，及四人各合得香、椒、牙幾何？\end{unitzh}
\begin{uniten}Question: A seagoing ship went to the office (wu), and the levy (chou) was completed. Apart from the goods paid in to the owner's house, there are aloeswood (chenxiang) 5088 liang (兩), pepper 10430 packs (bao 包) [Note: a pack is 40 jin (觔)], and ivory 212 he (合) [Note: a large and a small one are taken as a he; the jin and liang are all equal], which the four persons Jia, Yi, Bing and Ding obtained in trade with pooled capital. Because, when they lately put their capital together, they borrowed from one another, Jia, going separately to the authorities, stated in deposition: Jia's capital: gold 200 liang, salt 4 bags (dai 袋), vouchers (chao 鈔) 10 dao (道); Yi's capital: silver 800 liang, salt 3 bags, vouchers 88 dao; Bing's capital: silver 1670 liang, ordination certificates (dudie 度蝶) 15 dao; Ding's capital: ordination certificates 52 dao, gold 58 liang 8 zhu (銖). All the above are valued at 424000 guan. Jia borrowed Yi's vouchers, Yi borrowed Bing's silver, Bing borrowed Ding's ordination certificates, Ding borrowed Jia's gold. Now the borrowed items should be transferred back under the names of their original owners, to serve as the rates (lü), and the above goods divided fairly. One wishes to know the original prices of the gold, the silver, the bagged salt and the ordination certificates (度牒), and how much aloeswood, pepper and ivory each of the four persons should get.\end{uniten}

\begin{unitzh}\unitlabel{shushu-063.7}按：題意謂甲金、乙鹽、丙銀、丁牒原本不同，互借為同本，買得香、椒、牙三色。今有互借各物及同本貫數，求原本以分所買之物，葢方程而亷衰分之法也。甲、乙二條内鹽、鈔二色，寔即一色：先言鹽袋數，乃一鈔之數，以鈔數乘之，始為鹽數。是鈔數既贅，又不明言其故，皆故為隠晦也。\end{unitzh}
\begin{uniten}Note: The intent of the problem is that Jia's gold, Yi's salt, Bing's silver and Ding's certificates differed as original capitals, and by borrowing from one another they were made equal capitals, with which the three kinds, aloeswood, pepper and ivory, were bought. Now, given the items borrowed from one another and the number of guan of the equal capitals, one seeks the original capitals in order to divide the goods bought; this is indeed the method of rectangular arrays (fangcheng) and lian [Translator's note: 亷 is probably an error for 兼, "combined with"] distribution in proportion (cuifen). Within the two items of Jia and Yi, the two kinds, salt and vouchers, are in fact one kind: first the number of bags of salt is stated, which is the amount for one voucher; only when it is multiplied by the number of vouchers does it give the amount of salt. Thus the number of vouchers is superfluous, and moreover the reason for it is not stated clearly; all this is deliberately made obscure.\end{uniten}

\begin{unitzh}\unitlabel{shushu-063.8}答曰：甲金每兩四百八十貫文，本一十二萬四十貫文，合得沉香一千四百八十八兩，胡椒三千五十包一十一觔五兩五十三分兩之七，象牙六十二合。乙鹽每袋二百五十貫文，本七萬六千貫文，合得沉香九百一十二兩，胡椒一千八百六十九包二十一觔二兩五十三分兩之六，象牙三十八合。丙銀每兩五十貫文，本一十二萬三千五百貫文，合得沉香一千四百八十二兩，胡椒三千三十七包三十九觔五兩五十三分兩之二十三，象牙六十一合四分合之三。丁度牒每道一千五百貫文，本一十萬五百貫文，合得沉香一千二百六兩，胡椒二千四百七十二包八觔三兩五十三分兩之十七，象牙五十合分合之一。\end{unitzh}
\begin{uniten}Answer: Jia's gold, per liang 480 guan wen; capital 120040 guan wen; should get aloeswood 1488 liang, pepper 3050 packs 11 jin 5 7/53 liang, ivory 62 he. Yi's salt, per bag 250 guan wen; capital 76000 guan wen; should get aloeswood 912 liang, pepper 1869 packs 21 jin 2 6/53 liang, ivory 38 he. Bing's silver, per liang 50 guan wen; capital 123500 guan wen; should get aloeswood 1482 liang, pepper 3037 packs 39 jin 5 23/53 liang, ivory 61 3/4 he. Ding's ordination certificates, per dao 1500 guan wen; capital 100500 guan wen; should get aloeswood 1206 liang, pepper 2472 packs 8 jin 3 17/53 liang, ivory 50 he and 1 part of a he [Translator's note: the number of parts is missing].\end{uniten}

\begin{unitzh}\unitlabel{shushu-063.9}術曰：以方程求之，衰分入之，正負入之。置共錢，以人數約之，得數列入人數，各為行積。次置諸色各物數為叚子，對本色。有分者通之，可約者約之，為定率。以第一行為右，以第二行為副，以第三行為次，第四行為左。每以下位互遍乘之，每騐其積，以少减多。如同名相减，異名相加，正無人負之，負無人正之；如同名相加，異名相减，正無人正之，負無人負之。得一叚為法，以餘積為寔，除之，各得諸價。以諸價列右行，以各物數列左行，以兩行對乘，各得本率。以諸色求等約之，得列衰。併諸衰為總法。以列衰遍乘各物諸數，各為寔。諸寔並如總法而一，各得其物。除不盡者，以觔兩通而除之，或又分母命之。\end{unitzh}
\begin{uniten}Method: Solve it by rectangular arrays (fangcheng), bringing in distribution in proportion (cuifen) and bringing in positive and negative. Set up the total money and reduce it by the number of persons; arrange the number obtained according to the number of persons, each as the total (ji) of a column (hang). Next set up the amounts of the items of the various kinds as the segments (duanzi), opposite their own kinds. Where there are fractions, convert them (tong); where they can be reduced, reduce them; these make the fixed rates (dinglü). Take the first column as the right, the second column as the secondary (fu), the third column as the next (ci), and the fourth column as the left. Each time multiply mutually throughout by the lower places; each time examine their totals and subtract the smaller from the larger. If same names subtract one another and different names add to one another, a positive with nothing opposite (wu ren) is made negative and a negative with nothing opposite is made positive [Translator's note: 無人, literally "no person", is evidently for 無入, "nothing entered against"; so also below]; if same names add to one another and different names subtract one another, a positive with nothing opposite is made positive and a negative with nothing opposite is made negative. When one segment is obtained, take it as the divisor (fa) and the remaining total as the dividend (shi); divide, and obtain each of the prices. Arrange the prices as the right column and the amounts of each item as the left column; multiply the two columns against each other, each obtaining its own rate (benlü). Find the common measure (deng) of the various kinds and reduce by it, obtaining the arranged proportions (liecui). Add the proportions together as the total divisor (zongfa). Multiply the amounts of each item throughout by the arranged proportions, each as a dividend. Divide all the dividends by the total divisor, each obtaining its goods. Where the division does not come out exactly, convert by jin and liang and divide; or else name it by the denominator.\end{uniten}

\begin{unitzh}\unitlabel{shushu-063.10}草曰：置估值四十二萬四千貫，以四人約之，得一十萬六千貫，為各積，以人數列四位。次置甲金二百兩於右上；以四袋乘鈔一十道，得四十袋{\zhsmall 按：此條以四袋為一鈔，相乘為塩數，是二色止為一色，不如即用塩為正}，於右副，為右行。次置乙鈔八十八道，以三袋乘之，得鹽二百六十四袋{\zhsmall 按：此條以三袋為一鈔}，及銀八百兩，為副行。次置丙銀一千六百七十兩、度牒一十五道，為次行。次置丁度牒五十二道、金五十八兩八銖，為左行。驗得首圖左行上叚金帶八銖，是三分兩之一，乃以分母遍乘左行諸數，只以分子一内入左上金内，其左積得三十一萬八千貫，左金得一百七十五兩，左度牒得一百五十六道，為次圖。驗次圖四行皆可求等：右行求得四十，約之；副行求得八，約之；次行求得五，約之；左行求得一，約之。各得數為定率圖。\end{unitzh}
\begin{uniten}Working: Set up the valuation, 424000 guan; reduce it by the 4 persons, obtaining 106000 guan as the total of each; arrange it in 4 places according to the number of persons. Next place Jia's gold, 200 liang, at the top of the right; with 4 bags multiply the vouchers, 10 dao, obtaining 40 bags [Siku editors: This item takes 4 bags as one voucher; multiplied together they give the amount of salt. Thus the two kinds are only one kind; it would be better simply to use salt as the proper one], at the right's secondary (fu); this makes the right column. Next place Yi's vouchers, 88 dao; with 3 bags multiply them, obtaining salt 264 bags [Siku editors: This item takes 3 bags as one voucher], and silver 800 liang; this makes the secondary column. Next place Bing's silver, 1670 liang, and ordination certificates, 15 dao; this makes the next column. Next place Ding's ordination certificates, 52 dao, and gold, 58 liang 8 zhu; this makes the left column. Examining, one finds that in the first diagram the gold in the upper segment of the left column carries 8 zhu, which is 1/3 liang; so multiply all the numbers of the left column throughout by the denominator, only adding the numerator 1 into the gold at the top of the left: its left total becomes 318000 guan, the left gold 175 liang, the left ordination certificates 156 dao; this makes the next diagram. Examining the next diagram, a common measure can be found in all four columns: for the right column one finds 40 and reduces by it; for the secondary column one finds 8 and reduces by it; for the next column one finds 5 and reduces by it; for the left column one finds 1 and reduces by it. The numbers obtained in each make the fixed-rate diagram (dinglü tu).\end{uniten}

\begin{unitzh}\unitlabel{shushu-063.11}定率圖：右積得二千六百五十貫，金五兩，鹽一袋；副積得一萬三千二百五十貫，鹽三十三袋，銀一百兩；次積得二萬一千二百貫，銀三百三十四兩，度牒三道；左積得三十一萬八千貫，金一百七十五兩，度牒一百五十六道。乃以定圖次行度牒三因左行，左積得九十五萬四千貫，金五百二十五兩，度牒四百六十八道。次以定圖左下度牒一百五十六乗次行，積得三百三十萬七千二百貫，銀五萬二千一百四兩{\zhsmall 按：此下落「度牒四百六十八道」八字}。\end{unitzh}
\begin{uniten}Fixed-rate diagram: the right total comes to 2650 guan, gold 5 liang, salt 1 bag; the secondary total comes to 13250 guan, salt 33 bags, silver 100 liang; the next total comes to 21200 guan, silver 334 liang, ordination certificates 3 dao; the left total comes to 318000 guan, gold 175 liang, ordination certificates 156 dao. Then, with the ordination certificates, 3, of the next column of the fixed diagram, multiply (yin) the left column: the left total becomes 954000 guan, gold 525 liang, ordination certificates 468 dao. Next, with the ordination certificates, 156, at the bottom left of the fixed diagram, multiply the next column: the total becomes 3307200 guan, silver 52104 liang [Siku editors: below this the eight characters 度牒四百六十八道, "ordination certificates 468 dao", have dropped out].\end{uniten}

\begin{unitzh}\unitlabel{shushu-063.12}定率圖 維圖\end{unitzh}
\begin{uniten}Fixed-rate diagram. Wei diagram (維圖).\end{uniten}

\begin{unitzh}\unitlabel{shushu-063.13}乃驗維圖左及次行之下度牒等，當相減之。以積為端，當以左之少積減次之多積，按術曰「同名相減」。其次行之金空，而左行之金五百二十五兩有，為正；次空，為無，按術曰「正無人負之」，即以左行之金正加入次行金位，為負名，成音圖。仍置定圖左行諸數。乃驗音圖次行積得二百三十五萬三千二百貫正，金五百二十五兩負，銀五萬二千一百四兩為正，餘三行皆正。\end{unitzh}
\begin{uniten}Then examine the Wei diagram: the ordination certificates at the bottom of the left and the next columns are equal; they should be subtracted from one another. Taking the totals as the starting point (duan), the smaller total of the left should be subtracted from the larger total of the next; the method says "same names subtract one another". The gold of the next column is empty, while the gold of the left column, 525 liang, is present, being positive; the next, being empty, is "nothing" (wu); the method says "a positive with nothing opposite is made negative", so add the positive gold of the left column into the gold place of the next column, as a negative name; this makes the Yin diagram (音圖). Still set up the numbers of the left column of the fixed diagram. Then, examining the Yin diagram, the total of the next column comes to 2353200 guan, positive; gold 525 liang, negative; silver 52104 liang, positive; the remaining three columns are all positive.\end{uniten}

\begin{unitzh}\unitlabel{shushu-063.14}音圖 駁圖\end{unitzh}
\begin{uniten}Yin diagram. Bo diagram (駁圖).\end{uniten}

\begin{unitzh}\unitlabel{shushu-063.15}今驗音次行之負金，當以右行之正金補之，而其數不等。先以右金五約次金五百二十五，得一百五，以乗音圖右行。畢，其右積得二十七萬八千二百五十貫，金二百五十五兩正，鹽一百五袋正；其副、次、左三行如音圖故，乃成爻圖。今視爻圖右行之金正與次行之金負適等，即用右行直加次行，按術以「同名相加」；乃以右之金正減其次之金負，為空，按術以「異名相減」之。其次鹽空，為無人，按術以「正無人正之」。乃以爻圖右積二十七萬八千二百五十貫加次積二百三十五萬三千二百貫内，得二百六十三萬一千四百五十貫；其次金空，次鹽一百五袋正，次銀五萬二千一百四兩正。仍置定圖右行數，而成政圖。\end{unitzh}
\begin{uniten}Now, examining the negative gold in the next column of Yin, it should be made up by the positive gold of the right column, but their numbers are not equal. First, with the right gold, 5, reduce the next gold, 525, obtaining 105, and with this multiply the right column of the Yin diagram. When this is done, its right total comes to 278250 guan, gold 255 liang positive, salt 105 bags positive; its secondary, next and left, the three columns, are as in the Yin diagram, as before; this then makes the Yao diagram (爻圖). Now one sees that in the Yao diagram the positive gold of the right column and the negative gold of the next column are exactly equal, so directly add the right column to the next column, following the method's "same names add to one another"; then subtract the right's positive gold from the next's negative gold, which becomes empty, following the method's "different names subtract one another". The next's salt is empty, which is "nothing opposite" (wu ren), following the method's "a positive with nothing opposite is made positive". Then add the right total of the Yao diagram, 278250 guan, into the next total, 2353200 guan, obtaining 2631450 guan; its next gold is empty, the next salt 105 bags positive, the next silver 52104 liang positive. Still set up the numbers of the right column of the fixed diagram, and so form the Zheng diagram (政圖).\end{uniten}

\begin{unitzh}\unitlabel{shushu-063.16}政圖 卜圖\end{unitzh}
\begin{uniten}Zheng diagram. Bu diagram (卜圖).\end{uniten}

\begin{unitzh}\unitlabel{shushu-064.1}今視政圖從者，乃擇其諸行本色可求等：首金可，鹽亦可；盖金多鹽少，乃以政圖副、次兩行鹽數三十三與一百五求等，得三。故以三約三十三，得一十一，以乗次行；又以三約一百五，得三十五，以乗副行。畢，其副積得四十六萬三千七百五十貫，鹽一千一百五十五袋，銀三千五百兩；次積二千八百九十四萬五千九百五十貫，鹽一千一百五十五袋，銀五十七萬三千一百四十四兩，皆正。列成卜圖。\end{unitzh}
\begin{uniten}Now look at the zheng (政) diagram. As for what follows, select among its rows items of the same kind (bense) for which a common measure (deng) can be sought. Gold would do first, and salt would also do. Since the gold is much and the salt little, take the salt numbers of the two rows of the zheng diagram, the secondary row (fu) and the next row (ci), 33 and 105, and seek their common measure, obtaining 3. Therefore reduce 33 by 3, obtaining 11, and multiply the next row by it; again reduce 105 by 3, obtaining 35, and multiply the secondary row by it. When this is done, the secondary row's total (ji) is 463,750 guan (貫), salt 1,155 bags, silver 3,500 liang (兩); the next row's total is 28,945,950 guan, salt 1,155 bags, silver 573,144 liang; all are positive. Set out, these form the bu (卜) diagram.\end{uniten}

\begin{unitzh}\unitlabel{shushu-064.2}宫圖\end{unitzh}
\begin{uniten}The gong (宫) diagram\end{uniten}

\begin{unitzh}\unitlabel{shushu-064.3}及視上圖，副行積少，次行積多，即以副行求減次行，皆是同名相減之。既畢，仍置定圖副行數，其次行乃得積二千八百四十八萬二千二百貫，銀得五十六萬九千六百四十四兩，列為宫圖。驗宫圖次行下，只有銀五十六萬九千六百四十四兩獨一數，以為法；以次積二千八百四十八萬二千二百貫為寔。寔如法而一，得五十貫，為銀一兩價，而成千圖。\end{unitzh}
\begin{uniten}Next, look at the diagram above. The secondary row's total is small and the next row's total is large, so use the secondary row to subtract from the next row; in every case this is subtraction of like names (tongming xiang jian). When this is done, still set out the numbers of the secondary row of the settled diagram (ding tu 定圖). The next row then has a total of 28,482,200 guan, and its silver is 569,644 liang; set out, this is the gong diagram. Examining the next row of the gong diagram, below there is only the silver, 569,644 liang, a single number: take it as the divisor (fa). Take the next row's total, 28,482,200 guan, as the dividend (shi). Dividing the dividend by the divisor gives 50 guan, which is the price of 1 liang of silver; this forms the qian (千) diagram.\end{uniten}

\begin{unitzh}\unitlabel{shushu-064.4}千圖\end{unitzh}
\begin{uniten}The qian diagram\end{uniten}

\begin{unitzh}\unitlabel{shushu-064.5}乃以千圖副行銀一百兩乗兩價五十貫，得五千貫，以減千圖副行之積一萬三千二百五十貫。訖，副積餘八千二百五十貫，其下鹽得三十三袋，銀空，而成曜圖。\end{unitzh}
\begin{uniten}Then with the silver of the secondary row of the qian diagram, 100 liang, multiply the price per liang, 50 guan, obtaining 5,000 guan, and with it subtract from the total of the secondary row of the qian diagram, 13,250 guan. When this is done, the secondary total has 8,250 guan remaining; below it the salt is 33 bags and the silver is empty; this forms the yao (曜) diagram.\end{uniten}

\begin{unitzh}\unitlabel{shushu-064.6}曜圖 支圖\end{unitzh}
\begin{uniten}The yao diagram; the zhi (支) diagram\end{uniten}

\begin{unitzh}\unitlabel{shushu-064.7}乃以曜圖副行之積八千二百五十貫為鹽寔，以其下鹽三十三袋為法，除之，得二百五十貫，為鹽一袋價，而成支圖。乃以支圖右行鹽一袋徧乗副行，畢，其副積只得二百五十貫。次以副行直減右行，畢，右積餘二千四百貫，金五兩，鹽空，而成閏圖。乃以閏圖右積二千四百貫為寔，金五兩為法，除之，得四百八十貫，為金一兩價，成定圖。次以左金一百七十五兩徧乗右行，直減左行，訖，左積得二十三萬四千貫，度牒一百五十六道，左金空，而成定圖。\end{unitzh}
\begin{uniten}Then take the total of the secondary row of the yao diagram, 8,250 guan, as the salt dividend, and the salt below it, 33 bags, as the divisor. Dividing gives 250 guan, which is the price of 1 bag of salt; this forms the zhi diagram. Then with the salt of the right row of the zhi diagram, 1 bag, multiply the secondary row throughout (bian cheng); when this is done, the secondary total is only 250 guan. Next use the secondary row to subtract directly (zhi jian) from the right row. When this is done, the right total has 2,400 guan remaining, gold 5 liang, the salt empty; this forms the run (閏) diagram. Then take the right total of the run diagram, 2,400 guan, as the dividend, and the gold, 5 liang, as the divisor. Dividing gives 480 guan, which is the price of 1 liang of gold; this forms the settled diagram. Next with the left gold, 175 liang, multiply the right row throughout and subtract it directly from the left row. When this is done, the left total is 234,000 guan, ordination certificates (dudie 度牒) 156, the left gold empty; this forms the settled diagram.\end{uniten}

\begin{unitzh}\unitlabel{shushu-064.8}閏圖 定圖\end{unitzh}
\begin{uniten}The run diagram; the settled diagram\end{uniten}

\begin{unitzh}\unitlabel{shushu-064.9}今驗定圖左積二十三萬四十貫為寔，以左下度牒一百五十六道為法，除之，得一千五百貫，為度牒一道價，而成終圖。既得金銀每兩、鈔鹽每袋、度牒每各色之價，次列甲乙丙丁四人乗之。\end{unitzh}
\begin{uniten}Now examine the settled diagram. Take the left total, 230,040 guan [Translator's note: so in the text; above, 234,000 guan], as the dividend, and the ordination certificates below on the left, 156, as the divisor. Dividing gives 1,500 guan, which is the price of 1 ordination certificate; this forms the final diagram (zhong tu 終圖). Having obtained the prices of each kind — gold and silver per liang, voucher salt (chaoyan 鈔鹽) per bag, ordination certificates each — next set out the four persons Jia, Yi, Bing and Ding and multiply them.\end{uniten}

\begin{unitzh}\unitlabel{shushu-064.10}終圖\end{unitzh}
\begin{uniten}The final diagram\end{uniten}

\begin{unitzh}\unitlabel{shushu-064.11}後以首圖右金二百兩併左金五十八兩八銖，得二百五十八兩。以八銖為三分兩之一，通分内子，得七百七十五，於左甲。其右價四百八十貫，乃以左甲母三約之，為一百六十貫，於左甲。次以右鹽四千袋併副鹽二百六十四袋，得三百四袋，於左乙；以鹽價二百五十貫於右乙。次以副銀八百兩併銀一千六百七十兩，得二千四百七十兩，於左丙；以銀價五十貫於右丙。又以次行度牒一十五道併左度牒五十二道，得六十七道，於左丁；以度牒價一千五百貫於右丁。兩行對乗之。\end{unitzh}
\begin{uniten}Afterwards take the right gold of the first diagram (shou tu 首圖), 200 liang, and combine it with the left gold, 58 liang 8 zhu (銖), obtaining 258 liang. Taking 8 zhu as 1/3 liang, reduce to the common denominator and add in the numerator (tongfen nei zi), obtaining 775, at left Jia. Its right price is 480 guan; reduce it by the denominator of left Jia, 3, making 160 guan, at left Jia [Translator's note: so in the text; below, this 160 is the right Jia]. Next take the right salt, 4,000 bags, and combine it with the secondary salt, 264 bags, obtaining 304 bags [Translator's note: so in the text; 4,000 + 264 is not 304], at left Yi; put the salt price, 250 guan, at right Yi. Next take the secondary silver, 800 liang, and combine it with the silver 1,670 liang, obtaining 2,470 liang, at left Bing; put the silver price, 50 guan, at right Bing. Again take the ordination certificates of the next row, 15, and combine them with the left ordination certificates, 52, obtaining 67, at left Ding; put the price of the ordination certificates, 1,500 guan, at right Ding. Multiply the two rows against each other (dui cheng).\end{uniten}

\begin{unitzh}\unitlabel{shushu-064.12}以右甲一百六十乗左甲七百七十五兩，得一十二萬四千貫，為甲原本；以右乙二百五十貫乗左乙三百四袋，得七萬六千貫，為乙原本；以右丙五十貫乗左丙二千四百七十兩，得一十二萬三千五百貫，為丙原本；以右丁一千五百貫乗左丁六十七道，得一十萬五百貫，為丁原本。列四人各得原本，求等，得五百貫，皆以五百貫為法除之：甲行二百四十八，乙得一百五十二，丙得二百四十七，丁得二百一，各為列衰於右行。併右行列衰，得八百四十八，為總法。次置愽到沉香五千八十八兩，遍乗列衰，各為沉香寔；次置胡椒一萬四百三十包，亦遍乗列衰，為椒寔；次置象牙四百二十四條，以大小各半之，得二百一十二合，亦遍乗列衰，為牙寔。\end{unitzh}
\begin{uniten}With the right Jia, 160, multiply the left Jia, 775 liang, obtaining 124,000 guan, which is Jia's original capital (yuanben). With the right Yi, 250 guan, multiply the left Yi, 304 bags, obtaining 76,000 guan, which is Yi's original capital. With the right Bing, 50 guan, multiply the left Bing, 2,470 liang, obtaining 123,500 guan, which is Bing's original capital. With the right Ding, 1,500 guan, multiply the left Ding, 67 certificates, obtaining 100,500 guan, which is Ding's original capital. Set out the original capitals that each of the four persons obtains and seek their common measure, obtaining 500 guan; divide them all by 500 guan as the divisor: the Jia row 248, Yi gets 152, Bing gets 247, Ding gets 201; each is a listed proportion (liecui), in the right row. Add the listed proportions of the right row, obtaining 848, as the total divisor (zongfa). Next set out the agarwood (chenxiang 沉香) acquired in trade, 5,088 liang, and multiply the listed proportions throughout, each being an agarwood dividend. Next set out the pepper, 10,430 packages (bao 包), and likewise multiply the listed proportions throughout, as the pepper dividends. Next set out the ivory, 424 tusks; taking large and small as halves each gives 212 pairs (he 合); likewise multiply the listed proportions throughout, as the ivory dividends.\end{uniten}

\begin{unitzh}\unitlabel{shushu-064.13}甲得一百二十六萬一千八百二十四，乙得七十七萬三千三百七十六，丙得一百二十五萬六千七百二十六，丁得一百二萬三千六百八十八，各為沉香寔。以總法八百四十八除之：甲得沉香一千四百八十八兩，乙得沉香九百一十二兩，丙得沉香一千四百八十二兩，丁得沉香一千二百六兩。\end{unitzh}
\begin{uniten}Jia gets 1,261,824, Yi gets 773,376, Bing gets 1,256,726, Ding gets 1,023,688 [Translator's note: the figures for Bing and Ding are so in the text and do not divide exactly by 848], each being an agarwood dividend. Divide them by the total divisor 848: Jia gets agarwood 1,488 liang, Yi gets agarwood 912 liang, Bing gets agarwood 1,482 liang, Ding gets agarwood 1,206 liang.\end{uniten}

\begin{unitzh}\unitlabel{shushu-064.14}甲得二百五十八萬六千六百四十，乙得一百五十八萬五千三百六十，丙得二百五十七萬六千二百一十，丁得二百九萬六千四百三十，各為椒寔。以總法八百四十八除之：甲得三千五十包，不盡二百四十包，以包率四十觔乘之，得九千六百觔；又以法除之，得一十一觔，不盡二百七十二觔，以十六兩通之，得四千三百五十二兩；又以法除之，得五兩，不盡一百一十二，求等得一十六，約之，得五十三分兩之七。約甲得椒三千五十包一十一觔五兩五十三分兩之七。乙得一千八百六十九包，不盡四百四十八包，以四十觔乘之，得一萬七千九百二十；又以法除之，得二十一觔，不盡一百一十二觔，以十六兩通之，得一千七百九十二兩；又以法除，得二兩，不盡九十六兩，求等得十六，約之，得五十三分兩之六，為乙合得椒一千八百六十九包二十一觔二兩五十三分兩之六。\end{unitzh}
\begin{uniten}Jia gets 2,586,640, Yi gets 1,585,360, Bing gets 2,576,210, Ding gets 2,096,430, each being a pepper dividend. Divide them by the total divisor 848. Jia gets 3,050 packages, with a remainder (bu jin) of 240 packages; multiply it by the package rate (lü), 40 jin (觔), obtaining 9,600 jin. Again divide by the divisor, obtaining 11 jin, with a remainder of 272 jin; convert it (tong) by 16 liang, obtaining 4,352 liang. Again divide by the divisor, obtaining 5 liang, with a remainder of 112; seeking the common measure gives 16; reducing gives 7/53 liang. Reduced, Jia gets pepper 3,050 packages 11 jin 5 7/53 liang. Yi gets 1,869 packages, with a remainder of 448 packages; multiply it by 40 jin, obtaining 17,920. Again divide by the divisor, obtaining 21 jin, with a remainder of 112 jin; convert it by 16 liang, obtaining 1,792 liang. Again divide by the divisor, obtaining 2 liang, with a remainder of 96 liang; seeking the common measure gives 16; reducing gives 6/53 liang. This is what Yi is due (he de): pepper 1,869 packages 21 jin 2 6/53 liang.\end{uniten}

\begin{unitzh}\unitlabel{shushu-064.15}丙得三千三十七包，不盡八百三十四，以四十觔通之，得三萬三千三百六十觔；又以法除之，得三十九觔，不盡二百八十八，以十六兩通之，得四千六百八兩；又以法除之，得五兩，不盡三百六十八兩，求等得十六，約之，得五十三分兩之二十三，為丙合得椒三千三十七包三十九觔五兩五十三分兩之二十三。丁得二千四百七十二包，不盡一百七十四，以四十觔通之，得六千九百六十觔；又以法除之，得八觔，不盡一百七十六，以十六兩通之，得二千八百一十六；又以法除之，得三兩，不盡二百七十二，求等得十六，約之，得五十三分兩之一十七，為丁合得椒二千四百七十二包八觔三兩五十三分兩之一十七。\end{unitzh}
\begin{uniten}Bing gets 3,037 packages, with a remainder of 834; convert it by 40 jin, obtaining 33,360 jin. Again divide by the divisor, obtaining 39 jin, with a remainder of 288; convert it by 16 liang, obtaining 4,608 liang. Again divide by the divisor, obtaining 5 liang, with a remainder of 368 liang; seeking the common measure gives 16; reducing gives 23/53 liang. This is what Bing is due: pepper 3,037 packages 39 jin 5 23/53 liang. Ding gets 2,472 packages, with a remainder of 174; convert it by 40 jin, obtaining 6,960 jin. Again divide by the divisor, obtaining 8 jin, with a remainder of 176; convert it by 16 liang, obtaining 2,816. Again divide by the divisor, obtaining 3 liang, with a remainder of 272; seeking the common measure gives 16; reducing gives 17/53 liang. This is what Ding is due: pepper 2,472 packages 8 jin 3 17/53 liang.\end{uniten}

\begin{unitzh}\unitlabel{shushu-064.16}甲得五萬二千五百七十六合，乙得三萬二千二百二十四合，丙得五萬二千三百六十四合，丁得四萬二千六百一十二合，各為牙寔。皆以總法八百四十八除之：甲合得牙六十二合，乙合得牙三十八合，丙合得牙六十一合，不盡六百三十六，求等得二百一十二，約之，得四分合之三；丁合得牙五十合，不盡二百十二，求等得二百一十二，約之，得四分合之一。\end{unitzh}
\begin{uniten}Jia gets 52,576 pairs, Yi gets 32,224 pairs, Bing gets 52,364 pairs, Ding gets 42,612 pairs, each being an ivory dividend. Divide them all by the total divisor 848. Jia is due 62 pairs of ivory, and Yi is due 38 pairs of ivory. Bing is due 61 pairs of ivory, with a remainder of 636; seeking the common measure gives 212; reducing gives 3/4 pair. Ding is due 50 pairs of ivory, with a remainder of 212; seeking the common measure gives 212; reducing gives 1/4 pair.\end{uniten}

\begin{unitzh}\unitlabel{shushu-064.17}按此條於方程正負之用、通分乗除之變，多所發明；歩筭雖繁，寔有條而不紊也。\end{unitzh}
\begin{uniten}Note: this item sheds much light on the use of positive and negative in rectangular arrays (fangcheng) and on the variations of reducing to a common denominator (tongfen), multiplication and division; although the step-by-step reckoning is complicated, it truly has order and is not confused.\end{uniten}

\begin{unitzh}\unitlabel{shushu-064.18}推求物價\end{unitzh}
\begin{uniten}Deducing the Prices of Goods\end{uniten}

\begin{unitzh}\unitlabel{shushu-065.1}問：推貨務三次支物，凖錢各一百四十七萬貫文。先撥沉香三千五百裹，瑇瑁二千二百觔，乳香三百七十五套；次撥沉香二千九百七十裹，瑇瑁二千一百三十觔，乳香三千五十六套四分套之一；後撥沉香三千二百裹，瑇瑁一千五百觔，乳香三千七百五十套。欲求沉、乳、瑇瑁裹、觔、套各價幾何？答曰：沉香每裹三百貫文；乳香毎套六十四貫文；瑇瑁每斤一百八十貫文。\end{unitzh}
\begin{uniten}Question: The Goods-Disposal Office (tuihuowu) paid out goods on three occasions, each equivalent to 1,470,000 guan wen (貫文) in cash. The first allocation was agarwood 3,500 guo (裹), tortoiseshell 2,200 jin (觔), frankincense 375 tao (套). The next allocation was agarwood 2,970 guo, tortoiseshell 2,130 jin, frankincense 3,056 1/4 tao. The last allocation was agarwood 3,200 guo, tortoiseshell 1,500 jin, frankincense 3,750 tao. We wish to find the price of each guo, jin and tao of agarwood, frankincense and tortoiseshell. What is each? Answer: Agarwood, 300 guan wen per guo; frankincense, 64 guan wen per tao; tortoiseshell, 180 guan wen per jin (斤).\end{uniten}

\begin{unitzh}\unitlabel{shushu-065.2}術曰：以方程求之，正負入之。列積及物數於下，布行數，各對本色，有分者通之，可約者約之，為定率。積、列數每以下項互徧乘之，毎視其積，以少减多，其下物數各随積正負之類：如同名相减，異名相加，正無人負之，負無人正之；其如同名相加，異名相减，正無人正之，負無人負之。使其下項物數得一數者為法，其積為寔，寔如法而一。所得，不計遍損或益諸積，各得法寔除之。餘倣此。\end{unitzh}
\begin{uniten}Method: Solve it by rectangular arrays (fangcheng), entering positive and negative. Lay out the totals (ji), and the numbers of goods below them, arranging the numbers in columns (hang), each facing its own kind. Where there are fractions, convert them (tong); where they can be reduced, reduce them (yue). These are the fixed rates (dinglü). Each time, cross-multiply (hu bian cheng) the totals and the laid-out numbers throughout by the bottom items. Each time, look at their totals and subtract the smaller from the larger. The numbers of goods below each follow the class of positive and negative of the totals. If same names (tongming) are subtracted from each other and different names (yiming) added to each other, a positive with no counterpart is made negative and a negative with no counterpart is made positive. If instead same names are added to each other and different names subtracted from each other, a positive with no counterpart is made positive and a negative with no counterpart is made negative. Make the numbers of goods in the bottom items come to a single number and take it as the divisor (fa); its total is the dividend (shi); divide the dividend by the divisor. With what is obtained, regardless of whether it diminishes or increases the various totals throughout, each obtains its divisor and dividend and is divided. The rest follow this.\end{uniten}

\begin{unitzh}\unitlabel{shushu-065.3}草曰：置凖錢一百四十七萬貫，為三次撥錢，為三行積數。次置先撥沉香三千五百裏，瑇瑁二千二百觔，乳香三百七十五觔，為右行物數。又列次撥沉香二千九百七十裹，瑇瑁二千一百三十觔，乳香三千五十六套四分套之一，為中行物。次列沉香三千二百裹，瑇瑁一千五百觔，乳香三千七百五十套，為左行之物。各以本色相對列之。\end{unitzh}
\begin{uniten}Working: Set out the equivalent cash, 1,470,000 guan, as the cash of the three allocations; these are the totals of the three columns. Next set out the first allocation, agarwood 3,500 guo [Translator's note: the text writes 裏 for 裹 here], tortoiseshell 2,200 jin, frankincense 375 jin, as the numbers of goods of the right column. Then lay out the next allocation, agarwood 2,970 guo, tortoiseshell 2,130 jin, frankincense 3,056 1/4 tao, as the goods of the middle column. Next lay out agarwood 3,200 guo, tortoiseshell 1,500 jin, frankincense 3,750 tao, as the goods of the left column. Lay out each one facing its own kind in the other columns.\end{uniten}

\begin{unitzh}\unitlabel{shushu-065.4}率圖 定率圖\end{unitzh}
\begin{uniten}Diagram of rates (lü tu). Diagram of fixed rates (dinglü tu).\end{uniten}

\begin{unitzh}\unitlabel{shushu-065.5}其中行乳香有四分套之一，便以母四通中行諸數，只内子一入乳香叚内：積得五百八十八萬貫，沉香得一萬一千八百八十裹，瑇瑁得八千五百二十觔，乳香得一萬二千二百二十五套。以右行求等，得二十五，俱約之：積得五萬八千八百貫，沉香得一百四十裹，瑇瑁得八十八觔，乳香得一十五套。以中行求等，得一十五，約之：積得三十九萬二千貫，沉香得七百九十二裹，瑇瑁得五百六十八觔，乳香得八百一十五套。以左行求等，得五，約之：積得二萬九千四百貫，沉香得六十四裹，瑇瑁得三十觔，乳香得七十五套。列為定率圖三行，副置求之。\end{unitzh}
\begin{uniten}In the middle column the frankincense has 1/4 tao. So use the denominator (mu) 4 to convert all the numbers of the middle column, putting the numerator (zi) 1 into the frankincense section (duan) only. The total becomes 5,880,000 guan, agarwood becomes 11,880 guo, tortoiseshell becomes 8,520 jin, frankincense becomes 12,225 tao. Seeking the common measure (deng) of the right column gives 25; reduce them all by it. The total becomes 58,800 guan, agarwood becomes 140 guo, tortoiseshell becomes 88 jin, frankincense becomes 15 tao. Seeking the common measure of the middle column gives 15; reducing by it, the total becomes 392,000 guan, agarwood becomes 792 guo, tortoiseshell becomes 568 jin, frankincense becomes 815 tao. Seeking the common measure of the left column gives 5 [Translator's note: the values that follow correspond to a common measure of 50]. Reducing by it, the total becomes 29,400 guan, agarwood becomes 64 guo, tortoiseshell becomes 30 jin, frankincense becomes 75 tao. Lay these out as the three columns of the diagram of fixed rates; set out a duplicate (fuzhi) and work on it.\end{uniten}

\begin{unitzh}\unitlabel{shushu-065.6}今先欲去定圖下位乳香套數：一十五與左下七十五互乘左右兩行，右積得四百四十一萬貫，沉香一萬五百裹，瑇瑁得六千六百觔，乳香得一千一百二十五套；左積得四十四萬一千貫，沉香得九百六十一裹，瑇瑁得四百五十觔，乳香得一千一百二十五套。驗左積少，右積多，當以左行直減右行。畢，仍置定圖左行數。\end{unitzh}
\begin{uniten}Now we first wish to eliminate the tao-numbers of frankincense in the bottom position of the fixed diagram. The 15 and the lower-left 75 cross-multiply the left and right columns. The right total becomes 4,410,000 guan, agarwood 10,500 guo, tortoiseshell becomes 6,600 jin, frankincense becomes 1,125 tao. The left total becomes 441,000 guan, agarwood becomes 961 guo [Translator's note: 64 × 15 is 960], tortoiseshell becomes 450 jin, frankincense becomes 1,125 tao. Examining: the left total is smaller and the right total larger, so one should directly subtract (zhijian) the left column from the right column. When done, again set out the numbers of the left column of the fixed diagram.\end{uniten}

\begin{unitzh}\unitlabel{shushu-065.7}維圖{\zhsmall 按：以上三圖名舊本未載，今依前題補之，以便檢閲。}\end{unitzh}
\begin{uniten}The wei (維) diagram. [Siku editors: The names of the three diagrams above are not recorded in the old edition; they are now supplied following the preceding problem, for convenience of consultation.]\end{uniten}

\begin{unitzh}\unitlabel{shushu-065.8}右積得三百九十六萬九千貫，沉香得九千五百四十裹，瑇瑁得六千一百五十觔。次驗中左兩行各有下位叚，又以左下七十五互乗中行，乃以中行下八百一十五互乗左行。畢，中積得二千九百四十萬貫，沉香得五萬九千四百裹，瑇瑁得四萬二千六百觔，乳香得六萬一千一百二十五套；左積得二千三百九十六萬一千貫，沉香得五萬二千一百六十觔，瑇瑁得二萬四千四百五十觔，乳香得六萬一千一百二十五套。驗左積少，中積多，以左行同名直減中行。畢，仍置定圖左行數。\end{unitzh}
\begin{uniten}The right total becomes 3,969,000 guan, agarwood becomes 9,540 guo, tortoiseshell becomes 6,150 jin. Next, examining: the middle and left columns each have a section in the bottom position. So again use the lower-left 75 to cross-multiply the middle column, and then use the 815 at the bottom of the middle column to cross-multiply the left column. When done, the middle total becomes 29,400,000 guan, agarwood becomes 59,400 guo, tortoiseshell becomes 42,600 jin, frankincense becomes 61,125 tao. The left total becomes 23,961,000 guan, agarwood becomes 52,160 jin, tortoiseshell becomes 24,450 jin, frankincense becomes 61,125 tao. Examining: the left total is smaller and the middle total larger. With same names, directly subtract the left column from the middle column. When done, again set out the numbers of the left column of the fixed diagram.\end{uniten}

\begin{unitzh}\unitlabel{shushu-065.9}卜圖 宫圖\end{unitzh}
\begin{uniten}The bu (卜) diagram. The gong (宫) diagram.\end{uniten}

\begin{unitzh}\unitlabel{shushu-065.10}中積得五百四十三萬九千貫，沉香得七千二百四十裹，瑇瑁得一萬八千一百五十觔。今驗右中兩行數多，又求約之：其右行求得三十，約之，右積得一十三萬二千三百貫，沉香得三百一十八裹，瑇瑁得二百五觔；中行求得一十，約之，中積得五十四萬三千九百貫，沉香得七百二十四裹，瑇瑁得一千八百一一十五觔。\end{unitzh}
\begin{uniten}The middle total becomes 5,439,000 guan, agarwood becomes 7,240 guo, tortoiseshell becomes 18,150 jin. Now, examining: the numbers of the right and middle columns are large, so again seek [the common measure] and reduce them. For the right column, seeking gives 30; reducing by it, the right total becomes 132,300 guan, agarwood becomes 318 guo, tortoiseshell becomes 205 jin. For the middle column, seeking gives 10; reducing by it, the middle total becomes 543,900 guan, agarwood becomes 724 guo, tortoiseshell becomes 1,815 jin [Translator's note: the text writes 一千八百一一十五, with 一 doubled].\end{uniten}

\begin{unitzh}\unitlabel{shushu-065.11}今又欲去中左行之瑇瑁，乃以中行瑇瑁一千八百一十五互乗右行：右積得二億四千一十二萬四千五百貫，沉香得五十七萬七千一百七十裹，瑇瑁得三十七萬二千七十五套。以卜右瑇瑁二百五觔互乗中行：中積一億一千一百四十九萬九千五百貫，沉香得一十四萬八千四百二十裹，瑇瑁得三十七萬二千七十五觔。今驗宫右積多，中積少，乃以中行直減右行。畢，仍置卜圖中行數。\end{unitzh}
\begin{uniten}Now we again wish to eliminate the tortoiseshell of the middle and left columns. So use the middle column's tortoiseshell, 1,815, to cross-multiply the right column. The right total becomes 240,124,500 guan, agarwood becomes 577,170 guo, tortoiseshell becomes 372,075 tao. Use the tortoiseshell 205 jin on the right of the bu [diagram] to cross-multiply the middle column. The middle total is 111,499,500 guan, agarwood becomes 148,420 guo, tortoiseshell becomes 372,075 jin. Now, examining the gong [diagram]: the right total is larger and the middle total smaller, so directly subtract the middle column from the right column. When done, again set out the numbers of the middle column of the bu diagram.\end{uniten}

\begin{unitzh}\unitlabel{shushu-065.12}干圖 支圖\end{unitzh}
\begin{uniten}The gan (干) diagram. The zhi (支) diagram.\end{uniten}

\begin{unitzh}\unitlabel{shushu-065.13}今驗干圖右行叚數，只有沉香四十二萬八千七百五十裹，以為法；以右上積一億二千八百六十二萬五千貫為寔，寔如法而一，得三百貫，為沉香一裹價。便以中行沉香七百二十四乗三百貫，得二十一萬七千二百貫，減中行五十四萬三千九百貫，餘三十二萬六千七百貫，為中積，便減去中行沉香叚之數。次以左上沉香六十四乗三百貫，得一萬九千二百貫，減左積二萬九千四百貫，餘一萬二百貫，為左積，便減左上沉香裹數去之。今驗支圖中行，其下只有瑇瑁一千八百一十五觔，以為法；中積三十二萬六千七百貫為寔，寔如法而一，得一百八十貫，為瑇瑁價。\end{unitzh}
\begin{uniten}Now, examining the sections of the right column of the gan diagram: there is only agarwood, 428,750 guo; take it as the divisor. Take the upper-right total, 128,625,000 guan, as the dividend; divide the dividend by the divisor, giving 300 guan as the price of 1 guo of agarwood. Then multiply 300 guan by the middle column's agarwood, 724, giving 217,200 guan. Subtract it from the middle column's 543,900 guan, leaving 326,700 guan as the middle total; then subtract away the number of the agarwood section of the middle column. Next multiply 300 guan by the upper-left agarwood, 64, giving 19,200 guan. Subtract it from the left total, 29,400 guan, leaving 10,200 guan as the left total; then subtract the guo-number of the upper-left agarwood, removing it. Now, examining the middle column of the zhi diagram: below there is only tortoiseshell, 1,815 jin; take it as the divisor. The middle total, 326,700 guan, is the dividend. Divide the dividend by the divisor, giving 180 guan as the price of tortoiseshell.\end{uniten}

\begin{unitzh}\unitlabel{shushu-065.14}閏圖{\zhsmall 按：此圖舊本亦未載名，今補。}\end{unitzh}
\begin{uniten}The run (閏) diagram. [Siku editors: The name of this diagram is likewise not recorded in the old edition; it is now supplied.]\end{uniten}

\begin{unitzh}\unitlabel{shushu-065.15}今驗閏圖左行有瑇瑁三十觔，以乘價一百八十貫，得五千四百貫，减左積一萬二百貫，餘四千八百貫，為左積。其下只有乳香七十五套，以為法；以積四千八百貫為寔，寔如法而一，得六十四貫，為乳香套價。此題並係俱正，補草。\end{unitzh}
\begin{uniten}Now, examining the left column of the run diagram: there is tortoiseshell, 30 jin. Multiply the price, 180 guan, by it, giving 5,400 guan. Subtract it from the left total, 10,200 guan, leaving 4,800 guan as the left total. Below there is only frankincense, 75 tao; take it as the divisor. Take the total, 4,800 guan, as the dividend. Divide the dividend by the divisor, giving 64 guan as the price per tao of frankincense. In this problem everything is positive throughout; working supplied.\end{uniten}

\begin{unitzh}\unitlabel{shushu-065.16}數學九章巻九下\end{unitzh}
\begin{uniten}Mathematical Treatise in Nine Sections (Shuxue jiuzhang), juan 9, part 2.\end{uniten}
\clearpage\section*{Sources}\addcontentsline{toc}{section}{Sources}
The Chinese text is Siku quanshu (四庫全書) edition, titled 數學九章 there, as transcribed on Chinese Wikisource. The text of the work is in the public domain; the Wikisource transcription is credited here with the pages and revisions used. The punctuation, the division into units and the English translation are AI-produced and dedicated to the public domain (CC0 1.0).

\begin{itemize}
\item \href{https://zh.wikisource.org/wiki/\%E6\%95\%B8\%E5\%AD\%B8\%E4\%B9\%9D\%E7\%AB\%A0_\%28\%E5\%9B\%9B\%E5\%BA\%AB\%E5\%85\%A8\%E6\%9B\%B8\%E6\%9C\%AC\%29}{數書九章}, revision 619703
\item \href{https://zh.wikisource.org/wiki/\%E6\%95\%B8\%E5\%AD\%B8\%E4\%B9\%9D\%E7\%AB\%A0_\%28\%E5\%9B\%9B\%E5\%BA\%AB\%E5\%85\%A8\%E6\%9B\%B8\%E6\%9C\%AC\%29/\%E5\%8D\%B71\%E4\%B8\%8A}{卷1上}, revision 783952
\item \href{https://zh.wikisource.org/wiki/\%E6\%95\%B8\%E5\%AD\%B8\%E4\%B9\%9D\%E7\%AB\%A0_\%28\%E5\%9B\%9B\%E5\%BA\%AB\%E5\%85\%A8\%E6\%9B\%B8\%E6\%9C\%AC\%29/\%E5\%8D\%B71\%E4\%B8\%8B}{卷1下}, revision 619667
\item \href{https://zh.wikisource.org/wiki/\%E6\%95\%B8\%E5\%AD\%B8\%E4\%B9\%9D\%E7\%AB\%A0_\%28\%E5\%9B\%9B\%E5\%BA\%AB\%E5\%85\%A8\%E6\%9B\%B8\%E6\%9C\%AC\%29/\%E5\%8D\%B72\%E4\%B8\%8A}{卷2上}, revision 619669
\item \href{https://zh.wikisource.org/wiki/\%E6\%95\%B8\%E5\%AD\%B8\%E4\%B9\%9D\%E7\%AB\%A0_\%28\%E5\%9B\%9B\%E5\%BA\%AB\%E5\%85\%A8\%E6\%9B\%B8\%E6\%9C\%AC\%29/\%E5\%8D\%B72\%E4\%B8\%8B}{卷2下}, revision 619671
\item \href{https://zh.wikisource.org/wiki/\%E6\%95\%B8\%E5\%AD\%B8\%E4\%B9\%9D\%E7\%AB\%A0_\%28\%E5\%9B\%9B\%E5\%BA\%AB\%E5\%85\%A8\%E6\%9B\%B8\%E6\%9C\%AC\%29/\%E5\%8D\%B73\%E4\%B8\%8A}{卷3上}, revision 1497476
\item \href{https://zh.wikisource.org/wiki/\%E6\%95\%B8\%E5\%AD\%B8\%E4\%B9\%9D\%E7\%AB\%A0_\%28\%E5\%9B\%9B\%E5\%BA\%AB\%E5\%85\%A8\%E6\%9B\%B8\%E6\%9C\%AC\%29/\%E5\%8D\%B73\%E4\%B8\%8B}{卷3下}, revision 619675
\item \href{https://zh.wikisource.org/wiki/\%E6\%95\%B8\%E5\%AD\%B8\%E4\%B9\%9D\%E7\%AB\%A0_\%28\%E5\%9B\%9B\%E5\%BA\%AB\%E5\%85\%A8\%E6\%9B\%B8\%E6\%9C\%AC\%29/\%E5\%8D\%B74\%E4\%B8\%8A}{卷4上}, revision 783958
\item \href{https://zh.wikisource.org/wiki/\%E6\%95\%B8\%E5\%AD\%B8\%E4\%B9\%9D\%E7\%AB\%A0_\%28\%E5\%9B\%9B\%E5\%BA\%AB\%E5\%85\%A8\%E6\%9B\%B8\%E6\%9C\%AC\%29/\%E5\%8D\%B74\%E4\%B8\%8B}{卷4下}, revision 619679
\item \href{https://zh.wikisource.org/wiki/\%E6\%95\%B8\%E5\%AD\%B8\%E4\%B9\%9D\%E7\%AB\%A0_\%28\%E5\%9B\%9B\%E5\%BA\%AB\%E5\%85\%A8\%E6\%9B\%B8\%E6\%9C\%AC\%29/\%E5\%8D\%B75\%E4\%B8\%8A}{卷5上}, revision 619681
\item \href{https://zh.wikisource.org/wiki/\%E6\%95\%B8\%E5\%AD\%B8\%E4\%B9\%9D\%E7\%AB\%A0_\%28\%E5\%9B\%9B\%E5\%BA\%AB\%E5\%85\%A8\%E6\%9B\%B8\%E6\%9C\%AC\%29/\%E5\%8D\%B75\%E4\%B8\%8B}{卷5下}, revision 619683
\item \href{https://zh.wikisource.org/wiki/\%E6\%95\%B8\%E5\%AD\%B8\%E4\%B9\%9D\%E7\%AB\%A0_\%28\%E5\%9B\%9B\%E5\%BA\%AB\%E5\%85\%A8\%E6\%9B\%B8\%E6\%9C\%AC\%29/\%E5\%8D\%B76\%E4\%B8\%8A}{卷6上}, revision 619685
\item \href{https://zh.wikisource.org/wiki/\%E6\%95\%B8\%E5\%AD\%B8\%E4\%B9\%9D\%E7\%AB\%A0_\%28\%E5\%9B\%9B\%E5\%BA\%AB\%E5\%85\%A8\%E6\%9B\%B8\%E6\%9C\%AC\%29/\%E5\%8D\%B76\%E4\%B8\%8B}{卷6下}, revision 619687
\item \href{https://zh.wikisource.org/wiki/\%E6\%95\%B8\%E5\%AD\%B8\%E4\%B9\%9D\%E7\%AB\%A0_\%28\%E5\%9B\%9B\%E5\%BA\%AB\%E5\%85\%A8\%E6\%9B\%B8\%E6\%9C\%AC\%29/\%E5\%8D\%B77\%E4\%B8\%8A}{卷7上}, revision 619689
\item \href{https://zh.wikisource.org/wiki/\%E6\%95\%B8\%E5\%AD\%B8\%E4\%B9\%9D\%E7\%AB\%A0_\%28\%E5\%9B\%9B\%E5\%BA\%AB\%E5\%85\%A8\%E6\%9B\%B8\%E6\%9C\%AC\%29/\%E5\%8D\%B77\%E4\%B8\%8B}{卷7下}, revision 619691
\item \href{https://zh.wikisource.org/wiki/\%E6\%95\%B8\%E5\%AD\%B8\%E4\%B9\%9D\%E7\%AB\%A0_\%28\%E5\%9B\%9B\%E5\%BA\%AB\%E5\%85\%A8\%E6\%9B\%B8\%E6\%9C\%AC\%29/\%E5\%8D\%B78\%E4\%B8\%8A}{卷8上}, revision 619693
\item \href{https://zh.wikisource.org/wiki/\%E6\%95\%B8\%E5\%AD\%B8\%E4\%B9\%9D\%E7\%AB\%A0_\%28\%E5\%9B\%9B\%E5\%BA\%AB\%E5\%85\%A8\%E6\%9B\%B8\%E6\%9C\%AC\%29/\%E5\%8D\%B78\%E4\%B8\%8B}{卷8下}, revision 619695
\item \href{https://zh.wikisource.org/wiki/\%E6\%95\%B8\%E5\%AD\%B8\%E4\%B9\%9D\%E7\%AB\%A0_\%28\%E5\%9B\%9B\%E5\%BA\%AB\%E5\%85\%A8\%E6\%9B\%B8\%E6\%9C\%AC\%29/\%E5\%8D\%B79\%E4\%B8\%8A}{卷9上}, revision 619697
\item \href{https://zh.wikisource.org/wiki/\%E6\%95\%B8\%E5\%AD\%B8\%E4\%B9\%9D\%E7\%AB\%A0_\%28\%E5\%9B\%9B\%E5\%BA\%AB\%E5\%85\%A8\%E6\%9B\%B8\%E6\%9C\%AC\%29/\%E5\%8D\%B79\%E4\%B8\%8B}{卷9下}, revision 783961
\end{itemize}
\end{document}
